\documentclass[10pt]{amsart}
\usepackage{amsmath,amssymb,amsthm,mathtools,mathrsfs}
\usepackage[a4paper,hmargin=0.76in,vmargin=0.60in]{geometry}
\usepackage{booktabs}
\usepackage{array}
\usepackage{longtable}
\usepackage{tabularx}
\usepackage[hypertexnames=false]{hyperref}
\hypersetup{hidelinks}
\usepackage[final,expansion=false]{microtype}
\usepackage{url}
\emergencystretch=2.5em

\newtheorem{theorem}{Theorem}[section]
\newtheorem{proposition}[theorem]{Proposition}
\newtheorem{lemma}[theorem]{Lemma}
\newtheorem{corollary}[theorem]{Corollary}
\theoremstyle{definition}
\newtheorem{definition}[theorem]{Definition}
\newtheorem{problem}[theorem]{Problem}
\theoremstyle{remark}
\newtheorem{remark}[theorem]{Remark}

\newcommand{\C}{\mathbb C}
\newcommand{\R}{\mathbb R}
\newcommand{\N}{\mathbb N}
\newcommand{\Z}{\mathbb Z}
\newcommand{\Q}{\mathbb Q}
\newcommand{\Rel}{\operatorname{Re}}
\newcommand{\Iml}{\operatorname{Im}}
\newcommand{\supp}{\operatorname{supp}}
\newcommand{\TV}{\operatorname{TV}}
\newcommand{\Kp}{\mathcal K^{\mathrm{prim}}_{2T}}
\newcommand{\Kw}{\mathcal K_{2T}}
\newcommand{\PT}{\mathcal P_T}
\newcommand{\IT}{I_T}
\newcommand{\lamP}{\lambda^{\mathrm{prim}}_T}
\newcommand{\lamW}{\lambda^{\mathrm W}_T}
\newcommand{\lamA}{\lambda_{\mathrm A}(T)}
\newcommand{\sT}{\sigma_T}
\newcommand{\gh}[1]{\widehat g_{#1}}
\newcommand{\nuh}{\widehat\nu}
\newcommand{\wh}{\widehat w}
\newcommand{\Fh}{\widehat F}
\newcommand{\cF}{c_F}
\newcommand{\Lev}{\mathfrak E}
\newcommand{\Wcl}{\mathcal W_T}
\newcommand{\NT}{\mathcal N_T}
\newcommand{\Aconst}{\mathsf A_T}
\newcommand{\Ainf}{\mathcal A_\infty}
\newcommand{\PiT}{\Pi_T}
\newcommand{\mT}{\mathfrak m_T}
\newcommand{\laminf}{\lambda^{\infty}_{\min}}
\newcommand{\laminfP}{\lambda^{\infty,\mathrm P}_{\min}}
\newcolumntype{Y}{>{\raggedright\arraybackslash}X}

% macros of the series: every macro has in all nine papers the meaning it had in the single
% paper; the definitions needed before their place in the source files are collected here
\providecommand{\anCP}{C_P}
\providecommand{\anPa}{\Phi_a}
\providecommand{\anPb}{\Phi_b}
\providecommand{\anPl}{\Phi}
\providecommand{\ancl}{\mathrm{cl}}
\providecommand{\aneu}{\mathrm{eu}_d}
\providecommand{\anrd}{\mathrm{rd}}
\providecommand{\auAz}{\mathcal A_0}
\providecommand{\auDX}{D_X}
\providecommand{\auF}{\mathbb F}
\providecommand{\auG}{\overline G}
\providecommand{\auGL}{\mathrm{GL}}
\providecommand{\auP}{\mathcal P}
\providecommand{\auPl}{\Phi}
\providecommand{\auS}{\mathcal S}
\providecommand{\auSram}{\mathcal S^{\mathrm{ram}}}
\providecommand{\auT}{\mathcal T}
\providecommand{\auU}{\mathrm U}
\providecommand{\aubar}[1]{\overline{#1}}
\providecommand{\audT}{d_{\auT}}
\providecommand{\aupp}{\mathfrak p}
\providecommand{\pbDid}{\mathcal D_{\mathrm{id}}}
\providecommand{\pbDmain}{\mathcal D_{\mathrm{main}}}
\providecommand{\pbR}{\overline R_0}
\providecommand{\pbrpsi}{r_\psi}
\providecommand{\pcoB}{\mathsf B}
\providecommand{\pffK}{K_T}
\providecommand{\pffKM}{K_T(M)}
\providecommand{\pffZ}{\mathsf Z}
\providecommand{\pffm}{\varkappa}
\providecommand{\pfrM}{\mathsf M}
\providecommand{\pfrU}{\mathsf U}
\providecommand{\pfrX}{\Xi}
\providecommand{\pfrell}{\ell}
\providecommand{\pfrip}[2]{\langle#1,#2\rangle}
\providecommand{\pfrm}{\mathsf m}
\providecommand{\psHS}{\mathcal H}
% statements quoted from the other papers of the series (generated appendix)
\newenvironment{quotedthm}[2]{\par\medskip\noindent\begingroup
  \textbf{#1}#2.\ \itshape\ignorespaces}{\endgroup\par\medskip}
\newenvironment{quoteddef}[2]{\par\medskip\noindent\begingroup
  \textbf{#1}#2.\ \ignorespaces}{\endgroup\par\medskip}
\begin{document}
\addtolength{\abovedisplayskip}{-3.4pt}
\addtolength{\belowdisplayskip}{-3.4pt}
\addtolength{\belowdisplayshortskip}{-2pt}
\title[The Spanning Theorem: the second moment]{The Spanning Theorem: the second moment of a Poincar\'e map}
\author{David Alejandro Trejo Pizzo}
\date{}
\begin{abstract}
The Spanning Theorem is the second of three theorems on integral algebraic modular forms from
which the main paper of this series deduces the Riemann hypothesis for Dirichlet $L$-functions.
It says that a few translates, under a finite group acting through the level, of the image of a
Poincar\'e map from algebraic modular forms on a definite unitary group to cusp forms of discrete
series type on a larger unitary group span almost all of the target. We prove it through an
estimate for the second moment of the map over the group, in which the main term has to be
obtained with its exact constant. The estimate comes from a relative trace formula for Poincar\'e
series of Bessel type: the main term is extracted from the open cell by Siegel's formula and
repeated Poisson summation, and the remaining terms are bounded through the decay of the Bessel
functions of discrete series with large parameters and a count of pairs of terms. We also prove
that Hecke eigensystems occur with the same multiplicities on unitary groups that are isometric
at all finite places, and Harish-Chandra's formula for the elliptic orbital integrals of the
coefficients used.
\end{abstract}
\maketitle
\setcounter{tocdepth}{1}\tableofcontents

\section{Introduction}\label{B1:intro}

The main paper~\cite{P46} of this series deduces the Riemann hypothesis for Dirichlet
$L$-functions from three theorems on integral algebraic modular forms on a chain of theta lifts
between unitary groups over a CM field $E$: the Short Forms Theorem, the Spanning Theorem and the
Covolume Theorem. The first gives a supply of integral forms that are short at every conjugate,
namely suitable multiples of the images of integral forms under a Poincar\'e map $\auP$; the
third bounds from below the covolume of every family of independent forms of almost full rank. To compare the two one needs to know that the short forms span almost all of the
space, and this is the Spanning Theorem. In~\cite{P46} it is deduced from two propositions, an
estimate for the second moment of $\auP$ ({\cite[Proposition~14.32]{P46}}) and an equality of
multiplicities of Hecke eigensystems ({\cite[Proposition~14.36]{P46}}), which is also used for the
Covolume Theorem. We prove both here, except for three local statements whose proofs
are in~\cite{P52}. The Short Forms Theorem is proved in \cite{P49} and~\cite{P50}, the Covolume
Theorem in \cite{P53}, \cite{P54}, \cite{P55} and~\cite{P56}.

We recall the objects from {\cite[\S\S14.1--14.3]{P46}}. Let $E$ be the CM field
of~\cite{P46}, with maximal totally real subfield $E^+$ and set $\auPl$ of $f$ infinite places,
and let $Y_0=X_0\perp\mathbb H^{\varrho}\perp\langle-N\rangle$ be the hermitian space over $E$ which
is the sum of the standard positive definite space $X_0$ of rank $k$, of $\varrho$ hyperbolic
planes and of a line generated by a vector $z$ of norm $-N$; its rank is $r=k+2\varrho+1$ and its
signature at every infinite place is $(k+\varrho,\varrho+1)$. Let $U$ be the unipotent radical of
the parabolic subgroup of $\auU(Y_0)$ that stabilises a full flag of isotropic subspaces of
$\mathbb H^{\varrho}$, with a generic character $\psi_U$ that is invariant under $\auU(X_0)$, and
let $R_0=U\rtimes\auU(X_0)$ be the Bessel subgroup. The source $\auS$ is a space of algebraic
modular forms on the definite group $\auU(X_0)$, of dimension $\auDX$. The space $\auAz$ consists
of the bounded automorphic functions on $\auU(Y_0)$, of the principal level $J_M$ at the primes
of $E$ dividing an element $\pi_M$ of norm $M=p_0p_1$, that are reproduced by convolution with
the coefficient of the lowest $K$-type of a discrete series at every infinite place; the
parameters of these discrete series are large and far from the walls. The finite group
$\auG=\prod_w\auGL_r(\auF_w)$, the product over the $2f$ primes $w$ dividing $\pi_M$, acts
unitarily on $\auAz$ by right translation; we write $\varrho(s)$ for the action of $s\in\auG$.
The Poincar\'e map $\auP\colon\auS\to\auAz$ is the adjoint of the Bessel period over $R_0$, read
at a point $g_0$ of the real group and through a linear form $\mathsf t$ on the values of the
source forms. Finally $e_0$ is a central idempotent of the group algebra of $\auG$, the sum of
the idempotents of certain characters that are cuspidal at some prime above each of $p_0$ and
$p_1$, and the target is $\auT=\varrho(e_0)\auAz$, of dimension $\audT$.

The Spanning Theorem ({\cite[Theorem~14.14]{P46}}) states that $\audT>0$ and that there is a set
$\mathsf t\subseteq\auG$ such that, with $\auSram\subseteq\auS$ the subspace of the forms that
are ramified at the prime $N$ of the level,
\[
\#\mathsf t\cdot\dim\auSram\le5\,\audT\log(1+2b^2)\qquad\text{and}\qquad
\dim\sum_{s\in\mathsf t}\varrho(s)\varrho(e_0)\auP(\auSram)\ \ge\ \Bigl(1-\frac4{b^2}\Bigr)\audT ,
\]
where $b$ is one of the scales of~\cite{P46}. The estimate from which it follows is this
(Theorem~\ref{pb:thmB}, which is {\cite[Proposition~14.32]{P46}}). Let $\deg$ be the product over
$\tau\in\auPl$ of the formal degrees of the discrete series, for the Haar measures that
correspond to the measure of mass one on the compact form, and let $\mathrm{mass}_r$ be the mass
of the genus of the standard hermitian lattice of rank $r$ over the integers of $E$. Then
\emph{there are a point $g_0$, a linear form $\mathsf t$, an inner product on $\auS$ and a scalar
$c_P>0$ such that $P=c_P\auP$ satisfies, with Hilbert--Schmidt norms,}
\begin{gather}
0<\audT\le\#\auG\cdot\deg\cdot\mathrm{mass}_r,\label{B1:eq-i}\\
\|\varrho(e_0)\,P\,\mathrm{pr}_{\mathsf U}\|^2\ge\dim\mathsf U-5b^{-4}\auDX\qquad\text{\emph{for
every subspace} }\mathsf U\subseteq\auS,\label{B1:eq-ii}\\
\sum_{s\in\auG}\|P^*\varrho(s)P\|^2\le\bigl(1+9b^{-4}\bigr)\,\frac{\auDX^2}{\deg\cdot\mathrm{mass}_r}.
\label{B1:eq-iii}
\end{gather}
By \eqref{B1:eq-ii} the map $\varrho(e_0)P$ is almost an isometry. By \eqref{B1:eq-i} the right
side of \eqref{B1:eq-iii} is at most $(1+9b^{-4})\,\#\auG\,\auDX^2/\audT$. For an isometry $P$
of a space of dimension $\auDX$ into a unitary representation of $\auG$ of dimension $\audT$ the
left side of \eqref{B1:eq-iii} is at least $\#\auG\,\auDX^2/\audT$, since it equals $(\#\auG)^{-1}$
times the square of the Hilbert--Schmidt norm of $\sum_s\varrho(s)PP^*\varrho(s)^{-1}$, an operator
of trace $\#\auG\,\auDX$; equality holds when this operator is a scalar, that is, when the
translates of the image are evenly spread. So \eqref{B1:eq-i}--\eqref{B1:eq-iii} say that the
translates of the image of $P$ are almost evenly spread over $\auT$. An elementary argument
({\cite[Lemmas~14.30 and~14.31]{P46}}) converts this into the Spanning Theorem: by the
Cauchy--Schwarz inequality, for every subspace $F$ of $\auT$, of dimension $\alpha\,\audT$, there
is a translate of $\varrho(e_0)P$ under which the preimage of $F$ has dimension at most about
$\tfrac12(1+\alpha)\auDX$; adding the image of this translate to $F$ lowers the codimension by a
fixed proportion of $(1-\alpha)\auDX$, and repeating this fills $\auT$ up to the proportion
$4/b^2$. The argument has no room for a constant: with a factor $C>1$ in place of $1+9b^{-4}$ in
\eqref{B1:eq-iii} it would say nothing about subspaces $F$ of dimension at least $\audT/C$. This is why the main term of
\eqref{B1:eq-iii} has to be computed exactly, with the formal degree, the mass and all
normalisations of measures, uniformly in the rank $r$, which in~\cite{P46} grows with the other
parameters. Estimates of second moments by relative trace formulas are classical; we do not know
one of this uniformity in the literature.

The second result is Proposition~\ref{pm:mult} (the precise form of
{\cite[Proposition~14.36]{P46}}). Let $W_j$, $1\le j\le m$, be the hermitian spaces of rank $r$
of~\cite{P46}, which are isometric to $Y_0$ at all finite places and have other signatures at the
infinite places; $W_m$ is positive definite. Then every Hecke operator $T$ at the places of $E^+$ that
split in $E$ and are prime to $M$ has the same trace on the space of forms on $\auU(W_j)$ defined
like $\auT$, with the discrete series of the same infinitesimal characters, as on $\auT$; the same
holds for the whole spherical Hecke algebra away from $M$ (Remark~\ref{pm:rem-full}). In
particular the two spaces have the same dimension; for $j=m$ the group is compact at the infinite
places, the dimension is bounded by means of the mass formula, and this gives \eqref{B1:eq-i}.

\subsection*{The four inequalities}
The proof of \eqref{B1:eq-ii} and \eqref{B1:eq-iii} is organised in three layers, described at
the beginning of Section~\ref{pb:sec}. The first is finite group theory. The group $R_0$
determines a subgroup $\pbR$ of $\auG$, the image of the integral points of $R_0$, and a unitary
representation $\pbrpsi$ of $\pbR$ on $\auS$, in which the unipotent part acts through $\psi_U$
and the part coming from $\auU(X_0)$ by translation; the Poincar\'e map intertwines $\pbrpsi$ with
$\varrho$ (Lemma~\ref{pb:equiv}). Put $P_s=P^*\varrho(s)P$, and extend $\pbrpsi$ by zero to
$\auG$. Theorem~\ref{pb:est} asserts four inequalities. The first says that $P_s$ is close to
$\pbrpsi(s)$ on the set of the $s$ whose components above one of the primes $p_0$, $p_1$ are
those of an element of $\pbR$: the sum of $\|P_s-\pbrpsi(s)\|^2$ over this set is at most
$b^{-12}\,\pbDid$, where $\pbDid=\#\pbR\cdot\auDX$ is what $\pbrpsi$ itself contributes to the
second moment. The second and the third bound the sum of $\|P_s\|^2$ over the elements of $\auG$
that fix prescribed vectors of the span of the frame $e_1,\dots,e_\varrho,z$ modulo a set
$\mathsf A$ of primes above $p_0$: by $b^{-6kf}\,\pbDmain$ if a vector is fixed at each of the $f$
primes, and by
\[
\bigl(1+4b^{-4}\bigr)\,b^{-6k\#\mathsf A}\,\pbDmain,\qquad\pbDmain=\frac{\auDX^2}{\deg\cdot
\mathrm{mass}_r},
\]
for the sum over the \emph{open} elements, those for which, at some prime above $p_0$, the
leading minors of the matrix of products of the frame with its image do not vanish. The fourth
says that the character of $\auG$ induced from $\psi_U$ is almost free of multiplicities. From
these, \eqref{B1:eq-iii} follows by an induction on $\mathsf A$, in which the elements that are
not open are shown to be few by a lemma on functions of positive type (\S\ref{pb:ss-iii}), and
\eqref{B1:eq-ii} follows from the first and the fourth inequality and from the fact that the
characters excluded by $e_0$ have small Plancherel measure, by Green's description of the
characters of $\auGL_r(\auF_q)$~\cite{Green1955} (\S\ref{pb:ss-ii}).

\subsection*{The relative trace formula}
The second layer makes $P_s$ explicit. We show (Theorem~\ref{pbe:bessel}) that, for a suitable
point $g_0=(x_\tau)_\tau$ and form $\mathsf t=\bigotimes_\tau v_\tau$, the map $\auP$ is a
constant times the Poincar\'e series of a Bessel function $W=\bigotimes_\tau W_\tau$ on the real
group, where $W_\tau$ is a positive multiple of
\[
y\longmapsto\int_{R_\tau}c_\tau(x_\tau^{-1}s\,y)\,\pi_R(s)^{-1}v_\tau\,ds ,
\]
$c_\tau$ is the coefficient of the lowest $K$-type of the discrete series at $\tau$, $R_\tau$ the
group of real points of $R_0$ at $\tau$ and $\pi_R$ its representation given by the character and
by the weight of the source forms. Unfolding gives a formula of Petersson type
(Proposition~\ref{pb:unfold}): the matrix coefficients of $P_s$ are sums, over the double cosets
of the rational points $\gamma$ of $\auU(Y_0)$ that are integral and reduce to $s$, modulo the
integral points of $U$ on both sides, of the orbital integrals $\int W(y)W(\gamma^{-1}y)^*\,dy$.
The terms with $\gamma\in R_0$ give exactly $\pbrpsi(s)$ (Proposition~\ref{pb:besselterms}). The
two constants $\pbDid$ and $\pbDmain$ are then compared: $b^{11kf}\,\pbDid\le\pbDmain\le
(M/b^6)^{kf}\,\pbDid$ (Lemma~\ref{pb:sizes}). This is where the level $N$ is tuned: the ratio
contains $N^{kf}$, and the upper bound uses the upper end of the window in which~\cite{P46}
places $\log N$, whose lower end is needed for the Short Forms Theorem.

The existence of $W$ with the properties needed rests on four properties of $c_\tau$
(Theorem~\ref{pbe:ds}): the reproducing property, with the normalisation of the Haar measure by
the formal degree; the decay $|c_\tau(g)|\le(2/\|g\|)^{j+9r^2}$ for all $j\le r^7$; a comparison of
the formal degree with the volume of balls; and a lower bound for one Bessel period of $c_\tau$,
which shows that $W_\tau$ is not zero and controls its normalisation. The first three are proved
in Section~\ref{pds:sec}, the decay from the explicit integral formula for $c_\tau$ proved
in~\cite{P49} ({\cite[Theorem~3.17]{P49}}; here the gap $2r^8$ between the entries of the parameter
attached to coordinates of different signs is used). The lower bound for the period is proved in
Section~\ref{pbp:sec} by Laplace's method: the point $x_\tau$ is chosen so that on
$x_\tau^{-1}R_\tau x_\tau$ the phase of the character cancels the phases of the two coefficients
in the integrand to first order; near the identity the integrand is then close to a positive
number, and away from it the integrand is either small or equal to a value taken near the
identity. With the decay we cut $W$ off at a fixed distance from $R_\tau$, at the cost of a
negligible error in $L^2$ (Proposition~\ref{pbe:cut}).

\subsection*{The terms outside the Bessel subgroup}
The third layer bounds the sums of squares of the amplitudes. Squaring the Petersson formula
produces pairs $(\gamma,\gamma')$ of terms over the same $s$. To each $\gamma$ we attach its
\emph{link}, a subspace of dimension $k+1$ of two copies of the space $X_0\perp Ez$, which records
how $\gamma$ acts on the vectors orthogonal to the isotropic space $A$ of the frame and to
$\gamma A$; it is the graph of an isometry (Lemma~\ref{pbe:link}). Part (X) of
Theorem~\ref{pb:amp} says that the pairs with different links, and the terms $\gamma$ that map a
non-zero vector of the span of the frame into that span, contribute little, so that the sum of
$\|P_s-\pbrpsi(s)\|^2$ over a set of elements $s$ is governed by its \emph{link moment}, the sum
over $s$ and over the links of the squares of the amplitudes of the terms with a given link. Part
(L1) bounds the link moment of a set of elements $s$ each of which maps, modulo primes of total
norm at least $\mathsf z$, a non-zero vector of the span of the frame into that span, by
$(b^{2f}/\mathsf z)^k\,\pbDmain$ in the normalisation of the formula. Both parts are deduced in
\S\ref{pbe:ss-XL1} from three estimates, stated in exact form in Theorem~\ref{ppi:pillars}. The
\emph{size} of the contribution of one orbit of terms is bounded by Schur's test, with a weight
computed at each real place, and, on the open cell, by a bound for the pairing of two Bessel
functions (\S\S\ref{pbs:sec} and~\ref{pbo:sec}). The \emph{gain} is an inequality between the
gauges that govern these sizes and the discriminant of the lattice common to two terms, which improves by a factor $(M/2)^{f}$ for
each unit of codimension of the intersection of the two links (\S\ref{pbg:sec}). The \emph{count} bounds the number of pairs with a given common lattice
and sums over the lattices, by the geometry of numbers (\S\ref{pbc:sec}). A structural
statement on the elements with vanishing discriminant, formulated with a \emph{cell rank}, is
proved in Section~\ref{ppi:sec}.

What remains is the link moment of the open elements that fix prescribed vectors, and this is
the main term (Theorem~\ref{pbm:L2}, part (L2) of Theorem~\ref{pb:amp}). Over an open $s$ every $\gamma$
lies in the big cell and has a unique decomposition $\gamma=u\,w_p\,\sigma\,u'$ with $u,u'\in U$,
a Weyl element $w_p$ given by a \emph{head} $p=(p_0,\dots,p_{\varrho-1})$ of non-zero elements of
$E$, and $\sigma$ in the unitary group of $X_0\perp Ez$ (Proposition~\ref{pbm:cell}); the link is
the graph of $\sigma^{-1}$, and the sum over $u,u'$ of the contributions is a Kloosterman sum times
an orbital integral (Proposition~\ref{pbm:orbit}). Two operations then lead to the constant. The
elements $\sigma$ with a given last entry $t=h(\sigma z,z)$ correspond, through the vector
$\sigma z$, to the rational points of a sphere in $X_0$, and the sum of a non-negative function
over the rational points of the sphere and over the classes of the genus is bounded by the
constant of Siegel's formula times an adelic integral; this is \emph{Siegel's step}. The result
is a sum, over $t$ and over pairs of heads $(p,q)$, of a product of local hermitian forms in
$(p,q)$, and the sums over the entries of the heads are carried out one after the other by a
one-sided form of Poisson summation (Lemma~\ref{pbm:poisson}), which needs only positivity and
lower semicontinuity: at each step the zero frequency replaces the form by its integral over one
entry, and the other frequencies and the pairs with unequal entries must be shown to be small.
We axiomatise what is needed of the local forms in the notion of a \emph{sieve system}
(Definition~\ref{pbm:system}), with a \emph{mass}, bounding the form at the bottom, and a
\emph{decay} in the frequency, and prove the resulting inequality (Theorem~\ref{pbm:sieve}). The
existence of sieve systems with the right masses at the real places and at the finite places, and
Siegel's step, are the three statements of Theorem~\ref{pbm:inputs}; they are proved
in~\cite{P52}. The masses are $(1+b^{-10})$ times the inverse of the formal degree at a real
place, and $(1+b^{-6})\,b^{-6k\#\mathsf A}$ at the finite places, and the product of the constant
of Siegel's step and of the constants of the $\varrho+1$ Poisson summations, divided by the
normalisations of the orbital integrals, is bounded in Lemma~\ref{pbm:const}: it is at most $1+b^{-9}$ times
$\auDX^2/\mathrm{mass}_r$ times the square of the covolume $\nu$ of the integral points of $U$ in
its real points. This bound uses the exact form of the mass formula
of Gan, Hanke and Yu~\cite{GHY}, and the exact value of $\nu$, which is computed in
\S\ref{pbv:sec} (Theorem~\ref{pbv:covol}). The fourth inequality, on the induced character,
is independent of all this; it is proved in \S\ref{pbm:ss-E4} by two recursions on the rank
(Theorem~\ref{pbm:E4}).

\subsection*{Equality of multiplicities}
Proposition~\ref{pm:mult} is a comparison of trace formulas for two unitary groups with the same
finite adelic points, in a case in which it can be carried out by hand; the general comparison is
due to Labesse~\cite{Labesse}. At the infinite places we use as test function not a
pseudo-coefficient but the kernel $d_\pi\,\aubar{\langle\pi(\cdot)v,v\rangle}$ itself, which is
integrable because the parameter is far from the walls and which projects exactly onto the forms
of type $\pi$. The trace of a Hecke operator is then the integral over the diagonal of a kernel
whose partial sums are bounded, so that no truncation is needed (Propositions~\ref{pm:bound}
and~\ref{pm:trace}), and it is a sum over the characteristic polynomials of the rational points.
The idempotent $e_0$ provides a test function with a supercuspidal component at a split place,
for which the terms of the polynomials with two coprime factors vanish
(Proposition~\ref{pm:red}); the proper powers are removed by a condition at an auxiliary place, as
in the simple trace formula of Flicker and Kazhdan~\cite{FlickerKazhdan}, and the condition is
lifted afterwards (\S\ref{pm:ss-aux}). For an irreducible polynomial the centraliser is the norm
one torus of a field, the classes in the two groups are compared by class field theory, and the
sum of the orbital integrals of the kernel over a stable class is the value of a character of the
compact form, multiplied by a constant that depends only on the signature; this gives the equality
of the terms of the two groups (Proposition~\ref{pm:ell}). The archimedean statement is
Harish-Chandra's formula~\cite{HC66} for the
elliptic orbital integrals of a discrete series coefficient. We prove it in
\S\ref{pbr:sec} (Theorem~\ref{pbr:main}) for the groups $\auU(p,q)$, for parameters with a
large gap and for the coefficient of the lowest $K$-type, from the explicit formula for this
coefficient, Schmid's theorem on the $K$-types of a discrete series~\cite{Schmid75} and the theory
of compact Lie groups, without the theory of characters; the formal degree is obtained on the way
(Corollary~\ref{pbr:degree}).

The notation and the data are those of~\cite{P46}, and the deduction of the Spanning Theorem from
\eqref{B1:eq-i}--\eqref{B1:eq-iii} is given there. This paper and~\cite{P52} quote each other,
in the following order. Theorem~\ref{pbm:inputs} is stated here and proved in~\cite{P52}: its three parts are
{\cite[Theorems~2.17, 5.28 and~5.29]{P52}} of that paper. We use it only in the
proof of the main term (Theorem~\ref{pbm:L2}), that is, for part (L2) of Theorem~\ref{pb:amp}. We also use a lemma on the Haar measure of the real
points of $U$ ({\cite[Lemma~4.1]{P52}}), whose proof in~\cite{P52} does not refer to this paper. In the other direction \cite{P52} takes from this paper the
setting and the definition of a sieve system, the Bessel functions with their properties
(Theorems~\ref{pbe:ds} and~\ref{pbe:bessel}, Proposition~\ref{pbe:local}), the description of the
open cell and of the orbit sums (Propositions~\ref{pbm:cell} and~\ref{pbm:orbit}), lemmas on
measures (Lemmas~\ref{pbv:haar}, \ref{pbo:GF} and~\ref{pbo:append}), the structure of the level (Lemma~\ref{pb:KR}) and Witt's theorem in the form of Lemma~\ref{pbp:tb-witt}; their proofs use neither
Theorem~\ref{pbm:inputs} nor the main term. So these statements come first, then the three theorems of~\cite{P52}, then the main term. With~\cite{P49} the exchange is similar.
Two elementary lemmas of Section~\ref{pds:sec} (Lemmas~\ref{pds:blockdet} and~\ref{pds:fischer}), the estimate in the proof of Theorem~\ref{pds:decay} and
part (c) of Lemma~\ref{pm:kernel}, none of which uses the formula for the coefficient of
the lowest $K$-type, are used there to prove this formula. From~\cite{P49} we then use the formula with its consequences
({\cite[Theorem~3.17]{P49}}, {\cite[Proposition~3.21]{P49}}, {\cite[Lemma~3.25]{P49}}); on them rest Theorems~\ref{pm:arch} and~\ref{pds:decay}, parts (a) and (b) of Lemma~\ref{pm:kernel}, Lemma~\ref{pm:cusp} and Corollary~\ref{pbr:degree}. Finally we use the properties of
the forms of discrete series type ({\cite[Proposition~8.4]{P49}}), whose proofs use the statements just named
(in particular our decay bound and our formula for the formal degree) but
nothing that depends on that proposition. All these statements are reproduced in
Appendix~\ref{B1:quoted}. The dependencies of the whole proof, with the published results used,
are listed in \S\ref{pb:ss-remains}.

The paper is organised as follows. Section~\ref{pb:sec} contains the setting, the four
inequalities, the deduction of \eqref{B1:eq-i}--\eqref{B1:eq-iii} and of the Spanning Theorem from
them, the Petersson formula and the reduction of the first three inequalities to statements on
amplitudes. Section~\ref{B1:g2} computes the covolume $\nu$ (\S\ref{pbv:sec}), constructs the
Bessel functions and their Poincar\'e series and proves the expansion of the squares from the
estimates for pairs of terms (\S\ref{pbe:sec}), and proves the inequality for the induced
character, sets up sieve systems and the sieve, and obtains the main term from the local
statements (\S\ref{pbm:sec}). Section~\ref{B1:g3} proves the equality of multiplicities
(\S\ref{pm:sec}) and the formula for the elliptic orbital integrals (\S\ref{pbr:sec}).
Section~\ref{pds:sec} proves the properties of the discrete series coefficient at a real place,
except the lower bound for the Bessel period, which is the subject of Section~\ref{pbp:sec}.
Sections~\ref{ppi:sec} and~\ref{B1:g7} treat the pairs of terms: cell ranks and the exact
statement in Section~\ref{ppi:sec}, then the gain (\S\ref{pbg:sec}), the count
(\S\ref{pbc:sec}), the size of a block (\S\ref{pbs:sec}) and the pairing on the open cell
(\S\ref{pbo:sec}).


% ===== B-second-moment.tex =====
% proofs/B-second-moment.tex
% The second moment of the Poincare map: Proposition au:moment and Theorem B of
% sec-automorphic.tex. Labels of this file start with pb:.
\providecommand{\auO}{\mathcal O}
\providecommand{\auPl}{\Phi}
\providecommand{\auPa}{\Phi_a}
\providecommand{\auPb}{\Phi_b}
\providecommand{\auS}{\mathcal S}
\providecommand{\auSram}{\mathcal S^{\mathrm{ram}}}
\providecommand{\auT}{\mathcal T}
\providecommand{\auAz}{\mathcal A_0}
\providecommand{\auG}{\overline G}
\providecommand{\auP}{\mathcal P}
\providecommand{\auDX}{D_X}
\providecommand{\audT}{d_{\auT}}
\providecommand{\auNm}{\operatorname{Nm}}
\providecommand{\auU}{\mathrm U}
\providecommand{\auGL}{\mathrm{GL}}
\providecommand{\auA}{\mathbb A}
\providecommand{\auF}{\mathbb F}
\providecommand{\aurd}{\mathrm{rd}_E}
\providecommand{\aubar}[1]{\overline{#1}}
\providecommand{\auf}[1]{\texttt{#1}}
% macros of this file (prefix \pb)
\providecommand{\pbR}{\overline R_0}
\providecommand{\pbUb}{\overline U}
\providecommand{\pbHb}{\overline H}
\providecommand{\pbtr}{\operatorname{tr}}
\providecommand{\pbvol}{\operatorname{vol}}
\providecommand{\pbDid}{\mathcal D_{\mathrm{id}}}
\providecommand{\pbDmain}{\mathcal D_{\mathrm{main}}}
\providecommand{\pbdeg}{\operatorname{deg}}
\providecommand{\pbmass}{\operatorname{mass}}
\providecommand{\pbFix}{\operatorname{Fix}}
\providecommand{\pbpr}{\mathrm{pr}}
\providecommand{\pbrpsi}{r_\psi}
\providecommand{\pbW}{\mathbb W}
\providecommand{\pbg}{\mathsf g}
\providecommand{\pbLM}{\mathrm{LM}}
\providecommand{\pbnat}{\natural}
\providecommand{\pbone}{\mathbf 1}
\providecommand{\pbHS}[1]{\|#1\|}
\providecommand{\pbip}[2]{\langle#1,#2\rangle}

\section{The second moment of the Poincar\'e map}\label{pb:sec}

This section concerns {\cite[Proposition~14.32]{P46}} (the second moment of the Poincar\'e map) and,
through it, the Spanning Theorem ({\cite[Theorem~14.14]{P46}}). We use the notation of {\cite[\S\S14.1--14.3]{P46}}.

\emph{Structure of the proof.} The proof has three layers.
\begin{itemize}
\item[(I)] \emph{From four inequalities to {\cite[Theorem~14.14]{P46}}} (\S\S\ref{pb:ss-setting}--\ref{pb:ss-thmB}).
The Poincar\'e map $\auP$ intertwines an explicit unitary representation $\pbrpsi$ of a subgroup
$\pbR$ of $\auG$ on the source with the action of $\auG$ on the forms on $Y_0$
(Lemma~\ref{pb:equiv}). We state four inequalities for the operators
$P_s=P^*\varrho(s)P$, $P=c_P\auP$: \textup{(E1)} $P_s$ is close to $\pbrpsi(s)$ near $\pbR$;
\textup{(E2)}, \textup{(E3)} two bounds for sums of $\pbHS{P_s}^2$ over subgroups of $\auG$ that fix
vectors of the frame; \textup{(E4)} a bound for the character induced from the unipotent
character. From these we deduce, with complete proofs,
{\cite[Proposition~14.32(ii)]{P46}} and (iii) (Propositions~\ref{pb:ii} and~\ref{pb:iii}). Part (i) is
deduced from {\cite[Proposition~14.36]{P46}} (Proposition~\ref{pb:i}). Together with the deduction of
{\cite[\S14.7]{P46}} this gives {\cite[Theorem~14.14]{P46}} (Theorem~\ref{pb:thmB}).
\item[(II)] \emph{From amplitudes to} (E1)--(E3) (\S\ref{pb:ss-layer2}). Assuming that $\auP$ is
a Poincar\'e series of a Bessel function with the usual properties
(Hypothesis~\ref{pb:bessel}), we prove the Petersson formula for $P_s$
(Proposition~\ref{pb:unfold}), identify the contribution of the Bessel subgroup with $\pbrpsi(s)$
(Proposition~\ref{pb:besselterms}), compare the two normalising constants $\pbDmain$ and $\pbDid$
(Lemma~\ref{pb:sizes}; this is where the window~{\cite[(37)]{P46}} for the level $N$ is used),
and deduce (E1)--(E3) from three statements about sums of amplitudes
(Proposition~\ref{pb:reduction}).
\item[(III)] \emph{The analytic core}: the existence of the Bessel functions, the three statements
about amplitudes, the inequality (E4), and {\cite[Proposition~14.36]{P46}}. They are stated precisely here
(Hypothesis~\ref{pb:bessel}, Theorem~\ref{pb:amp}, Theorem~\ref{pb:est}(E4), {\cite[Proposition~14.36]{P46}})
and proved in \S\S\ref{pbe:sec}--\ref{pm:sec}, \S\S\ref{pds:sec}--\ref{pbo:sec} and~{\cite[\S\S2--6]{P52}}: the Bessel functions in Theorem~\ref{pbe:bessel},
\textup{(X)} and \textup{(L1)} in Theorem~\ref{pbe:XL1}, \textup{(L2)} in Theorem~\ref{pbm:L2},
\textup{(E4)} in Theorem~\ref{pbm:E4}, and {\cite[Proposition~14.36]{P46}} in Proposition~\ref{pm:mult}.
\end{itemize}
So {\cite[Theorem~14.14]{P46}} is proved modulo the published results named in the proofs; the value of the
covolume $\nu$ (Lemma~\ref{pb:covol}) is proved in \S\ref{pbv:sec}. (Part~(b) of Theorem~\ref{pm:arch}, Harish-Chandra's formula for stable
elliptic orbital integrals of discrete series coefficients, is proved in \S\ref{pbr:sec}.) The complete dependency tree is in \S\ref{pb:ss-remains}.

\subsection{The setting}\label{pb:ss-setting}

\subsubsection*{The group and the frame}
Put $n=\varrho+1$, so that $r=k+2\varrho+1$. The space $Y_0=\mathbb H^{\varrho}\perp X_0\perp
\langle-N\rangle$ has the basis
\[
e_1,\dots,e_\varrho;\quad x_1,\dots,x_k,\ z;\quad e_1^*,\dots,e_\varrho^*,
\]
with $h(e_a,e_b^*)=\delta_{ab}$, $h(x_i,x_j)=\delta_{ij}$, $h(z,z)=-N$ and all other products
zero. The \emph{frame} is $(f_0,\dots,f_\varrho)=(e_1,\dots,e_\varrho,z)$; $A=\langle e_1,\dots,
e_\varrho\rangle$ is isotropic and $A^\perp/A=X_0\perp Ez$. Write $G=\auU(Y_0)$, a group over
$E^+$; $\Gamma=G(E^+)$; $G_\infty=G(E^+\otimes\R)=\prod_{\tau\in\auPl}G_\tau$;
$G_f=G(\auA_{E^+,f})$. Let $K\subset G_f$ be the stabiliser of the lattice $L$ of $Y_0$ described in the proof of
Lemma~\ref{pb:KR} (for its relation to the lattice $\mathsf M_0$ of {\cite[(35)]{P49}} see
Lemma~\ref{pb:lattice}) and $J=J_M
\subseteq K$ the principal congruence subgroup of level $\pi_M$; it is normal in $K$. Each chosen
place $w$ is split over $E^+$, so the component of $K$ at the place below $w$ is
$\auGL_r(\auO_{E,w})$, and reduction modulo the chosen places identifies $K/J$ with
$\auG=\prod_w\auGL_r(\auF_w)$. We write $\auG=G_0\times G_1$, where $G_i$ is the product over the
$f$ chosen places above $p_i$. For $s\in\auG$, $\varrho(s)$ is right translation by a lift of $s$
to $K$; it is a unitary operator of $\auAz$, for the inner product $\pbip y{y'}=\int_{\Gamma
\backslash G(\auA)}\aubar y\,y'$. \emph{In this section all inner products are conjugate-linear in
the first variable and linear in the second}; this is the convention for which the map $\auP$ of
{\cite[\S14.3]{P46}}, defined by $\pbip y{\auP\varphi}=\int_{[R_0]}\mathcal W_\varphi(r)\aubar{y(rg_0)}
\,dr$, is linear.

\subsubsection*{Three conventions}\label{pb:ss-conv}
The following conventions hold in all the sections on {\cite[Theorem~14.14]{P46}} (\S\S\ref{pb:sec}--\ref{pm:sec},
\S\S\ref{pds:sec}--\ref{pbo:sec} and~{\cite[\S\S2--6]{P52}}); sections that use another one say so and point here.
\begin{enumerate}
\item[(Cv1)] \emph{Hermitian forms.} Every hermitian form $h$ on a space over $E$ or over a completion
of $E$ ($Y_0$, $X_0$, the $W_j$, $Y_\tau$, \dots) is \emph{linear in the first variable} and
conjugate-linear in the second: $h(\lambda u,v)=\lambda h(u,v)=h(u,\aubar\lambda v)$, $h(v,u)=
\aubar{h(u,v)}$. Thus $h(v,e_a^*)$ is the coefficient of $e_a$ in $v$. (The inner products of
functions and of the spaces $\pbW_\tau$ above follow the opposite convention; they are not
hermitian forms on the spaces $Y_j$. For representation coefficients: $c_\tau(g)=\langle v_0,\pi(g)v_0\rangle$
of \S\ref{pbe:ss-ds} and~{\cite[\S3.2]{P49}}, with the inner product conjugate-linear in the first variable, is
the same function as $\langle\pi(g)v_0,v_0\rangle'$ for the inner product $\langle u,w\rangle'=\langle w,u\rangle$,
linear in the first variable, used in \S\S\ref{pm:sec}, \ref{pbr:sec}; no conjugation separates them.)
\item[(Cv2)] \emph{The additive character.} $\psi_E=\psi_\Q\circ\operatorname{Tr}_{E/\Q}$, where
$\psi_\Q$ is the standard character of $\auA_\Q/\Q$: $\psi_{\Q,\infty}(t)=e^{-2\pi it}$ and
$\psi_{\Q,p}(t)=e^{2\pi i\{t\}_p}$. So at every complex place $\tau$ of $E$ the component of $\psi_E$ is
$w\mapsto e^{-2\pi i(w+\aubar w)}$, and $\psi_E$ is trivial exactly on the inverse different
$\varpi^{-1}\widehat\auO_E$ at the finite places. (By Lemma~\ref{pbp:sign} the opposite sign at
infinity would change nothing in the results.)
\item[(Cv3)] \emph{The Haar measure of $U_\infty$.} On $U_\tau$, $\tau\in\auPl$, $du_\tau$ is
Lebesgue measure in the free entries of the rows $\tilde e_a$ of $u$ in the normalised basis of
\S\ref{pbe:ss-real} (the complex entries above the diagonal block that are not fixed by unitarity,
and the imaginary parts of the entries $(\tilde e_a,\tilde e_a^*)$); equivalently, $du_\tau$ is the Haar
measure with $\int_{U_\tau}G_\varrho(u^{-1})\,du_\tau=1$ for the Gaussian $G_\varrho$ of
{\cite[\S4.1]{P52}} ({\cite[Lemma~4.1]{P52}}). On $U_\infty$, $du=\prod_\tau du_\tau$, and on
$U_\tau\backslash G_\tau$ the quotient measure. This is the measure of \S\S\ref{pbe:ss-real},
\ref{pbs:ss-measures} and of {\cite[Lemma~2.19]{P52}}, and the constant $\nu$ of Lemma~\ref{pb:ortho}
is the covolume $\pbvol(U_c\backslash U_\infty)$ for it.
\end{enumerate}

Let $U\subset G$ be the unipotent radical of the parabolic subgroup $P$ of {\cite[\S14.3]{P46}}:
the elements that preserve the flag $\langle e_1\rangle\subset\langle e_1,e_2\rangle\subset\dots
\subset A$, act trivially on its graded pieces and act trivially on $A^\perp/A$. Let
$H=\auU(X_0)\subset G$, acting trivially on $A$, $z$ and the $e_a^*$, and $R_0=U\rtimes H$. For
$u\in U$ put
\[
\sigma(u)=\sum_{a=1}^{\varrho-1}h(ue_{a+1},e_a^*)+h(uz,e_\varrho^*),
\]
the sum of the entries of $u$ just above the diagonal of its block on the frame; $\sigma$ is
additive on $U$ and invariant under conjugation by $H$. The character $\psi_U$ of
{\cite[\S14.3]{P46}} is taken to be
\begin{equation}\label{pb:eq-psi}
\psi_U(u)=\psi_E\Bigl(-\frac{\sigma(u)}{M\varpi}\Bigr)\qquad(u\in U(\auA_{E^+})),
\end{equation}
where $\psi_E$ is the standard additive character of $\auA_E/E$ and $\varpi=\gamma\,\Phi'_{q_1}(\zeta)\sqrt{-d}\in\auO_E$ is the purely imaginary generator of the different of $E$ of Lemma~\ref{pds:kappa} ($\zeta$ a primitive $q_1$-th root of unity, $\gamma$ a unit with $\gamma\aubar\gamma=1$ making $\gamma\Phi'_{q_1}(\zeta)$ real). Thus
$\psi_U$ is trivial on $U(E^+)$, invariant under conjugation by $H(\auA_{E^+})$, and trivial on
$U(\auA_{E^+,f})\cap J$: for $u\in J$ the entries of $u-1$ in the basis of the lattice vanish
modulo every chosen place and its conjugate, so they lie in $M\widehat{\auO}_E$, and $\psi_E$ is
trivial on $\varpi^{-1}\widehat{\auO}_E$. We write $\psi_\infty$, $\psi_f$ for the restrictions of
$\psi_U$ to $U_\infty=U(E^+\otimes\R)$ and to $U(\auA_{E^+,f})$; $\psi_\infty(u)\psi_f(u)=1$ for
$u\in U(E^+)$. (This is the character with which $\auP$ is defined in {\cite[\S14.3]{P46}}, and it
is the character $\psi_0$ of {\cite[\S5.1.3]{P49}}; likewise the pair $(g_0,\mathsf t)$ with which $\auP$
is defined there is the pair of Theorem~\ref{pb:est} below.)

\subsubsection*{The finite groups}
Put $K_U=U(\auA_{E^+,f})\cap K$, $K_H=H(\auA_{E^+,f})\cap K$ and $K_R=R_0(\auA_{E^+,f})\cap K$.

\begin{lemma}\label{pb:KR}
$K_R=K_U\rtimes K_H$. The group $K_H$ consists of the elements of the stabiliser $K^\circ$ of
$\auO_E^k$ in $H(\auA_{E^+,f})$ whose last row is $\equiv(0,\dots,0,1)$ modulo $\pi_N$, and
$K_H\cap J=J_X$.
\end{lemma}
\begin{proof}
Let $g=uh\in K$ with $u\in U(\auA_f)$, $h\in H(\auA_f)$, and let $L$ be the lattice of $Y_0$. Then
$g$ preserves $L\cap A$ and $L\cap A^\perp$, hence the lattice $(L\cap A^\perp)/(L\cap A)$ of
$A^\perp/A$, on which $u$ acts trivially; so $h$ preserves this lattice. As $h$ fixes the $e_a$
and the $e_a^*$, and $L$ is the sum of its intersections with $A\oplus A^*$ and with $X_0\perp Ez$,
$h$ preserves $L$. So $h\in K_H$ and $u=gh^{-1}\in K_U$. The lattice of $X_0\perp Ez$ is spanned
by $x_1,\dots,x_{k-1}$, $x_k/\aubar{\pi_N}+z/N$ and $z$ (it is unimodular: the last two vectors
have the Gram matrix $\bigl(\begin{smallmatrix}0&-1\\-1&-N\end{smallmatrix}\bigr)$), and its
intersection with $X_0$ is $L_1=\auO x_1+\dots+\auO x_{k-1}+\pi_N\auO x_k$. An element $h$ of
$H(\auA_f)$, which fixes $z$, preserves this lattice if and only if $hL_1=L_1$ and $(hx_k-x_k)/
\aubar{\pi_N}\in L_1$. Away from $N$ this says $h\in K^\circ$. At a prime $w\mid\pi_N$ the group
$H$ is $\auGL_k(E_w)$, the component at $\aubar w$ being the inverse adjoint, and $\aubar{\pi_N}$
is a unit at $w$. Read at $w$, the second condition says $h_we_k\in e_k+L_{1,w}=\auO_w^{k-1}\oplus
(1+\pi_N\auO_w)e_k$, and then the first, applied to $e_1,\dots,e_{k-1}\in L_{1,w}$, says that the
entries $h_{ij}$ ($j<k$) are integral with $h_{kj}\in\pi_N\auO_w$. Thus $h_w\in\auGL_k(\auO_w)$ with
last row $\equiv(0,\dots,0,1)$ modulo $w$; conversely such an $h_w$ satisfies both conditions at $w$.
Read at $\aubar w$, the two conditions give the same congruence for the last row of $h_w^{-1}$,
which is equivalent.
Finally $K_H\cap J$ is defined by the congruence $\equiv1$ at the chosen places, which
gives the level $J_X$ of {\cite[\S14.3]{P46}}.
\end{proof}

Let $\pbUb$, $\pbHb$, $\pbR$ be the images of $K_U$, $K_H$, $K_R$ in $\auG$. By the lemma
$\pbR=\pbUb\rtimes\pbHb$, and each of the three groups is the product of its components at the
chosen places: $\pbHb\cong K_H/J_X\cong\prod_w\auGL_k(\auF_w)$, and $\pbUb_w$ is the unipotent
radical of the parabolic subgroup of $\auGL_r(\auF_w)$ with blocks $1,\dots,1,k+1,1,\dots,1$, of
order $(\auNm w)^{\varrho(2k+2\varrho+1)}$. (For $\pbUb$ one uses that $U$ is a split unipotent
group, so that $U(\auO_w)\to U(\auF_w)$ is onto.) Hence
\begin{equation}\label{pb:eq-orders}
\#\pbUb=M^{f\varrho(2k+2\varrho+1)},\qquad\#\pbR=\#\pbUb\cdot\#\pbHb .
\end{equation}
By the remarks after \eqref{pb:eq-psi}, $\psi_f$ induces a character of $\pbUb$, again written
$\psi_U$, which is invariant under conjugation by $\pbHb$. Let $\tilde\psi\colon\auG\to\C$ be
$\psi_U$ on $\pbUb$ and $0$ elsewhere, and
\[
\psi^{\pbnat}(s)=\frac1{\#\auG}\sum_{g\in\auG}\tilde\psi(gsg^{-1})\qquad(s\in\auG).
\]

\subsubsection*{The source and the representation \texorpdfstring{$\pbrpsi$}{r}}
The source $\auS=\mathcal A_{\lambda^X}(J_X,1)$ carries the inner product of
{\cite[\S14.3]{P46}}, $\pbip\varphi{\varphi'}=\sum_{c}\pbip{\varphi(h_c)}{\varphi'(h_c)}_{W_{\lambda^X}}$,
where $c$ runs over the set $\mathcal C=H(E^+)\backslash H(\auA_f)/J_X$ of \emph{classes} and $h_c$
is a representative of $c$ (the level is free by {\cite[Lemma~14.5]{P46}}, so there are no weights), and
where the hermitian form on $W_{\lambda^X}$ is invariant under $H_\infty=H(E^+\otimes\R)$. Since
the level is free, $\varphi\mapsto(\varphi(h_c))_c$ is an isomorphism $\auS\to W_{\lambda^X}^{\,
\mathcal C}$, and
\begin{equation}\label{pb:eq-DX}
\auDX=\#\mathcal C\cdot d,\qquad d=\dim W_{\lambda^X}=\prod_{\tau\in\auPa}\binom{a+k-1}{k-1}.
\end{equation}
The group $K_H$ normalises $J_X$, so $\pbHb=K_H/J_X$ acts on $\auS$ by right translation,
$(R(t)\varphi)(g)=\varphi(gk_t)$; this action is unitary. For $s\in\auG$ put
\[
\pbrpsi(s)=\begin{cases}\psi_U(\bar u)\,R(\bar t)&\text{if }s=\bar u\,\bar t\in\pbR,\
\bar u\in\pbUb,\ \bar t\in\pbHb,\\[2pt] 0&\text{if }s\notin\pbR.\end{cases}
\]
Because $\psi_U$ is a character of $\pbUb$ invariant under $\pbHb$, the restriction of $\pbrpsi$ to
$\pbR$ is a unitary representation of $\pbR$ on $\auS$, and
\begin{equation}\label{pb:eq-rext}
\pbrpsi(st)=\pbrpsi(s)\pbrpsi(t),\quad\pbrpsi(ts)=\pbrpsi(t)\pbrpsi(s),\quad\pbrpsi(s)^*=\pbrpsi
(s^{-1})\qquad(t\in\pbR,\ s\in\auG).
\end{equation}

\begin{lemma}[Equivariance]\label{pb:equiv}
For every choice of $g_0$ and $\mathsf t$ in the definition of $\auP$, and every $t\in\pbR$,
$\varrho(t)\circ\auP=\auP\circ\pbrpsi(t)$.
\end{lemma}
\begin{proof}
The group $R_0(\auA)$ is unimodular: $U$ is unipotent, and the determinant of the action of $H$
on $\operatorname{Lie}U$ is an algebraic character $H\to\mathbb G_m$ over $E^+$, hence trivial,
because the characters of the unitary group $H$ factor through $\det\colon H\to\auU(1)$ and
$\auU(1)$ is anisotropic. So conjugation by $H(\auA)$ preserves the Haar measure of $U(\auA)$, and
the measure on $[R_0]=R_0(E^+)\backslash R_0(\auA)$ is invariant under right
translations. Let $k=u_kh_k\in K_R$ be a lift of $t$, $y\in\auAz$ and $\varphi\in\auS$. As $k$ is
a finite adele and $g_0$ is real, they commute, and the substitution $r\mapsto rk$ gives
\[
\pbip y{\varrho(k)\auP\varphi}=\pbip{\varrho(k)^{-1}y}{\auP\varphi}=\int_{[R_0]}\mathcal W_\varphi(r)
\,\aubar{y(rg_0k^{-1})}\,dr=\int_{[R_0]}\mathcal W_\varphi(rk)\,\aubar{y(rg_0)}\,dr .
\]
For $r=uh$ we have $rk=u\cdot(hu_kh^{-1})\cdot hh_k$, so, by the invariance of $\psi_U$ under $H$,
$\mathcal W_\varphi(rk)=\psi_U(u)\psi_U(u_k)\,\mathsf t\bigl(h_\infty^{-1}\varphi(h_fh_k)\bigr)=
\psi_U(u_k)\,\mathcal W_{R(h_k)\varphi}(r)$. Hence $\pbip y{\varrho(k)\auP\varphi}=\psi_U(u_k)\pbip y
{\auP R(h_k)\varphi}=\pbip y{\auP\,\pbrpsi(t)\varphi}$, as $\auP$ is linear and the inner product is
linear in the second variable.
\end{proof}

\begin{lemma}[Trace of $\pbrpsi$]\label{pb:trace}
Let $s\in\auG$ have trivial components at all chosen places above one of the primes $p_0$, $p_1$.
Then $\pbtr\pbrpsi(s)=\auDX\cdot\tilde\psi(s)$.
\end{lemma}
\begin{proof}
If $s\notin\pbR$ both sides vanish. Let $s=\bar u\bar t$. As $\pbR=\prod_w\pbUb_w\rtimes\pbHb_w$,
the components of $\bar t$ above that prime, $p_j$ say, are trivial. Let $k\in K_H$ be a lift of
$\bar t$. In the decomposition $\auS\cong W_{\lambda^X}^{\,\mathcal C}$, $(R(\bar t)\varphi)(h_c)=
\varphi(h_ck)=\eta\,\varphi(h_{c'})$ if $h_ck=\eta h_{c'}j$ with $\eta\in H(E^+)$, $j\in J_X$. So
$R(\bar t)$ permutes the blocks, and a block contributes to the trace only if $c'=c$. In that case
$\eta\in H(E^+)\cap h_cK_Hh_c^{-1}$ has finite order and $h_c^{-1}\eta h_c=kj^{-1}$ is $\equiv1$
modulo a chosen place above $p_j$. The argument of {\cite[Lemma~14.5]{P46}}, which uses only one chosen
place, gives $\eta=1$, so $k\in J_X$ and $\bar t=1$. Thus the trace is $0$ unless $\bar t=1$, that
is $s\in\pbUb$, and then $\pbrpsi(s)=\psi_U(s)\cdot1$.
\end{proof}

\subsubsection*{The constants}
For a positive definite standard space of rank $l$ let $\pbmass_l=\sum_{[g]}\#(\Gamma_l\cap
gK^\circ_lg^{-1})^{-1}$ be the mass of the genus of $\auO_E^l$, the sum being over $\auU_l(E^+)
\backslash\auU_l(\auA_f)/K_l^\circ$. If $J'\subseteq K_l^\circ$ is a subgroup of finite index which is free, then
\begin{equation}\label{pb:eq-classes}
\#\bigl(\Gamma_l\backslash\auU_l(\auA_f)/J'\bigr)=[K_l^\circ:J']\cdot\pbmass_l ,
\end{equation}
because the finite group $\Gamma_l\cap gK_l^\circ g^{-1}$ acts freely on $gK_l^\circ/J'$, by left
multiplication. For $\tau\in\auPl$ let
\[
\pbdeg_\tau=\prod_{1\le i<j\le r}\frac{\varsigma_{\tau,i}-\varsigma_{\tau,j}}{j-i},\qquad
\pbdeg=\prod_{\tau\in\auPl}\pbdeg_\tau ,
\]
where $\varsigma_\tau=(\varsigma_{\tau,1}>\dots>\varsigma_{\tau,r})$ is the parameter of
{\cite[\S14.3]{P46}}. By Weyl's formula $\pbdeg_\tau$ is the dimension of the irreducible
representation of the compact group $\auU(r)$ with infinitesimal character $\varsigma_\tau$; it is
the formal degree of the discrete series $\varsigma_\tau$ for the Haar measure that corresponds
to the measure of total mass $1$ on $\auU(r)$. This is the number $\deg$ of
{\cite[Proposition~14.32]{P46}}. Put
\begin{equation}\label{pb:eq-D}
\pbDid=\#\pbR\cdot\auDX,\qquad\pbDmain=\frac{\auDX^2}{\pbdeg\cdot\pbmass_r}.
\end{equation}

\subsubsection*{Vectors of the frame modulo the chosen places}
Let $w$ be a chosen place and $V_w=\auF_w^{\,r}$ the reduction modulo $w$ of the lattice of $Y_0$,
in a basis of the lattice containing the $e_a$, $z$ and the $e_a^*$. The hermitian form pairs the
completions at $w$ and at $\aubar w$; in this basis it gives a non-degenerate symmetric bilinear
form on $V_w$, with matrix $\mathsf J$. An element $s\in\auG$ acts on $V_w$ by $s_w$, and we let it
act on a second copy $V_{\aubar w}$ of $V_w$ by the contragredient $s_w^{\vee}=\mathsf J^{-1}\,
{}^ts_w^{-1}\mathsf J$; this is the action modulo $\aubar w$ of a lift of $s$ to $K$. For
$\mathfrak p\in\{w,\aubar w\}$ let $F_{\mathfrak p}\subset V_{\mathfrak p}$ be the span of the
frame, of dimension $n$. A \emph{frame vector at $w$} is a non-zero vector $\mathsf d$ of $F_w$ or
of $F_{\aubar w}$; $\pbFix(\mathsf d)\subseteq\auG$ is its stabiliser. For a set $\mathsf A$ of
chosen places above $p_0$ and a family $\mathsf d=(\mathsf d_w)_{w\mid p_0}$ of frame vectors put
$\pbFix(\mathsf A,\mathsf d)=\bigcap_{w\in\mathsf A}\pbFix(\mathsf d_w)$; it is a subgroup of
$\auG$, equal to $\auG$ if $\mathsf A=\emptyset$. For $g\in\auGL_r(\auF_w)$ let
\[
Z_w(g)=\bigl({}^tf_j\,\mathsf J\,g\,f_l\bigr)_{0\le j,l\le\varrho}\in M_n(\auF_w),\qquad
\delta_l(g)=\det\bigl(Z_w(g)_{j,l'}\bigr)_{0\le j,l'<l}\quad(0\le l\le n),
\]
the matrix of the products of the frame with its image, and its leading minors. An element
$s\in\auG$ is \emph{open} if there is a chosen place $w$ above $p_0$ with $\delta_l(s_w)\ne0$ for
$1\le l\le\varrho$.

\subsection{The four inequalities}\label{pb:ss-est}

All norms of linear maps are Hilbert--Schmidt norms. For a linear map $P\colon\auS\to\auAz$ and
$s\in\auG$ we write $P_s=P^*\varrho(s)P\in\operatorname{End}(\auS)$.

\begin{theorem}[The four inequalities]\label{pb:est}
There are a point $g_0\in G_\infty$, a linear form $\mathsf t$ on $W_{\lambda^X}$ and a scalar
$c_P>0$ such that the map $P=c_P\auP$, with $\auP$ defined by $g_0$, $\mathsf t$ and the character
\eqref{pb:eq-psi}, has the following properties.
\begin{enumerate}
\item[(E1)] For $i\in\{0,1\}$ let $\Omega_i$ be the set of the $s\in\auG$ whose components at the
chosen places above $p_i$ are those of an element of $\pbR$. Then
$\sum_{s\in\Omega_i}\pbHS{P_s-\pbrpsi(s)}^2\le b^{-12}\,\pbDid$.
\item[(E2)] For every family $\mathsf d$ of frame vectors at the $f$ chosen places above $p_0$,
$\sum_{s\in\pbFix(\mathrm{all},\mathsf d)}\pbHS{P_s}^2\le b^{-6kf}\,\pbDmain$.
\item[(E3)] For every set $\mathsf A$ of chosen places above $p_0$ and every family $\mathsf d$,
\[
\sum_{s\in\pbFix(\mathsf A,\mathsf d),\ s\text{ open}}\pbHS{P_s}^2\le\bigl(1+4b^{-4}\bigr)\,b^{-6k
\#\mathsf A}\,\pbDmain .
\]
\end{enumerate}
Moreover
\begin{enumerate}
\item[(E4)] $\sum_{s\in\auG}|\psi^{\pbnat}(s)|^2\le2$.
\end{enumerate}
\end{theorem}
\smallskip\noindent(E1)--(E3) are deduced in \S\ref{pb:ss-layer2} (Proposition~\ref{pb:reduction}) from
Hypothesis~\ref{pb:bessel} and Theorem~\ref{pb:amp}, proved in \S\S\ref{pbe:sec}, \ref{pbm:sec},
\S\S\ref{pds:sec}--\ref{pbo:sec} and~{\cite[\S\S2--6]{P52}}. (E4) is a statement
about the finite group $\auG$ and the character $\psi_U$ of $\pbUb$ only; it is proved in \S\ref{pbm:ss-E4} (Theorem~\ref{pbm:E4}), using $k\ge2\varrho^2+3$, $p_i\ge b^{10}$ and $18\varrho f\le b^8$ ({\cite[Lemma~13.4(S7)]{P46}}).\par\smallskip

\begin{remark}\label{pb:rem-E4}
(E4) says that the representation of $\auG$ induced from the character $\psi_U$ of $\pbUb$ is
almost as free of multiplicities as its dimension allows: a computation with double cosets gives
$\sum_s|\psi^{\pbnat}(s)|^2=(\#\pbUb)^2(\#\auG)^{-1}\cdot\dim\operatorname{End}_{\auG}
\bigl(\operatorname{Ind}_{\pbUb}^{\auG}\psi_U\bigr)$, and $\#\auG/(\#\pbUb)^2$ is the number of
double cosets of $\pbUb$ of full size, if all had full size.
\end{remark}

\subsection{Lemmas on finite groups}\label{pb:ss-finite}

\begin{lemma}[Central idempotents]\label{pb:central}
Let $\mathsf G$ be a finite group and $x=\sum_sx(s)s$ a central idempotent of $\C[\mathsf G]$.
\begin{enumerate}
\item[(a)] $x(s^{-1})=\aubar{x(s)}$, and $\sum_s|x(s)|^2=x(1)\in[0,1]$.
\item[(b)] For a unitary representation $\varrho$ of $\mathsf G$, $\varrho(x)=\sum_sx(s)\varrho(s)$
is an orthogonal projection that commutes with $\varrho(\mathsf G)$.
\item[(c)] If $\mathsf G$ is a subgroup of a finite group $\mathsf G'$ and $F$ is a function on
$\mathsf G'$ with values in a Hilbert space, then
$\sum_{s\in\mathsf G'}\bigl\|\sum_{t\in\mathsf G}x(t)F(st)\bigr\|^2\le\sum_{s\in\mathsf G'}\|F(s)
\|^2$.
\end{enumerate}
\end{lemma}
\begin{proof}
A central idempotent is a sum of the idempotents $e_\chi=\chi(1)(\#\mathsf G)^{-1}\sum_s
\aubar{\chi(s)}s$ over a set of irreducible characters $\chi$. This gives the first claim of (a)
and $x(1)=\sum_\chi\chi(1)^2/\#\mathsf G\in[0,1]$; and $x(1)=(x\cdot x)(1)=\sum_sx(s)x(s^{-1})=
\sum_s|x(s)|^2$. (b) $\varrho(x)^*=\sum_s\aubar{x(s)}\varrho(s^{-1})=\varrho(x)$ by (a), and
$\varrho(x)^2=\varrho(x)$. (c) Apply (b) to the unitary representation $(R_tF)(s)=F(st)$ of
$\mathsf G$ on the functions on $\mathsf G'$: $R(x)$ is an orthogonal projection.
\end{proof}

A real function $\beta$ on a finite group $\mathsf G$ is \emph{of positive type} if there are
vectors $u_a$ ($a\in\mathsf G$) of a complex inner product space with $\pbip{u_a}{u_b}=\beta(a^{-1}
b)$ for all $a,b$. Then $\beta(s)\le\beta(1)$ for all $s$, by the Cauchy--Schwarz inequality, as
$\|u_a\|^2=\beta(1)$. The restriction of $\beta$ to a subgroup is of positive type.

\begin{lemma}[Singular minors]\label{pb:singular}
Let $\mathsf G$ be a finite group, $\auF$ a field with $q$ elements, $\pi\colon\mathsf G\to\auGL_r
(\auF)$ a homomorphism, $A_0,A_1$ two $r\times l$ matrices over $\auF$, $\beta\ge0$ a function of
positive type on $\mathsf G$ and $Y\ge0$. Suppose that for every non-zero $v\in\auF^l$
\[
\sum_{s:\ \pi(s)A_0v=A_0v}\beta(s)\le Y\qquad\text{and}\qquad\sum_{s:\ {}^t\pi(s)^{-1}A_1v=A_1v}
\beta(s)\le Y .
\]
Then
\[
q\sum_{s:\ \det({}^tA_1\pi(s)A_0)=0}\beta(s)\ \le\ 2\sum_{s\in\mathsf G}\beta(s)+q^{\,l+r/2}\,Y .
\]
\end{lemma}
\begin{proof}
Let $\psi$ be a non-trivial additive character of $\auF$, $Z(s)={}^tA_1\pi(s)A_0$, and for $y,v\in
\auF^l$ let $\Sigma(y,v)=\sum_s\psi({}^tyZ(s)v)\beta(s)$. If $\det Z(s)=0$ the left kernel of
$Z(s)$ has at least $q$ elements, and by orthogonality of characters $q^{-l}\sum_v\psi({}^tyZ(s)
v)$ is $1$ if ${}^tyZ(s)=0$ and $0$ otherwise. So
\[
q\sum_{\det Z(s)=0}\beta(s)\le\sum_s\#\{y:{}^tyZ(s)=0\}\,\beta(s)=\sum_y\,q^{-l}\sum_v\Rel\Sigma(y,
v).
\]
If $y=0$ or $v=0$ then $\Sigma(y,v)=\sum_s\beta(s)$; this happens for $q^l$ pairs with $y=0$ and
for one $v$ for each of the other $y$, so these pairs contribute at most $2\sum_s\beta(s)$. It
remains to prove $\Rel\Sigma(y,v)\le q^{r/2}Y$ for $y,v\ne0$; then the remaining pairs contribute at
most $q^l\cdot q^{r/2}Y$.

Put $a=A_1y$, $z=A_0v$, so that ${}^tyZ(s)v={}^ta\,\pi(s)z$. Let $\mathcal H$ be the space of the
$u_a$ and $\ell^2(\auF^r)$ the space of functions on $\auF^r$, with the orthonormal bases
$(\epsilon_\xi)_\xi$ of delta functions and $(c_\xi)_\xi$, $c_\xi(\zeta)=q^{-r/2}\psi(-{}^t\xi
\zeta)$. Let
\[
X_1=\sum_{g\in\mathsf G}c_{{}^t\pi(g)^{-1}a}\otimes u_g,\qquad X_0=\sum_{g\in\mathsf G}\epsilon_{\pi(g)
z}\otimes u_g\qquad\text{in }\ell^2(\auF^r)\otimes\mathcal H .
\]
Since $\pbip{c_\xi}{\epsilon_\zeta}=\aubar{c_\xi(\zeta)}=q^{-r/2}\psi({}^t\xi\zeta)$ and ${}^t\bigl(
{}^t\pi(g)^{-1}a\bigr)\pi(g')z={}^ta\,\pi(g^{-1}g')z$, we get
\[
\pbip{X_1}{X_0}=q^{-r/2}\sum_{g,g'}\psi\bigl({}^ta\,\pi(g^{-1}g')z\bigr)\beta(g^{-1}g')=q^{-r/2}\,\#
\mathsf G\cdot\Sigma(y,v) .
\]
In the same way, by orthonormality, $\|X_0\|^2=\sum_{g,g'}[\pi(g)z=\pi(g')z]\,
\beta(g^{-1}g')=\#\mathsf G\sum_{s:\pi(s)z=z}\beta(s)\le\#\mathsf G\cdot Y$, and $\|X_1\|^2=\#
\mathsf G\sum_{s:{}^t\pi(s)^{-1}a=a}\beta(s)\le\#\mathsf G\cdot Y$. Hence $\Rel\Sigma(y,v)=q^{r/2}
(\#\mathsf G)^{-1}\Rel\pbip{X_1}{X_0}\le q^{r/2}(\#\mathsf G)^{-1}\|X_1\|\|X_0\|\le q^{r/2}Y$.
\end{proof}

\subsubsection*{Labels and bad types}
Let $\auF$ be a field with $q$ elements. For a label $c$ of degree $r$ over $\auF$
({\cite[\S14.3]{P46}}), with irreducible factors $c_1,\dots,c_h$ of degrees $d_1,\dots,d_h$, the
irreducible character $\Xi_c$ of $\auGL_r(\auF)$ named by $c$ is, in Green's
parametrisation~\cite[p.~403 and Thm.~13]{Green1955} (see also~\cite[Ch.~IV, (6.7)--(6.8)]{Macdonald}), the character induced from the
parabolic subgroup with blocks $d_1,\dots,d_h$ of the product of the cuspidal characters of
$\auGL_{d_j}(\auF)$ attached to the roots of $c_j$ (through a fixed embedding of $\overline{\auF}{}^
\times$ in $\C^\times$). We use the following facts of~\cite[Lemma~7.4 and Thms.~12--13]{Green1955}: the cuspidal characters of
$\auGL_d(\auF)$ have degree $\prod_{i=1}^{d-1}(q^i-1)$; $\Xi_c$ is irreducible, because the
$c_j$ are distinct; and $\Xi_c\ne\Xi_{c'}$ for $c\ne c'$. Counting the cosets of the parabolic
subgroup,
\begin{equation}\label{pb:eq-green}
\Xi_c(1)=\frac{\prod_{i=1}^r(q^i-1)}{\prod_{j=1}^h(q^{d_j}-1)} ,\qquad\mu(c):=\frac{\Xi_c(1)^2}{\#
\auGL_r(\auF)}=\frac{\prod_{i=1}^r(q^i-1)}{q^{r(r-1)/2}\prod_j(q^{d_j}-1)^2}\ \ge\ \frac{1-r/q}{
q^r} .
\end{equation}
The idempotent of $c$ is $e_c=\Xi_c(1)(\#\auGL_r(\auF))^{-1}\sum_g\aubar{\Xi_c(g)}g$, with
coefficient $\mu(c)$ at $1$.

\begin{lemma}[Bad types are rare]\label{pb:bad}
For $i\in\{0,1\}$ let $\mathsf g_i=\sum\bigotimes_{w\mid p_i}e_{c_w}\in\C[G_i]$, the sum over the
families $(c_w)_{w\mid p_i}$ of labels such that no $c_w$ has a root in $\auF_w$ that is a
$q_2$-th root of unity and at least one $c_w$ is irreducible with a root generating $\auF_{p_i^r}^
\times$. Then $\mathsf g_i$ is a central idempotent of $\C[G_i]$, $e_0=\mathsf g_0\mathsf g_1$ in
$\C[\auG]$, and the coefficient of $1-\mathsf g_i$ at $1$ is at most $b^{-8}/2$.
\end{lemma}
\begin{proof}
The $\bigotimes_we_{c_w}$ are the idempotents of distinct irreducible characters of $G_i$, so
$\mathsf g_i$ is a central idempotent; $e_0=\mathsf g_0\mathsf g_1$ is the definition of $e_0$.
Write $q=p_i$ and call a monic polynomial of degree $r$ over $\auF_q$ \emph{allowed} if it is a
label without a root that is a $q_2$-th root of unity, and \emph{primitive} if it is irreducible
with a root of order $q^r-1$. As $r>1$, primitive polynomials are allowed. Let $\mathsf a$ be the
number of allowed and $\mathsf p$ the number of primitive polynomials, and $m_0=(1-r/q)q^{-r}$. The
coefficient of $\mathsf g_i$ at $1$ is $\sum\prod_w\mu(c_w)$ over the good families, which by
\eqref{pb:eq-green} is at least $m_0^f\bigl(\mathsf a^f-(\mathsf a-\mathsf p)^f\bigr)$. So the
coefficient of $1-\mathsf g_i$ at $1$ is at most
\[
1-(m_0\mathsf a)^f+\bigl(m_0(q^r-\mathsf p)\bigr)^f .
\]
A monic polynomial of degree $r\ge2$ that is not allowed is not squarefree ($q^{r-1}$
polynomials), or vanishes at $0$ ($q^{r-1}$), or vanishes at one of at most $q_2$ roots of unity
($q_2q^{r-1}$). So $m_0\mathsf a\ge(1-r/q)(1-(2+q_2)/q)\ge1-\varepsilon$ with $\varepsilon=(1+q_2+
2r)/q$, and $1-(m_0\mathsf a)^f\le f\varepsilon$. The roots of the primitive polynomials are the
generators of $\auF_{q^r}^\times$, $r$ for each polynomial, so $\mathsf p=\varphi(q^r-1)/r$. For an
integer $\mathsf n\ge2$ with $\omega$ prime factors, $\varphi(\mathsf n)/\mathsf n\ge\prod_{j=2}^{
\omega+1}(1-1/j)=1/(\omega+1)\ge1/(2\log\mathsf n+1)$. With $\mathsf n=q^r-1\ge q^r/2$ this gives
$\mathsf p/q^r\ge\theta:=\bigl(2r(2r\log q+1)\bigr)^{-1}$, hence $m_0(q^r-\mathsf p)\le1-\theta\le
e^{-\theta}$. Altogether the coefficient is at most $f\varepsilon+e^{-f\theta}$. By the definition
of the data, $x^{11}\le f\le x^{12}$, $r\le3x^5$, $q_2\le b^{1/100}$, and
$b^{10}\le q\le b^{11}$; with {\cite[(36)]{P46}} this gives $f\varepsilon\le b^{-9}$ and
$f\theta\ge9\log b$, and $2b^{-9}\le b^{-8}/2$. Here $q_2\le b^{1/100}$ holds because
$q_2=(q_0d)^{s+6}$ with $s\le32\log x$ by (D3), $\log(q_0d)\le2\log C_P\le2\log\log\log b$ (as
$C_P\ge q_0+2d$ and $10^{12}C_P\le\log\log b$), and $\log x\le10^{-6}\log x_1\le10^{-6}(\log b)^{1/6}$
by (D2); so $\log q_2\le70\log x\cdot\log\log\log b\le10^{-4}(\log b)^{1/6}\log\log\log b\le
(\log b)/100$.
\end{proof}

\subsection{The sum over the group: \texorpdfstring{{\cite[Proposition~14.32(iii)]{P46}}}{[Row D, Proposition~14.32(iii)]}}\label{pb:ss-iii}

\begin{lemma}\label{pb:basic}
Let $P=c_P\auP$ be as in Theorem~\ref{pb:est}. Then:
\begin{enumerate}
\item[(a)] $P_{st}=P_s\,\pbrpsi(t)$ and $P_{ts}=\pbrpsi(t)\,P_s$ for $t\in\pbR$, $s\in\auG$;
\item[(b)] $\pbHS{P^*P-1}^2\le b^{-12}\auDX$ and $\pbHS{P^*P}^2\le2\auDX$;
\item[(c)] $\beta(s)=\pbHS{P_s}^2$ is of positive type on $\auG$, and $\beta(s)\le2\auDX$ for all
$s$.
\end{enumerate}
\end{lemma}
\begin{proof}
(a) By Lemma~\ref{pb:equiv}, $P_{st}=P^*\varrho(s)\varrho(t)P=P^*\varrho(s)P\pbrpsi(t)$. As
$\pbrpsi(t)$ is unitary,
\[
P_{ts}=(\varrho(t^{-1})P)^*\varrho(s)P=(P\pbrpsi(t^{-1}))^*\varrho(s)P=\pbrpsi(t)P_s .
\] (b) For $s\in\pbR$, by (a) and \eqref{pb:eq-rext}, $P_s-\pbrpsi(s)=(P^*P
-1)\pbrpsi(s)$, which has the norm of $P^*P-1$. As $\pbR\subseteq\Omega_0$, (E1) (proved in Proposition~\ref{pb:reduction} before the present lemma is applied there) gives $\#\pbR
\cdot\pbHS{P^*P-1}^2\le b^{-12}\#\pbR\cdot\auDX$. Then
$\pbHS{P^*P}\le\pbHS{P^*P-1}+\pbHS1\le(b^{-6}+1)\sqrt{\auDX}$. (c) Let $u_a=\varrho(a)PP^*\varrho(a)^{-1}$, in the space of endomorphisms of
$\auAz$ with the Hilbert--Schmidt inner product. Then
\[
\pbip{u_a}{u_b}=\pbtr\bigl(\varrho(a)PP^*\varrho(a^{-1}b)PP^*\varrho(b^{-1})\bigr)=\pbtr\bigl(P_{a^{-1}
b}\,P_{a^{-1}b}^{\,*}\bigr)=\beta(a^{-1}b) .
\]
So $\beta(s)\le\beta(1)=\pbHS{P^*P}^2\le2\auDX$.
\end{proof}

\begin{proposition}[{\cite[Proposition~14.32(iii)]{P46}}]\label{pb:iii}
Let $P$ be as in Theorem~\ref{pb:est}. Then $\sum_{s\in\auG}\pbHS{P^*\varrho(s)P}^2\le(1+9b^{-4})\,
\pbDmain$.
\end{proposition}
\begin{proof}
Let $\beta(s)=\pbHS{P_s}^2$ and, for a set $\mathsf A$ of chosen places above $p_0$ and a family
$\mathsf d$ of frame vectors, $m(\mathsf A,\mathsf d)=\sum_{s\in\pbFix(\mathsf A,\mathsf d)}
\beta(s)$. Put $\Theta=b^{6k}$ and choose $(\mathsf A,\mathsf d)$ such that $\Theta^{\#\mathsf A}
m(\mathsf A,\mathsf d)$ is largest. Since $m(\emptyset,\mathsf d)=\sum_{s\in\auG}\beta(s)$, it
suffices to prove
\begin{equation}\label{pb:eq-max}
m(\mathsf A,\mathsf d)\le(1+9b^{-4})\,\Theta^{-\#\mathsf A}\,\pbDmain .
\end{equation}
If $\mathsf A$ is the set of all chosen places above $p_0$, this is (E2). Otherwise let $w\notin
\mathsf A$ be a chosen place above $p_0$, $\Gamma_0=\pbFix(\mathsf A,\mathsf d)$ and $m=m(\mathsf A,
\mathsf d)$. For a frame vector $\mathsf d'$ at $w$, the group $\Gamma_0\cap\pbFix(\mathsf d')$ is
$\pbFix(\mathsf A\cup\{w\},\cdot)$ for the family with $\mathsf d'$ at $w$, so by the choice of
$(\mathsf A,\mathsf d)$
\begin{equation}\label{pb:eq-stab}
\sum_{s\in\Gamma_0\cap\pbFix(\mathsf d')}\beta(s)\le m/\Theta .
\end{equation}
Let $1\le l\le\varrho$. Apply Lemma~\ref{pb:singular} to $\mathsf G=\Gamma_0$, $\auF=\auF_w$ (so
$q=p_0$), $\pi(s)=s_w$, the matrix $A_0$ with columns $f_0,\dots,f_{l-1}$ and $A_1=\mathsf JA_0$,
and $Y=m/\Theta$. For $v\ne0$ the vector $A_0v$ is a non-zero vector of the span of the frame;
$\pi(s)A_0v=A_0v$ means $s\in\pbFix(\mathsf d')$ for $\mathsf d'=A_0v\in F_w$, and ${}^t\pi(s)^{-1}
A_1v=A_1v$ means $s_w^\vee A_0v=A_0v$, that is $s\in\pbFix(\mathsf d')$ for $\mathsf d'=A_0v\in
F_{\aubar w}$. So the hypotheses hold by \eqref{pb:eq-stab}, and since ${}^tA_1\pi(s)A_0$ is the
leading block of size $l$ of $Z_w(s_w)$,
\[
p_0\sum_{s\in\Gamma_0,\ \delta_l(s_w)=0}\beta(s)\le2m+p_0^{\,l+r/2}\,m/\Theta\le3m ,
\]
because $p_0^{\,l+r/2}\le b^{11(\varrho+r/2)}\le b^{6k}$ (as $11(4\varrho+1)\le k$). An element of
$\Gamma_0$ that is not open has $\delta_l(s_w)=0$ for some $1\le l\le\varrho$. So, with
$p_0\ge b^{10}$ and $3\varrho\le b^6$,
\[
\sum_{s\in\Gamma_0,\ s\text{ not open}}\beta(s)\le\varrho\cdot3m\,b^{-10}\le b^{-4}m .
\]
With (E3) for the open elements, $m\le(1+4b^{-4})\Theta^{-\#\mathsf A}\pbDmain+b^{-4}m$. As $(1+9
\epsilon)(1-\epsilon)\ge1+4\epsilon$ for $0\le\epsilon\le\frac1{10}$, \eqref{pb:eq-max} follows.
\end{proof}

\subsection{The good part: \texorpdfstring{{\cite[Proposition~14.32(ii)]{P46}}}{[Row D, Proposition~14.32(ii)]}}\label{pb:ss-ii}

\begin{lemma}[One prime]\label{pb:oneprime}
Let $P$ be as in Theorem~\ref{pb:est}, $i\in\{0,1\}$, and let $x$ be a central idempotent of
$\C[G_i]$ with $x(1)\le b^{-8}/2$. Then $\pbHS{\varrho(x)P}^2\le2b^{-4}\auDX$.
\end{lemma}
\begin{proof}
Regard $G_i$ as a subgroup of $\auG$. By Lemma~\ref{pb:central}(b), $\varrho(x)$ is an orthogonal
projection, so
\[
\pbHS{\varrho(x)P}^2=\pbtr\bigl(P^*\varrho(x)P\bigr)=\sum_{t\in G_i}x(t)\pbtr\pbrpsi(t)+\pbtr P'_x,
\qquad P'_x=\sum_{t\in G_i}x(t)\bigl(P_t-\pbrpsi(t)\bigr).
\]
\emph{The first term.} The elements of $G_i$ have trivial components above the other prime, so by
Lemma~\ref{pb:trace} the first term is $\auDX\sum_{t\in G_i}x(t)\tilde\psi(t)$. For
$g\in\auG$ and $t\in G_i$, $gtg^{-1}=g_itg_i^{-1}\in G_i$, and $x$ is a class function on $G_i$; so
the sum does not change if $\tilde\psi(t)$ is replaced by $\tilde\psi(gtg^{-1})$, hence by
$\psi^{\pbnat}(t)$. By the Cauchy--Schwarz inequality, Lemma~\ref{pb:central}(a) and (E4),
\[
\Bigl|\sum_{t\in G_i}x(t)\psi^{\pbnat}(t)\Bigr|^2\le x(1)\sum_{t\in\auG}|\psi^{\pbnat}(t)|^2
\le b^{-8} .
\]
\emph{The second term.} Put $P'_s=P_s-\pbrpsi(s)$. Let $i'$ be the other index. For $s\in\pbR$ and
$t\in G_i$ we have $st\in\Omega_{i'}$, and $P'_{st}=\pbrpsi(s)P'_t$ by Lemma~\ref{pb:basic}(a) and
\eqref{pb:eq-rext}. Let $F(s)=P'_s$ for $s\in\Omega_{i'}$ and $F(s)=0$ otherwise. Then, for
$s\in\pbR$, $\sum_tx(t)F(st)=\pbrpsi(s)P'_x$ has the norm of $P'_x$, and by
Lemma~\ref{pb:central}(c) and (E1)
\[
\#\pbR\cdot\pbHS{P'_x}^2\le\sum_{s\in\auG}\Bigl\|\sum_{t\in G_i}x(t)F(st)\Bigr\|^2\le\sum_{s\in
\Omega_{i'}}\pbHS{P'_s}^2\le b^{-12}\,\#\pbR\cdot\auDX .
\]
So $|\pbtr P'_x|\le\pbHS1\,\pbHS{P'_x}\le b^{-6}\auDX$. Together, $\pbHS{\varrho(x)P}^2\le(b^{-4}+
b^{-6})\auDX$.
\end{proof}

\begin{proposition}[{\cite[Proposition~14.32(ii)]{P46}}]\label{pb:ii}
Let $P$ be as in Theorem~\ref{pb:est}. Then $\pbHS{\varrho(e_0)P\,\pbpr_{\mathsf U}}^2\ge\dim
\mathsf U-5b^{-4}\auDX$ for every subspace $\mathsf U\subseteq\auS$. In particular $\audT>0$.
\end{proposition}
\begin{proof}
Let $\epsilon_i=\varrho(\mathsf g_i)$, with $\mathsf g_i$ as in Lemma~\ref{pb:bad}; these are
commuting orthogonal projections and $\varrho(e_0)=\epsilon_0\epsilon_1$. By Lemmas~\ref{pb:bad}
and~\ref{pb:oneprime} applied to $x=1-\mathsf g_i$, $\pbHS{(1-\epsilon_i)P}^2\le2b^{-4}\auDX$. Since
$1-\epsilon_0\epsilon_1=(1-\epsilon_0)+\epsilon_0(1-\epsilon_1)$ with orthogonal ranges,
$\pbHS{(1-\varrho(e_0))P}^2\le\pbHS{(1-\epsilon_0)P}^2+\pbHS{(1-\epsilon_1)P}^2\le4b^{-4}\auDX$. Now
\[
\pbHS{\varrho(e_0)P\pbpr_{\mathsf U}}^2=\pbHS{P\pbpr_{\mathsf U}}^2-\pbHS{(1-\varrho(e_0))P\pbpr_{
\mathsf U}}^2\ge\pbHS{P\pbpr_{\mathsf U}}^2-4b^{-4}\auDX ,
\]
and $\pbHS{P\pbpr_{\mathsf U}}^2=\pbtr(\pbpr_{\mathsf U}P^*P)=\dim\mathsf U+\pbtr\bigl(\pbpr_{\mathsf
U}(P^*P-1)\bigr)\ge\dim\mathsf U-\sqrt{\dim\mathsf U}\,\pbHS{P^*P-1}\ge\dim\mathsf U-b^{-6}\auDX$ by
Lemma~\ref{pb:basic}(b). For $\mathsf U=\auS$ this gives $\varrho(e_0)P\ne0$, so $\auT\ne0$.
\end{proof}

\subsection{The dimension of the target, and the Spanning Theorem}\label{pb:ss-thmB}

\begin{proposition}[{\cite[Proposition~14.32(i)]{P46}}]\label{pb:i}
If {\cite[Proposition~14.36]{P46}} holds for $j=m$, then $\audT\le\#\auG\cdot\pbdeg\cdot\pbmass_r$.
\end{proposition}
\begin{proof}
By {\cite[Proposition~14.36]{P46}}, $\audT=\dim e_0\mathcal A(W_m)\le\dim\mathcal A(W_m)$. The space
$W_m$ is the positive definite standard space of rank $r$ and $\auU(W_m)$ is compact at the
infinite places, where the discrete series with parameter $\varsigma_\tau$ is the irreducible
representation $\mathsf F_\tau$ of dimension $\pbdeg_\tau$. So $\mathcal A(W_m)$ is the space of
algebraic modular forms on $\auU(W_m)$ with values in $\bigotimes_\tau\mathsf F_\tau$ and level
$J_M$. This level is free ({\cite[Lemma~14.5]{P46}}, as $r\le2n_m$), so by \eqref{pb:eq-classes} its
dimension is $[K^\circ_r:J_M]\cdot\pbmass_r\cdot\pbdeg=\#\auG\cdot\pbmass_r\cdot\pbdeg$.
\end{proof}

\subsubsection*{The lattice of $Y_0$}
In this section and in all the sections on {\cite[Theorem~14.14]{P46}}, $L$ is the lattice of $Y_0$ with the
basis $e_a,e_a^*$ $(1\le a\le\varrho)$, $x_1,\dots,x_{k-1}$, $x_k/\aubar{\pi_N}+z/N$, $z$ (proof of
Lemma~\ref{pb:KR}), and $K$, $J$, $\auAz$, $\varrho$, $\auP$ are formed with it. The lattice of $Y_0$
with which the lower chain is defined is the lattice $\mathsf M_0$ of {\cite[(35)]{P49}}; in the
notation there the string is $f_{a-1}=e_a$, $f_\varrho=z$, $f_{2\varrho+1-a}=e_a^*$ $(1\le a\le\varrho)$
and $y=x_k$, the group $R_0$ and the character $\psi_0=\psi_U$ are those of the present section
({\cite[\S5.1.3]{P49}}), and $\mathsf M_0$ is spanned by $x_1,\dots,x_{k-1}$, $x_k/\aubar{\pi_N}+u_0z/N$,
$u_0z$, $u_0e_a$, $u_0e_a^*$, where $u_0$ is a finite idele with $u_0\aubar{u_0}=1$ that is not a unit
(at the chosen places, and at the places above $q_1$ where $\alpha/\aubar\alpha$ is not a unit). The
two lattices have different stabilisers, so the spaces of forms are different spaces of functions.
The next lemma shows that all the statements of {\cite[Proposition~14.32]{P46}} pass from one to the
other; it is applied with $\mathsf u=u_0$.

\begin{lemma}[Change of the lattice]\label{pb:lattice}
Let $\mathsf u$ be a finite idele of $E$ with $\mathsf u\aubar{\mathsf u}=1$, let $\mathsf h_{\mathsf u}$
be the element of the centre of $H(\auA_{E^+,f})$ that acts on $X_0$ by the scalar $\mathsf u^{-1}$
(and trivially on $A$, on $z$ and on the $e_a^*$), and let $L_{\mathsf u}$ be the lattice of $Y_0$
spanned over $\widehat\auO_E$ by $x_1,\dots,x_{k-1}$, $x_k/\aubar{\pi_N}+\mathsf uz/N$, $\mathsf uz$,
$\mathsf ue_a$, $\mathsf ue_a^*$ $(1\le a\le\varrho)$.
\begin{enumerate}
\item[(a)] $L_{\mathsf u}=\mathsf u\cdot\mathsf h_{\mathsf u}L$. So $L_{\mathsf u}$ is unimodular, its
stabiliser is $K_{\mathsf u}=\mathsf h_{\mathsf u}K\mathsf h_{\mathsf u}^{-1}$, the principal congruence
subgroup of level $\pi_M$ in it is $J_{\mathsf u}=\mathsf h_{\mathsf u}J\mathsf h_{\mathsf u}^{-1}$, and
the reduction of $k\in K_{\mathsf u}$ at the chosen places, in the basis of $L_{\mathsf u}$ just given, is
the image of $\mathsf h_{\mathsf u}^{-1}k\,\mathsf h_{\mathsf u}\in K$ in $\auG$.
\item[(b)] Let $\auAz^{\mathsf u}$ be the space $\auAz$ of {\cite[\S14.3]{P46}} formed with
$J_{\mathsf u}$ in place of $J$, with the action $\varrho_{\mathsf u}$ of $\auG=K_{\mathsf u}/J_{\mathsf u}$
by right translation, and let $\auP_{\mathsf u}\colon\auS\to\auAz^{\mathsf u}$ be the Poincar\'e map
defined by the same $g_0$, $\mathsf t$, $\psi_U$ and Haar measure of $[R_0]$ as $\auP$. Then
$(\mathsf j_{\mathsf u}y)(g)=y(g\,\mathsf h_{\mathsf u}^{-1})$ defines a unitary isomorphism
$\mathsf j_{\mathsf u}\colon\auAz^{\mathsf u}\to\auAz$ with $\mathsf j_{\mathsf u}\varrho_{\mathsf u}(s)=
\varrho(s)\mathsf j_{\mathsf u}$ for $s\in\auG$, the map $(R_{\mathsf u}\varphi)(h)=\varphi(h\,\mathsf
h_{\mathsf u}^{-1})$ is a unitary automorphism of $\auS$, and
\[
\auP_{\mathsf u}=\mathsf j_{\mathsf u}^{-1}\circ\auP\circ R_{\mathsf u} .
\]
\item[(c)] If {\rm(i)}, {\rm(ii)}, {\rm(iii)} of {\cite[Proposition~14.32]{P46}} hold for $P=c_P\auP$,
the space $\auAz$ and $\auT=\varrho(e_0)\auAz$, then they hold, with the same inner product on $\auS$
and the same $c_P$, for $P_{\mathsf u}=c_P\auP_{\mathsf u}$, the space $\auAz^{\mathsf u}$ and
$\auT^{\mathsf u}=\varrho_{\mathsf u}(e_0)\auAz^{\mathsf u}$; and $\dim\auT^{\mathsf u}=\dim\auT$.
\end{enumerate}
Another basis of $L_{\mathsf u}$ changes the identification of $K_{\mathsf u}/J_{\mathsf u}$ with $\auG$ by
an inner automorphism of $\auG$; the central idempotent $e_0$, the space $\auT^{\mathsf u}$ and the
statements {\rm(i)--(iii)} do not depend on it.
\end{lemma}
\begin{proof}
(a) $\mathsf h_{\mathsf u}L$ is spanned by $\mathsf u^{-1}x_i$ $(i<k)$, $\mathsf u^{-1}x_k/\aubar{\pi_N}+
z/N$, $z$, $e_a$, $e_a^*$; multiplying by $\mathsf u$ gives the basis of $L_{\mathsf u}$. As
$\mathsf u\aubar{\mathsf u}=1$, the scalar $\mathsf u$ and $\mathsf h_{\mathsf u}$ are isometries of
$Y_0\otimes\auA_{E,f}$, so $L_{\mathsf u}$ is unimodular; a scalar does not change the stabiliser of a
lattice, so $K_{\mathsf u}$ is the stabiliser of $\mathsf h_{\mathsf u}L$, which is $\mathsf h_{\mathsf u}
K\mathsf h_{\mathsf u}^{-1}$, and the congruence subgroups correspond. The matrix of $k$ in the basis
$\mathsf u\,\mathsf h_{\mathsf u}(\mathsf b)$, $\mathsf b$ the basis of $L$, is the matrix of $\mathsf
h_{\mathsf u}^{-1}k\,\mathsf h_{\mathsf u}$ in $\mathsf b$, the scalar $\mathsf u$ being central.

(b) If $y$ is right invariant under $J_{\mathsf u}$, then for $j\in J$, $(\mathsf j_{\mathsf u}y)(gj)=
y(g\,\mathsf h_{\mathsf u}^{-1}\cdot\mathsf h_{\mathsf u}j\mathsf h_{\mathsf u}^{-1})=(\mathsf j_{\mathsf u}y)
(g)$. Right translation by a finite adelic element preserves the invariance under $\Gamma$ on the
left, the boundedness, the equation $\int e_\infty(a)y(ga)\,da=y(g)$ (the two translations commute)
and the inner product $\int_{\Gamma\backslash G(\auA)}\aubar yy'$; so $\mathsf j_{\mathsf u}$ is a
unitary isomorphism, with inverse $y\mapsto y(\,\cdot\,\mathsf h_{\mathsf u})$. For $k\in K_{\mathsf u}$,
$(\mathsf j_{\mathsf u}\varrho_{\mathsf u}(k)y)(g)=y(g\,\mathsf h_{\mathsf u}^{-1}k)=(\varrho(\mathsf h_{\mathsf
u}^{-1}k\,\mathsf h_{\mathsf u})\mathsf j_{\mathsf u}y)(g)$, which by (a) is the stated intertwining.
The element $\mathsf h_{\mathsf u}$ is central in $H(\auA_{E^+,f})$, so $R_{\mathsf u}\varphi$ is again
right invariant under $J_X$ and $R_{\mathsf u}$ maps $\auS$ to itself; $c\mapsto c\,\mathsf h_{\mathsf
u}^{-1}$ permutes the classes, and if $h_c\mathsf h_{\mathsf u}^{-1}=\eta h_{c'}j$ with $\eta\in H(E^+)$,
$j\in J_X$, then $(R_{\mathsf u}\varphi)(h_c)=\eta\,\varphi(h_{c'})$ with $\eta$ acting unitarily on
$W_{\lambda^X}$; so $R_{\mathsf u}$ is unitary. Now let $y\in\auAz^{\mathsf u}$ and $\varphi\in\auS$. The
element $\mathsf h_{\mathsf u}$ is a finite adele and $g_0$ is real, so they commute, and
$y(rg_0)=(\mathsf j_{\mathsf u}y)(rg_0\mathsf h_{\mathsf u})=(\mathsf j_{\mathsf u}y)(r\,\mathsf h_{\mathsf u}g_0)$.
As $\mathsf h_{\mathsf u}\in R_0(\auA)$ and the measure of $[R_0]$ is invariant under right translations
(proof of Lemma~\ref{pb:equiv}), the substitution $r\mapsto r\,\mathsf h_{\mathsf u}^{-1}$ gives
\[
\pbip y{\auP_{\mathsf u}\varphi}=\int_{[R_0]}\mathcal W_\varphi(r)\,\aubar{y(rg_0)}\,dr=\int_{[R_0]}
\mathcal W_\varphi(r\,\mathsf h_{\mathsf u}^{-1})\,\aubar{(\mathsf j_{\mathsf u}y)(rg_0)}\,dr ,
\]
and for $r=uh$, $\mathcal W_\varphi(uh\,\mathsf h_{\mathsf u}^{-1})=\psi_U(u)\,\mathsf t\bigl(h_\infty^{-1}
\varphi(h_f\mathsf h_{\mathsf u}^{-1})\bigr)=\mathcal W_{R_{\mathsf u}\varphi}(uh)$. Hence $\pbip y{\auP_{\mathsf
u}\varphi}=\pbip{\mathsf j_{\mathsf u}y}{\auP R_{\mathsf u}\varphi}=\pbip y{\mathsf j_{\mathsf u}^{-1}\auP
R_{\mathsf u}\varphi}$.

(c) By (b), $\mathsf j_{\mathsf u}$ carries $\auT^{\mathsf u}$ onto $\auT$, which gives the equality of
dimensions and (i). For a subspace $\mathsf U\subseteq\auS$, $\varrho_{\mathsf u}(e_0)P_{\mathsf u}
\pbpr_{\mathsf U}=\mathsf j_{\mathsf u}^{-1}\,\varrho(e_0)P\,\pbpr_{R_{\mathsf u}\mathsf U}\,R_{\mathsf u}$,
which has the same Hilbert--Schmidt norm as $\varrho(e_0)P\,\pbpr_{R_{\mathsf u}\mathsf U}$, and
$\dim R_{\mathsf u}\mathsf U=\dim\mathsf U$; this gives (ii). Finally $P_{\mathsf u}^*\varrho_{\mathsf u}(s)
P_{\mathsf u}=R_{\mathsf u}^*\,P^*\varrho(s)P\,R_{\mathsf u}$ has the norm of $P^*\varrho(s)P$, which gives
(iii). The last assertion of the lemma is clear, the sum in (iii) being over all of $\auG$ and
$e_0$ being central.
\end{proof}

\noindent For $\mathsf u=u_0$ the lattice $L_{u_0}$ is $\mathsf M_0$. (Compare the remark in
\textup{(B4)} of Definition~\ref{pb:bessel}: a translation of the source by a central element of
$H(\auA_{E^+,f})$ is all that distinguishes the two Poincar\'e maps.) The scalar $\mathsf u$ on the
string alone, $d=\mathsf u\cdot\mathsf h_{\mathsf u}$, also carries $L$ to $L_{\mathsf u}$; it is the product
of an element of the centre of $G(\auA_{E^+,f})$ and of $\mathsf h_{\mathsf u}$, and no rational element
with the ideal of $\mathsf u$ is needed.

\begin{theorem}\label{pb:thmB}
Assume Theorem~\ref{pb:est} and {\cite[Proposition~14.36]{P46}} (for $j=m$). Then
{\cite[Proposition~14.32]{P46}} holds for the pair $(g_0,\mathsf t)$ and the scalar $c_P$ of
Theorem~\ref{pb:est}, the character \eqref{pb:eq-psi}, the inner product of
\S\ref{pb:ss-setting} on $\auS$, and $\deg=\pbdeg$; and {\cite[Theorem~14.14]{P46}} holds for the map
$\auP$ defined by these choices. Both statements hold for the lattice $L$ of this section and
for the lattice $\mathsf M_0$ of {\cite[(35)]{P49}}, with the same $(g_0,\mathsf t)$, $c_P$,
$\psi_U$ and inner product.
\end{theorem}
\begin{proof}
For the lattice $L$: parts (i), (ii), (iii) of {\cite[Proposition~14.32]{P46}} are Propositions~\ref{pb:i}, \ref{pb:ii}
and~\ref{pb:iii}, and $\audT>0$ by Proposition~\ref{pb:ii}; recall $\pbDmain=\auDX^2/(\pbdeg\cdot
\pbmass_r)$. For the lattice $\mathsf M_0=L_{u_0}$ they follow by Lemma~\ref{pb:lattice}(c), and
$\dim\auT^{u_0}=\audT>0$. {\cite[Theorem~14.14]{P46}} follows in either case by the argument of {\cite[\S14.7]{P46}}, which uses
{\cite[Proposition~14.32]{P46}} through {\cite[Proposition~14.34]{P46}} and
{\cite[Lemmas~14.29, 14.30 and~14.31]{P46}}; these are proved there and do not
depend on the choice of $(g_0,\mathsf t)$ or of $\psi_U$.
\end{proof}

The error $\frac32(8b^2)^{-1}$ in {\cite[Proposition~14.34]{P46}}, hence the numbers $4/b^2$ and
$5\log(1+2b^2)$ of {\cite[Theorem~14.14]{P46}}, are thus far from what the method gives: (ii) and (iii)
hold with $5b^{-4}$ and $9b^{-4}$, and these come from two places only, the factor $1+4b^{-4}$ of
(E3) and the weight $b^{-4}$ of the elements that are not open (the union over $\varrho$ minors of
sets of weight $3b^{-10}$), on one side, and the share $b^{-8}/2$ of the bad types on the other.
In the part of the argument that is proved here, the rank enters only through the inequalities
$11(4\varrho+1)\le k$ and $3\varrho\le b^6$ (Proposition~\ref{pb:iii}) and $1<r<p_i$,
$18\,r(22\,r\log b+1)\log b\le f$ (Lemma~\ref{pb:bad}).

\subsection{Second layer: the Poincar\'e series and the inequalities (E1)--(E3)}\label{pb:ss-layer2}

\subsubsection*{Bessel functions}
Let $H_\infty=\prod_\tau H_\tau$, $U_\infty=\prod_\tau U_\tau$, $R_\infty=U_\infty\rtimes H_\infty$.
Fix Haar measures on $G_\tau$ and $U_\tau$, hence invariant measures on $U_\tau\backslash G_\tau$;
they are invariant under left multiplication by $H_\tau$, which normalises $U_\tau$ and is
compact. Let $\pbW=\bigotimes_\tau\pbW_\tau$ be the hermitian space dual to $W_{\lambda^X}$, with
the action $\tau_\infty$ of $H_\infty$ and an orthonormal basis indexed by a set $I$, $\#I=d$, a
product of orthonormal bases of the $\pbW_\tau$. We write $\pbip w\xi=\sum_{i\in I}w_i\xi_i$ for
the pairing of $w\in\pbW$ and $\xi\in W_{\lambda^X}$, in dual bases; $\pbip{\tau_\infty(h)w}{h\xi}=
\pbip w\xi$.

\begin{definition}[Hypothesis on the Bessel functions]\label{pb:bessel}
We say that \emph{Bessel functions exist} for the data if there are a point $g_0\in G_\infty$, a
linear form $\mathsf t$ on $W_{\lambda^X}$, a constant $c_W>0$ and continuous functions
$W_\tau,W_\tau^\flat\colon G_\tau\to\pbW_\tau$ $(\tau\in\auPl)$ with the following properties, where
$W(y)=\bigotimes_\tau W_\tau(y_\tau)$ and $W^\flat(y)=\bigotimes_\tau W^\flat_\tau(y_\tau)$ for
$y\in G_\infty$.
\begin{enumerate}
\item[(B1)] $W_\tau(uhy)=\psi_\tau(u)\,\tau_\tau(h)\,W_\tau(y)$ for $u\in U_\tau$, $h\in H_\tau$,
where $\psi_\tau$ is the component of $\psi_\infty$; and the same for $W_\tau^\flat$.
\item[(B2)] $\|W_\tau\|$ is integrable and square-integrable on $U_\tau\backslash G_\tau$, and
$\int_{U_\tau\backslash G_\tau}\|W_\tau\|^2=\dim\pbW_\tau$.
\item[(B3)] For every class $c$ and $i\in I$ the series $\sum_{\gamma\in U_c\backslash\Gamma}
\bigl|W_i(\gamma_\infty g_\infty)\bigr|\,\pbone_{h_cJ}(\gamma g_f)$ is bounded on $G(\auA)$, where
$U_c=U(E^+)\cap h_cJh_c^{-1}$; and the function $P_{c,i}(g)=\sum_{\gamma\in U_c\backslash\Gamma}W_i(
\gamma_\infty g_\infty)\pbone_{h_cJ}(\gamma g_f)$ belongs to $\auAz$.
\item[(B4)] For every $\varphi\in\auS$, as functions on $G(\auA)$,
\[
\auP\varphi(g)=c_W\sum_{c\in\mathcal C}\ \sum_{\gamma\in U_c\backslash\Gamma}\pbip{W(\gamma_\infty
g_\infty)}{\varphi(h_c)}\,\pbone_{h_cJ}(\gamma g_f),
\]
where $\auP$ is defined by $g_0$, $\mathsf t$ and \eqref{pb:eq-psi}. (Replacing $\varphi(h_c)$ by
the value at $h_c$ of the translate of $\varphi$ by a fixed central element of $H(\auA_f)$ composes
$\auP$ with a unitary automorphism of $\auS$ that commutes with $\pbrpsi$, and changes nothing below.)
\item[(B5)] $W^\flat_\tau=\chi_\tau W_\tau$ with
$\chi_\tau(y)=\max\bigl(0,\min(1,2-2\ell_\tau(y)/b^{6n+2})\bigr)$, where
$\ell_\tau(y)=\inf_{s\in R_\tau}(\|sy\|_\tau+1)$ and $\|\cdot\|_\tau$ is the operator norm of
Definition~\ref{pbe:length}. So $\chi_\tau=1$ where $\ell_\tau\le\frac12b^{6n+2}$, $\chi_\tau=0$
where $\ell_\tau\ge b^{6n+2}$, and $\chi_\tau$ is continuous, $0\le\chi_\tau\le1$, and left
invariant under $R_\tau$; in particular the support of $W^\flat_\tau$ is compact modulo $R_\tau$.
\item[(B5$'$)] For each $\tau$ there is a compact set $C_\tau\subset G_\tau$ with
$\operatorname{supp}W^\flat_\tau\subseteq U_\tau C_\tau$. (This follows from (B5): $\chi_\tau(y)\ne0$ gives
$s=uh\in R_\tau$ with $\|sy\|_\tau<b^{6n+2}$, so $y\in h^{-1}u^{-1}B=u'h^{-1}B$ with $u'=h^{-1}u^{-1}h\in
U_\tau$ and $B$ the closed ball of radius $b^{6n+2}$ in $G_\tau$; take $C_\tau=H_\tau B$, compact as
$H_\tau$ is compact.)
\end{enumerate}
\end{definition}

The terms of the series in (B3), (B4) are well defined: by (B1), $W(uy)=\psi_\infty(u)W(y)$, and
$\psi_\infty(u)=\aubar{\psi_f(u)}=1$ for $u\in U_c$, because $\psi_f$ is invariant under
conjugation by $h_c\in H(\auA_f)$ and trivial on $U(\auA_f)\cap J$. By Theorem~\ref{pbe:bessel}
(with Theorem~\ref{pbe:ds}, proved in \S\ref{pds:sec}, {\cite[\S3.2]{P49}} and~\S\ref{pbp:sec}) Bessel functions
exist: $g_0=(x_\tau)_\tau$, $\mathsf t=\bigotimes_\tau v_\tau$, and $W_\tau$ is, up to a positive constant,
$y\mapsto\int_{R_\tau}c_\tau(x_\tau^{-1}sy)\,(\psi_\tau\otimes\tau_\tau)(s)^{-1}v_\tau\,ds$, where
$c_\tau=\langle v_0,\pi_\tau(\cdot)v_0\rangle$ is the coefficient of a unit highest weight vector of the
lowest $K$-type of the discrete series (Proposition~\ref{pbe:local}); (B4) is the unfolding of the Bessel
period against the reproducing kernel, (B2) and (B3) rest on the integrability of the discrete
series with a large gap and on the lower bound \textup{(DS4)} for the period, and $\chi_\tau$ is a cutoff
in the size of $y$ modulo $R_\tau$ at the radius $b^{6n+2}$.

\begin{lemma}\label{pb:ortho}
Assume \textup{(B1)}, \textup{(B2)}. Then $\int_{U_\infty\backslash G_\infty}W_i\,\aubar{W_{i'}}=
\delta_{ii'}$ for $i,i'\in I$. Moreover $\nu:=\pbvol(U_c\backslash U_\infty)$ is finite and does
not depend on $c$.
\end{lemma}
\begin{proof}
The hermitian matrix $\mathsf M_\tau=\int_{U_\tau\backslash G_\tau}W_\tau(y)W_\tau(y)^*\,dy$
satisfies $\tau_\tau(h)\mathsf M_\tau\tau_\tau(h)^*=\mathsf M_\tau$ for $h\in H_\tau$, by (B1) and
the invariance of the measure under $y\mapsto hy$. As $\pbW_\tau$ is irreducible, $\mathsf M_\tau$
is a scalar, and its trace is $\dim\pbW_\tau$ by (B2); so $\mathsf M_\tau=1$. Take the product
over $\tau$. For the second claim: $U$ is a split unipotent group, so $U(E^+)$ is discrete and
cocompact in $U(\auA)$ and $U(E^+)U_\infty$ is dense in $U(\auA)$ (strong approximation for the
additive group, and induction on a composition series). Hence, for the open compact subgroup
$\mathsf K_c=U(\auA_f)\cap h_cJh_c^{-1}$, $U_c\backslash U_\infty\cong U(E^+)\backslash U(\auA)/
\mathsf K_c$ and $\nu\cdot\pbvol(\mathsf K_c)=\pbvol(U(E^+)\backslash U(\auA))$; and $\pbvol(
\mathsf K_c)=\pbvol(U(\auA_f)\cap J)$ because conjugation by $h_c$ preserves the Haar measure of
$U(\auA_f)$ (proof of Lemma~\ref{pb:equiv}).
\end{proof}

The explicit value of $\nu$ enters the main term (\S\ref{pbm:ss-inputs}) and the open cell
(\S\ref{pbs:ss-open}). Put $\pbvol(\widehat\auO_E)=(2/\aurd)^f$ and, with $s_{\tau,a}=\kappa_\tau^{\varrho-a}/
\sqrt N$ as in \eqref{pbe:eq-scale},
\[
\nu_*=\frac{M^{f\varrho(2k+2\varrho+1)}}{N^{kf}}\cdot\frac{\aurd^{\,f\varrho/2}}{\pbvol(\widehat\auO_E)^{\varrho(\varrho+k)}\,
\Lambda},\qquad\Lambda=\prod_\tau\prod_{a=0}^{\varrho}s_{\tau,a}^{\,2(k+2(\varrho-a))}.
\]

\begin{lemma}[The covolume]\label{pb:covol}
For the Haar measure $du$ of \textup{(Cv3)}, $\nu=\nu_*$.
\end{lemma}
\begin{proof}
This is Theorem~\ref{pbv:covol} in \S\ref{pbv:sec}: in the coordinates of (Cv3) a product of fundamental
parallelotopes of $2\varrho$ explicit lattices is a fundamental domain for $U_c$, and the product of their
covolumes is $\nu_*$. Theorem~\ref{pbm:L2} and Lemma~\ref{pbs:open-arith} use $\nu=\nu_*$ (the three forms of
the constant agree, \S\ref{pbv:ss-uses}).
\end{proof}

\subsubsection*{The Petersson formula}
For classes $c,c'$ and $s\in\auG$ the \emph{terms over} $(c,s,c')$ are the elements of
\[
\Gamma_{c,s,c'}=\{\gamma\in\Gamma:\ h_c^{-1}\gamma h_{c'}\in K\text{ and its image in }\auG\text{ is
}s\}.
\]
This set is stable under $\gamma\mapsto u\gamma u'$ for $u\in U_c$, $u'\in U_{c'}$, because $J$ is
normal in $K$. A subset $T\subseteq\Gamma$ is \emph{stable} if $U(E^+)\,T\,U(E^+)=T$. For a stable
$T$ and $\flat\in\{\ ,\flat\}$ the \emph{amplitude sum} is the matrix, indexed by $I\times I$,
\begin{equation}\label{pb:eq-amp}
A^{\flat}_T(c,s,c')=\sum_{\gamma\in U_c\backslash(T\cap\Gamma_{c,s,c'})/U_{c'}}\ \int_{(U_c\cap\gamma
U_{c'}\gamma^{-1})\backslash G_\infty}W^{\flat}(y)\,W^{\flat}(\gamma^{-1}y)^*\,dy ,
\end{equation}
whenever the sum of the integrals of the norms converges. The integrand is invariant under
$U_c\cap\gamma U_{c'}\gamma^{-1}$ and the integral depends only on the double coset of $\gamma$,
since $\psi_\infty$ is trivial on $U_c$ and on $U_{c'}$. We write $A_T$ for the sum formed with
$W$ and $A_T^\flat$ for the one formed with $W^\flat$.

\begin{proposition}[Petersson formula]\label{pb:unfold}
Assume \textup{(B1)--(B4)}, and let $P=(\pbvol(J)\nu)^{-1/2}c_W^{-1}\,\auP$. Then the sums
\eqref{pb:eq-amp} for $W$ converge absolutely for every stable $T$, and for $\varphi,\varphi'\in\auS$
and $s\in\auG$
\[
\pbip{P\varphi}{\varrho(s)P\varphi'}=\frac1\nu\sum_{c,c'\in\mathcal C}\ \sum_{i,i'\in I}\aubar{\varphi(
h_c)_i}\,\varphi'(h_{c'})_{i'}\,\aubar{A_\Gamma(c,s,c')_{i,i'}} .
\]
Consequently $\pbHS{P_s}^2=\nu^{-2}\sum_{c,c'}\|A_\Gamma(c,s,c')\|^2$, with the Euclidean norm of
the matrices.
\end{proposition}
\begin{proof}
Let $F_{c,i}(g)=W_i(g_\infty)\pbone_{h_cJ}(g_f)$, so that $P_{c,i}(g)=\sum_{\gamma\in U_c\backslash
\Gamma}F_{c,i}(\gamma g)$ and $\auP\varphi=c_W\sum_{c,i}\varphi(h_c)_iP_{c,i}$ by (B4). By (B3) the
series converges absolutely, $|P_{c,i}|\le\sum_\gamma|F_{c,i}(\gamma g)|\le\mathsf m$ for a
constant $\mathsf m$, and $P_{c,i}\in\auAz$. Let $k\in K$ be a lift of $s$. Unfolding the first
series,
\[
\pbip{P_{c,i}}{\varrho(k)P_{c',i'}}=\int_{U_c\backslash G(\auA)}\aubar{F_{c,i}(g)}\,P_{c',i'}(gk)\,dg
=\pbvol(J)\int_{U_c\backslash G_\infty}\aubar{W_i(y)}\,P_{c',i'}(y,h_ck)\,dy ;
\]
this is justified because $\int_{U_c\backslash G(\auA)}|F_{c,i}|\,\mathsf m\,dg=\mathsf m\,\pbvol(J)
\,\nu\int_{U_\infty\backslash G_\infty}|W_i|<\infty$ by (B2), and we used that $g\mapsto P_{c',i'}
(gk)$ is right invariant under $J$. Next, $\gamma h_ck\in h_{c'}J$ if and only if $\gamma^{-1}\in
\Gamma_{c,s,c'}$, so $P_{c',i'}(y,h_ck)=\sum_{\gamma\in\Gamma_{c,s,c'}/U_{c'}}W_{i'}(\gamma^{-1}y)$,
and
\[
\int_{U_c\backslash G_\infty}\aubar{W_i(y)}\sum_{\gamma\in\Gamma_{c,s,c'}/U_{c'}}W_{i'}(\gamma^{-1}y)
\,dy=\sum_{\gamma\in U_c\backslash\Gamma_{c,s,c'}/U_{c'}}\int_{(U_c\cap\gamma U_{c'}\gamma^{-1})
\backslash G_\infty}\aubar{W_i(y)}\,W_{i'}(\gamma^{-1}y)\,dy ,
\]
by grouping the $\gamma$ in orbits of $U_c$; the sum of the integrals of the absolute values is at
most $\mathsf m\,\nu\int_{U_\infty\backslash G_\infty}|W_i|$. This proves the convergence (for
$T=\Gamma$, hence for every stable $T$) and $\pbip{P_{c,i}}{\varrho(s)P_{c',i'}}=\pbvol(J)\,\aubar{A_\Gamma
(c,s,c')_{i,i'}}$. The formula follows by sesquilinearity. For the last claim, the vectors with one
non-zero coordinate $\varphi(h_c)_i=1$ form an orthonormal basis of $\auS$.
\end{proof}

\begin{proposition}[The terms in the Bessel subgroup]\label{pb:besselterms}
Assume \textup{(B1)--(B4)} and let $P$ be as in Proposition~\ref{pb:unfold}. For all $s\in\auG$ and
$\varphi,\varphi'\in\auS$,
\[
\frac1\nu\sum_{c,c',i,i'}\aubar{\varphi(h_c)_i}\,\varphi'(h_{c'})_{i'}\,\aubar{A_{R_0(E^+)}(c,s,c')_{
i,i'}}=\pbip\varphi{\pbrpsi(s)\varphi'} .
\]
Hence $\pbHS{P_s-\pbrpsi(s)}^2=\nu^{-2}\sum_{c,c'}\|A_{\Gamma\smallsetminus R_0(E^+)}(c,s,c')\|^2$.
\end{proposition}
\begin{proof}
Let $\gamma=u_0\eta\in R_0(E^+)\cap\Gamma_{c,s,c'}$ with $u_0\in U(E^+)$, $\eta\in H(E^+)$. Then
$h_c^{-1}\gamma h_{c'}=(h_c^{-1}u_0h_c)(h_c^{-1}\eta h_{c'})\in K_R$, so by Lemma~\ref{pb:KR}
$s=\bar u\bar t\in\pbR$, where $\bar u$ is the image of $h_c^{-1}u_0h_c\in K_U$ and $\bar t$ that
of $k_t:=h_c^{-1}\eta h_{c'}\in K_H$. In particular there are no such terms if $s\notin\pbR$, and
both sides vanish. Let $s=\bar u\bar t\in\pbR$ and fix $c$. For a lift $k\in K_H$ of $\bar t$, $h_ck\in H(E^+)h_{c'}
J_X$ for one class $c'$, which gives a pair $(\eta,c')$ with $h_c^{-1}\eta h_{c'}\in K_H$ of image
$\bar t$. The pair is unique: if $(\eta_2,c'')$ is another one, then $h_{c'}^{-1}\eta^{-1}\eta_2h_{
c''}\in K_H\cap J=J_X$, so $c''=c'$ and $\eta^{-1}\eta_2\in H(E^+)\cap h_{c'}J_Xh_{c'}^{-1}=1$, the
level being free. The elements $u_0$ that occur form one coset
$U_cu_0$, non-empty by strong approximation for $U$. Finally $U_cu_0\eta=u_0\eta\,U_{c'}$, because
$K$ normalises $J$. So over $(c,s,\cdot)$ there is exactly one double coset of terms in
$R_0(E^+)$, it lies over $(c,s,c')$, and $U_c\cap\gamma U_{c'}\gamma^{-1}=U_c$. By (B1),
$W(\gamma^{-1}y)=\aubar{\psi_\infty(u_0)}\,\tau_\infty(\eta)^{-1}W(y)$, so by Lemma~\ref{pb:ortho}
\[
A_{R_0(E^+)}(c,s,c')=\psi_\infty(u_0)\int_{U_c\backslash G_\infty}W(y)W(y)^*\,dy\cdot\tau_\infty(\eta)
=\nu\,\psi_\infty(u_0)\,\tau_\infty(\eta) .
\]
Here $\aubar{\psi_\infty(u_0)}=\psi_f(u_0)=\psi_U(\bar u)$. The representation $\tau_\infty$ is
the contragredient of the unitary action of $H_\infty$ on $W_{\lambda^X}$, so in dual orthonormal
bases $\aubar{\tau_\infty(\eta)_{i,i'}}$ is the matrix of $\eta$ on $W_{\lambda^X}$; and $\eta\,
\varphi'(h_{c'})=\varphi'(\eta h_{c'})=\varphi'(h_ck_t)$. So the left side is
\[
\psi_U(\bar u)\sum_c\pbip{\varphi(h_c)}{\eta\,\varphi'(h_{c'})}_{W_{\lambda^X}}=\psi_U
(\bar u)\sum_c\pbip{\varphi(h_c)}{(R(\bar t)\varphi')(h_c)}_{W_{\lambda^X}}=\pbip\varphi{\pbrpsi(s)
\varphi'} .
\]
The last claim follows from Proposition~\ref{pb:unfold} and $A_\Gamma=A_{R_0(E^+)}+A_{\Gamma
\smallsetminus R_0(E^+)}$.
\end{proof}

\subsubsection*{The two normalising constants}

\begin{lemma}[Ratio of masses]\label{pb:massratio}
For $l\ge1$,
\[
\frac{\pbmass_{l+1}}{\pbmass_l}=\Bigl(\frac{\aurd^{\,l+1/2}\;l!}{(2\pi)^{l+1}}\Bigr)^{f}\cdot L\bigl(
l+1,\chi_{E/E^+}^{\,l+1}\bigr),
\]
where $\chi_{E/E^+}$ is the quadratic character of $E/E^+$ and $L(s,\chi^j_{E/E^+})$ means
$\zeta_{E^+}(s)$ for even $j$. In particular $\bigl|\log L(l+1,\cdot)\bigr|\le f\,2^{1-l}$ for
$l\ge2$.
\end{lemma}
\begin{proof}
This is the mass formula for the genus of the unimodular lattice $\auO_E^l$ in the form of
Gan, Hanke and Yu~\cite[Propositions~4.4 and 4.5]{GHY}, $\pbmass_l=\mathsf c^{\,l}\prod_{j=1}^l|L(1-j,\chi_{E/E^+}^{\,j})|$
up to a factor independent of $l$, with all local factors equal to $1$ because $E/E^+$ is
unramified at the finite places and the lattice is unimodular; together with the functional
equations of $\zeta_{E^+}$ and $L(s,\chi_{E/E^+})$, which give $|L(1-j,\chi^j)|=\bigl(2\,\aurd^{\,j-
1/2}(j-1)!/(2\pi)^j\bigr)^f L(j,\chi^j)$ (the discriminant of $E^+$ is $\aurd^{\,f}$ and the
conductor of $\chi_{E/E^+}$ is trivial); the formula of the lemma corresponds to $\mathsf c=2^{-f}$.
The bound for $\log L$ follows from $|\log L(s,\cdot)|\le f\log\zeta(s)$. In the notation of~\cite[(4.2), (4.3), Propositions~4.4--4.5]{GHY},
$\mathrm{Mass}=2^{-lf}\,\tau(G)\,L(M)\prod_{v\in S}\lambda_v$ with $\tau(G)=2$, $L(M)=\prod_{j\le l}L(1-j,\chi^j)$ and $S=\emptyset$; hence $\mathsf c=2^{-f}$. A constant $\mathsf c=C^f$ with another
absolute $C$ would multiply the ratio by a factor $(2C)^f$, which changes nothing in the proof of Lemma~\ref{pb:sizes}. The exact constant $\mathsf c=2^{-f}$ is used in Lemma~\ref{pbm:const}\,\textup{(C0)}, where the ratio $\pbmass_r/\pbmass_k$ has to be known up to a factor $1+b^{-10}$.
\end{proof}

\begin{lemma}[Sizes]\label{pb:sizes}
$b^{11kf}\,\pbDid\le\pbDmain\le(M/b^6)^{kf}\,\pbDid$.
\end{lemma}
\begin{proof}
By Lemma~\ref{pb:KR} and \eqref{pb:eq-classes}, $\#\mathcal C=[K^\circ_k:J_X]\,\pbmass_k=\#\pbHb
\cdot(N^k-1)^f\cdot\pbmass_k$: the index of $K_H$ in $K^\circ_k$ is the number of possible last
rows modulo the $f$ primes dividing $\pi_N$. With \eqref{pb:eq-orders}, \eqref{pb:eq-DX} and
\eqref{pb:eq-D},
\[
\mathcal R:=\frac{\pbDmain}{\pbDid}=\frac{\auDX}{\#\pbR\cdot\pbdeg\cdot\pbmass_r}=\frac{d\,(N^k-1)^f
\,\pbmass_k}{M^{f\varrho(2k+2\varrho+1)}\,\pbdeg\,\pbmass_r} .
\]
\emph{The degree.} Put $\mathsf L=\sum_\tau\log\kappa(\tau)$; $0\le\mathsf L\le f\log b$. The first
$k-1$ entries of $\varsigma_\tau$ are consecutive, so the pairs $i<j\le k-1$ contribute $1$ to
$\pbdeg_\tau$, and $\log\pbdeg=\sum_\tau\sum_{(i,j)}\log(\varsigma_{\tau,i}-\varsigma_{\tau,j})-f
\sum_{l=k-1}^{r-1}\log l!$, the sum over the pairs $i<j$ with $j\ge k$. There are $m(k-1)$ such
pairs with $i\le k-1$, each with a difference between $\kappa(\tau)$ and $\kappa(\tau)+m\Delta+k\le
\kappa(\tau)(1+1/x)$, and the other pairs have a difference between $\Delta$ and $m\Delta\le\kappa(
\tau)$; there are at most $mr$ pairs. So
\[
m(k-1)\mathsf L\ \le\ \log\pbdeg+f\sum_{l=k-1}^{r-1}\log l!\ \le\ mr(\mathsf L+f/x).
\]
\emph{The masses.} By Lemma~\ref{pb:massratio}, with $\mathsf S=\sum_{l=k}^{r-1}(l+\frac12)=(m-1)(k+
\frac{m-1}2)$,
\[
\log\frac{\pbmass_r}{\pbmass_k}=f\Bigl(\mathsf S\log\aurd+\sum_{l=k}^{r-1}\log l!-\bigl(\mathsf S+
\tfrac{m-1}2\bigr)\log2\pi\Bigr)+\theta,\qquad|\theta|\le1 .
\]
\emph{The level.} By {\cite[(37)]{P46}}, $kf\log N=km\mathsf L+kf(m-1)\log\frac{M\aurd}{2\pi}-
\frac{15}2kf\log b+\theta'$ with $0\le\theta'\le kf\sqrt{\log b}$, and $kf\log N-1\le f\log(N^k-1)
\le kf\log N$. Also $\varrho(2k+2\varrho+1)=(m-2)k+\frac{(m-2)(m-1)}2$.

Inserting these, the factorials cancel except $(k-1)!^f$, the coefficient of $\log M$ is $kf-f
\frac{(m-2)(m-1)}2$, that of $\log\aurd$ is $-f\frac{(m-1)^2}2$, that of $\log2\pi$ is $f\frac{m(m-
1)}2$, and
\[
\log\mathcal R=\log\bigl(d\,(k-1)!^f\bigr)+kf\bigl(\log M-\tfrac{15}2\log b\bigr)+\mathrm{Err},
\]
\[
-m^2\mathsf L-\frac{mrf}x-\frac{fm^2}2\log(M\aurd)-2\ \le\ \mathrm{Err}\ \le\ m\mathsf L+kf\sqrt{
\log b}+\frac{fm^2}2\log2\pi+1 .
\]
As $m^2\le4x$, $r\le3x^5$, $k\ge x^5/2$, $\log(M\aurd)\le24\log b$ and $\log b$ is large,
$|\mathrm{Err}|\le\frac14kf\log b$ (a comparison of sizes). By \eqref{pb:eq-DX}, $1\le d\le(a+k)^{
kf}$, so $0\le\log(d\,(k-1)!^f)\le kf(\log(2xb)+\log k)\le\frac54kf\log b$. Finally $20\log b\le
\log M$. Hence $\log\mathcal R\ge kf(20-\frac{15}2-\frac14)\log b\ge11kf\log b$ and $\log\mathcal
R\le kf(\log M-\frac{15}2\log b+\frac54\log b+\frac14\log b)=kf(\log M-6\log b)$.
\end{proof}

Lemma~\ref{pb:sizes} is the place where the level $N$ is tuned: $\pbDmain$ is the square of the
dimension of the source divided by (a bound for) the dimension of the target per element of
$\auG$, $\pbDid$ is what the terms in the Bessel subgroup contribute to the second moment, and the
window~{\cite[(37)]{P46}} puts their ratio between two powers of $b^{kf}$.

\subsubsection*{Links, meeting places, and the three statements on amplitudes}
For $\gamma\in\Gamma$ let $Z(\gamma)=\bigl(h(\gamma f_l,f_j)\bigr)_{0\le j,l\le\varrho}\in M_n(E)$,
with leading minors $\delta_l(\gamma)$ ($0\le l\le n$), and
\[
\pbg(\gamma)=\delta_n(\gamma)\aubar{\delta_n(\gamma)}-N^2\,\delta_{n-1}(\gamma)\aubar{\delta_{n-1}
(\gamma)}\in E^+ .
\]
The \emph{link} of $\gamma$ is the subspace of $(A^\perp/A)^2$ of the pairs $([t],[\gamma^{-1}t])$
for $t\in Y_0$ orthogonal to $A$ and to $\gamma A$. Both $\pbg$ and the link are invariant under
$\gamma\mapsto u\gamma u'$, $u,u'\in U(E^+)$, so that for a subspace $\Lambda$ the set
$T_\Lambda=\{\gamma:\pbg(\gamma)\ne0,\ \text{the link of }\gamma\text{ is }\Lambda\}$ is stable. For
$\Omega\subseteq\auG$ the \emph{link moment} is
\[
\pbLM(\Omega)=\sum_{c,c'\in\mathcal C}\ \sum_{s\in\Omega}\ \sum_\Lambda\bigl\|A^\flat_{T_\Lambda}(c,
s,c')\bigr\|^2\in[0,\infty].
\]
An element $s\in\auG$ \emph{meets} at a place $\mathfrak p\in\{w,\aubar w\}$, $w$ chosen, if
$F_{\mathfrak p}\cap sF_{\mathfrak p}\ne0$; its \emph{meeting norm} $\mathsf{mn}(s)$ is the product
of $\auNm\mathfrak p$ over the places at which it meets. A \emph{scalar coset} is a set $\zeta\pbR$,
where $\zeta\in\auO_E$, $\zeta\aubar\zeta=1$, acts on each $V_w$ by its residue; $\Omega$
\emph{avoids the scalar cosets} if $\Omega\cap\zeta\pbR=\emptyset$ for all such $\zeta\ne1$.

\begin{theorem}[Amplitudes]\label{pb:amp}
Bessel functions exist for the data (Definition~\ref{pb:bessel}), and they can be chosen such
that, with $\nu$ as in Lemma~\ref{pb:ortho} and $P$ as in Proposition~\ref{pb:unfold}:
\begin{enumerate}
\item[(X)] \emph{(Expansion of the squares.)} If $\Omega$ avoids the scalar cosets and $\pbLM(
\Omega)\le C\nu^2$ with $C\ge0$, then
\[
\Bigl(\sum_{s\in\Omega}\pbHS{P_s-\pbrpsi(s)}^2\Bigr)^{1/2}\le\sqrt C+b^{-15}\sqrt{\pbDid} .
\]
\item[(L1)] \emph{(Links and meeting places.)} If $\mathsf z\ge1$ and $\mathsf{mn}(s)\ge\mathsf z$
for all $s\in\Omega$, then $\pbLM(\Omega)\le\bigl(b^{2f}/\mathsf z\bigr)^k\,\pbDmain\,\nu^2$.
\item[(L2)] \emph{(The main term.)} For every set $\mathsf A$ of chosen places above $p_0$ and every
family $\mathsf d$ of frame vectors,
$\pbLM\bigl(\{s\in\pbFix(\mathsf A,\mathsf d):s\text{ open}\}\bigr)\le(1+b^{-4})\,b^{-6k\#\mathsf A}\,
\pbDmain\,\nu^2$.
\end{enumerate}
\end{theorem}
\smallskip\noindent The existence of the Bessel functions is Theorem~\ref{pbe:bessel}; {\rm(X)} and
{\rm(L1)} are Theorem~\ref{pbe:XL1} (from Theorems~\ref{pbe:ds} and~\ref{pbe:pillars}, proved in
\S\S\ref{pds:sec}--\ref{pbo:sec}); {\rm(L2)} is Theorem~\ref{pbm:L2}, deduced in \S\ref{pbm:sec} from
Theorem~\ref{pbm:inputs}, which is proved in {\cite[\S\S2--6]{P52}}. All three statements hold
for the same Bessel functions, those of Theorem~\ref{pbe:bessel}.\par\smallskip

\begin{lemma}\label{pb:elem}
\begin{enumerate}
\item[(a)] At most $M^f\,\#\pbR$ elements of $\auG$ lie in a scalar coset.
\item[(b)] An element of $\pbFix(\mathsf d')$, $\mathsf d'$ a frame vector in $F_{\mathfrak p}$,
meets at $\mathfrak p$. Every $s\in\Omega_i$ meets at the $2f$ places $w,\aubar w$ with $w\mid
p_i$ chosen.
\item[(c)] No element of a scalar coset is open.
\item[(d)] $\Omega_i$ avoids the scalar cosets.
\end{enumerate}
\end{lemma}
\begin{proof}
(a) $\zeta\pbR$ depends only on the residues of $\zeta$ at the $2f$ chosen places, at most $\prod_i
p_i^f=M^f$ possibilities. (b) The first claim is clear. An element of $R_0$ fixes $e_1=f_0$; so
does its adjoint inverse. Hence an element of $\pbR$ fixes $f_0$ in $V_w$ and in $V_{\aubar w}$,
and $s\in\Omega_i$ does so for $w\mid p_i$. (c) For $s\in\zeta\pbR$, $s_wf_0=\zeta f_0$, and $f_0=
e_1$ is orthogonal to the frame as $\varrho\ge1$; so the first column of $Z_w(s_w)$ vanishes and
$\delta_1(s_w)=0$ for every $w$. (d) Let $s\in\Omega_i\cap\zeta\pbR$. Comparing the action on $f_0$
modulo the chosen places $w\mid p_i$, by (b), gives $\zeta\equiv1$ modulo $w$. As $\zeta$ is an
algebraic integer of absolute value $1$ at every infinite place of the CM field $E$, it is a root
of unity (Kronecker), and a root of unity $\ne1$ that is $\equiv1$ modulo a prime above $p_i$ has
order divisible by $p_i$ (the product of the $1-\zeta^j$, $0<j<\operatorname{ord}\zeta$, is
$\operatorname{ord}\zeta$), which is impossible in $E$ because $p_i$ is unramified in $E$. So
$\zeta=1$.
\end{proof}

\begin{proposition}[From amplitudes to (E1)--(E3)]\label{pb:reduction}
Theorem~\ref{pb:amp} implies \textup{(E1)}, \textup{(E2)}, \textup{(E3)} of Theorem~\ref{pb:est},
with $c_P=(\pbvol(J)\nu)^{-1/2}c_W^{-1}$.
\end{proposition}
\begin{proof}
Write $\mathsf E(\Omega)=\sum_{s\in\Omega}\pbHS{P_s-\pbrpsi(s)}^2$.

(E1) By Lemma~\ref{pb:elem}(b),(d), $\Omega_i$ avoids the scalar cosets and $\mathsf{mn}\ge p_i^{2f}
\ge(M/b)^f$ on $\Omega_i$ (as $p_{1-i}\le b^{11}\le b\,p_i$). By (L1) and Lemma~\ref{pb:sizes},
\[
\pbLM(\Omega_i)\,\nu^{-2}\le\Bigl(\frac{b^{3f}}{M^f}\Bigr)^k\pbDmain\le\Bigl(\frac{b^{3}}{M}\Bigr)^{
kf}\Bigl(\frac M{b^6}\Bigr)^{kf}\pbDid=b^{-3kf}\pbDid\le b^{-30}\pbDid .
\]
By (X), $\sqrt{\mathsf E(\Omega_i)}\le2b^{-15}\sqrt{\pbDid}$, so $\mathsf E(\Omega_i)\le4b^{-30}\pbDid
\le b^{-12}\pbDid$.

In particular Lemma~\ref{pb:basic} applies, since its proof uses, of the three inequalities to be proved, only (E1), which is now established: $\beta(s)=\pbHS{P_s}^2\le2\auDX$ for all $s$.

(E3) Let $\Omega=\{s\in\pbFix(\mathsf A,\mathsf d):s\text{ open}\}$ and $\mathsf T=b^{-6k\#\mathsf
A}\pbDmain$. By Lemma~\ref{pb:elem}(c), $\Omega$ meets no scalar coset; in particular $\Omega\cap\pbR
=\emptyset$, so $\pbrpsi$ vanishes on $\Omega$ and $\sum_{s\in\Omega}\beta(s)=\mathsf E(\Omega)$. By
(L2) and (X), $\sqrt{\mathsf E(\Omega)}\le\sqrt{(1+b^{-4})\mathsf T}+b^{-15}\sqrt{\pbDid}$. By
Lemma~\ref{pb:sizes}, $\pbDid\le b^{-11kf}\pbDmain\le\mathsf T$. So $\sqrt{\mathsf E(\Omega)}\le(1+
\frac12b^{-4}+b^{-15})\sqrt{\mathsf T}\le(1+b^{-4})\sqrt{\mathsf T}$ and $\mathsf E(\Omega)\le(1+
b^{-4})^2\mathsf T\le(1+4b^{-4})\mathsf T$.

(E2) Let $\mathsf T=b^{-6kf}\pbDmain$ and split $\pbFix(\mathrm{all},\mathsf d)$ into the elements
that lie in a scalar coset and the set $\Omega$ of the others. By Lemma~\ref{pb:elem}(a) and
$\beta\le2\auDX$, the first part contributes at most $2M^f\,\#\pbR\,\auDX=2M^f\pbDid\le2b^{23f}
b^{-11kf}\pbDmain\le\frac14\mathsf T$. The set $\Omega$ meets no scalar coset, so $\sum_{s\in\Omega}
\beta(s)=\mathsf E(\Omega)$ as before, and $\mathsf{mn}\ge p_0^f\ge b^{10f}$ on $\Omega$ by
Lemma~\ref{pb:elem}(b). By (L1), $\pbLM(\Omega)\nu^{-2}\le(b^{2f}/b^{10f})^k\pbDmain=b^{-2kf}\mathsf
T$, and by (X) and $\pbDid\le\mathsf T$,
\[
\mathsf E(\Omega)\le2\bigl(b^{-2kf}\mathsf T+b^{-30}\pbDid\bigr)\le\tfrac12\mathsf T .
\]
So the sum over $\pbFix(\mathrm{all},\mathsf d)$ is at most $\mathsf T$.
\end{proof}

\subsection{The dependency tree}\label{pb:ss-remains}

{\cite[Theorem~14.14]{P46}} is proved above from the following statements; we list where each is proved.

\begin{enumerate}
\item[(R1)] \emph{The discrete series coefficient}, Theorem~\ref{pbe:ds}: \textup{(DS1)}--\textup{(DS3)} in
\S\ref{pds:sec}, the explicit formula \eqref{pds:eq-FJ} behind \textup{(DS2)} in {\cite[Theorem~3.17]{P49}}, and
\textup{(DS4)} in \S\ref{pbp:sec} (Theorem~\ref{pbp:main}). The Bessel functions and the Poincar\'e series
are treated in \S\S\ref{pbe:ss-real}--\ref{pbe:ss-global} (Theorem~\ref{pbe:bessel}).
\item[(R2)] \emph{Size, gain and count}, Theorem~\ref{pbe:pillars} in the form of
Theorem~\ref{ppi:pillars}: \textup{(R)} in Proposition~\ref{ppi:cell}, \textup{(G$'$)} in
\S\ref{pbg:sec}, \textup{(C1)}, \textup{(C2)} in \S\ref{pbc:sec}, \textup{(S)} in \S\ref{pbs:sec} from the
open-cell bound \textup{(O1)} of \S\ref{pbo:sec} and Lemma~\ref{pb:covol}. With them \textup{(X)} and
\textup{(L1)} follow (Theorem~\ref{pbe:XL1}).
\item[(R4)] \emph{The main term}, Theorem~\ref{pb:amp}\textup{(L2)}: Theorem~\ref{pbm:L2}, from
Theorem~\ref{pbm:inputs} and Lemma~\ref{pb:covol}; \textup{(R)} is {\cite[Theorem~2.17]{P52}}
({\cite[\S\S2, 3 and~4]{P52}}), \textup{(F)} and \textup{(S)} are {\cite[Theorems~5.28 and~5.29]{P52}} ({\cite[\S\S5 and~6]{P52}}); the constant is evaluated in Lemma~\ref{pbm:const}.
\item[(R5)] \emph{The induced character}, Theorem~\ref{pb:est}\textup{(E4)}: Theorem~\ref{pbm:E4}.
\item[(R6)] \emph{Equality of multiplicities}, {\cite[Proposition~14.36]{P46}}: Proposition~\ref{pm:mult}, from
Theorem~\ref{pm:arch}, whose part (a) is Remark~\ref{pm:rem-other}(1) and part (b) is
Theorem~\ref{pbr:main}.
\end{enumerate}
Every statement in this tree is proved in this paper, in~\cite{P49} or in~\cite{P52}, at the place named; the exact value of the covolume in
Lemma~\ref{pb:covol} is Theorem~\ref{pbv:covol}. (Part~(b) of Theorem~\ref{pm:arch} is proved in \S\ref{pbr:sec}.) In addition, Lemma~\ref{pb:massratio} rests on the mass formula of Gan, Hanke
and Yu~\cite[Propositions~4.4--4.5]{GHY}, and Lemma~\ref{pb:bad} on Green's description of
the characters of $\auGL_r(\auF_q)$.



\section{Covolumes, Poincar\'e series and the main term}\label{B1:g2}
% ===== B-covolume.tex =====
% proofs/B-covolume.tex
% The covolume of the unipotent congruence groups (Lemma pb:covol), and the completion of
% some short steps of Theorem B. Labels of this file start with pbv:.
\providecommand{\auO}{\mathcal O}
\providecommand{\auPl}{\Phi}
\providecommand{\auU}{\mathrm U}
\providecommand{\auGL}{\mathrm{GL}}
\providecommand{\auA}{\mathbb A}
\providecommand{\aurd}{\mathrm{rd}_E}
\providecommand{\aubar}[1]{\overline{#1}}
\providecommand{\pbvol}{\operatorname{vol}}
\providecommand{\pbUb}{\overline U}
\providecommand{\auNm}{\operatorname{Nm}}
\providecommand{\pbg}{\mathsf g}
\providecommand{\pffOT}{\auO_T}
\providecommand{\pfrG}{\mathsf G}
\providecommand{\pfrN}{\mathsf N}
\providecommand{\pgnD}{\mathfrak d}
\providecommand{\pgnd}{\operatorname{d}}
% macros of this file (prefix \pbv)
\providecommand{\pbvO}{\widehat{\auO}}
\providecommand{\pbvL}{\mathsf L}
\providecommand{\pbvcov}{\operatorname{covol}}
\providecommand{\pbvE}{\mathsf E}

\subsection{The covolume of the unipotent congruence groups}\label{pbv:sec}

This section proves Lemma~\ref{pb:covol}: for the Haar measure $du$ of convention (Cv3) of
\S\ref{pb:ss-setting}, the covolume $\nu=\pbvol(U_c\backslash U_\infty)$ of Lemma~\ref{pb:ortho} equals
the number $\nu_*$ defined before Lemma~\ref{pb:covol}, for every class $c$. The computation is
elementary: the coordinates of (Cv3) (the free entries of the rows $\tilde e_a$ and the imaginary
parts of the entries $(\tilde e_a,\tilde e_a^*)$) are additive along a triangular order
(Lemma~\ref{pbv:add}); the group $U_c$ translates each coordinate, once the earlier ones vanish, by an
explicit lattice, and every point of that lattice is attained by an Eichler transformation
(Lemma~\ref{pbv:latt}); so a product of fundamental parallelotopes is a fundamental domain
(Lemma~\ref{pbv:box}), and $\nu$ is the product of the covolumes of $2\varrho$ lattices: $\varrho$
lattices of linear forms on the spaces $Y_a$ (one for each floor of the tower of Heisenberg groups
of $U$), and $\varrho$ copies, rescaled, of the lattice of skew-hermitian integers of $E$
(Lemmas~\ref{pbv:vol-a}, \ref{pbv:vol-skew}). In \S\ref{pbv:ss-uses} we check that the three forms
of the constant used in the text are the same number. In \S\ref{pbv:ss-weak} we write out in detail four short
steps of the Spanning Theorem ({\cite[Theorem~14.14]{P46}}).

\emph{What is proved.} Lemma~\ref{pb:covol} is proved in full (Theorem~\ref{pbv:covol}), using two
classical facts quoted with their hypotheses: Minkowski's formula for the covolume of the ring of
integers of a number field, and the tower formula for discriminants (used only to see that
$|D_E|=D_{E^+}^2$, since $E/E^+$ is unramified at all finite places).

\subsubsection{Coordinates on the unipotent group}\label{pbv:ss-coord}

We use the notation of \S\ref{pb:ss-setting} and convention (Cv1): $h$ is linear in the first
variable, and $h(v,e_a^*)$ is the coefficient of $e_a$ in $v$. Let $F$ be either $E$, or $\C$ with
$Y_F=Y_\tau=Y_0\otimes_{E,\tau}\C$ for some $\tau\in\auPl$ (with the form $h_\tau(v\otimes1,w\otimes
1)=\tau(h(v,w))$). Write $\aubar x$ for the conjugation of $F$ and $F^+=\{x:\aubar x=x\}$. Let
$U(F)$ be the group of $F$-linear isometries of $Y_F$ that preserve the flag $\langle e_1\rangle
\subset\langle e_1,e_2\rangle\subset\dots\subset A$, act trivially on its graded pieces and act
trivially on $A^\perp/A$; so $U(E)=U(E^+)$ and $U(\C)=U_\tau$ in the notation of
\S\ref{pb:ss-setting}. For $0\le a\le\varrho$ put
\[
A_a=\langle e_1,\dots,e_a\rangle,\qquad Y_a=\langle e_b,e_b^*\ (b>a);\ x_1,\dots,x_k;\ z\rangle,
\qquad m_a=\dim Y_a=2(\varrho-a)+k+1 .
\]
So $A_a$ is isotropic, $A_a^\perp=A_a\oplus Y_a$, $A_{a-1}^\perp=A_a^\perp\oplus Fe_a^*$, $A_0^\perp=Y_F$,
$A_\varrho=A$, and $Y_a$ is non-degenerate. Every $u\in U(F)$ satisfies $(u-1)e_b\in A_{b-1}$ for all
$b$, preserves each $A_a$, hence each $A_a^\perp$, and $(u-1)A_\varrho^\perp\subseteq A_\varrho$.

\paragraph*{The coordinates}
For $u\in U(F)$ and $1\le a\le\varrho$ let the \emph{row form} of $u$ at $a$ be the linear form
\[
\mathsf r_a(u)\colon Y_a\to F,\qquad \mathsf r_a(u)(v)=h(uv,e_a^*)=h((u-1)v,e_a^*),
\]
the coefficient of $e_a$ in $uv$ (the two expressions agree because $Y_a\perp e_a^*$), and let
$\mathsf y_a(u)=h(ue_a^*,e_a^*)$, the coefficient of $e_a$ in $ue_a^*$. In a basis of $Y_a$, the values
of $\mathsf r_a(u)$ on the basis vectors are the entries of the row of $e_a$ of the matrix of $u$ in the
columns of $Y_a$; these are the columns $e_b,e_b^*$ $(b>a)$, $x_i$, $z$, that is, exactly the entries
of the row of $e_a$ above the diagonal block that are not fixed by unitarity. (The other entries of
that row are: $1$ at $e_a$ and $0$ at $e_b$, $b\ne a$, $b\le a$, by unipotence; the entry
$\mathsf y_a(u)$ at $e_a^*$, whose real part is determined by the other entries through
$h(ue_a^*,ue_a^*)=0$ (see the proof of Lemma~\ref{pbv:rows}); and the entries at $e_b^*$, $b<a$, which are determined by the rows $e_{b'}$,
$b'>b$.)

At a place $\tau$ we use the normalised basis of \S\ref{pbe:ss-real}: $\tilde e_a=\lambda_{\tau,a}e_a$,
$\tilde e_a^*=\lambda_{\tau,a}^{-1}e_a^*$, $x_i$, $\tilde z=s_{\tau,\varrho}z$, where
\begin{equation}\label{pbv:eq-lambda}
\lambda_{\tau,a}=\frac{\kappa_\tau^{\,\varrho-a+1}}{\sqrt N}=s_{\tau,a-1}\quad(1\le a\le\varrho),\qquad
s_{\tau,\varrho}=\frac1{\sqrt N},
\end{equation}
with $s_{\tau,a}=\kappa_\tau^{\varrho-a}/\sqrt N$ as in \eqref{pbe:eq-scale} (so $\tilde f_a=s_{\tau,a}f_a$).
Let $v_1,\dots,v_{m_a}$ be the basis $e_b,e_b^*$ $(b>a)$, $x_1,\dots,x_k$, $z$ of $Y_a$ and $\tilde
v_j=\lambda_{\tau,j}v_j$ the corresponding normalised vectors ($\lambda_{\tau,j}=\lambda_{\tau,b}$,
$\lambda_{\tau,b}^{-1}$, $1$, $s_{\tau,\varrho}$ respectively). The coordinates of (Cv3) of $u\in U_\tau$
are
\begin{equation}\label{pbv:eq-coord}
\rho_a(u)=\bigl(h(u\tilde v_j,\tilde e_a^*)\bigr)_{j\le m_a}=\Bigl(\frac{\lambda_{\tau,j}}{\lambda_{\tau,a}}
\,\mathsf r_a(u)(v_j)\Bigr)_j\in\C^{m_a},\qquad t_a(u)=\Iml h(u\tilde e_a^*,\tilde e_a^*)=
\lambda_{\tau,a}^{-2}\Iml\mathsf y_a(u)\in\R
\end{equation}
$(1\le a\le\varrho)$, and $du_\tau$ is the measure that corresponds to Lebesgue measure under
$c_\tau=(\rho_a,t_a)_{a\le\varrho}\colon U_\tau\to V:=\prod_{a=1}^{\varrho}(\C^{m_a}\times\R)$, Lebesgue
measure on $\C$ being $d\Rel\,d\Iml$. That $c_\tau$ is a bijection is proved in
Corollary~\ref{pbv:bij} below; it is continuous (polynomial), so it is a Borel isomorphism (a
continuous injection of a Polish space is a Borel isomorphism onto its image, by the Lusin--Souslin
theorem), and $du_\tau(S)=\pbvol(c_\tau(S))$ for Borel $S$. We order the $2\varrho$ coordinates as
\begin{equation}\label{pbv:eq-order}
\rho_\varrho<t_\varrho<\rho_{\varrho-1}<t_{\varrho-1}<\dots<\rho_1<t_1 ,
\end{equation}
and we call $u$ \emph{low for} a coordinate if all coordinates of $u$ before it vanish. Since the
$\lambda$'s are non-zero reals, the vanishing of $\rho_a(u)$ means $\mathsf r_a(u)=0$, and the vanishing of
$t_a(u)$ means $\mathsf y_a(u)\in\R$; the same definitions are used over $E$, with
\emph{$\mathsf r_a=0$ and $\mathsf y_a\in E^+$} in place of $\rho_a=0$, $t_a=0$.

\begin{lemma}[Vanishing rows]\label{pbv:rows}
Let $u\in U(F)$ and $1\le a\le\varrho$.
\begin{enumerate}
\item[(i)] If $\mathsf r_d(u)=0$ and $\mathsf y_d(u)\in F^+$ for all $d$ with $a<d\le\varrho$, then
$(u-1)A_a^\perp\subseteq A_a$.
\item[(ii)] If moreover $\mathsf r_a(u)=0$, then $(u-1)A_a^\perp\subseteq A_{a-1}$, $(u-1)e_a^*\in A_a$, and
$\mathsf y_a(u)+\aubar{\mathsf y_a(u)}=0$.
\item[(iii)] If $\mathsf r_d(u)=0$ and $\mathsf y_d(u)\in F^+$ for all $d\le\varrho$, then $u=1$.
\end{enumerate}
\end{lemma}
\begin{proof}
We show by descending induction on $d$ that $(u-1)A_d^\perp\subseteq A_d$ for $a\le d\le\varrho$, and
along the way the claims of (ii). For $d=\varrho$ this is the definition of $U$. Let $a<d\le\varrho$, or
$d=a$ in the situation of (ii), and assume $(u-1)A_d^\perp\subseteq A_d$ and $\mathsf r_d(u)=0$. For $v\in
A_d^\perp=A_d\oplus Y_d$ the vector $(u-1)v$ lies in $A_d$, and its coefficient at $e_d$ is $0$: for
$v=e_b$, $b\le d$, because $(u-1)e_b\in A_{b-1}\subseteq A_{d-1}$; for $v\in Y_d$, because it is
$\mathsf r_d(u)(v)=0$. Hence $(u-1)A_d^\perp\subseteq A_{d-1}$. Put $y=(u-1)e_d^*$. For $v\in A_d^\perp$
write $uv=v+\alpha$ with $\alpha\in A_{d-1}$; as $u$ is an isometry and $\alpha\perp e_d^*$,
\[
h(e_d^*,v)=h(ue_d^*,uv)=h(e_d^*,v)+h(e_d^*,\alpha)+h(y,uv)=h(e_d^*,v)+h(y,uv),
\]
so $h(y,uv)=0$. Since $uA_d^\perp=A_d^\perp$, $y\in(A_d^\perp)^\perp=A_d$. The coefficient of $e_d$ in $y$ is
$h(y,e_d^*)=\mathsf y_d(u)$ (as $h(e_d^*,e_d^*)=0$), and, $A_d$ being isotropic,
\[
0=h(ue_d^*,ue_d^*)=h(e_d^*+y,e_d^*+y)=h(e_d^*,y)+h(y,e_d^*)=\aubar{\mathsf y_d(u)}+\mathsf y_d(u).
\]
This proves (ii) when $d=a$. If $d>a$, then $\mathsf y_d(u)\in F^+$ by hypothesis, so $\mathsf y_d(u)=0$, $y\in
A_{d-1}$, and $(u-1)A_{d-1}^\perp\subseteq A_{d-1}$ because $A_{d-1}^\perp=A_d^\perp+Fe_d^*$. This completes the
induction and proves (i). For (iii), the case $d=1$ of the induction gives $(u-1)A_0^\perp\subseteq
A_0=0$, and $A_0^\perp=Y_F$.
\end{proof}

\begin{lemma}[Additivity]\label{pbv:add}
Let $x,u\in U(F)$ and $1\le a\le\varrho$, and assume $\mathsf r_d(u)=0$ and $\mathsf y_d(u)\in F^+$ for
$a<d\le\varrho$. Then $\mathsf r_a(xu)=\mathsf r_a(x)+\mathsf r_a(u)$. If moreover $\mathsf r_a(u)=0$, then
$\mathsf y_a(xu)=\mathsf y_a(x)+\mathsf y_a(u)$. In particular, at a place $\tau$, the coordinate of
$xu$ for which $u$ is low is the sum of the coordinates of $x$ and of $u$.
\end{lemma}
\begin{proof}
Every element of $U(F)$ maps $e_b$ into $e_b+A_{b-1}$; so for $\alpha\in A_a$ the coefficient of $e_a$ in
$x\alpha$ is that in $\alpha$. Let $v\in Y_a\subseteq A_a^\perp$. By Lemma~\ref{pbv:rows}(i), $uv=v+\alpha$ with
$\alpha\in A_a$, whose coefficient at $e_a$ is $\mathsf r_a(u)(v)$; and $xuv=xv+x\alpha$. Taking
coefficients at $e_a$ gives the first claim. If $\mathsf r_a(u)=0$, then by Lemma~\ref{pbv:rows}(ii)
$ue_a^*=e_a^*+y$ with $y\in A_a$ of coefficient $\mathsf y_a(u)$ at $e_a$, and $xue_a^*=xe_a^*+xy$; take
coefficients at $e_a$. The last sentence follows from \eqref{pbv:eq-coord} and the order
\eqref{pbv:eq-order}: $u$ low for $\rho_a$ means $\mathsf r_d(u)=0$, $\mathsf y_d(u)\in\R$ for $d>a$;
low for $t_a$ means this and $\mathsf r_a(u)=0$.
\end{proof}

\begin{lemma}[Eichler transformations]\label{pbv:eichler}
Let $1\le a\le\varrho$, $w\in Y_a$ \textup{(}a vector of $Y_F$\textup{)} and $q\in F$ with $q+\aubar q=
-h(w,w)$. Then
\[
\pbvE_a(w,q)\colon v\longmapsto v+h(v,w)\,e_a-h(v,e_a)\,w+q\,h(v,e_a)\,e_a
\]
lies in $U(F)$, and: $\mathsf r_a(\pbvE_a(w,q))=h(\cdot,w)|_{Y_a}$, $\mathsf y_a(\pbvE_a(w,q))=q$, and
$\mathsf r_d(\pbvE_a(w,q))=0$, $\mathsf y_d(\pbvE_a(w,q))=0$ for $d\ne a$. Every linear form on $Y_a$ is
$h(\cdot,w)$ for exactly one $w\in Y_a$.
\end{lemma}
\begin{proof}
Write $\gamma=\pbvE_a(w,q)$, and for $v\in Y_F$ put $\alpha=h(v,w)$, $\beta=h(v,e_a)$, so $\gamma v=v+
\alpha e_a-\beta w+q\beta e_a$. Using $h(e_a,e_a)=h(e_a,w)=0$,
\[
h(\gamma v,\gamma v)-h(v,v)=(\aubar\alpha\beta+\alpha\aubar\beta)-(\aubar\beta\alpha+\beta\aubar\alpha)+
(\aubar q+q)\beta\aubar\beta+\beta\aubar\beta\,h(w,w)=0 .
\]
By polarisation ($2\ne0$ in $F$) $\gamma$ preserves $h$, so it is injective, hence an isometry of $Y_F$. For
$v\in A^\perp$ we have $\beta=0$, so $(\gamma-1)v=\alpha e_a\in A$; for $v=e_b$, $\alpha=h(e_b,w)$ vanishes
if $b\le a$ (as $w\in Y_a$), so $(\gamma-1)e_b\in Fe_a\subseteq A_{b-1}$ for $b>a$ and $=0$ for $b\le a$.
Thus $\gamma$ preserves the flag, acts trivially on its graded pieces and on $A^\perp/A$: $\gamma\in U(F)$.
For $v\in Y_a$, $\beta=0$ and the coefficient of $e_a$ in $\gamma v$ is $\alpha=h(v,w)$; and $\gamma e_a^*=e_a^*-w+
qe_a$ ($\alpha=h(e_a^*,w)=0$, $\beta=1$), with coefficient $q$ at $e_a$. Let $d\ne a$. The coefficient of $e_d$
in $(\gamma-1)v$ is $-\beta\,h(w,e_d^*)$. If $d<a$ it is $0$, as $w\in Y_a$ has no component at $e_d$. If $d>a$
and $v\in Y_d$ or $v=e_d^*$, then $\beta=h(v,e_a)=0$. So $\mathsf r_d(\gamma)=0$ and $\mathsf y_d(\gamma)=0$.
The last claim holds because $h$ is non-degenerate on $Y_a$ and conjugate-linear in the second variable.
\end{proof}

\subsubsection{A box lemma}\label{pbv:ss-box}

\begin{lemma}\label{pbv:box}
Let $G$ be a group, $\Gamma\subseteq G$ a subgroup, and $c_i\colon G\to V_i$ $(1\le i\le m)$ maps into abelian
groups with $c_i(1)=0$; call $g$ \emph{low for $i$} if $c_j(g)=0$ for all $j<i$. Let $L_i\subseteq V_i$ be
subgroups and $F_i\subseteq V_i$ subsets such that for every $v\in V_i$ exactly one $l\in L_i$ has $l+v\in
F_i$. Assume:
\begin{enumerate}
\item[(T)] if $\gamma\in\Gamma$ is low for $i$, then $c_i(x\gamma)=c_i(\gamma)+c_i(x)$ for all $x\in G$;
\item[(M)] if $\gamma\in\Gamma$ is low for $i$, then $c_i(\gamma)\in L_i$;
\item[(A)] for every $i$ and $l\in L_i$ there is $\gamma\in\Gamma$, low for $i$, with $c_i(\gamma)=l$;
\item[(K)] if $\delta\in\Gamma$ and $c_i(\delta)=0$ for all $i$, then $\delta=1$.
\end{enumerate}
Then for every $x\in G$ there is exactly one $\gamma\in\Gamma$ with $c_i(x\gamma)\in F_i$ for all $i$.
\end{lemma}
\begin{proof}
\emph{Existence.} Put $\gamma_0=1$. Given $\gamma_{i-1}\in\Gamma$ with $c_j(x\gamma_{i-1})\in F_j$ for $j<i$,
choose $l\in L_i$ with $l+c_i(x\gamma_{i-1})\in F_i$ and, by (A), $\delta\in\Gamma$ low for $i$ with
$c_i(\delta)=l$, and put $\gamma_i=\gamma_{i-1}\delta$. For $j\le i$, $\delta$ is low for $j$, so by (T)
$c_j(x\gamma_i)=c_j(\delta)+c_j(x\gamma_{i-1})$; this is $c_j(x\gamma_{i-1})\in F_j$ for $j<i$ (as
$c_j(\delta)=0$) and $l+c_i(x\gamma_{i-1})\in F_i$ for $j=i$. Then $\gamma=\gamma_m$ works.

\emph{Uniqueness.} Let $\gamma,\gamma'$ both work and $\delta=\gamma^{-1}\gamma'$. Suppose $c_j(\delta)=0$ for
$j<i$. By (T), $c_i(x\gamma')=c_i(\delta)+c_i(x\gamma)$, and by (M) $c_i(\delta)\in L_i$; as both
$0+c_i(x\gamma)$ and $c_i(\delta)+c_i(x\gamma)$ lie in $F_i$, uniqueness of $l$ gives $c_i(\delta)=0$. By
induction all $c_i(\delta)$ vanish, and $\delta=1$ by (K).
\end{proof}

\begin{corollary}\label{pbv:bij}
For every $\tau$, $c_\tau=(\rho_a,t_a)_a\colon U_\tau\to V$ is bijective.
\end{corollary}
\begin{proof}
Apply Lemma~\ref{pbv:box} to $G=\Gamma=U_\tau$, the coordinates in the order \eqref{pbv:eq-order}, $L_i=V_i$ and
$F_i=\{z_i\}$ for a given $z\in V$, and $x=1$; it gives exactly one $\gamma$ with $c_\tau(\gamma)=z$. (T) is
Lemma~\ref{pbv:add}; (M) is empty; (K) is Lemma~\ref{pbv:rows}(iii). (A): for $\rho_a$ and a target
$(\rho_j)_j\in\C^{m_a}$ take $\gamma=\pbvE_a(w,-\frac12h(w,w))$ with $w\in Y_a\otimes_{E,\tau}\C$ the unique vector such that $h(v_j,w)=
\lambda_{\tau,a}\lambda_{\tau,j}^{-1}\rho_j$ for all $j$ (it exists and is unique by the last claim of Lemma~\ref{pbv:eichler}, as $(v_j)$ is a basis of $Y_a$); for $t_a$ and a target
$t\in\R$ take $\gamma=\pbvE_a(0,i\lambda_{\tau,a}^2t)$. By Lemma~\ref{pbv:eichler} and \eqref{pbv:eq-coord}
these have the required coordinate, and all their rows other than the row $a$ vanish; the second has
also $\rho_a=0$. So they are low for the coordinate in question.
\end{proof}

\begin{lemma}[The measure of \textup{(Cv3)} is a Haar measure]\label{pbv:haar}
For every $\tau$ and every $x\in U_\tau$ the left translation $u\mapsto xu$ preserves $du_\tau$. So $du_\tau$
is a Haar measure of $U_\tau$ (a nilpotent Lie group, hence unimodular), and it is the Haar measure of
{\cite[Lemma~4.1]{P52}}\,\textup{(c)}: the two descriptions of $du_\tau$ in \textup{(Cv3)} agree.
\end{lemma}
\begin{proof}
Let $x,u\in U_\tau$ and $1\le a\le\varrho$. As $u$ preserves $A_a^\perp=A_a\oplus Y_a$ and $A_a$, for $v\in Y_a$
we have $uv=\bar u_av+\alpha$ with $\bar u_av\in Y_a$, $\alpha\in A_a$, where $\bar u_a$ is the map induced by
$u$ on $A_a^\perp/A_a\cong Y_a$, and the coefficient of $e_a$ in $\alpha$ is $\mathsf r_a(u)(v)$. Since $x$ maps
$e_b$ into $e_b+A_{b-1}$, the coefficient of $e_a$ in $x\alpha$ is that of $\alpha$; so
\begin{equation}\label{pbv:eq-left1}
\mathsf r_a(xu)(v)=\mathsf r_a(u)(v)+\mathsf r_a(x)(\bar u_av)\qquad(v\in Y_a).
\end{equation}
Next, $ue_a^*\in A_{a-1}^\perp=A_a\oplus Y_a\oplus\C e_a^*$, and its coefficient at $e_a^*$ is $1$, because
$h(ue_a^*,e_a)=h(e_a^*,u^{-1}e_a)=1$ ($u^{-1}e_a\in e_a+A_{a-1}$ and $A_{a-1}\perp e_a^*$). So $ue_a^*=e_a^*+y'+
\alpha'$ with $y'\in Y_a$ and $\alpha'\in A_a$ of coefficient $\mathsf y_a(u)$ at $e_a$, and taking the
coefficient of $e_a$ in $xue_a^*=xe_a^*+xy'+x\alpha'$,
\begin{equation}\label{pbv:eq-left2}
\mathsf y_a(xu)=\mathsf y_a(u)+\mathsf y_a(x)+\mathsf r_a(x)(y').
\end{equation}
We claim that $\bar u_a$ is a function of the coordinates of $u$ that precede $\rho_a$ in the order
\eqref{pbv:eq-order}, and $y'$ a function of those that precede $t_a$. By the proof of
Lemma~\ref{pbv:box} (existence, with $x=1$) and of Corollary~\ref{pbv:bij}, $u=\delta_1\delta_2\cdots\delta_{2\varrho}$,
where $\delta_i$ is the Eichler transformation $\pbvE_d(w,q)$ of the floor $d$ of the $i$-th coordinate, and
is determined by the coordinates of $u$ up to the $i$-th. For $v\in A_a^\perp$ and $d\le a$, $h(v,e_d)=0$ and
$(\pbvE_d(w,q)-1)v=h(v,w)e_d\in A_a$; so the factors of floor $d\le a$ act trivially on $A_a^\perp/A_a$, and
$\bar u_a$ is induced by the product of the factors of floor $d>a$, which come first: this is the first
claim. For $v\in Y_a$, the isometry $u$ gives $0=h(e_a^*,v)=h(ue_a^*,uv)=h(e_a^*+y'+\alpha',\bar u_av+\alpha)=
\aubar{\mathsf r_a(u)(v)}+h(y',\bar u_av)$ ($A_a$ is isotropic and orthogonal to $Y_a$, and $h(e_a^*,\alpha)$ is
the conjugate of the coefficient of $e_a$ in $\alpha$). As $\bar u_a$ is bijective and $h$ is non-degenerate
on $Y_a$, $y'$ is determined by $\bar u_a$ and $\mathsf r_a(u)$: the second claim.

By \eqref{pbv:eq-coord}, \eqref{pbv:eq-left1} and \eqref{pbv:eq-left2}, in the coordinates $c_\tau$ the map
$u\mapsto xu$ is therefore $c_i\mapsto c_i+\phi_i$, where $\phi_i$ is a Borel function of $x$ and of the coordinates
$c_j$, $j<i$ (for $\rho_a$ the scaled values of $\mathsf r_a(x)\circ\bar u_a$; for $t_a$, $t_a(x)+\lambda_{\tau,a}^{-2}
\Iml\mathsf r_a(x)(y')$). Such a map preserves Lebesgue measure of $V$ (Fubini, integrating the last
coordinate first). So $du_\tau$ is left invariant. Finally, in the normalised basis the Gram matrix
$\mathsf J_1$ is a signed permutation matrix exchanging $\tilde e_a$ and $\tilde e_a^*$, and $u^{-1}=\mathsf J_1
u^*\mathsf J_1$; hence the entries of $u^{-1}$ in the column $\tilde e_a^*$ and the rows $\tilde v_j$ are, up to
signs, conjugates and a permutation of $j$, the entries $\rho_a(u)_j$, and $\Iml(u^{-1})_{\tilde e_a,\tilde e_a^*}=
-t_a(u)$. So the Gaussian $G_\varrho$ of {\cite[\S4.1]{P52}} satisfies $G_\varrho(u^{-1})=\exp(-\pi|c_\tau(u)|^2)$,
and $\int_{U_\tau}G_\varrho(u^{-1})\,du_\tau=\int_Ve^{-\pi|c|^2}dc=1$. The Haar measure of {\cite[Lemma~4.1\,(c)]{P52}}
has the same integral ({\cite[Lemma~4.1\,(b)]{P52}}, and invariance under $u\mapsto u^{-1}$); two Haar measures
with the same integral of a positive function coincide.
\end{proof}

\subsubsection{The lattices of a class}\label{pbv:ss-lattices}

Let $L\subseteq Y_0$ be the lattice of \S\ref{pb:ss-setting}: $L=\bigoplus_{a\le\varrho}(\auO e_a\oplus\auO
e_a^*)\perp L^{\mathrm{mid}}$, where $L^{\mathrm{mid}}\subseteq X_0\perp Ez$ is spanned by $x_1,\dots,x_{k-1}$,
$x_k/\aubar{\pi_N}+z/N$ and $z$ (proof of Lemma~\ref{pb:KR}); $L$ is integral and unimodular. For a class
$c$ with representative $h_c\in H(\auA_f)$ let $L_c=Y_0\cap h_c(L\otimes\pbvO)$. As $h_c$ fixes the $e_a$,
the $e_a^*$ and $z$ and preserves $X_0$,
\[
L_c=\bigoplus_{a\le\varrho}(\auO e_a\oplus\auO e_a^*)\perp L_c^{\mathrm{mid}},\qquad L_c^{\mathrm{mid}}=
(X_0\perp Ez)\cap h_c(L^{\mathrm{mid}}\otimes\pbvO),
\]
an integral unimodular lattice. For $1\le a\le\varrho$ put $L_{c,a}=L_c\cap Y_a=\bigoplus_{b>a}(\auO e_b
\oplus\auO e_b^*)\perp L_c^{\mathrm{mid}}$; it is an orthogonal direct summand of $L_c$, hence unimodular,
and we let
\[
L_{c,a}^*=\{\ell\in\operatorname{Hom}_E(Y_a,E):\ \ell(L_{c,a})\subseteq\auO\}=\{h(\cdot,w)|_{Y_a}:\ w\in
L_{c,a}\}
\]
(the second description because $L_{c,a}$ is unimodular). Finally $\auO^-=\{\sigma\in\auO:\aubar\sigma=
-\sigma\}$.

\begin{lemma}[The congruence group]\label{pbv:Uc}
$U_c=\{\gamma\in U(E^+):\ (\gamma-1)L_c\subseteq M L_c\}$.
\end{lemma}
\begin{proof}
By definition $U_c=U(E^+)\cap h_cJh_c^{-1}$, and $h_cJh_c^{-1}$ is the principal congruence subgroup of
level $\pi_M$ of the stabiliser of $h_c(L\otimes\pbvO)$. At a place $v$ of $E^+$ not below a chosen place,
the condition is $\gamma(L_c\otimes\auO_v)=L_c\otimes\auO_v$, and $M$ is a unit at $v$. At a place $v$ below a
chosen place $w$ (above $p_0$ or $p_1$; every place of $E$ above $p_0p_1$ is a chosen place or the
conjugate of one, and $M$ has valuation $1$ at all of them), $G(E^+_v)=\auGL_r(E_w)$, and the condition
$\equiv1$ modulo $w$ at $w$ is equivalent, for an isometry, to $\equiv1$ modulo $\aubar w$ at $\aubar w$ (the
component at $\aubar w$ is the inverse adjoint); together, $(\gamma-1)(L_c\otimes\auO_v)\subseteq M(L_c
\otimes\auO_v)$. Conversely, if $(\gamma-1)L_c\subseteq ML_c$ for $\gamma\in U(E^+)$, then $\gamma L_c\subseteq L_c$
and, as $\gamma-1$ is nilpotent, $\gamma^{-1}-1=\sum_{i\ge1}(1-\gamma)^i$ also maps $L_c$ into $ML_c$; so
$\gamma$ stabilises $L_c$ and is $\equiv1$ modulo $M$, that is, $\gamma\in h_cJh_c^{-1}$.
\end{proof}

\begin{lemma}[Rows of the congruence group]\label{pbv:latt}
Let $1\le a\le\varrho$.
\begin{enumerate}
\item[(M1)] For $\gamma\in U_c$, $M^{-1}\mathsf r_a(\gamma)\in L_{c,a}^*$.
\item[(M2)] If $\gamma\in U_c$, $\mathsf r_d(\gamma)=0$ and $\mathsf y_d(\gamma)\in E^+$ for $d>a$, and $\mathsf r_a(\gamma)=0$,
then $\mathsf y_a(\gamma)\in M\auO^-$.
\item[(A1)] For every $\ell\in L_{c,a}^*$ there is $\gamma\in U_c$ with $\mathsf r_a(\gamma)=M\ell$ and $\mathsf
r_d(\gamma)=0$, $\mathsf y_d(\gamma)=0$ for all $d\ne a$.
\item[(A2)] For every $\sigma\in\auO^-$ there is $\gamma\in U_c$ with $\mathsf r_a(\gamma)=0$, $\mathsf y_a(\gamma)=M\sigma$
and $\mathsf r_d(\gamma)=0$, $\mathsf y_d(\gamma)=0$ for all $d\ne a$.
\end{enumerate}
\end{lemma}
\begin{proof}
(M1) For $v\in L_{c,a}$, $(\gamma-1)v\in ML_c$ by Lemma~\ref{pbv:Uc}, and $e_a^*\in L_c$; as $L_c$ is
integral, $\mathsf r_a(\gamma)(v)=h((\gamma-1)v,e_a^*)\in M\auO$. (M2) Similarly $\mathsf y_a(\gamma)=
h((\gamma-1)e_a^*,e_a^*)\in M\auO$ (here $h(e_a^*,e_a^*)=0$), and $\mathsf y_a(\gamma)+\aubar{\mathsf y_a
(\gamma)}=0$ by Lemma~\ref{pbv:rows}(ii) over $F=E$.

(A1), (A2) The field $E$ contains a primitive sixth root of unity $\zeta_6$ ($6\mid5040$), and $\tau_0=\zeta_6
\in\auO$ satisfies $\tau_0+\aubar{\tau_0}=\zeta_6+\zeta_6^{-1}=1$ (as $\zeta_6^2-\zeta_6+1=0$). Let
$w\in L_{c,a}$ and $\sigma\in\auO^-$, and put
\[
\gamma=\pbvE_a(Mw,q),\qquad q=-M^2\tau_0\,h(w,w)+M\sigma .
\]
As $h(w,w)\in E^+$ and $M\in\Z$, $q+\aubar q=-M^2h(w,w)=-h(Mw,Mw)$, so $\gamma\in U(E^+)$ by
Lemma~\ref{pbv:eichler}, with $\mathsf r_a(\gamma)=Mh(\cdot,w)$, $\mathsf y_a(\gamma)=q$ and $\mathsf
r_d(\gamma)=\mathsf y_d(\gamma)=0$ for $d\ne a$. Moreover
\[
M^{-1}(\gamma-1)v=h(v,w)\,e_a-h(v,e_a)\,w+\bigl(-M\tau_0h(w,w)+\sigma\bigr)h(v,e_a)\,e_a ,
\]
which lies in $L_c$ for $v\in L_c$, since $w,e_a\in L_c$, $L_c$ is integral and $-M\tau_0h(w,w)+\sigma\in
\auO$. By Lemma~\ref{pbv:Uc}, $\gamma\in U_c$. For (A1) take $w$ with $\ell=h(\cdot,w)$ and $\sigma=0$
(then $\mathsf r_a(\gamma)=M\ell$); for (A2) take $w=0$ (then $q=M\sigma$ and $\mathsf r_a(\gamma)=0$).
\end{proof}

Let $\iota_c\colon U_c\to U_\infty=\prod_\tau U_\tau$ be the diagonal embedding. By \eqref{pbv:eq-coord}, for
$\gamma\in U_c$ the coordinates of $\iota_c(\gamma)$ are $\rho_a(\gamma)=\bigl((\lambda_{\tau,j}/\lambda_{\tau,a})
\,\tau(\mathsf r_a(\gamma)(v_j))\bigr)_{\tau,j}$ and $t_a(\gamma)=\bigl(\lambda_{\tau,a}^{-2}\Iml\tau(\mathsf y_a
(\gamma))\bigr)_\tau$. Accordingly define the subgroups
\begin{align*}
\Lambda_a&=\Bigl\{\Bigl(\frac{M\lambda_{\tau,j}}{\lambda_{\tau,a}}\,\tau(\ell(v_j))\Bigr)_{\tau\in\auPl,\,j\le m_a}:\
\ell\in L_{c,a}^*\Bigr\}\subseteq\prod_\tau\C^{m_a},\\
\Lambda'_a&=\Bigl\{\Bigl(\frac M{\lambda_{\tau,a}^2}\,\Iml\tau(\sigma)\Bigr)_{\tau\in\auPl}:\ \sigma\in\auO^-
\Bigr\}\subseteq\R^{\auPl}.
\end{align*}

\begin{lemma}[An adelic covolume]\label{pbv:adelic}
Give $\auA_{E,f}$ the Haar measure with $\pbvol(\pbvO)=(2/\aurd)^f$ and $E_\infty=\prod_{\tau\in\auPl}\C$ the
Lebesgue measure $\prod_\tau d\Rel\,d\Iml$. Let $m\ge1$ and let $\widehat W\subseteq\auA_{E,f}^m$ be an open
compact subgroup. Then $E^m\cap\widehat W$, embedded in $E_\infty^m$ by $(\tau)_\tau$, is a lattice of
covolume $\pbvol(\widehat W)^{-1}$.
\end{lemma}
\begin{proof}
For $\widehat W=\pbvO^m$ the group is $\auO^m$. The image of $\auO$ in $E_\infty$ has
covolume $2^{-f}|D_E|^{1/2}$ for the Lebesgue measure (Lemma~\ref{pbc:tb-covol}; see also \cite[Chapter~I, (5.1) and Proposition~(5.2)]{Neukirch}); as
$|D_E|=\aurd^{2f}$ this is $(\aurd/2)^f=\pbvol(\pbvO)^{-1}$, and the claim follows for $\pbvO^m$. In general
put $\widehat W_0=\widehat W\cap\pbvO^m$, open of finite index in both. As $E$ is dense in $\auA_{E,f}$,
every coset of $\widehat W_0$ in $\widehat W$ (an open set) meets $E^m$, so $[E^m\cap\widehat W:E^m\cap
\widehat W_0]=[\widehat W:\widehat W_0]$, and likewise with $\pbvO^m$ in place of $\widehat W$. Hence
$E^m\cap\widehat W$ is a lattice and
\[
\pbvcov(E^m\cap\widehat W)=\pbvcov(\auO^m)\,\frac{[\pbvO^m:\widehat W_0]}{[\widehat W:\widehat W_0]}=
\pbvol(\pbvO)^{-m}\,\frac{\pbvol(\pbvO^m)}{\pbvol(\widehat W)}=\pbvol(\widehat W)^{-1}.\qedhere
\]
\end{proof}

\begin{lemma}[The lattices of the rows]\label{pbv:vol-a}
$\Lambda_a$ is a lattice in $\prod_\tau\C^{m_a}$ of covolume
$\displaystyle\Bigl(\prod_\tau\frac{M^2}{\lambda_{\tau,a}^2}\Bigr)^{m_a}\pbvol(\pbvO)^{-m_a}$.
\end{lemma}
\begin{proof}
Identify $\ell$ with $y=(\ell(v_j))_j\in E^{m_a}$. Let $B\in\auGL_{m_a}(\auA_{E,f})$ be the matrix whose
columns are the coordinates, in the basis $(v_j)$, of an $\pbvO$-basis of $L_{c,a}\otimes\pbvO$. Since
$\auO=E\cap\pbvO$, $\ell\in L_{c,a}^*$ if and only if $yB\in\pbvO^{m_a}$, so $L_{c,a}^*$ corresponds to
$E^{m_a}\cap\widehat W$ with $\widehat W=\pbvO^{m_a}B^{-1}$, and $\pbvol(\widehat W)=|\det B|_f^{-1}\pbvol(
\pbvO)^{m_a}$. One may take $B=\operatorname{diag}(1,h_cB^{\mathrm{mid}})$, where $1$ is the identity on the
coordinates $e_b,e_b^*$ and $B^{\mathrm{mid}}$ has the columns $x_1,\dots,x_{k-1}$, $x_k/\aubar{\pi_N}+z/N$,
$z$ in the basis $x_1,\dots,x_k,z$; it is triangular with diagonal $1,\dots,1,\aubar{\pi_N}^{\,-1},1$. Now
$|\det h_c|_f=1$, because $\det h_c\cdot\aubar{\det h_c}=1$ and $|\aubar x|_f=|x|_f$; and $|\aubar{\pi_N}^{\,-1}|_f
=|\operatorname N_{E/\Q}(\pi_N)|=N^f$ by the product formula ($\pi_N\aubar{\pi_N}=N$ gives $\operatorname N_{E/\Q}
(\pi_N)^2=N^{2f}$). So $\pbvol(\widehat W)=N^{-f}\pbvol(\pbvO)^{m_a}$, and by Lemma~\ref{pbv:adelic} the
image of $L_{c,a}^*$ in $E_\infty^{m_a}$ is a lattice of covolume $N^f\pbvol(\pbvO)^{-m_a}$. The map to
$\Lambda_a$ multiplies the coordinate $(\tau,j)$ by the non-zero real number $M\lambda_{\tau,j}/\lambda_{\tau,a}$,
so it multiplies volumes by $\prod_{\tau,j}(M\lambda_{\tau,j}/\lambda_{\tau,a})^2$. For each $\tau$,
$\prod_j\lambda_{\tau,j}^2=s_{\tau,\varrho}^2=N^{-1}$: the factors of $e_b$ and $e_b^*$ cancel, those of the
$x_i$ are $1$, and that of $z$ is $s_{\tau,\varrho}^2$. Hence the covolume is $\prod_\tau\bigl((M^2/
\lambda_{\tau,a}^2)^{m_a}N^{-1}\bigr)\cdot N^f\pbvol(\pbvO)^{-m_a}$, as claimed.
\end{proof}

\begin{lemma}[The skew lattices]\label{pbv:vol-skew}
$\Lambda'_a$ is a lattice in $\R^{\auPl}$ of covolume $\displaystyle\prod_\tau\frac M{\lambda_{\tau,a}^2}\cdot
|D_E|^{1/4}=\prod_\tau\frac M{\lambda_{\tau,a}^2}\cdot\aurd^{\,f/2}$.
\end{lemma}
\begin{proof}
$E$ contains $i=\zeta_4$ ($4\mid5040$), a unit with $\aubar i=-i$. So $\sigma\in\auO^-$ if and only if
$\sigma=i\alpha$ with $\alpha\in\auO\cap E^+=\auO_{E^+}$, and $\Iml\tau(i\alpha)=\epsilon_\tau\tau(\alpha)$ with
$\epsilon_\tau=\pm1$ ($\tau(i)=\pm i$, $\tau(\alpha)\in\R$). The restrictions of the $\tau\in\auPl$ to $E^+$ are
the $f$ real embeddings of $E^+$, each once. By Minkowski's theorem for the totally real field $E^+$,
$\{(\tau(\alpha))_\tau:\alpha\in\auO_{E^+}\}$ is a lattice of covolume $|D_{E^+}|^{1/2}$. The tower formula
for discriminants, $|D_E|=D_{E^+}^{2}\,\operatorname N_{E^+/\Q}(\mathfrak d_{E/E^+})$, and the triviality of the
relative discriminant $\mathfrak d_{E/E^+}$ ($E/E^+$ is unramified at all finite places,
{\cite[\S14.1]{P46}}) give $|D_{E^+}|^{1/2}=|D_E|^{1/4}=\aurd^{f/2}$. Scaling the coordinate $\tau$ by
$\epsilon_\tau M\lambda_{\tau,a}^{-2}$ multiplies the covolume by $\prod_\tau M\lambda_{\tau,a}^{-2}$.
\end{proof}

\subsubsection{The value of the covolume}\label{pbv:ss-value}

\begin{theorem}[Lemma~\ref{pb:covol}]\label{pbv:covol}
For every class $c$, and for the Haar measure $du=\prod_\tau du_\tau$ of \textup{(Cv3)},
\[
\nu=\pbvol(U_c\backslash U_\infty)=\prod_{a=1}^{\varrho}\pbvcov(\Lambda_a)\,\pbvcov(\Lambda'_a)=
\frac{M^{f\varrho(2k+2\varrho+1)}}{N^{kf}}\cdot\frac{\aurd^{\,f\varrho/2}}{\pbvol(\pbvO)^{\varrho(\varrho+k)}\,
\Lambda}=\nu_*,
\]
with $\Lambda=\prod_\tau\prod_{a=0}^{\varrho}s_{\tau,a}^{\,2(k+2(\varrho-a))}$ and $\pbvol(\pbvO)=(2/\aurd)^f$. (The constant $\Lambda$ is a positive real number and has no relation to the lattices $\Lambda_a,\Lambda'_a$.)
\end{theorem}
\begin{proof}
\emph{A fundamental domain.} Apply Lemma~\ref{pbv:box} to $G=U_\infty$, $\Gamma=\iota_c(U_c)$, the
$2\varrho$ coordinates $\rho_a=(\rho_a^{(\tau)})_\tau$, $t_a=(t_a^{(\tau)})_\tau$ in the order
\eqref{pbv:eq-order}, with values in $\prod_\tau\C^{m_a}$ and $\R^{\auPl}$, the subgroups $L_{\rho_a}=\Lambda_a$,
$L_{t_a}=\Lambda'_a$, and for $F_i$ a half-open fundamental parallelotope of the lattice $L_i$ (spanned by
a $\Z$-basis; Lemmas~\ref{pbv:vol-a} and \ref{pbv:vol-skew}), which meets every coset of $L_i$ exactly once.
Here: (T) is Lemma~\ref{pbv:add}, applied at each $\tau$ (for $\gamma\in U_c$, being low for a coordinate
at all $\tau$ is the same as the corresponding condition over $E$: $\mathsf r_d(\gamma)=0$ and $\mathsf
y_d(\gamma)\in E^+$, as $\Iml\tau(y)=0$ for all $\tau$ means $y\in E^+$). (M) is Lemma~\ref{pbv:latt}(M1)
for $\rho_a$ (it needs no lowness) and (M2) for $t_a$. (A) is Lemma~\ref{pbv:latt}(A1), (A2): the elements
given there have vanishing rows $d\ne a$, and the element of (A2) has also $\mathsf r_a=0$, so they are low
for $\rho_a$, respectively $t_a$. (K) is Lemma~\ref{pbv:rows}(iii) over $E$ (and $\iota_c$ is injective).
Hence $B=c^{-1}(\prod_iF_i)$, $c=(c_\tau)_\tau$, meets every coset $x\Gamma$ in exactly one point.

\emph{The volume.} By Corollary~\ref{pbv:bij} and the definition of $du$, $du(B)=\prod_i\pbvol(F_i)=
\prod_a\pbvcov(\Lambda_a)\pbvcov(\Lambda'_a)$. The group $U_\infty$ is nilpotent, hence unimodular, and $du$ is a
Haar measure (Lemma~\ref{pbv:haar}), so it is invariant under $u\mapsto u^{-1}$; $B^{-1}$ meets every coset $\Gamma x$ exactly
once, and $\nu=\pbvol(U_c\backslash U_\infty)=du(B^{-1})=du(B)$. (This also shows again that $\nu$ does not
depend on $c$.)

\emph{The product.} By Lemmas~\ref{pbv:vol-a} and \ref{pbv:vol-skew},
\[
\nu=\prod_{a=1}^{\varrho}\Bigl[\Bigl(\prod_\tau\frac{M^2}{\lambda_{\tau,a}^2}\Bigr)^{m_a}\pbvol(\pbvO)^{-m_a}\cdot
\prod_\tau\frac M{\lambda_{\tau,a}^2}\cdot\aurd^{\,f/2}\Bigr].
\]
The power of $M$ is $f\sum_{a=1}^{\varrho}(2m_a+1)=f\sum_{a=1}^\varrho\bigl(2k+3+4(\varrho-a)\bigr)=f\varrho(2k+2\varrho
+1)$; that of $\pbvol(\pbvO)$ is $-\sum_am_a=-\varrho(k+1)-\varrho(\varrho-1)=-\varrho(\varrho+k)$; that of $\aurd$ is
$f\varrho/2$. For the scales, $\lambda_{\tau,a}=s_{\tau,a-1}$ and $m_a+1=k+2(\varrho-(a-1))$, so
\[
\prod_{a=1}^{\varrho}\lambda_{\tau,a}^{-2(m_a+1)}=\prod_{a'=0}^{\varrho-1}s_{\tau,a'}^{-2(k+2(\varrho-a'))}=
s_{\tau,\varrho}^{2k}\prod_{a'=0}^{\varrho}s_{\tau,a'}^{-2(k+2(\varrho-a'))}=N^{-k}\prod_{a'=0}^{\varrho}
s_{\tau,a'}^{-2(k+2(\varrho-a'))},
\]
and the product over $\tau$ is $N^{-kf}\Lambda^{-1}$. Collecting the factors gives $\nu_*$.
\end{proof}

\begin{remark}[Structure behind the computation]\label{pbv:rem-tower}
The coordinates $(\rho_a,t_a)$ are those of the tower $U=N_1\ltimes(N_2\ltimes(\dots\ltimes N_\varrho))$, where
$N_a$ is the Heisenberg group of the floor $a$ (the unipotent radical of the stabiliser of $e_a$ in
$\auU(A_{a-1}^\perp/A_{a-1})$), with centre the skew line $\{e_a^*\mapsto e_a^*+te_a\}$ and abelian quotient
$Y_a$; by Lemma~\ref{pbv:rows}(i) the elements low for $\rho_a$ act trivially on $A_a^\perp/A_a$, that is, they lie in $N_1\cdots N_a$, and
Lemma~\ref{pbv:eichler} gives the elements of $N_a$. Thus $\nu$ is the product over the floors of the
covolume of a lattice in the vector space $Y_a\otimes_{E^+}\R$ (the lattice $M L_{c,a}^*$, which carries the
level $M$, the scales $\lambda_{\tau,j}/\lambda_{\tau,a}$ of the normalised basis and the factor $N^f$ of the
non-trivial basis of $L^{\mathrm{mid}}$ at $N$) and of the covolume of the skew-hermitian lattice $M\auO^-$
of the centre (which carries the different through $|D_E|^{1/4}$, i.e.\ through $\varpi$: $|D_E|=\prod_\tau
|\tau(\varpi)|^2$). We did not need the tower explicitly.
\end{remark}

\subsubsection{The three forms of the constant}\label{pbv:ss-uses}

The constant appears in three places, written in three ways; they are the same number, so that
Theorem~\ref{pbv:covol} gives $\nu$ in each of them.
\begin{enumerate}
\item[(1)] Before Lemma~\ref{pb:covol}: $\nu_*=M^{f\varrho(2k+2\varrho+1)}N^{-kf}\aurd^{\,f\varrho/2}\pbvol(\pbvO)^{-\varrho
(\varrho+k)}\Lambda^{-1}$ with $\Lambda=\prod_\tau\prod_{a=0}^\varrho s_{\tau,a}^{2(k+2(\varrho-a))}$, $s_{\tau,a}=
\kappa_\tau^{\varrho-a}/\sqrt N$, $\pbvol(\pbvO)=(2/\aurd)^f$. This is the value of Theorem~\ref{pbv:covol}.
\item[(2)] In \eqref{pbm:eq-consts} (used in Theorem~\ref{pbm:L2} and Lemma~\ref{pbm:const}):
$\nu_*=\#\pbUb\,N^{-kf}\aurd^{\,f\varrho/2}\mathsf v^{-\varrho(\varrho+k)}\Lambda_*^{-1}$ with $\mathsf v=(2/\aurd)^f$ and
$\Lambda_*=\prod_\tau\prod_as_{\tau,a}^{k+2(\varrho-a)}$, where in \S\ref{pbm:sec} the letter $s_{\tau,a}$ denotes
$\kappa_\tau^{2(\varrho-a)}/N$, the \emph{square} of the scale \eqref{pbe:eq-scale}. So $\Lambda_*=\Lambda$, and
$\#\pbUb=M^{f\varrho(2k+2\varrho+1)}$ by \eqref{pb:eq-orders}: the two expressions coincide.
\item[(3)] In Remark~\ref{pbs:rem-nu} (used in Lemmas~\ref{pbs:O2} and \ref{pbs:open-arith}):
$\nu_{\mathrm{src}}=(M^{2d+\varrho}/N^k)^f\pbvol(\pbvO)^{-d}|D_E|^{\varrho/4}/\Lambda$ with $d=\varrho(\varrho+k)$ and
$\Lambda=\prod_\tau\prod_a(\kappa_\tau^{2(\varrho-a)}/N)^{k+2(\varrho-a)}$, which is the $\Lambda$ of (1). Here
$2d+\varrho=\varrho(2k+2\varrho+1)$ and $|D_E|^{\varrho/4}=\aurd^{\,2f\varrho/4}=\aurd^{\,f\varrho/2}$, so $\nu_{\mathrm{src}}=\nu_*$; and
the identity used in Step~2 of Lemma~\ref{pbs:O2},
\[
\nu_{\mathrm{src}}\Lambda=\#\pbUb\,N^{-kf}\pbvol(\pbvO)^{-\varrho(\varrho+k)}|D_E|^{\varrho/4},
\]
is the definition multiplied by $\Lambda$.
\end{enumerate}
Thus $\nu=\nu_*=\nu_{\mathrm{src}}$ for the measure (Cv3), which is the measure for which the local inputs
({\cite[Theorem~2.17]{P52}} and~Theorem~\ref{pbo:O1}) are proved. In the factorisation of Theorem~\ref{pbv:covol},
$\#\pbUb=\prod_aM^{f(2m_a+1)}$ collects the level $M$ in the $2m_a+1$ real dimensions of each floor, $N^{-kf}$ comes from the scale
$s_{\tau,\varrho}^{2k}$ of $\tilde z$ left over in $\Lambda$ (the factor $N^f$ of $L^{\mathrm{mid}}$ at the
primes above $N$ cancels the scale $s_{\tau,\varrho}^2$ of $\tilde z$ in each row), $\pbvol(\pbvO)^{-\varrho(\varrho+k)}$
from the $\sum_am_a$ coordinates in $E$, and $|D_E|^{\varrho/4}$ from the $\varrho$ skew coordinates.

\subsubsection{Four short steps of the Spanning Theorem}\label{pbv:ss-weak}

\paragraph*{(i) The rearrangement of the Heisenberg measures ({\cite[Lemma~2.41]{P52}})}
The last step of the proof of {\cite[Lemma~2.41]{P52}} is the identity
\begin{equation}\label{pbv:eq-rearr}
\eta_j*(\eta_0*\dots*\eta_{j-1})=\eta_0*\dots*\eta_{j-1}*\eta_j ,
\end{equation}
where $(\mu*\mu')(\Phi)=\iint\Phi(gg')\,d\mu(g)\,d\mu'(g')$. It follows from (N3), without any statement
about Haar measures of products. Put $\nu'=\eta_0*\dots*\eta_{j-1}$. For $n\in\pfrN_j\subseteq\pfrG^{(j)}$
and $a<j$, conjugation by $n$ maps $\pfrN_a$ to itself by $n\,n_{\omega,\beta}\,n^{-1}=n_{n\omega,\beta}$ and
preserves $\eta_a$: in the coordinates $(\omega,\beta)$ it is $\omega\mapsto n\omega$, a linear automorphism of the
$\omega$-space (the part of $V_{a+1}$ orthogonal to $f_a,f_a^*$, which $n$ preserves) whose modulus is
$|\det n|$, equal to $1$ because $\det n\cdot\aubar{\det n}=1$, with $\beta$ unchanged. Since $x\mapsto nxn^{-1}$ is
a group automorphism, it maps $\nu'$, the image of $\eta_0\otimes\dots\otimes\eta_{j-1}$ under the product
map, to itself. Hence, for $\Phi\ge0$ measurable, by Tonelli (all measures are $\sigma$-finite),
\[
\begin{aligned}\int\Phi\,d(\eta_j*\nu')&=\iint\Phi(nx)\,d\nu'(x)\,d\eta_j(n)=\iint\Phi\bigl((nxn^{-1})\,n\bigr)\,d\nu'(x)\,
d\eta_j(n)\\&=\iint\Phi(xn)\,d\nu'(x)\,d\eta_j(n)=\int\Phi\,d(\nu'*\eta_j),\end{aligned}
\]
which is \eqref{pbv:eq-rearr} (associativity of convolution gives the right side).

\paragraph*{(ii) The basis \textup{(L1)} at $T$ ({\cite[Lemma~5.13]{P52}})}
The explicit basis used in the proof of {\cite[Lemma~5.13]{P52}} is as follows (forms linear in the first variable, (Cv1); $c,\tau\in\pffOT$ with
$c\aubar c=-1$, $\tau+\aubar\tau=1$ as in the proof; note $\aubar c^{\,-1}=-c$).
\begin{itemize}
\item From $v=\aubar{\pi_N}^{\,-1}x_{k_0}+N^{-1}z$ and $z$ ($h(v,v)=0$, $h(v,z)=-1$, $h(z,z)=-N$): the pair
$p_0=-v$, $q_0=z-N\tau v$, with $h(p_0,q_0)=1$, $h(q_0,q_0)=-N+N\aubar\tau+N\tau=0$; and
$\bigl(\begin{smallmatrix}p_0&q_0\end{smallmatrix}\bigr)=\bigl(\begin{smallmatrix}v&z\end{smallmatrix}\bigr)
\bigl(\begin{smallmatrix}-1&-N\tau\\0&1\end{smallmatrix}\bigr)$, determinant $-1$; $z=q_0-N\tau p_0$.
\item From two orthonormal vectors $x,x'$: $p=x+cx'$, $q=\aubar\tau x-c\tau x'$. Then $h(p,p)=1+c\aubar c=0$,
$h(q,q)=\tau\aubar\tau+c\aubar c\,\tau\aubar\tau=0$, $h(p,q)=\tau-c\aubar c\,\aubar\tau=\tau+\aubar\tau=1$; and
$\bigl(\begin{smallmatrix}p&q\end{smallmatrix}\bigr)=\bigl(\begin{smallmatrix}x&x'\end{smallmatrix}\bigr)
\bigl(\begin{smallmatrix}1&\aubar\tau\\c&-c\tau\end{smallmatrix}\bigr)$, determinant $-c(\tau+\aubar\tau)=-c$, a unit.
\end{itemize}
Take $m'+1$ disjoint pairs $(x_{i},x_{i'})$ among the $x_i$, $i\ne k_0$, and keep the other $\iota$ of them.
The basis $e_1,\dots,e_\varrho;\ p_0,q_0;\ p_1,q_1,\dots,p_{m'+1},q_{m'+1};\ x_i\ (\iota\text{ of them});\
e_1^*,\dots,e_\varrho^*$ has the Gram matrix with blocks $\bigl(\begin{smallmatrix}0&1\\1&0\end{smallmatrix}\bigr)$ on
each pair $(e_a,e_a^*)$, $(p_l,q_l)$ and $1$ on the kept $x_i$; the change of basis from the lattice basis is
block diagonal with the blocks above and identity blocks, integral with determinant $\pm\prod_l c\in
\pffOT^\times$, so its inverse is integral. This is the basis of (L1) (the frame vector $z$ lies in the span
of the pair $p_0,q_0$ with integral coefficients, as displayed; $k+1=\iota+2(m'+1)+2$).

\paragraph*{(iii) The open cell}
The three steps of \S\ref{pbo:sec} used here are proved there: Lemma~\ref{pbo:tube} rests on the gauge $N(u)$ with $\sqrt{N(uu')}\le\sqrt{N(u)}+\sqrt{N(u')}$ and $N(u^{-1})=N(u)$, the scaling of the
balls $B(\rho)$ under $a_t$, and the bound $(1+1/Q)^Q\le e\le3$; Lemma~\ref{pbo:sobolev} on
the one-variable inequality, the recursion for $I_j(w)$ over the letters, and the binomial sum; and
Lemma~\ref{pbo:sizes}(c) contains the explicit polynomial comparison $\|M(Z,W)\|\le15P^6\le P^{12}$ for $P\ge8$.
We add one remark on (c): the Cauchy estimate is applied on the product of the two
operator-norm balls of radius $T^{-6}$ about $Z$ and about $Z^*$ (on which $M$ is holomorphic in $(Z,W)$ and
bounded by $(2T)^{12}$ by (b)), in one complex variable along each direction $H$ with $\|H\|=1$; the real
derivative of $Z\mapsto M(Z,Z^*)$ in the direction $H$ is $\partial_ZM[H]+\partial_WM[H^*]$, each term of norm
at most $(2T)^{12}T^6$, which gives the bound $2(2T)^{12}T^6\le T^{25}$ stated there.

\paragraph*{(iv) The convention for $h$ in \S\ref{pbg:sec}}
The sentence at the beginning of \S\ref{pbg:sec} is correct, with the following argument. The proofs of
\S\ref{pbg:sec} do use forms linear in the \emph{second} variable: the expansion $x=\sum_kR^\#_kh(R_k,x)$ in the
proof of Lemma~\ref{pbg:jacobi} and the formula $\mathfrak p=\aubar{\mathfrak a(L)}\mathfrak a(L')\det q(\ell,\ell')$
in Lemma~\ref{pbg:basic}(c) require it. Given the form $h$ of (Cv1), put $h'(x,y)=h(y,x)=\aubar{h(x,y)}$, which is
linear in the second variable. Then: orthogonality, integrality ($h(x,L)\subseteq\auO$), dual lattices and
unimodularity are the same for $h$ and $h'$, since $\aubar\auO=\auO$; for a tuple $\ell$, $h'(\ell,\ell)$ is the
transpose of $h(\ell,\ell)$, so $\pgnD_{h'}(L)=\pgnD_h(L)$ and $\pgnd_{h'}=\pgnd_h$; mixed determinants satisfy
$\det h'(\ell,\ell')=\aubar{\det h(\ell,\ell')}$, so the ideal $\mathfrak p$ of Lemma~\ref{pbg:basic}(c) is replaced by
$\aubar{\mathfrak p}$, of the same norm, and $\mathfrak p\subseteq\auO$ iff $\aubar{\mathfrak p}\subseteq\auO$; positivity at the
complex places is unchanged; the matrix $Z(\gamma)$ is replaced by its entrywise conjugate, so the minors
$\delta_l$ by $\aubar{\delta_l}$, and $\pbg(\gamma)$, the norms $\auNm\delta_l$ and all indices are unchanged; in
Lemma~\ref{pbg:jacobi} the form $\tilde h$ built from $(h',c)$ is the conjugate of the one built from $(h,\aubar c)$,
and $\auNm c=\auNm\aubar c$ (the lemma is applied with $c=N\in\Z$). Hence every statement of \S\ref{pbg:sec},
being about lattices, indices and absolute norms, holds verbatim for $h$ in convention (Cv1).



% ===== B-bessel-expansion.tex =====
% proofs/B-bessel-expansion.tex
% Bessel functions, the Poincare series, and the expansion of the squares:
% Hypothesis pb:bessel and Theorem pb:amp (X), (L1) of proofs/B-second-moment.tex.
% Labels of this file start with pbe:.
\providecommand{\auO}{\mathcal O}
\providecommand{\auPl}{\Phi}
\providecommand{\auPa}{\Phi_a}
\providecommand{\auS}{\mathcal S}
\providecommand{\auAz}{\mathcal A_0}
\providecommand{\auG}{\overline G}
\providecommand{\auP}{\mathcal P}
\providecommand{\auDX}{D_X}
\providecommand{\auNm}{\operatorname{Nm}}
\providecommand{\auU}{\mathrm U}
\providecommand{\auGL}{\mathrm{GL}}
\providecommand{\auA}{\mathbb A}
\providecommand{\auF}{\mathbb F}
\providecommand{\aurd}{\mathrm{rd}_E}
\providecommand{\aubar}[1]{\overline{#1}}
\providecommand{\auf}[1]{\texttt{#1}}
\providecommand{\panot}[1]{\par\smallskip{\raggedright\noindent\textbf{[Not proved here:} #1\textbf{]}\par}\smallskip}
\providecommand{\pbR}{\overline R_0}
\providecommand{\pbUb}{\overline U}
\providecommand{\pbHb}{\overline H}
\providecommand{\pbvol}{\operatorname{vol}}
\providecommand{\pbDid}{\mathcal D_{\mathrm{id}}}
\providecommand{\pbDmain}{\mathcal D_{\mathrm{main}}}
\providecommand{\pbdeg}{\operatorname{deg}}
\providecommand{\pbrpsi}{r_\psi}
\providecommand{\pbW}{\mathbb W}
\providecommand{\pbg}{\mathsf g}
\providecommand{\pbLM}{\mathrm{LM}}
\providecommand{\pbone}{\mathbf 1}
\providecommand{\pbHS}[1]{\|#1\|}
\providecommand{\pbip}[2]{\langle#1,#2\rangle}
% macros of this file (prefix \pbe)
\providecommand{\pbel}{\ell}
\providecommand{\pbels}{\tilde\ell}
\providecommand{\pbeL}{\mathsf L}
\providecommand{\pbeq}{\mathsf q}
\providecommand{\pbeB}{\mathsf B}
\providecommand{\pbeN}{\mathcal N}
\providecommand{\pbet}{\beta_*}
\providecommand{\pbenp}{\mathsf n_{\mathrm p}}
\providecommand{\pbesz}{\operatorname{size}}
\providecommand{\pbeGs}{\Gamma^*}
\providecommand{\pbeAp}{A^{+}}
\providecommand{\pbeCL}{\mathrm{CL}}

\subsection{Bessel functions, Poincar\'e series and the expansion of the squares}\label{pbe:sec}

This section concerns three statements of \S\ref{pb:sec} that were left without proof there:
Hypothesis~\ref{pb:bessel} (Bessel functions exist and $\auP$ is their Poincar\'e series) and parts
\textup{(X)} and \textup{(L1)} of Theorem~\ref{pb:amp}. We use the notation of \S\ref{pb:sec}:
$n=\varrho+1$, $r=k+2\varrho+1$, the frame $(f_0,\dots,f_\varrho)=(e_1,\dots,e_\varrho,z)$,
$A=\langle e_1,\dots,e_\varrho\rangle$, $U$, $H$, $R_0=U\rtimes H$, the classes $c\in\mathcal C$ with
representatives $h_c\in H(\auA_f)$, the groups $U_c=U(E^+)\cap h_cJh_c^{-1}$, the terms
$\Gamma_{c,s,c'}$ and the amplitude sums \eqref{pb:eq-amp}. Inner products are conjugate-linear in the
first variable.

\emph{What is proved.}
\begin{itemize}
\item[(1)] \emph{Bessel functions} (\S\S\ref{pbe:ss-real}--\ref{pbe:ss-global}).
Hypothesis~\ref{pb:bessel} is proved, with \textup{(B5)} in its exact form and with the
quantitative properties that the later estimates use (Theorem~\ref{pbe:bessel}), \emph{from four
properties \textup{(DS1)--(DS4)} of the discrete series coefficient at a real place}
(Theorem~\ref{pbe:ds}). \textup{(DS1)} is Harish-Chandra's theory; \textup{(DS2)}--\textup{(DS4)}
are quantitative statements about one explicit function (Flensted-Jensen's coefficient of the lowest
$K$-type), namely its decay, the size of the formal degree against the volume of balls, and a lower
bound for one Bessel period. \textup{(DS1)}--\textup{(DS3)} are proved in \S\ref{pds:sec}, the explicit
formula behind \textup{(DS2)} in {\cite[Theorem~3.17]{P49}}, and \textup{(DS4)} in \S\ref{pbp:sec}. Given
them, everything is proved here: convergence of the local period, equivariance, the reproducing
property, the identity ``the norm is a period'', the decay below an invariant majorant, the cost
of the cut in $L^2$, the count of the terms of a Poincar\'e series (Proposition~\ref{pbe:count}),
its boundedness, and the identification with $\auP$.
\item[(2)] \emph{Expansion of the squares, first two steps} (\S\S\ref{pbe:ss-char}--\ref{pbe:ss-cut}).
Vanishing by character (Lemma~\ref{pbe:vanish}), the structure of the agreeing terms that preserve
$A$ (Lemma~\ref{pbe:scalar}), and the cost of the cut (Proposition~\ref{pbe:cut}) are proved in
full, the last one given (1).
\item[(3)] \emph{Pairs of terms} (\S\S\ref{pbe:ss-pillars}--\ref{pbe:ss-XL1}).
The third step of \textup{(X)} and the whole of \textup{(L1)} are \emph{reduced}, by complete
arguments, to three estimates for the terms with $\pbg\ne0$ or with $\gamma A\ne A$: a bound for one
block by sizes of orbits \textup{(S)}, a bound for the discriminant of the lattice common to two
terms \textup{(G)}, and a count of pairs by their common lattice \textup{(C)}, together with a structural property \textup{(R)}
(Theorem~\ref{pbe:pillars}). These estimates are proved in \S\S\ref{ppi:sec}, \ref{pbg:sec},
\ref{pbc:sec}, \ref{pbs:sec}, \ref{pbo:sec}: (R) in Proposition~\ref{ppi:cell}, (G) from (G$'$) of
\S\ref{pbg:sec}, (C) from (C1), (C2) of \S\ref{pbc:sec}, and (S) in \S\ref{pbs:sec} from the open-cell
bound of \S\ref{pbo:sec}. The linear algebra
of links that the reduction needs is proved (Lemma~\ref{pbe:link}).
\end{itemize}
\S\ref{pbe:ss-remains} lists the statements on which this section rests.

\paragraph*{Sizes}
Besides {\cite[(36) and~(38)]{P46}} we use the following consequences of the
definition of the data, all of them in {\cite[Lemma~13.4(S7)]{P46}}: $n=m/2\ge23$ and $12n^3+500\le k$
((S7)(a)); for Theorem~\ref{pbe:pillars}, $\varrho\ge10^5$ and $k\ge1000(2n)^9$ \eqref{ppi:eq-sizes}, and
for \S\ref{pbs:ss-open}, $(200fk\,\aurd)^{1000}\le b$ ((S7)(d)); $r/2\le k<r$ (S2), so that $(r+k)^2\le
4r^2$, and $r\ge100$, $3r^3f\le r^7$, $kf\ge40$ ((S7)(b)); $n,r,f,\sqrt d\le b$, and $\sqrt d\le\aurd$
($\sqrt d$ is the root discriminant of $K\subseteq E$) ((S7)(e)); $b^{20}\le M\le b^{22}$ (S5);
$N\le b^{60n}$; $d_\infty:=\dim W_{\lambda^X}\le b^{2kf}$ ((S7)(f)); and, for \S\ref{pbo:sec},
$8r\le b$ and $(r+1)^4(L_\tau+r+1+2n^2)\le b^5$ with $L_\tau=\max_i|2\varsigma_{\tau,i}|\le b^2$
((S7)(b), (g)). One more inequality concerns the generator $\varpi$ of the
different in \eqref{pb:eq-psi}: for every $\tau\in\auPl$ the real number
\begin{equation}\label{pbe:eq-kappa}
\kappa_\tau:=\frac{M\,\tau(\varpi)}{2\pi i}\qquad\text{satisfies}\qquad1\le|\kappa_\tau|\le b^{23}.
\end{equation}
(It is real because $\varpi$ is purely imaginary. The inequality is false for some purely imaginary generators; it holds for the generator $\varpi=\gamma\Phi'_{q_1}(\zeta)\sqrt{-d}$ of Lemma~\ref{pds:kappa}, which is the one fixed in \eqref{pb:eq-psi}; the two conditions on the data used there are in {\cite[Lemma~13.4(S7)]{P46}}.)

\subsubsection{One real place: norms, lengths and the cut}\label{pbe:ss-real}

Fix $\tau\in\auPl$. Let $Y_\tau=Y_0\otimes_{E,\tau}\C$ and $G_\tau=\auU(Y_\tau)\cong\auU(k+\varrho,
\varrho+1)$. The \emph{normalised basis} of $Y_\tau$ is
\[
\tilde e_a=\frac{\kappa_\tau^{\,\varrho-a+1}}{\sqrt N}\,e_a,\qquad x_1,\dots,x_k,\qquad\tilde z=
\frac z{\sqrt N},\qquad\tilde e_a^*=\frac{\sqrt N}{\kappa_\tau^{\,\varrho-a+1}}\,e_a^*\qquad(1\le
a\le\varrho).
\]
In it the hermitian form has the matrix $\mathsf J_1$ with $h(\tilde e_a,\tilde e_b^*)=\delta_{ab}$,
$h(x_i,x_j)=\delta_{ij}$, $h(\tilde z,\tilde z)=-1$ and all other products zero; $\mathsf J_1$ is
hermitian and unitary. The \emph{normalised frame} is $(\tilde f_0,\dots,\tilde f_\varrho)=(\tilde
e_1,\dots,\tilde e_\varrho,\tilde z)$, so that $\tilde f_a=s_a f_a$ with
\begin{equation}\label{pbe:eq-scale}
s_a=\kappa_\tau^{\,\varrho-a}/\sqrt N\quad(0\le a\le\varrho),\qquad|s_a/s_{a'}|=|\kappa_\tau|^{a'-a}
\le1\ \text{ for }a'\le a .
\end{equation}
We give $Y_\tau$ the Euclidean norm $\|\cdot\|=\|\cdot\|_\tau$ for which the normalised basis is
orthonormal, and $\operatorname{End}(Y_\tau)$ the operator norm. For $g\in G_\tau$ we have $g^{-1}=
\mathsf J_1g^*\mathsf J_1$, hence
\begin{equation}\label{pbe:eq-norm}
\|g^{-1}\|=\|g\|\ge1,\qquad|h(v,w)|=|\pbip{\mathsf J_1w}{v}|\le\|v\|\,\|w\|\qquad(g\in G_\tau,\ v,w\in
Y_\tau).
\end{equation}
The groups $U_\tau$, $H_\tau$, $R_\tau=U_\tau\rtimes H_\tau$ are the real points of $U$, $H$, $R_0$
at $\tau$. An element of $R_\tau$ maps $\tilde f_a$ to $\tilde f_a$ plus a combination of the
$\tilde f_{a'}$ with $a'<a$, and $\tilde e_a$ to $\tilde e_a$ plus a combination of the $\tilde
e_{a'}$ with $a'<a$ ($U$ is unitriangular on the flag and trivial on $A^\perp/A$, and $H$ fixes the
$e_a$ and $z$).

\begin{definition}[Length and cut]\label{pbe:length}
For $y\in G_\tau$ the \emph{length} is $\pbel_\tau(y)=\inf_{s\in R_\tau}\bigl(\|sy\|+1\bigr)\ge2$,
and the \emph{scaled length} is $\pbels_\tau(y)=\max\bigl(1,\pbel_\tau(y)/b^{6n}\bigr)$. The
\emph{cutoff} is
\begin{equation}\label{pbe:eq-chi}
\chi_\tau(y)=\max\Bigl(0,\ \min\Bigl(1,\ 2-\frac{2\,\pbel_\tau(y)}{b^{6n+2}}\Bigr)\Bigr),
\end{equation}
so that $\chi_\tau(y)=1$ if $\pbel_\tau(y)\le\frac12b^{6n+2}$ and $\chi_\tau(y)=0$ if $\pbel_\tau(y)
\ge b^{6n+2}$. For $y=(y_\tau)\in G_\infty$ put $\pbeL(y)=\prod_\tau\pbels_\tau(y_\tau)\ge1$.
\end{definition}

\begin{lemma}\label{pbe:len}
Let $y,g\in G_\tau$ and $s\in R_\tau$.
\begin{enumerate}
\item[(a)] $\pbel_\tau(sy)=\pbel_\tau(y)$, and $\pbel_\tau(yg)\le\|g\|\,\pbel_\tau(y)$. Hence
$\pbel_\tau$, $\pbels_\tau$ and $\chi_\tau$ are continuous and left invariant under $R_\tau$.
\item[(b)] If $\chi_\tau(y)<1$ then $\pbels_\tau(y)\ge b$.
\item[(c)] Let $w_a=y^{-1}\tilde f_a$ $(0\le a\le\varrho)$ and let $o_0,\dots,o_\varrho$ be the
Gram--Schmidt orthogonalisation of $w_0,\dots,w_\varrho$, so that $w_a=o_a+\sum_{a'<a}t_{a'a}o_{a'}$.
The vectors $o_a$ do not change if $y$ is replaced by $sy$, and $\|o_a\|<\pbel_\tau(y)$.
\end{enumerate}
\end{lemma}
\begin{proof}
(a) The first claim is clear. For the second, $\|s yg\|+1\le\|sy\|\,\|g\|+1\le\|g\|(\|sy\|+1)$ by
\eqref{pbe:eq-norm}. Applied to $g$ and $g^{-1}$ this gives $\|g\|^{-1}\pbel_\tau(y)\le\pbel_\tau(yg)
\le\|g\|\pbel_\tau(y)$, hence the continuity. (b) $\chi_\tau(y)<1$ means $\pbel_\tau(y)>\frac12b^{6n+
2}$, so $\pbels_\tau(y)>b^2/2\ge b$. (c) $(sy)^{-1}\tilde f_a=y^{-1}s^{-1}\tilde f_a$ is $w_a$ plus
a combination of the $w_{a'}$, $a'<a$; so the flags of the two families agree and the $o_a$ are the
same. Hence $\|o_a\|\le\|(sy)^{-1}\tilde f_a\|\le\|(sy)^{-1}\|=\|sy\|$ for every $s$, and the
infimum over $s$ is $\pbel_\tau(y)-1$.
\end{proof}

\paragraph*{Measures}
We fix on $G_\tau$ the Haar measure $dg$ for which the formula \textup{(DS1)(iii)} below holds, on
the compact group $H_\tau$ the Haar measure of mass $1$, and on $U_\tau$ the Haar measure $du$ of
convention (Cv3) of \S\ref{pb:ss-setting} (Lebesgue measure in the free entries of the rows $\tilde e_a$,
as in {\cite[\S4.1]{P52}}). The
group $R_\tau$ is unimodular with Haar measure $ds=du\,dh$ ($s=uh$), because conjugation by the
compact group $H_\tau$ preserves $du$. With the quotient measure on $U_\tau\backslash G_\tau$,
\begin{equation}\label{pbe:eq-unfoldR}
\int_{U_\tau\backslash G_\tau}\int_{R_\tau}F(sy)\,ds\,dy=\int_{G_\tau}F(g)\,dg\qquad(F\ge0\text{
measurable}),
\end{equation}
since $\int_{U_\tau\backslash G_\tau}\int_{U_\tau}F(uhy)\,du\,dy=\int_{G_\tau}F(hg)\,dg=\int_{G_\tau}F$
for each $h\in H_\tau$.

\subsubsection{The discrete series coefficient: what is used}\label{pbe:ss-ds}

Let $\pi_\tau$ be the discrete series representation of $G_\tau$ with Harish-Chandra parameter
$\varsigma_\tau$ ({\cite[\S14.3]{P46}}) \emph{in the following chamber}: in a basis of $Y_\tau$ in
which the form is diagonal with $k+\varrho$ signs $+$ and $\varrho+1$ signs $-$, the $k-1$ entries
$\frac k2-1,\dots,1-\frac k2$ and the entries $-\kappa(\tau)-j\Delta$ with $j$ odd are attached to
coordinates of sign $+$, and the entries $-\kappa(\tau)-j\Delta$ with $j$ even to coordinates of sign
$-$. (The infinitesimal character $\varsigma_\tau$ alone does not determine the representation; this
is the choice made in the definition of $\auAz$ in {\cite[\S14.3]{P46}}.) Entries attached to coordinates of opposite signs differ by at least
$\Delta=2r^8\ge500r^3+r^7$. Let $K_\tau=G_\tau\cap\auU(r)$ and let $v_0$ be a unit vector of highest
weight in the lowest $K_\tau$-type of $\pi_\tau$. Put
\[
c_\tau(g)=\pbip{v_0}{\pi_\tau(g)v_0},\qquad\pbdeg_\tau=\prod_{i<j}\frac{\varsigma_{\tau,i}-
\varsigma_{\tau,j}}{j-i},\qquad e_\tau=\pbdeg_\tau\cdot\aubar{c_\tau},\qquad\pbeq=9r^2 .
\]
Let $\pbW_\tau$ be the representation of $H_\tau=\auU(k)$ of \S\ref{pb:ss-layer2} (irreducible, of
dimension $d_\tau$; it is trivial if $\tau\notin\auPa$), with its action $\tau_\tau$, and let
$\psi_\tau$ be the component of $\psi_\infty$. Then
\[
\pi_R(uh)=\psi_\tau(u)\,\tau_\tau(h)\qquad(u\in U_\tau,\ h\in H_\tau)
\]
is a unitary representation of $R_\tau$ on $\pbW_\tau$, because $\psi_\tau$ is invariant under
conjugation by $H_\tau$. For $x\in G_\tau$ let $\pbeB_x=\{s\in R_\tau:\|x^{-1}sx-1\|\le1\}$, a
compact neighbourhood of $1$ in $R_\tau$ (the map $s\mapsto x^{-1}sx$ is a homeomorphism of $R_\tau$
onto a closed subgroup of $G_\tau$); so $0<\pbvol(\pbeB_x)<\infty$.

\begin{theorem}[Properties of the discrete series coefficient]\label{pbe:ds}
For every $\tau\in\auPl$:
\begin{enumerate}
\item[(DS1)] \textup{(i)} $c_\tau$ is continuous, $c_\tau(g^{-1})=\aubar{c_\tau(g)}$ and $|c_\tau|
\le1$. \textup{(ii)} If $c_\tau$ is integrable, then $\int_{G_\tau}e_\tau(a)\,c_\tau(ga)\,da=c_\tau(g)$
for all $g\in G_\tau$, for the Haar measure $da$ for which $\int_{G_\tau}|c_\tau|^2=\pbdeg_\tau^{-1}$.
\textup{(iii)} This Haar measure is the one that corresponds to the measure of mass $1$ on the
compact form $\auU(r)$.
\item[(DS2)] $|c_\tau(g)|\le\bigl(2/\|g\|\bigr)^{\,j+\pbeq}$ for all $g\in G_\tau$ and all integers
$0\le j\le r^7$.
\item[(DS3)] $\pbdeg_\tau\cdot\int_{G_\tau}\bigl(32/\|g\|\bigr)^{\pbeq}\,dg\le b^{99r^2}$.
\item[(DS4)] There are a point $x_\tau\in G_\tau$ with $\pbel_\tau(x_\tau)\le b^{5n}$ and a unit
vector $v_\tau\in\pbW_\tau$ such that
\[
\Bigl|\int_{R_\tau}c_\tau(x_\tau^{-1}s\,x_\tau)\,\pbip{v_\tau}{\pi_R(s)^{-1}v_\tau}\,ds\Bigr|\ \ge\
b^{-60(r+k)^2}\,\pbvol(\pbeB_{x_\tau}).
\]
\end{enumerate}
\end{theorem}
\smallskip\noindent\textup{(DS1)}--\textup{(DS3)} are proved in \S\ref{pds:sec}: \textup{(DS1)(i),(ii)} by Schur orthogonality, \textup{(DS1)(iii)} from Corollary~\ref{pbr:degree} (Proposition~\ref{pds:haar} and~{\cite[Proposition~2.19]{P49}}), \textup{(DS2)} from the integral formula \eqref{pds:eq-FJ} for $c_\tau$ (Theorem~\ref{pds:decay} and~{\cite[Theorem~3.17]{P49}}), \textup{(DS3)} from \textup{(DS1)(ii)} and the Casimir values (Theorem~\ref{pds:tail}). \textup{(DS4)} is proved in \S\ref{pbp:sec} (Theorem~\ref{pbp:main}).\par\smallskip

\begin{lemma}[Covering]\label{pbe:cover}
For all $x,y\in G_\tau$, $\displaystyle\int_{R_\tau}\Bigl(\frac2{\|x^{-1}sy\|}\Bigr)^{\pbeq}ds\le16^{
\pbeq}\,\pbvol(\pbeB_x)$.
\end{lemma}
\begin{proof}
Let $m(s)=x^{-1}sy$ and, for $T\ge1$, $S_T=\{s\in R_\tau:\|m(s)\|\le T\}$. Choose a maximal family
$s_1,\dots,s_\nu$ in $S_T$ with $\|x^{-1}s_is_j^{-1}x-1\|>\frac12$ for $i\ne j$. For $i\ne j$,
$m(s_i)-m(s_j)=(x^{-1}s_is_j^{-1}x-1)\,m(s_j)$, so by \eqref{pbe:eq-norm} $\|m(s_i)-m(s_j)\|\ge
\frac12\|m(s_j)^{-1}\|^{-1}\ge(2T)^{-1}$. The $m(s_i)$ lie in the ball of radius $T$ of the space of
$r\times r$ matrices, of real dimension $2r^2$; comparing volumes of balls, $\nu\le(1+4T^2)^{2r^2}
\le(5T^2)^{2r^2}$. By maximality $S_T\subseteq\bigcup_i\pbeB_x s_i$, and the Haar measure is right
invariant, so $\pbvol(S_T)\le5^{2r^2}T^{4r^2}\pbvol(\pbeB_x)$. As $\|m(s)\|\ge1$,
\[
\int_{R_\tau}\Bigl(\frac2{\|m(s)\|}\Bigr)^{\pbeq}ds=2^{\pbeq}\int_1^\infty\pbeq\,T^{-\pbeq-1}\pbvol
(S_T)\,dT\le2^{\pbeq}5^{2r^2}\frac{\pbeq}{\pbeq-4r^2}\pbvol(\pbeB_x)\le16^{\pbeq}\pbvol(\pbeB_x).
\qedhere
\]
\end{proof}

\subsubsection{The Bessel function at a real place}\label{pbe:ss-local}

In this subsection $\tau$ is fixed and dropped from the notation where possible; $x=x_\tau$ and
$v=v_\tau$ are as in \textup{(DS4)}.

\begin{proposition}[The local Bessel function]\label{pbe:local}
Assume Theorem~\ref{pbe:ds}. For $y\in G_\tau$ the integral
\[
W^\circ(y)=\int_{R_\tau}c_\tau(x^{-1}sy)\,\pi_R(s)^{-1}v\,ds\ \in\pbW_\tau
\]
converges absolutely, and:
\begin{enumerate}
\item[(a)] $W^\circ$ is continuous, $\|W^\circ(y)\|\le16^{\pbeq}\pbvol(\pbeB_x)$, and $W^\circ(sy)=
\pi_R(s)W^\circ(y)$ for $s\in R_\tau$.
\item[(b)] $c_\tau$ is integrable on $G_\tau$, $\int_{G_\tau}e_\tau(a)W^\circ(ya)\,da=W^\circ(y)$ for
all $y$, and $\int_{U_\tau\backslash G_\tau}\|W^\circ\|\le\int_{G_\tau}|c_\tau|$.
\item[(c)] \emph{(The norm is a period.)} $\displaystyle\pbdeg_\tau\int_{U_\tau\backslash G_\tau}\|W^
\circ(y)\|^2dy=\pbip v{W^\circ(x)}=\int_{R_\tau}c_\tau(x^{-1}sx)\pbip v{\pi_R(s)^{-1}v}\,ds$. In
particular this number is real and, by \textup{(DS4)}, at least $b^{-60(r+k)^2}\pbvol(\pbeB_x)>0$.
\end{enumerate}
Let $c_0>0$ be such that $W_\tau:=c_0W^\circ$ has $\int_{U_\tau\backslash G_\tau}\|W_\tau\|^2=d_\tau$,
and $W_\tau^\flat:=\chi_\tau W_\tau$ with $\chi_\tau$ as in \eqref{pbe:eq-chi}. Put
$\mathsf A_\tau(y)=c_0\int_{R_\tau}(2/\|x^{-1}sy\|)^{\pbeq}ds$ and $\mathsf a_\tau=c_0\,16^{\pbeq}\pbvol(
\pbeB_x)$. Then:
\begin{enumerate}
\item[(d)] $\mathsf A_\tau$ is continuous and left invariant under $R_\tau$, $\mathsf A_\tau\le\mathsf
a_\tau$, and
\[
\pbels_\tau(y)^{3r^3f}\,\|W_\tau(y)\|\le\mathsf A_\tau(y)\quad(y\in G_\tau),\qquad\mathsf a_\tau\int_{
U_\tau\backslash G_\tau}\mathsf A_\tau\le b^{400r^2}\,d_\tau .
\]
\item[(e)] $\|W^\flat_\tau\|\le\|W_\tau\|$, $W^\flat_\tau(y)=0$ unless $\pbel_\tau(y)\le b^{6n+2}$, and
$\displaystyle\int_{U_\tau\backslash G_\tau}\|W_\tau^\flat-W_\tau\|^2\le b^{-2r^3f}\,d_\tau$.
\end{enumerate}
\end{proposition}
\begin{proof}
(a) By \textup{(DS2)} with $j=0$ and Lemma~\ref{pbe:cover}, $\int_{R_\tau}|c_\tau(x^{-1}sy)|ds\le16^{
\pbeq}\pbvol(\pbeB_x)$; as $\pi_R$ is unitary and $\|v\|=1$ this gives the convergence and the bound.
For $y'$ in a compact neighbourhood of $y$, $\|x^{-1}sy\|\le\|x^{-1}sy'\|\,\|y'^{-1}y\|$, so
$|c_\tau(x^{-1}sy')|\le C(2/\|x^{-1}sy\|)^{\pbeq}$ with a constant $C$, an integrable function of $s$;
the continuity follows by dominated convergence. The substitution $s\mapsto ss'^{-1}$ gives
$W^\circ(s'y)=\int c_\tau(x^{-1}sy)\pi_R(ss'^{-1})^{-1}v\,ds=\pi_R(s')W^\circ(y)$.

(b) By \textup{(DS2)} with $j=0$ and \textup{(DS3)}, $\int_{G_\tau}|c_\tau|<\infty$. By (a) the
double integral of $|e_\tau(a)|\,|c_\tau(x^{-1}sya)|$ over $R_\tau\times G_\tau$ is at most
$\|e_\tau\|_1\,16^{\pbeq}\pbvol(\pbeB_x)$, so Fubini applies, and by \textup{(DS1)(ii)}
\[
\int_{G_\tau}e_\tau(a)W^\circ(ya)\,da=\int_{R_\tau}\Bigl(\int_{G_\tau}e_\tau(a)c_\tau(x^{-1}sya)\,da
\Bigr)\pi_R(s)^{-1}v\,ds=W^\circ(y).
\]
By \eqref{pbe:eq-unfoldR}, $\int_{U_\tau\backslash G_\tau}\|W^\circ\|\le\int_{U_\tau\backslash G_\tau}
\int_{R_\tau}|c_\tau(x^{-1}sy)|\,ds\,dy=\int_{G_\tau}|c_\tau(x^{-1}g)|\,dg$.

(c) Insert the definition of $W^\circ$ in the first variable of $\pbip{W^\circ(y)}{W^\circ(y)}$ and
use (a):
\[
\pbip{W^\circ(y)}{W^\circ(y)}=\int_{R_\tau}\aubar{c_\tau(x^{-1}sy)}\,\pbip{\pi_R(s)^{-1}v}{W^\circ(y)}
\,ds=\int_{R_\tau}\aubar{c_\tau(x^{-1}sy)}\,\pbip v{W^\circ(sy)}\,ds .
\]
Let $\Phi(g)=e_\tau(x^{-1}g)\pbip v{W^\circ(g)}$; it is integrable on $G_\tau$, since $e_\tau$ is and
$W^\circ$ is bounded. By the identity above and \eqref{pbe:eq-unfoldR} (applied to the positive and
negative parts of the real and imaginary parts of $\Phi$),
\[
\pbdeg_\tau\int_{U_\tau\backslash G_\tau}\|W^\circ\|^2=\int_{U_\tau\backslash G_\tau}\int_{R_\tau}\Phi(
sy)\,ds\,dy=\int_{G_\tau}\Phi(g)\,dg=\int_{G_\tau}e_\tau(a)\pbip v{W^\circ(xa)}\,da=\pbip v{W^\circ
(x)},
\]
by (b). The last expression of (c) is the definition of $W^\circ(x)$.

(d) The invariance of $\mathsf A_\tau$ under $y\mapsto s'y$ is the right invariance of $ds$; the
continuity is proved as in (a); $\mathsf A_\tau\le\mathsf a_\tau$ is Lemma~\ref{pbe:cover}. Let $s\in
R_\tau$ and $g=x^{-1}sy$. By Lemma~\ref{pbe:len}(a), $\pbel_\tau(y)=\pbel_\tau(sy)=\pbel_\tau(xg)\le
\|g\|\pbel_\tau(x)\le\|g\|b^{5n}$, so $\pbel_\tau(y)/b^{6n}\le\|g\|b^{-n}\le\|g\|/2$ and
$\pbels_\tau(y)\le\max(1,\|g\|/2)$. If $\pbels_\tau(y)=1$,
\textup{(DS2)} with $j=0$ gives $\pbels_\tau(y)^{3r^3f}|c_\tau(g)|\le(2/\|g\|)^{\pbeq}$. Otherwise
$\pbels_\tau(y)\le\|g\|/2$, and \textup{(DS2)} with $j=3r^3f\le r^7$ gives $\pbels_\tau(y)^{j}|c_\tau(
g)|\le(\|g\|/2)^j(2/\|g\|)^{j+\pbeq}=(2/\|g\|)^{\pbeq}$. Integrating over $s$ and multiplying by
$c_0$ we get the first inequality of (d). For the second, by \eqref{pbe:eq-unfoldR}, (c), the
definition of $c_0$ and \textup{(DS3)},
\begin{align*}
\mathsf a_\tau\int_{U_\tau\backslash G_\tau}\mathsf A_\tau&=c_0^2\,16^{\pbeq}\pbvol(\pbeB_x)\int_{G_
\tau}\Bigl(\frac2{\|g\|}\Bigr)^{\pbeq}dg\le c_0^2\,b^{60(r+k)^2}\pbdeg_\tau\!\int_{U_\tau\backslash
G_\tau}\!\!\|W^\circ\|^2\cdot\int_{G_\tau}\Bigl(\frac{32}{\|g\|}\Bigr)^{\pbeq}dg\\
&=d_\tau\,b^{60(r+k)^2}\,\pbdeg_\tau\int_{G_\tau}\Bigl(\frac{32}{\|g\|}\Bigr)^{\pbeq}dg\le d_\tau\,
b^{240r^2+99r^2}\le d_\tau\,b^{400r^2}.
\end{align*}

(e) The first two claims follow from $0\le\chi_\tau\le1$ and \eqref{pbe:eq-chi}. Where $W_\tau^\flat
\ne W_\tau$ we have $\pbels_\tau\ge b$ by Lemma~\ref{pbe:len}(b), so by (d)
$b^{3r^3f}\|W_\tau^\flat-W_\tau\|\le b^{3r^3f}\|W_\tau\|\le\mathsf A_\tau\le\mathsf a_\tau$ everywhere,
hence $b^{6r^3f}\|W^\flat_\tau-W_\tau\|^2\le\mathsf a_\tau\mathsf A_\tau$. Integrating,
$\int\|W_\tau^\flat-W_\tau\|^2\le b^{-6r^3f}b^{400r^2}d_\tau\le b^{-2r^3f}d_\tau$, because $400r^2\le4r^3
\le4r^3f$.
\end{proof}

\begin{remark}\label{pbe:rem-local}
The whole local theory is thus four lines: the average of the coefficient over $R_\tau$ is
reproduced by the kernel; therefore its $L^2$-norm on $U_\tau\backslash G_\tau$ is its own value at
$x$, a period; the period is not small by \textup{(DS4)}; and the decay of the coefficient passes
to the average through the covering Lemma~\ref{pbe:cover}.
\end{remark}

\subsubsection{The Poincar\'e series}\label{pbe:ss-global}

Put $W=\bigotimes_\tau W_\tau$, $W^\flat=\bigotimes_\tau W_\tau^\flat$ on $G_\infty$ with values in
$\pbW=\bigotimes_\tau\pbW_\tau$, $g_0=(x_\tau)_\tau$, $e_\infty=\bigotimes_\tau e_\tau$, $\mathsf
A(y)=\prod_\tau\mathsf A_\tau(y_\tau)$ and $\mathsf a_\infty=\prod_\tau\mathsf a_\tau$. By
Proposition~\ref{pbe:local}(d), for $y\in G_\infty$,
\begin{equation}\label{pbe:eq-decay}
\|W(y)\|\le\mathsf A(y)\,\pbeL(y)^{-3r^3f}\le\mathsf a_\infty\,\pbeL(y)^{-3r^3f},\qquad\mathsf a_\infty
\int_{U_\infty\backslash G_\infty}\mathsf A\le b^{400r^2f}\,d_\infty .
\end{equation}

\paragraph*{The space $\auAz$}
By the definition of {\cite[\S14.3]{P46}}, $\auAz$ is the space of the bounded continuous functions $y$ on $G(\auA)$, left invariant under
$\Gamma$ and right invariant under $J$, with
\begin{equation}\label{pbe:eq-repro}
\int_{G_\infty}e_\infty(a)\,y(ga)\,da=y(g)\qquad(g\in G(\auA)).
\end{equation}
This is the meaning of ``the archimedean component is the vector singled out by the reproducing
kernel''. The quotient $\Gamma\backslash G(\auA)$ has finite volume~\cite[Theorem~5.6(i)]{BorelFin} (the group has no rational characters, its centre being anisotropic; see also \cite[Theorem~9.4]{BorelHC1962} for arithmetic subgroups of the real group), so $\auAz$
carries the inner product of \S\ref{pb:ss-setting}; as in {\cite[\S14.3]{P46}} we use that $\auAz$
is finite-dimensional. The linear form $\varphi\mapsto\int_{[R_0]}\mathcal W_\varphi(r)\aubar{y(rg_0)}
\,dr$ is defined for $y\in\auAz$ because $[R_0]$ is compact ($U$ is unipotent and $H$ is
anisotropic) and $y$ is continuous.

\paragraph*{Counting the terms of a series}

\begin{lemma}\label{pbe:eichler}
Let $c\in\mathcal C$. For every strictly upper triangular matrix $(C_{a'a})_{0\le a',a\le\varrho}$
with entries in $M\auO_E$ there is $u\in U_c$ with $uf_a=f_a+\sum_{a'<a}C_{a'a}f_{a'}$ for $0\le a\le
\varrho$.
\end{lemma}
\begin{proof}
Let $L'=h_cL$, where $L$ is the lattice of $Y_0$; as $h_c\in H(\auA_f)$, $L'=L_{\mathrm{hyp}}\oplus
L'_{\mathrm{mid}}$ with $L_{\mathrm{hyp}}=\bigoplus_a(\widehat{\auO}e_a\oplus\widehat{\auO}e_a^*)$ and
$L'_{\mathrm{mid}}$ a unimodular lattice of $X_0\perp Ez$ containing $z$, and $U_c$ consists of the
$u\in U(E^+)$ with $(u-1)L'\subseteq ML'$ (the condition at the conjugates of the chosen places is
the same condition, read on the adjoint).
\emph{First}, let $a_1$ be the unitriangular endomorphism of $A$ with $a_1e_{a}=e_a+\sum_{a'<a}C_{a'-1,
a-1}e_{a'}$ and $\mathsf m(a_1)\in G(E^+)$ the element that acts by $a_1$ on $A$, by the inverse
adjoint of $a_1$ on $\langle e_a^*\rangle$ and trivially on $X_0\perp Ez$. It lies in $U(E^+)$,
commutes with $h_c$, and $\mathsf m(a_1)-1$ maps $L_{\mathrm{hyp}}$ into $ML_{\mathrm{hyp}}$, since
$a_1^{-1}-1$ is a polynomial without constant term in $a_1-1$. So $\mathsf m(a_1)\in U_c$.
\emph{Second}, let $v\in M\bigoplus_a\auO_Ee_a$. The map $q\mapsto h(q,z)$ from $L'_{\mathrm{mid}}\cap
(X_0\perp Ez)$ to $\auO_E$ is onto: its image is an ideal whose completion is $h(L'_{\mathrm{mid}},z)
=\widehat{\auO}$, because $L'_{\mathrm{mid}}$ is unimodular and $z$ is primitive in it ($h(x_k/\aubar{
\pi_N}+z/N,z)=-1$, and $h_c$ fixes $z$). Choose $q$ with $h(z,q)=1$, and $\omega\in\auO_E$ with
$\omega+\aubar\omega=1$ ($E/E^+$ is unramified at all finite places, so the trace is onto). Then
\[
\mathsf E(t)=t+h(t,q)\,v-h(t,v)\,q-\omega\,h(q,q)\,h(t,v)\,v
\]
is an isometry of $Y_0$ (a direct computation, using $h(v,v)=h(v,q)=0$); it fixes $A$ pointwise and
maps $t\in A^\perp$ to $t+h(t,q)v\in t+A$, so $\mathsf E\in U(E^+)$; $\mathsf E(z)=z+v$; and $(\mathsf
E-1)L'\subseteq ML'$ because $q\in L'$, $v\in ML'$. So $\mathsf E\in U_c$. The element $u=\mathsf
m(a_1)\mathsf E$ with $v=a_1^{-1}\sum_{a'<\varrho}C_{a'\varrho}e_{a'+1}$ has the required property.
\end{proof}

\begin{proposition}[Count of terms]\label{pbe:count}
Let $c\in\mathcal C$, $s\in\auG$ and $y\in G_\infty$. Then
\[
\sum_{c'\in\mathcal C}\ \sum_{\gamma\in\Gamma_{c,s,c'}/U_{c'}}\pbeL(\gamma^{-1}y)^{-r^3f}\ \le\ 2\,b^{
12n^3rf}.
\]
\end{proposition}
The summand depends only on the coset $\gamma U_{c'}$, by Lemma~\ref{pbe:len}(a).
\begin{proof}
For $R\ge1$ let $\pbeN_R$ be the set of the pairs $(c',\gamma U_{c'})$ with $\gamma\in\Gamma_{c,s,c'}$
and $\pbeL(\gamma^{-1}y)\le R$. We show
\begin{equation}\label{pbe:eq-pack}
\#\pbeN_R\le\bigl(b^{6n^2}R^2\bigr)^{D},\qquad D=2nrf .
\end{equation}
Then, with $i=r^3f\ge4D$ (as $r^2\ge8n$), the sum of the proposition is
$\int_1^\infty iT^{-i-1}\#\pbeN_T\,dT\le b^{6n^2D}\,i/(i-2D)\le2b^{12n^3rf}$.

Put $\rho=b^{37n}R$. For $(c',\gamma U_{c'})\in\pbeN_R$ and every $\tau$ we have $\pbels_\tau\le\pbeL$,
hence
\begin{equation}\label{pbe:eq-lenR}
\pbel_\tau\bigl((\gamma^{-1}y)_\tau\bigr)\le b^{6n}R .
\end{equation}
\emph{Step 1: reduced representatives.} Each element of $\pbeN_R$ has a representative $\gamma$ with
$\|y_\tau^{-1}\gamma\tilde f_a\|\le\rho$ for all $\tau$ and $a$. Indeed let $z=\gamma^{-1}y$, and at
each $\tau$ let $w_a=z_\tau^{-1}\tilde f_a$, $o_a$ and $t_{a'a}$ be as in Lemma~\ref{pbe:len}(c). For
$u\in U_{c'}$ with $uf_a=f_a+\sum_{a'<a}C_{a'a}f_{a'}$, replacing $\gamma$ by $\gamma u$ replaces
$w_a$ by $z_\tau^{-1}u\tilde f_a=w_a+\sum_{a'<a}\tau(C_{a'a})(s_a/s_{a'})w_{a'}$, does not change the
$o_a$, and replaces $t_{a'a}(s_{a'}/s_a)$ by
\[
t_{a'a}\frac{s_{a'}}{s_a}+\tau(C_{a'a})+\sum_{a'<a''<a}t_{a'a''}\frac{s_{a'}}{s_{a''}}\,\tau(C_{a''a}).
\]
The lattice $M\auO_E$ of $\prod_\tau\C$ contains $M\auO_K[\zeta_{q_1}]$, which has the $\Z$-basis
$M\zeta_{q_1}^i$, $M\omega_K\zeta_{q_1}^i$ ($0\le i<f$, $\omega_K=\frac12(1+\sqrt{-d})$ or $\sqrt{-d}$)
of vectors with all coordinates of absolute value at most $M\sqrt d$; so every point of
$\prod_\tau\C$ is at distance at most $fM\sqrt d\le b^{26}$, in every coordinate, from a point of $M
\auO_E$. By Lemma~\ref{pbe:eichler} we can therefore choose the $C_{a'a}\in M\auO_E$, by induction
on $a-a'$, such that after the replacement $|t_{a'a}(s_{a'}/s_a)|\le b^{26}$ at every $\tau$; then
$|t_{a'a}|\le b^{26}$ by \eqref{pbe:eq-scale}, and by Lemma~\ref{pbe:len}(c) and \eqref{pbe:eq-lenR}
\[
\|y_\tau^{-1}\gamma\tilde f_a\|=\|w_a\|\le\sum_{a'\le a}|t_{a'a}|\,\|o_{a'}\|\le n\,b^{26}\,b^{6n}R\le\rho .
\]
\emph{Step 2: distinct elements have distinct frames.} Let $\gamma\in\Gamma_{c,s,c'}$, $\gamma'\in
\Gamma_{c,s,c''}$ with $\gamma f_a=\gamma'f_a$ for all $a$. Then $g=\gamma^{-1}\gamma'$ fixes the
frame, so $g=u\eta$ with $\eta\in H(E^+)$ and $u\in U(E^+)$ fixing the frame, and
$h_{c'}^{-1}gh_{c''}=(h_{c'}^{-1}\gamma^{-1}h_c)(h_c^{-1}\gamma'h_{c''})$ lies in $K$ and has reduction
$s^{-1}s=1$; so it lies in $J\cap R_0(\auA_f)$. By Lemma~\ref{pb:KR} and $\pbUb\cap\pbHb=1$,
$J\cap R_0(\auA_f)=(J\cap U(\auA_f))\rtimes J_X$. Hence $h_{c'}^{-1}\eta h_{c''}\in J_X$, so $c''=c'$
and $\eta\in H(E^+)\cap h_{c'}J_Xh_{c'}^{-1}=1$, the level being free; and then $u\in U_{c'}$. So
$\gamma'U_{c'}=\gamma U_{c'}$.

\emph{Step 3: rigidity.} Let $\gamma,\gamma'$ be reduced representatives, as in Step~1, of two
elements of $\pbeN_R$, and suppose that
\begin{equation}\label{pbe:eq-close}
\|y_\tau^{-1}\gamma'\tilde f_a-y_\tau^{-1}\gamma\tilde f_a\|<(N\rho)^{-1}\qquad\text{for all }\tau,a .
\end{equation}
We claim $\gamma'f_a=\gamma f_a$ for all $a$. Let $g=\gamma^{-1}\gamma'$; then $h_{c'}^{-1}gh_{c''}\in
K$, and the $h_c$ fix the vectors $e_a$, $e_a^*$, $z$ of the lattice, so the numbers $h(gf_a,f_b)$
and $h(gf_a,e_x^*)$ are algebraic integers of $E$. First, by \eqref{pbe:eq-norm}, \eqref{pbe:eq-close}
and $|s_a|\ge N^{-1/2}$,
\[
\bigl|\tau\bigl(h(gf_a,f_b)-h(f_a,f_b)\bigr)\bigr|=\frac{|h(\gamma'\tilde f_a-\gamma\tilde f_a,\gamma
\tilde f_b)|}{|s_a||s_b|}\le N\,\|y_\tau^{-1}(\gamma'\tilde f_a-\gamma\tilde f_a)\|\,\|y_\tau^{-1}
\gamma\tilde f_b\|<1
\]
for every $\tau$; an algebraic integer all of whose conjugates have absolute value less than $1$ is
$0$, so $h(gf_a-f_a,f_b)=0$ for all $a,b$. Thus $gf_a-f_a$ is orthogonal to $\pbeAp=A+Ez$, that is,
$gf_a=f_a+\alpha_a+\xi_a$ with $\alpha_a\in A$, $\xi_a\in X_0$. As $g$ is an isometry and $A\oplus
X_0$ is orthogonal to the frame, $h(f_a,f_b)=h(gf_a,gf_b)=h(f_a,f_b)+h(\xi_a,\xi_b)$; $X_0$ is
definite, so $\xi_a=0$. Second, write $\alpha_a=\sum_xc_{xa}e_x$ with $c_{xa}=h(gf_a-f_a,e_x^*)\in\auO_E$
and suppose $\alpha_a\ne0$; let $x$ be the largest index with $c_{xa}\ne0$. At a place $\tau$,
$g\tilde f_a-\tilde f_a=\sum_{x'}\tau(c_{x'a})(s_a/s'_{x'})\tilde e_{x'}=:\tilde\alpha$, where
$s'_{x'}$ is the scale of $e_{x'}$ and $|s_a/s'_{x'}|\ge|\kappa_\tau|^{-\varrho}\ge b^{-23\varrho}$
by \eqref{pbe:eq-kappa}. Let $z=\gamma^{-1}y$. For $s\in R_\tau$ the vector $s\tilde\alpha$ has
the same coefficient at $\tilde e_x$ as $\tilde\alpha$, and $s\tilde\alpha=(sz_\tau)\,z_\tau^{-1}\tilde
\alpha$, so
\[
b^{-23\varrho}|\tau(c_{xa})|\le\|s\tilde\alpha\|\le\|sz_\tau\|\,\|z_\tau^{-1}(g\tilde f_a-\tilde f_a)\|
<\|sz_\tau\|\,(N\rho)^{-1}
\]
by \eqref{pbe:eq-close}. Taking the infimum over $s$ and using \eqref{pbe:eq-lenR}, $|\tau(c_{xa})|<
b^{23\varrho+6n}R\,(Nb^{37n}R)^{-1}<1$ for every $\tau$, so $c_{xa}=0$, a contradiction. This proves
the claim.

\emph{Step 4: packing.} To each element of $\pbeN_R$ attach the point $(y_\tau^{-1}\gamma\tilde f_a)_{
\tau,a}$ of $V=(\C^r)^{n\times f}$, for a reduced representative $\gamma$; $V$ has real dimension
$D$ and we give it the norm $\max_{\tau,a}\|\cdot\|$. By Step~1 the points lie in the ball of radius
$\rho$; by Steps~2 and~3 the points of two different elements are at distance at least $\delta=(N
\rho)^{-1}$. Comparing volumes, $\#\pbeN_R\le(1+2\rho/\delta)^D=(1+2N\rho^2)^D\le(3b^{134n}R^2)^D
\le(b^{6n^2}R^2)^D$, as $n\ge23$. This is \eqref{pbe:eq-pack}.
\end{proof}

\paragraph*{Existence of the Bessel functions}

\begin{theorem}[Hypothesis~\ref{pb:bessel} holds]\label{pbe:bessel}
Assume Theorem~\ref{pbe:ds}. Let $g_0=(x_\tau)_\tau$, let $\mathsf t=\bigotimes_\tau v_\tau\in\pbW$,
regarded as a linear form on $W_{\lambda^X}$, and let $W_\tau$, $W_\tau^\flat$, $\chi_\tau$ be as in
Proposition~\ref{pbe:local}. Then \textup{(B1)--(B4)} of Definition~\ref{pb:bessel} hold, with
\[
c_W=\frac{\pbvol(R_0(\auA_f)\cap J)}{\pbvol(J)}\prod_\tau\frac{\pbdeg_\tau}{c_{0,\tau}}>0,
\]
and \textup{(B5)} holds in the following exact form:
\begin{enumerate}
\item[(B5)] $W_\tau^\flat=\chi_\tau W_\tau$, where $\chi_\tau(y)=\max\bigl(0,\min(1,2-2\pbel_\tau(y)/
b^{6n+2})\bigr)$ and $\pbel_\tau(y)=\inf_{s\in R_\tau}(\|sy\|_\tau+1)$ is the length of
Definition~\ref{pbe:length}.
\end{enumerate}
Moreover:
\begin{enumerate}
\item[(B6)] \emph{(Decay.)} \eqref{pbe:eq-decay} holds, with $\mathsf A$ continuous and left
invariant under $R_\infty$.
\item[(B7)] \emph{(Reproduction.)} $\int_{G_\tau}e_\tau(a)W_\tau(ya)\,da=W_\tau(y)$ for all $\tau$
and $y\in G_\tau$.
\end{enumerate}
\end{theorem}
\begin{proof}
(B1), (B2), (B5), (B6), (B7) are Proposition~\ref{pbe:local} and Lemma~\ref{pbe:len}(a) (for
$W_\tau^\flat$ in (B1)).

(B3) Fix $g=(g_\infty,g_f)$. If no $\gamma_1\in\Gamma$ has $\gamma_1g_f\in h_cJ$, the series is
empty. Otherwise the $\gamma$ that occur are those of $\Gamma_c\gamma_1$, $\Gamma_c=\Gamma\cap h_cJh_c
^{-1}=\Gamma_{c,1,c}$, and with $y=\gamma_{1,\infty}g_\infty$, \eqref{pbe:eq-decay} and
Proposition~\ref{pbe:count} (for $s=1$, keeping only $c'=c$, and $\gamma\mapsto\gamma^{-1}$),
\[
\sum_{\gamma\in U_c\backslash\Gamma}\|W(\gamma_\infty g_\infty)\|\,\pbone_{h_cJ}(\gamma g_f)=\sum_{
\gamma\in\Gamma_c/U_c}\|W(\gamma^{-1}y)\|\le\mathsf a_\infty\sum_{\gamma\in\Gamma_c/U_c}\pbeL(\gamma^{
-1}y)^{-r^3f}\le2\,\mathsf a_\infty\,b^{12n^3rf}.
\]
So the series is bounded, and $P_{c,i}$ is a bounded function, left invariant under $\Gamma$ and
right invariant under $J$. By (B7), Fubini and the bound just proved (with $\|e_\infty\|_1<\infty$),
$P_{c,i}$ satisfies \eqref{pbe:eq-repro}; and a bounded function that satisfies \eqref{pbe:eq-repro}
is continuous, because $g_\infty\mapsto e_\infty(g_\infty^{-1}\,\cdot\,)$ is continuous into
$L^1(G_\infty)$. So $P_{c,i}\in\auAz$.

(B4) Let $y\in\auAz$ and $\varphi\in\auS$. The function $F(r)=\mathcal W_\varphi(r)\aubar{y(rg_0)}$ on
$R_0(\auA)$ is left invariant under $R_0(E^+)$ and right invariant under $K_R'=R_0(\auA_f)\cap J$
(for $\mathcal W_\varphi$ this is the computation in the proof of Lemma~\ref{pb:equiv}, with
$\psi_f$ trivial on $U(\auA_f)\cap J$ and $\varphi$ right invariant under $J_X$). By strong
approximation for $U$ and the definition of the classes, $R_0(\auA)=\bigsqcup_cR_0(E^+)\,(R_\infty
\times h_cK_R')$, and the stabiliser $R_0(E^+)\cap h_cK_R'h_c^{-1}$ is $U_c$ (as in Step~2 of the
proof of Proposition~\ref{pbe:count}). Hence
\[
\pbip y{\auP\varphi}=\int_{[R_0]}F=\pbvol(K_R')\sum_c\int_{U_c\backslash R_\infty}\mathcal W_\varphi(s,
h_c)\,\aubar{y(sg_0,h_c)}\,ds .
\]
By \eqref{pbe:eq-repro} and the substitution $a=g_0^{-1}s^{-1}g$, $\aubar{y(sg_0,h_c)}=\int_{G_\infty}
\aubar{e_\infty(g_0^{-1}s^{-1}g)}\,\aubar{y(g,h_c)}\,dg$. For $s=uh$, $\mathcal W_\varphi(s,h_c)=
\psi_\infty(u)\,\mathsf t(h^{-1}\varphi(h_c))=\pbip{\pi_R(s)\mathsf t}{\varphi(h_c)}$ in the pairing of
$\pbW$ with $W_{\lambda^X}$ (here $\pi_R=\bigotimes_\tau\pi_{R,\tau}$ and $\tau_\infty$ is the
contragredient action). As $y$ is bounded, $e_\infty$ integrable and $y(ug,h_c)=y(g,h_c)$ for $u\in
U_c$, we may exchange the integrals and fold the integral over $G_\infty$:
\[
\pbip y{\auP\varphi}=\pbvol(K_R')\sum_c\int_{U_c\backslash G_\infty}\aubar{y(g,h_c)}\,\Bigl\langle\int_{
R_\infty}\aubar{e_\infty(g_0^{-1}s^{-1}g)}\,\pi_R(s)\mathsf t\,ds,\ \varphi(h_c)\Bigr\rangle\,dg .
\]
Since $\aubar{e_\tau}=\pbdeg_\tau c_\tau$, the substitution $s\mapsto s^{-1}$ shows that the inner
integral is $\bigotimes_\tau\pbdeg_\tau W_\tau^\circ(g_\tau)=\prod_\tau(\pbdeg_\tau/c_{0,\tau})\cdot
W(g)$. On the other hand, by (B3) the series $P_\varphi=\sum_{c,i}\varphi(h_c)_iP_{c,i}$ lies in
$\auAz$, and unfolding it as in the proof of Proposition~\ref{pb:unfold} gives $\pbip y{P_\varphi}=
\pbvol(J)\sum_c\int_{U_c\backslash G_\infty}\aubar{y(g,h_c)}\pbip{W(g)}{\varphi(h_c)}\,dg$. So
$\pbip y{\auP\varphi}=c_W\pbip y{P_\varphi}$ for all $y\in\auAz$, and $\auP\varphi=c_WP_\varphi$.
\end{proof}

\subsubsection{Vanishing by character, and the terms that preserve \texorpdfstring{$A$}{A}}
\label{pbe:ss-char}

For $\gamma\in\Gamma$ let $V_\gamma=U\cap\gamma U\gamma^{-1}$, a unipotent group over $E^+$, and
$\chi_\gamma(u)=\psi_\infty(u)\,\aubar{\psi_\infty(\gamma^{-1}u\gamma)}$, a character of $V_\gamma(
E^+\otimes\R)$. We say that $\gamma$ \emph{agrees} if $\chi_\gamma$ is trivial on $V_\gamma(E^+)$;
since $\psi_\infty\psi_f=1$ on $U(E^+)$, this means $\psi_f(u)=\psi_f(\gamma^{-1}u\gamma)$ for all
$u\in U(E^+)\cap\gamma U(E^+)\gamma^{-1}$. As $\psi_\infty$ is a character of $U_\infty$, $\chi_{u_1
\gamma u_2}(u)=\chi_\gamma(u_1^{-1}uu_1)$ for $u_1,u_2\in U(E^+)$, so the set of agreeing elements is
stable.

\begin{lemma}[Vanishing by character]\label{pbe:vanish}
Assume \textup{(B1)--(B4)}. Let $\flat\in\{\ ,\flat\}$ and let $\gamma\in\Gamma_{c,s,c'}$ not agree.
Then the double coset of $\gamma$ contributes $0$ to \eqref{pb:eq-amp}:
\[
\int_{(U_c\cap\gamma U_{c'}\gamma^{-1})\backslash G_\infty}W^{\flat}(y)\,W^{\flat}(\gamma^{-1}y)^*\,dy=0.
\]
Consequently $A_T^\flat(c,s,c')=0$ for every stable $T$ without agreeing elements.
\end{lemma}
\begin{proof}
Let $\Lambda=U_c\cap\gamma U_{c'}\gamma^{-1}$ and $V_\infty=V_\gamma(E^+\otimes\R)$. The integral
converges absolutely (Proposition~\ref{pb:unfold} and $\|W^\flat\|\le\|W\|$). The group $\Lambda$
is an arithmetic subgroup of the unipotent group $V_\gamma$, hence discrete and cocompact in
$V_\infty$, and $V_\infty$ is unimodular. By (B1) and Lemma~\ref{pbe:len}(a) the integrand $F(y)$
satisfies $F(uy)=\chi_\gamma(u)F(y)$ for $u\in V_\infty$, and it is invariant under $\Lambda$
(\S\ref{pb:ss-layer2}). If $F=0$ there is nothing to prove. Otherwise $\chi_\gamma$ is trivial on
$\Lambda$, and the integral equals $\int_{V_\infty\backslash G_\infty}F(y)\bigl(\int_{\Lambda\backslash
V_\infty}\chi_\gamma(u)\,du\bigr)dy$. The measure of $\Lambda\backslash V_\infty$ is invariant under
right translations, so $\int_{\Lambda\backslash V_\infty}\chi_\gamma(u)\,du=\chi_\gamma(u_0)\int_{
\Lambda\backslash V_\infty}\chi_\gamma(u)\,du$ for every $u_0\in V_\infty$; and $\chi_\gamma(u_0)\ne1$
for some $u_0\in V_\gamma(E^+)\subseteq V_\infty$, as $\gamma$ does not agree. So the inner integral
is $0$.
\end{proof}

\begin{lemma}[Agreeing terms that preserve $A$]\label{pbe:scalar}
Let $\gamma\in\Gamma$ agree and $\gamma A=A$. Then $\gamma=\zeta\gamma_0$ with $\gamma_0\in R_0(E^+)$
and $\zeta\in E$, $\zeta\aubar\zeta=1$. If moreover $\gamma\in\Gamma_{c,s,c'}$,
then $\zeta\in\auO_E$ and $s\in\zeta\pbR$.
\end{lemma}
\begin{proof}
Let $P_A$ be the stabiliser of $A$, with unipotent radical $N_A$ (the elements acting trivially on
$A$, $A^\perp/A$ and $Y_0/A^\perp$). Then $N_A\subseteq U$, $\gamma$ normalises $N_A$, and $\gamma$
induces $a\in\auGL(A)$ and an isometry $\mathsf m$ of $Q=A^\perp/A=X_0\perp Ez$. Write $A_j=\langle
e_1,\dots,e_j\rangle$ and $e_j^*$ also for the linear form $h(\cdot,e_j^*)$ on $A$. We use the
following remark. Let $\psi_{E,f}$ be the finite part of $\psi_E$, so that $\psi_f(u)=\psi_{E,f}(-
\sigma(u)/(M\varpi))$. If $\mathsf l_1,\mathsf l_2$ are $E$-linear forms on an $E$-vector space with
$\psi_{E,f}(\mathsf l_1(\phi)/(M\varpi))=\psi_{E,f}(\mathsf l_2(\phi)/(M\varpi))$ for all $\phi$, then
$\mathsf l_1=\mathsf l_2$; for otherwise $\mathsf l_1-\mathsf l_2$ would be onto $E$ and $\psi_{E,f}$
would be trivial on $E$, which is dense in $\auA_{E,f}$.

\emph{The radical.} An element $\nu\in N_A(E^+)$ induces $\phi_\nu\in\operatorname{Hom}_E(Q,A)$,
$\phi_\nu(t)=\nu t-t$, every $\phi$ occurs, $\sigma(\nu)=e_\varrho^*(\phi_\nu(z))$, and $\phi_{\gamma^{
-1}\nu\gamma}=a^{-1}\phi_\nu\mathsf m$. As $\gamma$ agrees, the remark gives $e_\varrho^*(\phi(z))=e_
\varrho^*(a^{-1}\phi(\mathsf mz))$ for all $\phi$. With $\phi=\lambda(\cdot)v$, $\lambda$ a linear
form on $Q$ and $v\in A$: $\lambda(z)e_\varrho^*(v)=\lambda(\mathsf mz)e_\varrho^*(a^{-1}v)$. If
$\mathsf mz\notin Ez$ we could take $\lambda(z)=0\ne\lambda(\mathsf mz)$ and get $e_\varrho^*\circ a^{
-1}=0$, which is absurd. So $\mathsf mz=\zeta z$, $\zeta\aubar\zeta=1$ because $h(z,z)\ne0$, and
$e_\varrho^*\circ a^{-1}=\zeta^{-1}e_\varrho^*$ on $A$: $a$ preserves $A_{\varrho-1}$ and acts by
$\zeta$ on $A/A_{\varrho-1}$.

\emph{Descent on the flag.} Suppose, for some $1\le j<\varrho$, that $a$ preserves $A_i$ and acts
by $\zeta$ on $A_{i+1}/A_i$ for $j\le i<\varrho$. For $v\in A_j$ let $u_v\in U(E^+)$ act on $A$ by
$1+\phi$, $\phi(e_{j+1})=v$, $\phi(e_i)=0$ for $i\ne j+1$, by the inverse adjoint on $\langle
e_a^*\rangle$ and trivially on $X_0\perp Ez$. Then $\gamma^{-1}u_v\gamma$ is conjugate under $N_A
\subseteq U$ to the element of the same kind that acts on $A$ by $1+a^{-1}\phi a$, and $a^{-1}\phi a$
maps $A$ into $A_j$ and kills $A_j$; so this element lies in $U$, $u_v\in V_\gamma(E^+)$, and, $\psi$
being invariant under conjugation by $U$, $\sigma(\gamma^{-1}u_v\gamma)$ may be computed on it:
\[
\sigma(u_v)=e_j^*(v),\qquad\sigma(\gamma^{-1}u_v\gamma)=\sum_{i<\varrho}e_i^*\bigl(a^{-1}\phi(ae_{i+1})
\bigr)=e_j^*\bigl(a^{-1}\phi(ae_{j+1})\bigr)=\zeta\,e_j^*(a^{-1}v).
\]
(For $i>j$ the term vanishes because $a^{-1}\phi(\cdot)\in A_j$; for $i<j$ because $ae_{i+1}\in A_j$
is killed by $\phi$; and $ae_{j+1}\in\zeta e_{j+1}+A_j$.) By the remark, $e_j^*\circ a^{-1}=\zeta^{-1}
e_j^*$ on $A_j$: $a$ preserves $A_{j-1}$ and acts by $\zeta$ on $A_j/A_{j-1}$. By descending
induction, $\zeta^{-1}a$ is unitriangular on the flag.

So $\gamma_0=\zeta^{-1}\gamma$ is unitriangular on $A$ and induces on $Q$ an isometry that fixes
$z$, that is, an element $\eta$ of $H(E^+)$; then $\gamma_0\eta^{-1}\in U(E^+)$ and $\gamma_0\in
R_0(E^+)$. If $\gamma\in\Gamma_{c,s,c'}$, then $k=h_c^{-1}\gamma h_{c'}\in K$ and $ke_1=\zeta e_1$, so
$\zeta$ and $\zeta^{-1}$ are integral; $\zeta^{-1}k=h_c^{-1}\gamma_0h_{c'}\in K\cap R_0(\auA_f)$ has
its reduction in $\pbR$, and $s$ is this reduction times the residues of $\zeta$.
\end{proof}

Let
\[
\pbeGs=\{\gamma\in\Gamma:\pbg(\gamma)\ne0\}\ \cup\ \{\gamma\in\Gamma:\ \gamma\text{ agrees},\ \pbg(
\gamma)=0,\ \gamma A\ne A\},
\]
a stable set (for the last condition note $u\gamma u'A=u\gamma A$ and $uA=A$).

\begin{corollary}\label{pbe:star}
Assume \textup{(B1)--(B4)}. Let $\Omega$ avoid the scalar cosets, $s\in\Omega$ and $\flat\in\{\ ,
\flat\}$. Then $A^\flat_{\Gamma\smallsetminus R_0(E^+)}(c,s,c')=A^\flat_{\pbeGs}(c,s,c')$ for all
$c,c'$.
\end{corollary}
\begin{proof}
If $\gamma A=A$ then the first $\varrho$ columns of $Z(\gamma)$ vanish, so $\pbg(\gamma)=0$; hence
$\pbeGs\cap R_0(E^+)=\emptyset$. The set $(\Gamma\smallsetminus R_0(E^+))\smallsetminus\pbeGs$ is the
union of the stable set $T_1$ of its elements that do not agree, which contributes $0$ by
Lemma~\ref{pbe:vanish}, and the set $T_2$ of the $\gamma\notin R_0(E^+)$ that agree and have
$\gamma A=A$. By Lemma~\ref{pbe:scalar} an element of $T_2\cap\Gamma_{c,s,c'}$ lies in $\zeta R_0(E^+)$
with $s\in\zeta\pbR$; as $s\in\Omega$, $\zeta=1$, a contradiction. So $T_2\cap\Gamma_{c,s,c'}=\emptyset$.
\end{proof}

\subsubsection{The cost of the cut}\label{pbe:ss-cut}

\begin{proposition}[Cost of the cut]\label{pbe:cut}
Assume Theorem~\ref{pbe:ds} and let $W,W^\flat$ be as in Theorem~\ref{pbe:bessel}. For every stable
$T\subseteq\Gamma$ and every $\Omega\subseteq\auG$,
\[
\sum_{c,c'\in\mathcal C}\ \sum_{s\in\Omega}\bigl\|A_T(c,s,c')-A_T^\flat(c,s,c')\bigr\|^2\le b^{-40}\,
\pbDid\,\nu^2 .
\]
\end{proposition}
\begin{proof}
Fix $c$ and $s$. Grouping the terms in orbits of $U_c$ as in the proof of
Proposition~\ref{pb:unfold},
\[
\sum_{c'}\|A_T(c,s,c')-A_T^\flat(c,s,c')\|\le\int_{U_c\backslash G_\infty}\ \sum_{c'}\ \sum_{\gamma\in
\Gamma_{c,s,c'}/U_{c'}}\bigl\|W(y)W(\gamma^{-1}y)^*-W^\flat(y)W^\flat(\gamma^{-1}y)^*\bigr\|\,dy .
\]
The difference in the integrand is $(1-\chi(y)\chi(\gamma^{-1}y))\,W(y)W(\gamma^{-1}y)^*$ with $\chi=
\prod_\tau\chi_\tau$. It vanishes unless $\chi_\tau<1$ at some component of $y$ or of $\gamma^{-1}y$,
in which case $\pbeL(y)\pbeL(\gamma^{-1}y)\ge b$ by Lemma~\ref{pbe:len}(b); and by \eqref{pbe:eq-decay}
its norm is at most
\[
\mathsf A(y)\,\mathsf a_\infty\,\bigl(\pbeL(y)\pbeL(\gamma^{-1}y)\bigr)^{-2r^3f}\,\pbeL(\gamma^{-1}y)^{
-r^3f}\le b^{-2r^3f}\,\mathsf a_\infty\,\mathsf A(y)\,\pbeL(\gamma^{-1}y)^{-r^3f}.
\]
By Proposition~\ref{pbe:count}, the invariance of $\mathsf A$ under $U_\infty$, $\pbvol(U_c\backslash
U_\infty)=\nu$ and \eqref{pbe:eq-decay},
\[
\sum_{c'}\|A_T(c,s,c')-A_T^\flat(c,s,c')\|\le b^{-2r^3f}\cdot2b^{12n^3rf}\cdot\nu\,\mathsf a_\infty\int_{
U_\infty\backslash G_\infty}\mathsf A\le\nu\,b^{-2r^3f+12n^3rf+400r^2f+2kf+1}\le\nu\,b^{-k^3f};
\]
the last step is $12n^3r+400r^2+2k+1+k^3\le2r^3$, which holds because $k<r$ and $12n^3+500\le r$.
The sum of the squares over $c'$ is at most the square of the sum. Summing over $c$ and over all
$s\in\auG$, with $\#\auG\le M^{r^2f}\le b^{22r^2f}$, $\#\mathcal C\le\pbDid$ and $2k^3f\ge22r^2f+40$ (as
$k\ge r/2$, $r\ge100$), we get the proposition.
\end{proof}

\subsubsection{Links, and the three estimates for pairs of terms}\label{pbe:ss-pillars}

Let $Q=A^\perp/A=X_0\perp Ez$ and let $Q^+=X_0\perp Ez^+$ be the same vector space with the positive
definite form in which $z^+$ has norm $N$. Write $\pbeAp=A+Ez$, the span of the frame. Recall that
the link of $\gamma$ is $\Lambda_\gamma=\{([t],[\gamma^{-1}t]):t\in(A+\gamma A)^\perp\}\subseteq Q^2$.

\begin{lemma}[Links]\label{pbe:link}
Let $\gamma,\gamma'\in\Gamma$.
\begin{enumerate}
\item[(a)] $\pbg(\gamma)$ is, up to sign, the Gram determinant of the $2n$ vectors $f_0,\dots,
f_\varrho,\gamma f_0,\dots,\gamma f_\varrho$, and $\pbg(\gamma)=0$ if and only if $\pbeAp\cap\gamma
\pbeAp\ne0$. Both $\pbg$ and the link are invariant under $\gamma\mapsto u\gamma u'$, $u,u'\in
U(E^+)$.
\item[(b)] $\dim\Lambda_\gamma=k+1$, and $\Lambda_\gamma$ is the set of the pairs $(y_0+\lambda'z,\
y_0'+\lambda z)$ with $\Upsilon_\gamma(y_0+\lambda z^+)=y_0'+\lambda'z^+$, for an isometry
$\Upsilon_\gamma$ of $Q^+$. So $\Lambda_\gamma=\Lambda_{\gamma'}$ if and only if $\Upsilon_\gamma=
\Upsilon_{\gamma'}$.
\item[(c)] Let $X_0^\gamma=X_0\cap\Upsilon_\gamma^{-1}X_0$. Then $\dim X_0^\gamma=k-1+\dim(\pbeAp\cap
\gamma\pbeAp)-\dim(A\cap\gamma A)$. In particular $\dim X_0^\gamma=k-1$ if $\pbg(\gamma)\ne0$, and
$\dim X_0^\gamma=k$ if $\pbg(\gamma)=0$ and $\pbeAp\cap\gamma\pbeAp\not\subseteq A$.
\item[(d)] Let $\pbeCL(\gamma,\gamma')$ be the space of the $x\in X_0$ such that $(x,x')\in\Lambda_
\gamma\cap\Lambda_{\gamma'}$ for some $x'\in X_0$, and $i(\gamma,\gamma')=k+1-\dim(\Lambda_\gamma\cap
\Lambda_{\gamma'})$. Then $k-\dim\pbeCL(\gamma,\gamma')\le i(\gamma,\gamma')+k-\dim X_0^\gamma$.
\end{enumerate}
\end{lemma}
\begin{proof}
(a) The Gram matrix of the $2n$ vectors is $\bigl(\begin{smallmatrix}\mathsf D&Z^*\\Z&\mathsf D
\end{smallmatrix}\bigr)$ with $Z=Z(\gamma)$ up to transposition and $\mathsf D$ the Gram matrix of
the frame, whose only non-zero entry is $-N$ in the last place. For square matrices $\mathsf B,
\mathsf C$ one has $\det\bigl(\begin{smallmatrix}\mathsf B&\mathsf D\\\mathsf D&\mathsf C
\end{smallmatrix}\bigr)=\det\mathsf B\det\mathsf C-N^2\delta_{n-1}(\mathsf B)\delta_{n-1}(\mathsf C)$
(for invertible $\mathsf B$ this is $\det\mathsf B\cdot\det(\mathsf C-\mathsf D\mathsf B^{-1}\mathsf
D)$, where $\mathsf D\mathsf B^{-1}\mathsf D$ has the only entry $N^2\delta_{n-1}(\mathsf B)/\det
\mathsf B$ in the last place; the general case follows, both sides being polynomials). So the Gram
determinant is $\pm(\det Z\,\aubar{\det Z}-N^2\delta_{n-1}\aubar{\delta_{n-1}})=\pm\pbg(\gamma)$. If $\pbeAp\cap\gamma\pbeAp\ne0$ the $2n$ vectors are dependent and
the determinant is $0$. Conversely, let it be $0$ with $\pbeAp\cap\gamma\pbeAp=0$; then the span
$\mathsf V$ of the $2n$ vectors is degenerate: there is $v\ne0$ in $\mathsf V\cap\mathsf V^\perp$.
At a real place $\pbeAp$ is a negative semidefinite subspace of dimension $n$, the largest
dimension possible in signature $(k+\varrho,n)$; $v$ is orthogonal to $\pbeAp$ and to itself, so
$\pbeAp+\C v$ is negative semidefinite and $v\in\pbeAp$. In the same way $v\in\gamma\pbeAp$, a
contradiction. For the invariance: $uf$ and $u'f$ are obtained from $f$ by unitriangular matrices,
so the Gram determinant of $(f,u\gamma u'f)$, which is that of $(u^{-1}f,\gamma u'f)$, is that of
$(f,\gamma f)$; and for $t\perp A+u\gamma A$ the vector $t_1=u^{-1}t$ is orthogonal to $A+\gamma A$,
$[t_1]=[t]$ and $[(u\gamma u')^{-1}t]=[\gamma^{-1}t_1]$ in $Q$.

(b) Let $I=A\cap\gamma A$. The space $(A+\gamma A)^\perp$ has dimension $r-2\varrho+\dim I=k+1+\dim
I$ and contains $I$, which is the kernel of $t\mapsto([t],[\gamma^{-1}t])$. So $\dim\Lambda_\gamma=k+1$.
Since $h(t,t)=h(\gamma^{-1}t,\gamma^{-1}t)$, a pair $(y_0+\lambda'z,y_0'+\lambda z)$ of $\Lambda_\gamma$
satisfies $|y_0|^2-N|\lambda'|^2=|y_0'|^2-N|\lambda|^2$, that is, $y_0+\lambda z^+$ and $y_0'+\lambda'
z^+$ have the same norm in $Q^+$. So the two projections of the $(k+1)$-dimensional space of the
pairs $(y_0+\lambda z^+,y_0'+\lambda'z^+)$ to $Q^+$ are injective, hence bijective, and this space is
the graph of an isometry.

(c) $X_0^\gamma$ is the set of the first members of the pairs of $\Lambda_\gamma$ with $\lambda=
\lambda'=0$, that is, with $t\perp z$ and $\gamma^{-1}t\perp z$; so it is the image of $(\pbeAp+
\gamma\pbeAp)^\perp$ under $t\mapsto[t]$, whose kernel is $I$ (note $I\perp\pbeAp+\gamma\pbeAp$; and if $[t]=0$ then
$[\gamma^{-1}t]=0$ by (b): the pair $([t],[\gamma^{-1}t])$ has $y_0=0$ and $\lambda=-h(t,\gamma z)/N=0$ because $t\perp\gamma\pbeAp$, so it is $0$ by the injectivity of the first projection; hence $t\in I$).
Hence $\dim X_0^\gamma=r-2n+\dim(\pbeAp\cap\gamma\pbeAp)-\dim I$, and $r-2n=k-1$. The last claims
follow from (a) and from $\dim X_0^\gamma\le k$.

(d) In the graph of $\Upsilon_\gamma$, of dimension $k+1$, the graph of $\Upsilon_\gamma|X_0^\gamma$
and the image of $\Lambda_\gamma\cap\Lambda_{\gamma'}$ meet in a space of dimension at least $\dim
X_0^\gamma+(k+1-i)-(k+1)$, and the first projection maps it injectively into $\pbeCL(\gamma,
\gamma')$.
\end{proof}

Put
\[
\pbet=b^{kf/4},\qquad\pbenp=\frac{\#\pbR\cdot\#\mathcal C}{N^{kf}},\qquad\mathcal D(\gamma)=\begin{cases}
\pbDmain&\text{if }\pbg(\gamma)\ne0,\\ \pbet^{-16}\,\pbDid&\text{if }\pbg(\gamma)=0.\end{cases}
\]
Let $U^0\subseteq U$ be the subgroup that fixes the frame and $U_c^0=U_c\cap U^0$.

\begin{theorem}[Size, gain and count]\label{pbe:pillars}
Let $W^\flat$ be as in Theorem~\ref{pbe:bessel}. For all classes $c,c'$ there are a set $\pbeN(c,c')
\subseteq\pbeGs$ that contains at most one element of each double coset $U_c^0\gamma U_{c'}^0$, a
function $\mathsf G=\mathsf G_{c,c'}\colon\pbeN(c,c')\to(0,b^{n^3f}]$ (the \emph{gauge}), and for
$\gamma,\gamma'\in\pbeN(c,c')$ over a common block a real number $\mathsf d(\gamma,\gamma')=\mathsf d(
\gamma',\gamma)\ge1$ (the discriminant of the common lattice), with the following properties. Put
$\pbesz(\gamma)=\pbet^2\bigl(\mathcal D(\gamma)/(\pbenp\,\mathsf G(\gamma)^k)\bigr)^{1/2}$.
\begin{enumerate}
\item[(R)] If $\gamma\in\pbeN(c,c')$ and $\pbg(\gamma)=0$, then $\pbeAp\cap\gamma\pbeAp\not\subseteq
A$.
\item[(S)] \emph{(Size.)} For every stable $T\subseteq\pbeGs$ and every block $(c,s,c')$,
\[
\bigl\|A_T^\flat(c,s,c')\bigr\|\le\nu\sum_{\gamma\in\pbeN(c,c')\cap T\cap\Gamma_{c,s,c'}}\pbesz(\gamma).
\]
\item[(G)] \emph{(Gain.)} For $\gamma,\gamma'\in\pbeN(c,c')\cap\Gamma_{c,s,c'}$, with $i=i(\gamma,
\gamma')$,
\[
\bigl(\mathsf d(\gamma,\gamma')\cdot\mathsf x\bigr)^2\le\Bigl(\frac2M\Bigr)^{2if}\,\mathsf G(\gamma)\,
\mathsf G(\gamma'),\qquad\mathsf x=\begin{cases}\mathsf{mn}(s)&\text{if }i=0\text{ and }\pbg(\gamma)
\ne0,\\1&\text{otherwise.}\end{cases}
\]
\item[(C)] \emph{(Count.)} For every integer $0\le l\le n^3$,
\[
\sum_{c,c'\in\mathcal C}\ \sum_{(\gamma,\gamma')}\mathsf d(\gamma,\gamma')^{-k}\le\pbet^{\,l+2}\,\pbenp ,
\]
the inner sum over the pairs of elements of $\pbeN(c,c')$ over a common block with $k-\dim\pbeCL(
\gamma,\gamma')=l$ and $\mathsf d(\gamma,\gamma')\le b^{n^3f}$.
\end{enumerate}
\end{theorem}
\smallskip\noindent The exact form of the theorem, with the set $\pbeN(c,c')$ restricted to agreeing blocks and with the gauge, the common lattice and nearness defined, is Theorem~\ref{ppi:pillars}; every use made below of the present theorem is covered by that form: Proposition~\ref{ppi:check} shows that Lemmas~\ref{pbe:rank} and~\ref{pbe:pairs}, Proposition~\ref{pbe:two} and Theorem~\ref{pbe:XL1} hold with Theorem~\ref{ppi:pillars} in its place, with the same proofs. In particular \textup{(S)} is used only for the stable sets $T_0$, $T_\Lambda$ and $\pbeGs$, and in that form the sum in \textup{(S)} runs over elements of agreeing blocks only; the elements of a non-agreeing block contribute $0$ to $A^\flat_T$ by Lemma~\ref{pbe:vanish}. \textup{(R)} is proved in Proposition~\ref{ppi:cell}, \textup{(G)} in \S\ref{pbg:sec}, \textup{(C)} in \S\ref{pbc:sec}, and \textup{(S)} in \S\ref{pbs:sec} from the open-cell bound (Theorem~\ref{pbo:O1}).\par\smallskip

\subsubsection{Proof of \textup{(X)} and \textup{(L1)} from Theorem~\ref{pbe:pillars}}\label{pbe:ss-XL1}

We say that $\gamma,\gamma'$ are \emph{in the same link class} if $\pbg(\gamma)\ne0$, $\pbg(\gamma')
\ne0$ and $\Lambda_\gamma=\Lambda_{\gamma'}$.

\begin{lemma}\label{pbe:rank}
Assume Theorem~\ref{pbe:pillars}. Let $\gamma,\gamma'\in\pbeN(c,c')\cap\Gamma_{c,s,c'}$, $i=i(\gamma,
\gamma')$ and $l=k-\dim\pbeCL(\gamma,\gamma')$. Then:
\begin{enumerate}
\item[(a)] $l\le i+1$, and $l\le i$ if $\pbg(\gamma)=0$;
\item[(b)] if $i=0$, then $\pbg(\gamma)=0$ if and only if $\pbg(\gamma')=0$;
\item[(c)] $\mathsf d(\gamma,\gamma')\le b^{n^3f}$ and $l\le n^3$.
\end{enumerate}
\end{lemma}
\begin{proof}
(a) By Lemma~\ref{pbe:link}(c),(d) and (R), $k-\dim X_0^\gamma$ is $1$ if $\pbg(\gamma)\ne0$ and $0$
otherwise. (b) If $i=0$ then $\Upsilon_\gamma=\Upsilon_{\gamma'}$ and $X_0^\gamma=X_0^{\gamma'}$, and
the dimension of this space decides whether $\pbg$ vanishes. (c) By (G), with $\mathsf x\ge1$ and
$\mathsf G\le b^{n^3f}$, $\mathsf d^2\le(2/M)^{2if}b^{2n^3f}$. So $\mathsf d\le b^{n^3f}$, and, as
$\mathsf d\ge1$ and $M/2\ge b^{18}$, $18i\le n^3$; hence $l\le i+1\le n^3$.
\end{proof}

\begin{lemma}[Sums over pairs]\label{pbe:pairs}
Assume Theorem~\ref{pbe:pillars}. For all $c,c'$ let $\mathcal Q(c,c')$ be a set of pairs $(\gamma,
\gamma')$ of elements of $\pbeN(c,c')$ over a common block, and let $B\ge0$. Suppose that for every
such pair there are $\eta>0$ and $\mathsf D\ge\max(\mathcal D(\gamma),\mathcal D(\gamma'))$ with
\[
(\mathsf d(\gamma,\gamma')\,\eta)^2\le\mathsf G(\gamma)\mathsf G(\gamma')\qquad\text{and}\qquad\pbet^{\,
l+6}\,\mathsf D\,\eta^{-k}\le B,\quad l=k-\dim\pbeCL(\gamma,\gamma').
\]
Then $\sum_{c,c'}\sum_{(\gamma,\gamma')\in\mathcal Q(c,c')}\pbesz(\gamma)\pbesz(\gamma')\le\pbet\,B$.
\end{lemma}
\begin{proof}
For such a pair, $\pbesz(\gamma)\pbesz(\gamma')\le\pbet^4\mathsf D\,\pbenp^{-1}(\mathsf G(\gamma)
\mathsf G(\gamma'))^{-k/2}\le\pbet^4\mathsf D\eta^{-k}\pbenp^{-1}\mathsf d^{-k}\le B\,\pbet^{-l-2}
\pbenp^{-1}\,\mathsf d(\gamma,\gamma')^{-k}$. By Lemma~\ref{pbe:rank}(c) every pair is counted in (C)
for its $l\le n^3$, so the pairs with a given $l$ contribute at most $B$, and the total is at most
$(n^3+1)B\le b^4B\le\pbet B$.
\end{proof}

\begin{proposition}\label{pbe:two}
Assume Theorem~\ref{pbe:pillars}.
\begin{enumerate}
\item[(a)] The sum of $\pbesz(\gamma)\pbesz(\gamma')$ over all $c,c'$ and all pairs of elements of
$\pbeN(c,c')$ over a common block that are \emph{not} in the same link class is at most $b^{-kf}
\pbDid$.
\item[(b)] Let $\mathsf z\ge1$. The same sum over the pairs \emph{in} the same link class and over a
block $(c,s,c')$ with $\mathsf{mn}(s)\ge\mathsf z$ is at most $(b^{2f}/\mathsf z)^k\,\pbDmain$.
\end{enumerate}
\end{proposition}
\begin{proof}
We apply Lemma~\ref{pbe:pairs}. Recall $\pbet^4=b^{kf}$, $b^{11kf}\pbDid\le\pbDmain\le(M/b^6)^{kf}\pbDid$
(Lemma~\ref{pb:sizes}); in particular $\mathcal D(\gamma)\le\pbDmain$ always.

(a) Let $B=\pbet^{-10}\pbDid$; the bound will be $\pbet^{-9}\pbDid\le\pbet^{-4}\pbDid$. Let $(\gamma,
\gamma')$ be a pair, $i=i(\gamma,\gamma')$. \emph{If $i=0$}, the links are equal, so by
Lemma~\ref{pbe:rank}(b) and the hypothesis $\pbg(\gamma)=\pbg(\gamma')=0$, and $l=0$ by
Lemma~\ref{pbe:rank}(a). Take $\eta=1$ (by (G)) and $\mathsf D=\pbet^{-16}\pbDid$: $\pbet^6\mathsf D=B$.
\emph{If $i\ge1$}, take $\eta=(M/2)^{if}$ (by (G)) and $\mathsf D=\pbDmain$; $l\le i+1$, and with
$\mathsf y=\pbet(2/M)^{kf}\le1$,
\[
\pbet^{\,l+6}\,\pbDmain\,\eta^{-k}\le\pbet^7\mathsf y^i\pbDmain\le\pbet^7\mathsf y\Bigl(\frac M{b^6}
\Bigr)^{kf}\pbDid=\pbet^8\Bigl(\frac2{b^6}\Bigr)^{kf}\pbDid\le\pbet^{8-20}\pbDid\le B .
\]
(b) Here $i=0$ and $\pbg(\gamma)\ne0$, so (G) gives $(\mathsf d\,\mathsf z)^2\le(\mathsf d\,\mathsf{
mn}(s))^2\le\mathsf G(\gamma)\mathsf G(\gamma')$, and $l\le1$. Take $\eta=\mathsf z$, $\mathsf D=
\pbDmain$, $B=\pbet^7\pbDmain\mathsf z^{-k}$; the bound is $\pbet^8\mathsf z^{-k}\pbDmain=(b^{2f}/
\mathsf z)^k\pbDmain$.
\end{proof}

\begin{theorem}[Theorem~\ref{pb:amp}\textup{(X)} and \textup{(L1)}]\label{pbe:XL1}
Assume Theorems~\ref{pbe:ds} and~\ref{pbe:pillars}. Then, for the Bessel functions of
Theorem~\ref{pbe:bessel}, parts \textup{(X)} and \textup{(L1)} of Theorem~\ref{pb:amp} hold.
\end{theorem}
\begin{proof}
(L1) Let $\mathsf{mn}\ge\mathsf z$ on $\Omega$. The sets $T_\Lambda$ are stable subsets of $\pbeGs$.
By (S),
\[
\|A^\flat_{T_\Lambda}(c,s,c')\|^2\le\nu^2\sum_{\gamma,\gamma'\in\pbeN(c,c')\cap T_\Lambda\cap\Gamma_{c,
s,c'}}\pbesz(\gamma)\pbesz(\gamma').
\]
The pairs that occur are in the same link class and over the block $(c,s,c')$, and a pair determines
$s$ and $\Lambda$; so, summing over $c,c'$, $s\in\Omega$ and $\Lambda$, $\pbLM(\Omega)\le\nu^2(b^{2f}/
\mathsf z)^k\pbDmain$ by Proposition~\ref{pbe:two}(b).

(X) Let $\Omega$ avoid the scalar cosets and $\pbLM(\Omega)\le C\nu^2$. Fix a block $(c,s,c')$ with
$s\in\Omega$. The set $\pbeGs$ is the disjoint union of the stable sets $T_\Lambda$ and $T_0=\{\gamma
\in\pbeGs:\pbg(\gamma)=0\}$. The amplitude sums converge absolutely, so $A^\flat_{\pbeGs}=\sum_\kappa
A^\flat_{T_\kappa}$ over the pieces $\kappa\in\{0\}\cup\{\Lambda\}$, and
\[
\|A^\flat_{\pbeGs}(c,s,c')\|^2\le\Bigl(\sum_\kappa\|A^\flat_{T_\kappa}\|\Bigr)^2=\sum_\Lambda\|A^\flat_{
T_\Lambda}\|^2+\sum_{(\kappa,\kappa')\ne(\Lambda,\Lambda)}\|A^\flat_{T_\kappa}\|\,\|A^\flat_{T_{\kappa'}}\|.
\]
By (S) each product in the last sum is at most $\nu^2\sum\pbesz(\gamma)\pbesz(\gamma')$ over $\gamma\in
\pbeN(c,c')\cap T_\kappa\cap\Gamma_{c,s,c'}$, $\gamma'\in\pbeN(c,c')\cap T_{\kappa'}\cap\Gamma_{c,s,c'}$;
these pairs are not in the same link class, and a pair determines $s$, $\kappa$ and $\kappa'$. (Here (S) is applied to $T_\kappa$, which may contain elements of non-agreeing blocks; by Lemma~\ref{pbe:vanish} these contribute $0$ to $A^\flat_{T_\kappa}$, so only elements of $\pbeN(c,c')$, which lie over agreeing blocks, enter the bound.)
Summing over the blocks and using Proposition~\ref{pbe:two}(a),
\[
\sum_{c,c'}\sum_{s\in\Omega}\|A^\flat_{\pbeGs}(c,s,c')\|^2\le\pbLM(\Omega)+\nu^2b^{-kf}\pbDid\le\nu^2\bigl(C
+b^{-kf}\pbDid\bigr).
\]
By Proposition~\ref{pb:besselterms}, Corollary~\ref{pbe:star}, the triangle inequality in $\ell^2$
and Proposition~\ref{pbe:cut} (for $T=\pbeGs$),
\[
\Bigl(\sum_{s\in\Omega}\pbHS{P_s-\pbrpsi(s)}^2\Bigr)^{1/2}=\frac1\nu\Bigl(\sum_{c,c',s}\|A_{\pbeGs}\|^2
\Bigr)^{1/2}\le b^{-20}\sqrt{\pbDid}+\sqrt{C+b^{-kf}\pbDid}\le\sqrt C+b^{-15}\sqrt{\pbDid},
\]
as $kf\ge40$.
\end{proof}

\subsubsection{The inputs of this section}\label{pbe:ss-remains}

Hypothesis~\ref{pb:bessel} and Theorem~\ref{pb:amp}\textup{(X)}, \textup{(L1)} rest on the following
statements, proved elsewhere:
\begin{enumerate}
\item[(i)] Theorem~\ref{pbe:ds}: proved in \S\ref{pds:sec}, {\cite[\S3.2]{P49}} and~\S\ref{pbp:sec}, modulo the published results cited there.
\item[(ii)] Theorem~\ref{pbe:pillars}: \textup{(O1)} (Theorem~\ref{pbo:O1}, for \textup{(S)}), and the statements \textup{(G$'$)}, \textup{(C1)}, \textup{(C2)} of \S\ref{ppi:sec}, proved in \S\S\ref{pbo:sec}, \ref{pbg:sec}, \ref{pbc:sec}; \textup{(S)} also uses Lemma~\ref{pb:covol}.
\item[(iii)] The inequality \eqref{pbe:eq-kappa} for the generator $\varpi$ (Lemma~\ref{pds:kappa}), and the
finite dimension of $\auAz$ ({\cite[Proposition~8.4(a)]{P49}}); $\auAz$ is by definition the space of bounded
solutions of \eqref{pbe:eq-repro} ({\cite[\S14.3]{P46}}).
\end{enumerate}



% ===== B-main-term.tex =====
% proofs/B-main-term.tex
% The induced character (Theorem pb:est (E4)) and the main term (Theorem pb:amp (L2)).
% Labels of this file start with pbm:.
\providecommand{\auO}{\mathcal O}
\providecommand{\auPl}{\Phi}
\providecommand{\auS}{\mathcal S}
\providecommand{\auG}{\overline G}
\providecommand{\auDX}{D_X}
\providecommand{\auNm}{\operatorname{Nm}}
\providecommand{\auU}{\mathrm U}
\providecommand{\auGL}{\mathrm{GL}}
\providecommand{\auA}{\mathbb A}
\providecommand{\auF}{\mathbb F}
\providecommand{\aurd}{\mathrm{rd}_E}
\providecommand{\aubar}[1]{\overline{#1}}
\providecommand{\auf}[1]{\texttt{#1}}
\providecommand{\panot}[1]{\par\smallskip{\raggedright\noindent\textbf{[Not proved here:} #1\textbf{]}\par}\smallskip}
\providecommand{\pbR}{\overline R_0}
\providecommand{\pbUb}{\overline U}
\providecommand{\pbHb}{\overline H}
\providecommand{\pbvol}{\operatorname{vol}}
\providecommand{\pbDid}{\mathcal D_{\mathrm{id}}}
\providecommand{\pbDmain}{\mathcal D_{\mathrm{main}}}
\providecommand{\pbdeg}{\operatorname{deg}}
\providecommand{\pbmass}{\operatorname{mass}}
\providecommand{\pbFix}{\operatorname{Fix}}
\providecommand{\pbg}{\mathsf g}
\providecommand{\pbLM}{\mathrm{LM}}
\providecommand{\pbnat}{\natural}
\providecommand{\pbip}[2]{\langle#1,#2\rangle}
% macros of this file (prefix \pbm)
\providecommand{\pbmD}{\mathsf D}
\providecommand{\pbmPhi}{\varphi}
\providecommand{\pbmGf}{G^{\flat}}
\providecommand{\pbmXp}{X^{+}}
\providecommand{\pbmOrb}{\mathcal O}
\providecommand{\pbmKl}{\mathsf{Kl}}
\providecommand{\pbmT}{\mathcal T}
\providecommand{\pbmAT}{\mathbb A_T}
\providecommand{\pbmRT}{R_T}
\providecommand{\pbmE}{\mathsf E}
\providecommand{\pbmn}{\mathsf n}
\providecommand{\pbmden}{\mathsf{den}}
\providecommand{\pbmvol}{\mathsf v}
\providecommand{\pbmdec}{\mathsf d}
\providecommand{\pbmaS}{\mathsf a_*}
\providecommand{\pbmrS}{\mathsf r_*}
\providecommand{\pbmnuS}{\nu_*}
\providecommand{\pbmcG}{c_G}
\providecommand{\pbmTh}{\Theta}
\providecommand{\pbmSt}{\operatorname{Stab}}

\subsection{The induced character and the main term}\label{pbm:sec}

This section concerns two of the statements to which \S\ref{pb:sec} reduced the Spanning Theorem ({\cite[Theorem~14.14]{P46}}):
the inequality \textup{(E4)} of Theorem~\ref{pb:est}, and the statement \textup{(L2)} of
Theorem~\ref{pb:amp}. We use the notation of \S\ref{pb:ss-setting}.

\emph{What is proved.}
\begin{itemize}
\item[(I)] \emph{The inequality \textup{(E4)} is proved in full} (\S\ref{pbm:ss-E4},
Theorem~\ref{pbm:E4}). The proof is short and does not use the analysis of the big cell: the
function that counts the elements of $\pbUb\cap h^{-1}\pbUb h$ on which $\psi_U$ and its conjugate
by $h$ agree is of positive type, so Lemma~\ref{pb:singular} applies to it; two exact recursions,
which remove one hyperbolic pair, or one vector of the frame, at a time, and a crude count do the
rest.
\item[(II)] \emph{The statement \textup{(L2)}} is deduced, by complete arguments, from one precisely
stated theorem on local objects (Theorem~\ref{pbm:inputs}, proved in {\cite[\S\S2--6]{P52}})
and the value of the covolume $\nu$ (Lemma~\ref{pb:covol}): the
existence, at each real place and at each finite set of finite places, of a system of forms on
pairs of heads with explicit bounds, and an inequality (\emph{Siegel's step}) that bounds sums of
squares of class sums by the top forms. We prove the description of the terms over
open elements (big cell, heads, links; Propositions~\ref{pbm:cell} and~\ref{pbm:classsum}); the
identification of an orbit sum with a Kloosterman sum times an orbital integral
(Proposition~\ref{pbm:orbit}); the sieve by one-sided Poisson summation, with all its constants
(Lemma~\ref{pbm:poisson}, Theorem~\ref{pbm:sieve}); the evaluation of the constant of the main
term, which is where $\pbDmain=\auDX^2/(\pbdeg\cdot\pbmass_r)$ comes from (Lemma~\ref{pbm:const});
and the deduction of \textup{(L2)} (Theorem~\ref{pbm:L2}).
\end{itemize}
\subsubsection{The induced character: proof of \textup{(E4)}}\label{pbm:ss-E4}

\paragraph*{Chain data}
Let $\auF$ be a finite field with $q$ elements and $\psi_1,\psi_2$ two non-trivial additive
characters of $\auF$. Let $i,i'\ge0$ and $m\ge1$ be integers. A \emph{chain datum of type $(i,i')$}
is a vector space $V$ over $\auF$ with a basis
\[
e_1,\dots,e_i;\qquad\text{a basis of a subspace }X'\text{ of dimension }m;\qquad e^*_{i'},\dots,e^*_1,
\]
together with a vector $z\in X'$, $z\ne0$, and a linear form $\zeta\ne0$ on $X'$, extended by zero
on the $e_a$ and the $e_a^*$. Write $r=\dim V=i+m+i'$, $\epsilon_a$ and $\epsilon_a^*$ for the
coordinate forms of $e_a$ and $e_a^*$, $A_a=\langle e_1,\dots,e_a\rangle$ $(0\le a\le i)$, $A=A_i$
and $B_a=A\oplus X'\oplus\langle e^*_{i'},\dots,e^*_{a+1}\rangle$ $(0\le a\le i')$, so that
$B_{i'}=A\oplus X'$ and $B_0=V$. Let $U=U_\pbmD$ be the group of the $u\in\auGL(V)$ with
\[
(u-1)A_a\subseteq A_{a-1}\ (1\le a\le i),\qquad(u-1)B_{i'}\subseteq A,\qquad(u-1)B_a\subseteq B_{a+1}\
(0\le a<i');
\]
it is the unipotent radical of the parabolic subgroup with blocks $1,\dots,1,m,1,\dots,1$. The
function
\[
\psi_\pbmD(u)=\psi_1\Bigl(\sum_{a=1}^{i-1}\epsilon_a(ue_{a+1})+\epsilon_i(uz)\Bigr)\,\psi_2\Bigl(\sum_{a=1}^{
i'-1}\epsilon^*_{a+1}(ue^*_a)+\zeta(ue^*_{i'})\Bigr)
\]
(the first factor is $1$ if $i=0$, the second is $1$ if $i'=0$) is a character of $U$: its
arguments are sums of entries of $u$ just above the diagonal blocks, and these are additive on
$U$. We call $e_1$ (or $z$, if $i=0$) the \emph{first vector of the left chain}. For $h\in\auGL(V)$
put
\[
f_\pbmD(h)=\sum_{u\in U\cap h^{-1}Uh}\psi_\pbmD(huh^{-1})\,\aubar{\psi_\pbmD(u)},\qquad\pbmPhi(\pbmD)=\frac1{\#
\auGL(V)}\sum_{h\in\auGL(V)}f_\pbmD(h).
\]

\begin{lemma}\label{pbm:pos}
Let $\pbmD$ be a chain datum, $U=U_\pbmD$, $\psi=\psi_\pbmD$, $f=f_\pbmD$ and $\mathsf G=\auGL(V)$.
\begin{enumerate}
\item[(a)] $f(h)$ is $\#(U\cap h^{-1}Uh)$ or $0$; in particular $f\ge0$. Moreover $f(uhu')=f(h)$
for $u,u'\in U$ and $f(h^{-1})=f(h)$.
\item[(b)] $f$ is of positive type on $\mathsf G$ (in the sense of \S\ref{pb:ss-finite}).
\item[(c)] Let $\tilde\psi$ be $\psi$ on $U$ and $0$ on $\mathsf G\smallsetminus U$, and
$\psi^{\pbnat}(s)=(\#\mathsf G)^{-1}\sum_{g\in\mathsf G}\tilde\psi(gsg^{-1})$. Then
$\sum_{s\in\mathsf G}|\psi^{\pbnat}(s)|^2=\pbmPhi(\pbmD)$.
\end{enumerate}
\end{lemma}
\begin{proof}
(a) $u\mapsto\psi(huh^{-1})\aubar{\psi(u)}$ is a character of the group $U\cap h^{-1}Uh$, and the
sum of a character over a finite group is the order of the group or $0$. The substitutions
$u\mapsto u'^{-1}uu'$ and the invariance of $\psi$ under conjugation by $U$ give $f(uhu')=f(h)$.
With $u=h^{-1}u_1h$ we get $f(h^{-1})=\sum_{u_1\in U\cap h^{-1}Uh}\psi(u_1)\aubar{\psi(hu_1h^{-1})}=
\aubar{f(h)}=f(h)$.
(b) In the space of functions on $\mathsf G$, with the orthonormal basis of the delta functions
$\delta_g$, let $v_a=\sum_{u\in U}\aubar{\psi(u)}\,\delta_{aua^{-1}}$ for $a\in\mathsf G$. Then, with
$s=b^{-1}a$,
\[
\pbip{v_a}{v_b}=\sum_{u,u'\in U}\psi(u)\aubar{\psi(u')}\,[aua^{-1}=bu'b^{-1}]=\sum_{u\in U\cap s^{-1}Us}
\psi(u)\aubar{\psi(sus^{-1})}=\aubar{f(b^{-1}a)}=f(a^{-1}b)
\]
by (a). (c) Expanding the square and substituting $s=g'^{-1}ug'$ with $u\in U$,
\[
\sum_s|\psi^{\pbnat}(s)|^2=\frac1{(\#\mathsf G)^2}\sum_{g,g'}\sum_{u\in U}\tilde\psi\bigl(gg'^{-1}u\,g'g^{
-1}\bigr)\aubar{\psi(u)}=\frac1{\#\mathsf G}\sum_{h\in\mathsf G}\ \sum_{u\in U,\ huh^{-1}\in U}\psi(huh^{-1})
\aubar{\psi(u)} .\qedhere
\]
\end{proof}

\begin{lemma}[Removing a hyperbolic pair]\label{pbm:pair}
Let $\pbmD$ be a chain datum of type $(i,i')$ with $i,i'\ge1$ on $V$, and let $\pbmD'$ be the chain
datum of type $(i-1,i'-1)$ on $V'=\langle e_2,\dots,e_i\rangle\oplus X'\oplus\langle e^*_{i'},\dots,
e^*_2\rangle$ with the same $z$, $\zeta$, $\psi_1$, $\psi_2$. Let $\Omega$ be the set of the
$h\in\auGL(V)$ with $\epsilon_1^*(he_1)\ne0$ and $\epsilon_1^*(h^{-1}e_1)\ne0$. Then
\[
\sum_{h\in\Omega}f_\pbmD(h)=\#\Omega\cdot\pbmPhi(\pbmD').
\]
\end{lemma}
\begin{proof}
Let $H_1=\ker\epsilon_1^*=B_1$, let $P_1$ be the stabiliser of the flag $\auF e_1\subset H_1$,
$N_1$ its unipotent radical (the $n$ with $(n-1)V\subseteq H_1$, $(n-1)H_1\subseteq\auF e_1$,
$ne_1=e_1$) and $L_1=\auF^\times\times\auGL(V')\times\auF^\times$ the subgroup of the $l=(t,g',t')$
with $le_1=te_1$, $l|_{V'}=g'$, $le_1^*=t'e_1^*$; so $P_1=N_1\rtimes L_1$. Let $w_1$ exchange $e_1$
and $e_1^*$ and fix $V'$; it normalises $L_1$, and $P_1\cap w_1P_1w_1^{-1}=L_1$.

\emph{$\Omega=N_1w_1L_1N_1$, and $(n,l,n')\mapsto nw_1ln'$ is a bijection $N_1\times L_1\times N_1\to
\Omega$.} For $h\in\Omega$ we have $he_1\notin H_1$ and $e_1\notin hH_1$, so $V=\auF e_1\oplus(H_1
\cap hH_1)\oplus\auF he_1$ and $H_1=\auF e_1\oplus(H_1\cap hH_1)$. Let $p\in\auGL(V)$ fix $e_1$, map
$H_1\cap hH_1$ onto $V'$ and $he_1$ to $e_1^*$. Then $p\in P_1$, and $phe_1=w_1e_1$, $phH_1=V'\oplus
\auF e_1^*=w_1H_1$, so that $w_1^{-1}ph\in P_1$. Hence $\Omega\subseteq P_1w_1P_1=N_1w_1L_1N_1$; the
other inclusion is clear, and $\#(P_1w_1P_1)=(\#P_1)^2/\#L_1=(\#N_1)^2\,\#L_1$ gives the bijection.

\emph{$U_\pbmD=N_1\rtimes U_{\pbmD'}$}, where $\auGL(V')$ is embedded in $\auGL(V)$ by letting it
fix $e_1$ and $e_1^*$. Indeed $N_1\subseteq U_\pbmD$ by the definitions ($(n-1)$ maps $V$ into
$H_1=B_1$ and $H_1$ into $\auF e_1=A_1$), every $u\in U_\pbmD$ lies in $P_1$, and its component in
$L_1$ is $(1,u',1)$ with $u'=n^{-1}u\in U_\pbmD\cap\auGL(V')=U_{\pbmD'}$. Moreover
$\psi_\pbmD|_{U_{\pbmD'}}=\psi_{\pbmD'}$: for $u'\in\auGL(V')$ the entries $\epsilon_1(u'e_2)$ (or
$\epsilon_1(u'z)$) and $\epsilon_2^*(u'e_1^*)$ (or $\zeta(u'e_1^*)$) vanish.

\emph{$f_\pbmD(nw_1ln')=f_{\pbmD'}(g')$ for $l=(t,g',t')$.} By Lemma~\ref{pbm:pos}(a) we may take
$n=n'=1$; put $h=w_1l$. Let $u=\nu u'$ with $\nu\in N_1$, $u'\in U_{\pbmD'}$. Then $huh^{-1}=(h\nu
h^{-1})(g'u'g'^{-1})$ with $h\nu h^{-1}\in w_1N_1w_1^{-1}$, the unipotent radical opposite to $N_1$,
and $g'u'g'^{-1}\in\auGL(V')\subseteq P_1$. If $huh^{-1}\in U_\pbmD\subseteq P_1$ then $h\nu h^{-1}\in
w_1N_1w_1^{-1}\cap P_1=1$, so $\nu=1$ and $g'u'g'^{-1}\in U_\pbmD\cap\auGL(V')=U_{\pbmD'}$. Hence
$f_\pbmD(h)=\sum_{u'\in U_{\pbmD'}\cap g'^{-1}U_{\pbmD'}g'}\psi_{\pbmD'}(g'u'g'^{-1})\aubar{\psi_{\pbmD'}(u')}=
f_{\pbmD'}(g')$.

Summing over $\Omega$ gives $(\#N_1)^2(q-1)^2\sum_{g'}f_{\pbmD'}(g')=(\#N_1)^2\,\#L_1\cdot\pbmPhi(\pbmD')$.
\end{proof}

\begin{lemma}[Fixing the first vector of the left chain]\label{pbm:first}
Let $\pbmD$ be a chain datum of type $(i,i')$ with $i\ge1$ on $V$, of dimension $r$, and let
$\pbmD_1$ be the chain datum of type $(i-1,i')$ on $V_1=\langle e_2,\dots,e_i\rangle\oplus X'\oplus
\langle e^*_{i'},\dots,e^*_1\rangle$ with the same $z$, $\zeta$, $\psi_1$, $\psi_2$. Let $v_2$ be
the first vector of the left chain of $\pbmD_1$ ($v_2=e_2$ if $i\ge2$, $v_2=z$ if $i=1$). Then
\[
\sum_{h\in\auGL(V),\ he_1=e_1}f_\pbmD(h)=q^{2(r-1)}\sum_{g\in\auGL(V_1),\ gv_2=v_2}f_{\pbmD_1}(g).
\]
\end{lemma}
\begin{proof}
Let $S$ be the stabiliser of $e_1$ and $N=\{1+e_1\otimes\xi:\xi\in V^*,\ \xi(e_1)=0\}\cong V_1^*$, of
order $q^{r-1}$; then $S=N\rtimes\auGL(V_1)$, where $\auGL(V_1)$ fixes $e_1$. As in the proof of
Lemma~\ref{pbm:pair}, $N\subseteq U_\pbmD\subseteq S$ and $U_\pbmD=N\rtimes U_{\pbmD_1}$, so that an
element of $S$ lies in $U_\pbmD$ if and only if its component in $\auGL(V_1)$ lies in
$U_{\pbmD_1}$; and $\psi_\pbmD|_{U_{\pbmD_1}}=\psi_{\pbmD_1}$, $\psi_\pbmD(1+e_1\otimes\xi)=\psi_1(\xi(
v_2))$. Every element of $U_{\pbmD_1}$ fixes $v_2$.

Let $h=ng\in S$ and $u=n_1u_1\in U_\pbmD$, with $n,n_1\in N$, $g\in\auGL(V_1)$, $u_1\in U_{\pbmD_1}$,
$n_1=1+e_1\otimes\xi_1$. Then $huh^{-1}\in U_\pbmD$ if and only if $u_1':=gu_1g^{-1}\in U_{\pbmD_1}$,
and in that case
\[
huh^{-1}=n\cdot gn_1g^{-1}\cdot u_1'n^{-1}u_1'^{-1}\cdot u_1',\qquad\psi_\pbmD(huh^{-1})=\psi_1\bigl(
\xi_1(g^{-1}v_2)\bigr)\,\psi_{\pbmD_1}(u_1'),
\]
because $gn_1g^{-1}=1+e_1\otimes\xi_1g^{-1}$ and because $\psi_\pbmD(u_1'n^{-1}u_1'^{-1})=\psi_\pbmD(n)^{-1}$,
as $u_1'$ fixes $v_2$. Summing over $\xi_1\in V_1^*$, by the orthogonality of the characters of
$V_1^*$ ($\psi_1$ is not trivial),
\[
f_\pbmD(ng)=\sum_{\xi_1}\psi_1\bigl(\xi_1(g^{-1}v_2-v_2)\bigr)\cdot f_{\pbmD_1}(g)=q^{r-1}\,[gv_2=v_2]\,
f_{\pbmD_1}(g).
\]
The sum over $n\in N$ gives another factor $q^{r-1}$.
\end{proof}

\begin{lemma}[A crude bound]\label{pbm:crude}
Let $\pbmD$ be a chain datum of type $(0,i')$ on $W$, with $\dim W=n=m+i'$ and $m>i'$. Then
\[
\sum_{g\in\auGL(W),\ gz=z}f_\pbmD(g)\ \le\ 16\,(i'+1)\,q^{\,i'(i'+1)}\cdot\frac{\#\auGL(W)}{q^n-1}.
\]
\end{lemma}
\begin{proof}
Here $A=0$, so every $u\in U=U_\pbmD$ is the identity on $X'$. If $u$ and $gug^{-1}$ lie in $U$,
then $u-1$ vanishes on $X'+g^{-1}X'$; hence, by Lemma~\ref{pbm:pos}(a),
\[
f_\pbmD(g)\le\#(U\cap g^{-1}Ug)\le q^{\,c(g)\,n},\qquad c(g)=\dim W-\dim(X'+g^{-1}X')=\dim(X'\cap g^{-1}
X')-(m-i').
\]
Let $g$ run over the stabiliser $S_z$ of $z$, of order $\#\auGL(W)/(q^n-1)$. Both $X'$ and
$g^{-1}X'$ contain $z$; let $\bar X=X'/\auF z\subset\bar W=W/\auF z$ and let $\bar g$ be the map
induced by $g$. The homomorphism $S_z\to\auGL(\bar W)$ is onto, and $c(g)\ge c$ means that
$\bar X\cap\bar g^{-1}\bar X$ contains a subspace $D$ of dimension $d=m-1-i'+c$. For a fixed $D$ of
dimension $d$ in $\bar X$, the subspace $\bar gD$ is equidistributed among the subspaces of
dimension $d$ of $\bar W$. With the Gaussian binomial coefficients, which satisfy $q^{d(a-d)}\le
\binom ad_q\le4q^{d(a-d)}$, the proportion of the $g\in S_z$ with $c(g)\ge c$ is therefore at most
\[
\binom{m-1}d_q^{\,2}\Big/\binom{n-1}d_q\ \le\ 16\,q^{\,2d(m-1-d)-d(n-1-d)}=16\,q^{-c\,(m-1-i'+c)} .
\]
As $0\le c(g)\le i'$,
\[
\sum_{g\in S_z}f_\pbmD(g)\le\#S_z\sum_{c=0}^{i'}16\,q^{-c(m-1-i'+c)}\,q^{\,c(m+i')}=16\,\#S_z\sum_{c=0}^{i'}
q^{\,c(2i'+1-c)}\le16\,(i'+1)\,q^{\,i'(i'+1)}\,\#S_z .\qedhere
\]
\end{proof}

\begin{proposition}\label{pbm:Phi}
Let $\pbmD$ be a chain datum of type $(\varrho,\varrho)$ with $m\ge2\varrho^2+4$ and $q\ge8$. Then
\[
1\le\pbmPhi(\pbmD)\le\Bigl(1-\frac4q\Bigr)^{-\varrho}\Bigl(1+\frac{128\,\varrho(\varrho+1)}{q^2}\Bigr).
\]
\end{proposition}
\begin{proof}
The lower bound is the term $u=1$. For $0\le j\le\varrho$ let $\pbmD^{(j)}$ be the chain datum of
type $(j,j)$ obtained from $\pbmD$ by removing the vectors $e_a,e_a^*$ with $a\le\varrho-j$, on the
space $V^{(j)}$ of dimension $r_j=m+2j$, with group $\mathsf G_j=\auGL(V^{(j)})$. For $j=0$ the
group $U$ is trivial and $\pbmPhi(\pbmD^{(0)})=1$. Let $1\le j\le\varrho$, write $f=f_{\pbmD^{(j)}}$ and
let $e_1,e_1^*$ denote the first vectors of the two chains of $\pbmD^{(j)}$. By
Lemma~\ref{pbm:pair}, with its set $\Omega$,
\begin{equation}\label{pbm:eq-rec}
\#\mathsf G_j\cdot\pbmPhi(\pbmD^{(j)})\le\#\Omega\cdot\pbmPhi(\pbmD^{(j-1)})+\sum_{\epsilon_1^*(he_1)=0}f(h)
+\sum_{\epsilon_1^*(h^{-1}e_1)=0}f(h),
\end{equation}
and the last two sums are equal, because $f(h^{-1})=f(h)$. By Lemma~\ref{pbm:pos}, $f\ge0$ is of
positive type. We apply Lemma~\ref{pb:singular} with $\pi$ the identity of $\mathsf G_j$, $l=1$,
$A_0$ the column $e_1$ and $A_1$ the column of the linear form $\epsilon_1^*$, so that
${}^tA_1hA_0=\epsilon_1^*(he_1)$:
\begin{equation}\label{pbm:eq-fourier}
\sum_{\epsilon_1^*(he_1)=0}f(h)\le\frac2q\sum_{h\in\mathsf G_j}f(h)+q^{\,r_j/2}\,Y,\qquad
Y=\max\Bigl(\sum_{he_1=e_1}f(h),\ \sum_{\epsilon_1^*\circ h=\epsilon_1^*}f(h)\Bigr)
\end{equation}
(the stabiliser of a non-zero multiple of $e_1$, or of $\epsilon_1^*$, is that of $e_1$, or of
$\epsilon_1^*$).

\emph{The bound for $Y$.} By Lemma~\ref{pbm:first}, applied $j$ times (to the data of types $(j,j),
(j-1,j),\dots,(1,j)$ on spaces of dimensions $r_j,r_j-1,\dots,r_j-j+1$), and by
Lemma~\ref{pbm:crude} for the datum of type $(0,j)$ on a space of dimension $n=m+j$,
\[
\sum_{he_1=e_1}f(h)=q^{\,2\sum_{t=1}^j(r_j-t)}\sum_{g\in\auGL_n(\auF),\ gz=z}f_{(0,j)}(g)\le q^{\,2\sum_{t=1}^j(
r_j-t)}\cdot16(j+1)q^{\,j(j+1)}\,\frac{\#\auGL_n(\auF)}{q^n-1}.
\]
Since $\#\mathsf G_j=\#\auGL_n(\auF)\prod_{t=0}^{j-1}(q^{r_j-t}-1)q^{r_j-t-1}$, and $\prod_{s\ge2}(1-
q^{-s})^{-1}\le2$, $(q^n-1)^{-1}\le2q^{-n}$,
\[
\frac1{\#\mathsf G_j}\sum_{he_1=e_1}f(h)\le16(j+1)q^{\,j(j+1)}\cdot2q^{-n}\cdot\prod_{t=0}^{j-1}\frac{q^{-1}}{1-
q^{-(r_j-t)}}\le64\,(j+1)\,q^{\,j^2-m-j}.
\]
The second sum in $Y$ is of the same kind for the \emph{dual datum}: let $\auGL(V)$ act on $V^*$ by
$a\mapsto a\circ h^{-1}$, and give $V^*$ the basis $\epsilon_1^*,\dots,\epsilon_j^*$; the dual of
the basis of $X'$; $\epsilon_j,\dots,\epsilon_1$, with the vector $\zeta\in X'^*$ and the linear form
``evaluation at $z$'' on $X'^*$. In this basis $U$ is again the group of a chain datum of type
$(j,j)$, and, since $\epsilon(u^{-1}y)=-\epsilon(uy)$ for the entries just above the diagonal
blocks, $\psi_{\pbmD^{(j)}}$ is the character of this datum for the pair of characters $(\psi_2^{-1},
\psi_1^{-1})$. Its function $f$ is the same function on the same group, and the stabiliser of the
form $\epsilon_1^*$ is the stabiliser of the first vector of its left chain. So the bound holds
for both sums, and, as $r_j/2=m/2+j$ and $m\ge2j^2+4$,
\[
q^{\,r_j/2}\,Y\le64\,(j+1)\,q^{\,j^2-m/2}\cdot\#\mathsf G_j\le64\,(\varrho+1)\,q^{-2}\cdot\#\mathsf G_j .
\]
\emph{Conclusion.} By \eqref{pbm:eq-rec}, \eqref{pbm:eq-fourier} and $\#\Omega\le\#\mathsf G_j$,
\[
\Bigl(1-\frac4q\Bigr)\pbmPhi(\pbmD^{(j)})\le\pbmPhi(\pbmD^{(j-1)})+\frac{128\,(\varrho+1)}{q^2},
\]
and the proposition follows by induction on $j$.
\end{proof}

\begin{lemma}[The character $\psi_U$ at a chosen place]\label{pbm:chain}
Let $w$ be a chosen place, $\auF=\auF_w$, and $V_w$, $\mathsf J$ as in \S\ref{pb:ss-setting}, with
the basis $e_1,\dots,e_\varrho$; the reduction $X'$ of the lattice of $X_0\perp Ez$;
$e^*_\varrho,\dots,e^*_1$. Then $\pbUb_w$ is the group, and the component $\psi_{U,w}$ of $\psi_U$ on
$\pbUb_w$ the character, of the chain datum of type $(\varrho,\varrho)$ on $V_w$ with this basis,
$m=k+1$, the vector $z$, the linear form $\zeta(y)={}^tz\,\mathsf Jy$ on $X'$, and two non-trivial
characters $\psi_1,\psi_2$ of $\auF_w$.
\end{lemma}
\begin{proof}
The group $\pbUb_w$ is the unipotent radical of the parabolic subgroup with blocks $1,\dots,1,k+1,1,
\dots,1$ that stabilises the reduction of the flag $A_1\subset\dots\subset A\subset A^\perp\subset
A_{\varrho-1}^\perp\subset\dots$; in the basis above this is $U_\pbmD$ (\S\ref{pb:ss-setting}). The
vector $z$ lies in the lattice, $z\ne0$ in $X'$, and $\zeta(z)=h(z,z)=-N$ is not zero modulo $w$,
as $w\nmid N$. Let $u\in K_U$ have the components $u_w$ and $u_{\aubar w}=\mathsf J^{-1}\,{}^tu_w^{-1}
\mathsf J$ modulo $w$ and $\aubar w$. By \eqref{pb:eq-psi}, $\psi_U(u)=\prod_v\psi_{E,v}(-\sigma(u)_v/
(M\varpi))$; the factors at the places $v\nmid M$ are $1$, because $\sigma(u)_v$ is integral and
$M\varpi$ generates the different at $v$. At $v\in\{w,\aubar w\}$ the place is unramified over $\Q$,
$M$ is a uniformiser and $\varpi$ a unit, so $x\mapsto\psi_{E,v}(-x/(M\varpi))$ is a non-trivial
character of $\auO_v/v=\auF_w$; call these $\psi_1$ (at $w$) and $\psi_2'$ (at $\aubar w$). The
component of $\sigma(u)$ at $w$ is $\sum_{a<\varrho}\epsilon_a(u_we_{a+1})+\epsilon_\varrho(u_wz)$.
At $\aubar w$, with $\langle y,y'\rangle={}^ty\,\mathsf Jy'$ and $\langle u_{\aubar w}y,y'\rangle=
\langle y,u_w^{-1}y'\rangle$, it is
\[
\sum_{a<\varrho}\langle e_{a+1},u_w^{-1}e_a^*\rangle+\langle z,u_w^{-1}e_\varrho^*\rangle=\sum_{a<\varrho}
\epsilon^*_{a+1}(u_w^{-1}e_a^*)+\zeta(u_w^{-1}e^*_\varrho)=-\Bigl(\sum_{a<\varrho}\epsilon^*_{a+1}(u_we_a^*)+
\zeta(u_we^*_\varrho)\Bigr),
\]
because $u_w^{-1}=1-(u_w-1)+(u_w-1)^2-\dots$ and $(u_w-1)^2$ contributes nothing to these entries.
So $\psi_{U,w}=\psi_\pbmD$ with $\psi_2=\psi_2'^{-1}$.
\end{proof}

\begin{theorem}[Theorem~\ref{pb:est}\,\textup{(E4)}]\label{pbm:E4}
Assume $b\ge3$, $k\ge2\varrho^2+3$, $\auNm w\ge b^{10}$ for the chosen places $w$, and $18\,\varrho f\le b^8$
(these hold for the data of \S\ref{pb:ss-setting}: $\varrho<2\sqrt x$, $k\ge x^5/2$, $f\le x^{12}$, $p_i\ge b^{10}$, and $b\ge1000$ by {\cite[Lemma~13.4(S1)]{P46}}). Then
$\sum_{s\in\auG}|\psi^{\pbnat}(s)|^2\le2$.
\end{theorem}
\begin{proof}
$\auG$, $\pbUb$ and $\psi_U$ are products over the $2f$ chosen places $w$ (above $p_0$ and $p_1$)
of $\auGL_r(\auF_w)$, $\pbUb_w$ and $\psi_{U,w}$; hence $\tilde\psi$ and $\psi^{\pbnat}$ are products
of the corresponding functions of the factors, and by Lemmas~\ref{pbm:pos}(c)
and~\ref{pbm:chain}, $\sum_{s\in\auG}|\psi^{\pbnat}(s)|^2=\prod_w\pbmPhi(\pbmD_w)$ for chain data
$\pbmD_w$ of type $(\varrho,\varrho)$ with $m=k+1$ over fields with $q_w\ge b^{10}$ elements. By
Proposition~\ref{pbm:Phi}, and $-\log(1-4/q)\le8/q$, and $128(\varrho+1)\le128(b^8/18+1)\le b^{10}\le q_w$ (as $b\ge3$),
\[
\log\pbmPhi(\pbmD_w)\le\frac{8\varrho}{q_w}+\frac{128\,\varrho(\varrho+1)}{q_w^2}\le\frac{9\varrho}{b^{10}},
\qquad\sum_w\log\pbmPhi(\pbmD_w)\le\frac{18\,\varrho f}{b^{10}}\le b^{-2}\le\log2 .\qedhere
\]
\end{proof}

\begin{remark}\label{pbm:rem-E4}
(a) The proof does not use $\zeta(z)\ne0$. (b) The same method bounds the sum of $f_\pbmD$ over
the stabiliser of any non-zero vector of the frame. The gain $b^{-6k}$ for each fixed vector in
Theorem~\ref{pbm:inputs}(F) is proved differently ({\cite[Theorem~5.24]{P52}}): the stabiliser of a vector
has index $p_0^r-1$, which makes its volume small, while its mass integral is controlled by its main
part. (c) \textup{(E4)} could also be deduced from a bound for $\int|\int_K\tilde\psi(yxy^{-1})\,dy|^2dx$
over the whole group at the level primes, which is part of the local theory behind \textup{(L2)}; the
argument above is independent of it.
\end{remark}

\subsubsection{The terms over open elements}\label{pbm:ss-cell}

We turn to \textup{(L2)}. Put $n=\varrho+1$. We index the frame by $0\le a\le\varrho$:
$f_a=e_{a+1}$ for $a<\varrho$ and $f_\varrho=z$, and write $f_a^*=e_{a+1}^*$ for $a<\varrho$. Let
$\pbmXp=X_0\perp Ez$, the orthogonal complement of $A\oplus A^*$, where $A^*=\langle e_1^*,\dots,
e_\varrho^*\rangle$, and let $\pbmGf=\auU(\pbmXp)\subset G$ be the group of the elements that fix
the $e_a$ and the $e_a^*$; so $H\subset\pbmGf$ is the stabiliser of $z$. A \emph{head} is a tuple
$p=(p_0,\dots,p_{\varrho-1})$ of non-zero elements of $E$; its \emph{Weyl element} $w_p\in\Gamma$
is defined by
\[
w_pf_a=p_a\,f_a^*,\qquad w_pf_a^*=\aubar{p_a}^{\,-1}f_a\quad(a<\varrho),\qquad w_p=1\text{ on }\pbmXp .
\]
For $\sigma\in\pbmGf(E^+)$ the \emph{last entry} is $t(\sigma)=h(\sigma z,z)$. As $\sigma z=x-
(t/N)z$ with $x\in X_0$ and $h(\sigma z,\sigma z)=-N$,
\begin{equation}\label{pbm:eq-last}
h(x,x)=\frac{t\aubar t-N^2}N,\qquad\text{so}\qquad|\tau(t)|\ge N\ \text{ for all }\tau\in\auPl ,
\end{equation}
with equality for one $\tau$ if and only if $x=0$. For $\gamma=w_p\sigma$ we have $h(\gamma f_l,f_j)
=p_l\delta_{jl}$ for $j,l<\varrho$, $h(\gamma z,z)=t(\sigma)$ and $h(\gamma f_l,z)=h(\gamma z,f_j)=0$,
so
\begin{equation}\label{pbm:eq-Z}
Z(w_p\sigma)=\operatorname{diag}(p_0,\dots,p_{\varrho-1},t(\sigma)),\qquad\delta_l(w_p\sigma)=p_0\cdots
p_{l-1}\ (l\le\varrho),\qquad\pbg(w_p\sigma)=|\delta_\varrho|^2\bigl(t\aubar t-N^2\bigr),
\end{equation}
where $|\delta|^2=\delta\aubar\delta$.

\begin{proposition}[The big cell]\label{pbm:cell}
Let $\gamma$ be a point of $G$ with values in $E^+$ or in a completion $E^+_v$; its matrix has
entries in $F=E$ or $F=E\otimes_{E^+}E^+_v$. Suppose that $\delta_l(\gamma)$ is invertible in $F$ for
$1\le l\le\varrho$. Then there are unique $u,u'\in U$, a
head $p$ and $\sigma\in\pbmGf$ with $\gamma=u\,w_p\,\sigma\,u'$; moreover $p_a=\delta_{a+1}(\gamma)/
\delta_a(\gamma)$ and $t(\sigma)=\delta_n(\gamma)/\delta_\varrho(\gamma)$. In particular $U\cap\gamma
U\gamma^{-1}=1$. The maps $\gamma\mapsto u,u',p,\sigma$ are continuous on the open set where the
$\delta_l$, $l\le\varrho$, are invertible. If $v$ is a finite place, $\gamma$ stabilises the lattice
of $Y_0$ at $v$ and the $\delta_l(\gamma)$, $l\le\varrho$, are units at $v$, then $u$, $u'$,
$\sigma$ stabilise the lattice at $v$.
\end{proposition}
\begin{proof}
The leading minors $\delta_l$ are invariant under $\gamma\mapsto u\gamma u'$: $u'f_l$ is $f_l$ plus
a combination of the $f_{l'}$ with $l'<l$, and the same holds for $u^{-1}f_j$, so $Z(u\gamma u')=
\mathsf L\,Z(\gamma)\,\mathsf R$ with $\mathsf L$ lower and $\mathsf R$ upper unitriangular. Let
$U_Q\subseteq U$ be the subgroup of the elements that act trivially on $A$, on $A^\perp/A$ and on
$Y_0/A^\perp$, and $N_A\subseteq U$ the subgroup of the elements that preserve $A^*$ and $\pbmXp$;
then $U=U_Q\rtimes N_A$, and $N_A$ acts on $A$ by all the unitriangular matrices of the basis
$(e_a)$. The group $U_Q$ acts simply transitively on the totally isotropic subspaces $A'$ of
dimension $\varrho$ with $A'\cap A^\perp=0$: such an $A'$ has a unique basis $(e_a')$ with $h(e_b,
e_a')=\delta_{ab}$, and $e_a'-e_a^*\in A^\perp$; the linear map $u$ with $ue_a=e_a$, $ue_a^*=e_a'$
and $ux=x-\sum_a\aubar{h(e_a',x)}\,e_a$ for $x\in\pbmXp$ is unitary (the correction makes $ux$
orthogonal to the $e_a'$, and $A$ is orthogonal to $A^\perp$), lies in $U_Q$ and maps $A^*$ onto
$A'$; an element of $U_Q$ with this property must map $e_a^*$ to $e_a'$ and $x$ to the element of
$x+A$ orthogonal to $A'$, so it is $u$.

\emph{Existence.} As $\delta_\varrho(\gamma)\ne0$, the pairing of $A$ with $\gamma A$ is perfect, so
$\gamma A\cap A^\perp=0$ and there is $u_1\in U_Q$ with $u_1A^*=\gamma A$. Then $\gamma_1=u_1^{-1}
\gamma$ maps $A$ onto $A^*$; the pairing of $A$ with $\gamma_1^{-1}A$ is the adjoint of that of
$\gamma_1A=A^*$ with $A$, hence perfect, and there is $u_2\in U_Q$ with $u_2A^*=\gamma_1^{-1}A$. Now
$\gamma_2=\gamma_1u_2$ maps $A$ onto $A^*$ and $A^*$ onto $A$, hence preserves $\pbmXp$. The block
$(h(\gamma_2f_l,f_j))_{j,l<\varrho}$ has the leading minors $\delta_l(\gamma)\ne0$, so it is
$\mathsf L\,\mathsf D\,\mathsf R$ with unique unitriangular $\mathsf L$, $\mathsf R$ and diagonal
$\mathsf D=\operatorname{diag}(\delta_{a+1}/\delta_a)$. Multiplying $\gamma_2$ on both sides by the
elements of $N_A$ that induce $\mathsf L^{-1}$ and $\mathsf R^{-1}$ we get $\gamma_3$ with $\gamma_3
f_a=p_af_a^*$; by unitarity $\gamma_3f_a^*=\aubar{p_a}^{-1}f_a$, so $\gamma_3=w_p\sigma$ with
$\sigma=\gamma_3|_{\pbmXp}\in\pbmGf$. The value of $t(\sigma)$ follows from \eqref{pbm:eq-Z}.

\emph{Uniqueness.} Let $u\,w_p\sigma\,u'=w_q\sigma'$. The minors give $p=q$. Applying both sides to
$A$ gives $uA^*=A^*$, so $u\in N_A$, and in the same way $u'\in N_A$; then $\mathsf L\mathsf D\mathsf
R=\mathsf D$ forces $u=u'=1$ and $\sigma=\sigma'$. For $u\in U\cap\gamma U\gamma^{-1}$ apply this to
$u\gamma=\gamma u'$.

\emph{Continuity and integrality.} Every step is the solution of a linear system whose determinant
is a product of the $\delta_l$; this gives the continuity, and, over the ring of integers at $v$
when the $\delta_l$ are units and $\gamma$ preserves the lattice (which is the sum of its
intersections with $A\oplus A^*$ and $\pbmXp$, with the basis $(e_a,e_a^*)$ of the first), the same
construction gives integral $u_1,u_2,\mathsf L,\mathsf R$, hence integral $u$, $u'$, $\sigma$.
\end{proof}

\paragraph*{Orbit sums}
We use the cut Bessel functions $W^\flat$ of Definition~\ref{pb:bessel}, and assume in addition:
\begin{enumerate}
\item[(B5$'$)] for each $\tau$ there is a compact set $C_\tau\subset G_\tau$ with $\operatorname{supp}
W_\tau^\flat\subseteq U_\tau C_\tau$.
\end{enumerate}
(This holds for the cutoff of Theorem~\ref{pbe:bessel}: there $W^\flat_\tau$ vanishes unless
$\|sy\|_\tau<b^{6n+2}$ for some $s\in R_\tau=U_\tau\rtimes H_\tau$, and $H_\tau$ is compact.) For
$\gamma\in\Gamma$ in the big cell, that is, with $\delta_l(\gamma)\ne0$ for $1\le l\le\varrho$, the
\emph{orbital integral} is
\[
O(\gamma)=\int_{G_\infty}W^\flat(y)\,W^\flat(\gamma^{-1}y)^*\,dy .
\]
The integrand is continuous with compact support: if it does not vanish at $y$, then for each
$\tau$, $y_\tau=u_1c_1$ and $\gamma^{-1}y_\tau=u_2c_2$ with $u_i\in U_\tau$, $c_i\in C_\tau$, so that
$u_1^{-1}\gamma u_2=c_1c_2^{-1}$ lies in the compact set $C_\tau C_\tau^{-1}$; by
Proposition~\ref{pbm:cell} the set $U_\tau\gamma U_\tau$ is closed in $G_\tau$ (it is defined by
the equations $\delta_l=\delta_l(\gamma)$, $l\le\varrho$, and $\sigma(\cdot)=\sigma(\gamma)$, where
$\sigma(\cdot)$ is continuous on the big cell) and $(u_1,u_2)\mapsto u_1^{-1}\gamma u_2$ is a
homeomorphism of $U_\tau^2$ onto it; hence $u_1$ stays in a compact set, and so does $y_\tau$. By
Proposition~\ref{pb:unfold} and $\|W^\flat\|\le\|W\|$ the amplitude sums \eqref{pb:eq-amp} for
$W^\flat$ converge absolutely. For classes
$c,c'$, $s\in\auG$ and $\gamma_0$ in the big cell the set $U(E^+)\gamma_0U(E^+)$ is stable; its amplitude sum is the \emph{orbit sum}
$\pbmOrb_{c,s,c'}(\gamma_0)=A^\flat_{U(E^+)\gamma_0U(E^+)}(c,s,c')$.

\begin{proposition}[Orbit sums are Kloosterman sums]\label{pbm:orbit}
Let $\gamma_0$ be in the big cell and
\[
\pbmKl_{c,s,c'}(\gamma_0)=\sum_{\substack{u\in U(E^+)/U_c,\ u'\in U(E^+)/U_{c'}\\ u^{-1}\gamma_0u'\in\Gamma_{c,s,
c'}}}\psi_f(u)\,\aubar{\psi_f(u')} .
\]
If $O(\gamma_0)\ne0$ the sum is finite and $\pbmOrb_{c,s,c'}(\gamma_0)=\pbmKl_{c,s,c'}(\gamma_0)\cdot
O(\gamma_0)$; if $O(\gamma_0)=0$ then $\pbmOrb_{c,s,c'}(\gamma_0)=0$.
\end{proposition}
\begin{proof}
By Proposition~\ref{pbm:cell}, $(u,u')\mapsto u^{-1}\gamma_0u'$ is a bijection of $U(E^+)^2$ onto
the orbit, the stabilisers $U_c\cap\gamma U_{c'}\gamma^{-1}$ are trivial, and the double cosets
$U_c\gamma U_{c'}$ in the orbit correspond to the pairs $(uU_c,u'U_{c'})$. For $\gamma=u^{-1}\gamma_0
u'$ the substitution $y\mapsto u^{-1}y$ and (B1) give
\[
O(\gamma)=\int_{G_\infty}W^\flat(u^{-1}y)\,W^\flat(u'^{-1}\gamma_0^{-1}y)^*\,dy=\aubar{\psi_\infty(u)}\,\psi_\infty(
u')\,O(\gamma_0)=\psi_f(u)\,\aubar{\psi_f(u')}\,O(\gamma_0),
\]
since $\psi_\infty\psi_f=1$ on $U(E^+)$. All terms of the amplitude sum have the norm $\|O(\gamma_0)
\|$, and the sum of the norms converges; so the number of terms is finite if $O(\gamma_0)\ne0$.
\end{proof}

By strong approximation for $U$, $U(E^+)/U_c=U(\auA_f)/(U(\auA_f)\cap h_cJh_c^{-1})$; so
$\pbmKl_{c,s,c'}(\gamma_0)$ is a sum over the pairs $(u,u')$ of finite adelic points with $h_c^{-1}
u^{-1}\gamma_0u'h_{c'}\in K$ of reduction $s$: a Kloosterman sum for the group $U$, the character
$\psi_f$ and the level $J$.

\paragraph*{Heads, links and class sums}
By (B5$'$) and the invariance of the $\delta_l$ there is a constant $C_0$ such that
\begin{equation}\label{pbm:eq-bdd}
O(\gamma)\ne0\ \Longrightarrow\ |\tau(\delta_l(\gamma))|\le C_0\quad(\tau\in\auPl,\ 1\le l\le n):
\end{equation}
if $O(\gamma)\ne0$ there are $y_\tau=u_1c_1$ and $\gamma^{-1}y_\tau=u_2c_2$ with $u_i\in U_\tau$,
$c_i\in C_\tau$, so $\tau(\delta_l(\gamma))=\delta_l(c_1c_2^{-1})$, which is bounded on the compact
set $C_\tau C_\tau^{-1}$. Let $\mathsf B$ be the finite set of the non-zero integers of $E$ all of
whose conjugates have absolute value at most $C_0$, and
\[
\pbmT=\{\delta/\delta':\delta,\delta'\in\mathsf B\}\subset E^\times ,
\]
the \emph{set of entries}; it is finite and depends only on the data and on the $C_\tau$.

\begin{proposition}[Class sums over open elements]\label{pbm:classsum}
Let $s\in\auG$ be open, $c,c'$ classes and $\Lambda$ a link.
\begin{enumerate}
\item[(a)] Every $\gamma\in\Gamma_{c,s,c'}$ lies in the big cell, and $\delta_l(\gamma)\in\auO_E$ for
$l\le n$.
\item[(b)] For $\gamma=u\,w_p\,\sigma\,u'$ as in Proposition~\ref{pbm:cell} the link of $\gamma$ is
the graph $\{(x,\sigma^{-1}x):x\in\pbmXp\}$ of $\sigma^{-1}$, and $\pbg(\gamma)\ne0$ if and only if
$|\tau(t(\sigma))|>N$ for all $\tau$.
\item[(c)] If $A^\flat_{T_\Lambda}(c,s,c')\ne0$, there is a unique $\sigma=\sigma_\Lambda\in\pbmGf(E^+)$
whose inverse has the graph $\Lambda$; its last entry $t=t(\sigma)$ lies in $\pbmT$ and
satisfies $|\tau(t)|>N$ for all $\tau$; and
\[
A^\flat_{T_\Lambda}(c,s,c')=\sum_{p\in\pbmT^\varrho}\pbmOrb_{c,s,c'}(w_p\sigma).
\]
Moreover, if $v$ is a finite place at which all elements of $\mathsf B$ are units, then
$h_c^{-1}\sigma h_{c'}$ stabilises the lattice at $v$.
\end{enumerate}
\end{proposition}
\begin{proof}
(a) Let $\kappa=h_c^{-1}\gamma h_{c'}\in K$. As $h_c,h_{c'}\in H(\auA_f)$ are unitary and fix the
frame, $h(\gamma f_l,f_j)=h(\kappa f_l,f_j)$ is integral at all finite places (the lattice is
unimodular and contains the frame); so $Z(\gamma)$ has entries in $\auO_E$. Its reduction modulo a
chosen place $w$ is the matrix $Z_w(s_w)$ of \S\ref{pb:ss-setting} (up to transposition, which does
not change leading minors). As $s$ is open, $\delta_l(s_w)\ne0$ for $1\le l\le\varrho$ at some $w$,
hence $\delta_l(\gamma)\ne0$. (b) The link is invariant under $U(E^+)$ on both sides, so we may take
$\gamma=w_p\sigma$. Then $\gamma A=A^*$, the vectors orthogonal to $A$ and $\gamma A$ are those of
$\pbmXp$, and $\gamma^{-1}x=\sigma^{-1}x$ for $x\in\pbmXp$. The second claim is \eqref{pbm:eq-Z},
\eqref{pbm:eq-last} and the invariance of $\pbg$: $t\aubar t-N^2\in E^+$ is totally non-negative.
(c) By (a), (b) and Proposition~\ref{pbm:cell}, $T_\Lambda\cap\Gamma_{c,s,c'}$ is the union, over
the heads $p$, of the sets $U(E^+)w_p\sigma U(E^+)\cap\Gamma_{c,s,c'}$ for the unique $\sigma$ with
graph of $\sigma^{-1}$ equal to $\Lambda$ (if there is none, or if $|\tau(t(\sigma))|=N$, the set is
empty). So the amplitude sum is the sum of the orbit sums, each orbit being stable. If
$\pbmOrb_{c,s,c'}(w_p\sigma)\ne0$, the orbit contains a $\gamma\in\Gamma_{c,s,c'}$ with $O(\gamma)\ne
0$; by (a) and \eqref{pbm:eq-bdd} its minors $\delta_l(\gamma)$, $1\le l\le n$, lie in $\mathsf B$
(note $\delta_n=\delta_\varrho\,t\ne0$), so $p_a=\delta_{a+1}/\delta_a\in\pbmT$ and $t=\delta_n/
\delta_\varrho\in\pbmT$. Finally let $v$ be as in the last claim and $\gamma$ as before. Then
$\kappa_v=(h_c^{-1}\gamma h_{c'})_v$ stabilises the lattice and has unit minors, so $\kappa_v=u_v
w_p\sigma_vu_v'$ with $\sigma_v$ integral, by Proposition~\ref{pbm:cell}. As $h_c$, $h_{c'}$ commute
with $w_p$ and normalise $U$, $\gamma=(h_cu_vh_c^{-1})\,w_p\,(h_c\sigma_vh_{c'}^{-1})\,(h_{c'}u_v'h_{c'}^{
-1})$, and the uniqueness in Proposition~\ref{pbm:cell} gives $h_c^{-1}\sigma h_{c'}=\sigma_v$.
\end{proof}

\subsubsection{Sieve systems and one-sided Poisson summation}\label{pbm:ss-sieve}

\paragraph*{Sieve systems}
A \emph{place} is a quadruple $\mathfrak P=(C,dc,\chi,\iota)$: a locally compact ring $C$ with a
Haar measure $dc$ such that $C^\times$ is open with complement of measure $0$, a continuous unitary
character $\chi$ of $(C,+)$, and a map $\iota\colon C^\times\to C$. A \emph{head of length $i$} at
$\mathfrak P$ is an element of $(C^\times)^i$; for heads $p,q$ of length $i$ and $c\in C^\times$,
$(p,c)$ is the head of length $i+1$.

\begin{definition}[Sieve system]\label{pbm:system}
Let $\mathfrak P$ be a place, $\pbmE\ge0$, $\beta\colon C^\times\to[0,\infty)$ and $\Lambda_i\colon
(C^\times)^i\to[0,\infty)$ for $0\le i\le\varrho$. A \emph{sieve system} at $\mathfrak P$ with mass
$\pbmE$, weights $\Lambda$ and decay $\beta$ is a family of functions $J_i\colon(C^\times)^i\times
(C^\times)^i\to\C$, $0\le i\le n=\varrho+1$, such that:
\begin{enumerate}
\item[(S1)] for $0\le i\le\varrho$, heads $p,q$ of length $i$ and $\xi\in C$, the function $c\mapsto
\chi(\xi c)J_{i+1}((p,c),(q,c))$ is integrable on $C$; $\int_CJ_{i+1}((p,c),(q,c))\,dc=J_i(p,q)$; and
$\bigl|\int_C\chi(\xi c)J_{i+1}((p,c),(q,c))\,dc\bigr|\le\beta(\xi)\Lambda_i(p)\Lambda_i(q)$ for
$\xi\in C^\times$;
\item[(S2)] for heads $p^{(1)},\dots,p^{(\kappa)}$ of length $\varrho$ and $t\in C^\times$ the matrix
$\bigl(J_n((p^{(a)},t),(p^{(b)},t))\bigr)_{a,b}$ is hermitian and positive semidefinite;
\item[(S3)] $J_n$ is continuous;
\item[(S4)] $|J_0|\le\pbmE$;
\item[(S5)] for $1\le i\le\varrho$ and heads $p,q$ of length $i$ such that $\xi=\iota(p_{i-1})-\iota(
q_{i-1})$ is a unit, $|J_i(p,q)|\le\beta(\xi)\Lambda_i(p)\Lambda_i(q)$.
\end{enumerate}
\end{definition}

\begin{lemma}\label{pbm:sysprop}
\begin{enumerate}
\item[(a)] \emph{(Products.)} Let $(J^v_i)$ be sieve systems at places $\mathfrak P_v=(C_v,dc_v,\chi_v,
\iota_v)$, $v$ in a finite set, with masses $\pbmE_v$, weights $\Lambda^v$ and decays $\beta_v$. Then
$J_i(p,q)=\prod_vJ^v_i(p_v,q_v)$ is a sieve system at the product place $(\prod C_v,\prod dc_v,\prod
\chi_v,\prod\iota_v)$ with mass $\prod\pbmE_v$, weights $\prod_v\Lambda^v_i(p_v)$ and decay $\prod_v
\beta_v(\xi_v)$.
\item[(b)] \emph{(Positivity and semicontinuity.)} Let $(J_i)$ be a sieve system, $0\le i\le\varrho$
and $\mathcal P$ a finite set of heads of length $i$. Then $D(x)=\sum_{p,q\in\mathcal P}J_{i+1}((p,x),
(q,x))$ is real, non-negative and lower semicontinuous on $C^\times$.
\end{enumerate}
\end{lemma}
\begin{proof}
(a) The units of the product are the families of units. (S1) follows from Fubini's theorem: the
function is a product of integrable functions of the separate variables. (S2) is Schur's theorem
that the entrywise product of positive semidefinite hermitian matrices is positive semidefinite.
(S3)--(S5) are clear. (b) For $0\le j\le\varrho-i$ and $c\in(C^\times)^j$ let $D_j(x,c)=\sum_{p,q\in
\mathcal P}J_{i+1+j}((p,x,c),(q,x,c))$. For $j=\varrho-i$ this is the sum of all entries of a matrix
as in (S2), so it is $\ge0$, and it is continuous by (S3). By (S1), $D_j(x,c)=\int_CD_{j+1}(x,c,c')
\,dc'$; so, by descending induction on $j$ and Fatou's lemma, $D_j\ge0$ is lower semicontinuous in
$(x,c)$. The case $j=0$ is the claim.
\end{proof}

\paragraph*{One-sided Poisson summation}

\begin{lemma}\label{pbm:poisson}
Let $\mathbb A$ be a locally compact abelian group with a Haar measure, $L\subset\mathbb A$ a
discrete subgroup with compact quotient of volume $\mathsf c_L$, and $L^\perp$ the group of the
characters of $\mathbb A$ trivial on $L$. Let $F\ge0$ be integrable on $\mathbb A$ and lower
semicontinuous at every point of a finite set $V\subseteq L$. Then
\[
\sum_{t\in V}F(t)\le\frac1{\mathsf c_L}\sum_{\eta\in L^\perp}\Bigl|\int_{\mathbb A}F(x)\,\eta(x)\,dx\Bigr| .
\]
\end{lemma}
\begin{proof}
Let $\mathsf S$ be the right side; we may assume $\mathsf S<\infty$. The function $PF(x)=\sum_{t\in L}
F(x+t)$ is integrable on the compact group $\mathbb A/L$, whose dual is $L^\perp$, and its Fourier
coefficients, for the measure of mass $1$, are $\mathsf c_L^{-1}\int_{\mathbb A}F\aubar\eta$. They are
absolutely summable, so $PF$ agrees almost everywhere with a continuous function of absolute value
at most $\mathsf S$ (two integrable functions with the same Fourier coefficients agree almost
everywhere). Let $\epsilon>0$. By lower semicontinuity there is a neighbourhood $\mathsf N$ of $0$
with $F(t+x)>F(t)-\epsilon$ for $t\in V$, $x\in\mathsf N$. As $\mathsf N$ has positive measure,
there is $x\in\mathsf N$ with $\mathsf S\ge PF(x)\ge\sum_{t\in V}F(t+x)\ge\sum_{t\in V}F(t)-\epsilon\,
\#V$.
\end{proof}

\paragraph*{The global place}
Let $T$ be a finite set of finite places of $E$, stable under conjugation, containing the places
above $M$. Let $E_T=\prod_{w\in T}E_w$, with integers $\auO_T$ and the Haar measure for which
$\auO_T$ has volume $1$; $|\cdot|_T=\prod_{w\in T}|\cdot|_w$ is the module, and for $\xi\in E_T^\times$
\[
\pbmn(\xi)=|\xi|_T^{-1},\qquad\pbmden(\xi)=[\auO_T:\auO_T\cap\xi^{-1}\auO_T]=\prod_{w\in T}\max(1,|\xi|_w).
\]
Let $\pbmRT\subset E$ be the ring of the elements integral outside $T$, and $\psi_T(x)=\psi_{E,f}(-x/
\varpi)$ for $x\in E_T$, where $\psi_{E,f}$ is the finite part of $\psi_E$. Put
\[
\pbmAT=\C^{\auPl}\times E_T,\qquad\chi_T(x)=\prod_{\tau\in\auPl}e^{x_\tau-\aubar{x_\tau}}\cdot\psi_T(x_T),
\]
with Lebesgue measure on each $\C$. Recall $\kappa_\tau=M\tau(\varpi)/(2\pi i)\in\R$ and put
\begin{equation}\label{pbm:eq-scales}
s_{\tau,a}=\frac{\kappa_\tau^{\,2(\varrho-a)}}N\quad(0\le a\le\varrho),\qquad[t]_a=\bigl((s_{\tau,a}\tau(
t))_\tau,\ t\bigr)\in\pbmAT\quad(t\in E),
\end{equation}
the element $t$ \emph{read as an entry at position $a$} ($s_{\tau,a}$ is the square of the scale of
$f_a$ in the normalised frame of \S\ref{pbe:ss-real}; thus this $s_{\tau,a}$ is the square of the $s_{\tau,a}$ of \eqref{pbe:eq-scale} and of Lemma~\ref{pb:covol}, which is why $\Lambda_*=\Lambda$ below). For a tuple $p=(p_0,\dots,p_{i-1})$ of
non-zero elements of $E$ the head $[p]=([p_0]_0,\dots,[p_{i-1}]_{i-1})$ has components $[p]_\tau\in
(\C^\times)^i$ and $[p]_T=p\in(E_T^\times)^i$. Let
\[
\pbmvol=(2/\aurd)^f,\qquad c_a=\frac{\pbmvol}{\prod_\tau s_{\tau,a}^{\,2}}\quad(0\le a\le\varrho).
\]

\begin{lemma}[Poisson summation in one entry]\label{pbm:poissonT}
Let $F\ge0$ be integrable on $\pbmAT$ and lower semicontinuous at the points $[t]_a$, $t\in V$, for
a finite set $V\subset\pbmRT$. Then
\[
\sum_{t\in V}F([t]_a)\le c_a\sum_{y\in\pbmRT}\Bigl|\int_{\pbmAT}F(x)\,\chi_T(\xi_a(y)x)\,dx\Bigr|,\qquad
\xi_a(y)=\Bigl(\Bigl(\frac{2\pi i}{\tau(\varpi)\,s_{\tau,a}}\,\tau(y)\Bigr)_\tau,\ y\Bigr).
\]
\end{lemma}
\begin{proof}
Embed $\pbmRT$ in $\pbmAT$ by $t\mapsto((\tau(t))_\tau,t)$. It is discrete with compact quotient:
$\pbmAT\times\prod_{v\notin T}\auO_v$ is an open subgroup of $\auA_E$ that meets $E$ in $\pbmRT$ and,
by strong approximation for the additive group, generates $\auA_E$ together with $E$. So $\pbmAT/
\pbmRT$ is the quotient of $\auA_E/E$ by the image of $\prod_{v\notin T}\auO_v$, and its characters
are the characters $x\mapsto\psi_E(y'x)$, $y'\in E$, of $\auA_E/E$ that are trivial on $\prod_{v
\notin T}\auO_v$, that is, those with $y'\in\varpi^{-1}\pbmRT$ (the different is generated by
$\varpi$). As the components of $\psi_E$ at the infinite places are $x\mapsto e^{-2\pi i(x+\aubar x)}$
(convention (Cv2) of \S\ref{pb:ss-setting}), these are, with $y=-\varpi y'\in\pbmRT$, the characters $x\mapsto\prod_\tau e^{\omega_\tau-\aubar{
\omega_\tau}}\cdot\psi_T(yx_T)$ with $\omega_\tau=2\pi i\,\tau(y/\varpi)\,x_\tau$, up to the sign of
$y$. The covolume of $\auO_E$ in $\C^{\auPl}$ is
$2^{-f}|\operatorname{disc}E|^{1/2}=\pbmvol^{-1}$, and $\pbmRT+\auO_T=E_T$, $\pbmRT\cap\auO_T=\auO_E$; so
$\mathsf c_L=\pbmvol^{-1}$. Apply Lemma~\ref{pbm:poisson} to $F_a(x)=F((s_{\tau,a}x_\tau)_\tau,x_T)$;
the substitution $x_\tau\mapsto x_\tau/s_{\tau,a}$ has the Jacobian $\prod_\tau s_{\tau,a}^{-2}$ and
turns the character of $y$ into $\chi_T(\xi_a(y)\,\cdot\,)$.
\end{proof}

\paragraph*{The sieve}
We consider sieve systems at the place $(\pbmAT,dx,\chi_T,\iota)$ with
\begin{equation}\label{pbm:eq-iota}
\iota\bigl((x_\tau)_\tau,x_T\bigr)=\bigl((x_\tau^{-1})_\tau,\ M^{-1}x_T^{-1}\bigr).
\end{equation}

\begin{theorem}[The sieve, abstract form]\label{pbm:abstract}
Let $(J_i)$ be a sieve system at $(\pbmAT,dx,\chi_T,\iota)$ with mass $\pbmE$, weights $\Lambda$ and
decay $\beta$. Let $\pbmT\subset E^\times$ be a finite set with $t,t^{-1}\in\pbmRT$ for $t\in\pbmT$,
and $\epsilon,\Gamma\ge0$ such that:
\begin{enumerate}
\item[(D)] for every $a\le\varrho$, every family $(c_\tau)$ of complex numbers with $|c_\tau|^{-1}\le
|\kappa_\tau|/M$ and every finite set $X\subset\pbmRT\smallsetminus\{0\}$,
$\sum_{y\in X}\beta\bigl((c_\tau\tau(y)/s_{\tau,a})_\tau,\,y\bigr)\le\epsilon$;
\item[(W)] $\bigl(\sum_{p\in\pbmT^i}\Lambda_i([p])\bigr)^2\le\Gamma\prod_{a<i}c_a$ for $0\le i\le
\varrho$.
\end{enumerate}
Assume $|\kappa_\tau|\ge1$ for all $\tau$. Then for every subset $\mathsf U\subseteq\pbmT$
\[
\sum_{t\in\mathsf U}\ \sum_{p,q\in\pbmT^\varrho}\Rel J_n\bigl([p,t],[q,t]\bigr)\le\prod_{a=0}^{\varrho}c_a
\cdot\bigl(\pbmE+2n\,\Gamma\,\epsilon\bigr).
\]
\end{theorem}
\begin{proof}
For $0\le i\le\varrho$ and $x\in\pbmAT^\times$ let $D_i(x)=\sum_{p,q\in\pbmT^i}J_{i+1}(([p],x),([q],x))$,
and put $\Sigma_0=\Rel J_0$, $\Sigma_{i+1}=\sum_{t\in\pbmT}D_i([t]_i)$ and $W_i=\sum_{p\in\pbmT^i}\Lambda_i
([p])$. By Lemma~\ref{pbm:sysprop}(b), $D_i\ge0$ is lower semicontinuous at the points $[t]_i$, so
the left side of the theorem is at most $\Sigma_n$. We show
\begin{equation}\label{pbm:eq-step}
\Sigma_{i+1}\le c_i\bigl(\Sigma_i+2\epsilon W_i^2\bigr)\qquad(0\le i\le\varrho);
\end{equation}
with (W) and (S4) this gives $\Sigma_i\le\prod_{a<i}c_a\cdot(\pbmE+2i\Gamma\epsilon)$ by induction.

\emph{Poisson summation.} By (S1), $D_i$ is integrable, $\int D_i=\sum_{p,q\in\pbmT^i}J_i([p],[q])$,
and $|\int D_i(x)\chi_T(\xi x)dx|\le\beta(\xi)W_i^2$ for units $\xi$. For $y\in\pbmRT$, $y\ne0$, the
frequency $\xi_i(y)$ of Lemma~\ref{pbm:poissonT} is a unit and has the form in (D) with $c_\tau=2\pi
i/\tau(\varpi)$, for which $|c_\tau|^{-1}=|\kappa_\tau|/M$. By Lemma~\ref{pbm:poissonT} (the sums
over finite sets of frequencies being bounded by (D)),
\[
\Sigma_{i+1}\le c_i\Bigl(\sum_{p,q\in\pbmT^i}\Rel J_i([p],[q])+\epsilon W_i^2\Bigr).
\]
\emph{Unequal last entries.} For $i=0$ the sum is $\Sigma_0$. Let $i\ge1$. The pairs $(p,q)$ with
$p_{i-1}=q_{i-1}$ contribute $\Sigma_i$, by definition. For a pair with $p_{i-1}=t\ne\tilde t=q_{i-1}$ put
$y=M^{-1}(t^{-1}-\tilde t^{-1})\in\pbmRT\smallsetminus\{0\}$; then $\iota([t]_{i-1})-\iota([\tilde t]_{i-1})=((M\tau(
y)/s_{\tau,i-1})_\tau,y)$ is a unit of the form in (D) with $c_\tau=M$, and $M^{-1}\le|\kappa_\tau|/
M$. By (S5) and (D) with $X=\{y\}$, $\Rel J_i([p],[q])\le\epsilon\Lambda_i([p])\Lambda_i([q])$. The
sum over these pairs is at most $\epsilon W_i^2$, which proves \eqref{pbm:eq-step}.
\end{proof}

\paragraph*{The explicit weights and decay}
Let $\ell_K(r)=r^K\min(1,r)^6$, let $\pbmTh>0$ be a constant and
\begin{equation}\label{pbm:eq-beta}
\begin{gathered}
\beta^\infty(\xi)=\pbmTh\,b^{2000n^9}\,\ell_{k-n^7}(|\xi|^{-1})\quad(\xi\in\C^\times),\qquad
\beta^T(\xi)=\frac{\pbmn(\xi)^{(k-n^7)/2}}{\pbmden(\xi)^3}\,b^{\,n^9f}\quad(\xi\in E_T^\times),\\
\Lambda^\infty_i(p)=\prod_{l=1}^i\bigl[\,|\delta_l|\le b^{60n^2}\bigr]\max(1,|\delta_l|^{-1})^{n^6},\qquad
\Lambda^T_i(p)=\prod_{l=1}^i\bigl[\delta_l\in\auO_T\bigr]\,\pbmn(\delta_l)^{n^6/2},
\end{gathered}
\end{equation}
where $\delta_l=p_0\cdots p_{l-1}$ and $[\,\cdot\,]$ is $1$ or $0$ according as the condition holds
or not. Let
\[
\pbmdec=\prod_{\tau\in\auPl}\pbmTh\Bigl(\frac{|\kappa_\tau|^{2\varrho+1}}{MN}\Bigr)^{k-n^7}.
\]

\begin{lemma}[Sums over frequencies]\label{pbm:freq}
For $y\in E^\times$ let $\mathfrak b_y=\{a\in\auO_E:ay\in\auO_E\}$. For every finite $X\subset E^\times$
\[
\sum_{y\in X}\ \prod_{\tau\in\auPl}\max(1,|\tau(y)|)^{-6}\cdot\auNm(\mathfrak b_y)^{-3}\le162^f .
\]
\end{lemma}
\begin{proof}
The $y$ with $\mathfrak b_y=\mathfrak b$ lie in $\mathfrak b^{-1}$, which is the union of $\auNm
\mathfrak b$ classes modulo $\auO_E$. Cut each $\C$ into the half-open squares $Q_{\mathsf m}$ of
side $\frac12$ with corners in $\frac12\Z[i]$. Two elements of one class that lie in the same
product of squares differ by an integer all of whose conjugates have absolute value $<1$, so they
are equal. Hence the sum over one class is at most $C_1^f$ with $C_1=\sum_{\mathsf m}\sup_{Q_{\mathsf
m}}\max(1,|z|)^{-6}\le1+\sum_{j\ge1}8j\max(1,\frac{j-1}2)^{-6}<54$. So the left side is at most
$C_1^f\sum_{\mathfrak b}\auNm(\mathfrak b)^{-2}\le(54\,\zeta(2)^2)^f$.
\end{proof}

\begin{theorem}[The sieve]\label{pbm:sieve}
Assume $n\ge3$, $b\ge1000$, $2n\le b$, $1\le|\kappa_\tau|\le b^{23}$ for all $\tau$, $N\le b^{60n}$
and $\aurd\le b$. Let $T$ and $\pbmT$ be as in Theorem~\ref{pbm:abstract}, with $M^{-1}\in\pbmRT$. For
each $\tau\in\auPl$ let $(J^\tau_i)$ be a sieve system at $(\C,\text{Lebesgue},e^{c-\aubar c},w\mapsto
w^{-1})$ with mass $\pbmE_\tau$, weights $\Lambda^\infty$ and decay $\beta^\infty$, and let $(J^T_i)$
be a sieve system at $(E_T,dx,\psi_T,w\mapsto M^{-1}w^{-1})$ with mass $\pbmE_T$, weights $\Lambda^T$
and decay $\beta^T$. Then for every $\mathsf U\subseteq\pbmT$
\[
\sum_{t\in\mathsf U}\ \sum_{p,q\in\pbmT^\varrho}\Rel\Bigl(\prod_\tau J^\tau_n\bigl([p,t]_\tau,[q,t]_\tau
\bigr)\cdot J^T_n\bigl((p,t),(q,t)\bigr)\Bigr)\le\prod_{a=0}^{\varrho}c_a\cdot\Bigl(\prod_\tau\pbmE_\tau
\cdot\pbmE_T+b^{2403\,n^9f}\,\pbmdec\Bigr).
\]
\end{theorem}
\begin{proof}
By Lemma~\ref{pbm:sysprop}(a) the products $J_i=\prod_\tau J^\tau_i\cdot J^T_i$ form a sieve system
at $(\pbmAT,dx,\chi_T,\iota)$ with mass $\prod\pbmE_\tau\cdot\pbmE_T$, weights $\Lambda_i(p)=\prod_\tau
\Lambda^\infty_i(p_\tau)\Lambda^T_i(p_T)$ and decay $\beta(\xi)=\prod_\tau\beta^\infty(\xi_\tau)\beta^T(
\xi_T)$. We verify (D) with $\epsilon=\pbmdec\,b^{2002n^9f}$ and (W) with $\Gamma=b^{400n^9f}$;
Theorem~\ref{pbm:abstract} and $2n\le b^{n^9f}$ then give the claim. We use the product formula in
the form
\begin{equation}\label{pbm:eq-pf}
\pbmn(y)^{1/2}\le\prod_\tau|\tau(y)|\qquad(y\in\pbmRT,\ y\ne0),
\end{equation}
which holds because $|y|_w\le1$ for $w\notin T$.

\emph{(D).} Let $K=k-n^7$, $x_\tau=|\kappa_\tau|^{2\varrho+1}/(MN)$ and $B=b^{23(2\varrho+1)}\ge\max(1,
x_\tau)$. Let $a$, $(c_\tau)$ be as in (D) and $y\in\pbmRT$, $y\ne0$. Put $r_\tau=s_{\tau,a}/(|c_\tau|
|\tau(y)|)$; then $r_\tau|\tau(y)|\le x_\tau$, as $s_{\tau,a}\le|\kappa_\tau|^{2\varrho}/N$, and
$\min(1,r_\tau)\le B\max(1,|\tau(y)|)^{-1}$. With \eqref{pbm:eq-pf} and $\pbmden(y)=\auNm(\mathfrak
b_y)$,
\[
\prod_\tau\ell_K(r_\tau)\cdot\frac{\pbmn(y)^{K/2}}{\pbmden(y)^3}\le\prod_\tau\bigl(r_\tau|\tau(y)|\bigr)^K
\min(1,r_\tau)^6\cdot\auNm(\mathfrak b_y)^{-3}\le\prod_\tau x_\tau^KB^6\max(1,|\tau(y)|)^{-6}\cdot\auNm(
\mathfrak b_y)^{-3}.
\]
By Lemma~\ref{pbm:freq} the sum over $y\in X$ of $\beta$ is at most $\pbmdec\cdot(b^{2000n^9}\cdot
b^{n^9}\cdot162B^6)^f\le\pbmdec\,b^{2002n^9f}$, as $162\,b^{138(2\varrho+1)}\le b^{n^9}$ for $n\ge3$.

\emph{(W).} Let $i\le\varrho$, $p\in\pbmT^i$ with $\Lambda_i([p])\ne0$, and $\delta_l=p_0\cdots p_{l-1}
\in\pbmRT$ for $1\le l\le i$. Then $\delta_l\in\auO_T\cap\pbmRT=\auO_E$, and $|\sigma_{\tau,l}\tau(\delta_l
)|\le b^{60n^2}$ with $\sigma_{\tau,l}=\prod_{a<l}s_{\tau,a}\ge N^{-l}\ge b^{-60n^2}$; so $|\tau(\delta_l)|
\le R:=b^{120n^2}$. By \eqref{pbm:eq-pf} the factor of $\Lambda_i([p])$ belonging to $l$ is at most
\[
\Bigl(\prod_\tau\max\bigl(1,|\sigma_{\tau,l}\tau(\delta_l)|^{-1}\bigr)|\tau(\delta_l)|\Bigr)^{n^6}=\Bigl(
\prod_\tau\max\bigl(|\tau(\delta_l)|,\sigma_{\tau,l}^{-1}\bigr)\Bigr)^{n^6}\le R^{\,fn^6},
\]
so $\Lambda_i([p])\le R^{fn^7}=b^{120n^9f}$. The map $p\mapsto(\delta_1,\dots,\delta_i)$ is injective,
and by the argument with squares of Lemma~\ref{pbm:freq} there are at most $(4R+3)^{2f}$ integers of
$E$ with all conjugates of absolute value at most $R$; so at most $(4R+3)^{2fn}\le b^{242n^3f}$
heads have a non-zero weight. Hence $W_i\le b^{130n^9f}$. On the other hand $s_{\tau,a}\le b^{46
\varrho}$ and $\aurd\le b$ give $c_a\ge b^{-93nf}$ and $\prod_{a<i}c_a\ge b^{-93n^2f}$. So $W_i^2\le
b^{400n^9f}\prod_{a<i}c_a$.
\end{proof}

\subsubsection{The local systems and Siegel's step: statement}\label{pbm:ss-inputs}

\paragraph*{The constants}
Let $d=\dim W_{\lambda^X}$, $\pbmcG=\pbvol\auU(k)^2/(\pbvol\auU(r)\pbvol\auU(k-1))$ with $\pbvol
\auU(j)=\prod_{l=1}^j2\pi^l/(l-1)!$ (the product of the areas of the unit spheres of $\C^l$), and
$\pbmTh=\pi^{(2\varrho+1)k}\pbvol\auU(k)/\pbvol\auU(r)$ in \eqref{pbm:eq-beta}. Put
$\Lambda_*=\prod_\tau\prod_{a=0}^{\varrho}s_{\tau,a}^{\,k+2(\varrho-a)}$ and
\begin{equation}\label{pbm:eq-consts}
\begin{gathered}
\pbmaS=\#\pbR\cdot\#\pbUb\cdot\#\mathcal C\cdot\Bigl(\frac{\pi^k}{(k-1)!\,N^k}\Bigr)^f\aurd^{\,f/2}\,\pbmvol^{
\,k},\qquad\pbmrS=\frac1d\prod_\tau\Bigl(\sqrt{\pbmcG}\prod_{a=0}^{\varrho}s_{\tau,a}^{\,k-1+2(\varrho-a)}
\Bigr),\\
\pbmnuS=\frac{\#\pbUb}{N^{kf}}\cdot\frac{\aurd^{\,f\varrho/2}}{\pbmvol^{\,\varrho(\varrho+k)}\,\Lambda_*} .
\end{gathered}
\end{equation}
The statement \textup{(L2)} is homogeneous in the Haar measure of $U_\infty$: if this measure is
multiplied by $\lambda$, then, to keep (B2), $W$ and $W^\flat$ are multiplied by $\lambda^{1/2}$,
the orbital integrals over $G_\infty$ by $\lambda$, and $\nu$ by $\lambda$; so both sides of
\textup{(L2)} are multiplied by $\lambda^2$. (They do not depend on the Haar measure of $G_\infty$.)
The local inputs (R) and (S) below are, however, proved for the measure $du$ of convention (Cv3) of
\S\ref{pb:ss-setting}, which we use; for it $\nu=\pbmnuS$ (Lemma~\ref{pb:covol}; $\pbmnuS$ is the
constant $\nu_*$ there, as $\Lambda_*=\Lambda$).

\begin{theorem}[Local systems and Siegel's step]\label{pbm:inputs}
Bessel functions exist for the data (Definition~\ref{pb:bessel}), with \textup{(B5$'$)}, and they
can be chosen, together with functions $J^\tau_i$ $(\tau\in\auPl)$ and $J^T_i$ $(0\le i\le n$; the
$J^T_i$ depend on $T$, $\mathsf A$, $\mathsf d$ as in \textup{(F)}$)$, such that the following
holds, for the entry set $\pbmT$ of \S\ref{pbm:ss-cell}.
\begin{enumerate}
\item[(R)] \emph{(Real places.)} For every $\tau\in\auPl$, $(J^\tau_i)_{0\le i
\le n}$ is a sieve system at the place $(\C,\text{Lebesgue},e^{c-\aubar c},w\mapsto w^{-1})$ with mass $(1+b^{-10})/
\pbdeg_\tau$, weights $\Lambda^\infty$ and decay $\beta^\infty$; and its top form is the Gram form
$J^\tau_n(x,x')=\langle\mathsf N(x),\mathsf N(x')\rangle$ of the normalised orbital integrals of
\S\ref{pbm:ss-remains} ({\cite[Theorem~2.17\,(R-b)]{P52}}).
\item[(F)] \emph{(Finite places.)} Let $T$ be a finite set of finite places of $E$, stable under
conjugation, that contains the places above $MN$ and those at which some element of the set
$\mathsf B$ of \S\ref{pbm:ss-cell} is not a unit. For every set $\mathsf A$ of chosen places above
$p_0$ and every family $\mathsf d$ of frame vectors, $(J^{T}_i)_{0\le i\le
n}$ is a sieve system at the place $(E_T,dx,\psi_T,w\mapsto M^{-1}w^{-1})$ with mass $(1+b^{-6})\,b^{-6k\#\mathsf A}$,
weights $\Lambda^T$ and decay $\beta^T$.
\item[(S)] \emph{(Siegel's step.)} Let $T$, $\mathsf A$, $\mathsf d$ be as in \textup{(F)}, and let $J^T_n$
be the top form of \textup{(F)} for the same $T$, $\mathsf A$, $\mathsf d$ (it is built from a group
$\mathsf Z(\mathsf A,\mathsf d)$ that contains the components at $T$ of the integral elements whose
reductions lie in $\pbFix(\mathsf A,\mathsf d)$; see {\cite[Theorem~5.28]{P52}}). Let
$(c_j,c_j',s_j,\sigma_j)_{j\in\mathsf I}$ be a finite family of pairwise distinct quadruples of
classes $c_j,c_j'$, elements $s_j\in\pbFix(\mathsf A,\mathsf d)$ and $\sigma_j\in\pbmGf(E^+)$ such
that the last entry $t_j=t(\sigma_j)$ lies in $\pbmT$, $|\tau(t_j)|>N$ for all $\tau$, and
$h_{c_j}^{-1}\sigma_jh_{c_j'}$ stabilises the lattice at every finite place outside $T$. Then
\[
\pbmrS^{\,2}\sum_{j\in\mathsf I}\Bigl\|\sum_{p\in\pbmT^\varrho}\pbmOrb_{c_j,s_j,c_j'}(w_p\sigma_j)\Bigr\|^2\le
(1+b^{-6})\,\pbmaS\sum_{t\in\{t_j\}}\ \sum_{p,q\in\pbmT^\varrho}\Rel\Bigl(\prod_\tau J^\tau_n\bigl([p,t]_\tau,
[q,t]_\tau\bigr)\,J^T_n\bigl((p,t),(q,t)\bigr)\Bigr).
\]
\end{enumerate}
\end{theorem}
\smallskip\noindent The existence of the Bessel functions with \textup{(B5$'$)} is
Theorem~\ref{pbe:bessel}. \textup{(R)} is {\cite[Theorem~2.17]{P52}}, proved in {\cite[\S\S2, 3 and~4]{P52}} (its hypotheses (H0)--(H5) are checked for the Bessel functions of
Theorem~\ref{pbe:bessel} in {\cite[\S4]{P52}}); \textup{(F)} and \textup{(S)} are {\cite[Theorems~5.28 and~5.29]{P52}}, proved in {\cite[\S\S5 and~6]{P52}}. \S\ref{pbm:ss-remains} describes the forms. The proofs of these three theorems use nothing that depends on the present theorem, which serves only for the main term below.\par\smallskip

\subsubsection{The constant of the main term}\label{pbm:ss-const}

\begin{lemma}[Constants]\label{pbm:const}
Assume $k\ge512000\,n^9$, $n\ge100$, $b\ge1000$, $f\le b/1000$, $1\le|\kappa_\tau|$ for all
$\tau$, $M\le b^{23}$, $N\le b^{60n}$, $d\le b^{2kf}$ and $\prod_\tau|\kappa_\tau|=(M\aurd/2\pi)^f$.
Then:
\begin{enumerate}
\item[(C0)] $\displaystyle\pbmaS\prod_{a=0}^{\varrho}c_a\le(1+b^{-9})\,\pbmrS^{\,2}\,\pbdeg\cdot\pbDmain\,
\pbmnuS^{\,2}$;
\item[(C1)] $\displaystyle\pbmaS\prod_{a=0}^{\varrho}c_a\cdot b^{2403\,n^9f}\,\pbmdec\le b^{-10}\,b^{-6kf}\,
\pbmrS^{\,2}\,\pbDmain\,\pbmnuS^{\,2}$.
\end{enumerate}
\end{lemma}
\begin{proof}
From the definitions, $\prod_ac_a/\pbmrS^{\,2}=\pbmvol^{\,n}d^2/(\pbmcG^{\,f}\Lambda_*^2)$. Recall
$\pbDmain=(\#\mathcal C\,d)^2/(\pbdeg\,\pbmass_r)$, $\#\pbR=\#\pbUb\,\#\pbHb$ and $\#\mathcal C=\#\pbHb
(N^k-1)^f\pbmass_k$ (proof of Lemma~\ref{pb:sizes}).

(C0) After cancelling $d^2\,\#\pbUb^2\,\#\pbHb\,\#\mathcal C\,\Lambda_*^{-2}$ the inequality reads
\[
\Bigl(\frac{\pi^k}{(k-1)!\,N^k}\Bigr)^f\aurd^{\,f/2}\,\pbmvol^{\,k+n+2\varrho(\varrho+k)}\le(1+b^{-9})\,\pbmcG^{
\,f}\,\frac{(N^k-1)^f}{N^{2kf}}\cdot\frac{\pbmass_k}{\pbmass_r}\,\aurd^{\,f\varrho}.
\]
By Lemma~\ref{pb:massratio}, $\pbmass_r/\pbmass_k=\bigl(\prod_{l=k}^{r-1}\aurd^{\,l+1/2}\,l!/(2\pi)^{l+1}
\bigr)^f\cdot\mathsf L$ with $\mathsf L=\prod_{l=k}^{r-1}L(l+1,\chi^{\,l+1}_{E/E^+})\le e^{f2^{2-k}}$. We
have $\sum_{l=k}^{r-1}(l+\frac12)=(2\varrho+1)(k+\varrho+\frac12)$ and $\sum_{l=k}^{r-1}l=(2\varrho+1)k+
\varrho(2\varrho+1)$, and
\[
\pbmcG\,\frac{(k-1)!}{\pi^k}=\frac{2\,\pbvol\auU(k)}{\pbvol\auU(r)}=2\prod_{l=k}^{r-1}\frac{l!}{2\pi^{l+1}}=
2^{\,1+\sum_{l=k}^{r-1}l}\,\prod_{l=k}^{r-1}\frac{l!}{(2\pi)^{l+1}} .
\]
Inserting $\pbmvol=(2/\aurd)^f$, the powers of $\aurd$ on the two sides are $\aurd^{f(1/2-k-n-2\varrho(
\varrho+k))}$ and $\aurd^{f(\varrho-(2\varrho+1)(k+\varrho+1/2))}$, which are equal; the powers of $2$
are $2^{f(k+n+2\varrho(\varrho+k))}$ and $2^{f(1+(2\varrho+1)k+\varrho(2\varrho+1))}$, which are equal;
the factorials, the powers of $\pi$ and one factor $N^{-kf}$ cancel. What remains is $\mathsf L\le
(1+b^{-9})(1-N^{-k})^f$, which holds because $\log\mathsf L+2fN^{-k}\le f2^{3-k}\le b^{-10}$. (Unlike
Lemma~\ref{pb:sizes}, this uses the exact constant of the mass formula in
Lemma~\ref{pb:massratio}, that is, $\pbmass_{l+1}/\pbmass_l=2^{-f}|L(-l,\chi_{E/E^+}^{\,l+1})|$.)

(C1) By the same computation, with $\mathsf s=\aurd^{1/2}$, $\mathsf Q=M\mathsf s^2/2\pi$, $x=
\mathsf Q^{2\varrho+1}/(MN)$, so that $\pbmdec=(\pbmTh\,x^{k-n^7})^f$, and $\mathsf D=\varrho(2k+2
\varrho+1)$,
\[
\pbmaS\,\frac{\prod_ac_a}{\pbmrS^{\,2}}\,\pbmdec=\mathsf c\cdot\mathsf A^f,\qquad d\cdot\pbDid\cdot\pbmnuS^{
\,2}=\mathsf c\cdot\mathsf B^f,\qquad\mathsf c=\#\pbR\,\#\mathcal C\,d^2\Lambda_*^{-2},
\]
\[
\mathsf A=M^{\mathsf D}\,\frac{\pi^k\,\mathsf s}{(k-1)!\,N^k}\Bigl(\frac2{\mathsf s^2}\Bigr)^{k+n}\frac{
\pbmTh}{\pbmcG}\,x^{k-n^7},\qquad\mathsf B=\Bigl(\frac{M^{\mathsf D}}{N^k}\Bigl(\frac{\mathsf s^2}2
\Bigr)^{\varrho(k+\varrho)}\mathsf s^{\varrho}\Bigr)^2 .
\]
As $\pbmTh/\pbmcG=\pi^{(2\varrho+1)k}(k-1)!/(2\pi^k)$, a direct computation gives
\[
\mathsf A=\mathsf B\cdot\Bigl(\frac2{\mathsf s^2}\Bigr)^{n+2\varrho^2}\frac{\mathsf s^{1-2\varrho}}{2\,M^{
\varrho(2\varrho+1)}}\cdot x^{-n^7}\le\mathsf B\cdot b^{\,n+2\varrho^2}\,(MN)^{n^7}\le\mathsf B\cdot b^{\,n^9},
\]
because $\mathsf s\ge1$, $\mathsf Q\ge1$ and $MN\le b^{23+60n}$. Hence the left side of (C1) is at
most $b^{2404n^9f}\,d\,\pbDid\,\pbmrS^{\,2}\pbmnuS^{\,2}$. By Lemma~\ref{pb:sizes}, $\pbDid\le b^{-11kf}
\pbDmain$, and $d\le b^{2kf}$; as $2404n^9f+10\le3kf$, (C1) follows.
\end{proof}

The identity behind \textup{(C0)} is the reason for the shape of the main term: the constant of
Siegel's step, times the constants of the $n$ Poisson summations, divided by the normalisations of
the orbital integrals, is $\auDX^2\nu^2/\pbmass_r$ up to the factor $\mathsf L(1-N^{-k})^{-f}$; the
factor $1/\pbdeg$ is the mass of the local systems at the real places.

\subsubsection{Proof of \textup{(L2)} from Theorem~\ref{pbm:inputs}}\label{pbm:ss-L2}

\begin{theorem}[Theorem~\ref{pb:amp}\,\textup{(L2)}]\label{pbm:L2}
Assume Theorem~\ref{pbm:inputs} and the inequalities on the data listed in Theorem~\ref{pbm:sieve}
and Lemma~\ref{pbm:const}. Then, for the Bessel functions of Theorem~\ref{pbm:inputs}, every set
$\mathsf A$ of chosen places above $p_0$ and every family $\mathsf d$ of frame vectors,
\[
\pbLM\bigl(\{s\in\pbFix(\mathsf A,\mathsf d):s\text{ open}\}\bigr)\le(1+b^{-4})\,b^{-6k\#\mathsf A}\,\pbDmain
\,\nu^2 .
\]
\end{theorem}
\begin{proof}
By Lemma~\ref{pb:covol}, $\nu=\pbmnuS$. The left side is a sum of non-negative terms $\|A^\flat_{
T_\Lambda}(c,s,c')\|^2$ over the \emph{blocks} $(c,s,c',\Lambda)$ with $s\in\pbFix(\mathsf A,\mathsf d)$
open; it suffices to bound the sum over a finite set $\mathsf I$ of blocks with non-zero terms.
Let $T$ be as in Theorem~\ref{pbm:inputs}(F); then $M^{-1},N^{-1}\in\pbmRT$ and $t,t^{-1}\in\pbmRT$
for $t\in\pbmT$. By Proposition~\ref{pbm:classsum}(c) each block $j=(c_j,s_j,c_j',\Lambda_j)$ of
$\mathsf I$ has an element $\sigma_j\in\pbmGf(E^+)$ that determines $\Lambda_j$, so the quadruples
$(c_j,c_j',s_j,\sigma_j)$ are pairwise distinct; $t_j=t(\sigma_j)\in\pbmT$, $|\tau(t_j)|>N$,
$h_{c_j}^{-1}\sigma_jh_{c_j'}$ is integral outside $T$, and the term of the block is the square of
the norm of $\sum_{p\in\pbmT^\varrho}\pbmOrb_{c_j,s_j,c_j'}(w_p\sigma_j)$. By Siegel's step (S) and then
Theorem~\ref{pbm:sieve}, applied to the systems of (R) and (F) and to $\mathsf U=\{t_j\}$,
\[
\pbmrS^{\,2}\sum_{j\in\mathsf I}\|A^\flat_{T_{\Lambda_j}}(c_j,s_j,c_j')\|^2\le(1+b^{-6})\,\pbmaS\prod_{a=0}^{
\varrho}c_a\cdot\Bigl(\frac{(1+b^{-10})^f}{\pbdeg}\,(1+b^{-6})\,b^{-6k\#\mathsf A}+b^{2403n^9f}\,\pbmdec
\Bigr).
\]
By Lemma~\ref{pbm:const}, $(1+b^{-10})^f\le e^{fb^{-10}}\le1+b^{-9}$ and $\#\mathsf A\le f$, the right
side is at most
\[
(1+b^{-6})\bigl((1+b^{-9})^2(1+b^{-6})+b^{-10}\bigr)\,b^{-6k\#\mathsf A}\,\pbmrS^{\,2}\,\pbDmain\,\pbmnuS^{\,2}
\le(1+b^{-4})\,b^{-6k\#\mathsf A}\,\pbmrS^{\,2}\,\pbDmain\,\pbmnuS^{\,2},
\]
since $(1+u^6)((1+u^9)^2(1+u^6)+u^{10})\le1+u^4$ for $0\le u\le10^{-3}$. Divide by $\pbmrS^{\,2}>0$.
\end{proof}

The inequalities on the data used in Theorems~\ref{pbm:sieve}, \ref{pbm:L2} and
Lemma~\ref{pbm:const} are: $n=\varrho+1\ge100$ (so $m=2n$ in the notation of {\cite[Lemma~13.4]{P46}}); $k\ge512000\,n^9$; $b\ge1000$, $f\le b/1000$,
$\aurd\le b$, $N\le b^{60n}$; $M\le b^{23}$; $d\le b^{2kf}$ (all in {\cite[Lemma~13.4(S1)]{P46}}, (S5),
(S7), with (S7)(a), (e), (f) for $n\ge100$, $M\le b^{23}$ and $d\le b^{2kf}$); for
Theorem~\ref{pbm:inputs}\,(R), the hypotheses (H0)--(H5) of {\cite[Theorem~2.17]{P52}}, among them
$\varrho\ge10^8$, $k\ge n^7$, $k>10^9n^3+3n+6$ ({\cite[Lemma~13.4(S7)(a)]{P46}}) and $b^{6n+2}\ge8r$
((S7)(b)); and the condition \eqref{pbe:eq-kappa} on the generator $\varpi$, $1\le|\kappa_\tau|\le b^{23}$
(Lemma~\ref{pds:kappa}), together with $\prod_\tau|\kappa_\tau|=(M\aurd/2\pi)^f$, which holds because
$\prod_\tau|\tau(\varpi)|^2$ is the norm of the different, $\aurd^{\,2f}$.

\subsubsection{The forms of Theorem~\ref{pbm:inputs}}\label{pbm:ss-remains}

After this section, \textup{(L2)} rests on Theorem~\ref{pbm:inputs}, on Lemma~\ref{pb:covol}, and,
through Lemma~\ref{pbm:const}, on the mass formula in the form of Lemma~\ref{pb:massratio}. We
describe the forms of the three statements and how they are proved.

\emph{The forms.} At a real place $\tau$ write $\tilde G=G_\tau$ in the normalised basis of
\S\ref{pbe:ss-real}, in which a head $x=(x_0,\dots,x_\varrho)$ of complex numbers defines a Weyl
element $w_x$ and $x_\varrho$ is a last entry, $|x_\varrho|>1$. For $\sigma\in\pbmGf_\tau$ with last
entry $x_\varrho$ the top form is the Gram matrix of normalised orbital integrals,
\begin{gather*}
J^\tau_n(x,x')=\pbip{\mathsf N(x)}{\mathsf N(x')},\\
\mathsf N(x)=\frac{\sqrt{\pbmcG}\,(|x_\varrho|^2-1)^{(k-1)/2}}{\dim\mathbb W_\tau}\prod_{a<\varrho}|x_a|^{
k+1+2(\varrho-a-1)}\int_{G_\tau}W^\flat_\tau(y)\otimes\aubar{W^\flat_\tau((w_x\sigma)^{-1}y)}\,dy,
\end{gather*}
which depends only on $x$, $x'$ when $x_\varrho=x'_\varrho$ (by (B1) for $H_\tau$). At $T$ the
Bessel function is replaced by the vector of the translates, by the cosets of the principal
congruence subgroup in a group $\mathsf Z$ between it and the stabiliser of the lattice, of the
function $\chi_M$ that is $\psi_f(u)$ on $u(K\cap J)$, $u\in U(E_T)$, and $0$ elsewhere; its
orbital integrals are the Kloosterman sums of Proposition~\ref{pbm:orbit}, and $J^T_n(p,q)$ is the
integral, over the points with last entry $p_\varrho$ of the quadric of the vectors of norm $-N$
of $\pbmXp\otimes E_T$, of the pairing of two such orbital integrals, divided by $\pbvol(\mathsf Z)
\pbvol(K\cap J)\,\#\pbUb^{\,2}$ in a suitable normalisation. The group $\mathsf Z$ is the group of
the elements whose reductions lie in a subgroup of $\auG$ containing $\pbFix(\mathsf A,\mathsf d)$.
The lower forms are defined by integration over the entries, as (S1) requires: $J_i$ is the form
$J_n$ integrated over the last $n-i$ entries; by Fourier inversion along the Heisenberg groups of
the frame it is an integral, over a position $g$ and two frames, of the product of two
correlations of the Bessel function, and $J_0$ is the \emph{mass} $\int|\int_{\mathsf Z}\chi_M(ygy^{
-1})dy|^2dg$ (compare Lemma~\ref{pbm:pos}(c)).

\begin{enumerate}
\item[(R)] For the cut Bessel function of a discrete series whose parameter has a large gap: the
absolute convergence of the forms; the mass is at most $(1+b^{-10})/\pbdeg_\tau$ (the positions near the
identity give $1/\pbdeg_\tau$, by (B2) and the reproducing property; the remote positions are small);
the decay $\beta^\infty$ in the frequency, with the gain $(\pi/|\xi|)^{k-n^7}$, and the same gain for two
heads with different entries. This is {\cite[Theorem~2.17]{P52}}, proved in {\cite[\S\S2, 3 and~4]{P52}} (sizes and lengths, Plancherel's formula on quadrics, depth of frame vectors, separation
of the two frames by a weight made of depths, the gains at one position, the sum over dyadic shells
of positions), with the bounds for the Bessel function on the open cell of \S\ref{pbo:sec}.
\item[(F)] The same at the finite places, for the function $\chi_M$, with the decay $\beta^T$, and
in addition: the mass of the stabiliser of the lattice is at most $1+b^{-6}$ times its main part
({\cite[Corollary~5.12]{P52}}; the positions outside $K(M)$ are handled by the sum over dyadic shells, and
Theorem~\ref{pbm:E4} is not used), and the mass of $\mathsf Z$ is smaller by $b^{-6k}$ for each fixed
vector ({\cite[Theorem~5.24]{P52}}, through the index of the stabiliser of a vector). This is
{\cite[Theorem~5.28]{P52}}.
\item[(S)] Siegel's step. The quadruples $(c,c',s,\sigma)$ with a given last entry $t$ are
parametrised, through $\sigma\mapsto\sigma z$, by classes and rational points $x$ of the sphere
$h(x,x)=(t\aubar t-N^2)/N$ of $X_0$, see \eqref{pbm:eq-last}; by Proposition~\ref{pbm:orbit} the
squares of the class sums are values at these points of a non-negative function on the adelic
sphere, which is a pairing of orbital integrals at $T$ times the Gram matrix of orbital integrals
at the real places. The sum of such a function over the rational points of the sphere, and over the
classes, is bounded by $(1+b^{-6})$ times the constant $\pbmaS$ times its integral: Siegel's formula
for the representations of a number by the genus of the hermitian lattice $\auO_E^k$, derived with
all constants from Weil's Siegel--Weil theorem. This is {\cite[Theorem~5.29]{P52}}, with
{\cite[Proposition~6.35]{P52}}.
\end{enumerate}



\section{Multiplicities and orbital integrals}\label{B1:g3}
% ===== B-multiplicities.tex =====
% proofs/B-multiplicities.tex
% Equality of multiplicities: Proposition au:mult of sec-automorphic.tex.
% Labels of this file start with pm:.
\providecommand{\auO}{\mathcal O}
\providecommand{\auPl}{\Phi}
\providecommand{\auT}{\mathcal T}
\providecommand{\auAz}{\mathcal A_0}
\providecommand{\auG}{\overline G}
\providecommand{\audT}{d_{\auT}}
\providecommand{\auU}{\mathrm U}
\providecommand{\auGL}{\mathrm{GL}}
\providecommand{\auA}{\mathbb A}
\providecommand{\auF}{\mathbb F}
\providecommand{\aubar}[1]{\overline{#1}}
\providecommand{\auf}[1]{\texttt{#1}}
% macros of this file (prefix \pm)
\providecommand{\pmA}{\mathcal A}
\providecommand{\pmtr}{\operatorname{tr}}
\providecommand{\pmvol}{\operatorname{vol}}
\providecommand{\pmsupp}{\operatorname{supp}}
\providecommand{\pmN}{\operatorname{N}}
\providecommand{\pmTr}{\operatorname{Tr}}
\providecommand{\pmAd}{\operatorname{Ad}}
\providecommand{\pmch}{\operatorname{ch}}
\providecommand{\pmS}{\mathfrak S}
\providecommand{\pmn}{\mathfrak n}
\providecommand{\pmg}{\mathfrak g}
\providecommand{\pmGa}{\Gamma}
\providecommand{\pmell}{\mathrm{ell}}
\providecommand{\pmQ}{[G]}
\providecommand{\pmP}{\mathsf P}
\providecommand{\pmE}{\mathfrak E}

\subsection{Equality of multiplicities}\label{pm:sec}

This section proves {\cite[Proposition~14.36]{P46}}, in the precise form of
Proposition~\ref{pm:mult}, from two statements on discrete series of the real groups
$\auU(p,q)$ (Theorem~\ref{pm:arch}), known theorems of Harish-Chandra and of
Trombi--Varadarajan and Mili\v{c}i\'c which are proved here for the coefficients we use
(Remark~\ref{pm:rem-other}(1) and \S\ref{pbr:sec}), and from the fact that the irreducible labels
name cuspidal characters (\S\ref{pm:ss-mult}, (G)). Everything else is proved here or quoted from
named results.

\emph{The argument.} Let $V$, $V'$ be two hermitian spaces over $E$ of the same rank that are
isometric at all finite places; the groups $\auU(V)$ and $\auU(V')$ then have the same finite
adelic points $G_f$. At the infinite places we use, instead of a pseudo-coefficient, the kernel
$e=d_\pi\,\aubar{\langle\pi(\cdot)v,v\rangle}$ of a discrete series: it is integrable, because the
parameter is far from the walls, and $R(e)$ is the projection onto the forms of type $\pi$, so that
no other representation with the same infinitesimal character interferes. For a test function
$\varphi$ on $G_f$ the operator $R(e\otimes\varphi)$ has finite rank and its trace is the
integral over the diagonal of the kernel $\sum_\gamma(e\otimes\varphi)(x^{-1}\gamma y)$
(Proposition~\ref{pm:trace}); no truncation is needed, because $e$ is a cusp form in the sense of
Harish-Chandra and this makes all partial sums of the kernel bounded
(Proposition~\ref{pm:bound}). The trace is thus a sum over the characteristic polynomials $R$
of the rational points. If $\varphi$ has a supercuspidal component at a split place, the terms
of the $R$ that are products of two coprime factors vanish (Proposition~\ref{pm:red}; an
argument of Deligne--Kazhdan type), and the terms of the proper powers are
removed by a condition at an auxiliary place (cf.~\cite[Theorem~1]{FlickerKazhdan}, where a
condition at an auxiliary split place removes the regular non-elliptic terms), which is afterwards lifted by linear algebra in a
parahoric Hecke algebra (\S\ref{pm:ss-aux}). For irreducible $R$ the centraliser is the norm one
torus of a field; the conjugacy classes with polynomial $R$ in $\auU(V)$ and in $\auU(V')$ are
described by class field theory, and the sum of the orbital integrals of $e$ over the classes
in a stable class is, by Harish-Chandra, a character of the compact form, which does not see the
signs (Proposition~\ref{pm:ell}). There is no endoscopic term: a torus of this kind lies in no
proper endoscopic group. This is the comparison of Labesse~\cite{Labesse} in a case in which it can
be done by hand.

\subsubsection{Setting and statements}\label{pm:ss-setting}

\paragraph*{Groups}
$E$ is a CM field, $F=E^+$, $f=[F:\Q]$, and $\auPl$ is the set of infinite places of $F$, each
with a fixed extension $\tau\colon E\to\C$. Let $(V,h)$ be a non-degenerate hermitian space over
$E$ of dimension $n\ge2$, $G=G_V=\auU(V)$, a reductive group over $F$ with anisotropic centre,
$\pmGa=G(F)$, $G_\infty=\prod_{\tau}G_\tau$ with $G_\tau=\auU(V\otimes_{E,\tau}\C)$, and
$G_f=G(\auA_{F,f})$, the restricted product for the stabilisers of an $\auO_E$-lattice of $V$.
We write $\pmQ=\pmGa\backslash G(\auA_F)$. A finite place $v$ of $F$ that splits in $E$ comes
with a place of $E$ above it; then $G(F_v)=\auGL(V\otimes_EF_v)$, which we write $H_v$, and a basis
identifies $H_v$ with $\auGL_n(F_v)$. For $y\in G_f$, $y_v$ is its component at $v$.
The characteristic polynomial of $\gamma\in\pmGa$, an element of $E[X]$, is $\pmch(\gamma)$; for
split $v$, $\pmch(y_v)\in F_v[X]$ and $E\subset F_v$.

We fix a sign vector $s\in\{\pm1\}^n$ with $p$ entries $+1$ and $q$ entries $-1$, the same for all
$\tau$, and for each $\tau$ an orthogonal basis of $V_\tau=V\otimes_{E,\tau}\C$ on which $h$ is
$\operatorname{diag}(s)$. This identifies $G_\tau$ with $G_s=\{g\in\auGL_n(\C):g^*I_sg=I_s\}$,
$I_s=\operatorname{diag}(s)$; a change of the basis changes the identification by an inner
automorphism. Let $K_s=G_s\cap\auU(n)$, $T\subset K_s$ the diagonal torus, $\|g\|$ the operator
norm of $g\in G_s$ on $\C^n$ (so $\|g\|=\|g^{-1}\|\ge1$, as $g^{-1}=I_sg^*I_s$, and $\|kgk'\|=\|g\|$
for $k,k'\in K_s$), and for $g\in G_\infty$ put $\|g\|=\max_\tau\|g_\tau\|$ and
$\pmP(g)=\prod_\tau\|g_\tau\|$. The dimension of $G_s$ is $n^2$.

\paragraph*{Measures}
On each $G_\tau$, on $G_f$ and on every other unimodular group we fix a Haar measure; compact
groups at the infinite places have mass $1$. $G(\auA_F)$ has the product measure, $\pmGa$ the
counting measure, and $\pmQ$ the quotient measure, which is finite~\cite[Theorem~5.6(i)]{BorelFin}. The right
regular representation of $G(\auA_F)$ on $L^2(\pmQ)$ is $R$, and
$R(f)\psi(x)=\int_{G(\auA_F)}f(g)\psi(xg)\,dg$ for $f\in L^1$; thus $R(f_1*f_2)=R(f_1)R(f_2)$ with
$(f_1*f_2)(g)=\int f_1(h)f_2(h^{-1}g)\,dh$, and $R(f)^*=R(f^*)$ with $f^*(g)=\aubar{f(g^{-1})}$.
$C_c^\infty(G_f)$ is the space of locally constant functions of compact support (\emph{test
functions}).

\paragraph*{The input at the infinite places}
A \emph{parameter} is a vector $\lambda\in(\tfrac{n-1}2+\Z)^n$ with distinct entries. The pair
$(\lambda,s)$ satisfies the \emph{gap condition} if
\begin{equation}\label{pm:eq-gap}
|\lambda_i-\lambda_j|\ \ge\ 500\,n^3+10\,n^5\qquad\text{whenever }s_i\ne s_j .
\end{equation}
For a parameter $\lambda$ let $\chi_\lambda$ be the character of the irreducible representation of
the compact group $\auU(n)$ with Harish-Chandra parameter $\lambda$; on
$t=\operatorname{diag}(z_1,\dots,z_n)$ with distinct $z_i$ it is $\det(z_i^{\lambda_j})/\det(z_i^{
\varrho_j})$, up to the sign that makes it the character of a representation, where
$\varrho=(\tfrac{n-1}2,\dots,\tfrac{1-n}2)$. It is a symmetric function of $z_1,\dots,z_n$. For
$\sigma\in S_n$ put $t^\sigma=\operatorname{diag}(z_{\sigma(1)},\dots,z_{\sigma(n)})$.

\begin{theorem}[Harish-Chandra; Trombi--Varadarajan, Mili\v{c}i\'c]\label{pm:arch}
Let $(\lambda,s)$ satisfy \eqref{pm:eq-gap}. There are an irreducible unitary representation
$(\pi,\mathsf H)=\pi(\lambda,s)$ of $G_s$, square-integrable with formal degree $d_\pi$, and a
unit vector $v\in\mathsf H$ contained in a finite-dimensional $K_s$-stable subspace, such that
$c(g)=\langle\pi(g)v,v\rangle$ has the following properties.
\begin{enumerate}
\item[(a)] $\sup_{g\in G_s}\|g\|^{5n^3}\,|c(g)|<\infty$.
\item[(b)] For every $t\in T$ with distinct eigenvalues the function $x\mapsto c(x^{-1}tx)$ is
integrable on $G_s$, and
\[
\sum_{\sigma\in S_n}d_\pi\int_{G_s}\aubar{c(x^{-1}t^\sigma x)}\,dx=p!\,q!\,(-1)^{pq}\,
\aubar{\chi_\lambda(t)} .
\]
\end{enumerate}
\end{theorem}
\noindent Here $\pi(\lambda,s)$ is the discrete series of $G_s\cong\auU(p,q)$ with Harish-Chandra
parameter $\lambda$, for the positive system in which $e_i-e_j$ is positive when
$\lambda_i>\lambda_j$, $v$ is a unit highest weight vector of its lowest $K_s$-type
$\Lambda=\lambda+\rho_{\mathrm n}-\rho_{\mathrm c}$, and inner products are linear in the first
variable, so that $c$ is the function $c_\lambda$ of {\cite[\S3.2]{P49}} (written there with the
opposite convention); if $q=0$ it is the representation with character $\chi_\lambda$.
\emph{Proof.} (a) is proved in item (1) of Remark~\ref{pm:rem-other} from {\cite[Theorem~3.17]{P49}} and {\cite[Lemma~3.25]{P49}}, which do not depend on the present theorem; items (2) and (3) of that remark are comments made under the theorem and are not used in its proof.
(b) is Theorem~\ref{pbr:main}, which proves it, together with the unstable form
$d_\pi\int_{G_s}c(x^{-1}tx)\,dx=(-1)^{pq}\sum_{w\in S_p\times S_q}\operatorname{sgn}(w)t^{\kappa\circ w}
/\prod_{\lambda_a>\lambda_b}(t_a-t_b)$, $\kappa=\lambda+\frac{n-1}2\mathbf 1$, from
{\cite[Theorem~3.17]{P49}}, Schur orthogonality, admissibility, Schmid's theorem on the $K$-types of a
discrete series, and the theory of compact Lie groups; Harish-Chandra's character theory is not used.
The product $d_\pi\,dx$ does not depend on the choice of the Haar measure $dx$.

\begin{remark}\label{pm:rem-arch}
(1) By Harish-Chandra, $\Theta_\pi(t)=(-1)^{pq}\sum_{w\in S_p\times S_q}\operatorname{sgn}(w)\,
t^{w\lambda}/D_\lambda(t)$ on the regular set of $T$, with $D_\lambda$ the Weyl denominator of the
positive system of $\lambda$; for the orbital integrals of $c$ this is \eqref{pbr:eq-unstable},
proved in Theorem~\ref{pbr:main}. As $D_\lambda(t^\sigma)=\operatorname{sgn}(\sigma)D_\lambda(t)$, the sum
of $\Theta_\pi(t^\sigma)$ over $\sigma\in S_n$ is $p!q!(-1)^{pq}\sum_{w\in S_n}\operatorname{sgn}(w)
t^{w\lambda}/D_\lambda(t)=p!q!(-1)^{pq}\chi_\lambda(t)$. This is the form (b). (2) The elements
$t^\sigma$, $\sigma\in S_n/(S_p\times S_q)$, represent the conjugacy classes of $G_s$ in the set of
elements with the eigenvalues of $t$: two of them are conjugate in $G_s$ if and only if the same
eigenvalues occur on the positive coordinates.
\end{remark}

\begin{remark}[Relation to other sections]\label{pm:rem-other}
(1) Part (a) of Theorem~\ref{pm:arch} holds for every pair $(\lambda,s)$ satisfying
\eqref{pm:eq-gap}: by {\cite[Theorem~3.17]{P49}} (whose hypothesis $h\ge n^2+1$ follows from
\eqref{pm:eq-gap}) and {\cite[Lemma~3.25]{P49}}, $|c(g)|\le(2/\|g\|)^{c}$ for every integer $c$ up to the
gap, which exceeds $5n^3$. Here $c$ is the coefficient of a unit highest weight vector of the lowest
$K$-type, with $\pi(\lambda,s)$ the discrete series of parameter $\lambda$; with inner products
linear in the first variable, $c(g)=\langle\pi(g)v,v\rangle$ is exactly this function $c_\lambda(g)=
\langle v,\pi_\lambda(g)v\rangle$ of {\cite[\S3.2]{P49}}, where the opposite convention is used.
(2) The space $\pmA(V;K)$ of Definition~\ref{pm:def-forms} below is defined in $L^2$. By
Proposition~\ref{pm:trace}(a) its elements are bounded and continuous, so for $V=Y_0$ and $K=J_M$
it is the space of forms of type $e_0$ and level $J_M$ of {\cite[Definition~8.1]{P49}}, with the
same sign vector \eqref{pm:eq-signs} as in {\cite[(74)]{P49}}; this gives the boundedness that
Proposition~\ref{pds:A0} asks for, under Theorem~\ref{pm:arch}. (3) Lemma~\ref{pm:cusp} proves
the cusp property (A2) of {\cite[\S8.1]{P49}} for the coefficient $c$ of Theorem~\ref{pm:arch} from
the bound (a).
\end{remark}

In the rest of this section we \emph{assume Theorem~\ref{pm:arch}}. Let $\lambda=(\lambda_\tau)_{
\tau\in\auPl}$ be parameters such that every $(\lambda_\tau,s)$ satisfies \eqref{pm:eq-gap}; let
$\pi_\tau$, $v_\tau$, $c_\tau$ be as in the theorem, transported to $G_\tau$, and
\begin{equation}\label{pm:eq-e}
e_\tau=d_{\pi_\tau}\,\aubar{c_\tau},\qquad e=e_V=\bigotimes_{\tau\in\auPl}e_\tau\colon G_\infty\to\C,
\qquad B=5n^3 .
\end{equation}

\begin{definition}\label{pm:def-forms}
For a compact open subgroup $K\subset G_f$, the space of \emph{forms of type $(\lambda,s)$ and
level $K$} on $V$ is
$\pmA(V;K)=\{\psi\in L^2(\pmQ):R(e)\psi=\psi,\ R(k)\psi=\psi\text{ for }k\in K\}$.
\end{definition}

By Lemma~\ref{pm:kernel}(a), $R(e)$ is an orthogonal projection. For $\psi=R(e)\psi$ the map
$\sum_ia_i\pi_\infty(g_i)v\mapsto\sum_ia_iR(g_i)\psi$, with $\pi_\infty=\bigotimes_\tau\pi_\tau$ and
$v=\bigotimes_\tau v_\tau$, is well defined and is $\|\psi\|$ times an isometry, because
$\langle R(g)\psi,\psi\rangle=\langle R(e)R(g)R(e)\psi,\psi\rangle$ and $e*\delta_g*e=
\langle\pi_\infty(g)v,v\rangle\,e$ by the orthogonality relations. Hence
$\pmA(V;K)\cong\operatorname{Hom}_{G_\infty}(\pi_\infty,L^2(\pmQ)^K)$, $\psi=A(v)\leftrightarrow A$:
the space, with the action of the Hecke algebra of $(G_f,K)$, does not depend on the choice of
$v$, and it is the space of the automorphic forms on $V$ of level $K$ that generate $\pi_\infty$
at infinity, evaluated at the vector $v$. This is the space meant in {\cite[\S14.3]{P46}}.

\paragraph*{The comparison}
A \emph{comparison datum} is $(V,V',\iota)$ with $V$, $V'$ hermitian spaces of dimension $n$, with
sign vectors $s$, $s'$ at every infinite place, and $\iota\colon V\otimes_F\auA_{F,f}\to V'\otimes_F
\auA_{F,f}$ an isometry of $\auA_{E,f}$-modules that carries a lattice to a lattice at almost all
places. It identifies $G_{V'}(\auA_{F,f})$ with $G_f=G_V(\auA_{F,f})$, and $H_v$ for $V'$ with
$H_v$ for $V$. The parameters $\lambda_\tau$ are the same for $V$ and $V'$, and we assume
\eqref{pm:eq-gap} for $(\lambda_\tau,s)$ and for $(\lambda_\tau,s')$; objects of $V'$ carry a
prime. A function $\beta\in C_c^\infty(H_v)$ is a \emph{supercusp function} if
$\int_N\beta(xny)\,dn=0$ for all $x,y\in H_v$ and the unipotent radical $N$ of every proper
parabolic subgroup of $H_v$.

\begin{theorem}\label{pm:main}
Assume Theorem~\ref{pm:arch} and let $f=[F:\Q]$ be even. Let $(V,V',\iota)$ be a comparison datum,
$v_0$ a finite place of $F$ that splits in $E$, and $\varphi=\varphi^{v_0}\otimes\beta\in
C_c^\infty(G_f)$ with $\beta$ a supercusp function on $H_{v_0}$. Then $R(e\otimes\varphi)$ and
$R'(e'\otimes\varphi)$ have finite rank and
\[
\pmtr\bigl(R(e_V\otimes\varphi)\mid L^2(\pmQ)\bigr)=\pmtr\bigl(R'(e_{V'}\otimes\varphi)\mid
L^2([G'])\bigr).
\]
If $\varphi$ is bi-invariant under a compact open $K\subset G_f$, these are the traces of
$R(\varphi)$ on $\pmA(V;K)$ and of $R'(\varphi)$ on $\pmA(V';K)$.
\end{theorem}

The proof occupies \S\S\ref{pm:ss-kernel}--\ref{pm:ss-aux}. We now state the consequence for the
data. In the notation of {\cite[\S14.3]{P46}}, $n=r$, and the parameter is $\varsigma_\tau$, whose
entries we name: $\varsigma_\tau(i)=\tfrac k2-i$ for $1\le i\le k-1$ (the \emph{small} entries) and
$\varsigma_\tau(k-1+a)=-\kappa(\tau)-a\Delta$ for $1\le a\le m$. The sign vectors are
\begin{equation}\label{pm:eq-signs}
s_{W_j}(k-1+a)=\begin{cases}+1,&a\le j,\\-1,&a>j,\end{cases}\qquad
s_{Y_0}(k-1+a)=(-1)^{a-1},\qquad s(i)=+1\ \text{for }i\le k-1 ,
\end{equation}
so that $W_j$ has $m-j$ and $Y_0$ has $m/2=\varrho+1$ negative signs, as in
{\cite[\S14.3]{P46}}; the discrete series at $\tau$ is $\pi(\varsigma_\tau,s)$. The spaces
$\auAz$ and $\pmA(W_j)$ of {\cite[\S14.3]{P46}} are $\pmA(Y_0;J_M)$ and $\pmA(W_j;J_M)$. Let
$\mathcal H^M$ be the algebra of the functions in $C_c^\infty(G_f)$ of the form
$\varphi_S\otimes1_{J_M^{S}}$, with $S$ a finite set of places of $F$ that split in $E$ and do not
divide $M$, $K^\circ\supset J_M$ the stabiliser of the lattice, $\varphi_S$ bi-invariant under
$K^\circ_S=J_{M,S}$, and $J_M^S$ the part of $J_M$ outside $S$; it is commutative (the spherical
Hecke algebras of the groups $\auGL_r(F_v)$ are), stable under $\varphi\mapsto\varphi^*$, and it acts
on $\pmA(V;J_M)$ by $\varphi\mapsto\pmvol(J_M)^{-1}R(\varphi)$, commuting with $\auG$. These are the Hecke operators
at the split places prime to $M$.

\begin{proposition}[{\cite[Proposition~14.36]{P46}}, precise form]\label{pm:mult}
Assume Theorem~\ref{pm:arch} and the property \textup{(G)} of \S\ref{pm:ss-mult}. Assume that $f$
is even and that for each of the sign vectors \eqref{pm:eq-signs} and each $\tau$ the entries of
$\varsigma_\tau$ on coordinates of different signs differ by at least $500r^3+10r^5$. (Both hold
for the data: $f=\#(\Z/q_1)^\times$, and these entries differ by at least $\Delta=2r^8\ge500r^3+10r^5$.) Then for
$1\le j\le m$ and every $T\in\mathcal H^M$
\[
\pmtr\bigl(T\mid e_0\pmA(W_j)\bigr)=\pmtr\bigl(T\mid\auT\bigr).
\]
In particular $\dim e_0\pmA(W_j)=\audT$, and every character of $\mathcal H^M$ has the same
multiplicity in $e_0\pmA(W_j)$ and in $\auT$.
\end{proposition}

\subsubsection{The kernel at infinity}\label{pm:ss-kernel}

\begin{lemma}\label{pm:kernel}
Let $e_\tau$ be as in \eqref{pm:eq-e}, on $G_\tau\cong G_s$.
\begin{enumerate}
\item[(a)] $e_\tau$ is continuous, $e_\tau^*=e_\tau$, $|e_\tau(g)|\le C\|g\|^{-B}$,
$e_\tau\in L^1$, and $e_\tau*e_\tau=e_\tau$.
\item[(b)] There is $\alpha_\tau\in C_c^\infty(G_\tau)$ with $\alpha_\tau*e_\tau=e_\tau$. Hence
$e_\tau=e_\tau*\alpha_\tau^*$, $e_\tau$ is smooth, and $|(De_\tau)(g)|\le C_D\|g\|^{-B}$ for every
left-invariant differential operator $D$.
\item[(c)] $\int_{G_\tau}\|g\|^{-B'}dg<\infty$ for $B'>n^2$.
\end{enumerate}
The same holds for $e$ on $G_\infty$ with $\alpha=\bigotimes\alpha_\tau$, with $\pmP(g)^{-B}$ in place
of $\|g\|^{-B}$.
\end{lemma}
\begin{proof} We begin with (c). It concerns the group alone: its proof uses neither $e_\tau$ nor Theorem~\ref{pm:arch}, and it is this part, and no other, that enters the proof of {\cite[Lemma~3.25]{P49}}, on which Theorem~\ref{pm:arch}(a) rests. Parts (a) and (b) then use Theorem~\ref{pm:arch}.
(c) In the decomposition $G_s=K_sA^+K_s$, with $A\cong\R^{\min(p,q)}$ and coordinates
$t_1\ge\dots\ge t_{\min(p,q)}\ge0$ on $A^+$, one has $\|kak'\|=e^{t_1}$ and the Haar measure is
$\delta(a)\,dk\,da\,dk'$ with $\delta(a)\le e^{2\rho(\log a)}\le e^{n(n-1)t_1}$
(\cite[Ch.~I, \S5, no.~2]{Helgason}, \cite[Lemma~2.4.2]{WallachRRG1}: up to a constant, $\delta(a)$ is the product of $\sinh\alpha(\log a)$ over the positive restricted roots with multiplicities). (a) Continuity and $e^*=e$ are clear, the bound is
Theorem~\ref{pm:arch}(a), integrability follows from (c), and $e*e=e$ is the orthogonality relation
$\int c(h)c(h^{-1}g)\,dh=d_\pi^{-1}c(g)$ for square-integrable
representations~\cite[Proposition~9.6, p.~284]{Knapp}. (b) Let $\mathsf H_0$ be the sum of the $K_\tau$-isotypic
subspaces of the types occurring in the span of $K_\tau v$; it is finite-dimensional, since
irreducible unitary representations are admissible \cite[Theorem~8.1, p.~205]{Knapp}, and the projection
onto it is $\pi(\chi)$ for a central function $\chi$ on $K_\tau$. The subspace
$\{\pi(\chi)\pi(\alpha_0)v:\alpha_0\in C_c^\infty(G_\tau)\}$ of $\mathsf H_0$ contains $v$ in its
closure, so it contains $v$: $\pi(\alpha_1)v=v$ with $\alpha_1=\chi*\alpha_0\in C_c^\infty$. Put
$\alpha=\alpha_1$. Then, as $e(h^{-1}y)=d_\pi\aubar{\langle\pi(y)v,\pi(h)v\rangle}$ and the inner product is conjugate-linear in the second variable, $(\alpha*e)(y)=d_\pi\aubar{\langle\pi(y)v,\pi(\alpha_1)v\rangle}=e(y)$.
Taking adjoints, $e=e*\alpha^*$, so $De=e*D\alpha^*$, and $|(e*\beta)(g)|\le
\int|e(gh^{-1})||\beta(h)|dh\le C\|g\|^{-B}\int\|h\|^B|\beta(h)|dh$.
\end{proof}

\begin{lemma}[$e$ is a cusp form]\label{pm:cusp}
Let $N'$ be the unipotent radical of a proper parabolic subgroup of $G_\tau$. Then
$\int_{N'}|e_\tau(x\nu y)|\,d\nu<\infty$ and $\int_{N'}e_\tau(x\nu y)\,d\nu=0$ for all $x,y\in G_\tau$.
\end{lemma}
\begin{proof}
This is Harish-Chandra's theorem that the coefficients of the discrete series are cusp forms; we
give his proof under the bound of Lemma~\ref{pm:kernel}. Write $e=e_\tau$ and $P'=M'A'N'$ for the
Langlands decomposition of the parabolic subgroup, with $A'=\exp\mathfrak a'\ne1$ consisting of
positive hermitian matrices, and $\mathfrak n'$ for the Lie algebra of $N'$. For
$\nu=\exp X\in N'$ we have $(\nu-1)^{n}=0$, so $|X|\le C\|\nu\|^{n-1}$ and
$\int_{N'}\|\nu\|^{-B/4}d\nu\le C\int_{\mathfrak n'}(1+|X|)^{-B/(4(n-1))}dX<\infty$, as
$B/(4(n-1))>n^2\ge\dim\mathfrak n'$. For $a\in A'$ the eigenvalues of $a\nu$ are those of $a$, so
$\|a\nu\|\ge\|a\|$, and $\|\nu\|\le\|a^{-1}\|\|a\nu\|\le\|a\nu\|^2$; hence
$\|a\nu\|\ge\|a\|^{1/2}\|\nu\|^{1/4}$. Fix $x$ and put $\Phi(g)=\int_{N'}e(x\nu g)\,d\nu$. As
$\|a\nu\|\le\|x\|\|xa\nu g\|\|g\|$, the substitution $\nu\mapsto a\nu a^{-1}$ gives
\begin{equation}\label{pm:eq-cuspdecay}
|\Phi(ag)|=\delta(a)\Bigl|\int_{N'}e(xa\nu g)\,d\nu\Bigr|\le C\|x\|^B\|g\|^B\,\|a\|^{2n^2-B/2}
\int_{N'}\|\nu\|^{-B/4}d\nu ,
\end{equation}
where $\delta(a)=|\det\pmAd(a)|_{\mathfrak n'}|\le\|a\|^{2n^2}$; in particular the integral defining
$\Phi$ converges absolutely, and by Lemma~\ref{pm:kernel}(b) it may be differentiated under the
integral sign. Thus $\Phi$ is smooth, left invariant under $N'$, and an eigenfunction of the
centre $Z(\pmg)$ of the enveloping algebra acting on the right, because $e$ is ($\pi$ is
irreducible, so $Z(\pmg)$ acts by scalars on the smooth vector $v$). For central $z$ the right
action of $z$ is the left action of its transpose, which is central. By Harish-Chandra, every $z\in Z(\pmg)$ lies in $\mu(z)+U(\pmg)\mathfrak n'$ with
$\mu(z)\in Z(\mathfrak m'\oplus\mathfrak a')$, and $Z(\mathfrak m'\oplus\mathfrak a')$ is a finitely
generated module over $\mu(Z(\pmg))$. (In \cite[proof of Proposition~8.22, p.~225]{Knapp}, $Z(\pmg)\subseteq\mathfrak n'U(\pmg)\oplus U(\mathfrak m'\oplus\mathfrak a')$, and the projection $\mu$ to the second summand satisfies $\gamma'_{\mathfrak a'\oplus\mathfrak m'}\circ\mu=\gamma'$, (8.34) there; applying the principal antiautomorphism of $U(\pmg)$, which preserves $Z(\pmg)$ and $U(\mathfrak m'\oplus\mathfrak a')$, gives the form with $U(\pmg)\mathfrak n'$. For the finiteness: via these maps, with a Cartan subalgebra $\mathfrak h$ of $\mathfrak m'\oplus\mathfrak a'$, $\mu(Z(\pmg))$ and $Z(\mathfrak m'\oplus\mathfrak a')$ correspond to the images of $U(\mathfrak h)^W$ and $U(\mathfrak h)^{W_{M'}}$ under translation automorphisms $\sigma_1,\sigma_2$ of $U(\mathfrak h)$ \cite[Theorem~8.18]{Knapp}; $U(\mathfrak h)$ is a finitely generated module over $U(\mathfrak h)^W$ \cite[Theorem~8.19]{Knapp}, hence over the Noetherian ring $\sigma_1(U(\mathfrak h)^W)$, so its submodule $\sigma_2(U(\mathfrak h)^{W_{M'}})$ is finitely generated.) Elements of $U(\pmg)\mathfrak n'$, acting on the left, kill
left $N'$-invariant functions. So $\Phi$ is an eigenfunction of the left action of
$\mu(Z(\pmg))$, and so are its left derivatives along $H\in\mathfrak a'$, which are again left
$N'$-invariant; as $H\in Z(\mathfrak m'\oplus\mathfrak a')$ is integral over $\mu(Z(\pmg))$, the
function $u(r)=\Phi(\exp(rH)g)$ satisfies a linear differential equation with constant
coefficients: it is an exponential polynomial. By \eqref{pm:eq-cuspdecay} and $2n^2<B/2$ it tends
to $0$ for $r\to\pm\infty$ if $H\ne0$, so $u=0$ and $\Phi(g)=u(0)=0$.
\end{proof}

\subsubsection{Reduction theory and bounds for the kernel}\label{pm:ss-bound}

Let $\varrho_V$ be the Witt index of $V$. If $\varrho_V=0$ then $\pmQ$ is compact~\cite[Theorem~5.6(ii)]{BorelFin} and
everything in this subsection is trivial or void; we assume $\varrho_V\ge1$ and fix a Witt
decomposition $V=\langle e_1,\dots,e_{\varrho_V}\rangle\oplus W\oplus\langle e_1',\dots,
e'_{\varrho_V}\rangle$, with $h(e_i,e'_j)=\delta_{ij}$, $e_i$, $e_i'$ isotropic and orthogonal to the
anisotropic space $W$. Let $V_i=\langle e_1,\dots,e_i\rangle$, $P_i$ the stabiliser of $V_i$ (the
maximal standard parabolic subgroups), $N_i$ its unipotent radical, $P_0=\bigcap_iP_i$. Every proper
parabolic $F$-subgroup of $G$ is conjugate under $\pmGa$ to a subgroup of some $P_i$, and an element
of $\pmGa$ lies in a proper parabolic $F$-subgroup if and only if it stabilises a non-zero isotropic
subspace of $V$. For $a=(a_1,\dots,a_{\varrho_V})\in\R_{>0}^{\varrho_V}$ we also write $a$ for the
element of $G_\infty$ that, at every $\tau$, multiplies $e_i$ by $a_i$, $e'_i$ by $a_i^{-1}$ and is
the identity on $W$. Put $\beta_i(a)=a_i/a_{i+1}$ for $i<\varrho_V$, $\beta_{\varrho_V}(a)=
a_{\varrho_V}$, $\beta(a)=\max_i\beta_i(a)$. We choose the bases at infinity (that is, the maximal
compact subgroup $K_\infty=\prod K_\tau$ and the norms) so that $e_i,e'_i$ and an orthogonal basis
of $W_\tau$ are orthonormal for the norm $|\cdot|_\tau$ of $\C^n$; then $a$ is positive and
diagonal, $|h_\tau(u,u')|\le|u|_\tau|u'|_\tau$, and $\|a\|=\max_i\max(a_i,a_i^{-1})$. Let $K_f^\circ$
be the stabiliser of a lattice $\Lambda\subset V$ containing the Witt basis.

\emph{Siegel sets} (\cite[Th\'eor\`eme~7]{Godement63}; \cite[\S I.2.1, p.~20]{MW}). There are a compact set
$\Omega\subset P_0(\auA_F)$ and $t\in(0,1]$ such that $G(\auA_F)=\pmGa\,\pmS$ with
\[
\pmS=\{\omega ak:\ \omega\in\Omega,\ a\in\R_{>0}^{\varrho_V},\ \beta_i(a)\ge t\ (1\le i\le\varrho_V),\
k\in K_\infty K_f^\circ\}.
\]
For such $a$: $a_i\ge t^{n}$ for all $i$, and $\|a\|\le t^{-n}\max(1,\beta(a))^{n/2}$. For
$x=\omega ak\in\pmS$ we let $i_0=i_0(x)$ be an index with $\beta_{i_0}(a)=\beta(a)$ and put
$P_x=P_{i_0}$, $N_x=N_{i_0}$. In the estimates below, constants depend on $V$, $\Omega$, $t$, the
function $e$ and the test function, but not on $x$.

\begin{lemma}\label{pm:siegel}
Let $x=\omega ak\in\pmS$, $P=P_x$, $N=N_x$.
\begin{enumerate}
\item[(a)] Let $\Omega_f\subset G_f$ be compact. There is $c>0$ such that
$\pmP\bigl((x^{-1}\gamma x)_\infty\bigr)\ge c\,\beta(a)^{f}$ for every $\gamma\in\pmGa\smallsetminus
P(F)$ with $(x^{-1}\gamma x)_f\in\Omega_f$.
\item[(b)] Let $K'\subset G_f$ be compact open. There is $C$ such that $N(\auA_F)\,x\subset
N(F)\,x\,(Z_\delta\times K')$ with $\delta=C\beta(a)^{-1}$, where $Z_\delta=\{z\in G_\infty:
\|z_\tau-1\|\le\delta\text{ for all }\tau\}$.
\item[(c)] Let $B'>n^2$ and $M_{B'}(y)=\pmP(y_\infty)^{-B'}1_{\Omega_f}(y_f)$. There is $C$ such that
$\sum_{\gamma\in\pmGa}M_{B'}(g^{-1}\gamma x)\le C\|a\|^{2fn^2}$ for all $g\in G(\auA_F)$. For $x$ in a compact subset of $G(\auA_F)$, not necessarily in $\pmS$, the series converges
uniformly in $(g,x)$.
\end{enumerate}
\end{lemma}
\begin{proof}
Let $\mathcal B$ be the basis of $V$ formed by the $e_i$, a basis of $W$ in $\Lambda$ and the
$e'_i$, and for $b\in\mathcal B$ let $\chi_b(a)\in\{a_i,1,a_i^{-1}\}$ be the eigenvalue of $a$. Split
$\mathcal B=\mathcal B_1\sqcup\mathcal B_2\sqcup\mathcal B_3$ with $\mathcal B_1=\{e_l:l\le i_0\}$,
a basis of $V_{i_0}$, $\mathcal B_3=\{e'_l:l\le i_0\}$, and $\mathcal B_1\sqcup\mathcal B_2$ a basis
of $V_{i_0}^\perp$. Put $\theta=t^{-n}a_{i_0}^{-1}$ and $\Theta=t^{-n}\max(1,a_{i_0+1})$, with
$a_{\varrho_V+1}:=1$. As $a_l\ge t^na_{i_0}$ for $l\le i_0$ and $t^n\le a_l\le t^{-n}a_{i_0+1}$ for
$l>i_0$,
\begin{equation}\label{pm:eq-contract}
\chi_b^{-1}\le\theta\ (b\in\mathcal B_1),\qquad\chi_b\le\theta\ (b\in\mathcal B_3),\qquad
\chi_b^{\pm1}\le\Theta\ (b\in\mathcal B_2),\qquad\theta\max(\theta,\Theta)\le t^{-5n}\beta(a)^{-1},
\end{equation}
the last because $a_{i_0}=\beta_{i_0}(a)a_{i_0+1}\ge t^n\beta(a)$ and $\beta(a)\ge t$.

(a) As $\gamma V_{i_0}\not\subset V_{i_0}=(V_{i_0}^\perp)^\perp$, there are $b\in\mathcal B_1$ and
$b'\in\mathcal B_1\sqcup\mathcal B_2$ with $\xi=h(b',\gamma b)\ne0$. Let $\Omega^{\mathrm{fin}}$ be
the set of finite components of $\Omega$ and $\Omega'_f=\Omega^{\mathrm{fin}}K_f^\circ$, a compact
set containing $x_f$. The finite component of $\gamma$ lies in the compact set
$\Omega'_f\Omega_f\Omega_f'^{-1}$, so $\gamma\Lambda\subset d^{-1}\Lambda$ and $\xi\in d^{-1}\auO_E$
for an integer $d\ge1$ independent of $x$ and $\gamma$; hence $\prod_\tau|\tau(\xi)|\ge d^{-f}$. On
the other hand
$|\tau(\xi)|=|h_\tau(x_\tau^{-1}b',(x^{-1}\gamma x)_\tau x_\tau^{-1}b)|\le|x_\tau^{-1}b'||x_\tau^{-1}b|
\,\|(x^{-1}\gamma x)_\tau\|$, and $x_\tau^{-1}b=k_\tau^{-1}a^{-1}\omega_\tau^{-1}b$, where
$\omega_\tau^{-1}b$ is a bounded combination of the vectors of $\mathcal B_1$, since $P_0$
preserves $V_{i_0}$; likewise $\omega_\tau^{-1}b'$ is a bounded combination of the vectors of
$\mathcal B_1\sqcup\mathcal B_2$. By \eqref{pm:eq-contract},
$|x_\tau^{-1}b'||x_\tau^{-1}b|\le C\theta\max(\theta,\Theta)\le C'\beta(a)^{-1}$, and the product over
$\tau$ gives (a).

(b) $N(\auA_F)=N(F)\,(C_\infty\times U_f)$ for every open compact subgroup $U_f$ of
$N(\auA_{F,f})$ and a suitable compact $C_\infty\subset N(F\otimes\R)$ (strong approximation for
unipotent groups). Take $U_f\subset\bigcap_{y\in\Omega'_f}yK'y^{-1}$. For $\nu\in N(F)$,
$c\in C_\infty$, $u\in U_f$ we have $\nu cu\,x=\nu x\cdot(x^{-1}cux)$, the finite component of
$x^{-1}cux$ lies in $K'$, and its infinite component is $k^{-1}a^{-1}n'ak$ with
$n'=\omega_\infty^{-1}c\,\omega_\infty$ in a compact subset of $N(F\otimes\R)$. As $n'-1$ maps
$V$ into $V_{i_0}^\perp$, $V_{i_0}^\perp$ into $V_{i_0}$ and $V_{i_0}$ to $0$, its matrix in the basis
$\mathcal B$ has a non-zero entry from $b$ to $b'$ only if $b'\in\mathcal B_1$,
$b\in\mathcal B_2\sqcup\mathcal B_3$, or $b'\in\mathcal B_2$, $b\in\mathcal B_3$; conjugation by
$a^{-1}$ multiplies this entry by $\chi_b(a)\chi_{b'}(a)^{-1}\le\theta\max(\theta,\Theta)$, by
\eqref{pm:eq-contract}. So $\|a^{-1}n'a-1\|\le C\beta(a)^{-1}$, and $k$ is unitary.

(c) Let $U_x=\{u\in G_\infty:\|u_\tau-1\|\le r_0\|a\|^{-2}\}\times K''$ with $r_0\le1$ and a
compact open subgroup $K''$. Then $xU_xU_x^{-1}x^{-1}$ is contained in $\bigcup_{\omega}\omega U_0U_0^{-1}\omega^{-1}\times
\bigcup_{y}yK''y^{-1}$, the unions over the infinite components $\omega$ of the elements of $\Omega$
and over $y\in\Omega'_f$, where $U_0=\{\|u_\tau-1\|\le r_0\}$, because
$a\,u\,a^{-1}-1=a(u-1)a^{-1}$ has norm at most $\|a\|^2\|u-1\|$. These sets shrink to $\{1\}$ with
$r_0$ and $K''$, and $\pmGa$ is discrete; so for small $r_0$, $K''$ the sets $g^{-1}\gamma xU_x$,
$\gamma\in\pmGa$, are disjoint for all $g$ and all $x\in\pmS$. As $\pmP(yu)\le2^f\pmP(y)$ for
$u\in U_x$, we get $M_{B'}(y)\le2^{fB'}M'(yu)$ with $M'=\pmP^{-B'}\otimes1_{\Omega_fK''}$, and
\[
\sum_\gamma M_{B'}(g^{-1}\gamma x)\le\frac{2^{fB'}}{\pmvol(U_x)}\sum_\gamma\int_{U_x}M'(g^{-1}\gamma
xu)\,du\le\frac{2^{fB'}}{\pmvol(U_x)}\int_{G(\auA_F)}M'.
\]
The integral is finite by Lemma~\ref{pm:kernel}(c), and $\pmvol(U_x)\ge c(r_0\|a\|^{-2})^{fn^2}$,
since $\dim G_\infty=fn^2$. The same computation, with $U_x$ replaced by a fixed small neighbourhood of $1$ and for the sum
over the complement of a finite set of $\gamma$, shows that the series converges uniformly for
$g\in G(\auA_F)$ and $x$ in any compact subset of $G(\auA_F)$.
\end{proof}

\begin{lemma}[Decay in a cusp]\label{pm:decay}
For every integer $j\ge0$ there is $\varepsilon_0\in(0,1]$ with the following property. Let $N$
be the unipotent radical of a proper parabolic $F$-subgroup of $G$, $x\in G(\auA_F)$,
$K'\subset G_f$ compact open and $\delta>0$ with $N(\auA_F)x\subset N(F)\,x\,(Z_\delta\times K')$.
Let $\mathsf F$ be a continuous function on $G(\auA_F)$ with $|\mathsf F|\le T$, left invariant
under $N(F)$, right invariant under $K'$, and with $\mathsf F(g)=\int_{G_\infty}\mathsf F(gh)e(h)\,dh$
for all $g$. Then $|\mathsf F(x)|\le(\delta/\varepsilon_0)^j\,T$.
\end{lemma}
\begin{proof}
\emph{The constant term vanishes.} Let $\mathsf F_N(g)=\int_{N(F)\backslash N(\auA_F)}\mathsf F(ng)\,
dn$ (mass $1$). It is bounded, left invariant under $N(\auA_F)$, and again
$\mathsf F_N(g)=\int\mathsf F_N(gh)e(h)dh$. Fix $\tau$ and write $h=h^\tau h_\tau$. For fixed $g$,
$h^\tau$ the function $\Psi(h_\tau)=\mathsf F_N(gh^\tau h_\tau)$ on $G_\tau$ is left invariant under
$N'=g_\tau^{-1}N(F_\tau)g_\tau$, the unipotent radical of a proper parabolic subgroup of $G_\tau$, so
$\int_{G_\tau}\Psi(h_\tau)e_\tau(h_\tau)dh_\tau=\int_{N'\backslash G_\tau}\Psi(h)\int_{N'}e_\tau(nh)\,dn\,dh
=0$ by Lemma~\ref{pm:cusp}. Hence $\mathsf F_N=0$.

\emph{Oscillation.} By Lemma~\ref{pm:kernel}(b), $\mathsf F(g)=\int\mathsf F(gh)\alpha(h)\,dh$.
Let $\Omega_\alpha=Z_1\pmsupp\alpha$, $c_\alpha=\sup_{\Omega_\alpha}\|h\|$, and
$\omega(\varepsilon)=\sup_{z\in Z_\varepsilon}\int|\alpha(z^{-1}h)-\alpha(h)|dh$, which tends to $0$
with $\varepsilon$. For $z\in Z_\varepsilon$, $\varepsilon\le1$,
\begin{equation}\label{pm:eq-osc}
|\mathsf F(gz)-\mathsf F(g)|=\Bigl|\int\mathsf F(gh)\bigl(\alpha(z^{-1}h)-\alpha(h)\bigr)dh\Bigr|
\le\omega(\varepsilon)\sup_{h\in\Omega_\alpha}|\mathsf F(gh)| .
\end{equation}
Choose $\varepsilon_0$ with $\omega(\varepsilon_0)c_\alpha^{2j}\le\tfrac12$. Let
$m(h)=|\mathsf F(xh)|$ for $h\in G_\infty$ and $C^*=\sup_hm(h)\|h\|^{-2j}\le T$. For $h\in G_\infty$,
$N(\auA_F)xh\subset N(F)\,xh\,(Z_{\delta\|h\|^2}\times K')$, as $h^{-1}Z_\delta h\subset
Z_{\delta\|h\|^2}$. If $\delta\|h\|^2\le\varepsilon_0$, then, since $\mathsf F_N=0$ and by
\eqref{pm:eq-osc},
\[
m(h)=|\mathsf F(xh)-\mathsf F_N(xh)|\le\sup_{z\in Z_{\delta\|h\|^2}}|\mathsf F(xhz)-\mathsf F(xh)|
\le\omega(\varepsilon_0)\sup_{h'\in\Omega_\alpha}m(hh')\le\tfrac12C^*\|h\|^{2j}.
\]
If $\delta\|h\|^2>\varepsilon_0$, then $m(h)\le T\le T(\delta/\varepsilon_0)^j\|h\|^{2j}$. Therefore
$C^*\le\max(\tfrac12C^*,T(\delta/\varepsilon_0)^j)$, that is, $C^*\le T(\delta/\varepsilon_0)^j$, and
$|\mathsf F(x)|=m(1)\le C^*$.
\end{proof}

A subset $S\subset\pmGa$ is \emph{admissible} if it is stable under conjugation by $\pmGa$ and
$N(F)\cdot(S\cap P(F))=S\cap P(F)$ for every proper parabolic $F$-subgroup $P$ with unipotent
radical $N$. For a set $\mathcal R$ of polynomials, $\pmGa_{\mathcal R}=\{\gamma:\pmch(\gamma)\in
\mathcal R\}$ is admissible, because $\gamma\in P(F)$ and $\nu\gamma$, $\nu\in N(F)$, induce the
same maps on the graded pieces of the flag stabilised by $P$. Let $\pmGa_\pmell$ be the set of
the $\gamma$ that stabilise no non-zero isotropic subspace; it contains every $\gamma$ with
irreducible $\pmch(\gamma)$.

\begin{proposition}\label{pm:bound}
Let $\varphi'\in C_c^\infty(G_f)$ and $\mathsf f'=e\otimes\varphi'$.
\begin{enumerate}
\item[(a)] The series $k_{\mathsf f'}(x,y)=\sum_{\gamma\in\pmGa}\mathsf f'(x^{-1}\gamma y)$ converges absolutely and
uniformly on compact sets, and is continuous on $\pmQ\times\pmQ$.
\item[(b)] $\sup_{x\in G(\auA_F)}\sum_{\gamma\in\pmGa_\pmell}|\mathsf f'(x^{-1}\gamma x)|<\infty$.
\item[(c)] There is $C$ such that $\bigl|\sum_{\gamma\in S}\mathsf f'(x^{-1}\gamma x)\bigr|\le C$ for all
$x\in G(\auA_F)$ and all admissible $S$.
\end{enumerate}
\end{proposition}
\begin{proof}
Let $\Omega_f=\pmsupp\varphi'$ and let $K'$ be a compact open subgroup with
$\varphi'(k^{-1}y)=\varphi'(y)$ for $k\in K'$. By Lemma~\ref{pm:kernel}, $|\mathsf f'|\le CM_B$, and
$M_B\le c^{-j}\beta^{-fj}M_{B-j}$ on the set where $\pmP(y_\infty)\ge c\beta^f$. (a) follows from
Lemma~\ref{pm:siegel}(c). The sums in (b), (c) are left invariant under $\pmGa$ as functions of
$x$, so we may take $x=\omega ak\in\pmS$. If $\beta(a)\le1$, then $x$
is in a compact set and the bound is (a). Let $\beta(a)>1$ and $P=P_x$, $N=N_x$. Take
$j=n^3$; then $B-j>n^2$, and by Lemma~\ref{pm:siegel}(a),(c)
\begin{equation}\label{pm:eq-far}
\sum_{\gamma\in\pmGa\smallsetminus P(F)}|\mathsf f'(x^{-1}\gamma x)|\le C\beta(a)^{-fn^3}\sum_\gamma M_{B-j}
(x^{-1}\gamma x)\le C'\beta(a)^{-fn^3}\|a\|^{2fn^2}\le C'' ,
\end{equation}
as $\|a\|\le t^{-n}\beta(a)^{n/2}$. This proves (b), since $\pmGa_\pmell\cap P(F)=\emptyset$. For
(c) it remains to bound $\mathsf F(x)$, where $\mathsf F(g)=\sum_{\gamma\in S\cap P(F)}\mathsf f'(g^{-1}\gamma
x)$. By Lemma~\ref{pm:siegel}(c), $\mathsf F$ is continuous and $|\mathsf F|\le T=C\|a\|^{2fn^2}\le
C_1\beta(a)^{fn^3}$. It is right invariant under $K'$, left invariant under $N(F)$ because $S$ is
admissible, and $\int\mathsf F(gh)e(h)dh=\mathsf F(g)$ because $e*e=e$ (the exchange of sum and
integral is allowed by the bound for $\sum|\mathsf f'|$). By Lemma~\ref{pm:siegel}(b) and
Lemma~\ref{pm:decay} with $j=fn^3$, $|\mathsf F(x)|\le(C\beta(a)^{-1}/\varepsilon_0)^{fn^3}C_1
\beta(a)^{fn^3}$, a constant.
\end{proof}

\subsubsection{Forms, the trace, and the terms of the characteristic polynomials}\label{pm:ss-trace}

\begin{proposition}\label{pm:trace}
\begin{enumerate}
\item[(a)] For every compact open $K\subset G_f$ the space $\pmA(V;K)$ is finite-dimensional, and
its elements are bounded continuous functions.
\item[(b)] For $\varphi\in C_c^\infty(G_f)$ and $\mathsf f=e\otimes\varphi$, $R(\mathsf f)$ has finite rank and
$\pmtr R(\mathsf f)=\int_{\pmQ}k_{\mathsf f}(x,x)\,dx$. If $\varphi$ is bi-invariant under $K$, then
$\pmtr R(\mathsf f)=\pmtr(R(\varphi)\mid\pmA(V;K))$.
\item[(c)] For a polynomial $R$ put $J_R(\varphi)=J_R^V(\varphi)=\int_{\pmQ}\sum_{\pmch(\gamma)=R}
\mathsf f(x^{-1}\gamma x)\,dx$. This is defined for every $\varphi\in C_c^\infty(G_f)$; $J_R$ is linear,
$J_R(\varphi^g)=J_R(\varphi)$ for $g\in G_f$ and $\varphi^g(y)=\varphi(g^{-1}yg)$; for split $v$,
$J_R(\varphi)=0$ if $\varphi$ vanishes on $\{y\in G_f:\pmch(y_v)=R\}$; and
$\pmtr R(\mathsf f)=\sum_RJ_R(\varphi)$, the sum over all polynomials being unconditionally convergent.
\end{enumerate}
\end{proposition}
\begin{proof}
(a) Let $\mathsf f_0=\pmvol(K)^{-1}e\otimes1_K$. Then $\mathsf f_0^*=\mathsf f_0=\mathsf f_0*\mathsf f_0$, and $R(\mathsf f_0)$ is the orthogonal
projection onto $\pmA(V;K)$. For $\psi\in C_c(\pmQ)$, unfolding gives $R(\mathsf f_0)\psi(x)=\int_{\pmQ}
k_{\mathsf f_0}(x,y)\psi(y)\,dy$, and $\int_{\pmQ}|k_{\mathsf f_0}(x,y)|^2dy=k_{\mathsf f_0*\mathsf f_0^*}(x,x)=k_{\mathsf f_0}(x,x)\le C$ by
Proposition~\ref{pm:bound}(c) with $S=\pmGa$. So $|R(\mathsf f_0)\psi(x)|\le C^{1/2}\|\psi\|_2$; by
continuity every $\psi\in\pmA(V;K)$ is represented by $x\mapsto\langle\psi,\aubar{k_{\mathsf f_0}(x,\cdot)}
\rangle$, a continuous function bounded by $C^{1/2}\|\psi\|_2$. For orthonormal $\psi_1,\dots,\psi_d
\in\pmA(V;K)$, Bessel's inequality gives $\sum_i|\psi_i(x)|^2\le k_{\mathsf f_0}(x,x)\le C$, so
$d\le C\pmvol(\pmQ)$.

(b) Let $\varphi$ be invariant under $K$ on both sides (shrink $K$). Then $R(\mathsf f)=R(\mathsf f)R(\mathsf f_0)=R(\mathsf f_0)
R(\mathsf f)$, so $R(\mathsf f)$ has finite rank, with image and coimage in $\pmA(V;K)$; with an orthonormal
basis $(\psi_i)$ of $\pmA(V;K)$, the function $\sum_i(R(\mathsf f)\psi_i)(x)\aubar{\psi_i(y)}$ is a
continuous kernel of $R(\mathsf f)$. So is $k_{\mathsf f}$ (Proposition~\ref{pm:bound}(a) and unfolding). Two
continuous kernels of the same operator are equal; on the diagonal, $\int k_{\mathsf f}(x,x)dx=\sum_i\langle
R(\mathsf f)\psi_i,\psi_i\rangle=\pmtr R(\mathsf f)$. The last assertion follows from $R(\mathsf f)=R(\varphi)R(e)$, as
$R(e)$ and $R(\varphi)$ commute and $R(\varphi)$ preserves $L^2(\pmQ)^K$, on whose orthogonal complement
it vanishes.

(c) The integrand of $J_R$ is bounded by Proposition~\ref{pm:bound}(c), and $\pmQ$ has finite
volume. Invariance is the substitution $x\mapsto xg$, and the support condition holds because
$\pmch((x^{-1}\gamma x)_v)=\pmch(\gamma)$. Finally $k_{\mathsf f}(x,x)$ is the sum over $R$ of the integrands,
the partial sums over finite sets $\mathcal R$ of polynomials are bounded by the constant of
Proposition~\ref{pm:bound}(c), and dominated convergence applies.
\end{proof}

\subsubsection{Polynomials that are not irreducible}\label{pm:ss-red}

For a finite split place $u$ of $F$ let $Z_u\subset H_u\cong\auGL_n(F_u)$ be the set of matrices
with entries in $\auO_u$ whose reduction has characteristic polynomial with a \emph{simple} root
$0$. A test function $\varphi$ \emph{satisfies $(Z)$ at $u$} if $y_u\in Z_u$ for all
$y\in\pmsupp\varphi$.

\begin{proposition}\label{pm:red}
Let $\varphi=\varphi^{v_0}\otimes\beta$ with $\beta$ a supercusp function on $H_{v_0}$.
\begin{enumerate}
\item[(a)] If $R=AB$ with $A,B\in E[X]$ monic, coprime and not constant, then $J_R(\varphi)=0$.
\item[(b)] If $\varphi$ satisfies $(Z)$ at some $u$ and $R=Q^j$ with $j\ge2$, then $J_R(\varphi)=0$.
\end{enumerate}
Hence, if $\varphi$ satisfies $(Z)$ at some $u$, then $J_R(\varphi)=0$ for every monic $R$ that
is not irreducible.
\end{proposition}
\begin{proof}
(b) If $g\in Z_u$ had $\pmch(g)=Q^j$ over $F_u$, then $Q$ would be monic with coefficients in
$\auO_u$ (Gauss) and every root of the reduction of $\pmch(g)$ would have multiplicity divisible
by $j$. So $\varphi$ vanishes on $\{y:\pmch(y_u)=R\}$, and Proposition~\ref{pm:trace}(c) applies.

(a) Write $H=H_{v_0}$, $\mathsf C=\{h\in H:\pmch(h)=AB\}$, a closed subset, and
$D(\beta')=J_R(\varphi^{v_0}\otimes\beta')$ for $\beta'\in C_c^\infty(H)$. By
Proposition~\ref{pm:trace}(c), $D$ is invariant under conjugation by $H$ and vanishes on the
functions that vanish on $\mathsf C$; so it is an invariant distribution on $\mathsf C$
(restriction $C_c^\infty(H)\to C_c^\infty(\mathsf C)$ is surjective~\cite[\S1.8, Proposition]{BZ76}). We show
that every such distribution kills supercusp functions. Let $d=\deg A$, $\epsilon$ the projector
of $F_{v_0}^n$ onto the span of the first $d$ basis vectors, $M\cong\auGL_d\times\auGL_{n-d}$ its
centraliser in $H$, $N$ the group of the $1+\epsilon z(1-\epsilon)$, $P=MN$, $K_0=\auGL_n(
\auO_{v_0})$, and $\mathsf X=\{(m_1,m_2)\in M:\pmch(m_1)=A,\ \pmch(m_2)=B\}$. Choose
$a,b\in F_{v_0}[X]$ with $aA+bB=1$. For $h\in\mathsf C$, $(bB)(h)$ is the projector onto
$\ker A(h)$ along $\ker B(h)$, of rank $d$, and it depends continuously on $h$. Hence the map
$\varpi\colon H\times\mathsf X\to\mathsf C$, $(g,m)\mapsto gmg^{-1}$, is surjective, its fibres are
the orbits of $M$ acting by $(g,m)\mu=(g\mu,\mu^{-1}m\mu)$, and it has continuous local sections
(because $H\to H/M$, $g\mapsto g\epsilon g^{-1}$, has: $H/M$ is homeomorphic to the orbit of
$\epsilon$ under conjugation \cite[\S1.6]{BZ76}, and for a projector $e$ in that orbit near
$\epsilon$ the element $s(e)=e\epsilon+(1-e)(1-\epsilon)$ is invertible, depends continuously on
$e$, and satisfies $s(e)\epsilon=e\epsilon=e\,s(e)$, so $s(e)\epsilon s(e)^{-1}=e$; translating by $H$
gives local sections everywhere). For
$\Psi\in C_c^\infty(H\times\mathsf X)$ put $\varpi_*\Psi(gmg^{-1})=\int_M\Psi(g\mu,\mu^{-1}m\mu)d\mu$,
an element of $C_c^\infty(\mathsf C)$. The functional $\Psi\mapsto D(\varpi_*\Psi)$ is invariant
under left translations of the first variable; as $C_c^\infty(H\times\mathsf X)=C_c^\infty(H)\otimes
C_c^\infty(\mathsf X)$ and a left-invariant functional on $C_c^\infty(H)$ is a multiple of the
Haar integral, it is of the form $\Psi\mapsto\nu\bigl(m\mapsto\int_H\Psi(g,m)dg\bigr)$. Thus
\begin{equation}\label{pm:eq-fibre}
\int_H\Psi(g,m)\,dg=0\ \text{ for all }m\in\mathsf X\quad\Longrightarrow\quad D(\varpi_*\Psi)=0 .
\end{equation}
Now let $\beta$ be a supercusp function. Replacing it by its mean over the conjugates by $K_0$
changes neither $D(\beta)$ nor the supercusp property, so let $\beta$ be invariant under
conjugation by $K_0$, and put $\Psi(g,m)=1_{K_0N}(g)\,\beta(gmg^{-1})$. For $g=k\nu$ with $k\in K_0$,
$\nu\in N$, $\Psi(g,m)=\beta(\nu m\nu^{-1})=\beta(m\cdot\nu_m)$, $\nu_m=m^{-1}\nu m\nu^{-1}$. For
$m\in\mathsf X$ the map $\nu\mapsto\nu_m$ is a homeomorphism of $N$ which multiplies the Haar
measure by a constant, and it depends continuously on $m$: in the coordinate $z$ it is
$z\mapsto m_1^{-1}zm_2-z$, an invertible linear map because $m_1$, $m_2$ have no common eigenvalue.
Hence $\Psi$ is locally constant with compact support ($K_0N$ is open and closed, and on the
support of $\Psi$ the element $m\nu_m$ stays in the compact set $\pmsupp\beta$, so $m$, $\nu_m$,
and $\nu$ modulo $N\cap K_0$, stay in compact sets), and $\int_H\Psi(g,m)dg$ is a constant times
$\int_N\beta(mn)\,dn=0$. Moreover $\varpi_*\Psi=\pmvol(M\cap K_0)\cdot\beta|_{\mathsf C}$: for
$g=k\nu\mu_0$ (Iwasawa decomposition $H=K_0NM$) the set of $\mu\in M$ with $g\mu\in K_0N$ is
$\mu_0^{-1}(M\cap K_0)$, because $P\cap K_0N=(M\cap K_0)N$. By \eqref{pm:eq-fibre}, $D(\beta)=0$.

For the last statement, a monic $R$ that is neither irreducible nor a proper power has two
distinct monic irreducible factors, and is of the form $AB$ as in (a).
\end{proof}

\subsubsection{Irreducible polynomials}\label{pm:ss-ell}

Let $R\in E[X]$ be monic and irreducible of degree $n$ with $\pmGa_R=\{\gamma\in\pmGa:
\pmch(\gamma)=R\}\ne\emptyset$ (for $V$ or for $V'$). Let $L=E[X]/(R)$, $y$ the class of $X$. For
$\gamma\in\pmGa_R$ the adjoint involution of $V$ induces on $E[\gamma]\cong L$ the involution
$*$ with $y^*=y^{-1}$ extending the conjugation of $E$; it depends only on $R$. Let $L^+$ be its
fixed field, $\pmN(z)=zz^*$ the norm of $L/L^+$, and $T_R$ the torus over $F$ with
$T_R(F)=\{z\in L^\times:\pmN(z)=1\}$; it is anisotropic, so $T_R(F)\backslash T_R(\auA_F)$ is
compact. The centraliser of $\gamma$ in $G$ is $T_R$, through $L\cong E[\gamma]$.

\emph{The invariant.} $V$ is a line over $E[\gamma]$. Choosing $v_0\ne0$ we get
$h(a(\gamma)v_0,b(\gamma)v_0)=\pmTr_{L/E}(ab^*c)$ for a unique $c\in L^{+\times}$, and changing
$v_0$ multiplies $c$ by a norm: $\operatorname{inv}(\gamma)\in L^{+\times}/\pmN(L^\times)$ is well
defined. Write $V_c$ for the space $L$ with the form $\pmTr_{L/E}(ab^*c)$, on which $y$ is an
isometry. Two elements $\gamma_1,\gamma_2\in\pmGa_R$ are conjugate in $\pmGa$ if and only if
$\operatorname{inv}(\gamma_1)=\operatorname{inv}(\gamma_2)$: an isometry $g$ of $V$ with
$g\gamma_1g^{-1}=\gamma_2$ is the same as $z\in L^\times$ with $c_1=c_2\pmN(z)$, by
$g(a(\gamma_1)v_1)=a(\gamma_2)zv_2$. The same holds over every completion $F_w$ of $F$, with
$L\otimes_FF_w$ in place of $L$, and over $\auA_{F,f}$ if $z$ is taken in $(L\otimes\auA_{F,f})^\times$.

\emph{Signs.} Assume that for every $\tau$ all roots of $\tau(R)$ have absolute value $1$. Then
$*$ is the complex conjugation in every embedding, $L$ is a CM field and $L^+$ is totally real;
its real places $u$ correspond to the pairs $(\tau,z)$ with $z$ a root of $\tau(R)$, and
$h_\tau$ has the sign of $u(c)$ on the $z$-eigenline of $\gamma$ in $V_\tau$. We write
$\varepsilon(\gamma)=(\operatorname{sgn}u(c))_u$ and let $\pmE_q$ be the set of the families
$\varepsilon=(\varepsilon_u)\in\{\pm1\}^{\{u\}}$ with exactly $q$ entries $-1$ above each $\tau$.
Thus $\varepsilon(\gamma)\in\pmE_q$ for $\gamma\in\pmGa_R$.

\begin{lemma}\label{pm:classes}
Let $f$ be even, $R$ as above with all roots of all $\tau(R)$ of absolute value $1$, and let
$X_R(V)$ be the set of $\pmGa$-conjugacy classes in $\pmGa_R$. The map
\[
X_R(V)\longrightarrow\{\text{conjugacy classes of }G_f\}\times\pmE_q,\qquad
[\gamma]\mapsto\bigl([\gamma]_{G_f},\varepsilon(\gamma)\bigr)
\]
is injective, and its image is $\mathsf C_R\times\pmE_q$ for a set $\mathsf C_R$ of conjugacy
classes of $G_f$. The set $\mathsf C_R$ is the same for $V$ and for $V'$ (with $G_{V'}(\auA_{F,f})$
identified with $G_f$ by $\iota$).
\end{lemma}
\begin{proof}
\emph{Class field theory.} (i) An element of $L^{+\times}$ that is a norm from $L\otimes_{L^+}L^+_u$
at every place $u$ of $L^+$ is a norm from $L$ (Hasse's norm theorem for the quadratic extension
$L/L^+$, \cite[Chapter~VI, Corollary~(4.5)]{Neukirch}). (ii) If $\eta_u\in\{\pm1\}$ is given at every real place of
$L^+$ with $\prod_u\eta_u=1$, there are $x\in L^{+\times}$ and $z\in(L\otimes\auA_{\Q,f})^\times$
with $\operatorname{sgn}u(x)=\eta_u$ for all real $u$ and $x\,\pmN(z)=1$ in $L^+\otimes\auA_{\Q,f}$:
the idele of $L^+$ with the components $\eta_u$ at the real places and $1$ elsewhere has trivial
Artin symbol for $L/L^+$, so it is the product of an element $x$ of $L^{+\times}$ and a norm
$\pmN(z')$ of an idele of $L$ (the kernel of the global norm residue symbol is $L^{+\times}\pmN(\auA_L^\times)$, \cite[Chapter~VI, Theorem~(5.5)]{Neukirch}); take for $z$ the finite part of $z'$; note that $\pmN(z'_u)>0$ at the real places. (iii) Two hermitian
spaces over $E$ that are isometric at all places of $F$ are isometric
(Landherr~\cite{Landherr}; \cite[Chapter~10, Theorem~6.1]{Scharlau}).

\emph{Injectivity.} Let $\gamma_1,\gamma_2\in\pmGa_R$ have the same image, with invariants
$c_1,c_2$. Then $c_1/c_2$ is positive at all real places and, being $\pmN(z)$ for some
$z\in(L\otimes\auA_{F,f})^\times$, a local norm at all finite places. By (i) it is a norm, and
$\gamma_1,\gamma_2$ are conjugate.

\emph{Change of signs and of space.} Let $\gamma\in\pmGa_R$ with invariant $c$, and
$\varepsilon'\in\pmE_{q'}$. We claim: there is $\gamma'\in G_{V'}(F)$ with $\pmch(\gamma')=R$,
$\varepsilon(\gamma')=\varepsilon'$, and $\gamma'$ conjugate in $G_f$ to $\iota\gamma\iota^{-1}$.
Indeed $\prod_u\varepsilon(\gamma)_u\varepsilon'_u=(-1)^{fq}(-1)^{fq'}=1$, as $f$ is even; let
$x$, $z$ be as in (ii) for $\eta=\varepsilon(\gamma)\varepsilon'$. The space $V_{cx}$ has the
signature of $V'$ at every $\tau$, and multiplication by $z$ is an isometry $V_c\otimes\auA_{F,f}
\to V_{cx}\otimes\auA_{F,f}$ commuting with $y$, because $\pmN(z)x=1$; so $V_{cx}$ is isometric to $V_c\cong V$, hence to
$V'$, at every finite place. By (iii) there is an isometry $j\colon V_{cx}\to V'$ over $E$, and
$\gamma'=jyj^{-1}$ has the three properties: the element $j\circ z\circ(\text{the isometry }V\cong
V_c)\circ\iota^{-1}$ of $G_{V'}(\auA_{F,f})$ conjugates $\iota\gamma\iota^{-1}$ into $\gamma'$.

Applied to $V'=V$, $\iota=1$, the claim shows that the image of the map is a product
$\mathsf C_R(V)\times\pmE_q$. Applied to $(V,V')$ and to $(V',V,\iota^{-1})$ it shows
$\mathsf C_R(V)=\mathsf C_R(V')$.
\end{proof}

\begin{lemma}[Selberg's principle]\label{pm:selberg}
Let $R$ be irreducible and suppose that some $\tau(R)$ has a root of absolute value $\ne1$. Then
$J_R(\varphi)=0$ for all $\varphi\in C_c^\infty(G_f)$.
\end{lemma}
\begin{proof}
By Proposition~\ref{pm:bound}(b), sum and integral in $J_R(|\varphi|)$ (with $|e|$) converge, and
unfolding gives, with $\mathsf f=e\otimes\varphi$, $J_R(\varphi)=\sum_{[\gamma]\in X_R(V)}\int_{T_R(F)\backslash G(\auA_F)}
\mathsf f(x^{-1}\gamma x)dx$, each integral being absolutely convergent. Fix $\gamma$. The integral is
$\pmvol(T_R(F)\backslash T_R(\auA_F))$ times the product of $\int_{T_R(\auA_{F,f})\backslash G_f}
\varphi(x^{-1}\gamma x)dx$ and of the integrals $I_{\tau'}=\int_{T_R(F_{\tau'})\backslash G_{\tau'}}
e_{\tau'}(x^{-1}\gamma x)\,dx$, and the same holds with absolute values, in $[0,\infty]$; so either one of the factors with
absolute values is $0$, and then the term of $\gamma$ is $0$, or every $I_{\tau'}$ converges
absolutely. In the second case we show $I_\tau=0$. The eigenvalues of
$\gamma$ on $V_\tau$ are distinct, and eigenvectors for $z,z'$ are orthogonal unless
$z\aubar{z'}=1$. So the span $V_>$ of the eigenlines with $|z|>1$ is isotropic and not zero, its
stabiliser $P_1=M_1N_1$ is a proper parabolic subgroup of $G_\tau$, where we take for $M_1$ the
stabiliser of the decomposition $V_\tau=V_>\oplus V_=\oplus V_<$ into the spans of the eigenlines
with $|z|>1$, $|z|=1$, $|z|<1$; this is a Levi factor of $P_1$, since $V_>^\perp=V_>\oplus V_=$ and
the hermitian form puts $V_<$ in duality with $V_>$. The element $\gamma$ and its centraliser
$T_R(F_\tau)$ preserve every eigenline of $\gamma$, hence lie in $M_1$; and $\pmAd(\gamma)$ has no eigenvalue $1$ on $\pmn_1$ (its eigenvalues
are quotients $z/z'$ of eigenvalues of $\gamma$ with $|z|>|z'|$). By the Iwasawa decomposition $G_\tau=M_1N_1K_\tau$,
\[
I_\tau=\int_{T_R(F_\tau)\backslash M_1}\int_{K_\tau}\int_{N_1}e_\tau\bigl(k^{-1}n^{-1}\gamma_mn\,k\bigr)
\,dn\,dk\,\delta(m)dm,\qquad\gamma_m=m^{-1}\gamma m,
\]
and $n\mapsto\gamma_m^{-1}n^{-1}\gamma_mn$ is a bijection of $N_1$ onto itself that multiplies the
Haar measure by $|\det(\pmAd(\gamma_m^{-1})-1)|_{\pmn_1}|\ne0$ (induction on the descending
central series of $N_1$). So the inner integral is a multiple of $\int_{N_1}e_\tau(k^{-1}\gamma_m
nk)\,dn$, which is $0$ by Lemma~\ref{pm:cusp}.
\end{proof}

\begin{proposition}\label{pm:ell}
Let $f$ be even and $R$ monic irreducible of degree $n$. Then
$J_R^V(\varphi)=J_R^{V'}(\varphi)$ for every $\varphi\in C_c^\infty(G_f)$.
\end{proposition}
\begin{proof}
If $\pmGa_R$ and $\pmGa'_R$ are empty, both sides vanish. If some $\tau(R)$ has a root off the
unit circle, both vanish by Lemma~\ref{pm:selberg}. Otherwise, as in the proof of that lemma,
\[
J_R^V(\varphi)=v_R\sum_{[\gamma]\in X_R(V)}O_{[\gamma]_{G_f}}(\varphi)\prod_{\tau}\int_{G_\tau}
e_\tau(x^{-1}\gamma x)\,dx,
\]
where $v_R=\pmvol(T_R(F)\backslash T_R(\auA_F))$ and $O_\kappa(\varphi)=\int_{T_R(\auA_{F,f})
\backslash G_f}\varphi(x^{-1}\gamma_\kappa x)dx$ for $\gamma_\kappa$ in the class $\kappa$; the
measure of $T_R(\auA_{F,f})$ is fixed once and transported to the centraliser of $\gamma_\kappa$
through $L$ (this is compatible with conjugation), and $T_R(F\otimes\R)$, a product of circles,
has mass $1$. So $v_R$ and $O_\kappa$ are the same for $V$ and $V'$. In the identification
$G_\tau\cong G_s$ the element $\gamma$ is conjugate to $t^\sigma$, where
$t=\operatorname{diag}(z_1,\dots,z_n)$ lists the roots of $\tau(R)$ and $\sigma$ is any
permutation such that $s_i$ is the sign of $h_\tau$ on the eigenline of $z_{\sigma(i)}$; for each
$\varepsilon_\tau$ there are $p!q!$ such $\sigma$. Write $I_\tau(\varepsilon_\tau)$ for the integral over $G_\tau$ when $\gamma$ has the signs
$\varepsilon_\tau$ above $\tau$. By Lemma~\ref{pm:classes} the sum over $X_R(V)$ is a sum over
$\mathsf C_R\times\pmE_q$, absolutely convergent also with $|\varphi|$ and $|e_\tau|$, and by
Theorem~\ref{pm:arch}(b)
\[
\mathsf X:=\sum_{\varepsilon\in\pmE_q}\prod_\tau I_\tau(\varepsilon_\tau)=\prod_\tau\Bigl(\frac1{p!q!}
\sum_{\sigma\in S_n}\int_{G_s}e_\tau(x^{-1}t_\tau^\sigma x)\,dx\Bigr)=(-1)^{pqf}\prod_\tau
\aubar{\chi_{\lambda_\tau}(t_\tau)} .
\]
If the products $\prod_\tau I_\tau(\varepsilon_\tau)$, $\varepsilon\in\pmE_q$, do not all vanish,
then $\sum_{\kappa\in\mathsf C_R}O_\kappa(|\varphi|)<\infty$ and $J_R^V(\varphi)=v_R\bigl(\sum_\kappa
O_\kappa(\varphi)\bigr)\mathsf X$; if they all vanish, $J_R^V(\varphi)=0$ and $\mathsf X=0$. As $f$
is even, $\mathsf X=\prod_\tau\aubar{\chi_{\lambda_\tau}(t_\tau)}$ is the same number for $V$ and
for $V'$. If it is $0$, both $J_R^V(\varphi)$ and $J_R^{V'}(\varphi)$ vanish. If not, the first case
holds for both spaces, with the same $v_R$, $\mathsf C_R$ and $O_\kappa$.
\end{proof}

\begin{remark}\label{pm:rem-even}
The hypothesis that $f$ is even is used twice, in the proof of Lemma~\ref{pm:classes} and at the
end of the last proof, to make $(-1)^{f(q+q')}$ and $(-1)^{pqf}(-1)^{p'q'f}$ equal to $1$. Both hold
for every comparison datum: the quotient of the discriminants of $V$ and $V'$ is an element of
$F^\times$ that is a norm from $E$ at all finite places and has the sign $(-1)^{q-q'}$ at all
$f$ real places, so $f(q-q')$ is even by the product formula for the norm residue symbol of
$E/F$; and $pq\equiv(n-1)q\pmod 2$. We keep the hypothesis because it holds for the data.
\end{remark}

\begin{corollary}\label{pm:simple}
Theorem~\ref{pm:main} holds for the test functions $\varphi=\varphi^{v_0}\otimes\beta$ that
satisfy $(Z)$ at some finite split place $u$.
\end{corollary}
\begin{proof}
By Proposition~\ref{pm:trace}, $\pmtr R(e\otimes\varphi)=\sum_RJ_R^V(\varphi)$. The $R$ that are
not monic of degree $n$ do not occur; those that are not irreducible contribute $0$ by
Proposition~\ref{pm:red}; the others have $J^V_R(\varphi)=J^{V'}_R(\varphi)$ by
Proposition~\ref{pm:ell}.
\end{proof}

\subsubsection{The auxiliary place: end of the proof of Theorem~\ref{pm:main}}\label{pm:ss-aux}

\begin{lemma}\label{pm:line}
Let $\mathsf B$, $\mathsf B'$ be finite-dimensional Hilbert spaces with operators $(A,T,\epsilon)$
and $(A',T',\epsilon')$ such that $\epsilon^2=\epsilon$, $TT^*=d(1+c\,\epsilon)$ with real $d,c>0$,
$A$ commutes with $T$ and $T^*$, and likewise for the primed operators with the same $d$, $c$. If
$\pmtr(A^lT^i)=\pmtr(A'^lT'^i)$ for all $l,i\ge1$, then $\pmtr(A\epsilon)=\pmtr(A'\epsilon')$.
\end{lemma}
\begin{proof}
$T$ is invertible, and $\epsilon$ is a polynomial in $TT^*$, so $A$ commutes with $\epsilon$. For
$\alpha\in\C$ let $\mathsf B_\alpha$ be the generalised eigenspace of $A$; it is stable under $T$,
$T^*$, $\epsilon$, and $T^*|\mathsf B_\alpha$ is the adjoint of $T|\mathsf B_\alpha$. On
$\mathsf B_\alpha$, $A-\alpha$ is nilpotent and commutes with $T$, so $\pmtr(A^lT^i)=\sum_\alpha
\alpha^l\pmtr(T^i|\mathsf B_\alpha)$. As the sequences $(\alpha^l)_{l\ge1}$, $\alpha\ne0$, are
linearly independent, $\pmtr(T^i|\mathsf B_\alpha)=\pmtr(T'^i|\mathsf B'_\alpha)$ for all
$\alpha\ne0$, $i\ge1$. The power sums of positive index of the eigenvalues of an invertible
operator determine these eigenvalues with their multiplicities; hence
$\dim\mathsf B_\alpha=\dim\mathsf B'_\alpha$ and $|\det(T|\mathsf B_\alpha)|=|\det(T'|\mathsf B'_
\alpha)|$. Now $|\det(T|\mathsf B_\alpha)|^2=\det(TT^*|\mathsf B_\alpha)=d^{\dim\mathsf B_\alpha}
(1+c)^{\operatorname{rk}(\epsilon|\mathsf B_\alpha)}$, so $\epsilon$ and $\epsilon'$ have the same
rank on $\mathsf B_\alpha$ and $\mathsf B'_\alpha$, for $\alpha\ne0$. Finally
$\pmtr(A\epsilon)=\sum_\alpha\alpha\operatorname{rk}(\epsilon|\mathsf B_\alpha)$.
\end{proof}

\begin{lemma}\label{pm:hecke}
Let $u$ be a finite split place of $F$, $\mathbb F$ its residue field, $\mathsf q=\#\mathbb F$,
$K_u=\auGL_n(\auO_u)\subset H_u$, $J\subset K_u$ the subgroup of the matrices whose reduction maps
the first basis vector to a multiple of itself, and $t=\operatorname{diag}(\varpi,1,\dots,1)$ with a
uniformiser $\varpi$. Then
\begin{enumerate}
\item[(a)] $1_{JtJ}*1_{Jt^{-1}J}=\pmvol(J)\,\mathsf q^{n-2}\bigl(1_J+(\mathsf q-1)1_{K_u}\bigr)$;
\item[(b)] for $i\ge1$ the $i$-fold convolution power of $1_{JtJ}$ is supported in $Z_u$.
\end{enumerate}
\end{lemma}
\begin{proof}
(a) $tJt^{-1}$ is the subgroup $J'$ of the matrices in $K_u$ whose first row is
$(\ast,0,\dots,0)$ modulo $\varpi$. So $JtJt^{-1}J=JJ'J\subset K_u$; as $J$ is the preimage of a
maximal parabolic subgroup $\mathsf P$ of $\auGL_n(\mathbb F)$, $\mathsf P\backslash\auGL_n(
\mathbb F)/\mathsf P$ has two elements, and $J'\not\subset J$, we get $JJ'J=K_u$. The convolution
is bi-invariant under $J$, so it is $x_11_J+x_21_{K_u\smallsetminus J}$. Its value at $1$ is
$\pmvol(JtJ)=\pmvol(J)[J:J\cap J']=\pmvol(J)\mathsf q^{n-1}$, and its integral is
$\pmvol(JtJ)^2$; with $[K_u:J]=(\mathsf q^n-1)/(\mathsf q-1)$ this gives $x_1=\pmvol(J)\mathsf
q^{n-1}$ and $x_2=\pmvol(J)\mathsf q^{n-2}(\mathsf q-1)$. (b) The reduction of an element of $JtJ$
kills the first basis vector and induces an invertible map of $\mathbb F^n$ modulo that vector.
Matrices with entries in $\auO_u$ with this property are stable under products, and the
characteristic polynomial of their reduction is $X$ times a polynomial with
non-zero constant term.
\end{proof}

\begin{proof}[Proof of Theorem~\ref{pm:main}]
The finiteness of the rank and the last assertion are Proposition~\ref{pm:trace}(b). Choose a
compact open subgroup $K=\prod_wK_w$ of $G_f$ such that $\varphi$ is bi-invariant under $K$, and a
finite place $u\ne v_0$ of $F$, split in $E$, with $K_u=\auGL_n(\auO_u)$ in a suitable basis and
$\varphi=\varphi^u\otimes1_{K_u}$; almost all split places qualify. Let $J$, $t$ be as in
Lemma~\ref{pm:hecke} and $K^\flat=K^u\times J$. On $\mathsf B=\pmA(V;K^\flat)$, which is
finite-dimensional, consider
\[
A=\frac{R(\varphi^u\otimes1_J)}{\pmvol(K^u)\pmvol(J)},\qquad T=\frac{R(1_{JtJ})}{\pmvol(J)},\qquad
\epsilon=\frac{R(1_{K_u})}{\pmvol(K_u)},
\]
where functions on $H_u$ act through $R$ restricted to $H_u$. Then $\epsilon$ is the orthogonal
projection of $\mathsf B$ onto $\pmA(V;K)$, $T^*=R(1_{Jt^{-1}J})/\pmvol(J)$, $A$ commutes with $T$
and $T^*$, and by Lemma~\ref{pm:hecke}(a), $TT^*=\mathsf q^{n-2}(1+(\mathsf q^n-1)\epsilon)$ on
$\mathsf B$, since $R(1_J)=\pmvol(J)$ there and $\pmvol(K_u)=\pmvol(J)(\mathsf q^n-1)/(\mathsf q-1)$.
For $l,i\ge1$ the operator $A^lT^i$ is a positive multiple, depending only on $l,i$ and the
measures, of $R(\varphi_{l,i})$ on $\mathsf B$, with $\varphi_{l,i}=(\varphi^u)^{*l}\otimes(1_{JtJ})^{*i}$.
This test function is bi-invariant under $K^\flat$, satisfies $(Z)$ at $u$ by
Lemma~\ref{pm:hecke}(b), and is of the form $\varphi_{l,i}^{v_0}\otimes\beta^{*l}$; a convolution
power of a supercusp function is a supercusp function, because
$\int_N(\beta*\beta')(xny)dn=\int\beta'(h)\int_N\beta(xnyh^{-1})\,dn\,dh$. By
Corollary~\ref{pm:simple} and Proposition~\ref{pm:trace}(b), $\pmtr(A^lT^i|\mathsf B)=\pmtr(A'^l
T'^i|\mathsf B')$, where the primed objects are those of $V'$. Lemma~\ref{pm:line} gives
$\pmtr(A\epsilon|\mathsf B)=\pmtr(A'\epsilon'|\mathsf B')$. On the image $\pmA(V;K)$ of $\epsilon$,
$A=\pmvol(K)^{-1}R(\varphi)$. Hence $\pmtr(R(\varphi)|\pmA(V;K))=\pmtr(R'(\varphi)|\pmA(V';K))$.
\end{proof}

\subsubsection{Proof of Proposition~\ref{pm:mult}}\label{pm:ss-mult}

The following property of Green's parametrisation~\cite{Green1955} is the only one used; we
derive it below from the statements of~\cite{Green1955}.
\begin{enumerate}
\item[(G)] If the label $c$ over $\auF$ is an irreducible polynomial, the character $\Xi_c$ of
$\auGL_r(\auF)$ named by $c$ is cuspidal: the sum of $\Xi_c(xuy)$ over $u\in\mathsf U$ is zero for
all $x,y$ and the unipotent radical $\mathsf U$ of every proper parabolic subgroup.
\end{enumerate}
Green does not use the word cuspidal, so we derive (G) from statements of~\cite{Green1955}. For
$c$ irreducible of degree $r$, $\Xi_c$ is Green's primary character $(g^{(1)})$ of the $r$-simplex
$g$ attached to $c$ \cite[Lemma~7.7 and Thm.~12, pp.~438--439]{Green1955}. Let $\pi$ realise
$\Xi_c$. Since $\Xi_c$ is a class function, $\sum_u\Xi_c(xuy)=\sum_u\Xi_c(yxu)=\#\mathsf U\,
\operatorname{tr}(\pi(yx)E)$, with $E$ the projection onto $\pi^{\mathsf U}$; as $\pi$ is
irreducible, the $\pi(h)$ span $\operatorname{End}(\pi)$, so (G) says $\pi^{\mathsf U}=0$. A larger
parabolic has a smaller unipotent radical and every parabolic is conjugate to a standard one, so it
suffices to take $\mathsf U$ the radical of the standard parabolic with Levi
$\auGL_{r_1}\times\auGL_{r_2}$, $r_1+r_2=r$, $r_i\ge1$. If $\pi^{\mathsf U}\ne0$, it contains an
irreducible representation $\sigma_1\otimes\sigma_2$ of the Levi, and by Frobenius reciprocity
$\Xi_c$ occurs in the character induced from the parabolic, which is Green's $\circ$-product
$\chi_1\circ\chi_2$ of the characters $\chi_i$ of $\sigma_i$ \cite[p.~403]{Green1955}. By
\cite[Thm.~13, p.~439]{Green1955}, $\chi_i=\prod_{g'}(g'^{\nu_i(g')})$ ($\circ$-product) with $\nu_i$
supported on simplexes of degree $\le r_i<r$. The $\circ$-product is commutative and associative,
and for each simplex $g'$ of degree $s$ the map $P_e\mapsto Z_e$ is a ring homomorphism from the
symmetric functions into Green's character algebra sending the Schur function $\{\lambda\}$ to
$\pm(g'^{\lambda})$ \cite[Lemma~7.2, p.~436, Def.~7.3, p.~437, and p.~439]{Green1955}. Since a
product of Schur functions of degrees $a,b$ is a linear combination of Schur functions of degree
$a+b$, $(g'^{\nu_1(g')})\circ(g'^{\nu_2(g')})$ is a combination of the $(g'^{\lambda})$ with
$|\lambda|=|\nu_1(g')|+|\nu_2(g')|$. Hence $\chi_1\circ\chi_2$ is a linear combination of the
characters $(\cdots g'^{\nu(g')}\cdots)$ of Theorem~13 with $\nu$ supported on simplexes of degree
$<r$. These are irreducible and distinct from $\Xi_c=(g^{(1)})$, whose function is supported on
$g$ of degree $r$ \cite[Thm.~13 and p.~440]{Green1955}; by orthogonality of irreducible characters
$\langle\Xi_c,\chi_1\circ\chi_2\rangle=0$, a contradiction. So $\pi^{\mathsf U}=0$, which is (G).

\begin{lemma}\label{pm:supercusp}
Let $v$ be a finite split place of $F$ with residue field $\auF$, and $\Xi$ a cuspidal class
function on $\auGL_n(\auF)$ in the sense of \textup{(G)}. The function $\beta$ on $H_v$ that is
$\Xi(\bar k)$ at $k\in K_v=\auGL_n(\auO_v)$, $\bar k$ the reduction of $k$, and $0$ outside $K_v$, is
a supercusp function.
\end{lemma}
\begin{proof}
Every parabolic subgroup is conjugate under $K_v$ to a standard one, $P=MN$, and $H_v=K_vP=PK_v$.
For $x=kp$, $y=p'k'$ we have $pnp'=\mu n'$ with $\mu\in M$ depending only on $p,p'$, and $n\mapsto
n'$ multiplies the Haar measure of $N$ by a constant. Now $k\mu n'k'\in K_v$ if and only if
$\mu\in M\cap K_v$ and $n'\in N\cap K_v$. So $\int_N\beta(xny)\,dn$ is $0$ if $\mu\notin K_v$, and
otherwise a multiple of $\int_{N\cap K_v}\Xi(\bar k\bar\mu\bar n'\bar k')\,dn'$, a multiple of the
sum of $\Xi$ over a coset of the unipotent radical of a proper parabolic subgroup of
$\auGL_n(\auF)$.
\end{proof}

\begin{proof}[Proof of Proposition~\ref{pm:mult}]
Let $V=Y_0$ and $V'=W_j$, with the sign vectors \eqref{pm:eq-signs} and $\lambda_\tau=\varsigma_\tau$;
\eqref{pm:eq-gap} holds by hypothesis. (For the data: two entries on coordinates of different
signs are either two entries $-\kappa(\tau)-a\Delta$ with different $a$, or one of these and a
small entry; in both cases they differ by at least $\Delta$.) Both spaces carry a lattice
that is unimodular at every finite place, and $E/F$ is unramified at all finite places; hence
at every finite place $w$ of $F$ the two lattices are isometric (for $w$ inert this is the
uniqueness of unimodular hermitian lattices over an unramified extension, i.e.\ the existence of an orthonormal basis (proof: the reduction of the lattice is a non-degenerate hermitian space over the quadratic extension of residue fields, so it has a vector $\bar x$ with $\bar h(\bar x,\bar x)\ne0$ (otherwise polarisation and the surjectivity of the residue trace give $\bar h=0$); a lift $x$ has $h(x,x)$ a unit of the base ring, hence the norm of a unit since the extension is unramified, so after scaling $h(x,x)=1$; then the lattice is $\mathcal O x\perp(x^\perp\cap\text{lattice})$ with unimodular second summand, and one inducts on the rank; see also~\cite{Jacobowitz}); for
$w$ split it is trivial). Choosing such isometries gives a comparison datum $(Y_0,W_j,\iota)$
for which $\iota$ identifies the stabilisers $K^\circ$ of the lattices, the levels $J_M$, the
groups $\auG=K^\circ/J_M=\prod_w\auGL_r(\auF_w)$ and their actions $\varrho$, and the algebra
$\mathcal H^M$. Each chosen place $w$ lies above a place of $F$ above $p_0$ or $p_1$ that splits
in $E$, and distinct chosen places lie above distinct places of $F$.

Let $(c_w)$ be a good family of labels and $T=\varphi_S\otimes1_{J_M^{S}}\in\mathcal H^M$. For
$s\in\auG$ the idempotent $\bigotimes_we_{c_w}$ has the coefficient $\prod_w\Xi_{c_w}(1)\Xi_{c_w}
(s_w^{-1})/\#\auGL_r(\auF_w)$. Hence $\pmvol(J_M)^{-1}R(T)\circ\varrho(\bigotimes_we_{c_w})$ is
a constant multiple, the same for $Y_0$ and $W_j$, of $R(\varphi)$ on $\pmA(V;J_M)$, where
\[
\varphi=\varphi_S\otimes\bigotimes_{w\text{ chosen}}\beta_w\otimes1_{K^{\circ}}^{\,\text{other places}},
\qquad\beta_w(k)=\Xi_{c_w}(\bar k^{-1})\ \text{on }\auGL_r(\auO_w),
\]
a test function bi-invariant under $J_M$. As the family is good, some $c_{w_0}$ is irreducible;
by (G) and Lemma~\ref{pm:supercusp}, $\beta_{w_0}$ is a supercusp function (if $\Xi$ is cuspidal,
so is $g\mapsto\Xi(g^{-1})$). Theorem~\ref{pm:main} gives
$\pmtr(T\varrho(\bigotimes e_{c_w})\mid\pmA(Y_0;J_M))=\pmtr(T\varrho(\bigotimes e_{c_w})\mid
\pmA(W_j;J_M))$. Summing over the good families, and using that $\varrho(e_0)$ is a projection
commuting with $T$, we get $\pmtr(T\mid\auT)=\pmtr(T\mid e_0\pmA(W_j))$.

For $T=1$ this is the equality of dimensions. $\mathcal H^M$ is commutative and stable under
adjoints, so both spaces are direct sums of eigenspaces $\mathsf X_\chi$ for characters $\chi$ of
$\mathcal H^M$, and $\pmtr(T|\mathsf X)=\sum_\chi\chi(T)\dim\mathsf X_\chi$. Distinct characters
are linearly independent, so the multiplicities are equal.
\end{proof}

\begin{remark}[The full spherical algebra]\label{pm:rem-full}
The trace identity holds for every $T=\varphi_S\otimes1_{J_M^S}$ with $S$ any finite set of finite places
of $F$ not below a chosen place, split or inert, and $\varphi_S$ bi-invariant under $K^\circ_S$. Indeed the
comparison datum of the proof identifies, through $\iota$, the stabilisers $K^\circ_w$ at \emph{every}
finite place (for inert $w$ by~\cite{Jacobowitz}), hence the bi-$K^\circ_S$-invariant functions on the two
groups; and Theorem~\ref{pm:main} allows an arbitrary function $\varphi^{v_0}$ away from $v_0$. The first two
paragraphs of the proof therefore apply verbatim. (The last paragraph, on multiplicities, uses only the
split places, where $\mathcal H^M$ is commutative.) This is the form of {\cite[Proposition~14.36]{P46}}.
\end{remark}



% ===== B-orbital.tex =====
% proofs/B-orbital.tex
% Proof of Theorem pm:arch(b) of proofs/B-multiplicities.tex: the stable sum of the elliptic
% orbital integrals of the lowest-K-type coefficient of a discrete series of U(p,q).
% Labels of this file start with pbr:. Only amsmath/amssymb/amsthm are used.
\providecommand{\auU}{\mathrm U}
\providecommand{\aubar}[1]{\overline{#1}}
\providecommand{\pcom}{\mathsf m}
\providecommand{\pcoD}{\mathsf D}
\providecommand{\pcoH}{\mathsf H}
% macros of this file (prefix \pbr)
\providecommand{\pbrB}{\mathsf B}
\providecommand{\pbrN}{\mathsf N}
\providecommand{\pbrn}{\mathfrak n}
\providecommand{\pbrK}{K}
\providecommand{\pbrW}{W_K}
\providecommand{\pbrlev}{\ell}
\providecommand{\pbrsgn}{\operatorname{sgn}}
\providecommand{\pbrAlt}{\mathsf A}
\providecommand{\pbrTp}{T^{+}}
\providecommand{\pbrTh}{\Theta}
\providecommand{\pbrE}{\mathcal E}
\providecommand{\pbrdc}{\delta_{\mathrm c}}
\providecommand{\pbrdD}{\delta_{D}}
\providecommand{\pbrtop}{\tau}
\providecommand{\pbrLo}{\Lambda^{\circ}}
\providecommand{\pbrNo}{N^{\circ}}
\providecommand{\pbrO}{\mathcal O}
\providecommand{\pbrre}{\operatorname{Re}}

\subsection{Elliptic orbital integrals of the discrete series coefficient}\label{pbr:sec}

This section proves part (b) of Theorem~\ref{pm:arch}. Part (a) is Remark~\ref{pm:rem-other}(1).
The result is Harish-Chandra's formula for the elliptic orbital integrals of a discrete series
coefficient, summed over a stable class; we prove it for $\auU(p,q)$, for parameters satisfying the
gap condition \eqref{pm:eq-gap}, and for the coefficient of the lowest $K$-type, from the explicit
formula of {\cite[Theorem~3.17]{P49}}. \emph{Dependencies:} the proofs use the published results listed
in \S\ref{pbr:ss-inputs} (Schur orthogonality and admissibility, already used in
\S\ref{pm:sec}; the two statements of Schmid's theorem on the $K$-types of a discrete series, of
which the first is (D3)(b) of {\cite[\S3.1.3]{P49}}; and the classical theory of compact connected Lie
groups). Harish-Chandra's character theory is not used, and neither is his formula for the formal
degree.

\subsubsection{Statement, conventions and inputs}\label{pbr:ss-statement}

\paragraph*{The group and the parameter}
Let $n\ge2$, $s=(\varepsilon_1,\dots,\varepsilon_n)\in\{\pm1\}^n$ with $p$ entries $+1$ and $q$
entries $-1$, $J=I_s=\operatorname{diag}(s)$, $G=G_s=\{g:g^*Jg=J\}$, $K=G\cap\auU(n)$ (the block
diagonal elements of $G$), and $T\subset K$ the diagonal torus. We write $x\in T$ as
$x=\operatorname{diag}(x_1,\dots,x_n)$, and for $\sigma\in S_n$, $x^\sigma=\operatorname{diag}(
x_{\sigma(1)},\dots,x_{\sigma(n)})$, as in \S\ref{pm:ss-setting}. Let $\lambda\in(\frac{n-1}2+\Z)^n$
have distinct entries, and let
\[
h_\lambda=\min\{|\lambda_a-\lambda_b|:\varepsilon_a\ne\varepsilon_b\}
\]
be its gap ($h_\lambda=\infty$ if $pq=0$). Throughout, $(\lambda,s)$ satisfies \eqref{pm:eq-gap}:
\begin{equation}\label{pbr:eq-gap}
h_\lambda\ \ge\ h_0:=500\,n^3+10\,n^5 .
\end{equation}
As in {\cite[\S3.2.1]{P49}}: the positive roots are the $e_a-e_b$ with $\lambda_a>\lambda_b$, compact
if $\varepsilon_a=\varepsilon_b$ and noncompact otherwise; $\rho,\rho_{\mathrm c},\rho_{\mathrm n}$
are the half sums; $\Lambda=\lambda+\rho_{\mathrm n}-\rho_{\mathrm c}$ and $\mu=\lambda+\rho$, both in
$\Z^n$. A weight $\nu\in\Z^n$ is \emph{$K$-dominant} if $\nu_a\ge\nu_b$ whenever $\varepsilon_a=
\varepsilon_b$ and $\lambda_a>\lambda_b$; the irreducible representations of $K\cong\auU(p)\times
\auU(q)$ are the $(\sigma_\nu,V_\nu)$ with highest weight $\nu$ for the positive compact roots,
$\nu$ $K$-dominant, with characters $\chi^K_\nu$. By {\cite[Lemma~3.18(c)]{P49}}, $\Lambda$ is
$K$-dominant. We also put
\[
\pbrLo=\lambda-\rho ,
\]
the highest weight, for the order of $\lambda$, of the representation of $\auU(n)$ with character
$\chi_\lambda$ (\S\ref{pm:ss-setting}); $\pbrLo$ is dominant: $\pbrLo_a-\pbrLo_b=(\lambda_a-
\lambda_b)-(\rho_a-\rho_b)\ge0$ for $\lambda_a>\lambda_b$, because $\rho_a-\rho_b-1$ is the number of
entries strictly between $\lambda_b$ and $\lambda_a$, which is at most $\lambda_a-\lambda_b-1$.

\paragraph*{The representation, the coefficient and the measure}
Inner products are linear in the \emph{first} variable. Let $(\pi,\mathsf H)=\pi(\lambda,s)$ be the
discrete series of $G$ with Harish-Chandra parameter $\lambda$ and $v$ a unit highest weight vector
of its lowest $K$-type $\Lambda$ (unique up to a scalar of modulus one by (D3)); put
\begin{equation}\label{pbr:eq-c}
c(g)=\langle\pi(g)v,v\rangle .
\end{equation}
In the convention of {\cite[\S3.2]{P49}} (inner products conjugate-linear in the first variable) this is
the function $c_\lambda(g)=\langle v,\pi_\lambda(g)v\rangle$, so by {\cite[Theorem~3.17]{P49}} (its
hypothesis $h_\lambda\ge n^2+1$ holds) $c=F$ is given by {\cite[(14)]{P49}}, and by
{\cite[Lemma~3.25]{P49}}
\begin{equation}\label{pbr:eq-decay}
|c(g)|\le(2/\|g\|)^{j}\qquad(g\in G,\ j\in\Z,\ 0\le j\le h_\lambda).
\end{equation}
Let $dg$ be the Haar measure on $G$ compatible with the Haar measure of mass $1$ on $\auU(n)$
(\S\ref{pds:ss-ds1}(iii), Remark~\ref{pds:cayley}), and $d=d_\pi$ the formal degree of $\pi$ for
$dg$: $\int_G|c|^2dg=d^{-1}$. For another Haar measure $c'dg$ the formal degree is $d/c'$, so the
measure $d\,dg$ and every statement below are independent of this choice. Compact groups carry
their Haar measure of mass $1$.

\begin{theorem}\label{pbr:main}
Assume \eqref{pbr:eq-gap}. For every $t\in T$ with distinct eigenvalues the function
$g\mapsto c(g^{-1}tg)$ is integrable on $G$, and
\begin{equation}\label{pbr:eq-unstable}
d\int_Gc(g^{-1}tg)\,dg=(-1)^{pq}\,\frac{\sum_{w\in S_n,\ s\circ w=s}\pbrsgn(w)\,t^{\kappa\circ w}}
{\prod_{\lambda_a>\lambda_b}(t_a-t_b)},\qquad\kappa=\lambda+\tfrac{n-1}2\mathbf 1,
\end{equation}
where $(\kappa\circ w)_a=\kappa_{w(a)}$ and $t^\gamma=\prod_at_a^{\gamma_a}$. Consequently
\begin{equation}\label{pbr:eq-stable}
\sum_{\sigma\in S_n}d\int_G\aubar{c(g^{-1}t^\sigma g)}\,dg=p!\,q!\,(-1)^{pq}\,\aubar{\chi_\lambda(t)} ,
\end{equation}
which is Theorem~\ref{pm:arch}(b).
\end{theorem}

Formula \eqref{pbr:eq-unstable} is Harish-Chandra's formula for the character of $\pi$ on $T$
(\cite[Theorem~16, p.~96]{HC66}; compare Remark~\ref{pm:rem-arch}(1)); we do not use it, we prove it.

\begin{remark}[The convention]\label{pbr:rem-conv}
For $q=0$ ($G=\auU(n)$, $\pi$ the representation with character $\chi_\lambda$, $d=\dim\pi$),
$\int_{\auU(n)}\pi(g^{-1}tg)\,dg=d^{-1}\chi_\lambda(t)$ by Schur's lemma, so $d\int c(g^{-1}tg)dg=
\chi_\lambda(t)$ and \eqref{pbr:eq-stable} reads $n!\,\aubar{\chi_\lambda(t)}$: the complex conjugate
in Theorem~\ref{pm:arch}(b) is correct for the convention \eqref{pbr:eq-c}, linear in the first
variable, in which $e=d\,\aubar c$ is the kernel of the projection onto the $\pi$-isotypic part.
With the opposite convention one has to drop both conjugations in \eqref{pbr:eq-stable}. Only the
fact that the right side of \eqref{pbr:eq-stable} does not depend on $s$ beyond the sign
$(-1)^{pq}$ is used in Proposition~\ref{pm:ell}, so either reading suffices there; the convention of
\S\ref{pm:sec} is stated in Theorem~\ref{pm:arch} and Remark~\ref{pm:rem-other}(1).
\end{remark}

\paragraph*{Published inputs}\label{pbr:ss-inputs}
\begin{enumerate}
\item[(R1)] \emph{Square-integrable representations.} Schur orthogonality $\int_G\langle\pi(g)a,b
\rangle\aubar{\langle\pi(g)a',b'\rangle}dg=d^{-1}\langle a,a'\rangle\aubar{\langle b,b'\rangle}$
(\cite[Proposition~9.6, p.~284]{Knapp}); admissibility: each $K$-isotypic subspace $\mathsf H_\nu$ of $\pi$ is
finite-dimensional (\cite[Theorem~8.1, p.~205]{Knapp}). These are used in \S\ref{pm:sec} already.
\item[(R2)] \emph{Compact groups} (\cite{KnappBeyond}: Schur orthogonality, Corollary~4.10; orthonormality of characters, Corollary~4.16; Peter--Weyl, Theorem~4.20 and Corollary~4.21; the remaining classical statements below without locator). For a
unitary representation $\sigma$ of $K$, $(\dim V_\nu)\int_K\aubar{\chi^K_\nu(k)}\sigma(k)\,dk$ is the
orthogonal projection onto the $\nu$-isotypic subspace. Peter--Weyl: the $\chi^K_\nu$ form an
orthonormal basis of the space of class functions in $L^2(K)$. Every element of $K$ is
$K$-conjugate to an element of $T$, and a class function is determined by its restriction to $T$,
which is invariant under $\pbrW=\{w\in S_n:s\circ w=s\}\cong S_p\times S_q$. Weyl's integration
formula: $\int_K\phi\,dk=|\pbrW|^{-1}\int_T\phi(x)|V_{\mathrm c}(x)|^2dx$ for continuous class
functions $\phi$, and Weyl's character formula $V_{\mathrm c}\chi^K_\nu=\pbrAlt_{\pbrW}(\nu+\pbrdc)$
on $T$, with the notation of \S\ref{pbr:ss-laurent}; the same two statements for $\auU(n)$, where the
character with highest weight $\pbrLo$ is $\chi_\lambda$ and $V\chi_\lambda=\pbrAlt_{S_n}(\kappa)$.
Schur orthogonality for $\auU(n)$, and the holomorphic extension of its representations to
$\mathrm{GL}_n(\C)$ (Weyl's unitary trick).
\item[(R3)] \emph{Schmid's theorem} (\cite[Lemma~(2.22), Lemma~(3.8), Corollary~(3.9), Corollary~(4.19) and the proof of Theorem~(1.3) on p.~560]{Schmid75}: $\pi$ is the irreducible quotient of a module whose $K$-types satisfy both assertions; valid for every discrete series, see the paragraph ``On (D3)'' in {\cite[\S3.1.3]{P49}}; see also
\cite[Theorem~9.20(b),(c), p.~310]{Knapp}, and \cite[Proposition~9.6, Theorem~8.1]{Knapp} in (R1): these are stated for linear connected semisimple groups and apply to
$\auU(p,q)=Z\cdot\mathrm{SU}(p,q)$, $Z\cong\auU(1)$ the compact centre, $Z\cap\mathrm{SU}(p,q)$ finite: an
irreducible unitary (square-integrable) representation of $\auU(p,q)$ restricts irreducibly to one of
$\mathrm{SU}(p,q)$, $Z$ acting by a character, $K=Z\cdot(K\cap\mathrm{SU}(p,q))$, so $K$-isotypic
subspaces, $K$-types with their multiplicities and the lowest $K$-type correspond, and formal degrees
agree up to the normalisation of Haar measures; see the paragraph ``On (D3)'' in
{\cite[\S3.1.3]{P49}}). The $K$-type $\Lambda$ occurs in $\pi$ with
multiplicity one, and every $K$-type $\nu$ of $\pi$ has the form $\nu=\Lambda+\sum_\beta n_\beta
\beta$ with integers $n_\beta\ge0$, the sum over the positive noncompact roots $\beta$. The first
half is (D3)(b) of {\cite[\S3.1.3]{P49}}.
\item[(R4)] \emph{Linear algebra.} Fischer's inequality $\Delta_{S\cup S'}(X)\le\Delta_S(X)\Delta_{S'}
(X)$ and Jacobi's identity $\Delta_S(X^{-1})=\Delta_{S^c}(X)/\det X$ (as in Lemma~\ref{pds:fischer});
for positive semidefinite $X\le Y$ (that is, $Y-X\ge0$), $\Delta_S(X)\le\Delta_S(Y)$
\cite[Corollary~7.7.4(e)]{HornJohnson}, applied to the principal submatrices on $S$; $|\Delta_S(X)|\le\|X\|^{|S|}$.
\end{enumerate}
Real-analytic functions on a connected real-analytic manifold that are not identically zero have
null zero sets; holomorphic functions on a connected open set agreeing on a set with an
accumulation point agree. We use the minors $\Delta_S$ and the functions $\pcom_w$ of
{\cite[\S3.2.1]{P49}}, and {\cite[Lemma~3.18]{P49}}.

\emph{Plan.} \S\ref{pbr:ss-int} proves the integrability, uniformly on $K$ after multiplication by a
power of the noncompact Weyl denominator $D$. \S\ref{pbr:ss-mult} shows, from (R1)--(R3), that
$V_{\mathrm c}D^{M}\cdot d\int_Gc(g^{-1}xg)dg$ is a Laurent polynomial $P$ whose leading part is
controlled. \S\S\ref{pbr:ss-flag}--\ref{pbr:ss-open} compute the orbital integral on an open part of
$T$ as a Weyl character of $\auU(n)$ minus an error $\pbrE$, using the formula {\cite[(14)]{P49}} and
an integral formula for densities on the flag variety. \S\ref{pbr:ss-error} shows that $\pbrE$ does
not affect the monomials of low level, and \S\ref{pbr:ss-proof} identifies $P$ by comparing those
monomials.

\subsubsection{Laurent polynomials, levels and alternants}\label{pbr:ss-laurent}

A \emph{Laurent polynomial} is $P=\sum_{\gamma\in\Z^n}P_\gamma x^\gamma$, finitely many $P_\gamma\ne0$,
$x^\gamma=\prod_ax_a^{\gamma_a}$; its \emph{exponents} are the $\gamma$ with $P_\gamma\ne0$. It is
determined by its values on $T$. A \emph{formal series} is any function $\gamma\mapsto t_\gamma$ on
$\Z^n$, and for a Laurent polynomial $Q$ we put $(Qt)_\gamma=\sum_\beta Q_\beta t_{\gamma-\beta}$.
For $w\in S_n$ and $\gamma\in\Z^n$ let $(\gamma\circ w)_a=\gamma_{w(a)}$; then
$P(x^w)=\sum_\beta P_{\beta\circ w}x^\beta$. We call $P$ (or a formal series) \emph{$\pbrW$-symmetric}
if $P_{\gamma\circ w}=P_\gamma$, and \emph{$\pbrW$-antisymmetric} if $P_{\gamma\circ w}=\pbrsgn(w)
P_\gamma$, for all $w\in\pbrW$ and $\gamma$; on $T$ this means $P(x^w)=P(x)$, resp.\ $\pbrsgn(w)P(x)$.
For a subgroup $W'\subseteq S_n$ the \emph{alternant} is $\pbrAlt_{W'}(\gamma)=\sum_{w\in W'}\pbrsgn(w)
x^{\gamma\circ w}$; if $\gamma$ has distinct entries, its coefficient at $\gamma\circ w$ is
$\pbrsgn(w)$ for $w\in W'$, and it has no other exponents.

\paragraph*{Notation}
For $a\in\{1,\dots,n\}$ put
\begin{gather*}
o_a=\#\{b:\lambda_b>\lambda_a\},\qquad \pbrdc(a)=\#\{b:\varepsilon_b=\varepsilon_a,\ \lambda_b<
\lambda_a\},\qquad \pbrdD(a)=\#\{b:\varepsilon_b\ne\varepsilon_a,\ \lambda_b>\lambda_a\},\\
\pbrlev(\gamma)=\sum_ao_a\gamma_a\quad(\text{the \emph{level} of }\gamma\in\Z^n),\qquad
N=4n^2,\qquad M=N+1,
\end{gather*}
and recall $\kappa=\lambda+\frac{n-1}2\mathbf 1\in\Z^n$. The vector $o$ strictly reverses the order
of $\lambda$: $\lambda_a>\lambda_b$ implies $o_a<o_b$; and $0\le o_a\le n-1$. Let
\[
V=\prod_{\lambda_a>\lambda_b}(x_a-x_b),\qquad V_{\mathrm c}=\prod_{\substack{\lambda_a>\lambda_b\\
\varepsilon_a=\varepsilon_b}}(x_a-x_b),\qquad V_{\mathrm n}=\prod_{\substack{\lambda_a>\lambda_b\\
\varepsilon_a\ne\varepsilon_b}}(x_a-x_b),\qquad D=\prod_{\substack{\lambda_a<\lambda_b\\
\varepsilon_a\ne\varepsilon_b}}(x_a-x_b).
\]

\begin{lemma}\label{pbr:basic}
\begin{enumerate}
\item[(a)] $V=\pbrAlt_{S_n}(\delta)$ with $\delta_a=\#\{b:\lambda_b<\lambda_a\}$, and $V_{\mathrm c}=
\pbrAlt_{\pbrW}(\pbrdc)$; $V$ is $S_n$-antisymmetric and $V_{\mathrm c}$ is $\pbrW$-antisymmetric.
\item[(b)] $V=V_{\mathrm c}V_{\mathrm n}$ and $D=(-1)^{pq}V_{\mathrm n}$. There is $\zeta\in\{\pm1\}$
with $D=\zeta\prod_{\varepsilon_a=1,\ \varepsilon_b=-1}(x_a-x_b)$; so $D$ is $\pbrW$-symmetric, and
$D(x)\ne0$ if the $x_a$ are distinct.
\item[(c)] For $k\ge0$, $D^k$ has coefficient $1$ at $k\pbrdD$, and every other exponent $\beta$ of
$D^k$ has $\pbrlev(\beta)<\pbrlev(k\pbrdD)$. The exponents of $V$ and of $D^k$ have nonnegative
entries, hence nonnegative level.
\item[(d)] $\Lambda+\pbrdc+\pbrdD=\kappa$ and $\kappa-\mu=o$.
\item[(e)] \textup{(Rearrangement with a gain.)} Let $b\in\Z^n$ be such that $b_j<b_i$ implies
$o_i<o_j$. Then $\pbrlev(b)+|b_i-b_{w(i)}|\le\pbrlev(b\circ w)$ for every $w\in S_n$ and every $i$.
\end{enumerate}
\end{lemma}
\begin{proof}
(a) List the indices as $a_1,\dots,a_n$ with $\lambda$ decreasing; Vandermonde's determinant gives
$\prod_{i<j}(x_{a_i}-x_{a_j})=\det(x_{a_i}^{n-j})_{i,j}$, which is $\pbrAlt_{S_n}(\delta)$ since
$\delta(a_i)=n-i$. The same in each of the two blocks of equal sign gives $V_{\mathrm c}$. The
antisymmetry is that of determinants under permutation of rows.
(b) The first identity sorts the pairs. $D$ is $V_{\mathrm n}$ with each of its $pq$ factors
$x_b-x_a$ ($\lambda_b>\lambda_a$) replaced by $x_a-x_b$. In $D$ the factor attached to a pair
$\varepsilon_a=1$, $\varepsilon_b=-1$ is $\pm(x_a-x_b)$, with a sign that depends only on the pair; the
product of the $x_a-x_b$ over these pairs is invariant under $\pbrW$, which permutes the indices of
each sign among themselves.
(c) Each factor $x_a-x_b$ of $D$ has $\lambda_a<\lambda_b$, so $o_a>o_b$: the monomial $x_a$ has the
larger level, and $\pbrdD$ is the sum of these $e_a$. Expanding $D^k$ as a product of $kpq$ factors,
a monomial is obtained by choosing $x_a$ or $-x_b$ in each factor; its level is at most
$\pbrlev(k\pbrdD)$, with equality only if $x_a$ is chosen every time, which happens for exactly one
choice, with coefficient $1$. The last sentence is clear since $o_a\ge0$.
(d) Write $\mathrm{sb},\mathrm{sa},\mathrm{ob},\mathrm{oa}$ for the numbers of $b\ne a$ with
$\varepsilon_b=\varepsilon_a$ and $\lambda_b<\lambda_a$, resp.\ $\varepsilon_b=\varepsilon_a$ and
$\lambda_b>\lambda_a$, resp.\ $\varepsilon_b\ne\varepsilon_a$, $\lambda_b<\lambda_a$, resp.\
$\varepsilon_b\ne\varepsilon_a$, $\lambda_b>\lambda_a$. Then $2\rho_{\mathrm c,a}=\mathrm{sb}-
\mathrm{sa}$, $2\rho_{\mathrm n,a}=\mathrm{ob}-\mathrm{oa}$, $\pbrdc(a)=\mathrm{sb}$, $\pbrdD(a)=
\mathrm{oa}$, so $\Lambda_a+\pbrdc(a)+\pbrdD(a)-\lambda_a=\frac12(\mathrm{sb}+\mathrm{sa}+\mathrm{ob}
+\mathrm{oa})=\frac{n-1}2$. Similarly $\kappa_a-\mu_a=\frac{n-1}2-\rho_a=\frac12(n-1-\mathrm{sb}-
\mathrm{ob}+\mathrm{sa}+\mathrm{oa})=\mathrm{sa}+\mathrm{oa}=o_a$.
(e) If $o'\in\Z^n$ satisfies $(o'_i-o'_j)(b_i-b_j)\le0$ for all $i,j$, then $\sum_ao'_ab_a\le\sum_a
o'_ab_{w(a)}$ for every $w$ (rearrangement inequality). Since $o$ is integral and strictly reverses
the order of $b$, so does weakly $o'=o+\eta e_i$ for $\eta=\pm1$. This gives $\pbrlev(b)+\eta b_i\le
\pbrlev(b\circ w)+\eta b_{w(i)}$ for both signs, which is the claim.
\end{proof}

\begin{lemma}\label{pbr:top}
Let $k\ge0$, $a\in\Z^n$, and let $t$ be a formal series with $t_a=1$ such that every $\gamma$ with
$t_\gamma\ne0$ is $\gamma=\gamma'\circ w$ with $w\in\pbrW$ and $\pbrlev(\gamma')\le\pbrlev(a)$. Put
$b=a+k\pbrdD$ and assume that $\lambda_j<\lambda_i$ implies $b_j<b_i$. Then $(D^kt)_b=1$, and for
every $\gamma$ with $(D^kt)_\gamma\ne0$ there is $w\in\pbrW$ with $\pbrlev(\gamma\circ w)\le
\pbrlev(b)$.
\end{lemma}
\begin{proof}
Let $(D^k)_\beta t_{\gamma-\beta}\ne0$, and write $\gamma-\beta=\gamma'\circ w$ as in the hypothesis.
Then $\pbrlev(\gamma\circ w^{-1})=\pbrlev(\beta\circ w^{-1})+\pbrlev(\gamma')\le\pbrlev(\beta\circ
w^{-1})+\pbrlev(a)$. By Lemma~\ref{pbr:basic}(b), $(D^k)_{\beta\circ w^{-1}}=(D^k)_\beta\ne0$, so by
(c) $\pbrlev(\beta\circ w^{-1})\le\pbrlev(k\pbrdD)$, with equality only if $\beta\circ w^{-1}=
k\pbrdD$. Hence
\begin{equation}\label{pbr:eq-toplevel}
\pbrlev(\gamma\circ w^{-1})+[\beta\circ w^{-1}\ne k\pbrdD]\le\pbrlev(b),
\end{equation}
which gives the second claim. For $\gamma=b$: $b$ has the order of $\lambda$, so $o$ strictly
reverses it and Lemma~\ref{pbr:basic}(e) gives $\pbrlev(b\circ w^{-1})\ge\pbrlev(b)+|b_i-b_{w^{-1}
(i)}|$ for all $i$. With \eqref{pbr:eq-toplevel}, $b_{w^{-1}(i)}=b_i$ for all $i$, so $w=1$ ($b$ has
distinct entries), and $\beta=k\pbrdD$. Thus the only term of $(D^kt)_b$ is $(D^k)_{k\pbrdD}
t_a=1$.
\end{proof}

\begin{lemma}\label{pbr:tau}
Put $\pbrtop=\kappa+N\pbrdD$. If $\lambda_j<\lambda_i$ then $\pbrtop_j<\pbrtop_i$. Moreover
\[
\pbrlev(o+N\pbrdD)<h_0,\qquad \pbrlev(N\pbrdD)<h_0 .
\]
\end{lemma}
\begin{proof}
$\kappa_i-\kappa_j=\lambda_i-\lambda_j>0$ and $0\le\pbrdD\le n-1$. If $\lambda_i-\lambda_j>Nn$, then
$\pbrtop_i-\pbrtop_j\ge(\lambda_i-\lambda_j)-N(n-1)>0$. Otherwise $\lambda_i-\lambda_j\le Nn=4n^3<h_0$,
so every $a$ with $\lambda_j\le\lambda_a\le\lambda_i$ has $\varepsilon_a=\varepsilon_i$ (else
$|\lambda_a-\lambda_i|\ge h_\lambda\ge h_0$); in particular $\varepsilon_j=\varepsilon_i$, and an index $b$
of the other sign with $\lambda_b>\lambda_j$ has $\lambda_b>\lambda_i$. So $\pbrdD(j)=\pbrdD(i)$ and
$\pbrtop_i-\pbrtop_j=\lambda_i-\lambda_j>0$. For the bounds, $o_a,\pbrdD(a)\le n-1$, so
$\pbrlev(o+N\pbrdD)\le n(n-1)^2(1+N)<n^3+4n^5<h_0$, and $\pbrlev(N\pbrdD)\le Nn(n-1)^2<h_0$.
\end{proof}

\subsubsection{Integrability}\label{pbr:ss-int}

Write $+$ and $-$ for the sets of indices of sign $+1$ and $-1$, $X_{++}$, $X_{--}$ for the
corresponding blocks of a matrix $X$, and $\Delta_+=\Delta_{+}$ for the principal minor on $+$.

\begin{lemma}\label{pbr:minor}
Let $u\in G$, $x\in T$ and $g=u^{-1}xu$. Then
$(\|u\|/2)^2\,\bigl|\prod_{a\in+,\,b\in-}(x_a-x_b)\bigr|\le(2\|g\|)^{pq}$.
\end{lemma}
\begin{proof}
If $pq=0$, $G$ is compact, $\|u\|=1$ and the claim is $\frac14\le1$. Let $\Pi=\prod_{b\in-}(x-x_b)$,
a diagonal matrix with $\Pi_{aa}=\prod_{b\in-}(x_a-x_b)$, which is $0$ for $a\in-$. Then
$u^{-1}\Pi u=\prod_{b\in-}(g-x_b)$, so $\|u^{-1}\Pi u\|\le(\|g\|+1)^q\le(2\|g\|)^q$ (as $\|g\|\ge1$),
and $|\Delta_+(u^{-1}\Pi u)|\le(2\|g\|)^{pq}$ by (R4). Since $\Pi$ vanishes off the block $++$,
$(u^{-1}\Pi u)_{++}=(u^{-1})_{++}\Pi_{++}u_{++}$, and $(u^{-1})_{++}=(Ju^*J)_{++}=(u_{++})^*$; hence
$|\Delta_+(u^{-1}\Pi u)|=|\det u_{++}|^2\prod_{a\in+}|\Pi_{aa}|$. By Jacobi's identity
$\Delta_+(u^{-1})=\det u_{--}/\det u$ and $|\det u|=1$, so $|\det u_{++}|=|\det u_{--}|$ and
$|\det u_{++}|^2=|\det u_{++}\det u_{--}|=|\det a(u)|\ge(\|u\|/2)^2$ by Lemma~\ref{pds:blockdet}.
\end{proof}

For $k\in K$ write $k=\operatorname{diag}(k_+,k_-)$ in the blocks, and put $\widetilde D(k)=
\zeta(-1)^{pq}\det\chi_{k_+}(k_-)$, where $\chi_{k_+}(z)=\det(z-k_+)$ and $\zeta$ is as in
Lemma~\ref{pbr:basic}(b). It is a continuous class function on $K$, and $\widetilde D(x)=
\zeta\prod_{b\in-}\prod_{a\in+}(x_a-x_b)=D(x)$ for $x\in T$.

\begin{proposition}\label{pbr:int}
The map $\Psi\colon G\to C(K)$, $\Psi(g)(k)=\widetilde D(k)^Mc(g^{-1}kg)$, is continuous, and
$\int_G\sup_K|\Psi(g)|\,dg<\infty$; so $\Psi$ is Bochner integrable in the Banach space $C(K)$ with the
supremum norm. For $x\in T$ with distinct entries, $g\mapsto c(g^{-1}xg)$ is integrable on $G$.
\end{proposition}
\begin{proof}
Continuity: $(g,k)\mapsto c(g^{-1}kg)$ is continuous and $K$ is compact. Let $k=yxy^{-1}$ with
$y\in K$, $x\in T$ (R2), and $u=y^{-1}g$, so $c(g^{-1}kg)=c(u^{-1}xu)$ and $\|u\|=\|g\|$. By
Lemma~\ref{pbr:minor} and \eqref{pbr:eq-decay} with $j=pqM$ (note $pqM\le\frac{n^2}4(4n^2+1)\le
h_0\le h_\lambda$),
\[
|\widetilde D(k)|^M|c(u^{-1}xu)|\le\bigl(2\|u^{-1}xu\|\bigr)^{pqM}(2/\|u\|)^{2M}\bigl(2/\|u^{-1}xu\|
\bigr)^{pqM}=4^{pqM}(2/\|g\|)^{2M}.
\]
Since $2M=8n^2+2>n^2$, $\int_G\|g\|^{-2M}dg<\infty$ by Lemma~\ref{pm:kernel}(c). For the last claim
evaluate at $k=x$, where $\widetilde D(x)=D(x)\ne0$ (Lemma~\ref{pbr:basic}(b)).
\end{proof}

\subsubsection{Multiplicities of $K$-types and the Laurent polynomial}\label{pbr:ss-mult}

For a $K$-dominant $\nu$ let $m(\nu)$ be the multiplicity of $V_\nu$ in $\pi|_K$, finite by (R1).

\begin{lemma}\label{pbr:mult}
Let $\nu$ be $K$-dominant and $R_\nu(g)=\int_K\aubar{\chi^K_\nu(k)}\,c(g^{-1}kg)\,dk$. Then
$R_\nu\ge0$, $R_\nu\in L^1(G)$ and $d\int_GR_\nu\,dg=m(\nu)$. Moreover $m(\Lambda)=1$, and
$m(\nu)\ne0$ implies $\pbrlev(\nu)\le\pbrlev(\Lambda)$.
\end{lemma}
\begin{proof}
$c(g^{-1}kg)=\langle\pi(k)\pi(g)v,\pi(g)v\rangle$, so by (R2) $R_\nu(g)=(\dim V_\nu)^{-1}\langle
E_\nu\pi(g)v,\pi(g)v\rangle$, with $E_\nu$ the projection onto $\mathsf H_\nu$. Let $e_1,\dots,e_r$ be
an orthonormal basis of $\mathsf H_\nu$, $r=m(\nu)\dim V_\nu$. Then $R_\nu(g)=(\dim V_\nu)^{-1}
\sum_i|\langle\pi(g)v,e_i\rangle|^2$, and by Schur orthogonality (R1) $d\int_G|\langle\pi(g)v,e_i
\rangle|^2dg=1$; so $d\int_GR_\nu=r/\dim V_\nu=m(\nu)$. By (R3), $m(\Lambda)=1$, and if $m(\nu)\ne0$ then
$\nu-\Lambda$ is a sum of positive noncompact roots $e_a-e_b$, $\lambda_a>\lambda_b$, each of level
$o_a-o_b<0$.
\end{proof}

For a continuous function $f$ on $T$ and $\gamma\in\Z^n$ let $\hat f(\gamma)=\int_Tf(x)x^{-\gamma}dx$.
A vector $\gamma\in\Z^n$ is \emph{$\pbrW$-regular} if $\gamma_a\ne\gamma_b$ whenever $a\ne b$ have
the same sign; each $\pbrW$-regular $\gamma$ is $(\nu+\pbrdc)\circ w$ for exactly one $K$-dominant
$\nu$ and one $w\in\pbrW$ (sort the entries in each block in the order of $\lambda$), and the
$\nu+\pbrdc$, $\nu$ $K$-dominant, are exactly the $\pbrW$-regular vectors that are decreasing in the
order of $\lambda$ on each block. Hence the formal series
\[
t=\sum_{\nu\ K\text{-dominant}}m(\nu)\,\pbrAlt_{\pbrW}(\nu+\pbrdc),\qquad t_\gamma=\pbrsgn(w)m(\nu)
\text{ if }\gamma=(\nu+\pbrdc)\circ w,\quad t_\gamma=0\text{ if $\gamma$ is not $\pbrW$-regular},
\]
is well defined and $\pbrW$-antisymmetric.

\begin{proposition}\label{pbr:laurent}
Let $h=d\int_G\Psi(g)\,dg\in C(K)$ (Proposition~\ref{pbr:int}), and for $x\in T$ such that
$g\mapsto c(g^{-1}xg)$ is integrable put $\pbrTh(x)=d\int_Gc(g^{-1}xg)\,dg$.
\begin{enumerate}
\item[(a)] $h$ is a class function on $K$, and $h(x)=D(x)^M\pbrTh(x)$ whenever $\pbrTh(x)$ is
defined.
\item[(b)] There is a $\pbrW$-antisymmetric Laurent polynomial $P$ with $P(x)=V_{\mathrm c}(x)h(x)$
for $x\in T$, and $P=D^Mt$.
\item[(c)] $t_{\Lambda+\pbrdc}=1$, and every $\gamma$ with $t_\gamma\ne0$ is $\gamma'\circ w$ with
$w\in\pbrW$ and $\pbrlev(\gamma')\le\pbrlev(\Lambda+\pbrdc)$.
\end{enumerate}
\end{proposition}
\begin{proof}
(a) Evaluation at $k$ is a continuous functional on $C(K)$, so $h(k)=d\int_G\widetilde D(k)^Mc(g^{-1}
kg)\,dg$; for $y\in K$, the substitution $g\mapsto yg$ gives $h(yky^{-1})=h(k)$. For $k=x$ with
$\pbrTh(x)$ defined, split off the constant $\widetilde D(x)^M=D(x)^M$.

\emph{Step 1: coefficients against characters.} Let $f$ be a continuous class function on $K$ and
$\langle\chi^K_\nu,f\rangle=\int_K\aubar{\chi^K_\nu}f\,dk$. The function $V_{\mathrm c}f$ on $T$ is
$\pbrW$-antisymmetric, so $\widehat{V_{\mathrm c}f}(\gamma\circ w)=\pbrsgn(w)\widehat{V_{\mathrm c}f}
(\gamma)$ for $w\in\pbrW$ (substitute $x\mapsto x^w$ in the integral, using
$x^{\gamma\circ w}=(x^{w^{-1}})^\gamma$); in particular $\widehat{V_{\mathrm c}f}(\gamma)=0$ if
$\gamma$ is not $\pbrW$-regular (use the transposition of two equal entries). By (R2), since
$\aubar{x^\gamma}=x^{-\gamma}$ on $T$,
\[
\langle\chi^K_\nu,f\rangle=\frac1{|\pbrW|}\int_T\aubar{V_{\mathrm c}\chi^K_\nu}\,V_{\mathrm c}f\,dx
=\frac1{|\pbrW|}\sum_{w\in\pbrW}\pbrsgn(w)\,\widehat{V_{\mathrm c}f}\bigl((\nu+\pbrdc)\circ w\bigr)
=\widehat{V_{\mathrm c}f}(\nu+\pbrdc).
\]

\emph{Step 2: the coefficients of $h$ are integers.} Let $\bar D^M$ be the Laurent polynomial with
$(\bar D^M)_\beta=(D^M)_{-\beta}$; as $D$ has integer coefficients, $\aubar{D(x)^M}=\bar D^M(x)$ on
$T$. Fix $\nu$ and let $f_\nu=\chi^K_\nu\,\aubar{\widetilde D}{}^M$, a continuous class function. By
Step 1 and (R2), for $K$-dominant $\nu'$,
\begin{equation}\label{pbr:eq-n}
n_{\nu\nu'}:=\langle\chi^K_{\nu'},f_\nu\rangle=\bigl(\pbrAlt_{\pbrW}(\nu+\pbrdc)\bar D^M\bigr)_{\nu'+
\pbrdc}=\sum_\beta(D^M)_\beta\,\pbrAlt_{\pbrW}(\nu+\pbrdc)_{\nu'+\pbrdc+\beta}\in\Z,
\end{equation}
and $n_{\nu\nu'}\ne0$ for finitely many $\nu'$ only. By Peter--Weyl, $f_\nu=\sum_{\nu'}n_{\nu\nu'}
\chi^K_{\nu'}$ (the difference is a continuous class function orthogonal to all characters), so
$\aubar{\chi^K_\nu}\widetilde D^M=\sum_{\nu'}n_{\nu\nu'}\aubar{\chi^K_{\nu'}}$. Pairing with
$\chi^K_\nu$ is a continuous functional on $C(K)$, hence commutes with the Bochner integral, and by
Lemma~\ref{pbr:mult}
\[
\langle\chi^K_\nu,h\rangle=d\int_G\int_K\aubar{\chi^K_\nu(k)}\widetilde D(k)^Mc(g^{-1}kg)\,dk\,dg=
\sum_{\nu'}n_{\nu\nu'}\,d\int_GR_{\nu'}\,dg=\sum_{\nu'}n_{\nu\nu'}m(\nu')\in\Z .
\]

\emph{Step 3: $h$ is a finite sum of characters.} By Bessel's inequality $\sum_\nu|\langle\chi^K_\nu,h
\rangle|^2\le\|h\|_{L^2}^2<\infty$; integers with a finite sum of squares vanish with finitely many
exceptions. By (a) and Peter--Weyl, $h=\sum_\nu\langle\chi^K_\nu,h\rangle\chi^K_\nu$, a finite sum, and
by Weyl's character formula $P:=\sum_\nu\langle\chi^K_\nu,h\rangle\pbrAlt_{\pbrW}(\nu+\pbrdc)$ is a
Laurent polynomial with $P=V_{\mathrm c}h$ on $T$. It is $\pbrW$-antisymmetric, and $P_\gamma=
\widehat{V_{\mathrm c}h}(\gamma)$.

\emph{Step 4: $P=D^Mt$.} Both sides are $\pbrW$-antisymmetric ($D^M$ is $\pbrW$-symmetric, so
$(D^Mt)_{\gamma\circ w}=\sum_\beta(D^M)_{\beta\circ w}t_{(\gamma-\beta)\circ w}=\pbrsgn(w)(D^Mt)_\gamma$),
so both vanish at the $\gamma$ that are not $\pbrW$-regular, and it suffices to compare them at
$\gamma=\nu+\pbrdc$, $\nu$ $K$-dominant. There $P_\gamma=\langle\chi^K_\nu,h\rangle$ by Steps 1 and 3,
which by Step 2 and \eqref{pbr:eq-n} is $\sum_{\nu'}m(\nu')\sum_{w\in\pbrW}\pbrsgn(w)(D^M)_{(\nu+
\pbrdc)\circ w-(\nu'+\pbrdc)}$. By the symmetry of $D^M$ under $w^{-1}$,
$(D^M)_{(\nu+\pbrdc)\circ w-(\nu'+\pbrdc)}=(D^M)_{(\nu+\pbrdc)-(\nu'+\pbrdc)\circ w^{-1}}$, and
$\sum_{w}\pbrsgn(w)(D^M)_{\gamma-(\nu'+\pbrdc)\circ w^{-1}}=\sum_\beta(D^M)_\beta\,\pbrAlt_{\pbrW}
(\nu'+\pbrdc)_{\gamma-\beta}$. Hence $P_\gamma=\sum_\beta(D^M)_\beta\sum_{\nu'}m(\nu')\pbrAlt_{\pbrW}
(\nu'+\pbrdc)_{\gamma-\beta}=(D^Mt)_\gamma$; the sums are finite (finitely many $\beta$, and for each
$\beta$ at most one $\nu'$ contributes).

(c) $t_{\Lambda+\pbrdc}=m(\Lambda)=1$, and if $t_\gamma\ne0$ then $\gamma=(\nu+\pbrdc)\circ w$ with
$m(\nu)\ne0$, so $\pbrlev(\nu+\pbrdc)\le\pbrlev(\Lambda+\pbrdc)$ by Lemma~\ref{pbr:mult}.
\end{proof}

\subsubsection{Densities on the flag variety}\label{pbr:ss-flag}

Let $\pbrB=\{b\in\mathrm{GL}_n(\C):b_{ij}=0\text{ whenever }\lambda_i>\lambda_j\}$ (upper triangular
when the indices are listed with $\lambda$ increasing), $\pbrN=\{\nu:\nu_{aa}=1,\ \nu_{ij}=0\text{ for
}i\ne j\text{ unless }\lambda_i>\lambda_j\}$, $\pbrn$ its Lie algebra, and for $Y\in M_n(\C)$ let
$Y_-\in\pbrn$ be the matrix of the entries $Y_{ij}$ with $\lambda_i>\lambda_j$. For $b\in\pbrB$ and
$w\in\Z^n$ put $b^w=\prod_ab_{aa}^{w_a}$, and recall $2\rho_a=\#\{b:\lambda_b<\lambda_a\}-\#\{b:
\lambda_b>\lambda_a\}$. The Haar measure $d\nu$ of the unipotent group $\pbrN$ is Lebesgue measure in
the entries $\nu_{ij}$, $\lambda_i>\lambda_j$. A \emph{density} is a Borel function $F$ on
$\mathrm{GL}_n(\C)$, with values in $\C$ or in $[0,\infty]$, such that $F(yb)=|b^{2\rho}|^2F(y)$ for
$y\in\mathrm{GL}_n(\C)$ and $b\in\pbrB$.

\begin{lemma}\label{pbr:jac}
For $b\in\pbrB$ and $X\in M_n(\C)$, $(bXb^{-1})_-=(bX_-b^{-1})_-$, and the $\C$-linear map
$Y\mapsto(bYb^{-1})_-$ of $\pbrn$ has determinant $b^{2\rho}$.
\end{lemma}
\begin{proof}
$X-X_-$ has its entries at $\lambda_i\le\lambda_j$, and so has $b(X-X_-)b^{-1}$, as $\pbrB$ is a group;
this gives the first claim. Let $r(a)$ be the rank of $\lambda_a$ among the entries of $\lambda$. For
$\lambda_i>\lambda_j$, $(bYb^{-1})_{ij}=\sum_{k,l}b_{ik}Y_{kl}(b^{-1})_{lj}$, and a nonzero term has
$\lambda_i\le\lambda_k$, $\lambda_l\le\lambda_j$, $\lambda_k>\lambda_l$, so $r(k)-r(l)\ge r(i)-r(j)$
with equality only for $(k,l)=(i,j)$. Ordering the pairs by decreasing $r(i)-r(j)$, the matrix of
the map is triangular with diagonal entries $b_{ii}(b^{-1})_{jj}=b_{ii}/b_{jj}$, and
$\prod_{\lambda_i>\lambda_j}b_{ii}/b_{jj}=\prod_ab_{aa}^{\#\{j:\lambda_j<\lambda_a\}-\#\{i:\lambda_i>
\lambda_a\}}=b^{2\rho}$.
\end{proof}

By Gaussian elimination, $g\in\mathrm{GL}_n(\C)$ lies in $\pbrN\pbrB$ if and only if the principal
minors $\Delta_{L_r}(g)$, $L_r$ the set of the $r$ indices with the smallest $\lambda_a$, are nonzero
for $1\le r\le n$; then $g=\nu(g)b(g)$ uniquely, with $\nu(g)$, $b(g)$ rational in $g$. A group
$H=\auU(J')=\{h:h^*J'h=J'\}$, $J'$ a diagonal sign matrix, carries the Haar measure compatible with
the Haar measure of mass $1$ on $\auU(n)$ (\S\ref{pds:ss-ds1}(iii)); for $J'=\pm1$ it is that measure.

\begin{proposition}\label{pbr:flag}
There is a constant $j>0$ such that for every $H=\auU(J')$ as above and every density $F$ which is
nonnegative, or such that $|F|$ is integrable on $H$,
\[
\int_HF(h)\,dh=j^{-1}\int_{\pbrN\cap H\pbrB}F(\nu)\,d\nu .
\]
\end{proposition}
\begin{proof}
(i) $H\cap\pbrB=T$: if $b\in H\cap\pbrB$ then $b^{-1}=J'b^*J'$ has its nonzero entries at
$\lambda_i\ge\lambda_j$ and also, being in $\pbrB$, at $\lambda_i\le\lambda_j$; so $b$ is diagonal,
and $b^*J'b=J'$ gives $|b_{aa}|=1$.

(ii) $H\smallsetminus\pbrN\pbrB$ is null: it is the zero set of the real-analytic function
$\prod_r\Delta_{L_r}$ on the connected group $H$, which is $1$ at $1$.

(iii) Let $\varphi(h)=\nu(h)$ on $H\cap\pbrN\pbrB$. Then $\varphi(ht)=\varphi(h)$ for $t\in T$, and if
$\varphi(h)=\varphi(h')$ then $h^{-1}h'=b(h)^{-1}b(h')\in H\cap\pbrB=T$. The image of $\varphi$ is
$\pbrN\cap H\pbrB$. For $X\in\mathfrak h=\operatorname{Lie}H$ and $h=\nu b$, $he^{sX}=\nu(1+sY+O(s^2))
b$ with $Y=bXb^{-1}$, and $1+sY=(1+sY_-)(1+s(Y-Y_-))+O(s^2)$ with $1+s(Y-Y_-)$ triangular, so
\[
d\varphi_h(hX)=\nu\cdot(bXb^{-1})_-=\nu\cdot(bX_-b^{-1})_-
\]
by Lemma~\ref{pbr:jac}. If this vanishes then $X_-=0$; then $X\in\mathfrak h$ has its nonzero entries
at $\lambda_i\le\lambda_j$, and $X=-J'X^*J'$ has them at $\lambda_i\ge\lambda_j$, so $X\in
\mathfrak t=\operatorname{Lie}T$. As $\dim_\R\mathfrak h-\dim\mathfrak t=n^2-n=\dim_\R\pbrn$, the map
$\bar\varphi\colon(H\cap\pbrN\pbrB)/T\to\pbrN$ is an injective local diffeomorphism onto the open set
$\pbrN\cap H\pbrB$.

(iv) Let $\mathfrak m\subset\mathfrak h$ be the subspace of matrices with zero diagonal, a complement
of $\mathfrak t$. The Haar measure of $H$ corresponds to a Lebesgue measure on $\mathfrak h$ through
the exponential map near $0$, and that of $T$ (mass $1$) to one on $\mathfrak t$; let $d_{\mathfrak m}$
be the quotient Lebesgue measure on $\mathfrak m\cong\mathfrak h/\mathfrak t$. The invariant measure
on $H/T$ with $\int_Hf=\int_{H/T}\int_Tf(ht)\,dt$ has, in the chart $X\mapsto h\exp(X)T$ of a
neighbourhood of $hT$, a continuous density with respect to $d_{\mathfrak m}$ equal to $1$ at $X=0$.
Let $j_H=|\det_\R(X\mapsto X_-)|$ for the map $\mathfrak m\to\pbrn$ and the measures
$d_{\mathfrak m}$, $d\nu$. By (iii), Lemma~\ref{pbr:jac} and the left invariance of $d\nu$, the
Jacobian of $\bar\varphi$ at $hT$ is $j_H|b(h)^{2\rho}|^2$ (it depends only on $hT$ since
$|t^{2\rho}|=1$). By the change of variables formula and (ii), for a density $F\ge0$, using
$F(ht)=F(h)$ and $F(h)=F(\nu(h)b(h))=|b(h)^{2\rho}|^2F(\nu(h))$,
\[
\int_HF\,dh=\int_{(H\cap\pbrN\pbrB)/T}F(h)\,d(hT)=\int_{\pbrN\cap H\pbrB}\frac{|b^{2\rho}|^2F(\nu)}
{j_H|b^{2\rho}|^2}\,d\nu=j_H^{-1}\int_{\pbrN\cap H\pbrB}F(\nu)\,d\nu .
\]
A complex density with $|F|$ integrable is a combination of four nonnegative densities
($(\pbrre F)^\pm$, $(\operatorname{Im}F)^\pm$).

(v) $j_H$ does not depend on $J'$. By definition of compatibility, the map $c\colon X\mapsto X_{\mathrm
d}+iX_{\mathrm o}$ from $\mathfrak h$ to $\mathfrak u(n)$, where $X_{\mathrm d}$ (resp.\ $X_{\mathrm o}$)
keeps the entries with $\varepsilon'_i=\varepsilon'_j$ (resp.\ $\ne$), carries the Lebesgue measure of
$\mathfrak h$ to that of $\mathfrak u(n)$. It is the identity on $\mathfrak t$ and maps $\mathfrak m$
onto the corresponding subspace of $\mathfrak u(n)$, so it carries $d_{\mathfrak m}$ to the measure
used for $\auU(n)$. Moreover $(cX)_-=\iota(X_-)$, where $\iota$ multiplies the entries $(i,j)$ with
$\varepsilon'_i\ne\varepsilon'_j$ by $i$, a $\C$-linear map with $|\det_\R\iota|=1$. Hence
$j_H=j_{\auU(n)}=:j$.
\end{proof}

\begin{corollary}\label{pbr:flagG}
The set $G\pbrB$ is open. Let $F$ be a density that is continuous on $\mathrm{GL}_n(\C)$. Then
$F|_G\in L^1(G)$ and
\[
\int_GF(g)\,dg=\int_{\auU(n)\cap G\pbrB}F(k)\,dk=\int_{\auU(n)}F(k)\,dk-\int_{\auU(n)\smallsetminus
G\pbrB}F(k)\,dk .
\]
\end{corollary}
\begin{proof}
The map $G\times\pbrB\to\mathrm{GL}_n(\C)$, $(g,b)\mapsto gb$, is a submersion: its differential at
$(g,b)$ has image $gb\,(b^{-1}\mathfrak gb+\mathfrak b)$, and $b^{-1}\mathfrak gb\cap\mathfrak b=
b^{-1}(\mathfrak g\cap\mathfrak b)b=b^{-1}\mathfrak tb$ has dimension $n$ (as in (i) above), so the
image has real dimension $n^2+(n^2+n)-n=2n^2$. So $G\pbrB$ is open. The function $F1_{G\pbrB}$ is a
density, since $G\pbrB\pbrB=G\pbrB$. By Gram--Schmidt, $\auU(n)\pbrB=\mathrm{GL}_n(\C)$, so
$\pbrN\cap\auU(n)\pbrB=\pbrN$. Proposition~\ref{pbr:flag} for $H=\auU(n)$ and for $H=G$ gives
$\int_{\auU(n)}|F|1_{G\pbrB}=j^{-1}\int_{\pbrN\cap G\pbrB}|F|=\int_G|F|$, finite because $\auU(n)$ is
compact; and then the same with $F$ in place of $|F|$.
\end{proof}

\subsubsection{The orbital integral on an open part of the torus}\label{pbr:ss-open}

Let $\pbrTp=\{x\in T:\pbrre(\varepsilon_ax_a)>0\text{ for all }a\}$, the set of $x\in T$ with $Jx$
accretive. For an accretive matrix $A$ and $y\in\mathrm{GL}_n(\C)$ put
\[
F_A(y)=\frac{\pcom_{-\pbrLo}(y^*Jy)}{\pcom_{-\mu}(y^*Ay)},\qquad \pbrE(A)=\int_{\auU(n)\smallsetminus
G\pbrB}F_A(k)\,dk .
\]

\begin{lemma}\label{pbr:FA}
$F_A$ is a density, continuous on $\mathrm{GL}_n(\C)$, and $A\mapsto\pbrE(A)$ is holomorphic on the
open set of accretive matrices.
\end{lemma}
\begin{proof}
$y^*Jy$ is invertible and $y^*Ay$ is accretive ({\cite[Lemma~3.18(d)]{P49}}), so the numerator is
defined and continuous (a polynomial in the entries divided by a power of $\det(y^*Jy)$) and the
denominator is continuous and nonzero. Let $b\in\pbrB$ and $w\in\{-\pbrLo,-\mu\}$. If $w_c<w_d$ then
$\lambda_c>\lambda_d$ (for $w=-\pbrLo$ because $\pbrLo$ is dominant, for $w=-\mu$ by
{\cite[Lemma~3.18(c)]{P49}}), so $b_{cd}=0$; likewise $(b^*)_{cd}=\aubar{b_{dc}}=0$ whenever $w_d<w_c$.
{\cite[Lemma~3.18(b)]{P49}} gives $\pcom_w(b^*Xb)=\aubar{b^w}\,b^w\pcom_w(X)$ for invertible $X$, so
$F_A(yb)=|b^{-\pbrLo}|^2|b^{-\mu}|^{-2}F_A(y)=|b^{2\rho}|^2F_A(y)$, as $\mu-\pbrLo=2\rho$. The integrand
of $\pbrE$ is jointly continuous in $(A,k)$ and holomorphic in $A$, and the domain is compact.
\end{proof}

Let $(\sigma,V_\sigma)$ be the irreducible representation of $\auU(n)$ with highest weight $\pbrLo$
(character $\chi_\lambda$), extended holomorphically to $\mathrm{GL}_n(\C)$, with an invariant inner
product and a unit highest weight vector $v^\circ$, and $\pbrNo=(\dim V_\sigma)^{-1}$.

\begin{lemma}\label{pbr:hw}
$\pcom_{\pbrLo}(g)=\langle\sigma(g)v^\circ,v^\circ\rangle$ for $g\in\mathrm{GL}_n(\C)$; hence
$\int_{\auU(n)}|\pcom_{\pbrLo}|^2=\pbrNo$ and $\int_{\auU(n)}\pcom_{\pbrLo}(y^{-1}xy)\,dy=\pbrNo
\chi_\lambda(x)$. For $z,a\in\mathrm{GL}_n(\C)$ with $a^{-1}z$ accretive,
\begin{equation}\label{pbr:eq-FJn}
\pcom_{\pbrLo}(z)=\frac1{\pbrNo}\int_{\auU(n)}\frac{\pcom_{\pbrLo}(au)\,\aubar{\pcom_{\pbrLo}(u)}}
{\pcom_{-\mu}(u^{-1}a^{-1}zu)}\,du .
\end{equation}
\end{lemma}
\begin{proof}
$\sigma(g)^*=\sigma(g^*)$: both sides are antiholomorphic in $g$ and agree on $\auU(n)$, and a
holomorphic function on $\mathrm{GL}_n(\C)$ vanishing on $\auU(n)$ vanishes (its derivatives at $1$
along $\mathfrak u(n)$, hence along $\mathfrak{gl}_n=\mathfrak u(n)\oplus i\mathfrak u(n)$, vanish).
Let $g=\nu't\eta$ with $\nu'$ unipotent with nonzero off-diagonal entries only at $\lambda_a<
\lambda_b$, $t$ diagonal, $\eta$ unipotent with nonzero off-diagonal entries only at $\lambda_a>
\lambda_b$; such products form a dense open subset. Since $v^\circ$ is killed by the $d\sigma(E_{ab})$,
$\lambda_a>\lambda_b$, $\sigma(\eta)v^\circ=v^\circ=\sigma(\nu'^*)v^\circ$, so $\langle\sigma(g)v^\circ,
v^\circ\rangle=\langle\sigma(t)v^\circ,\sigma(\nu'^*)v^\circ\rangle=t^{\pbrLo}$. By
{\cite[Lemma~3.18(b)]{P49}} (if $\pbrLo_a<\pbrLo_b$ then $\lambda_a<\lambda_b$, so $\eta_{ab}=0$;
similarly for $\nu'$), $\pcom_{\pbrLo}(g)=t^{\pbrLo}$. Both sides are holomorphic on
$\mathrm{GL}_n(\C)$, so they agree. The two integrals follow from Schur orthogonality on $\auU(n)$ and
from $\int\sigma(y^{-1}xy)\,dy=\pbrNo\chi_\lambda(x)\cdot1$.

{\cite[Proposition~3.21]{P49}} and its proof apply verbatim to the signature $(n,0)$, that is, to $J=1$,
$\pcoD=M_n(\C)$, $\pcoH=\mathrm{GL}_n(\C)$, $K=\auU(n)$; then $\rho_{\mathrm n}=0$, $\Lambda=\pbrLo$,
and $\mu$ is unchanged (no gap condition enters there). So the right side of \eqref{pbr:eq-FJn} does
not depend on $a$ with $a^{-1}z$ accretive. For $a=z$ it is $(\pbrNo)^{-1}\int\pcom_{\pbrLo}(zu)
\aubar{\pcom_{\pbrLo}(u)}\,du$ ($\pcom_{-\mu}(1)=1$), and with $\pcom_{\pbrLo}(zu)=\langle\sigma(u)
v^\circ,\sigma(z^*)v^\circ\rangle$ Schur orthogonality gives $(\pbrNo)^{-1}\cdot\pbrNo\,\aubar{\langle
\sigma(z^*)v^\circ,v^\circ\rangle}=\langle\sigma(z)v^\circ,v^\circ\rangle=\pcom_{\pbrLo}(z)$.
\end{proof}

\begin{lemma}\label{pbr:Un}
For $x\in\pbrTp$, $\int_{\auU(n)}F_{Jx}(k)\,dk=\pbrNo\chi_\lambda(x)$.
\end{lemma}
\begin{proof}
Write $m=\pcom_{\pbrLo}$ and $\psi(Y)=\pcom_{-\mu}(Y)^{-1}$. For $y\in\auU(n)$ apply
\eqref{pbr:eq-FJn} to $z=y^{-1}xy$ and $a=y^{-1}Jy$; then $a^{-1}z=y^*(Jx)y$ is accretive. With
Lemma~\ref{pbr:hw},
\[
\pbrNo\chi_\lambda(x)=\int m(y^{-1}xy)\,dy=\frac1{\pbrNo}\iint m(y^{-1}Jyk)\,\aubar{m(k)}\,
\psi\bigl((yk)^*Jx(yk)\bigr)\,dk\,dy .
\]
The substitution $(y,k)=(vu,u^{-1})$ preserves the product measure on $\auU(n)^2$; it gives
$y^{-1}Jyk=u^{-1}v^*Jv$, $\aubar{m(u^{-1})}=m(u)$ and $yk=v$. By Schur orthogonality, with
$b=v^*Jv$, $\int m(u^{-1}b)m(u)\,du=\int\langle\sigma(u)v^\circ,v^\circ\rangle\allowbreak \aubar{\langle\sigma(u)
v^\circ,\sigma(b)v^\circ\rangle}du=\pbrNo\,m(b)$. Hence $\pbrNo\chi_\lambda(x)=\int m(v^*Jv)\psi(v^*Jxv)
\,dv$, and $m(v^*Jv)=\pcom_{-\pbrLo}((v^*Jv)^{-1})=\pcom_{-\pbrLo}(v^*Jv)$ by
{\cite[Lemma~3.18(a)]{P49}}, since $(v^*Jv)^{-1}=v^*Jv$.
\end{proof}

\begin{proposition}\label{pbr:orbit}
Let $x\in\pbrTp$ and $\sigma_\lambda=\prod_a\varepsilon_a^{\Lambda_a+\pbrLo_a}\in\{\pm1\}$. Then
$g\mapsto c(g^{-1}xg)$ is integrable on $G$ and
\[
\sigma_\lambda\int_Gc(g^{-1}xg)\,dg=\pbrNo\chi_\lambda(x)-\pbrE(Jx).
\]
\end{proposition}
\begin{proof}
Let $g\in G$ and $y=g^{-1}xg\in G$. Since $g^{-1}=Jg^*J$, $J^{-1}y=g^*(Jx)g$ is accretive, so $J\in
A_y$ in the notation of {\cite[\S3.2.4]{P49}}. By {\cite[Theorem~3.17]{P49}}, $c(y)=F(y)=d_\Lambda I(y;k)$
with $y=kz$ the polar decomposition, and by {\cite[Proposition~3.21]{P49}}, $c(y)=d_\Lambda I(y;J)$. In
{\cite[(15)]{P49}}, $\pcom_\Lambda(Ju)=J^\Lambda\pcom_\Lambda(u)$ ({\cite[Corollary~3.19(b)]{P49}}) and
$u^{-1}J^{-1}yu=(gu)^*Jx(gu)$. On the other hand $(gu)^*J(gu)=J$ for $u\in K$, so the numerator of
$F_{Jx}(gu)$ is $\pcom_{-\pbrLo}(J)=J^{-\pbrLo}=J^{\pbrLo}$. Hence
\[
\sigma_\lambda\,c(g^{-1}xg)=\int_Kd_\Lambda|\pcom_\Lambda(u)|^2F_{Jx}(gu)\,du .
\]
By Lemma~\ref{pbr:FA} and Corollary~\ref{pbr:flagG}, $F_{Jx}\in L^1(G)$, and $\int_G|F_{Jx}(gu)|dg=
\int_G|F_{Jx}|$ for each $u$. By Tonelli the right side is integrable over $G$, and by Fubini and
$d_\Lambda\int_K|\pcom_\Lambda|^2=1$, $\sigma_\lambda\int_Gc(g^{-1}xg)\,dg=\int_GF_{Jx}$. Conclude with
Corollary~\ref{pbr:flagG} and Lemma~\ref{pbr:Un}.
\end{proof}

\subsubsection{The error term}\label{pbr:ss-error}

\begin{lemma}\label{pbr:gain}
Let $u\in\auU(n)$, $f\in\Z^n$, $D_f=\operatorname{diag}(2^{f_1},\dots,2^{f_n})$ and $X_f=u^*D_fu$. Let
$w\in\Z^n$, $m\in\Z$ and $h'\in\N$ be such that $f_a+|m|\le f_c$ whenever $w_c<w_a$, and let $i,j$ satisfy
$|w_i-w_j|\ge h'$. Then $|\pcom_w(D_f)|\,(|u_{ij}|^22^m)^{h'}\le|\pcom_w(X_f)|$.
\end{lemma}
\begin{proof}
If $w_i<w_j$, replace $(w,f,X_f)$ by $(-w,-f,X_f^{-1}=u^*D_{-f}u)$: the hypotheses are preserved (with
the same $i,j$), $\pcom_w(X_f)=\pcom_{-w}(X_f^{-1})$ by {\cite[Lemma~3.18(a)]{P49}}, and $|\pcom_{-w}
(D_{-f})|=|\pcom_w(D_f)|$. So assume $w_j\le w_i-h'$. Let $w_0=\min w$, $w_1=\max w$, $S_t=\{a:w_a\ge
t\}$. From the definition, $\pcom_w(X)=(\det X)^{w_0}\prod_{w_0<t\le w_1}\Delta_{S_t}(X)$ for invertible
$X$, and $\det X_f=\det D_f$; all minors involved are positive. Since $w_0\le w_j\le w_i-h'$, the $h'$
integers $t\in(w_i-h',w_i]$ lie in $(w_0,w_1]$, and it suffices to show, for $w_0<t\le w_1$, $S=S_t$,
\[
g_t\,\Delta_S(D_f)\le\Delta_S(X_f),\qquad g_t=|u_{ij}|^22^m\text{ if }w_i-h'<t\le w_i,\quad g_t=1
\text{ otherwise}.
\]
$S^c=\{a:w_a<t\}$ is nonempty; let $c\in S^c$ minimise $f$ on $S^c$. For $a\in S$, $w_a\ge t>w_c$, so
$f_a+|m|\le f_c$. Put $g=\min(f,f_c\mathbf 1)$ entrywise: $g_a=f_a$ on $S$ and $g_a=f_c$ on $S^c$. Let
$Q=\{j\}$ if $w_i-h'<t\le w_i$ (then $i\in S$, $j\in S^c$) and $Q=\emptyset$ otherwise, and $p=S^c
\smallsetminus Q$. (1) $g_t\le2^{f_c|Q|}\Delta_Q(X_g^{-1})$: for $Q=\{j\}$, $(X_g^{-1})_{jj}=(u^*D_{-g}
u)_{jj}=\sum_a|u_{aj}|^22^{-g_a}\ge|u_{ij}|^22^{-f_i}$ and $f_c-f_i\ge|m|\ge m$. (2) By Jacobi, Fischer
($Q^c=p\sqcup S$), and the monotonicity of minors ($X_g\le X_{f_c\mathbf 1}=2^{f_c}\cdot1$ and $X_g\le X_f$
since $g\le f_c\mathbf1$, $g\le f$),
\[
\Delta_Q(X_g^{-1})=\frac{\Delta_{Q^c}(X_g)}{\det X_g}\le\frac{\Delta_p(X_g)\Delta_S(X_g)}{\det X_g}
\le\frac{2^{f_c|p|}\Delta_S(X_f)}{2^{\sum_{a\in S}f_a+|S^c|f_c}} .
\]
(3) As $|p|+|Q|=|S^c|$ and $\Delta_S(D_f)=2^{\sum_{a\in S}f_a}$, (1) and (2) give the claim.
\end{proof}

\begin{lemma}\label{pbr:compact}
For every $h'\in\N$ there is $\eta>0$ such that every $k\in\auU(n)\smallsetminus G\pbrB$ has an entry
$k_{ij}$ with $\varepsilon_i\ne\varepsilon_j$ and $\eta\,|\pcom_{-\pbrLo}(k^*Jk)|<|k_{ij}|^{2h'}$.
\end{lemma}
\begin{proof}
$C=\auU(n)\smallsetminus G\pbrB$ is compact (Corollary~\ref{pbr:flagG}), and $\phi(k)=\pcom_{-\pbrLo}
(k^*Jk)$ is continuous on $\auU(n)$. A unitary $k$ with $k_{ij}=0$ whenever $\varepsilon_i\ne
\varepsilon_j$ lies in $K\subset G\pbrB$. So each $k_0\in C$ has such an entry $k_{0,ij}\ne0$, and
$\{k:|k_{ij}|^{2h'}>\eta_0|\phi(k)|\}$, with $\eta_0=|k_{0,ij}|^{2h'}/(1+2|\phi(k_0)|)$, is an open
neighbourhood of $k_0$. Take a finite subcover and the least $\eta_0$.
\end{proof}

\begin{proposition}\label{pbr:Ebound}
Let $e\in\Z^n$. There is $C_e$ such that $|\pbrE(\operatorname{diag}(2^{-(ro+e)}))|\le C_e\,2^{-r(
\pbrlev(\mu)+h_0)}$ for all large integers $r$.
\end{proposition}
\begin{proof}
Let $C\ge2\max_a|e_a|$ be an integer, $r>C$, $f=-(ro+e)$, $A=D_f$ (positive definite, so accretive).
We apply Lemma~\ref{pbr:gain} with $w=-\mu$, $m=r-C$, $h'=h_0$. If $w_c<w_a$ then $\mu_a<\mu_c$, so
$\lambda_a<\lambda_c$ and $o_a\ge o_c+1$, whence $f_c-f_a=r(o_a-o_c)+e_a-e_c\ge r-C=|m|$. If
$\varepsilon_i\ne\varepsilon_j$ then $|\mu_i-\mu_j|=|\lambda_i-\lambda_j|+|\rho_i-\rho_j|\ge h_0$ ($\rho$
is ordered like $\lambda$). Also $|\pcom_{-\mu}(D_f)|=2^{\sum_a\mu_a(ro_a+e_a)}=2^{r\pbrlev(\mu)+e\cdot
\mu}$. Let $k\in\auU(n)\smallsetminus G\pbrB$ and $i,j$ as in Lemma~\ref{pbr:compact} with $h'=h_0$
(so $k_{ij}\ne0$). Then
\[
|F_A(k)|\le\frac{\eta^{-1}|k_{ij}|^{2h_0}}{2^{r\pbrlev(\mu)+e\cdot\mu}(|k_{ij}|^22^{r-C})^{h_0}}=
\eta^{-1}2^{Ch_0-e\cdot\mu}\,2^{-r(\pbrlev(\mu)+h_0)},
\]
and the domain of integration has mass at most $1$.
\end{proof}

For $\gamma\in\Z^n$ write $s^\gamma=\prod_a\varepsilon_a^{\gamma_a}\in\{\pm1\}$.

\begin{lemma}\label{pbr:low}
Let $P'$, $R$ be Laurent polynomials with $P'(x)=R(x)\,\pbrE(Jx)$ for all $x\in\pbrTp$, and let
$\gamma\in\Z^n$ satisfy $\pbrlev(\gamma)<\pbrlev(\mu)+\pbrlev(\beta)+h_0$ for every exponent $\beta$ of
$R$. Then $P'_\gamma=0$.
\end{lemma}
\begin{proof}
\emph{Continuation.} Let $f\in\Z^n$, $f\ne0$, and $\Sigma=\{z=\varrho e^{i\theta}:\varrho>0,\ |\theta|<
\pi/(2\max|f_a|)\}$. For $z\in\Sigma$, $\operatorname{diag}(z^{f_a})$ is accretive, and both sides of
\[
\sum_{\gamma'}P'_{\gamma'}s^{\gamma'}z^{f\cdot\gamma'}=\Bigl(\sum_\beta R_\beta s^\beta z^{f\cdot\beta}
\Bigr)\pbrE\bigl(\operatorname{diag}(z^{f_1},\dots,z^{f_n})\bigr)
\]
are holomorphic on the connected set $\Sigma$ (Lemma~\ref{pbr:FA}). For $z=e^{i\theta}\in\Sigma$ the
point $x$ with $x_a=\varepsilon_az^{f_a}$ lies in $\pbrTp$ and $Jx=\operatorname{diag}(z^{f_a})$, so
the identity holds on an arc through $1$, hence on $\Sigma\supset(0,\infty)$; for $f=0$ it holds
trivially.

\emph{Vanishing.} Suppose $P'_\gamma\ne0$, and let $\beta_0$ be the least level of an exponent of $P'$;
then $\beta_0\le\pbrlev(\gamma)<\pbrlev(\mu)+\pbrlev(\beta)+h_0$ for all exponents $\beta$ of $R$. Fix
$e\in\Z^n$ and take $z=\frac12$, $f=ro+e$ ($r\in\N$), multiplied by $2^{r\beta_0}$:
\[
\sum_{\gamma'}P'_{\gamma'}s^{\gamma'}2^{-e\cdot\gamma'}2^{-r(\pbrlev(\gamma')-\beta_0)}=\Bigl(\sum_\beta
R_\beta s^\beta2^{-e\cdot\beta}2^{-r(\pbrlev(\beta)-\beta_0)}\Bigr)\pbrE\bigl(\operatorname{diag}(
2^{-(ro+e)})\bigr).
\]
As $r\to\infty$ the left side tends to $\sum_{\pbrlev(\gamma')=\beta_0}P'_{\gamma'}s^{\gamma'}2^{-e\cdot
\gamma'}$, and by Proposition~\ref{pbr:Ebound} the right side is $O\bigl(\sum_\beta2^{-r(\pbrlev(\beta)
+\pbrlev(\mu)+h_0-\beta_0)}\bigr)\to0$. So $\sum_{\pbrlev(\gamma')=\beta_0}P'_{\gamma'}s^{\gamma'}
2^{-e\cdot\gamma'}=0$ for all $e\in\Z^n$. The maps $e\mapsto2^{-e\cdot\gamma'}$ are distinct characters
of $\Z^n$, hence linearly independent, so $P'_{\gamma'}=0$ whenever $\pbrlev(\gamma')=\beta_0$,
contradicting the choice of $\beta_0$.
\end{proof}

\subsubsection{Proof of Theorem~\ref{pbr:main}}\label{pbr:ss-proof}

Let $P$ be the Laurent polynomial of Proposition~\ref{pbr:laurent}, $\pbrtop=\kappa+N\pbrdD$, and
\[
Q=D^N\pbrAlt_{\pbrW}(\kappa),\qquad Q^{(1)}=D^N\pbrAlt_{S_n}(\kappa).
\]

\emph{Step 1: leading terms.} By Lemma~\ref{pbr:basic}(d), $(\Lambda+\pbrdc)+M\pbrdD=\kappa+N\pbrdD=
\pbrtop$, and $\pbrtop$ is ordered like $\lambda$ (Lemma~\ref{pbr:tau}). Lemma~\ref{pbr:top} with
$k=M$, $a=\Lambda+\pbrdc$ and the series $t$ (Proposition~\ref{pbr:laurent}(b),(c)) gives $P_\pbrtop=1$,
and every exponent $\gamma$ of $P$ has $w\in\pbrW$ with $\pbrlev(\gamma\circ w)\le\pbrlev(\pbrtop)$.
Lemma~\ref{pbr:top} with $k=N$, $a=\kappa$ and $t=\pbrAlt_{\pbrW}(\kappa)$ (whose exponents are the
$\kappa\circ w$, $w\in\pbrW$, as $\kappa$ has distinct entries) gives the same two statements for
$Q$. Both $P$ and $Q$ are $\pbrW$-antisymmetric.

\emph{Step 2: the relation on $\pbrTp$.} Let $x\in\pbrTp$. By Proposition~\ref{pbr:orbit},
$\pbrTh(x)$ is defined and $\pbrTh(x)=d\sigma_\lambda(\pbrNo\chi_\lambda(x)-\pbrE(Jx))$. By
Proposition~\ref{pbr:laurent}(a),(b) and Lemma~\ref{pbr:basic}(b), $P(x)=V_{\mathrm c}D^M\pbrTh(x)=
(-1)^{pq}VD^N\pbrTh(x)$, and $V\chi_\lambda=\pbrAlt_{S_n}(\kappa)$ (R2). Hence
\[
P(x)-s_0\,Q^{(1)}(x)=R(x)\,\pbrE(Jx),\qquad s_0=(-1)^{pq}d\,\sigma_\lambda\pbrNo,\quad R=-(-1)^{pq}
d\,\sigma_\lambda\,VD^N .
\]

\emph{Step 3: low levels.} Let $\pbrlev(\gamma)\le\pbrlev(\pbrtop)$. The exponents $\beta$ of $R$ have
$\pbrlev(\beta)\ge0$ (Lemma~\ref{pbr:basic}(c)), and $\pbrtop-\mu=o+N\pbrdD$ (Lemma~\ref{pbr:basic}(d)),
so by Lemma~\ref{pbr:tau}
\[
\pbrlev(\gamma)\le\pbrlev(\mu)+\pbrlev(o+N\pbrdD)<\pbrlev(\mu)+h_0\le\pbrlev(\mu)+\pbrlev(\beta)+h_0 .
\]
By Lemma~\ref{pbr:low}, $P_\gamma=s_0Q^{(1)}_\gamma$. Next, $Q^{(1)}_\gamma=\sum_\beta(D^N)_\beta
\sum_w\pbrsgn(w)$, the inner sum over the $w\in S_n$ with $\kappa\circ w=\gamma-\beta$. If such a $w$
is not in $\pbrW$, some $i$ has $\varepsilon_{w(i)}\ne\varepsilon_i$, so $|\kappa_i-\kappa_{w(i)}|=
|\lambda_i-\lambda_{w(i)}|\ge h_0$, and Lemma~\ref{pbr:basic}(e) ($o$ strictly reverses the order of
$\kappa$) gives $\pbrlev(\kappa\circ w)\ge\pbrlev(\kappa)+h_0$; but $\pbrlev(\kappa\circ w)=\pbrlev
(\gamma)-\pbrlev(\beta)\le\pbrlev(\pbrtop)=\pbrlev(\kappa)+\pbrlev(N\pbrdD)<\pbrlev(\kappa)+h_0$ by
Lemmas~\ref{pbr:basic}(c) and \ref{pbr:tau}. So only $w\in\pbrW$ occur, and $Q^{(1)}_\gamma=Q_\gamma$.
Thus $P_\gamma=s_0Q_\gamma$ whenever $\pbrlev(\gamma)\le\pbrlev(\pbrtop)$.

\emph{Step 4: $P=Q$.} Let $\gamma\in\Z^n$ with $P_\gamma\ne0$ or $Q_\gamma\ne0$. By Step 1 there is
$w\in\pbrW$ with $\pbrlev(\gamma\circ w)\le\pbrlev(\pbrtop)$, and by antisymmetry and Step 3,
$\pbrsgn(w)P_\gamma=P_{\gamma\circ w}=s_0Q_{\gamma\circ w}=s_0\pbrsgn(w)Q_\gamma$. So $P=s_0Q$, and
comparing the coefficients at $\pbrtop$ (Step 1), $s_0=1$ and $P=Q$.

\emph{Step 5: conclusion.} Let $t\in T$ have distinct eigenvalues. By Proposition~\ref{pbr:int},
$g\mapsto c(g^{-1}tg)$ is integrable, and by Proposition~\ref{pbr:laurent} and Step 4,
$(-1)^{pq}V(t)D(t)^N\pbrTh(t)=P(t)=D(t)^N\pbrAlt_{\pbrW}(\kappa)(t)$. As $D(t)\ne0\ne V(t)$, this is
\eqref{pbr:eq-unstable}. For $\sigma\in S_n$, $t^\sigma$ again has distinct eigenvalues,
$(t^\sigma)^\gamma=t^{\gamma\circ\sigma^{-1}}$ and $V(t^\sigma)=\pbrsgn(\sigma)V(t)$, so
\[
\sum_{\sigma\in S_n}\pbrTh(t^\sigma)=\frac{(-1)^{pq}}{V(t)}\sum_{\sigma\in S_n}\sum_{w\in\pbrW}
\pbrsgn(w\sigma^{-1})\,t^{\kappa\circ(w\sigma^{-1})}=(-1)^{pq}|\pbrW|\,\frac{\pbrAlt_{S_n}(\kappa)(t)}
{V(t)}=(-1)^{pq}p!\,q!\,\chi_\lambda(t),
\]
since for fixed $w$, $\sigma\mapsto w\sigma^{-1}$ is a bijection of $S_n$, and by (R2). Taking complex
conjugates gives \eqref{pbr:eq-stable}. \qed

\begin{corollary}[Formal degree]\label{pbr:degree}
Under \eqref{pbr:eq-gap}, $d=\dim V_\sigma=\prod_{i<j}(\lambda_i-\lambda_j)/(j-i)$ (entries listed in
decreasing order), and $\sigma_\lambda=(-1)^{pq}$.
\end{corollary}
\begin{proof}
By Step 4, $1=s_0=(-1)^{pq}\sigma_\lambda\,d/\dim V_\sigma$ with $d>0$; Weyl's dimension formula (R2)
for the highest weight $\pbrLo$ gives $\dim V_\sigma=\prod_{i<j}(\kappa_i-\kappa_j)/(j-i)$, and
$\kappa_i-\kappa_j=\lambda_i-\lambda_j$.
\end{proof}
This is Proposition~\ref{pds:haar} for parameters satisfying \eqref{pbr:eq-gap}, now without
appeal to Harish-Chandra's formula.

\begin{remark}[What is used]\label{pbr:rem-used}
The proof uses: {\cite[Theorem~3.17]{P49}}, {\cite[Lemma~3.25]{P49}}, {\cite[Lemma~3.18]{P49}},
{\cite[Corollary~3.19]{P49}}, {\cite[Proposition~3.21]{P49}} (also in the signature $(n,0)$),
Lemmas~\ref{pds:blockdet} and \ref{pm:kernel}(c), and the inputs (R1)--(R4). The gap condition
enters through {\cite[Lemma~3.25]{P49}} (with exponent $pqM\le h_0$), through Lemma~\ref{pbr:tau}
($4n^3<h_0$, $n^3+4n^5<h_0$), and through Proposition~\ref{pbr:Ebound}. Schmid's theorem (R3) enters
only in Lemma~\ref{pbr:mult}, to place the $K$-types of $\pi$ in $\Lambda$ plus the cone of positive
noncompact roots; this is what makes $\Lambda+\pbrdc$ the leading exponent of the series $t$.
\end{remark}



% ===== B-discrete-series.tex =====
% proofs/B-discrete-series.tex
% The discrete series coefficient at a real place: Theorem pbe:ds of proofs/B-bessel-expansion.tex,
% the chamber of the discrete series, the bound for kappa_tau, and the two definitions of A_0.
% Labels of this file start with pds:.
\providecommand{\auO}{\mathcal O}
\providecommand{\auPl}{\Phi}
\providecommand{\auAz}{\mathcal A_0}
\providecommand{\auU}{\mathrm U}
\providecommand{\auA}{\mathbb A}
\providecommand{\aurd}{\mathrm{rd}_E}
\providecommand{\aubar}[1]{\overline{#1}}
\providecommand{\auf}[1]{\texttt{#1}}
\providecommand{\pbvol}{\operatorname{vol}}
\providecommand{\pbdeg}{\operatorname{deg}}
\providecommand{\pbip}[2]{\langle#1,#2\rangle}
\providecommand{\pbeq}{\mathsf q}
\providecommand{\pbW}{\mathbb W}
\providecommand{\pbel}{\ell}
% macros of this file (prefix \pds)
\providecommand{\pdsK}{\mathbf K}
\providecommand{\pdsg}{\mathfrak g}
\providecommand{\pdsk}{\mathfrak k}
\providecommand{\pdsp}{\mathfrak p}
\providecommand{\pdsm}{m}
\providecommand{\pdsD}{\Delta^{\mathrm{min}}}
\providecommand{\pdsP}{P}
\providecommand{\pdsHS}[1]{\|#1\|_{2}}

\section{The discrete series coefficient at a real place}\label{pds:sec}

This section treats Theorem~\ref{pbe:ds} of \S\ref{pbe:sec}, the four properties
\textup{(DS1)}--\textup{(DS4)} of the coefficient $c_\tau$, together with three points of the
set-up of \S\ref{pbe:ss-ds}: which discrete series is meant (\S\ref{pds:ss-chamber}), the
inequality \eqref{pbe:eq-kappa} for $\kappa_\tau$ (\S\ref{pds:ss-kappa}), and the agreement of two
definitions of $\auAz$ (\S\ref{pds:ss-A0}).

\emph{What is proved and what is not.}
\begin{itemize}
\item[(1)] \textup{(DS1)(i),(ii)} are Schur's orthogonality relations (published); see
\S\ref{pds:ss-ds1}.
\item[(2)] \textup{(DS1)(iii)} is \emph{equivalent} to Harish-Chandra's formula for the formal degree
in the normalisation by the compact dual (Proposition~\ref{pds:haar}); for parameters with gap at
least $500N^3+10N^5$ this formula is proved in Corollary~\ref{pbr:degree}, without Harish-Chandra's
theorem; all parameters that occur have this gap ({\cite[Proposition~2.19]{P49}}), so (DS1)(iii) is
proved in full.
\item[(3)] \textup{(DS2)} is proved in full (Theorem~\ref{pds:decay}) from the explicit
integral formula \eqref{pds:eq-FJ} for $c_\tau$, which is proved in {\cite[Theorem~3.17]{P49}}
(the gap $h\ge2r^8$ exceeds $N^2+1$).
\item[(4)] \textup{(DS3)} is proved in full (Theorem~\ref{pds:tail}) from \textup{(DS2)}, the
definition of the measure in \textup{(DS1)(ii)}, and the value of the Casimir operator on a discrete
series (published). It does not use Harish-Chandra's formula.
\item[(5)] \textup{(DS4)} is proved in \S\ref{pbp:sec} (Theorem~\ref{pbp:main}); \S\ref{pds:ss-ds4}
only records how it is used.
\item[(6)] The chamber of \S\ref{pbe:ss-ds} is confirmed from the definition of the signs and of the
parameter (Lemma~\ref{pds:chamber}). The inequality \eqref{pbe:eq-kappa} holds for the generator
$\varpi$ of the different chosen in Lemma~\ref{pds:kappa}, not for an arbitrary purely imaginary
generator. The space $\auAz$ consists of the bounded solutions of the reproducing equation
({\cite[\S14.3]{P46}}); they are square-integrable forms of the required type
(Proposition~\ref{pds:A0}).
\end{itemize}

\subsection{Set-up}\label{pds:ss-setup}

Fix $\tau\in\auPl$ and drop it from the notation. Put $N=r=k+2\varrho+1$. As in
\S\ref{pbe:ss-real}, $Y_\tau$ has the normalised basis, in which the hermitian form has the matrix
$\mathsf J_1$, hermitian and unitary, and the norm $\|\cdot\|$ for which that basis is orthonormal.
Since $\mathsf J_1$ is hermitian and unitary, there is a unitary matrix $p$ with $p^{-1}\mathsf J_1p=
J:=\operatorname{diag}(\varepsilon_1,\dots,\varepsilon_N)$, $\varepsilon_i=\pm1$, with $k+\varrho$
signs $+$. Conjugation by $p$ preserves operator norms and carries $G_\tau$ onto
\[
G=\auU(J)=\{g\in\mathrm{GL}_N(\C):g^*Jg=J\},
\]
and $K_\tau=G_\tau\cap\auU(N)$ onto $\pdsK=G\cap\auU(N)$, the block diagonal elements of $G$ (an
element of $G$ is unitary if and only if it commutes with $J$). We work in $G$ from now on;
$\|g\|$ is the operator norm, and $\|g\|=\|g^{-1}\|\ge1$ because $g^{-1}=Jg^*J$. For a matrix $x$ let
\[
a(x)=\tfrac12(x+JxJ)
\]
be its \emph{block diagonal part} (the entries $x_{ij}$ with $\varepsilon_i=\varepsilon_j$). The Lie
algebra is $\pdsg=\{X:X^*J+JX=0\}=\pdsk\oplus\pdsp$ with $\pdsk=\pdsg\cap\{X=JXJ\}$ (block diagonal,
skew-hermitian) and $\pdsp=\pdsg\cap\{X=-JXJ\}$ (off the diagonal blocks, hermitian).

\subsubsection*{Parameters}
Let $\lambda=\varsigma_\tau\in\R^N$ be the Harish-Chandra parameter, its entries attached to the
coordinates as in \S\ref{pbe:ss-ds}; all entries are integers ($k$ is even and $\kappa(\tau)$,
$\Delta$ are integers), they are distinct, and $N-1$ is even, so $\lambda\in\frac{N-1}2+\Z^N$. The
positive roots are the $e_i-e_j$ with $\lambda_i>\lambda_j$. Put
\[
\rho_i=\tfrac12\bigl(\#\{j:\lambda_j<\lambda_i\}-\#\{j:\lambda_j>\lambda_i\}\bigr),\qquad
\mu=\lambda+\rho\in\Z^N,
\]
and let $\rho_{\mathrm c}$, $\rho_{\mathrm n}$ be the same half sums restricted to the $j$ with
$\varepsilon_j=\varepsilon_i$, respectively $\varepsilon_j\ne\varepsilon_i$. The highest weight of the
lowest $\pdsK$-type is $\Lambda=\lambda+\rho_{\mathrm n}-\rho_{\mathrm c}\in\Z^N$ (Blattner's
formula for the lowest $K$-type; e.g.\ \cite[Theorem~9.20(b), p.~310]{Knapp}, stated there for
linear connected semisimple groups; it applies to $\auU(p,q)=Z\cdot\mathrm{SU}(p,q)$, $Z\cong\auU(1)$
the compact centre, $Z\cap\mathrm{SU}(p,q)$ finite: a discrete series of $\auU(p,q)$ restricts
irreducibly to a discrete series of $\mathrm{SU}(p,q)$ with $Z$ acting by a character,
$\pdsK=Z\cdot(\pdsK\cap\mathrm{SU}(p,q))$, and $K$-types, their multiplicities and the lowest $K$-type
correspond under restriction of highest weights to $\mathfrak t\cap\mathfrak{su}(p,q)$; see the
paragraph ``On (D3)'' in {\cite[\S3.1.3]{P49}}). The \emph{gap} is
\[
h=\min\{|\lambda_i-\lambda_j|:\varepsilon_i\ne\varepsilon_j\}\ \ge\ \Delta=2r^8,
\]
by \S\ref{pbe:ss-ds}. We put $L=\max_i|\lambda_i|$; then $L\le\kappa(\tau)+m\Delta\le2b$ (comparison
of sizes; $\kappa(\tau)\le b$), $|\rho_i|\le N/2$ and $|\Lambda_i|\le L+N$.

\subsubsection*{Products of minors}
For $x\in M_N(\C)$ and $S\subseteq\{1,\dots,N\}$ let $\Delta_S(x)=\det(x_{ij})_{i,j\in S}$
($\Delta_\emptyset=1$). For $w\in\Z^N$ and $t\in\Z$ let $S_t(w)=\{i:w_i\ge t\}$, and for invertible
$x$ put
\begin{equation}\label{pds:eq-m}
\pdsm_w(x)=\prod_{t\in\Z}\Delta_{S_t(w)}(x)\cdot(\det x)^{-[t\le0]},
\end{equation}
where $[t\le 0]$ is $1$ or $0$; almost all factors are $1$. Shifting $t$ shows
\begin{equation}\label{pds:eq-shift}
\pdsm_{w+c\mathbf 1}(x)=(\det x)^c\,\pdsm_w(x)\qquad(c\in\Z).
\end{equation}
For $B\in\N$ with $|w_i|<B$ for all $i$ we also use the finite product $\pdsP^B_w(x)=\prod_{-B<t\le B}
\Delta_{S_t(w)}(x)$; it equals $\pdsm_w(x)$ when $\det x=1$.

\subsection{The explicit formula}\label{pds:ss-FJ}

Every $g\in G$ is $g=kz$ with $k\in\pdsK$ and $z\in G$ positive definite: $z=(g^*g)^{1/2}$ lies in
$G$ because $J(g^*g)J=(g^*g)^{-1}$ forces $JzJ=z^{-1}$ (uniqueness of positive square roots), and
$k=gz^{-1}$ is unitary. Let $dk$ be the Haar measure of mass $1$ on $\pdsK$ and
$d_\Lambda=(\int_{\pdsK}|\pdsm_\Lambda|^2)^{-1}$.
\begin{equation}\label{pds:eq-FJ}
c_\tau(kz)=d_\Lambda\int_{\pdsK}\frac{\pdsm_\Lambda(ku)\,\aubar{\pdsm_\Lambda(u)}}{\pdsm_{-\mu}(u^{-1}zu)}
\,du\qquad(k\in\pdsK,\ z\in G\text{ positive definite}).
\end{equation}
(The denominator is positive by Lemma~\ref{pds:fischer} below.)

Formula \eqref{pds:eq-FJ} is {\cite[Theorem~3.17]{P49}}, applied with $n=N$; its hypothesis $h\ge N^2+1$
holds since $h\ge2r^8$. (This is Flensted-Jensen's description of the coefficient of the lowest
$K$-type~\cite{FJ}; the proof in {\cite[\S3.2]{P49}} does not rely on~\cite{FJ}.) Only \textup{(DS2)} uses
\eqref{pds:eq-FJ}; \textup{(DS1)}, \textup{(DS3)} use only that $c_\tau$ is the coefficient of a unit
vector of the lowest $K$-type.

\subsection{The chamber}\label{pds:ss-chamber}

\begin{lemma}\label{pds:chamber}
The signs of \eqref{pm:eq-signs} on $Y_0$ are $\varepsilon_i=(-1)^{\max(0,\,i+1-k)}$ for the
coordinate numbered $i=0,\dots,r-1$, and the doubled parameter is twice
$\lambda_i=\frac k2-1-i$ for $i<k-1$ and $\lambda_i=-\kappa(\tau)-(i+2-k)\Delta$ for $i\ge k-1$.
Hence the entries $\frac k2-1,\dots,1-\frac k2$ and the entries $-\kappa(\tau)-j\Delta$ with $j$ odd
sit at coordinates of sign $+$, those with $j$ even at coordinates of sign $-$; there are $k+\varrho$
signs $+$ and $\varrho+1$ signs $-$. On the $k-1$ coordinates of the first group $\Lambda=0$, and on
the other $m=2\varrho+2$ coordinates, in the order of decreasing $\lambda$, the signs are $+,-,+,
\dots,-$.
\end{lemma}
\begin{proof}
The formulas restate the signs \eqref{pm:eq-signs} and the definition of $\varsigma_\tau$ in
{\cite[\S14.3]{P46}} (the entry $i\ge k-1$ is $-\kappa(\tau)-j\Delta$ with $j=i+2-k$). With $j=i+2-k$
the sign is $(-1)^{j-1}$, which is the first claim;
the coordinates $i\le k-1$ give $k$ signs $+$ and the $m-1$ coordinates $i\ge k$ give $\varrho$
signs $+$ and $\varrho+1$ signs $-$. For $i<k-1$: the entries below $\lambda_i$ are the $k-2-i$
later entries of the first group and all $m$ entries of the second, the entries above are $i$; the
coordinates of sign $-$ are the $\varrho+1$ entries with $j$ even, all below. So $2\rho_{\mathrm n,i}
=\varrho+1$ and $2\rho_{\mathrm c,i}=(k-2-i)+(\varrho+1)-i$, and $2\Lambda_i=(k-2-2i)+(\varrho+1)-
(k-2-2i+\varrho+1)=0$. The last claim restates the first.
\end{proof}
This agrees with what \S\ref{pbe:ss-ds} states, and with the structure used for the
Laplace method of \textup{(DS4)} in \S\ref{pbp:sec} (weights $0$ on $k-1$ coordinates of sign $+$, then $2\varrho+2$
strictly decreasing weights in $[-3B,-B]$ with alternating signs starting with $+$, for
$B=\kappa(\tau)+\Delta/2$): indeed $\Lambda_i=\lambda_i+O(N)$ there, and $\kappa(\tau)+m\Delta+N\le
3B$ because $m\Delta\le\kappa(\tau)$ (comparison of sizes).

\subsection{(DS1)}\label{pds:ss-ds1}

\emph{(i)} $c_\tau(g)=\pbip{v_0}{\pi(g)v_0}$ is continuous because $\pi$ is strongly continuous;
$c_\tau(g^{-1})=\pbip{\pi(g)v_0}{v_0}=\aubar{c_\tau(g)}$; $|c_\tau|\le\|v_0\|^2=1$.

\emph{(ii)} $\pi$ is irreducible and square-integrable, so for any Haar measure $da$ there is
$d_\pi>0$ with $\int_G\pbip{u'}{\pi(a)v'}\allowbreak \aubar{\pbip{u}{\pi(a)v}}da=d_\pi^{-1}\pbip{u'}{u}
\pbip v{v'}$ for all $u,u',v,v'$ (Schur orthogonality, \cite[Proposition~9.6, p.~284]{Knapp}; written for inner products conjugate-linear in the first variable). With all four
vectors $v_0$, $d_\pi^{-1}=\int|c_\tau|^2$; so for the measure with $\int|c_\tau|^2=\pbdeg_\tau^{-1}$
we have $d_\pi=\pbdeg_\tau$, and with $u=v=v'=v_0$, $u'=\pi(g)^{-1}v_0$:
$\int_G\pbdeg_\tau\aubar{c_\tau(a)}\,c_\tau(ga)\,da=\pbip{v_0}{\pi(g)v_0}$, using $c_\tau(ga)=
\pbip{\pi(g)^{-1}v_0}{\pi(a)v_0}$. This is \textup{(DS1)(ii)}; the integrability of $c_\tau$ assumed
there is proved below (Theorem~\ref{pds:tail}).

\emph{(iii)} Let $\mathfrak u=\pdsk\oplus i\pdsp=\mathfrak u(N)$, the compact dual. A Haar measure on
$G$ (on $\auU(N)$) is determined by a Lebesgue measure on $\pdsg$ (on $\mathfrak u$) through the
exponential map near $0$; the measures are \emph{compatible} if the Lebesgue measures correspond
under $X+Y\mapsto X+iY$ ($X\in\pdsk$, $Y\in\pdsp$).

\begin{proposition}\label{pds:haar}
Let $dg$ be the Haar measure on $G$ compatible with the Haar measure of mass $1$ on $\auU(N)$. Then
$\int_G|c_\tau|^2dg=\pbdeg_\tau^{-1}$ with $\pbdeg_\tau=\prod_{i<j}(\lambda_i-\lambda_j)/(j-i)$
(entries in decreasing order). Thus the measures of \textup{(DS1)(ii)} and \textup{(DS1)(iii)} agree.
\end{proposition}
\begin{proof}
By (ii), $\int|c_\tau|^2dg=d_\pi^{-1}$ for the formal degree $d_\pi$ relative to $dg$, and
$d_\pi=\prod_{i<j}(\lambda_i-\lambda_j)/(j-i)$ by Corollary~\ref{pbr:degree}, which applies because
the parameter satisfies the gap condition \eqref{pbr:eq-gap} ({\cite[Proposition~2.19]{P49}}). This is
Harish-Chandra's formula $d_\pi=\prod_{\alpha>0}|\pbip\lambda\alpha|/\pbip\rho\alpha$ in the
normalisation by the compact dual; \cite[Theorem~16, p.~96]{HC66} gives it as $c^{-1}[W_G]|\varpi(\lambda)|$ with a constant $c$ that is not
computed there for this measure, and the proof here does not use it.
\end{proof}

\begin{remark}\label{pds:cayley}
The Haar measure of Proposition~\ref{pds:haar} can also be described by the Cayley chart
$y\mapsto2(1-\frac12J^{-1}S_y)^{-1}-1$ with density $(|\det J|\cdot|\det(1-\frac12J^{-1}S_y)|^2)^{-N}$,
divided by the constant that makes the mass $1$ for $J=1$; $S_y$ runs over the skew-hermitian matrices. The differential of the chart at $0$ is $X\mapsto2X$ on $\pdsg$ written as
$X=\frac12J^{-1}S$, and $S=2JX$; the map $X\mapsto JX$ from $\pdsg$ to $\mathfrak u(N)$ and the map
$X+Y\mapsto X+iY$ both act on the off-diagonal blocks by an isometry of the real coordinates and on
$\pdsk$ by $\pm1$ on each block, so both have Jacobian $1$. Hence this measure is the one of
Proposition~\ref{pds:haar}, and for this measure the identity $\pbdeg=(\int|c|^2)^{-1}=\prod_{i<j}
(\lambda_i-\lambda_j)/(j-i)$ is the same statement.
\end{remark}

\subsection{(DS2): decay}\label{pds:ss-ds2}

\begin{lemma}\label{pds:blockdet}
For $g\in G$: $a(g)$ is invertible, $|\det a(g)|\ge(\|g\|/2)^2$, and $|\det a(kgk')|=|\det a(g)|$ for
$k,k'\in\pdsK$. If $g$ is positive definite, so is $a(g)=\frac12(g+g^{-1})$.
\end{lemma}
\begin{proof}
Elements of $\pdsK$ commute with $J$, so $a(kgk')=k\,a(g)k'$ with $|\det k|=|\det k'|=1$. Write
$g=kz$ as in \S\ref{pds:ss-FJ}; then $\|g\|=\|z\|$ and $a(z)=\frac12(z+JzJ)=\frac12(z+z^{-1})$.
The spectrum of $z$ is that of $z^{-1}=JzJ$, so with $x_1\ge\dots\ge x_N>0$ the eigenvalues,
$x_N=x_1^{-1}$ and $x_1=\|z\|$. Thus $\det a(z)=\prod_i\frac12(x_i+x_i^{-1})$, a product of factors
$\ge1$; if $N\ge2$ and $x_1\ne1$ the two factors for $x_1,x_N$ give $\det a(z)\ge\frac14(x_1+
x_1^{-1})^2\ge(\|z\|/2)^2$, and if $x_1=1$ then $\det a(z)=1\ge\frac14$. The last claim: $g^{-1}=
JgJ$ when $g=g^*$, and $g^{-1}$ is positive definite.
\end{proof}

\begin{lemma}\label{pds:fischer}
Let $w\in G$ be positive definite and $0\le c\le h$ an integer. Then $\pdsm_{-\mu}(w)=\pdsm_\mu(w)>0$
and $\det a(w)^c\le\pdsm_\mu(w)^2$.
\end{lemma}
\begin{proof}
We use three facts. (F1) \emph{Fischer's inequality}: for positive definite $x$ and disjoint $T,T'$,
$\Delta_{T\cup T'}(x)\le\Delta_T(x)\Delta_{T'}(x)$ \cite[Theorem~7.8.5]{HornJohnson}, applied to the
principal submatrix of $x$ on $T\cup T'$.
(F2) \emph{Jacobi}: $\Delta_S(x^{-1})=\Delta_{S^c}(x)/\det x$. (F3) $\det w=1$ (it is positive and
$|\det w|=1$), and $\Delta_S(w^{-1})=\Delta_S(JwJ)=\Delta_S(w)$; with (F2), $\Delta_S(w)=
\Delta_{S^c}(w)$ for all $S$.

By (F3), $\pdsm_{-\mu}(w)=\prod_t\Delta_{\{\mu_i\le-t\}}(w)=\prod_s\Delta_{S_s(\mu)^c}(w)=\prod_s
\Delta_{S_s(\mu)}(w)=\pdsm_\mu(w)$, and this is positive since principal minors of positive definite
matrices are positive. Let $v=\mu+c\mathbf 1$, choose $B$ with $|\mu_i|,|v_i|<B$, and write $a=a(w)$,
positive definite by Lemma~\ref{pds:blockdet}. Fix $t$ and put $S=S_t(v)\supseteq T=S_t(\mu)$ and
$D=S\setminus T=\{i:\mu_i<t\le\mu_i+c\}$. \emph{All $i\in D$ have the same sign}: for $i,j\in D$,
$|\mu_i-\mu_j|\le c-1$; but $\rho_i-\rho_j$ has the sign of $\lambda_i-\lambda_j$, so $|\mu_i-\mu_j|
\ge|\lambda_i-\lambda_j|\ge h\ge c$ if $\varepsilon_i\ne\varepsilon_j$. Hence $a$ and $w$ have the same
entries on $D\times D$, and
\[
\Delta_S(a)\overset{\textup{(F1)}}\le\Delta_T(a)\Delta_D(a)=\Delta_T(a)\Delta_D(w)
\overset{\textup{(F3)}}=\Delta_T(a)\Delta_{T\cup S^c}(w)\overset{\textup{(F1)}}\le\Delta_T(a)
\Delta_T(w)\Delta_{S^c}(w).
\]
Multiply over $-B<t\le B$. On the left, $\pdsP^B_v(a)=(\det a)^c\pdsP^B_\mu(a)$, as in
\eqref{pds:eq-shift} (the $c$ new factors have $S_t(\mu)$ equal to everything, the $c$ lost ones
have $S_t(\mu)=\emptyset$). On the right, $\prod_t\Delta_{S_t(v)^c}(w)=\prod_t\Delta_{S_t(v)}(w)=
\pdsm_v(w)=\pdsm_\mu(w)$ by (F3) and \eqref{pds:eq-shift}. So $(\det a)^c\pdsP^B_\mu(a)\le
\pdsP^B_\mu(a)\,\pdsm_\mu(w)^2$, and $\pdsP^B_\mu(a)>0$.
\end{proof}

\begin{theorem}[\textup{(DS2)}]\label{pds:decay}
For every $g\in G$ and every integer $0\le c\le h$ (formula \eqref{pds:eq-FJ} holds by {\cite[Theorem~3.17]{P49}}),
\[
|c_\tau(g)|\le|\det a(g)|^{-c/2}\le\bigl(2/\|g\|\bigr)^c .
\]
In particular $|c_\tau(g)|\le(2/\|g\|)^{j+\pbeq}$ for $0\le j\le r^7$, $\pbeq=9r^2$, and for $g\in
G_\tau$ with the norm of \S\ref{pbe:ss-real}.
\end{theorem}
\begin{proof}
Write $g=kz$ and let $D=|\det a(g)|=\det a(z)$. For $u\in\pdsK$, $w_u=u^{-1}zu\in G$ is positive
definite and $\det a(w_u)=D$ (Lemma~\ref{pds:blockdet}), so $\pdsm_{-\mu}(w_u)\ge D^{c/2}$ by
Lemma~\ref{pds:fischer}. By \eqref{pds:eq-FJ} and $2|xy|\le|x|^2+|y|^2$,
\[
|c_\tau(g)|\le\frac{d_\Lambda}{D^{c/2}}\int_{\pdsK}\frac{|\pdsm_\Lambda(ku)|^2+|\pdsm_\Lambda(u)|^2}2du
=D^{-c/2},
\]
by the invariance of $du$. The second inequality is Lemma~\ref{pds:blockdet}. For the last
sentence, $j+9r^2\le r^7+9r^2\le\Delta\le h$, and conjugation by $p$ preserves the norm.
\end{proof}
The bound holds for all exponents up to the gap, which is more than \textup{(DS2)} asks; nothing
else in \S\ref{pbe:sec} needs to change.

\subsection{(DS3): the formal degree against the volume of balls}\label{pds:ss-ds3}

Let $dg$ be the Haar measure of \textup{(DS1)(ii)}, $\int_G|c_\tau|^2dg=\pbdeg_\tau^{-1}$. We use one
published fact. Let $E_{ab}$ be the matrix units, acting on the $\pdsK$-finite vectors of $\pi$
through the complexified differential $d\pi$, and $\Omega=\sum_{a,b}E_{ab}E_{ba}$,
$\Omega_{\pdsk}=\sum_{\varepsilon_a=\varepsilon_b}E_{ab}E_{ba}$. Then
\begin{equation}\label{pds:eq-casimir}
d\pi(\Omega)=|\lambda|^2-|\rho|^2,\qquad d\pi(\Omega_{\pdsk})v_0=\bigl(|\Lambda+\rho_{\mathrm c}|^2-
|\rho_{\mathrm c}|^2\bigr)v_0 .
\end{equation}
The first is the value of the Casimir operator on a representation with infinitesimal character
$\lambda$ (Harish-Chandra isomorphism; \cite[Theorem~8.18, (8.32)]{Knapp}); the second is the same computation for the irreducible representation of highest
weight $\Lambda$ of $\pdsK\cong\auU(k+\varrho)\times\auU(\varrho+1)$.

\begin{lemma}\label{pds:lip}
Put $C_0=3N(L+2N)$. If $g\in G$ and $\|g-1\|\le\delta:=1/(4C_0)$, then $|c_\tau(g)|\ge\frac12$.
\end{lemma}
\begin{proof}
\emph{Step 1: $\|d\pi(X)v_0\|\le C_0\|X\|$ for $X\in\pdsg$.} For $\varepsilon_a\ne\varepsilon_b$ the
matrices $E_{ab}+E_{ba}$ and $i(E_{ab}-E_{ba})$ lie in $\pdsg$, and $d\pi$ of an element of $\pdsg$ is
skew-adjoint; hence $d\pi(E_{ab})^*=-d\pi(E_{ba})$ on $\pdsK$-finite vectors, and
$\pbip{v_0}{d\pi(E_{ab}E_{ba})v_0}=-\|d\pi(E_{ba})v_0\|^2$. By \eqref{pds:eq-casimir},
\[
\sum_{\varepsilon_a\ne\varepsilon_b}\|d\pi(E_{ab})v_0\|^2=|\Lambda+\rho_{\mathrm c}|^2-
|\rho_{\mathrm c}|^2-|\lambda|^2+|\rho|^2\le N(L+2N)^2+N^3/4\le2N(L+2N)^2 .
\]
For $Y=\sum_{\varepsilon_a\ne\varepsilon_b}Y_{ab}E_{ab}\in\pdsp$, Cauchy--Schwarz gives $\|d\pi(Y)v_0\|
\le\pdsHS Y\,(2N)^{1/2}(L+2N)\le\sqrt2\,N(L+2N)\|Y\|$. For $X\in\pdsk$, $d\pi(X)v_0$ is the action of
$X$ in the $\pdsK$-type $\Lambda$; conjugating $X$ in $\pdsK$ to $\operatorname{diag}(i\theta)$,
$|\theta_a|\le\|X\|$, its eigenvalues there are $i\pbip\nu\theta$ for weights $\nu$ in the convex hull
of the Weyl group orbit of $\Lambda$, so its norm is at most $\sum_a|\Lambda_a|\,\|X\|\le N(L+N)\|X\|$.
A general $X\in\pdsg$ is $X_{\pdsk}+X_{\pdsp}$ with $X_{\pdsk}=a(X)$, and both parts have norm at
most $\|X\|$.

\emph{Step 2.} Let $\|g-1\|\le\delta\le\frac12$ and $X=\log g=\sum_{n\ge1}(-1)^{n+1}(g-1)^n/n$. Then
$\exp X=g$ and $\|X\|\le-\log(1-\delta)\le2\delta$. Also $X\in\pdsg$: $JX^*J=\log(Jg^*J)=\log(g^{-1})=
-X$, by the holomorphic functional calculus (the series for $\log(g^{-1})$ converges because
$\|g^{-1}-1\|=\|Jg^*J-1\|=\|g-1\|$). The vector $v_0$ is $\pdsK$-finite, hence smooth, so
$t\mapsto\pi(\exp tX)v_0$ has derivative $\pi(\exp tX)d\pi(X)v_0$ and
\[
|c_\tau(g)-1|=|\pbip{v_0}{\pi(g)v_0-v_0}|\le\int_0^1\|d\pi(X)v_0\|\,dt\le2C_0\delta=\tfrac12 .
\qedhere
\]
\end{proof}

\begin{lemma}\label{pds:balls}
For $T\ge1$ let $S_T=\{g\in G:\|g\|\le T\}$ and $B_\delta=\{g:\|g-1\|\le\delta\}$ with $\delta$ as in
Lemma~\ref{pds:lip}. Then $\pbvol(S_T)\le(5T^2/\delta)^{2N^2}\pbvol(B_\delta)\le4(5T^2/\delta)^{2N^2}
\pbdeg_\tau^{-1}$.
\end{lemma}
\begin{proof}
Choose a maximal family $g_1,\dots,g_\nu$ in $S_T$ with $\|g_ig_j^{-1}-1\|>\delta$ for $i\ne j$ (it is
finite by the estimate that follows, $S_T$ being compact). Since $g_i-g_j=(g_ig_j^{-1}-1)g_j$ and
$\|g_j^{-1}\|=\|g_j\|\le T$, $\|g_i-g_j\|>\delta/T$. The balls of radius $\delta/(2T)$ about the $g_i$
in $M_N(\C)$ (real dimension $2N^2$) are disjoint and lie in the ball of radius $2T$ about $0$, so
$\nu\le(4T^2/\delta)^{2N^2}$. By maximality every $g\in S_T$ has $gg_i^{-1}\in B_\delta$ for some $i$,
so $S_T\subseteq\bigcup_iB_\delta g_i$ and, $G$ being unimodular, $\pbvol(S_T)\le\nu\,\pbvol(B_\delta)$.
Finally $1\le4|c_\tau|^2$ on $B_\delta$ (Lemma~\ref{pds:lip}), so $\pbvol(B_\delta)\le4\int_G|c_\tau|^2
=4\pbdeg_\tau^{-1}$.
\end{proof}

\begin{theorem}[\textup{(DS3)}]\label{pds:tail}
For every real $\mathsf s>4N^2$, $\pbdeg_\tau\int_G\|g\|^{-\mathsf s}dg\le4(5/\delta)^{2N^2}\mathsf s/
(\mathsf s-4N^2)$. In particular $c_\tau$ is integrable, and
\[
\pbdeg_\tau\int_{G_\tau}\bigl(32/\|g\|\bigr)^{\pbeq}dg\le8\cdot32^{\pbeq}\bigl(60N(L+2N)\bigr)^{2N^2}
\le b^{99r^2}\qquad(\pbeq=9r^2).
\]
\end{theorem}
\begin{proof}
As $\|g\|\ge1$, $\|g\|^{-\mathsf s}=\int_1^\infty\mathsf sT^{-\mathsf s-1}[\|g\|\le T]\,dT$; by Fubini
and Lemma~\ref{pds:balls},
\[
\int_G\|g\|^{-\mathsf s}dg=\int_1^\infty\mathsf sT^{-\mathsf s-1}\pbvol(S_T)\,dT\le
\frac{4}{\pbdeg_\tau}\Bigl(\frac5\delta\Bigr)^{2N^2}\int_1^\infty\mathsf sT^{4N^2-\mathsf s-1}dT .
\]
Integrability of $c_\tau$ follows with Theorem~\ref{pds:decay} ($c=9N^2\le h$). For $\mathsf s=
\pbeq=9N^2$, $\mathsf s/(\mathsf s-4N^2)=9/5<2$ and $5/\delta=20C_0=60N(L+2N)$, which gives the
middle expression after multiplying by $32^{\pbeq}$. Lastly $N=r$, $L\le2b$ and $r\le b$ give
$60N(L+2N)\le b^3$, so the middle expression is at most $8\cdot2^{45r^2}b^{6r^2}\le b^{99r^2}$.
\end{proof}

\subsection{(DS4): the Bessel period at a stationary point}\label{pds:ss-ds4}

\smallskip\noindent\textup{(DS4)} is proved in \S\ref{pbp:sec} (Theorem~\ref{pbp:main}).\par\smallskip

\subsection{The bound for \texorpdfstring{$\kappa_\tau$}{kappa}}\label{pds:ss-kappa}

The inequality \eqref{pbe:eq-kappa} is not true for every purely imaginary generator $\varpi$ of the
different: $\varepsilon\varpi$, with $\varepsilon$ a unit of $E^+$ having one large conjugate,
violates it. It holds for the generator of Lemma~\ref{pds:kappa}, which is the one fixed in
\eqref{pb:eq-psi}. By {\cite[Lemma~13.4(S7)]{P46}}, $q_1=q_0t$ with every prime factor of $t$ dividing
$q_0$ (so that $\Phi_{q_1}(X)=\Phi_{q_0}(X^t)$ for the cyclotomic polynomials), and $(q_0/2)^{q_0}\le
q_0^{q_0}\le x$.

\begin{lemma}\label{pds:kappa}
Let $\zeta\in E$ be a primitive $q_1$-th root of unity and $\varpi=\gamma\,\Phi_{q_1}'(\zeta)\sqrt{-d}$,
where $\gamma\in\auO_E^\times$ satisfies $\gamma\aubar\gamma=1$ and $\aubar{\gamma\Phi'_{q_1}(\zeta)}=
\gamma\Phi'_{q_1}(\zeta)$ \textup{(}such $\gamma$ exist, e.g.\ $\gamma=\zeta^s$ with
$s=q_1/4+(q_1-\varphi(q_1))/2+1$; see the first paragraph of the proof\textup{)}. Then $\varpi$ is a purely imaginary generator of the different of $E$, and
for every $\tau\in\auPl$, $2\le|\tau(\varpi)|\le b$ and $1\le M|\tau(\varpi)|/2\pi\le b^{23}$.
\end{lemma}
\begin{proof}
\emph{Existence of $\gamma$.} Write $m=q_1$ and $u=\Phi_m'(\zeta)$, and let $g(X)=\prod_{e\mid m,\,e<m}
\Phi_e(X)$, so $X^m-1=\Phi_m(X)g(X)$, $g(\zeta)\ne0$, and differentiating at $\zeta$ gives
$u=m\zeta^{-1}/g(\zeta)$. Complex conjugation of the CM field $E$ sends $\zeta$ to $\zeta^{-1}$ and
fixes the rational coefficients, so $\aubar u=m\zeta/g(\zeta^{-1})$. Since $\Phi_1(1/X)=-X^{-1}\Phi_1(X)$
and $X^{\varphi(e)}\Phi_e(1/X)=\Phi_e(X)$ for $e\ge2$, we get $g(\zeta^{-1})=-\zeta^{-D}g(\zeta)$ with
$D=\deg g=m-\varphi(m)$, hence $\aubar u=-\zeta^{D+2}u=\zeta^{m/2+D+2}u=\zeta^{2s}u$ with
$s=m/4+D/2+1\in\mathbb Z$ (as $16\mid m$ and $\varphi(m)$ is even). For $\gamma=\zeta^s$:
$\gamma\aubar\gamma=1$ and $\aubar{\gamma u}=\zeta^{-s}\zeta^{2s}u=\gamma u$.

Since $d\equiv q_0/2-1\pmod{q_0}$ and $16\mid q_0$, $d\equiv-1$ modulo every prime factor of $q_0$ and
$d\equiv3\pmod 4$; so $d$ is prime to $q_1$, $\auO_K=\Z[\frac{1+\sqrt{-d}}2]$ has different
$(\sqrt{-d})$, and the discriminants of $K$ and $\Q(\zeta)$ are coprime. Hence $\auO_E=\auO_K[\zeta]$,
the different of $E/K$ is $(\Phi'_{q_1}(\zeta))$, and the different of $E/\Q$ is the product of those of
$E/K$ and $K/\Q$.

(\emph{$\auO_E=\auO_K[\zeta]$.} Let $F=\Q(\zeta)$; its ring of integers is $\Z[\zeta]$, with
$\Z$-basis $\beta_j=\zeta^{j}$, $0\le j<n=\varphi(q_1)$, and discriminant $d_F=\pm\auNm_{F/\Q}
\Phi'_{q_1}(\zeta)$~\cite[Chapter~I, \S2 and Proposition~(10.2)]{Neukirch}. Differentiating $X^{q_1}-1=\Phi_{q_1}(X)g(X)$ gives
$q_1\zeta^{q_1-1}=\Phi'_{q_1}(\zeta)g(\zeta)$, so $\Phi'_{q_1}(\zeta)$ divides $q_1$ and every prime
factor of $d_F$ divides $q_1$; as $d_K=-d$ is prime to $q_1$, $\gcd(d_K,d_F)=1$. Since $K$ is ramified
at $d$ and $F$ is not, $K\cap F=\Q$, so $[E:\Q]=2n$, $\Phi_{q_1}$ is irreducible over $K$, the
$\beta_j$ are a $K$-basis of $E$, and $\omega_1=1,\omega_2=\frac{1+\sqrt{-d}}2$ are an $F$-basis
of $E$. Let $x\in\auO_E$ and write $x=\sum_{i,j}c_{ij}\omega_i\beta_j$ with $c_{ij}\in\Q$. Put
$c_j=\sum_ic_{ij}\omega_i\in K$. The $n$ embeddings $\sigma$ of $E$ over $K$ restrict to the $n$
embeddings of $F$; solving $\sigma(x)=\sum_jc_j\sigma(\beta_j)$ by Cramer's rule gives
$c_j=\gamma_j/\delta$ with $\delta=\det(\sigma\beta_j)$, $\delta^2=d_F$, and $\gamma_j$ an algebraic
integer. So $d_Fc_j=\delta\gamma_j\in\auO_K$, i.e.\ $d_Fc_{ij}\in\Z$. The same argument with the
roles of $K$ and $F$ exchanged (embeddings of $E$ over $F$, basis $\omega_i$, $\det^2=d_K$) gives
$d_Kc_{ij}\in\Z$. Hence $c_{ij}\in\Z$, and $\auO_E=\auO_K[\zeta]$.

\emph{The different of $E/K$.} Let $A=\auO_K$, $B=\auO_E=A[\zeta]$, $f=\Phi_{q_1}$ (monic of degree
$n$, irreducible over $K$), and write $f(X)/(X-\zeta)=b_0+b_1X+\dots+b_{n-1}X^{n-1}$ with $b_k\in
B$. Lagrange interpolation at the $n$ distinct roots $\sigma(\zeta)$ gives, for $0\le i<n$,
$\sum_\sigma\frac{f(X)}{X-\sigma\zeta}\frac{\sigma(\zeta)^i}{f'(\sigma\zeta)}=X^i$, since both
sides have degree $<n$ and agree at every $\sigma(\zeta)$. Comparing coefficients of $X^k$:
$\mathrm{Tr}_{E/K}\bigl(\zeta^ib_k/f'(\zeta)\bigr)=\delta_{ik}$ (Euler's formula). So
$\{b_k/f'(\zeta)\}$ is the basis dual to $\{\zeta^i\}$ under the trace, and since $B$ is free over
$A$ on $\{\zeta^i\}$, $\{x:\mathrm{Tr}_{E/K}(xB)\subseteq A\}=\sum_kA\,b_k/f'(\zeta)$. Comparing
coefficients in $f(X)=(X-\zeta)\sum b_kX^k$ gives $b_{n-1}=1$ and $b_{k-1}=\zeta b_k+a_k$ ($a_k\in A$
the coefficients of $f$), so $b_{n-1-m}\in\zeta^m+\sum_{l<m}A\zeta^l$; the $b_k$ span $A[\zeta]=B$
over $A$. Hence $\mathfrak D_{E/K}^{-1}=f'(\zeta)^{-1}B$, i.e.\ $\mathfrak D_{E/K}=(\Phi'_{q_1}(\zeta))$.)

(Transitivity of the different: for a finite extension $B/A$ of rings of integers
of number fields, $\mathfrak D_{B/A}^{-1}=\{x:\mathrm{Tr}_{B/A}(xB)\subseteq A\}$, and an $A$-submodule $M$ of
$\mathrm{Frac}(A)$ satisfies $\mathrm{Tr}_{\mathrm{Frac}(A)/\Q}(M)\subseteq\Z$ iff $M\subseteq\mathfrak D_{A/\Z}^{-1}$. As
$\mathrm{Tr}_{E/\Q}=\mathrm{Tr}_{K/\Q}\circ\mathrm{Tr}_{E/K}$ and $\mathrm{Tr}_{E/K}(x\auO_E)$ is an
$\auO_K$-module, $x\in\mathfrak D_{E/\Q}^{-1}$ iff $\mathrm{Tr}_{E/K}(x\auO_E)\subseteq\mathfrak
D_{K/\Q}^{-1}$ iff $\mathrm{Tr}_{E/K}(x\mathfrak D_{K/\Q}\auO_E)\subseteq\auO_K$ iff $x\mathfrak
D_{K/\Q}\auO_E\subseteq\mathfrak D_{E/K}^{-1}$; so $\mathfrak D_{E/\Q}^{-1}=\mathfrak D_{E/K}^{-1}
\mathfrak D_{K/\Q}^{-1}\auO_E$, fractional ideals being invertible. See also \cite[Chapter~I, Proposition~(2.11), and Chapter~III, Propositions~(2.2) and~(2.4)]{Neukirch}.) As
$\gamma$ is a unit, $\varpi$ generates the different, and $\aubar\varpi=-\varpi$.

For $\tau\in\auPl$, $|\tau(\gamma)|=1$ and, by the chain rule, $|\tau\Phi_{q_1}'(\zeta)|=
t\,|\Phi_{q_0}'(\xi)|$ with $\xi=\tau(\zeta)^t$ a primitive $q_0$-th root of unity. Now
$|\Phi'_{q_0}(\xi)|=\prod_{\xi'}|\xi-\xi'|$ over the other primitive $q_0$-th roots, fewer than $q_0$
factors each in $[2\sin(\pi/q_0),2]\subseteq[4/q_0,2]$. So $|\tau(\varpi)|=A_\tau\,t\sqrt d$ with
$(4/q_0)^{q_0}\le A_\tau\le2^{q_0}$, and any two conjugates differ by a factor at most $(q_0/2)^{q_0}
\le x$. Since $\prod_{\tau\in\auPl}|\tau(\varpi)|^2=|\operatorname{N}_{E/\Q}\varpi|=|\operatorname{disc}
E|=\aurd^{2f}$, the geometric mean of the $|\tau(\varpi)|$ is $\aurd\in[x^{10},x^{12}]$
{\cite[(38)]{P46}}, so $x^9\le|\tau(\varpi)|\le x^{13}$, whence $2\le|\tau(\varpi)|\le b$. With $b^{20}
\le M=p_0p_1\le b^{22}$: $1\le M\cdot2/2\pi\le M|\tau(\varpi)|/2\pi\le b^{22}\cdot b/2\pi\le b^{23}$.
\end{proof}

\subsection{The two definitions of \texorpdfstring{$\auAz$}{A0}}\label{pds:ss-A0}

Let $e_\infty=\bigotimes_\tau e_\tau$, $e_\tau=\pbdeg_\tau\aubar{c_\tau}$, an integrable function on
$G_\infty$ (Theorem~\ref{pds:tail}), and $X=G_0(E^+)\backslash G_0(\auA_{E^+})/J_M$, of finite volume
\cite[Theorem~5.6(i)]{BorelFin}, see also \cite[Theorem~9.4]{BorelHC1962} (the centre of $G_0$ is anisotropic). Let $\auAz^{\mathrm b}$ be the space of bounded
continuous $y$ on $X$ with $\int_{G_\infty}e_\infty(a)y(ga)\,da=y(g)$ for all $g$ (this is $\auAz$, by
the definition of {\cite[\S14.3]{P46}}), and $\auAz^2$ the space of $y\in L^2(X)$ with $\mathcal Ey=y$,
where $\mathcal Ey(g)=\int e_\infty(a)y(ga)\,da$ (the $L^2$ variant of the definition; only the
condition $\mathcal Ey=y$ is used below).

\begin{proposition}\label{pds:A0}
$\auAz^{\mathrm b}\subseteq\auAz^2$. Equality holds if every $y\in\auAz^2$ is (a.e.\ equal to) a bounded
continuous function.
\end{proposition}
\begin{proof}
A bounded $y$ is in $L^2(X)$ since $X$ has finite volume, and $\mathcal Ey=y$ pointwise. Conversely, if
$y\in\auAz^2$ is bounded and continuous, $\mathcal Ey$ is defined everywhere (as $e_\infty\in L^1$) and
continuous (by dominated convergence and the continuity of translation in $L^1(G_\infty)$), and
equals $y$ a.e., hence everywhere.
\end{proof}

\begin{remark}\label{pds:rem-A0}
The converse inclusion is not used in this paper, since $\auAz=\auAz^{\mathrm b}$ by definition.
For orientation: it amounts to the boundedness of every $y\in\auAz^2$. Elements of $\auAz^2$ are $K_\infty$-finite and
$Z(\pdsg)$-finite, hence smooth automorphic forms (Harish-Chandra). Boundedness would follow from
cuspidality, since cusp forms are rapidly decreasing; cuspidality of square-integrable forms whose
archimedean component is a discrete series with parameter far from the walls is a theorem of
Wallach~\cite{WallachCT}. Nothing in \S\S\ref{pb:sec}--\ref{pm:sec} uses it.
\end{remark}

\subsection{The inputs of Theorem~\ref{pbe:ds}}\label{pds:ss-summary}
Of Theorem~\ref{pbe:ds}: \textup{(DS1)(i),(ii)} follow from Schur orthogonality; \textup{(DS1)(iii)}
from Corollary~\ref{pbr:degree} \textup{(}Proposition~\ref{pds:haar}, {\cite[Proposition~2.19]{P49}}\textup{)}; \textup{(DS2)} is proved in full from the explicit formula
\eqref{pds:eq-FJ}; \textup{(DS3)} is proved in full from \textup{(DS1)(ii)} and the Casimir values
\eqref{pds:eq-casimir}; \textup{(DS4)} is proved in \S\ref{pbp:sec}. Formula \eqref{pds:eq-FJ} is
{\cite[Theorem~3.17]{P49}}. Published results used: Schur orthogonality and the value of the Casimir operator on a discrete series.



% ===== B-pillars.tex =====
% proofs/B-pillars.tex
% Theorem pbe:pillars of proofs/B-bessel-expansion.tex: cell ranks (R), size (S), gain (G),
% count (C). Labels of this file start with ppi:.
\providecommand{\auO}{\mathcal O}
\providecommand{\auNm}{\operatorname{Nm}}
\providecommand{\auA}{\mathbb A}
\providecommand{\aubar}[1]{\overline{#1}}
\providecommand{\auf}[1]{\texttt{#1}}
\providecommand{\panot}[1]{\par\smallskip{\raggedright\noindent\textbf{[Not proved here:} #1\textbf{]}\par}\smallskip}
\providecommand{\pbR}{\overline R_0}
\providecommand{\pbDid}{\mathcal D_{\mathrm{id}}}
\providecommand{\pbDmain}{\mathcal D_{\mathrm{main}}}
\providecommand{\pbg}{\mathsf g}
\providecommand{\pbeN}{\mathcal N}
\providecommand{\pbet}{\beta_*}
\providecommand{\pbenp}{\mathsf n_{\mathrm p}}
\providecommand{\pbesz}{\operatorname{size}}
\providecommand{\pbeGs}{\Gamma^*}
\providecommand{\pbeAp}{A^{+}}
\providecommand{\pbeCL}{\mathrm{CL}}
% macros of this file (prefix \ppi)
\providecommand{\ppiF}{\mathsf F}
\providecommand{\ppiLc}{\Lambda^{X}_{c}}
\providecommand{\ppidisc}{\operatorname{disc}}

\section{Pairs of terms: cell ranks, size, gain and count}\label{ppi:sec}

This section concerns Theorem~\ref{pbe:pillars}. We use the notation of \S\S\ref{pb:sec}
and~\ref{pbe:sec}: the frame $f_0,\dots,f_\varrho$ with $f_a=e_{a+1}$ ($a<\varrho$) and $f_\varrho=z$,
$A=\langle f_0,\dots,f_{\varrho-1}\rangle$, $\pbeAp=A+Ez$, $Q=A^\perp/A=X_0\perp Ez$, the matrix $Z(
\gamma)$ with leading minors $\delta_l(\gamma)$, the discriminant $\pbg(\gamma)$, the link
$\Lambda_\gamma$, the isometry $\Upsilon_\gamma$ and the space $X_0^\gamma$ of
Lemma~\ref{pbe:link}, the set $\pbeGs$, and $\sigma$, $\psi_f$ as in Lemma~\ref{pbe:scalar}, so that
$\psi_f(u)=\psi_{E,f}(-\sigma(u)/(M\varpi))$ for $u\in U(E^+)$. For a coefficient vector $x=(x_0,
\dots,x_\varrho)\in E^n$ we write $fx=\sum_ax_af_a$ and $\gamma fx=\gamma(fx)$.

\emph{What is proved.}
\begin{itemize}
\item[(1)] The exact statement of Theorem~\ref{pbe:pillars} is fixed (Theorem~\ref{ppi:pillars}):
the set $\pbeN(c,c')$, the gauge, the discriminant and the size conditions under which it is
proved. Two size conditions on the data are needed that are not in \S\ref{pbe:sec}
(Remark~\ref{ppi:rem-sizes}).
\item[(2)] \textup{(R)}, and the existence of cell ranks with non-vanishing leading minor, are
\textbf{proved in full} (Proposition~\ref{ppi:cell}).
\item[(3)] The elementary parts of \textup{(G)} and \textup{(C)} are proved: the bound $1\le\mathsf
G\le b^{n^3f}$ for near terms (Lemma~\ref{ppi:gauge}), $\mathsf d\ge1$, the passage from the bound on
each common part to \textup{(G)} (Proposition~\ref{ppi:G}), and the reduction of \textup{(C)} to a
count of terms over one lattice and a sum over lattices (Proposition~\ref{ppi:C}).
The count over a lattice (C1) and the sum over lattices (C2) are proved in \S\ref{pbc:sec}. The bound for the common parts (G$'$) is proved in \S\ref{pbg:sec} (Theorem~\ref{pbg:main}). \textup{(S)} is proved in \S\ref{pbs:sec} given (O1) of \S\ref{pbo:sec}.
\item[(4)] Lemmas~\ref{pbe:rank}, \ref{pbe:pairs}, Proposition~\ref{pbe:two} and
Theorem~\ref{pbe:XL1} still follow from the corrected statement (\S\ref{ppi:ss-check}).
\end{itemize}

\subsection{The exact statement}\label{ppi:ss-statement}

\subsubsection*{Common vectors of two flags}
Let $\gamma\in\Gamma$, $0\le a\le\varrho$ and $\mu,\mu'\in E$. We say that \emph{$(\mu,\mu')$ occurs at
$a$} (for $\gamma$) if there are $x,y\in E^n$ with $fx=\gamma fy$, $x_b=y_b=0$ for $b>a$, and
$x_a=\mu$, $y_a=\mu'$. A common vector $fx=\gamma fy\ne0$ has a \emph{top index}: the largest $l$
with $(x_l,y_l)\ne(0,0)$; then $(x_l,y_l)$ occurs at $l$.

\begin{definition}[Cell rank]\label{ppi:def-cell}
Let $\gamma\in\Gamma$. If $\pbg(\gamma)\ne0$, the \emph{cell rank} of $\gamma$ is $j(\gamma)=\varrho+
1$. If $\pbg(\gamma)=0$, $j(\gamma)$ is the least $j\in\{0,\dots,\varrho+1\}$ such that, for one
$\zeta\in E$ with $\zeta\aubar\zeta=1$, $(\zeta,1)$ occurs at every $a$ with $j\le a\le\varrho$ (for
$j=\varrho+1$ the condition is empty, so $j(\gamma)$ exists).
\end{definition}

The cell rank is invariant under $\gamma\mapsto u\gamma u'$ with $u,u'$ in the group $U^0(E^+)$ of
the elements of $U(E^+)$ that fix every $f_a$ (if $fx=\gamma fy$ then $fx=ufx=u\gamma u'(u'^{-1}fy)=
u\gamma u'fy$). One may also require of a cell rank $\delta_{j}(\gamma)\ne0$
and $j\le\varrho\iff\pbg(\gamma)=0$; by Proposition~\ref{ppi:cell} below this holds for every
$\gamma\in\pbeGs$, so the two definitions agree on $\pbeGs$.

\subsubsection*{Near terms, labels, gauges, lattices}
Reduced and boxed points, and near elements, are defined in Definition~\ref{pbc:def-reduced}: $y$ is
reduced if the Gram--Schmidt coefficients of the vectors $y_\tau^{-1}\tilde f_a$, read in the frame
$f$, lie in a fixed Borel set of representatives of $\C^{\auPl}/M\auO_E$ with coordinates at most
$b^{26}$; boxed if moreover $\pbel_\tau(y_\tau)\le b^{6n+2}$ for all $\tau$; and $\gamma$ is near if
some boxed $y$ has $\gamma^{-1}y$ boxed. At a boxed point $\|y_\tau^{-1}\tilde f_a\|\le b^{37n}$
(Lemma~\ref{pbc:box}). For classes $c,c'$ let
$U_c^0=U_c\cap U^0$ and fix in each double coset $U_c^0\gamma U_{c'}^0$ one element, the \emph{chosen}
one.

For $\gamma\in\Gamma_{c,s,c'}$ write $k_\gamma=h_c^{-1}\gamma h_{c'}\in K$. Let $\Lambda=L$ be the
lattice of $Y_0$, $\ppiLc=h_cL\cap X_0$, and let $L_\gamma\subseteq\ppiLc$ be the set of the $v\in
\ppiLc$ such that some vector of $v+M(h_cL\cap A)$ is orthogonal to $\gamma f_0,\dots,\gamma
f_\varrho$. For $\gamma,\gamma'\in\Gamma_{c,s,c'}$ the \emph{common lattice} is $L_{\gamma\gamma'}=
L_\gamma\cap L_{\gamma'}\cap\pbeCL(\gamma,\gamma')$, and $\mathsf d(\gamma,\gamma')=\ppidisc
L_{\gamma\gamma'}$, where for a lattice $L'$ of rank $\rho$ in a subspace $V\subseteq X_0$ we put
$\ppidisc L'=[\ppiLc\cap V:L']\cdot\ppidisc(\ppiLc\cap V)$, and $\ppidisc(\ppiLc\cap V)$ is the square
root of the absolute norm of the ideal of $\auO_E$ generated by the Gram determinants of the $\rho$-tuples
of $\ppiLc\cap V$ (a generalised index).

Let $\gamma\in\Gamma_{c,s,c'}$ have cell rank $j$ and put $j'=\min(\varrho,j)$. The \emph{gauge} is
\begin{equation}\label{ppi:eq-gauge}
\mathsf G(\gamma)=M^{2j'f}\,\frac{\max\bigl(\auNm\pbg(\gamma),\ N^{f}\,\auNm\delta_j(\gamma)\,N^{-2(j-
\varrho)_+f}\bigr)}{i_\gamma\,i'_\gamma},
\end{equation}
where $\auNm x=\prod_{\tau\in\auPl}|\tau(x)|$ for $x\in E^+$ and $\auNm x=\prod_\tau|\tau(x)|^2$ for
$x\in E$, and $i_\gamma=[R_A:R_{<j'}]$, where $R_A$ (resp.\ $R_{<j'}$) is the $\auO_E$-span in $E^n$ of the
rows $a<\varrho$ (resp.\ $a<j'$) of $Z(\gamma)$, and $i'_\gamma$ is the same index for
$Z(\gamma^{-1})=Z(\gamma)^*$; both are finite and $\ge1$, and equal to $1$ if $j=\varrho+1$
(Remark~\ref{pbg:rem-gauge}). Equivalently, $\mathsf G(\gamma)=M^{2j'f}\max(\auNm\delta_*(\gamma)/
(i_\gamma i'_\gamma N^f),\auNm\pbg(\gamma))$ with the subframe minor $\delta_*$
(Lemma~\ref{pbg:subframe}).

Finally, $s\in\auG$ \emph{agrees} if $\psi_f(u)=\psi_f(u')$ for all integral $u,u'\in U$ whose
reductions satisfy $\bar us=s\bar u'$ (this notion, for $s$, is distinct from the agreement of an element $\gamma$ in \S\ref{pbe:ss-char}, which is the notion used in Lemma~\ref{ppi:climb} and Proposition~\ref{ppi:cell}). The \emph{label} of $\gamma\in\Gamma_{c,s,c'}$ is the component of $s$
in $\pbR$ for a fixed decomposition $\auG=\bigsqcup_\rho\rho\pbR$.

\begin{definition}\label{ppi:def-N}
$\pbeN(c,c')$ is the set of the $\gamma\in\pbeGs$ that are near and chosen and lie in
$\Gamma_{c,s,c'}$ for an agreeing $s$ (which $\gamma$ determines).
\end{definition}

Recall $\pbet=b^{kf/4}$, $\pbenp=\#\pbR\cdot\#\mathcal C/N^{kf}$, $\mathcal D(\gamma)$ and
$\pbesz(\gamma)=\pbet^2(\mathcal D(\gamma)/(\pbenp\mathsf G(\gamma)^k))^{1/2}$ from
\S\ref{pbe:ss-pillars}. Note that, for $\gamma\in\pbeGs$, $\mathcal D(\gamma)=\pbDmain$ if and only if $j(\gamma)=\varrho+1$ (Proposition~\ref{ppi:cell}).

\begin{theorem}[Size, gain and count; exact form]\label{ppi:pillars}
Assume the size conditions
\begin{equation}\label{ppi:eq-sizes}
\varrho\ge10^5,\qquad k\ge1000\,(2n)^9 .
\end{equation}
Let $W^\flat$ be as in Theorem~\ref{pbe:bessel}. With $\pbeN(c,c')$, $\mathsf G$ and $\mathsf d$ as
above, $\pbeN(c,c')$ contains at most one element of each double coset $U_c^0\gamma U_{c'}^0$,
$\mathsf G$ takes values in $(0,b^{n^3f}]$ on $\pbeN(c,c')$, $\mathsf d(\gamma,\gamma')=\mathsf d(\gamma',
\gamma)\ge1$, and:
\begin{enumerate}
\item[(R)] If $\gamma\in\pbeGs$ and $\pbg(\gamma)=0$, then $\pbeAp\cap\gamma\pbeAp\not\subseteq A$.
\item[(S)] For every stable $T\subseteq\pbeGs$ and every block $(c,s,c')$, $\|A^\flat_T(c,s,c')\|\le
\nu\sum_{\gamma\in\pbeN(c,c')\cap T\cap\Gamma_{c,s,c'}}\pbesz(\gamma)$. In particular $A^\flat_T(c,s,
c')=0$ if $s$ does not agree.
\item[(G)] For $\gamma,\gamma'\in\pbeN(c,c')\cap\Gamma_{c,s,c'}$ and $i=i(\gamma,\gamma')$,
$(\mathsf d(\gamma,\gamma')\,\mathsf x)^2\le(2/M)^{2if}\,\mathsf G(\gamma)\mathsf G(\gamma')$, with
$\mathsf x=\mathsf{mn}(s)$ if $i=0$ and $\pbg(\gamma)\ne0$, and $\mathsf x=1$ otherwise.
\item[(C)] For every integer $0\le l\le n^3$, $\sum_{c,c'}\sum_{(\gamma,\gamma')}\mathsf d(\gamma,
\gamma')^{-k}\le\pbet^{\,l+2}\pbenp$, the inner sum over the pairs of elements of $\pbeN(c,c')$ over a
common block with $k-\dim\pbeCL(\gamma,\gamma')=l$ and $\mathsf d(\gamma,\gamma')\le b^{n^3f}$.
\end{enumerate}
\end{theorem}

\begin{remark}[Changes against Theorem~\ref{pbe:pillars}]\label{ppi:rem-sizes}
(i) \eqref{ppi:eq-sizes} ($10^5\le\varrho$, $1000(2\varrho+2)^9\le k$) holds for the data by
{\cite[Lemma~13.4(S7)]{P46}}. It is used at least twice:
$167n+270n^2\le n^3$ in Lemma~\ref{ppi:gauge} (which needs $n\ge280$), and
$4000\,n^6(l+1)\le k/4$ for $l\le n^3$, i.e.\ $b^{4000n^6(l+1)f}\le\pbet$, in
Proposition~\ref{ppi:C} (which needs $k\ge32000\,n^9$). (ii) (R) is stated for all of $\pbeGs$, not only for $\pbeN(c,c')$. (iii) $\pbeN(c,c')$
is restricted to agreeing blocks, as (G) and (C) are proved only there; (S) is not
weakened by this, because the bound is proved with the sum restricted to agreeing $s$ (\S\ref{pbs:sec}). (iv)
The gauge, the common lattice and nearness are now defined. In (S), the constant in front of the sum
is $\nu=\pbvol(U_c\backslash U_\infty)$, the integral over $U_\infty$ of a cutoff for $U_c$, for the
Haar measure $du$ of convention (Cv3) of \S\ref{pb:ss-setting}; (S) as a statement does not depend on
the choice of $du$ (Remark~\ref{pbs:rem-nu}).
\end{remark}

\subsection{Cell ranks}\label{ppi:ss-cell}

In this subsection, as in \S\ref{pbg:sec}, hermitian forms are written linear in the second variable
and conjugate-linear in the first, the opposite of convention (Cv1) of \S\ref{pb:ss-setting}: we write
$h$ for the form $(x,y)\mapsto h(y,x)$. Isometries, orthogonality, isotropy, the matrix $Z(\gamma)$ up to
transposition and conjugation, the vanishing of the minors $\delta_l(\gamma)$ and of $\pbg(\gamma)$, and
$U(E^+)$ do not change under this replacement (\S\ref{pbv:ss-weak}(iv)), so the statements below hold
in either convention.

\begin{lemma}[Eichler elements]\label{ppi:eichler}
Let $0\le j<\varrho$, let $p\in\langle f_0,\dots,f_j\rangle$ and let $q\in Y_0$ be orthogonal to
$f_0,\dots,f_j$. Choose $\omega\in\auO_E$ with $\omega+\aubar\omega=1$ and put $\varsigma=\omega h(q,q)$.
Then
\[
\mathsf E_{p,q}(t)=t+h(q,t)\,p-h(p,t)\,(q+\varsigma p)
\]
lies in $U(E^+)$, and if $p=fx$ then $\sigma(\mathsf E_{p,q})=x_j\,h(q,f_{j+1})$. For $g\in G(E^+)$,
$g\mathsf E_{p,q}g^{-1}=\mathsf E_{gp,gq}$.
\end{lemma}
\begin{proof}
$p\in A$ is isotropic and $h(q,p)=0$. With these two facts and $\varsigma+\aubar\varsigma=h(q,q)$, a
direct expansion of $h(\mathsf E t,\mathsf Et')$ gives $h(t,t')$ (the Eichler transformation), so
$\mathsf E=\mathsf E_{p,q}$ is an isometry, defined over $E$. For $t\in A_i=\langle e_1,\dots,e_i
\rangle$ we have $h(p,t)=0$, and $h(q,t)=0$ if $i\le j+1$; so $\mathsf E$ fixes $A_{j+1}$ pointwise
and maps $t\in A_i$ to $t+h(q,t)p\in t+A_{j+1}$. For $t\in A^\perp$, $\mathsf Et=t+h(q,t)p\in t+A$.
So $\mathsf E$ preserves the flag of $A$, acts trivially on its graded pieces and on $A^\perp/A$:
$\mathsf E\in U(E^+)$. For $b<\varrho$, $\mathsf Ef_{b+1}-f_{b+1}=h(q,f_{b+1})p$, whose coefficient
at $f_b$ is $x_b\,h(q,f_{b+1})$; as $x_b=0$ for $b>j$ and $h(q,f_{b+1})=0$ for $b+1\le j$, only
$b=j$ remains in $\sigma(\mathsf E)$. The last claim holds because $g$ is an isometry.
\end{proof}

\begin{lemma}[Climbing]\label{ppi:climb}
Let $\gamma$ agree (in the sense of \S\ref{pbe:ss-char}) and let $(\mu,\mu')$ occur at $a\le\varrho$. If $a<\varrho$, then $(\mu,\mu')$
occurs at $a+1$. Hence it occurs at every $a'$ with $a\le a'\le\varrho$, and in all cases
$\mu\aubar\mu=\mu'\aubar{\mu'}$.
\end{lemma}
\begin{proof}
Let $fx=\gamma fy$ with $x,y$ supported on $\{0,\dots,a\}$, $x_a=\mu$, $y_a=\mu'$, and let $W_a=\langle
f_b,\gamma f_b:b\le a\rangle$. If $\mu f_{a+1}-\mu'\gamma f_{a+1}=\sum_{b\le a}(\alpha_bf_b-\beta_b
\gamma f_b)$, then $x'=\mu\,\epsilon_{a+1}-\sum_{b\le a}\alpha_b\epsilon_b$ and $y'=\mu'\epsilon_{a+1}-
\sum_{b\le a}\beta_b\epsilon_b$ (with $\epsilon_b$ the unit vectors) show that $(\mu,\mu')$ occurs at
$a+1$. So it suffices that every $E$-linear form $\lambda$ on $Y_0$ that vanishes on $W_a$ satisfies
$\Delta:=\mu\lambda(f_{a+1})-\mu'\lambda(\gamma f_{a+1})=0$. Write $\lambda=h(q,\cdot)$. Then $q\perp
f_b$ and $\gamma^{-1}q\perp f_b$ for $b\le a$. For $r\in E$ put $u=\mathsf E_{rfx,q}$ and $w=\mathsf
E_{rfy,\gamma^{-1}q}$ (Lemma~\ref{ppi:eichler}, with $h(\gamma^{-1}q,\gamma^{-1}q)=h(q,q)$, so the same
$\varsigma$). Both lie in $U(E^+)$, and $\gamma w\gamma^{-1}=\mathsf E_{r\gamma fy,q}=u$. So $u\in
U(E^+)\cap\gamma U(E^+)\gamma^{-1}$ and, $\gamma$ agreeing, $\psi_f(u)=\psi_f(w)$. By
Lemma~\ref{ppi:eichler}, $\sigma(u)=r\mu\lambda(f_{a+1})$ and $\sigma(w)=r\mu'h(\gamma^{-1}q,f_{a+1})=
r\mu'\lambda(\gamma f_{a+1})$. With $r=-M\varpi t$, $t\in E$, this gives $\psi_{E,f}(t\Delta)=1$ for
all $t\in E$. If $\Delta\ne0$, $\psi_{E,f}$ would be trivial on $E$, which is dense in $\auA_{E,f}$,
whereas the finite part of the standard character is not trivial. So $\Delta=0$.

By induction (trivially if $a=\varrho$) $(\mu,\mu')$ occurs at $\varrho$: $fx=\gamma fy$ with $x_\varrho=\mu$, $y_\varrho=\mu'$.
As $A$ is isotropic and orthogonal to $z$, $-N\mu\aubar\mu=h(fx,fx)=h(\gamma fy,\gamma fy)=h(fy,fy)=
-N\mu'\aubar{\mu'}$.
\end{proof}

\begin{proposition}[Cell ranks]\label{ppi:cell}
Let $\gamma\in\Gamma$ and $j=j(\gamma)$.
\begin{enumerate}
\item[(a)] If $\pbg(\gamma)\ne0$, then $\delta_n(\gamma)\ne0$.
\item[(b)] If $\gamma$ agrees and $\pbg(\gamma)=0$, then $j\le\varrho$, $\delta_j(\gamma)\ne0$ (with
$\delta_0=1$), and $\pbeAp\cap\gamma\pbeAp\not\subseteq A$.
\end{enumerate}
In particular \textup{(R)} holds, and every $\gamma\in\pbeGs$ has $\delta_{j(\gamma)}(\gamma)\ne0$ and
$j(\gamma)\le\varrho\iff\pbg(\gamma)=0$.
\end{proposition}
\begin{proof}
We use twice the following remark: \emph{if $y\in E^n$ and $\gamma fy$ is orthogonal to $f_0,\dots,
f_\varrho$, then $y_\varrho=0$ and $\gamma fy\in A$.} Indeed $(\pbeAp)^\perp=A\oplus X_0$, so $\gamma fy=
\alpha+\xi$ with $\alpha\in A$, $\xi\in X_0$, and $h(\xi,\xi)=h(\gamma fy,\gamma fy)=h(fy,fy)=-N|y_\varrho
|^2$; as $X_0$ is definite (positive at every real place), both sides vanish.

(a) If $\delta_n(\gamma)=0$ there is $y\ne0$ with $\gamma fy\perp f_a$ for all $a$; by the remark
$0\ne\gamma fy\in A\cap\gamma\pbeAp$, so $\pbg(\gamma)=0$ by Lemma~\ref{pbe:link}(a).

(b) By Lemma~\ref{pbe:link}(a) there is a common vector $fx=\gamma fy\ne0$. Let $l$ be its top index.
By Lemma~\ref{ppi:climb}, $|x_l|=|y_l|$, so both are non-zero; with $\zeta=x_l/y_l$, $(\zeta,1)$
occurs at $l$, hence at every $a\ge l$, and $\zeta\aubar\zeta=1$. So $j\le l\le\varrho$. At $a=\varrho\ge
j$ the definition of $j$ gives $fx=\gamma fy$ with $x_\varrho=\zeta\ne0$: this vector lies in $\pbeAp\cap
\gamma\pbeAp$ and not in $A$.

Let $j\ge1$ and suppose $\delta_j(\gamma)=0$. Then there is $y\ne0$, supported on $\{0,\dots,j-1\}$,
with $h(f_a,\gamma fy)=0$ for $a<j$. We show $h(f_a,\gamma fy)=0$ for $j\le a\le\varrho$ by induction
on $a$. By definition of $j$ there is $fx^{(a)}=\gamma fy^{(a)}$ with top index $a$, $x^{(a)}_a=\zeta$,
$y^{(a)}_a=1$. Then $h(fx^{(a)},\gamma fy)=h(fy^{(a)},fy)=-N\aubar{y^{(a)}_\varrho}y_\varrho=0$, as
$y_\varrho=0$; and $h(fx^{(a)},\gamma fy)=\sum_{b\le a}\aubar{x^{(a)}_b}h(f_b,\gamma fy)=\aubar\zeta\,
h(f_a,\gamma fy)$ by the choice of $y$ and the induction hypothesis. So $h(f_a,\gamma fy)=0$. By the
remark, $\gamma fy=fx\in A$ with $x\ne0$ (as $y\ne0$). Let $l$ be the top index of this common vector.
As above $x_l,y_l\ne0$ and $(x_l/y_l,1)$ occurs at every $a\ge l$ with $|x_l/y_l|=1$, so $l\ge j$ by
minimality; but $y_l\ne0$ forces $l<j$. This contradiction shows $\delta_j(\gamma)\ne0$.

The elements of $\pbeGs$ with $\pbg=0$ agree by definition, which gives the last sentence.
\end{proof}

\subsection{The gauge of a near term, and the gain}\label{ppi:ss-gain}

\begin{lemma}\label{ppi:gauge}
Let $\gamma\in\Gamma_{c,s,c'}$ be near. Then $|\tau(h(f_a,\gamma f_b))|\le b^{134n}$ for all $\tau,a,b$,
$\auNm\pbg(\gamma)\le b^{270n^2f}$, $\auNm\delta_l(\gamma)\le b^{270n^2f}$ for all $l$, and, if
$n\ge280$ and $\gamma\in\pbeGs$, $0<\mathsf G(\gamma)\le b^{n^3f}$.
\end{lemma}
\begin{proof}
Let $y$ and $\gamma^{-1}y$ be boxed; by Lemma~\ref{pbc:box}, $\|y_\tau^{-1}\tilde f_a\|,\|(\gamma^{-1}y)_\tau^{-1}\tilde f_a\|\le b^{37n}$. With the scales $s_a$ of \eqref{pbe:eq-scale}, $|s_a|\ge
N^{-1/2}$ (as $|\kappa_\tau|\ge1$), and by \eqref{pbe:eq-norm}
\[
|\tau(h(f_a,\gamma f_b))|=\frac{|h(y_\tau^{-1}\tilde f_a,(\gamma^{-1}y)_\tau^{-1}\tilde f_b)|}{|s_as_b|}
\le N\,b^{37n}\,b^{37n}\le b^{134n},
\]
using $N\le b^{60n}$. A minor of order $l\le n$ is a sum of $l!\le n^l\le b^l$ products of $l$
entries (as $n\le b$, {\cite[Lemma~13.4(S7)(e)]{P46}}), so $|\tau(\delta_l)|^2\le b^{2l+268nl}\le b^{2n+268n^2}\le b^{270n^2}$. For the discriminant,
\[
|\tau(\pbg)|\le|\tau(\delta_n)|^2+N^2|\tau(\delta_{n-1})|^2\le b^{2n+268n^2}+b^{120n}\,b^{2(n-1)+268n(n-1)}
\le2\,b^{2n+268n^2}\le b^{270n^2},
\]
since $120n+2(n-1)+268n(n-1)=268n^2-146n-2\le2n+268n^2$, and $2\le b\le b^{2n^2-2n}$ as $2\le n\le b$
({\cite[Lemma~13.4(S7)(e)]{P46}} and~\eqref{ppi:eq-sizes}). In particular (a weaker form used below and in \S\ref{pbs:sec}) $|\tau(\pbg)|\le b^{121n+270n^2}$. Taking products over the $f$ places,
and as $i_\gamma,i'_\gamma\ge1$, $N\le b^{60n}$ and $M\le b^{23}$,
\[
\mathsf G(\gamma)\le b^{46nf}\,b^{(121n+270n^2)f}\le b^{n^3f}\qquad(n\ge280).
\]
Positivity: by Proposition~\ref{ppi:cell}, $\delta_{j(\gamma)}(\gamma)\ne0$ for $\gamma\in\pbeGs$.
\end{proof}

\begin{proposition}[Gain, given the bound for the common parts]\label{ppi:G}
Assume the following statement \textup{(G$'$)}. Then \textup{(G)} holds, and $\mathsf d\ge1$.
\begin{enumerate}
\item[(G$'$)] Let $s$ agree, $\gamma,\gamma'\in\Gamma_{c,s,c'}\cap\pbeGs$, $C_\gamma=L_\gamma\cap
\pbeCL(\gamma,\gamma')$ and $t(\gamma)=[L\cap(A+k_\gamma A):L\cap A+k_\gamma(L\cap A)]=[h_cL\cap(A+\gamma A):h_cL\cap A+\gamma(h_cL\cap A)]$ (the
torsion index). Then $C_\gamma$ is a lattice of rank $\dim\pbeCL(\gamma,\gamma')$, $[C_\gamma:L_{\gamma
\gamma'}]$ divides $t(\gamma')$, and
\[
t(\gamma)\,\ppidisc C_\gamma\cdot t(\gamma')\,\ppidisc C_{\gamma'}\cdot\mathsf x^2\le(2/M)^{2if}\,
\mathsf G(\gamma)\mathsf G(\gamma').
\]
\end{enumerate}
\end{proposition}
\begin{proof}
$L_{\gamma\gamma'}\subseteq C_\gamma$ with index dividing $t(\gamma')$, so $L_{\gamma\gamma'}$ has the
same rank and $\mathsf d(\gamma,\gamma')=[C_\gamma:L_{\gamma\gamma'}]\ppidisc C_\gamma\le t(\gamma')
\ppidisc C_\gamma$; symmetrically ($L_{\gamma\gamma'}$ and $\pbeCL$ are symmetric in $\gamma,\gamma'$)
$\mathsf d\le t(\gamma)\ppidisc C_{\gamma'}$. Multiply and use (G$'$). The ideal generated by the
Gram determinants of $\ppiLc\cap V$ is integral, because $\ppiLc$ is an integral hermitian lattice;
so $\ppidisc(\ppiLc\cap V)\ge1$ and $\mathsf d\ge1$.
\end{proof}

Statement (G$'$) is Theorem~\ref{pbg:main}, proved in \S\ref{pbg:sec}.

\subsection{The count}\label{ppi:ss-count}

\begin{proposition}[Count, given two lattice estimates]\label{ppi:C}
Assume \eqref{ppi:eq-sizes}, \textup{(G$'$)} and the following two statements, for $0\le l\le n^3$,
$r=k-l$:
\begin{enumerate}
\item[(C1)] For a class $c$, a label $t\in\pbR$ and a lattice $L'\subseteq\ppiLc$ of rank $r$ with
$\ppidisc L'\le b^{n^3f}$, there are at most $b^{2000n^6(l+1)f}$ pairs $(c',\gamma)$ with $\gamma$ near,
chosen, of label $t$, in $\Gamma_{c,s,c'}\cap\pbeGs$ for some $s$, and with $L'\subseteq L_\gamma$.
\item[(C2)] $\displaystyle\sum_{c\in\mathcal C}\ \sum_{L'}\ (\ppidisc L')^{-k}\le\frac{\#\mathcal
C}{N^{kf}}\,\pbet^{\,l+1}$, the inner sum over the lattices $L'\subseteq\ppiLc$ of rank $r$ with
$\ppidisc L'\le b^{n^3f}$.
\end{enumerate}
Then \textup{(C)} holds.
\end{proposition}
\begin{proof}
Fix $c$. Send a pair $(\gamma,\gamma')$ counted in (C) for $c'$ to $(t,L_{\gamma\gamma'})$, where $t$
is the common label (both lie over $(c,s,c')$). By Proposition~\ref{ppi:G}, $L_{\gamma\gamma'}$ has
rank $\dim\pbeCL(\gamma,\gamma')=r$; it lies in $\ppiLc$, in $L_\gamma$ and in $L_{\gamma'}$, and
$\ppidisc L_{\gamma\gamma'}\le b^{n^3f}$. So both $(c',\gamma)$ and $(c',\gamma')$ lie in the set of
(C1) for $(t,L_{\gamma\gamma'})$, and the map $(c',\gamma,\gamma')\mapsto((c',\gamma),(c',\gamma'))$ is
injective. Hence each fibre has at most $b^{4000n^6(l+1)f}\le b^{8000n^9f}\le b^{kf/4}=\pbet$ elements, by
$l+1\le2n^3$ and \eqref{ppi:eq-sizes} ($32000\,n^9\le512000\,n^9\le k$). All pairs in a
fibre have the weight $(\ppidisc L_{\gamma\gamma'})^{-k}$, so
\[
\sum_{c'}\sum_{(\gamma,\gamma')}\mathsf d^{-k}\le\#\pbR\cdot\pbet\sum_{L'}(\ppidisc L')^{-k}.
\]
Summing over $c$ and using (C2), the total is at most $\#\pbR\cdot\pbet\cdot\#\mathcal CN^{-kf}
\pbet^{l+1}=\pbet^{l+2}\pbenp$.
\end{proof}

Statements (C1) and (C2) are proved in Theorems~\ref{pbc:C1} and~\ref{pbc:C2}; hence (C) holds, given (G$'$).

\subsection{Size}\label{ppi:ss-size}

\smallskip\noindent\textup{(S)} is Theorem~\ref{pbs:S}, proved in \S\ref{pbs:sec} from the open-cell bound \textup{(O1)} (Theorem~\ref{pbo:O1}).\par\smallskip

\subsection{The consequences still follow}\label{ppi:ss-check}

\begin{proposition}\label{ppi:check}
Lemmas~\ref{pbe:rank} and~\ref{pbe:pairs}, Proposition~\ref{pbe:two} and Theorem~\ref{pbe:XL1} hold
with Theorem~\ref{ppi:pillars} in place of Theorem~\ref{pbe:pillars}, with the same proofs.
\end{proposition}
\begin{proof}
The proofs in \S\ref{pbe:ss-XL1} use only the following, all contained in
Theorem~\ref{ppi:pillars}: (R) for $\gamma\in\pbeN(c,c')\subseteq\pbeGs$; $1\le\mathsf d$ and $\mathsf
G\le b^{n^3f}$; (G) with the same $\mathsf x$; (C) for every $l\le n^3$; and (S) for the stable sets
$T_\Lambda$, $T_0=\{\gamma\in\pbeGs:\pbg(\gamma)=0\}$ and $\pbeGs$, which are subsets of $\pbeGs$.
We check the points where the definitions enter. In Lemma~\ref{pbe:rank}(b), $\Upsilon_\gamma=
\Upsilon_{\gamma'}$ gives $X_0^\gamma=X_0^{\gamma'}$, and by Lemma~\ref{pbe:link}(c) and (R) $\dim
X_0^\gamma$ is $k-1$ or $k$ according as $\pbg\ne0$ or $\pbg=0$: unchanged. In
Lemma~\ref{pbe:rank}(c), $18i\le n^3$ gives $l\le i+1\le n^3$. In Lemma~\ref{pbe:pairs}, $\mathcal
D(\gamma)$ is $\pbDmain$ exactly at cell rank $\varrho+1$, so $\pbesz$ is the size used in Theorem~\ref{ppi:pillars}(S) (defined in \S\ref{pbe:ss-pillars}), and
$n^3+1\le b^4\le\pbet$. In Proposition~\ref{pbe:two}(a) with $i=0$, $\pbg(\gamma)=\pbg(\gamma')=0$ and
$\mathsf x=1$; with $i\ge1$, $\mathsf x\ge1$; in (b), $i=0$ and $\pbg\ne0$, so $\mathsf x=\mathsf{mn}(s)
\ge\mathsf z$. In Theorem~\ref{pbe:XL1}, a pair of elements of $\pbeN(c,c')\cap\Gamma_{c,s,c'}$ lies
over an agreeing block, as required by (G) and (C), and the sum in (S) runs over such elements only.
\end{proof}



% ===== B-period.tex =====
% proofs/B-period.tex
% Proof of (DS4) of Theorem pbe:ds: the lower bound for the local Bessel period of the discrete
% series coefficient at a stationary point. Labels of this file start with pbp:.
\providecommand{\auU}{\mathrm U}
\providecommand{\aubar}[1]{\overline{#1}}
\providecommand{\pbvol}{\operatorname{vol}}
\providecommand{\pbip}[2]{\langle#1,#2\rangle}
\providecommand{\pbW}{\mathbb W}
\providecommand{\pbel}{\ell}
\providecommand{\pbeB}{\mathsf B}
\providecommand{\pdsK}{\mathbf K}
\providecommand{\pdsg}{\mathfrak g}
\providecommand{\pdsm}{m}
\providecommand{\pdsHS}[1]{\|#1\|_{2}}
\providecommand{\pcoD}{\mathsf D}
% macros of this file (prefix \pbp)
\providecommand{\pbpF}[1]{\|#1\|_{\mathrm F}}
\providecommand{\pbptr}{\operatorname{tr}}
\providecommand{\pbpct}{\operatorname{ct}}
\providecommand{\pbpCT}{\operatorname{CT}}
\providecommand{\pbpQ}{\mathcal Q}
\providecommand{\pbpLam}{\widehat\Lambda}
\providecommand{\pbpst}{\tilde\sigma}
\providecommand{\pbpI}{\mathsf I}

\section{The Bessel period at a stationary point}\label{pbp:sec}

This section proves \textup{(DS4)} of Theorem~\ref{pbe:ds}, stated in \S\ref{pds:ss-ds4}.
\emph{Dependencies.} Theorem~\ref{pbp:main} below, which is \textup{(DS4)}, uses:
the explicit formula \eqref{pds:eq-FJ} ({\cite[Theorem~3.17]{P49}}) together with the independence of the
block diagonal factor ({\cite[Proposition~3.21]{P49}}) and the $K$-behaviour ({\cite[Lemma~3.23]{P49}}) of
{\cite[\S3.2]{P49}}; the decay Theorem~\ref{pds:decay} and the lower bound of Lemma~\ref{pds:lip}; Witt's
theorem for hermitian forms (proved in the form used as Lemma~\ref{pbp:tb-witt}; see also \cite[Chapter~7, Theorem~9.1]{Scharlau}); Schwarz's lemma; and elementary linear algebra. The normalisation of
the character (Lemmas~\ref{pbp:psi}, \ref{pbp:sign}), the weight vector (Lemma~\ref{pbp:W}) and the
integrality of $\kappa(\tau)$ (Lemma~\ref{pbp:sizes}) are treated in \S\ref{pbp:ss-data}.

The argument is Laplace's method. A point $x$ is chosen (\S\S\ref{pbp:ss-pos}--\ref{pbp:ss-point})
so that, on the group $x^{-1}R_\tau x$, the phase of the character cancels the phases of the two
coefficients to first order (``stationary position''); near $1$ each coefficient is then its phase
times a number close to $[0,1]$ (\S\S\ref{pbp:ss-minor}--\ref{pbp:ss-gauss}, ``Gaussian shape'');
away from $1$ the integrand is either small or equal to a value taken near $1$
(\S\ref{pbp:ss-loc}); so the integrand lies in a cone about the positive axis up to an
exponentially small error, and a covering argument compares the integral with the volume of
$\pbeB_x$ (\S\ref{pbp:ss-cone}).

\subsection{Statement, conventions and sizes}\label{pbp:ss-stat}

Fix $\tau\in\auPl$; it is dropped from the notation when no confusion arises. A word on letters: $h$ denotes the gap below, and also the hermitian form in
\S\ref{pbp:ss-data} and an element of $H_\tau$ (or of $\auU(k)$) in Lemmas~\ref{pbp:R}, \ref{pbp:W}, \ref{pbp:compact} and in the
Reduction at the end of \S\ref{pbp:ss-data}, where the form itself does not occur; the image of
$s\in R_\tau$ in $H_\tau$ is written $h_s$. Likewise $k$
denotes an integer (a rank or a coordinate index) or an element of $\pdsK$; the meaning is fixed in
each lemma and the two never occur in the same formula. We use the set-up of
\S\ref{pds:ss-setup} (with rank $r$): $G=\auU(J)$ with $J=\operatorname{diag}(\varepsilon_1,\dots,
\varepsilon_r)$, $\pdsK=G\cap\auU(r)$, the parameter $\lambda$, $\rho$, $\mu=\lambda+\rho$, the
highest weight $\Lambda$ of the lowest $\pdsK$-type, the gap $h\ge\Delta=2r^8$, $L=\max_i|\lambda_i|$,
the products of minors $\pdsm_w$ of \eqref{pds:eq-m}, the block diagonal part $a(x)$, and the unitary
matrix $p$ with $p^{-1}\mathsf J_1p=J$, through which $G$ and $G_\tau$ are identified; thus
$c_\tau(g)=c(p^{-1}gp)$ for $g\in G_\tau$, where $c$ is the coefficient on $G$ given by
\eqref{pds:eq-FJ}. We regard $p$ as a map from $\C^r$, with the form $B(v,w)=v^*Jw$, to $Y_\tau$ with
the coordinates of the normalised basis; it is an isometry both for the hermitian forms and for
the Euclidean norms.

\emph{Norms.} For a matrix $X$, $\|X\|$ is the operator norm and $\pbpF X=(\pbptr X^*X)^{1/2}$ the
Frobenius norm. We use $\|X\|\le\pbpF X$, $\pbpF{XY}\le\pbpF X\pbpF Y$, $\pbpF{XY}\le\|X\|\pbpF Y$,
$|\pbptr XY|\le\pbpF X\pbpF Y$, $|X_{ij}|\le\pbpF X$, and that $\pbpF{\cdot}$ is invariant under
multiplication by unitary matrices; for an orthogonal projection $P$, $\pbpF{PX}\le\pbpF X$.

\emph{Coordinates of $G$.} By Lemma~\ref{pds:chamber} and the remark after it, the coordinates
$1,\dots,r$ of $\C^r$ split into a set $X$ of $k-1$ coordinates, on which $\varepsilon=+1$ and
$\Lambda=0$, and $2n$ further coordinates ($n=\varrho+1$, so $2n=2\varrho+2=m$), which we number
$0,1,\dots,2n-1$ in the order of decreasing $\Lambda$; on the coordinate numbered $i$ the sign is
$(-1)^i$, and the values $\mu_i:=\Lambda_i$ satisfy
\begin{equation}\label{pbp:eq-steep}
-3B\le\mu_{2n-1}<\dots<\mu_1<\mu_0\le-B,\qquad B:=\kappa(\tau)+\Delta/2 .
\end{equation}
(The values are integers; $|\Lambda_i-\lambda_i|<r$, the entries $\lambda_i=-\kappa(\tau)-j\Delta$,
$1\le j\le m$, decrease in steps of $\Delta>2r$, $r\le\Delta/2$ and $m\Delta+r\le2\kappa(\tau)+
\Delta/2$; the last two are comparisons of sizes, Lemma~\ref{pbp:sizes}.) To shorten we write
$\mu$ also for the diagonal matrix $\operatorname{diag}(\Lambda)$ when it acts on vectors; the
vector $\lambda+\rho$ of \S\ref{pds:ss-setup} will only appear as $\lambda+\rho$.

\emph{Coordinates of $H$.} $H_\tau=\auU(k)$ acts on the span of $x_1,\dots,x_k$; its coordinates are
$X'=\{1,\dots,k-1\}$ and the coordinate $k$. We fix a bijection between $X'$ and $X$ and use it to
identify them. We write $a_\tau$ for the number of columns of the diagram of $\lambda^X$ at $\tau$:
$a_\tau=a=xb$ if $\tau\in\auPa$ and $a_\tau=0$ otherwise ({\cite[\S14.3]{P46}}), and
$\Lambda'=\operatorname{diag}(0,\dots,0,a_\tau)$.

\emph{Sizes.} Put
\begin{equation}\label{pbp:eq-T}
D=r+k,\qquad\beta=b^3,\qquad T=D^2\log\beta,\qquad h_0=\lceil10^{30}T^3\rceil,\qquad
\rho_0=(200\beta^5)^{-1}.
\end{equation}

\begin{lemma}\label{pbp:sizes}
\textup{(a)} $\kappa(\tau)$ is a positive integer, $1\le B$, $m\Delta+r\le\kappa(\tau)$,
$r\le\Delta/2$, and $3B+a_\tau\le b^2$. \textup{(b)} $L\le2b$, $\max_i|\Lambda_i|\le3b$, $8nr\le b$,
$r\le b$, $a_\tau\le b^2$. \textup{(c)} $\beta\ge10^{12}$, $\log\beta\ge1$, $r^2\le T$,
$18D^2\le h_0\le\Delta\le h$. \textup{(d)} $\rho_0\le(48b^2)^{-1}$ and $a_\tau\rho_0\le\frac12$.
\end{lemma}
\begin{proof}
(a) $\kappa(\tau)=b$ or $xb/x_1$, an integer because $q_1x_1$ divides $b$ ({\cite[Definition~13.3]{P46}},
(D4), (D6)); so $B\ge1$. By {\cite[(36)]{P46}}, $x\le b^{10^{-12}}$, $r\le3x^5$
and $m\le2\sqrt x$, so $m\Delta+r\le4\sqrt x\,(3x^5)^8+3x^5\le x^{42}$, while $\kappa(\tau)\ge b/x_1
\ge b^{1-10^{-6}}$; and $r\le r^8=\Delta/2$. Finally $3B+a_\tau\le3b+3r^8+xb\le b^2$.
(b) $L=\kappa(\tau)+m\Delta\le2\kappa(\tau)\le2b$; $|\Lambda_i|\le L+r\le3b$; $8nr\le8r^2\le
72x^{10}\le b$; $a_\tau\le xb\le b^2$.
(c) $\beta=b^3$ is enormous; $T\ge r^2\log\beta\ge r^2$. Since $r/2\le k<r$, $D\le2r$ and $T\le
12r^2\log b$, so $10^{30}T^3+1\le1.8\cdot10^{33}r^6(\log b)^3\le2r^8$, because $r\ge x^5\ge
(\log b)^{500}$ and $\log b\ge128$ give $r^2\ge(\log b)^{1000}\ge10^{34}(\log b)^3$. Also $18D^2\le
72r^2\le h_0$, and $h\ge\Delta$ by \S\ref{pbe:ss-ds}.
(d) $48b^2\le200b^{15}$; $a_\tau\rho_0\le b^2/(200b^{15})$.
\end{proof}

\begin{theorem}[\textup{(DS4)}]\label{pbp:main}
There are $x_\tau\in G_\tau$ with $\pbel_\tau(x_\tau)\le b^{5n}$ and a unit vector $v_\tau\in
\pbW_\tau$ such that
\[
\pbvol(\pbeB_{x_\tau})\le b^{60(r+k)^2}\,\Bigl|\int_{R_\tau}c_\tau(x_\tau^{-1}s\,x_\tau)\,
\pbip{v_\tau}{\pi_R(s)^{-1}v_\tau}\,ds\Bigr| .
\]
\end{theorem}

This is \textup{(DS4)} as stated in Theorem~\ref{pbe:ds}. The proof is completed in
\S\ref{pbp:ss-cone}.

\subsection{The local data: the group, the character and the weight vector}\label{pbp:ss-data}

Write $A=\langle\tilde e_1,\dots,\tilde e_\varrho\rangle$, and recall the normalised frame
$(\tilde f_0,\dots,\tilde f_\varrho)=(\tilde e_1,\dots,\tilde e_\varrho,\tilde z)$ of
\S\ref{pbe:ss-real}; $A^\perp=A\oplus\langle x_1,\dots,x_k,\tilde z\rangle$. The form $h$ is linear
in the first variable; $h(\tilde e_a,\tilde e_b^*)=\delta_{ab}$, $h(x_i,x_j)=\delta_{ij}$,
$h(\tilde z,\tilde z)=-1$, all other products of basis vectors vanish.

\begin{lemma}\label{pbp:R}
$R_\tau$ is the set of $g\in G_\tau$ such that $g\tilde f_a\in\tilde f_a+\langle\tilde f_{a'}:a'<a
\rangle$ for $0\le a\le\varrho$. Every $s=uh\in R_\tau$ \textup($u\in U_\tau$, $h\in H_\tau$\textup)
satisfies $s\,x_i\in\sum_{i'}h_{i'i}x_{i'}+A$, where $(h_{i'i})$ is the matrix of $h$ in the basis
$(x_i)$. The map $s\mapsto F(s)$, where $F(s)$ is the unitriangular matrix of $s$ on the frame, is a
homomorphism; hence
\[
\tilde\sigma(s)=\sum_{a=1}^{\varrho}F(s)_{a-1,a}\qquad(\text{the coefficient of }\tilde f_{a-1}
\text{ in }s\tilde f_a,\text{ summed})
\]
is additive on $R_\tau$, and $\tilde\sigma(h)=0$ for $h\in H_\tau$.
\end{lemma}
\begin{proof}
Elements of $R_\tau$ have the stated form (\S\ref{pbe:ss-real}); and $u\in U_\tau$ preserves $A$, hence
$A^\perp$, and is trivial on $A^\perp/A$, so $ux_i\in x_i+A$, while $hx_i=\sum h_{i'i}x_{i'}$.
Conversely let $g$ be unitriangular on the frame. Then $g$ preserves each $\langle\tilde e_1,\dots,
\tilde e_a\rangle$, hence $A$ and $A^\perp$, and induces an isometry $\bar g$ of the nondegenerate
space $A^\perp/A=\langle\bar x_1,\dots,\bar x_k,\bar z\rangle$ with $\bar g\bar z=\bar z$. So $\bar g$
preserves $\bar z^\perp=\langle\bar x_i\rangle$ and acts there by a unitary matrix $h$; let $h\in
H_\tau$ also denote the element acting by $h$ on the $x_i$ and trivially on $\tilde e_a,\tilde z,
\tilde e_a^*$. Then $u=gh^{-1}$ preserves the flag $\langle\tilde e_1\rangle\subset\dots\subset A$, acts
trivially on its graded pieces (as $g$ is unitriangular and $h$ fixes the $\tilde e_a$) and trivially
on $A^\perp/A$; so $u\in U_\tau$ and $g=uh\in R_\tau$. The span of the frame is $A\oplus\langle\tilde z
\rangle$ and is preserved by $R_\tau$, so $F(s_1s_2)=F(s_1)F(s_2)$; the entries just above the
diagonal of a product of unitriangular matrices are the sums of those of the factors.
\end{proof}

\begin{lemma}[The character]\label{pbp:psi}
Let $\sigma$ be as in \eqref{pb:eq-psi}, $\kappa_\tau$ as in \eqref{pbe:eq-kappa}. For $u\in U_\tau$,
$\tilde\sigma(u)=\sigma(u)/\kappa_\tau$, and
\[
\aubar{\psi_\tau(u)}=\exp\bigl(\epsilon_0\,(\aubar{\tilde\sigma(u)}-\tilde\sigma(u))\bigr),
\]
where $\epsilon_0=+1$ if the component of $\psi_E$ at the place $\tau$ is $w\mapsto e^{-2\pi i(w+
\bar w)}$ and $\epsilon_0=-1$ if it is $w\mapsto e^{2\pi i(w+\bar w)}$.
\end{lemma}
\begin{proof}
Here $N$ is the level of \S\ref{pbe:ss-real}. From $\tilde e_a=\kappa_\tau^{\varrho-a+1}e_a/\sqrt N$, $\tilde e_a^*=\sqrt N e_a^*/\kappa_\tau^{\varrho-a+1}$
and $\tilde z=z/\sqrt N$ ($\kappa_\tau$ is real), $h(u\tilde e_{a+1},\tilde e_a^*)=h(ue_{a+1},
e_a^*)/\kappa_\tau$ and $h(u\tilde z,\tilde e_\varrho^*)=h(uz,e_\varrho^*)/\kappa_\tau$. Since $h(\cdot,
\tilde e_a^*)$ is the coefficient of $\tilde e_a$, these are the entries $F(u)_{a-1,a}$, and
$\tilde\sigma(u)=\sigma(u)/\kappa_\tau$. At $\tau$, $\psi_\tau(u)$ is the component of $\psi_E$ at
$w=-\sigma(u)/(M\tau(\varpi))$, where now $\sigma(u)$ is computed in $Y_\tau$. As $M\tau(\varpi)=
2\pi i\kappa_\tau$, $2\pi iw=-\sigma(u)/\kappa_\tau=-\tilde\sigma(u)$ and $2\pi i\bar w=
\aubar{\tilde\sigma(u)}$. So $\psi_\tau(u)=\exp(\mp(\aubar{\tilde\sigma}-\tilde\sigma))$, with the upper
sign for $e^{-2\pi i(w+\bar w)}$; the exponent is imaginary, and conjugation changes its sign.
\end{proof}

The paper uses the normalisation of $\psi_E$ only through \S\ref{pb:ss-layer2}; the next lemma shows
that the sign $\epsilon_0$ does not affect \textup{(DS4)}.

\begin{lemma}[The sign]\label{pbp:sign}
Let $\delta\in\operatorname{End}(Y_\tau)$ be $\delta\tilde e_a=(-1)^a\tilde e_a$, $\delta\tilde e_a^*=
(-1)^a\tilde e_a^*$, $\delta\tilde z=(-1)^{\varrho+1}\tilde z$, $\delta x_i=x_i$. Then $\delta\in G_\tau$,
$\delta$ is an isometry for $\|\cdot\|$, $\delta R_\tau\delta^{-1}=R_\tau$, $s\mapsto\delta s\delta^{-1}$
preserves the Haar measure of $R_\tau$, $\tilde\sigma(\delta s\delta^{-1})=-\tilde\sigma(s)$, the
images of $s$ and of $\delta s\delta^{-1}$ in $H_\tau$ coincide, $\pbel_\tau(\delta^{-1}y)=\pbel_\tau(y)$
and $\pbvol(\pbeB_{\delta^{-1}y})=\pbvol(\pbeB_y)$ for $y\in G_\tau$.
\end{lemma}
\begin{proof}
$\delta$ is diagonal with entries $\pm1$ in the orthonormal normalised basis and preserves the
products of the basis vectors, so it is in $G_\tau$ and isometric. We have $\delta\tilde f_a=
(-1)^{a+1}\tilde f_a$ for all $0\le a\le\varrho$; so $\delta s\delta^{-1}$ is unitriangular on the frame
with $F(\delta s\delta^{-1})_{a'a}=(-1)^{a+a'}F(s)_{a'a}$, hence lies in $R_\tau$
(Lemma~\ref{pbp:R}), and the entries just above the diagonal change sign. As $\delta$ is trivial on
the $x_i$ and preserves $A$, the actions of $s$ and $\delta s\delta^{-1}$ on $A^\perp/A$ agree on the
$\bar x_i$. The automorphism $s\mapsto\delta s\delta^{-1}$ has order $2$, so it multiplies the Haar
measure by a positive number whose square is $1$. Finally $\|s\delta^{-1}y\|=\|\delta^{-1}(\delta s
\delta^{-1})y\|=\|(\delta s\delta^{-1})y\|$, and $\pbeB_{\delta^{-1}y}=\delta^{-1}\pbeB_y\delta$.
\end{proof}

\begin{lemma}[The weight vector]\label{pbp:W}
As a representation of $H_\tau=\auU(k)$ in the basis $(x_i)$, $\pbW_\tau$ is isomorphic to the space
of homogeneous polynomials of degree $a_\tau$ in $k$ variables, with $(\tau_\tau(h)P)(t)=P(h^{\mathrm t}
t)$, i.e.\ $\tau_\tau(h)$ sends the monomial $t_j^{a_\tau}$ to $(\sum_ih_{ij}t_i)^{a_\tau}$; its
weights are the $\sum_i\alpha_ie_i$ with $\alpha_i\ge0$, $\sum\alpha_i=a_\tau$, and the highest weight
is $(a_\tau,0,\dots,0)$. For the unit vector $v_\tau=t_k^{a_\tau}/\|t_k^{a_\tau}\|$ \textup(any
invariant inner product\textup),
\begin{equation}\label{pbp:eq-cu}
c'(h):=\pbip{v_\tau}{\tau_\tau(h)v_\tau}=h_{kk}^{\,a_\tau}\qquad(h\in H_\tau).
\end{equation}
\end{lemma}
\begin{proof}
By {\cite[\S14.3]{P46}}, $W_{\lambda^X}$ at $\tau$ is the dual $V_\tau$ of the space of forms of degree
$a_\tau$, in its model as the dual of the Schur module spanned by the products of minors ($1\times1$
minors here), on which $\mathrm{GL}_k(\C)$ acts polynomially by the Schur matrices; so as a
$\mathrm{GL}_k(\C)$-module $W_{\lambda^X,\tau}$ is the polynomial representation $\operatorname{Sym}^{a_\tau}$
of highest weight $(a_\tau,0,\dots,0)$. The group $H_\tau$ acts on it through $h\mapsto\tau(\bar h)=
{}^{\mathrm t}h^{-1}$ ({\cite[\S14.3]{P46}}). The space $\pbW_\tau$ is the dual of $W_{\lambda^X,\tau}$ for the
bilinear pairing $\pbip w\xi=\sum w_i\xi_i$ with $\pbip{\tau_\infty(h)w}{h\xi}=\pbip w\xi$
(\S\ref{pb:ss-layer2}); so $\tau_\tau(h)$ is the transpose inverse of $\operatorname{Sym}^{a_\tau}
({}^{\mathrm t}h^{-1})$, which is $\operatorname{Sym}^{a_\tau}(h)$. In the model of polynomials this is the
stated action. The monomials are weight vectors and are orthogonal for any $\auU(k)$-invariant inner
product (distinct weights); the coefficient of $t_k^{a_\tau}$ in $(\sum_ih_{ik}t_i)^{a_\tau}$ is
$h_{kk}^{a_\tau}$, which gives \eqref{pbp:eq-cu}.
\end{proof}

\begin{remark}[What is used of the data]\label{pbp:rem-doubts}
We record: (1) Sign and scaling: by
Lemma~\ref{pbp:psi} the character in the normalised frame is $\exp(\pm(\aubar{\tilde\sigma}-
\tilde\sigma))$ with $\tilde\sigma=\sigma/\kappa_\tau$, i.e.\ the scaling by $\kappa_\tau$ is exactly
absorbed by the normalised frame; the sign depends on the normalisation of $\psi_E$ at infinity and is
irrelevant by Lemma~\ref{pbp:sign}. (2) $\pbW_\tau\cong\operatorname{Sym}^{a_\tau}$ has the extremal weight
vector $t_k^{a_\tau}$ of weight $a_\tau e_k$, a single nonnegative entry (Lemma~\ref{pbp:W}); this
uses the model of $W_\lambda$ in {\cite[\S14.3]{P46}} (the Schur matrices are polynomial in $g$). (3) $\kappa(\tau)$ is an integer by (D4) and (D6) of
{\cite[Definition~13.3]{P46}} (Lemma~\ref{pbp:sizes}).
\end{remark}

\emph{Reduction.} Let $y\in G_\tau$ and $v_\tau$ be as in Lemma~\ref{pbp:W}. For $s=uh\in R_\tau$,
$\pi_R(s)^{-1}=\tau_\tau(h)^{-1}\psi_\tau(u)^{-1}$ and $\pbip{v_\tau}{\tau_\tau(h^{-1})v_\tau}=
\aubar{c'(h)}$, so by Lemmas~\ref{pbp:R} and \ref{pbp:psi}
\begin{equation}\label{pbp:eq-integrand}
c_\tau(y^{-1}sy)\pbip{v_\tau}{\pi_R(s)^{-1}v_\tau}=c_\tau(y^{-1}sy)\,\aubar{c'(h_s)}\,
\exp\bigl(\epsilon_0(\aubar{\tilde\sigma(s)}-\tilde\sigma(s))\bigr),
\end{equation}
$h_s$ the image of $s$ in $H_\tau$. By Lemma~\ref{pbp:sign}, the substitution $s\mapsto\delta s
\delta^{-1}$ turns the integral of \eqref{pbp:eq-integrand} for $\epsilon_0$ and $y$ into the integral
for $-\epsilon_0$ and $\delta^{-1}y$, without changing $\pbel_\tau$ and $\pbvol(\pbeB)$. \emph{It
therefore suffices to prove Theorem~\ref{pbp:main} with the integrand
\eqref{pbp:eq-integrand} for $\epsilon_0=+1$}, which we assume from now on.

\subsection{Positions and the two traces}\label{pbp:ss-pos}

For $l\in X$ let $e_l$ be the corresponding standard basis vector of $\C^r$.

\begin{definition}\label{pbp:def-pos}
A \emph{stationary position} is a bijective linear map $\iota\colon\C^r\to Y_\tau$ with
$h(\iota v,\iota w)=B(w,v)$ \textup(an isometry\textup) such that, for $\widehat\Lambda=\iota\mu
\iota^{-1}$,
\begin{enumerate}
\item[(S1)] $\widehat\Lambda\tilde f_a=\tilde f_{a+1}$ for $0\le a<\varrho$;
\item[(S2)] $h(\widehat\Lambda x_i,x_{i'})=\Lambda'_i\delta_{ii'}$ for $1\le i,i'\le k$;
\item[(S3)] $\iota e_l=x_l$ for $l\in X$.
\end{enumerate}
For $s\in R_\tau$ we then put $g_s=\iota^{-1}s\iota\in G$ and write $h_s\in\auU(k)$ for the image of
$s$ in $H_\tau$ (Lemma~\ref{pbp:R}), and
\[
\varphi(s)=\pbptr\bigl(\mu(g_s-1)\bigr)-\pbptr\bigl(\Lambda'(h_s-1)\bigr).
\]
\end{definition}

(The order of the arguments of $B$ only reflects that $h$ is linear in the first variable and $B$ in
the second.) Since $\mu$ is real diagonal and commutes with $J$, $\widehat\Lambda$ is self-adjoint for
$h$; and $h(sw,w')=h(w,s^{-1}w')$ for $s\in G_\tau$.

\begin{lemma}[Two traces]\label{pbp:traces}
Let $\iota$ be a stationary position and $s\in R_\tau$. Then
\begin{enumerate}
\item[(a)] $\pbptr(g_s-1)=\pbptr(h_s-1)$;
\item[(b)] $\varphi(s)=\tilde\sigma(s)-\aubar{\tilde\sigma(s)}$; in particular $\varphi$ is additive on
$R_\tau$ and purely imaginary;
\item[(c)] for every unitary $u$ on the coordinates $X$, the element $s_u\in H_\tau$ acting by $u$ on
$\langle x_l:l\in X\rangle$ and trivially on $x_k$, $\tilde e_a$, $\tilde z$, $\tilde e_a^*$ satisfies:
$g_{s_u}$ is $u$ on the coordinates $X$ and the identity on the others, $h_{s_u}=u\oplus1$, and
$\varphi(s_u)=0$.
\end{enumerate}
\end{lemma}
\begin{proof}
The coefficients of a vector $w$ on the basis $\tilde e_a,\tilde e_a^*,\tilde z,x_i$ are $h(w,\tilde
e_a^*)$, $h(w,\tilde e_a)$, $-h(w,\tilde z)$, $h(w,x_i)$, so for $T\in\operatorname{End}(Y_\tau)$
\begin{equation}\label{pbp:eq-trace}
\pbptr T=\sum_{a=1}^\varrho\bigl(h(T\tilde e_a,\tilde e_a^*)+h(T\tilde e_a^*,\tilde e_a)\bigr)
-h(T\tilde z,\tilde z)+\sum_{i=1}^kh(Tx_i,x_i).
\end{equation}
Let $Y=s-1$ and $Y^\dagger=s^{-1}-1$, so $h(Yw,w')=h(w,Y^\dagger w')$; $s^{-1}\in R_\tau$. By
Lemma~\ref{pbp:R}, $Y\tilde f_a\in\langle\tilde f_{a'}:a'<a\rangle$, $Y\tilde z\in A$ and $Yx_i\in
\sum_{i'}(h_s-1)_{i'i}x_{i'}+A$, and the same holds for $Y^\dagger$ with $h_s^{-1}$.

(b) Apply \eqref{pbp:eq-trace} to $T=\widehat\Lambda Y$; then $\pbptr T=\pbptr(\mu(g_s-1))$. First,
$Y\tilde e_a\in\langle\tilde e_1,\dots,\tilde e_{a-1}\rangle$ and $\widehat\Lambda\tilde e_b=\tilde e_{b+1}$
for $b<\varrho$, so $h(\widehat\Lambda Y\tilde e_a,\tilde e_a^*)$ is the coefficient of $\tilde e_{a-1}$ in
$Y\tilde e_a$, i.e.\ $F(s)_{a-2,a-1}$ ($2\le a\le\varrho$; zero for $a=1$). Second, write $Y\tilde z=
\sum_bc_b\tilde e_b$; $-h(\widehat\Lambda Y\tilde z,\tilde z)=-\sum_bc_bh(\widehat\Lambda\tilde e_b,
\tilde z)$, and $h(\widehat\Lambda\tilde e_b,\tilde z)$ is $h(\tilde e_{b+1},\tilde z)=0$ for $b<\varrho$
and $h(\tilde z,\tilde z)=-1$ for $b=\varrho$ (by (S1), $\widehat\Lambda\tilde f_{\varrho-1}=\tilde f_\varrho$);
so this term is $c_\varrho=F(s)_{\varrho-1,\varrho}$. The two give $\tilde\sigma(s)$. Third, $h(\widehat
\Lambda Yx_i,x_i)=h(Yx_i,\widehat\Lambda x_i)$; for $\alpha\in A$, $h(\alpha,\widehat\Lambda x_i)=
h(\widehat\Lambda\alpha,x_i)=0$ because $\widehat\Lambda A\subseteq\langle\tilde e_2,\dots,\tilde e_\varrho,
\tilde z\rangle\perp x_i$; and $h(x_{i'},\widehat\Lambda x_i)=\aubar{h(\widehat\Lambda x_i,x_{i'})}=
\Lambda'_i\delta_{ii'}$ by (S2). So the third sum is $\sum_i(h_s-1)_{ii}\Lambda'_i=\pbptr(\Lambda'
(h_s-1))$. Fourth, $h(\widehat\Lambda Y\tilde e_a^*,\tilde e_a)=h(\tilde e_a^*,Y^\dagger\widehat\Lambda
\tilde e_a)=\aubar{h(Y^\dagger\widehat\Lambda\tilde e_a,\tilde e_a^*)}$, and $\widehat\Lambda\tilde e_a=
\tilde f_a$; summing over $a$ gives $\aubar{\tilde\sigma(s^{-1})}=-\aubar{\tilde\sigma(s)}$
(Lemma~\ref{pbp:R}). Hence $\pbptr(\mu(g_s-1))=\tilde\sigma(s)+\pbptr(\Lambda'(h_s-1))-
\aubar{\tilde\sigma(s)}$.

(a) The same computation with $\widehat\Lambda$ replaced by the identity: $h(Y\tilde e_a,\tilde e_a^*)=0$,
$h(Y\tilde z,\tilde z)=0$ ($Y\tilde z\in A\perp\tilde z$), the third sum is $\pbptr(h_s-1)$, and
$h(Y\tilde e_a^*,\tilde e_a)=\aubar{h(Y^\dagger\tilde e_a,\tilde e_a^*)}=0$.

(c) By (S3), $g_{s_u}e_l=\iota^{-1}s_ux_l=\sum_{l'}u_{l'l}e_{l'}$ for $l\in X$. The isometry $\iota$ maps
the $B$-orthogonal complement of $\langle e_l:l\in X\rangle$, which is spanned by the other coordinate
vectors, onto the $h$-orthogonal complement of $\langle x_l:l\in X\rangle$, which $s_u$ fixes
pointwise. Finally $s_u\in H_\tau$ fixes the frame, so $\tilde\sigma(s_u)=0$ and $\varphi(s_u)=0$ by (b).
\end{proof}

\emph{The integrand.} If $\iota$ is a stationary position and $y=\iota p^{-1}\in G_\tau$, then
$c_\tau(y^{-1}sy)=c(p^{-1}y^{-1}syp)=c(g_s)$, and by Lemma~\ref{pbp:traces}(b) the integrand
\eqref{pbp:eq-integrand} (with $\epsilon_0=1$) is
\begin{equation}\label{pbp:eq-theta}
\theta(s)=c(g_s)\,\aubar{c'(h_s)}\,e^{-\varphi(s)},\qquad|\theta(s)|=|c(g_s)|\,|c'(h_s)|.
\end{equation}

\subsection{The stationary position}\label{pbp:ss-stationary}

Let $\mu_0>\dots>\mu_{2n-1}$ be as in \eqref{pbp:eq-steep}, $p(t)=\prod_{i=0}^{2n-1}(t-\mu_i)$ and
$q(t)=t-a_\tau$. As $a_\tau\ge0>\mu_i$, $q(\mu_i)<0$ and $p(a_\tau)>0$. For a polynomial $F$ let
$[F]\in\R^r\subset\C^r$ be the vector that vanishes on $X$ and has the entry
\[
[F]_i=F(\mu_i)\,\bigl|p'(\mu_i)q(\mu_i)\bigr|^{-1/2}
\]
on the coordinate numbered $i$. Then $\mu[F]=[tF]$. For real vectors $v,w$ vanishing on $X$,
$B(v,w)=\sum_i(-1)^iv_iw_i$.

\begin{lemma}[Residue sums]\label{pbp:residue}
Let $F$ be a polynomial of degree at most $2n-1$ and $c\in\R$ not a node. Then
$\sum_iF(\mu_i)/p'(\mu_i)$ is the coefficient of $t^{2n-1}$ in $F$, and $\sum_iF(\mu_i)/((\mu_i-c)
p'(\mu_i))=-F(c)/p(c)$.
\end{lemma}
\begin{proof}
By Lagrange interpolation $F(t)=\sum_iF(\mu_i)\prod_{i'\ne i}(t-\mu_{i'})/p'(\mu_i)$; compare the
coefficients of $t^{2n-1}$. For the second, $H(t)=(F(t)p(c)-F(c)p(t))/(t-c)$ is a polynomial of degree
at most $2n-1$ whose coefficient of $t^{2n-1}$ is $-F(c)$ ($p$ is monic of degree $2n$, $\deg F<2n$), and
$H(\mu_i)=F(\mu_i)p(c)/(\mu_i-c)$; apply the first statement to $H$.
\end{proof}

\begin{lemma}\label{pbp:gram}
\textup{(i)} For $a'+b'\le2\varrho$, $B([t^{a'}q],[t^{b'}q])=-1$ if $a'+b'=2\varrho$ and $0$ otherwise.
\textup{(ii)} $B([1],[t^bq])=0$ for $b\le2\varrho$. \textup{(iii)} $B([1],[1])=1/p(a_\tau)$.
\textup{(iv)} For $b\le2\varrho$ every entry of $[t^bq]$ has modulus at most $(3B+a_\tau)^{2n}$.
\end{lemma}
\begin{proof}
Among the $\mu_{i'}$, exactly $i$ exceed $\mu_i$, so $p'(\mu_i)=\prod_{i'\ne i}(\mu_i-\mu_{i'})$ has the
sign $(-1)^i$, and $|p'(\mu_i)q(\mu_i)|=(-1)^{i+1}p'(\mu_i)q(\mu_i)$. Hence
$B([F_1],[F_2])=-\sum_iF_1F_2(\mu_i)/(p'(\mu_i)q(\mu_i))$. (i) This is $-\sum_i\mu_i^{a'+b'}q(\mu_i)/
p'(\mu_i)$, minus the coefficient of $t^{2n-1}$ in $t^{a'+b'}q$, which has degree $a'+b'+1\le2n-1$.
(ii) It is $-\sum_i\mu_i^b/p'(\mu_i)=0$ since $b<2n-1$. (iii) It is $-\sum_i1/((\mu_i-a_\tau)
p'(\mu_i))=1/p(a_\tau)$ by Lemma~\ref{pbp:residue}. (iv) The $\mu_i$ are distinct integers, so
$|p'(\mu_i)|\ge1$, and $|[t^bq]_i|\le|\mu_i|^b|q(\mu_i)|^{1/2}\le(3B)^b(3B+a_\tau)^{1/2}$.
\end{proof}

\begin{lemma}[Witt's theorem in the form used]\label{pbp:tb-witt}
Let $F$ be a field of characteristic $\ne2$ with an involution $a\mapsto\aubar a$ \textup(possibly
trivial\textup), $F_0$ its fixed field, and $V,V'$ finite-dimensional $F$-spaces with non-degenerate
hermitian forms $h,h'$ \textup(linear in the first variable\textup) that are isometric.
\textup{(a)} If $e\in V$, $e'\in V$ and $h(e,e)=h(e',e')\ne0$, some isometry $g$ of $V$ has $ge=e'$.
\textup{(b)} Every injective linear isometry $\phi\colon S\to V'$ from a subspace $S\subseteq V$ extends to
a bijective isometry $V\to V'$.
\end{lemma}
\begin{proof}
For $u\in V$ with $H=h(u,u)\ne0$ and $\varepsilon\in F$ with $\varepsilon\aubar\varepsilon=1$, the map
$\sigma(x)=x-kh(x,u)u$, $k=(1-\varepsilon)/H$, satisfies $h(\sigma x,\sigma y)=h(x,y)+h(x,u)h(u,y)
(|k|^2H-k-\aubar k)=h(x,y)$, since $(1-\varepsilon)(1-\aubar\varepsilon)=(1-\varepsilon)+(1-\aubar
\varepsilon)$; it fixes $u^\perp$ and is injective, hence an isometry of $V$.

(a) Let $c=h(e,e)\in F_0$ and $b=h(e,e')$. For $u=e-e'$ put $z=h(e,u)=c-b$; then $h(u,u)=z+\aubar z$.
If $z+\aubar z\ne0$, take $\varepsilon=-\aubar z/z$; then $1-\varepsilon=(z+\aubar z)/z$, so
$\sigma(e)=e-u=e'$. Otherwise $b+\aubar b=2c$, and the same construction for the pair $(e,-e')$ applies,
because there $z+\aubar z=2c+b+\aubar b=4c\ne0$; it gives $g$ with $ge=-e'$, and $-g$ works.

(b) Composing with an isometry $V'\to V$ we may take $V'=V$. \emph{Cancellation.} A non-degenerate
space $X\ne0$ has an anisotropic vector (if $h(x,x)=0$ for all $x$ then $h(x,y)+h(y,x)=0$; replacing $y$
by $\lambda y$ with $\lambda\ne\aubar\lambda$, or using $h(x,y)=h(y,x)$ when the involution is trivial,
gives $h=0$), hence an orthogonal basis. If $X=Fe\perp U=Fe'\perp U'$ with $h(e,e)=h(e',e')\ne0$, then by
(a) $U'=e'^\perp=g(e^\perp)=g(U)\cong U$. By induction over an orthogonal basis: if $T,T'\subseteq V$ are
isometric non-degenerate subspaces, then $T^\perp\cong T'^\perp$. \emph{Non-degenerate $S$.} Then
$V=S\perp S^\perp=\phi S\perp(\phi S)^\perp$, and $\phi$ together with an isometry $S^\perp\to(\phi
S)^\perp$ is the extension. \emph{General $S$.} Let $R$ be the radical of $S$, with basis $r_1,\dots,
r_k$, and $S=R\oplus S_0$; $S_0$ is non-degenerate and $R\subseteq S_0^\perp$, which is non-degenerate.
So there are $y'_j\in S_0^\perp$ with $h(r_i,y'_j)=\delta_{ij}$; with $g_{ij}=h(y'_i,y'_j)$ the vectors
$y_j=y'_j-\tfrac12\sum_lg_{jl}r_l$ satisfy $h(r_i,y_j)=\delta_{ij}$ and $h(y_i,y_j)=g_{ij}-\tfrac12
g_{ij}-\tfrac12\aubar g_{ji}=0$. Then $\tilde S=S_0\oplus\langle r_i,y_i\rangle$ is non-degenerate.
Since $\phi(R)$ is the radical of $\phi(S)=\phi(R)\oplus\phi(S_0)$, the same construction gives $y''_j$
for $\phi(r_j)$ in $\phi(S_0)^\perp$, and $\phi$ extended by $y_j\mapsto y''_j$ is an injective isometry
of $\tilde S$; apply the non-degenerate case. (See also \cite[Chapter~7, Theorem~9.1]{Scharlau}.)
\end{proof}

\begin{proposition}\label{pbp:stationary}
There is a stationary position $\iota$ such that $\iota^{-1}\tilde f_a=\mu^av_1$ for $0\le a\le\varrho$,
where $v_1=[q]$; moreover, for $a,b\le2\varrho$, $\|\mu^av_1\|\le M_0:=\sqrt r\,(3B+a_\tau)^{2n}$ and
$B(\mu^av_1,\mu^bv_1)=-1$ if $a+b=2\varrho$, $0$ if $a+b<2\varrho$.
\end{proposition}
\begin{proof}
Put $v_j=[t^{j-1}q]=\mu^{j-1}v_1$ ($1\le j\le n$) and $w=p(a_\tau)^{1/2}[1]$. The vectors $v_1,\dots,
v_n,w$ are the evaluation vectors, scaled by nonzero factors, of the $n+1$ polynomials $t^{j-1}q$, $1$
of degree at most $n<2n$ at the $2n$ distinct nodes, so they are linearly independent; with the $e_l$,
$l\in X$, which have disjoint supports, they form a basis of a subspace $S\subseteq\C^r$. By
Lemma~\ref{pbp:gram} and disjointness of supports, the Gram matrix of $(v_1,\dots,v_n;e_l\,(l\in X);w)$
for $B$ is that of $(\tilde f_0,\dots,\tilde f_\varrho;x_l\,(l\in X);x_k)$ for $h$: all products vanish
except $B(v_n,v_n)=-1=h(\tilde z,\tilde z)$, $B(e_l,e_l)=1$ and $B(w,w)=1$ ((i) with $a'=b'=\varrho$,
(ii), (iii)). The second family is also linearly independent. By Witt's theorem (Lemma~\ref{pbp:tb-witt}(b),
with $V=\C^r$ and $V'=Y_\tau$, which are isometric via $p$), the linear map $S\to Y_\tau$ sending the
first family to the second extends to a bijective isometry $\iota\colon\C^r\to Y_\tau$. (S3) holds by
construction, and $\iota\mu^av_1=\iota v_{a+1}=\tilde f_a$. (S1): $\widehat\Lambda\tilde f_a=\iota\mu
v_{a+1}=\iota v_{a+2}=\tilde f_{a+1}$ for $a<\varrho$. (S2): $\iota^{-1}x_l=e_l$ and $\mu e_l=0$ for
$l\in X$; and $\mu w=p(a_\tau)^{1/2}[t]=p(a_\tau)^{1/2}[q]+a_\tau w=p(a_\tau)^{1/2}v_1+a_\tau w$. So
$h(\widehat\Lambda x_i,x_{i'})=B(\iota^{-1}x_{i'},\mu\iota^{-1}x_i)$ is $0$ for $i\in X$, and for $i=k$
it is $a_\tau B(\iota^{-1}x_{i'},w)+p(a_\tau)^{1/2}B(\iota^{-1}x_{i'},v_1)=a_\tau\delta_{i'k}$. The last
claims are Lemma~\ref{pbp:gram}(i),(iv) ($[t^bq]$ has at most $r$ entries).
\end{proof}

\subsection{The point and its length}\label{pbp:ss-point}

Let $\mathsf J_1$ also denote the operator on $Y_\tau$ with matrix $\mathsf J_1$ in the normalised basis,
so that $h(v,w)=\pbip{\mathsf J_1w}{v}$; it is unitary and self-adjoint, $\mathsf J_1\tilde e_a=\tilde
e_a^*$, $\mathsf J_1\tilde e_a^*=\tilde e_a$, $\mathsf J_1x_i=x_i$, $\mathsf J_1\tilde z=-\tilde z$, and
$g\mathsf J_1g^*=\mathsf J_1$ for $g\in G_\tau$ (because $g^{-1}=\mathsf J_1g^*\mathsf J_1$).

\begin{lemma}\label{pbp:frame}
For every unitriangular $(\varrho+1)\times(\varrho+1)$ matrix $T$ there is $s\in R_\tau$ with $F(s)=T$.
\end{lemma}
\begin{proof}
Define a linear map on $A^\perp=\langle\tilde f_0,\dots,\tilde f_\varrho,x_1,\dots,x_k\rangle$ by
$\tilde f_a\mapsto\sum_{a'}T_{a'a}\tilde f_{a'}$, $x_i\mapsto x_i$. It is injective and preserves $h$:
the images of $\tilde e_1,\dots,\tilde e_\varrho$ lie in $A$, which is isotropic and orthogonal to
$\tilde z$ and the $x_i$; the image of $\tilde z$ is $\tilde z+\alpha$ with $\alpha\in A$, of square
$-1$. By Witt's theorem (Lemma~\ref{pbp:tb-witt}(b)) it extends to some $s\in G_\tau$, which lies in $R_\tau$ by
Lemma~\ref{pbp:R}.
\end{proof}

\begin{proposition}[The point]\label{pbp:point}
Let $\iota$ be the stationary position of Proposition~\ref{pbp:stationary} and $y=\iota p^{-1}\in
G_\tau$. Then $\pbel_\tau(y)\le8(\varrho+1)M_0\le b^{5n}$.
\end{proposition}
\begin{proof}
\emph{Step 1.} Put $V_a=p\mu^av_1\in Y_\tau$ for $0\le a\le2\varrho$. Since $p$ is an isometry for both
structures, $\|V_a\|\le M_0$ and $h(V_a,V_b)=-1$ if $a+b=2\varrho$, $0$ if $a+b<2\varrho$
(Proposition~\ref{pbp:stationary}); and $W_a:=y^{-1}\tilde f_a=p\iota^{-1}\tilde f_a=V_a$ for
$a\le\varrho$.

\emph{Step 2.} Let $o_0,\dots,o_\varrho$ be the Gram--Schmidt orthogonalisation of $W_0,\dots,W_\varrho$
(Lemma~\ref{pbe:len}(c)), so that $o_a=W_a-\sum_{a'<a}t_{a'a}W_{a'}$ for some numbers $t_{a'a}$. Then
$\|o_a\|\le\|W_a\|\le M_0$, and $h(V_{2\varrho-a},o_a)=h(V_{2\varrho-a},V_a)-\sum_{a'<a}\aubar{t_{a'a}}
h(V_{2\varrho-a},V_{a'})=-1$; by \eqref{pbe:eq-norm}, $1\le\|V_{2\varrho-a}\|\,\|o_a\|$, so
$\|o_a\|\ge M_0^{-1}$.

\emph{Step 3.} By Lemma~\ref{pbp:frame} there is $s_1\in R_\tau$ with $s_1^{-1}\tilde f_a=\tilde f_a-
\sum_{a'<a}t_{a'a}\tilde f_{a'}$; then $y_1=s_1y$ has $y_1^{-1}\tilde f_a=o_a$.

\emph{Step 4.} For $1\le a\le\varrho$ put $b_a=\mathsf J_1o_{a-1}/\|o_{a-1}\|^2$, of norm $\|o_{a-1}\|^{-1}
\le M_0$. Then $h(o_c,b_a)=\pbip{o_{a-1}}{o_c}/\|o_{a-1}\|^2=\delta_{c,a-1}$, and $h(b_a,b_{a'})$ is a
multiple of $\aubar{h(o_{a'-1},o_{a-1})}=\aubar{h(\tilde e_{a'},\tilde e_a)}=0$. Hence $\beta_a=y_1b_a$
satisfies $h(\tilde e_c,\beta_a)=\delta_{ca}$, $h(\tilde z,\beta_a)=0$, $h(\beta_a,\beta_{a'})=0$. Define
$q\in\operatorname{End}(Y_\tau)$ by $q\tilde e_a=\tilde e_a$, $q\tilde z=\tilde z$, $q\tilde e_a^*=\beta_a$,
$qx_i=x_i-\sum_ah(x_i,\beta_a)\tilde e_a$. The images of the basis have the Gram matrix of the basis:
the only products to check are $h(qx_i,qx_j)=h(x_i,x_j)$ (the added vectors lie in $A$, which is
isotropic and orthogonal to the $x_j$) and $h(qx_i,\beta_b)=h(x_i,\beta_b)-\sum_ah(x_i,\beta_a)
\delta_{ab}=0$. So $q\in G_\tau$; it fixes the frame, so $q\in R_\tau$ (Lemma~\ref{pbp:R}).

\emph{Step 5.} Let $y_2=q^{-1}y_1\in R_\tau y$ and $\mathcal A=y_2^{-1}=y_1^{-1}q$. Then $\mathcal A\tilde
e_a=o_{a-1}$, $\mathcal A\tilde z=o_\varrho$, $\mathcal A\tilde e_a^*=b_a$, all of norm at most $M_0$. The
vectors $(1-\mathsf J_1)w$, $w$ in the normalised basis, are $\tilde e_a-\tilde e_a^*$, $\tilde e_a^*-
\tilde e_a$, $2\tilde z$ and $0$; so $\|\mathcal A(1-\mathsf J_1)\|\le\sum_w\|\mathcal A(1-\mathsf J_1)w\|
\le(4\varrho+2)M_0$. From $\mathcal A\mathsf J_1\mathcal A^*=\mathsf J_1$,
\[
\|\mathcal A\|^2=\|\mathcal A\mathcal A^*\|=\|\mathsf J_1+\mathcal A(1-\mathsf J_1)\mathcal A^*\|\le
1+(4\varrho+2)M_0\|\mathcal A\|,
\]
and as $\|\mathcal A\|\ge1$, $\|\mathcal A\|\le1+(4\varrho+2)M_0$. By Lemma~\ref{pbe:len}(a) and
\eqref{pbe:eq-norm}, $\pbel_\tau(y)=\pbel_\tau(y_2)\le\|y_2\|+1=\|\mathcal A\|+1\le8(\varrho+1)M_0$.
Finally $8nM_0=8n\sqrt r(3B+a_\tau)^{2n}\le8nr\,b^{4n}\le b^{5n}$ by Lemma~\ref{pbp:sizes}.
\end{proof}

From now on $\iota$ and $y$ are as in Propositions~\ref{pbp:stationary} and \ref{pbp:point}, $g_s=
\iota^{-1}s\iota$, and $\theta$ is \eqref{pbp:eq-theta}.

\subsection{One minor to second order}\label{pbp:ss-minor}

In this subsection all norms of matrices are Frobenius norms, $\|\cdot\|=\pbpF{\cdot}$, and the
dimension is arbitrary; no constant depends on it. Let $Q$ be an orthogonal projection. For matrices
$A,B$ put
\[
\pbpct_Q(A,B)=\pbptr\bigl(QA(1-Q)B\bigr)=\pbptr(QAB)-\pbptr(QAQB),\qquad c_Q(A)=QA(1-Q)+(1-Q)AQ,
\]
and $\pbpCT_Q(A_1,A_2,A_3)=\frac12\sum_i\pbpct_Q(A_i,A_i)+\sum_{i<j}\pbpct_Q(A_i,A_j)$. Let $U=2Q-1$;
then $U^2=1$, $UQ=QU$, $Uc_Q(A)U=-c_Q(A)$, $A-c_Q(A)=QAQ+(1-Q)A(1-Q)=\frac12(A+UAU)$ commutes
with $Q$, $\|A-c_Q(A)\|\le\|A\|$, and $c_Q(A)=U[Q,A]$, so $\|c_Q(A)\|=\|[Q,A]\|$.

\begin{lemma}[Third order]\label{pbp:third}
Let $\|A_i\|\le\rho\le\frac1{10}$ \textup($i=1,2,3$\textup). Then
\[
\det\bigl(1+Q(e^{A_1}e^{A_2}e^{A_3}-1)\bigr)=\exp\Bigl(\pbptr\bigl(Q(A_1+A_2+A_3)\bigr)+
\pbpCT_Q(A_1,A_2,A_3)+r\Bigr),\qquad|r|\le50\rho^3 .
\]
\end{lemma}
\begin{proof}
Expanding the three exponential series, $e^{A_1}e^{A_2}e^{A_3}-1=L+S+R$ with $L=\sum A_i$,
$S=\sum_i\frac12A_i^2+\sum_{i<j}A_iA_j$ and $R$ the sum of the monomials $A_1^{n_1}A_2^{n_2}A_3^{n_3}/
(n_1!n_2!n_3!)$ of degree $n_1+n_2+n_3\ge3$. Bounding each monomial by $\rho^{n_1+n_2+n_3}$ (the coefficients sum to those of
$e^{3\rho}$): $\|R\|\le e^{3\rho}-1-3\rho-\frac92\rho^2\le\frac{27}6\rho^3e^{0.3}\le6.1\rho^3$,
$\|Y-L\|\le e^{3\rho}-1-3\rho\le6.1\rho^2$ and $\|Y\|\le e^{3\rho}-1\le3.5\rho$, where $Y=L+S+R$; also
$\|L\|\le3\rho$. Each monomial of $R$ begins with some $A_i$, and $|\pbptr(QA_iW)|\le\|QA_i\|\,\|W\|\le
\rho\|W\|$; so $|\pbptr(QR)|\le6.1\rho^3$ as well. Let $M=QY$, $\|M\|\le0.35$. Then $\det(1+M)=
\exp\pbptr\log(1+M)$ with $\log(1+M)=\sum_{j\ge1}(-1)^{j+1}M^j/j$, and since $|\pbptr M^j|\le\|M\|^j$
for $j\ge2$,
\[
\pbptr\log(1+M)=\pbptr M-\tfrac12\pbptr M^2+r_1,\qquad|r_1|\le\frac{\|M\|^3}{3(1-\|M\|)}\le22\rho^3.
\]
Next $\pbptr M=\pbptr(QL)+\pbptr(QS)+\pbptr(QR)$ and $\pbptr M^2=\pbptr(QLQL)+r_2$ with $r_2=
\pbptr(Q(Y-L)QY)+\pbptr(QLQ(Y-L))$, $|r_2|\le6.1\rho^2(3.5\rho+3\rho)\le40\rho^3$. Finally, as
$\pbptr(QA_iQA_j)=\pbptr(QA_jQA_i)$,
\[
\pbptr(QS)-\tfrac12\pbptr(QLQL)=\sum_i\tfrac12\bigl(\pbptr(QA_i^2)-\pbptr(QA_iQA_i)\bigr)+\sum_{i<j}
\bigl(\pbptr(QA_iA_j)-\pbptr(QA_iQA_j)\bigr)=\pbpCT_Q(A_1,A_2,A_3).
\]
So $r=r_1+\pbptr(QR)-r_2/2$ and $|r|\le(22+6.1+20)\rho^3\le50\rho^3$.
\end{proof}

\begin{proposition}[Second order with crossing parts]\label{pbp:second}
Let $\|A_i\|\le\delta\le\frac1{20}$ and $\|[Q,A_i]\|\le\zeta$ \textup($i=1,2,3$\textup). Then the identity
of Lemma~\ref{pbp:third} holds with $|r|\le400\,\delta\zeta^2$.
\end{proposition}
\begin{proof}
If $\zeta\ge\delta$, Lemma~\ref{pbp:third} with $\rho=\delta$ gives $|r|\le50\delta^3\le50\delta\zeta^2$.
If $\zeta=0$, the $A_i$ commute with $Q$. Then $1+Q(e^{A}-1)=e^{QA}$ (as $(QA)^j=QA^j$), and
$1+Q(XY-1)=(1+Q(X-1))(1+Q(Y-1))$ when $X$ commutes with $Q$; so the determinant is
$\prod_ie^{\pbptr QA_i}$, while $\pbpct_Q(A_i,A_j)=\pbptr(A_iQ(1-Q)A_j)=0$; thus $r=0$.
Let $0<\zeta<\delta$ and put $B_i(t)=A_i-c_Q(A_i)+tc_Q(A_i)$ for $t\in\C$, and
\[
g(t)=\det\bigl(1+Q(e^{B_1(t)}e^{B_2(t)}e^{B_3(t)}-1)\bigr)\exp\Bigl(-\pbptr\bigl(Q\textstyle\sum_iB_i(t)
\bigr)-\pbpCT_Q(B_1(t),B_2(t),B_3(t))\Bigr),
\]
an entire function of $t$ with $g(1)=e^r$. For $|t|<R_0:=\delta/\zeta$, $\|B_i(t)\|\le\|A_i\|+|t|\,
\|[Q,A_i]\|\le2\delta\le\frac1{10}$, so by Lemma~\ref{pbp:third} $g(t)=e^{r(t)}$ with $|r(t)|\le
50(2\delta)^3=400\delta^3\le0.05$. Hence $f=\operatorname{Log}g$ (principal branch) is holomorphic on the
disc $|t|<R_0$, $f(t)=r(t)$ there, $|f|\le400\delta^3$, and $f(1)=r$. The $B_i(0)$ commute with $Q$, so
$f(0)=0$ by the case $\zeta=0$. Conjugation by $U$ sends $B_i(t)$ to $B_i(-t)$ and leaves the three
ingredients of $g$ unchanged: $\det(1+Q(UXU-1))=\det(U(1+Q(X-1))U)=\det(1+Q(X-1))$,
$\pbptr(QUBU)=\pbptr(QB)$, and $\pbpct_Q(UAU,UBU)=\pbptr(UQA(1-Q)BU)=\pbpct_Q(A,B)$. So $f$ is even,
$f(t)=\phi(t^2)$ with $\phi$ holomorphic on $|w|<R_0^2$, $\phi(0)=0$, $|\phi|\le400\delta^3$, and Schwarz's
lemma gives $|f(1)|=|\phi(1)|\le400\delta^3/R_0^2=400\delta\zeta^2$.
\end{proof}

\begin{lemma}[Size of one minor]\label{pbp:minorsize}
Let $k$ be unitary, $S$ a set of coordinates and $P$ the coordinate projection onto $S$. Then
$|\Delta_S(k)|^2\le\exp(-\|[P,k]\|^2/2)$.
\end{lemma}
\begin{proof}
Let $\mathcal M=(1-P)kP$ and $\mathcal Y=1-P+PkP$, so $\Delta_S(k)=\det\mathcal Y$. Using $k^*k=1$,
$\mathcal Y^*\mathcal Y=1-P+Pk^*PkP=1-\mathcal M^*\mathcal M$. The eigenvalues $\sigma_i\in[0,1]$ of
$\mathcal M^*\mathcal M$ give $|\det\mathcal Y|^2=\prod(1-\sigma_i)\le e^{-\sum\sigma_i}=e^{-\|\mathcal
M\|^2}$. Finally $\|[P,k]\|^2=\|Pk(1-P)\|^2+\|(1-P)kP\|^2$, and both terms equal $\pbptr P-
\pbptr(Pk^*PkP)$ ($k$ unitary), so $\|[P,k]\|^2=2\|\mathcal M\|^2$.
\end{proof}

\begin{proposition}[One minor near the identity]\label{pbp:minorconj}
Let $o$ be an injective real function on the coordinates, $S$ a set of coordinates with $o_i<o_j$
for $i\in S$, $j\notin S$, and $P$ the coordinate projection onto $S$. Let $Z\in\pdsg$ with
$\|Z\|\le\delta\le\frac1{20}$, and let $u_-$ \textup(resp.\ $u_+$\textup) keep the entries $Z_{ij}$ with
$\varepsilon_i=\varepsilon_j$ and $o_i<o_j$ \textup(resp.\ $o_i>o_j$\textup) and vanish elsewhere. Put
$\hat g=e^{-u_-}e^Ze^{-u_+}$. Then for every $k\in\pdsK$,
\[
\Delta_S(k^*\hat gk)=\exp\bigl(\pbptr(PZ)+\tfrac14\|[P,Z]\|^2+r\bigr),\qquad|r|\le10^3\delta\bigl(
\|[P,Z]\|^2+\|[P,k]\|^2\bigr),
\]
and $\det\hat g=e^{\pbptr Z}$.
\end{proposition}
\begin{proof}
$\det e^{-u_\pm}=e^{-\pbptr u_\pm}=1$. Let $Q=kPk^*$. As $\Delta_S(x)=\det(1+P(x-1))$ and $\det(1+
\mathcal A\mathcal B)=\det(1+\mathcal B\mathcal A)$, $\Delta_S(k^*\hat gk)=\det(1+Pk^*(\hat g-1)k)=\det(1+
Q(\hat g-1))$. Put $A_1=-u_-$, $A_2=Z$, $A_3=-u_+$, $x=\|[P,Z]\|$ and $\kappa=\|[P,k]\|$; then
$\|A_i\|\le\delta$, $Q-P=[k,P]k^*$ so $\|Q-P\|=\kappa$, and since $[P,\cdot]$ multiplies the entry
$(a,b)$ by $P_{aa}-P_{bb}$, $\|[P,A_i]\|\le x$. As $[Q,A]=[P,A]+[Q-P,A]$, $\|[Q,A_i]\|\le x+2\delta\kappa
=:\zeta\le x+\kappa$. Proposition~\ref{pbp:second} gives the formula with $\pbptr(Q(A_1+A_2+A_3))+
\pbpCT_Q(A_1,A_2,A_3)$ in place of $\pbptr(PZ)+x^2/4$ and an error at most $400\delta\zeta^2$. We compare.

(i) $W=A_1+Z+A_3$ has $W_{ab}=0$ for $a\ne b$ with $\varepsilon_a=\varepsilon_b$ (as $o$ is injective) and
$W_{aa}=Z_{aa}$, while $Q$ is block diagonal ($k\in\pdsK$ commutes with $J$). So $\pbptr(QW)-
\pbptr(PZ)=\sum_a(Q-P)_{aa}Z_{aa}$. Now $Q-P=-[P,k]k^*$, and $|[P,k]_{ab}k_{ab}|=|[P,k]_{ab}|^2$; so
$\sum_a|(Q-P)_{aa}|\le\kappa^2$ and $|\pbptr(QW)-\pbptr(PZ)|\le\delta\kappa^2$.

(ii) $\pbpCT_P(A_1,Z,A_3)=x^2/4$. Indeed $\pbpct_P(A,B)=\sum_{i\in S,j\notin S}A_{ij}B_{ji}$, and for
$i\in S$, $j\notin S$ we have $o_i<o_j$, so $(A_1)_{ij}=-Z_{ij}$ and $(A_3)_{ji}=-Z_{ji}$ if
$\varepsilon_i=\varepsilon_j$, while $(A_1)_{ji}=(A_3)_{ij}=0$, and $(A_1)_{ij}=(A_3)_{ji}=0$ if
$\varepsilon_i\ne\varepsilon_j$. As $Z\in\pdsg$, $Z_{ij}Z_{ji}=-|Z_{ji}|^2$ if $\varepsilon_i=
\varepsilon_j$ and $=|Z_{ji}|^2$ otherwise. With $C$ and $C'$ the sums of $|Z_{ji}|^2$ over $i\in S$,
$j\notin S$ with equal, resp.\ opposite, signs: $\pbpct_P(A_1,A_1)=\pbpct_P(A_3,A_3)=0$, $\pbpct_P(Z,Z)
=C'-C$, $\pbpct_P(A_1,Z)=C$, $\pbpct_P(A_1,A_3)=-C$, $\pbpct_P(Z,A_3)=C$; so $\pbpCT_P=\frac12(C+C')$.
And $x^2=\sum_{i\in S,j\notin S}(|Z_{ij}|^2+|Z_{ji}|^2)=2(C+C')$, since $|Z_{ij}|=|Z_{ji}|$.

(iii) For $R\in\{P,Q\}$, $\pbpct_R(A,B)=\pbptr(\mathcal X_R\mathcal Y_R)$ with $\mathcal X_R=RA(1-R)=
R[R,A](1-R)$ and $\mathcal Y_R=(1-R)BR$, so $\|\mathcal X_R\|\le\|[R,A]\|$, $\|\mathcal Y_R\|\le
\|[R,B]\|$; and $QA(1-Q)-PA(1-P)=(Q-P)A(1-Q)-PA(Q-P)$, so $\|\mathcal X_Q-\mathcal X_P\|\le2\delta\kappa$
for $\|A\|\le\delta$, and likewise for $\mathcal Y$. Hence $|\pbpct_Q(A_i,A_j)-\pbpct_P(A_i,A_j)|\le
2\delta\kappa\zeta+x\cdot2\delta\kappa\le4\delta\kappa\zeta$, and $|\pbpCT_Q-\pbpCT_P|\le\frac92\cdot
4\delta\kappa\zeta=18\delta\kappa\zeta$.

Altogether $|r|\le\delta\kappa^2+18\delta\kappa(x+\kappa)+400\delta(x+\kappa)^2\le10^3\delta(x^2+
\kappa^2)$.
\end{proof}

\subsection{The Gaussian shape of the coefficients}\label{pbp:ss-gauss}

Norms are again Frobenius norms unless stated otherwise. For a matrix $Y$ and an integer vector $w$
put $\pbpQ_w(Y)=\sum_{i,j}|w_i-w_j|\,|Y_{ij}|^2$. For $g\in G$ with $\|g-1\|\le\frac12$ let $\log g=
\sum_{j\ge1}(-1)^{j+1}(g-1)^j/j$; as in the proof of Lemma~\ref{pds:lip} (Step~2), $\log g\in\pdsg$,
$e^{\log g}=g$, and $\|\log g\|\le-\log(1-\|g-1\|)\le2\|g-1\|$. The same holds in $\auU(k)$, where
$\log h$ is skew-hermitian.

\begin{lemma}[Removing the unipotent factors]\label{pbp:reduced}
Let $Z\in\pdsg$ with $\|Z\|\le\frac1{20}$, and let $u_\pm$ and $\hat g$ be as in
Proposition~\ref{pbp:minorconj} for $o=\lambda$. Then
\[
c(e^Z)=d_\Lambda\int_{\pdsK}\frac{|\pdsm_\Lambda(k)|^2}{\pdsm_{-(\lambda+\rho)}(k^{-1}\hat gk)}\,dk .
\]
\end{lemma}
\begin{proof}
Let $x=e^Z$; $\|x-1\|\le e^{1/20}-1<1$ in operator norm, so $x$ is accretive and $1\in A_x$ in the
notation of {\cite[\S3.2.4]{P49}}. By {\cite[Theorem~3.17]{P49}} and {\cite[Proposition~3.21]{P49}}, $c(x)=
F(x)=d_\Lambda I(x;1)=d_\Lambda\mathcal J(1,1)$, with $\mathcal J$ and $U_x$ as in the proof of
{\cite[Proposition~3.21]{P49}}. By the second assertion of {\cite[Proposition~3.21]{P49}}, $\mathcal J$
is constant along every path in $U_x$. For $t\in[0,1]$,
$e^{tu_+}x^{-1}e^{tu_-}$ is within $e^{3/20}-1<1$ of $1$ in operator norm, hence accretive, so $(e^{tu_-},
e^{tu_+})\in U_x$, and
\[
\mathcal J(1,1)=\mathcal J(e^{u_-},e^{u_+})=\int_{\pdsK}\frac{\pdsm_\Lambda(e^{u_-}k)\,\pdsm_\Lambda(k^*
e^{u_+})}{\pdsm_{\lambda+\rho}(k^*e^{u_+}x^{-1}e^{u_-}k)}\,dk .
\]
The matrix $b=e^{u_-}$ has $b_{cd}=0$ unless $c=d$ or $\varepsilon_c=\varepsilon_d$ and $\lambda_c<
\lambda_d$ (the powers of $u_-$ have this shape), and $b_{cc}=1$; for such $(c,d)$, $\Lambda_c\le
\Lambda_d$ by {\cite[Lemma~3.18(c)]{P49}}. So {\cite[Lemma~3.18(b)]{P49}} (left version; $b$ is triangular
for the order of decreasing $\lambda$) gives $\pdsm_\Lambda(bk)=\pdsm_\Lambda(k)$. Likewise $e^{u_+}$
has its off-diagonal entries at $(c,d)$ with $\lambda_c>\lambda_d$, $\Lambda_c\ge\Lambda_d$, and
$\pdsm_\Lambda(k^*e^{u_+})=\pdsm_\Lambda(k^*)=\aubar{\pdsm_\Lambda(k)}$. Finally
$(k^*e^{u_+}x^{-1}e^{u_-}k)^{-1}=k^{-1}\hat gk$, and $\pdsm_{\lambda+\rho}(y^{-1})=\pdsm_{-(\lambda+\rho)}
(y)$ ({\cite[Lemma~3.18(a)]{P49}}).
\end{proof}

\begin{lemma}\label{pbp:dLambda}
$d_\Lambda\le e^{2T}$.
\end{lemma}
\begin{proof}
Let $M_\Lambda=\max(1,\max_i|\Lambda_i|)\le3b$ and $\varepsilon_0=1/(6M_\Lambda r)$. By
\eqref{pds:eq-m}, $\pdsm_\Lambda(k)=(\det k)^{-M_\Lambda}\prod_{t=1-M_\Lambda}^{M_\Lambda}\Delta_{S_t
(\Lambda)}(k)$. If $k\in\pdsK$ and $\|k-1\|_{\mathrm{op}}\le\varepsilon_0$, then $|\det k|=1$ and each
$\Delta_S(k)=\det(1+(k-1)_{SS})$ is a product of at most $r$ factors $1+\xi$ with $|\xi|\le\varepsilon_0
\le\frac12$, so $|\Delta_S(k)|\ge e^{-2r\varepsilon_0}$ and $|\pdsm_\Lambda(k)|\ge e^{-4M_\Lambda r
\varepsilon_0}\ge\frac12$. Choose a maximal family $k_1,\dots,k_\nu\in\pdsK$ with $\|k_ik_j^{-1}-1\|_{
\mathrm{op}}>\varepsilon_0$ for $i\ne j$. Then $\|k_i-k_j\|_{\mathrm{op}}>\varepsilon_0$, the operator norm
balls of radius $\varepsilon_0/2$ about the $k_i$ are disjoint and lie in the ball of radius $2$ about
$0$ in $M_r(\C)$, of real dimension $2r^2$, so $\nu\le(4/\varepsilon_0)^{2r^2}$; and by maximality
$\pdsK$ is covered by the right translates $\{k:\|k-1\|_{\mathrm{op}}\le\varepsilon_0\}k_i$. As $dk$ has
mass $1$, $d_\Lambda^{-1}=\int|\pdsm_\Lambda|^2\ge\frac14(\varepsilon_0/4)^{2r^2}$, so $d_\Lambda\le
4(72b^2)^{2r^2}\le b^{6r^2}\le e^{2T}$ ($T\ge3r^2\log b$).
\end{proof}

\begin{proposition}[Gaussian shape, noncompact factor]\label{pbp:gauss}
Let $g\in G$ with $\|g-1\|\le\delta_1$, $\delta_1T\le5\cdot10^{-13}$, and $Z=\log g$. There is
$\gamma\in[0,1]$ with $|c(g)e^{-\pbptr(\mu Z)}-\gamma|\le\gamma/10+e^{-100T}$.
\end{proposition}
\begin{proof}
Put $\delta=2\delta_1$, so $\|Z\|\le\delta$, $\delta T\le10^{-12}$ and $\delta\le10^{-12}$. By
Lemma~\ref{pbp:reduced}, $c(g)=d_\Lambda\int|\pdsm_\Lambda(k)|^2/\allowbreak \pdsm_{-(\lambda+\rho)}(k^{-1}\hat gk)\,dk$.
Let $w=-(\lambda+\rho)$ and $M'=\max_i|w_i|$; then $\pdsm_w(y)=(\det y)^{-M'}\prod_{t=1-M'}^{M'}
\Delta_{S_t(w)}(y)$. For $i\in S_t(w)$, $j\notin S_t(w)$ we have $(\lambda+\rho)_i<(\lambda+\rho)_j$,
hence $\lambda_i<\lambda_j$ ({\cite[Lemma~3.18(c)]{P49}}), so Proposition~\ref{pbp:minorconj} applies to
each factor with $o=\lambda$. Since exactly $M'+w_i$ of the $t\in(-M',M']$ satisfy $t\le w_i$, and exactly
$|w_i-w_j|$ separate $i$ from $j$, multiplying gives
\[
\pdsm_w(k^{-1}\hat gk)=\exp\bigl(-\pbptr((\lambda+\rho)Z)+\tfrac14\pbpQ_w(Z)+r(k)\bigr),\qquad|r(k)|\le
10^3\delta\bigl(\pbpQ_w(Z)+\pbpQ_w(k)\bigr).
\]
Now $\lambda+\rho=\Lambda+2\rho_{\mathrm c}$ with $|2\rho_{\mathrm c,i}|\le r$, so $|\pbptr(2\rho_{\mathrm c}Z)|
\le r^2\delta$ and $|w_i-w_j|\le|\Lambda_i-\Lambda_j|+2r$; with $\sum_{ij}|k_{ij}|^2=r$ this gives
$\pbpQ_w(k)\le\pbpQ_\Lambda(k)+2r^2$. Put $q=\pbpQ_w(Z)/4$. Then $\pdsm_w(k^{-1}\hat gk)=\exp(q-
\pbptr(\mu Z)-r'(k))$ with $|r'(k)|\le10^4\delta\,s(k)$, $s(k)=q+\pbpQ_\Lambda(k)+r^2$, and since
$d_\Lambda\int|\pdsm_\Lambda|^2=1$,
\[
c(g)e^{-\pbptr(\mu Z)}-e^{-q}=d_\Lambda\int_{\pdsK}|\pdsm_\Lambda(k)|^2e^{-q}\bigl(e^{r'(k)}-1\bigr)\,dk .
\]
By Lemma~\ref{pbp:minorsize}, multiplied over $t$, $|\pdsm_\Lambda(k)|^2\le e^{-\pbpQ_\Lambda(k)/2}$. If
$s(k)\le5\cdot10^4T$, then $|r'(k)|\le5\cdot10^{-4}$ and $|e^{r'}-1|\le\frac1{10}$. Otherwise $|r'(k)|\le
s(k)/4$ and $|\pdsm_\Lambda(k)|^2e^{-q}|e^{r'}-1|\le e^{-\pbpQ_\Lambda(k)/2-q+s(k)/4}\le e^{-s(k)/4+r^2/2}
\le e^{-5100T}$, as $r^2\le T$. Integrating, $|c(g)e^{-\pbptr(\mu Z)}-e^{-q}|\le e^{-q}/10+d_\Lambda
e^{-5100T}$, and $d_\Lambda\le e^{2T}$ (Lemma~\ref{pbp:dLambda}). Take $\gamma=e^{-q}$.
\end{proof}

\begin{lemma}[Gaussian shape, compact factor]\label{pbp:compact}
Let $h\in\auU(k)$ with $\|h-1\|\le\delta_1$, $\delta_1T\le5\cdot10^{-13}$, and $Z'=\log h$. There is
$\gamma'\in[0,1]$ with $|c'(h)e^{-\pbptr(\Lambda'Z')}-\gamma'|\le\gamma'/10+e^{-100T}$.
\end{lemma}
\begin{proof}
Let $\delta=2\delta_1\ge\|Z'\|$, $P$ the projection onto the coordinate $k$, and $x=\|[P,Z']\|$. By
Proposition~\ref{pbp:second} with $Q=P$, $A_1=A_3=0$, $A_2=Z'$, $\zeta=x$: $h_{kk}=\det(1+P(e^{Z'}-1))=
\exp(Z'_{kk}+\frac12\pbpct_P(Z',Z')+r)$ with $|r|\le400\delta x^2$, and $\pbpct_P(Z',Z')=\sum_{i\ne k}
Z'_{ki}Z'_{ik}=-\sum_{i\ne k}|Z'_{ik}|^2=-x^2/2$. By \eqref{pbp:eq-cu}, $c'(h)e^{-\pbptr(\Lambda'Z')}=
h_{kk}^{a_\tau}e^{-a_\tau Z'_{kk}}=\gamma'e^{a_\tau r}$ with $\gamma'=e^{-a_\tau x^2/4}$. If $a_\tau x^2\le
5\cdot10^4T$, $|a_\tau r|\le2\cdot10^{-5}$ and $\gamma'|e^{a_\tau r}-1|\le\gamma'/10$. Otherwise
$\gamma'|e^{a_\tau r}-1|\le\gamma'e^{|a_\tau r|}\le e^{-a_\tau x^2/8}\le e^{-100T}$.
\end{proof}

\subsection{Localisation}\label{pbp:ss-loc}

Norms are Frobenius norms. For $s\in R_\tau$ write $x_s=g_s-1$, $x'_s=h_s-1$, and $y_s=g_s-a(g_s)$ for
the \emph{noncompact part} of $g_s$ (the entries at $(i,j)$ with $\varepsilon_i\ne\varepsilon_j$). Let
$P_{\mathrm{off}}$ be the coordinate projection onto the coordinates outside $X$.

\begin{lemma}[The phase at the logarithms]\label{pbp:phaselog}
If $\|x_s\|,\|x'_s\|\le\frac12$, then $\pbptr(\mu\log g_s)-\pbptr(\Lambda'\log h_s)=\varphi(s)$.
\end{lemma}
\begin{proof}
For pairs of matrices put $\Phi(X,X')=\pbptr(\mu X)-\pbptr(\Lambda'X')$ (linear), multiply pairs
componentwise, and let $\xi=(x_s,x'_s)$, $\xi_t=(x_t,x'_t)$ for $t\in R_\tau$. We show $\Phi(\xi^i\xi
(1+\xi_t))=0^i\varphi(s)$ for all $i\ge0$ and $t$. As $(1+\xi)(1+\xi_t)=1+\xi_{st}$, $\xi(1+\xi_t)=\xi_{st}
-\xi_t$, so for $i=0$ the left side is $\varphi(st)-\varphi(t)=\varphi(s)$ (Lemma~\ref{pbp:traces}(b)).
And $\xi^{i+1}\xi(1+\xi_t)=\xi^i\xi(\xi_{st}-\xi_t)=\xi^i\xi(1+\xi_{st})-\xi^i\xi(1+\xi_t)$, which gives the
induction step. With $t=1$, $\Phi(\xi^{i+1})=0^i\varphi(s)$; insert this in the series for $\log$.
\end{proof}

\begin{lemma}[Rows and columns]\label{pbp:rows}
For $s\in R_\tau$: $\|x_s\|^2=2\|y_s\|^2-2\Rel\pbptr x_s$, $\|x'_s\|^2=-2\Rel\pbptr x'_s$,
$\|P_{\mathrm{off}}x_s\|^2\le6\|y_s\|^2$ and $\|x_sP_{\mathrm{off}}\|^2\le6\|y_s\|^2$.
\end{lemma}
\begin{proof}
For $g\in G$, $gJg^*=J$; its diagonal entry $i$ reads $\sum_j\varepsilon_i\varepsilon_j|g_{ij}|^2=1$, i.e.\
$\sum_j|g_{ij}|^2-2\sum_j|y_{ij}|^2=1$ with $y=g-a(g)$; with $x=g-1$, $\sum_j|x_{ij}|^2=\sum_j|g_{ij}|^2-
2\Rel x_{ii}-1$, so
\[
2\Rel x_{ii}=2\textstyle\sum_j|y_{ij}|^2-\sum_j|x_{ij}|^2 ,
\]
and for $h\in\auU(k)$ likewise $2\Rel x'_{ii}=-\sum_j|x'_{ij}|^2$. Summing over $i$ gives the two
identities. Multiply the first row identity by $-\Lambda_i\ge0$ and sum: $\sum_i(-\Lambda_i)\sum_j
|x_{ij}|^2=2\sum_i(-\Lambda_i)\sum_j|y_{ij}|^2+2\Rel\pbptr(\mu x)$. By \eqref{pbp:eq-steep} the left side is
at least $B\|P_{\mathrm{off}}x\|^2$ and the first term on the right at most $6B\|y\|^2$. Multiplying the
second row identity by $\Lambda'_i\ge0$ gives $\Rel\pbptr(\Lambda'x'_s)\le0$, and $\Rel\pbptr(\mu x_s)=
\Rel\pbptr(\Lambda'x'_s)$ because $\varphi(s)$ is imaginary. Hence $B\|P_{\mathrm{off}}x_s\|^2\le6B
\|y_s\|^2$. For the columns apply this to $s^{-1}$: $x_{s^{-1}}=Jx_s^*J$, $y_{s^{-1}}=Jy_s^*J$, and
$P_{\mathrm{off}}$ commutes with $J$.
\end{proof}

\begin{lemma}\label{pbp:leftX}
If $k_0\in\pdsK$ is the identity on the coordinates outside $X$, then $c(k_0g)=c(g)$ for $g\in G$.
\end{lemma}
\begin{proof}
By {\cite[(17)]{P49}} (applied with $k_1=k_0$ and with $k_1=1$, for the same $a\in A_g$) it suffices
that $\pdsm_\Lambda(k_0Y)=\pdsm_\Lambda(Y)$ for invertible $Y$. Let $S=S_t(\Lambda)$. If $t\ge1$, $S\cap X=
\emptyset$ ($\Lambda=0$ on $X$) and the rows of $k_0Y$ indexed by $S$ are those of $Y$, so $\Delta_S(k_0Y)
=\Delta_S(Y)$. If $t\le0$, $S\supseteq X$ and $(k_0)_{ij}=0$ for $i\in S$, $j\notin S$, so $(k_0Y)_{SS}=
(k_0)_{SS}Y_{SS}$ and $\Delta_S(k_0Y)=\det k_0\cdot\Delta_S(Y)$; this is compensated by the factor
$(\det k_0Y)^{-1}$ in \eqref{pds:eq-m}.
\end{proof}

\begin{lemma}[Far from $\pdsK$]\label{pbp:far}
For $g\in G$ with noncompact part $y$, $|c(g)|^4\le(1+\|y\|^2)^{-h_0}$.
\end{lemma}
\begin{proof}
From $g^*Jg=J$, with $J$ commuting with $a=a(g)$ and anticommuting with $y$, the block diagonal part
gives $a^*a-y^*y=1$. So $|\det a|^2=\det(1+y^*y)\ge1+\pbptr(y^*y)=1+\|y\|^2$, and
Theorem~\ref{pds:decay} with $c=h_0\le h$ (Lemma~\ref{pbp:sizes}) gives $|c(g)|^2\le|\det a|^{-h_0}$.
\end{proof}

\begin{lemma}[Localisation]\label{pbp:localise}
Let $0<\rho_1\le\frac16$ and $s\in R_\tau$ with $\|y_s\|\le\rho_1$. There is $t\in R_\tau$ with
$\theta(t)=\theta(s)$, $\|x_t\|\le4\rho_1$ and $\|x'_t\|\le4\rho_1$.
\end{lemma}
\begin{proof}
Let $\alpha=(g_s)_{XX}$ be the block of $g_s$ on $X$. The block on $X\times X$ of $g_s^*Jg_s=J$ is
$\alpha^*\alpha+\mathcal W^*J'\mathcal W=1$, where $\mathcal W$ consists of the rows outside $X$ and the
columns in $X$ of $g_s$, which are those of $x_s$, and $J'$ is the diagonal of signs outside $X$. So
$\|\alpha^*\alpha-1\|\le\|P_{\mathrm{off}}x_s\|^2\le6\rho_1^2\le\rho_1<1$ (Lemma~\ref{pbp:rows}); thus
$\alpha$ is invertible, $\alpha=k_1z_1$ with $k_1$ unitary and $z_1=(\alpha^*\alpha)^{1/2}$, and $\|z_1-1\|
\le\|z_1^2-1\|\le\rho_1$ (on eigenvalues $\nu\ge0$, $|\nu-1|\le|\nu^2-1|$). Let $u=k_1^{-1}$ and $t=s_us$
(Lemma~\ref{pbp:traces}(c)). Then $g_t=g_{s_u}g_s$ has the block $z_1$ on $X$; $c(g_t)=c(g_s)$ by
Lemma~\ref{pbp:leftX}; $(h_t)_{kk}=(h_s)_{kk}$ since $h_{s_u}=u\oplus1$, so $c'(h_t)=c'(h_s)$; and
$\varphi(t)=\varphi(s_u)+\varphi(s)=\varphi(s)$. So $\theta(t)=\theta(s)$. As $g_{s_u}$ is block diagonal and
unitary, $y_t=g_{s_u}y_s$ and $\|y_t\|=\|y_s\|$. Counting the entries of $x_t$ in the rows outside
$X$, in the columns outside $X$ and in the block on $X$, Lemma~\ref{pbp:rows} gives $\|x_t\|^2\le
6\rho_1^2+6\rho_1^2+\rho_1^2$, so $\|x_t\|\le4\rho_1$. Finally, by Lemma~\ref{pbp:rows} and
Lemma~\ref{pbp:traces}(a), $\|x'_t\|^2=-2\Rel\pbptr x'_t=-2\Rel\pbptr x_t=\|x_t\|^2-2\|y_t\|^2\le13\rho_1^2$.
\end{proof}

\begin{proposition}[The cone]\label{pbp:cone}
For every $s\in R_\tau$, $|\theta(s)|\le2\Rel\theta(s)+12e^{-100T}$.
\end{proposition}
\begin{proof}
Let $E=e^{-100T}$ and $\rho_1=10^{-13}/T\le\frac16$. If $\|y_s\|\ge\rho_1$, Lemma~\ref{pbp:far} gives
$|\theta(s)|^4\le|c(g_s)|^4\le(1+\rho_1^2)^{-h_0}$, and $h_0\log(1+\rho_1^2)\ge10^{30}T^3\rho_1^2/2=5000T$,
so $|\theta(s)|\le E\le2\Rel\theta(s)+12E$ (as $\Rel\theta\ge-|\theta|$). Otherwise take $t$ as in
Lemma~\ref{pbp:localise}, and $\delta_1=4\rho_1$, so $\delta_1T\le5\cdot10^{-13}$. Let $Z=\log g_t$,
$Z'=\log h_t$. The number $\pbptr(\Lambda'Z')$ is imaginary, so by Lemma~\ref{pbp:phaselog}
\[
\theta(t)=u_1\,u_2,\qquad u_1=c(g_t)e^{-\pbptr(\mu Z)},\quad u_2=\aubar{c'(h_t)e^{-\pbptr(\Lambda'Z')}} .
\]
By Proposition~\ref{pbp:gauss} and Lemma~\ref{pbp:compact} there are $\gamma,\gamma'\in[0,1]$ with
$|u_1-\gamma|\le\gamma/10+E$ and $|u_2-\gamma'|\le\gamma'/10+E$. Then, as $\gamma,\gamma',E\le1$,
\[
|u_1u_2-\gamma\gamma'|\le\gamma|u_2-\gamma'|+\gamma'|u_1-\gamma|+|u_1-\gamma||u_2-\gamma'|\le
0.21\gamma\gamma'+4E ,
\]
and for $z=\theta(s)=\theta(t)$ and $P=\gamma\gamma'\ge0$: $|z|-2\Rel z\le(P+|z-P|)-2(P-|z-P|)\le
-P+3(0.21P+4E)\le12E$.
\end{proof}

\subsection{Coverings and the proof of \texorpdfstring{\textup{(DS4)}}{(DS4)}}\label{pbp:ss-cone}

In this subsection $\|\cdot\|$ is the operator norm. Note $\|g_s\|=\|y^{-1}sy\|$ and $\|g_s-1\|=
\|y^{-1}sy-1\|$, since $p$ is unitary; so $\pbeB_y=\{s:\|g_s-1\|\le1\}$. The map $s\mapsto y^{-1}sy$ is a
homeomorphism of $R_\tau$ onto a closed subgroup of $G_\tau$, so the sets $\{s:\|g_s\|\le T'\}$ are
compact. Put
\[
\mathrm{Large}=\{s\in R_\tau:\|g_s\|\le\beta+1\},\qquad\mathrm{Small}=\{s\in R_\tau:\|g_s-1\|\le\rho_0,\
\|h_s-1\|\le\rho_0\}.
\]

\begin{lemma}\label{pbp:cover1}
For real $\mathsf c>4r^2$, $\int_{R_\tau}(2/\|g_s\|)^{\mathsf c}ds\le2^{\mathsf c}5^{2r^2}\frac{\mathsf c}{
\mathsf c-4r^2}\pbvol(\pbeB_y)$; for $\mathsf c=9D^2$ this is at most $16^{9D^2}\pbvol(\pbeB_y)$.
\end{lemma}
\begin{proof}
The proof of Lemma~\ref{pbe:cover} (with $x=y$ there) shows $\pbvol\{s:\|g_s\|\le T'\}\le5^{2r^2}
T'^{4r^2}\pbvol(\pbeB_y)$ for $T'\ge1$; integrate against $\mathsf cT'^{-\mathsf c-1}dT'$ as there. For
$\mathsf c=9D^2\ge9r^2$, $\frac{\mathsf c}{\mathsf c-4r^2}\le\frac95$ and $\frac95\cdot5^{2r^2}\le8^{9D^2}$.
\end{proof}

\begin{lemma}\label{pbp:cover2}
$\pbvol(\mathrm{Large})\le\beta^{16D^2}\pbvol(\mathrm{Small})$.
\end{lemma}
\begin{proof}
Let $\mathcal J(s)=g_s\oplus h_s\in\mathrm{GL}_D(\C)$; $\|\mathcal J(s)\|=\max(\|g_s\|,1)=\|g_s\|$,
$\|\mathcal J(s)^{-1}\|=\|\mathcal J(s)\|$, and $s\in\mathrm{Small}$ if and only if $\|\mathcal J(s)-1\|\le
\rho_0$. Let $T'=\beta+1$ and choose a maximal family $s_1,\dots,s_\nu$ in Large with $\|\mathcal J(s_i
s_j^{-1})-1\|>\rho_0$ for $i\ne j$. Since $\mathcal J(s_i)-\mathcal J(s_j)=(\mathcal J(s_is_j^{-1})-1)
\mathcal J(s_j)$, $\|\mathcal J(s_i)-\mathcal J(s_j)\|\ge\rho_0/T'$. The balls of radius $\rho_0/(2T')$ about
the $\mathcal J(s_i)$ in $M_D(\C)$ are disjoint and lie in the ball of radius $2T'$, so $\nu\le
(4T'^2/\rho_0)^{2D^2}=(800\beta^5(\beta+1)^2)^{2D^2}\le\beta^{16D^2}$. By maximality, Large is covered
by the sets $\mathrm{Small}\cdot s_i$, and the Haar measure of $R_\tau$ is right invariant.
\end{proof}

\begin{proof}[Proof of Theorem~\ref{pbp:main}]
By the reduction at the end of \S\ref{pbp:ss-data} we may assume $\epsilon_0=1$. Take $\iota$ and $y$
as in Propositions~\ref{pbp:stationary} and \ref{pbp:point}, $x_\tau=y$ and $v_\tau$ as in
Lemma~\ref{pbp:W}; then $\pbel_\tau(x_\tau)\le b^{5n}$, and the integrand is $\theta$ of
\eqref{pbp:eq-theta}. Put $G(s)=(2/\|g_s\|)^{9D^2}$ and $E=\beta^{-50D^2}$, so $E^2=e^{-100T}$.

\emph{Decay.} By Theorem~\ref{pds:decay} with the exponents $9D^2$ and $18D^2$ (both at most $h$,
Lemma~\ref{pbp:sizes}), $|\theta(s)|\le|c(g_s)|\le G(s)$ and $|\theta(s)|\le G(s)^2$. By
Lemma~\ref{pbp:cover1}, $G$ and hence the continuous function $\theta$ are integrable, and $\int G\le
16^{9D^2}\pbvol(\pbeB_y)$.

\emph{Cone with decaying error.} By Proposition~\ref{pbp:cone}, $\omega:=|\theta|-2\Rel\theta\le12E^2$;
also $\omega\le3|\theta|\le3G^2$. Hence $\omega\le6EG$: otherwise $36E^2G^2<\omega^2\le12E^2\cdot3G^2$.

\emph{Small ball.} For $s\in\mathrm{Small}$, $|c(g_s)|\ge\frac12$ by Lemma~\ref{pds:lip}, because
$\rho_0\le(48b^2)^{-1}\le(4C_0)^{-1}$ with $C_0=3r(L+2r)\le12b^2$ (Lemma~\ref{pbp:sizes}); and
$|c'(h_s)|=|(h_s)_{kk}|^{a_\tau}\ge(1-\rho_0)^{a_\tau}\ge1-a_\tau\rho_0\ge\frac12$. So $|\theta|\ge\frac14$
on Small.

\emph{Integration.} With $I=|\int_{R_\tau}\theta\,ds|$, $V=\pbvol(\mathrm{Large})$ and $W=\pbvol(\mathrm{Small})$,
and using $\pbeB_y\subseteq\mathrm{Large}$ ($\|g\|\le1+\|g-1\|\le2$),
\[
\tfrac14W\le\int|\theta|\le2\Rel\int\theta+6E\int G\le2I+6E\,16^{9D^2}\,V .
\]
By Lemma~\ref{pbp:cover2}, $V\le\beta^{16D^2}W\le8\beta^{16D^2}I+24\cdot16^{9D^2}\beta^{-34D^2}V$, and
$24\cdot16^{9D^2}\beta^{-34D^2}\le\frac12$ as $\beta\ge10^{12}$. As $V<\infty$, $V\le16\beta^{16D^2}I\le
\beta^{20D^2}I$. Finally $\pbvol(\pbeB_y)\le V$ and $\beta^{20D^2}=b^{60(r+k)^2}$.
\end{proof}

\subsection{Conclusion: the statement (DS4)}\label{pbp:ss-summary}
\textup{(DS4)} of Theorem~\ref{pbe:ds} holds (Theorem~\ref{pbp:main}), with $x_\tau$ the point of
Proposition~\ref{pbp:point} (composed with $\delta^{-1}$ of Lemma~\ref{pbp:sign} if the component of
$\psi_E$ at infinity is $e^{2\pi i(w+\bar w)}$) and $v_\tau$ the normalised monomial $t_k^{a_\tau}$ of
Lemma~\ref{pbp:W}. Published results used: Schwarz's lemma (Witt's theorem is proved as Lemma~\ref{pbp:tb-witt}; see also \cite[Chapter~7, Theorem~9.1]{Scharlau});
the rest is proved here or in {\cite[\S3.2]{P49}} and~\S\ref{pds:sec}.



\section{Gain, count and size for pairs of terms}\label{B1:g7}
% ===== B-gain.tex =====
% proofs/B-gain.tex
% Statement (G') of proofs/B-pillars.tex (Proposition ppi:G): torsion index times discriminant of
% the common parts. Labels of this file start with pbg:.
\providecommand{\auO}{\mathcal O}
\providecommand{\auNm}{\operatorname{Nm}}
\providecommand{\auA}{\mathbb A}
\providecommand{\aubar}[1]{\overline{#1}}
\providecommand{\pbg}{\mathsf g}
\providecommand{\pbeAp}{A^{+}}
\providecommand{\pbeCL}{\mathrm{CL}}
\providecommand{\ppiLc}{\Lambda^{X}_{c}}
\providecommand{\ppidisc}{\operatorname{disc}}
% macros of this file (prefix \pgn)
\providecommand{\pgnD}{\mathfrak d}
\providecommand{\pgnd}{\operatorname{d}}
\providecommand{\pgnN}{\mathrm N}
\providecommand{\pgnsrc}{\operatorname{src}}
\providecommand{\pgnlm}[2]{\langle #1,#2\rangle_{\Lambda}}
\providecommand{\pgnlf}[2]{\langle\!\langle #1,#2\rangle\!\rangle}

\subsection{The gain for the common parts}\label{pbg:sec}

This section proves statement (G$'$) of Proposition~\ref{ppi:G}, from which that proposition
deduces the gain (G) of Theorem~\ref{ppi:pillars}. We use the notation of \S\S\ref{pbe:sec}
and~\ref{ppi:sec}: the frame $f_0,\dots,f_\varrho$ ($f_a=e_{a+1}$ for $a<\varrho$, $f_\varrho=z$),
$A=\langle f_0,\dots,f_{\varrho-1}\rangle$, $Q=A^\perp/A=X_0\perp Ez$, $Q^+=X_0\perp Ez^+$, the matrix
$Z(\gamma)=(h(f_a,\gamma f_b))_{a,b}$ with leading minors $\delta_l(\gamma)$ ($\delta_0=1$), the
discriminant $\pbg(\gamma)=\delta_n\aubar{\delta_n}-N^2\delta_{n-1}\aubar{\delta_{n-1}}$, the link
$\Lambda_\gamma$ and the isometry $\Upsilon_\gamma$ of Lemma~\ref{pbe:link}, the cell rank
(Definition~\ref{ppi:def-cell}) and the gauge \eqref{ppi:eq-gauge}. Hermitian forms are written here
linear in the second variable, the opposite of convention (Cv1) of \S\ref{pb:ss-setting}; the statements
of this section concern lattices, indices and absolute norms of Gram determinants, which do not change
when $h$ is replaced by $\aubar h$ (\S\ref{pbv:ss-weak}(iv)). For $x\in E$ we write $\auNm x=\prod_{\tau}|\tau(x)|^2=|\pgnN_{E/\Q}(x)|$, and for $x\in
E^+$, $\auNm x=\prod_\tau|\tau(x)|$, as in \S\ref{ppi:sec}. Gram determinants, the constant $c$ of
Lemma~\ref{pbg:jacobi} and the integer $N$ are always normed as elements of $E$, $\auNm x=|\pgnN_{E/\Q}(x)|$;
only the discriminant $\pbg(\gamma)$ is normed as an element of $E^+$.

The lattice of a link is computed exactly through a bordered
Gram determinant (Proposition~\ref{pbg:L1}); passing to the level costs $M^{2j'f}$
(Proposition~\ref{pbg:Lplus}); the link lattice is cut along the intersection with a second link
(Proposition~\ref{pbg:cut}); the two complements are paired integrally after division by $M$
(Proposition~\ref{pbg:congr}); and the common lattice has index dividing a torsion index in the
common part (Proposition~\ref{pbg:index}).

\subsubsection{Lattices and discriminants}\label{pbg:ss-lattices}

Write $\auO=\auO_E$. A \emph{lattice} of rank $\rho$ in an $E$-subspace $V$ of dimension $\rho$ of a
hermitian space is a finitely generated $\auO$-submodule $L\subseteq V$ with $EL=V$. For a hermitian
form $q$ on $V$ (not necessarily the restriction of $h$) let $\pgnD_q(L)$ be the fractional ideal of
$\auO$ generated by the Gram determinants $\det(q(\ell_a,\ell_b))_{a,b\le\rho}$ of the
$\rho$-tuples of elements of $L$, and
\[
\pgnd_q(L)=\pgnN(\pgnD_q(L))^{1/2}\in[0,\infty),
\]
$\pgnN$ the absolute norm of fractional ideals (with $\pgnN(0)=0$). If $q$ is nondegenerate on $V$,
$\pgnd_q(L)>0$. We drop $q$ when it is $h$. For a lattice $L'$ in $V\subseteq X_0$ contained in
$\ppiLc$, $\ppidisc L'$ of \S\ref{ppi:ss-statement} equals $\pgnd(L')$, by Lemma~\ref{pbg:basic}(a)
applied to $L'\subseteq\ppiLc\cap V$. Every lattice has a \emph{pseudo-basis}: an $E$-basis $\ell=
(\ell_1,\dots,\ell_\rho)$ of $V$ and fractional ideals $\mathfrak a_l$ with $L=\bigoplus_l\mathfrak a_l
\ell_l$ (Steinitz). Writing $q(\ell,\ell')$ for the matrix $(q(\ell_a,\ell'_b))_{a,b}$ and $\mathfrak a(L)=
\prod_l\mathfrak a_l$, one has $\pgnD_q(L)=\mathfrak a(L)\aubar{\mathfrak a(L)}\det q(\ell,\ell)$, since
the determinants of the $\rho\times\rho$ matrices with columns in $\bigoplus\mathfrak a_l$ generate
$\mathfrak a(L)$. In particular $\pgnd_q(L)^2=\pgnN(\mathfrak a(L))^2\auNm\det q(\ell,\ell)$.

\begin{lemma}\label{pbg:basic}
Let $q$ be a hermitian form on $V$.
\begin{enumerate}
\item[(a)] If $L'\subseteq L$ are lattices of the same rank, then $\pgnd_q(L')=[L:L']\,\pgnd_q(L)$.
\item[(b)] Let $V=V_1\perp V_2$ for $q$, $L$ a lattice in $V$, $L_1=L\cap V_1$ (a lattice in $V_1$) and
$L_2$ the image of $L$ under the orthogonal projection $\pi_2$ to $V_2$. Then $L_2$ is a lattice in
$V_2$ and $\pgnD_q(L)=\pgnD_q(L_1)\pgnD_q(L_2)$.
\item[(c)] \textup{(Cauchy--Schwarz for Gram determinants.)} Let $q$ be positive definite at every complex place, and let $L$,
$L'$ be lattices of rank $r$ in subspaces of $V$. Let $\mathfrak p$ be the fractional ideal generated
by the determinants $\det q(\ell,\ell')$, $\ell$, $\ell'$ running over the $r$-tuples of $L$ and $L'$.
If $\mathfrak p\subseteq\auO$, $\mathfrak p\ne0$ and $x\le\pgnN\mathfrak p$, then $x^2\le\pgnd_q(L)^2\,
\pgnd_q(L')^2$.
\item[(d)] Let $q$ be positive definite at every complex place, $L'$ a lattice of rank $i$ and $u$ an
isometry of $V$ with $1-u$ injective on $EL'$. Then $\pgnd_q((1-u)L')\le4^{if}\pgnd_q(L')$.
\end{enumerate}
\end{lemma}
\begin{proof}
(a) Localising at the primes of $E^+$ (where $\auO$ is a principal ideal ring), $L'=L C$ for bases,
$\det q(\ell',\ell')=\aubar{\det C}\det C\det q(\ell,\ell)$, and $\pgnN(\det C)$ is the local index
(elementary divisors); multiply over the primes.

(b) $L/L_1\cong L_2$ is torsion free, hence projective; so $L=L_1\oplus L_2'$ with $\pi_2:L_2'\to L_2$
bijective. Take pseudo-bases $(\ell_1,\mathfrak a_1)$ of $L_1$ and $(\ell'_2,\mathfrak a_2)$ of $L'_2$;
then $(\ell_2,\mathfrak a_2)$ with $\ell_2=\pi_2\ell'_2$ is one of $L_2$, and $(\ell_1,\ell'_2)$ with
the ideals $\mathfrak a_1,\mathfrak a_2$ is one of $L$. Now $\ell'_2=\ell_2+\ell_1Y$ for a matrix $Y$
over $E$ (as $\ell'_2-\pi_2\ell'_2\in V_1$), and column operations give $\det q((\ell_1,\ell'_2),(\ell_1,
\ell'_2))=\det q((\ell_1,\ell_2),(\ell_1,\ell_2))=\det q(\ell_1,\ell_1)\det q(\ell_2,\ell_2)$, as
$\ell_1\perp\ell_2$.

(c) Take pseudo-bases $(\ell,\mathfrak a_l)$, $(\ell',\mathfrak a'_l)$; then $\mathfrak p=\aubar{\mathfrak
a(L)}\mathfrak a(L')\det q(\ell,\ell')$. The Gram matrix of $(\ell,\ell')$ is positive semidefinite at
every complex place $\tau$, so $|\tau\det q(\ell,\ell')|^2\le\tau\det q(\ell,\ell)\cdot\tau\det q(\ell',
\ell')$ (for a positive semidefinite $\bigl(\begin{smallmatrix}P&R\\R^*&S\end{smallmatrix}\bigr)$ one has
$|\det R|^2\le\det P\det S$: if $P>0$ this follows from $S-R^*P^{-1}R\ge0$, and in general by
continuity). The product over $\tau$ gives $\pgnN(\mathfrak p)^2\le\pgnd_q(L)^2\pgnd_q(L')^2$.

(d) With a pseudo-basis $(b,\mathfrak a_l)$ of $L'$, $((1-u)b,\mathfrak a_l)$ is one of $(1-u)L'$; put
$b'=ub$. At each complex place $G_\pm=q(b\pm b',b\pm b')$ are positive semidefinite and $G_-+G_+=
2q(b,b)+2q(b',b')=4q(b,b)$. The determinant is monotone on positive semidefinite matrices, so
$\tau\det G_-\le4^i\tau\det q(b,b)$. Take the product of the squares over $\tau$.
\end{proof}

\begin{lemma}[Jacobi's identity with a border]\label{pbg:jacobi}
Let $(\bar Q,h)$ be a unimodular hermitian lattice of rank $m$. Let $F=(F_1,\dots,F_\sigma)$ be
linearly independent vectors of $\bar Q$, let $p=(p_1,p_2)$ be vectors of $\bar Q$, and let $0\ne
c\in\auO$. Let $P^\perp=\{x\in\bar Q:h(F_a,x)=0\ \forall a\}$, of rank $m-\sigma$, and let $\iota=
[F^{\mathrm{sat}}:\textstyle\sum\auO F_a]$, where $F^{\mathrm{sat}}=\bar Q\cap\sum EF_a$. On $\bar
Q\otimes E$ put $\tilde h(x,y)=h(x,y)+c^{-1}\sum_{i=1,2}\aubar{h(p_i,x)}h(p_i,y)$. Then, with
\[
\mathsf G=\begin{pmatrix}h(F,F)&h(F,p)\\h(p,F)&h(p,p)+c\cdot1_2\end{pmatrix},
\qquad
\pgnd_{\tilde h}(P^\perp)^2\,\iota^2\,\auNm(c)^2=\auNm\det\mathsf G .
\]
\end{lemma}
\begin{proof}
Localise. Let $g$ be a basis of $F^{\mathrm{sat}}$, so that $F=gC$ with $\pgnN(\det C)=\iota$;
$F^{\mathrm{sat}}$ is a direct summand of $\bar Q$, so $g$ extends to a basis $R=(g,e)$ of $\bar Q$.
Let $\mathfrak G=h(R,R)$; $\det\mathfrak G$ is a unit. The dual basis $R^\#=R\mathfrak G^{-1}$
(with $h(R,R^\#)=1$) is a basis of $\bar Q^\#=\bar Q$; a vector $x=\sum_kR^\#_kh(R_k,x)$ of $\bar Q$
lies in $P^\perp$ if and only if its coefficients at $g$ vanish, so the last $m-\sigma$ vectors $n$ of
$R^\#$ form a basis of $P^\perp$, and $h(n,n)=(\mathfrak G^{-1})_{22}$ (the block of the indices of
$e$). Write $p=R\pi$; then $h(R^\#,p)=\pi$, so $h(n,p)=\pi_2$ (the rows of $\pi$ at $e$) and
$\tilde h(n,n)=(\mathfrak G^{-1})_{22}+c^{-1}\pi_2\pi_2^*$. Consider
\[
\tilde{\mathfrak G}=\begin{pmatrix}\mathfrak G&\mathfrak G\pi\\\pi^*\mathfrak G&\pi^*\mathfrak
G\pi+c\end{pmatrix}=\begin{pmatrix}1&0\\\pi^*&1\end{pmatrix}\begin{pmatrix}\mathfrak G&0\\0&c
\end{pmatrix}\begin{pmatrix}1&\pi\\0&1\end{pmatrix},\quad
\tilde{\mathfrak G}^{-1}=\begin{pmatrix}\mathfrak G^{-1}+c^{-1}\pi\pi^*&-c^{-1}\pi\\-c^{-1}\pi^*&c^{-1}
\end{pmatrix},
\]
with $\det\tilde{\mathfrak G}=c^2\det\mathfrak G$. Let $S_1$ be the indices of $g$ and of the last two
rows, $S_2$ those of $e$. Jacobi's theorem on complementary minors, $\det(\tilde{\mathfrak
G}^{-1})_{S_2S_2}=\det\tilde{\mathfrak G}_{S_1S_1}/\det\tilde{\mathfrak G}$, gives $\det\tilde
h(n,n)=\det\hat{\mathsf G}/(c^2\det\mathfrak G)$ with $\hat{\mathsf G}=\tilde{\mathfrak G}_{S_1S_1}$,
which is $\mathsf G$ with $g$ in place of $F$. Finally $\det\mathsf G=\aubar{\det C}\det C\det\hat{
\mathsf G}$. Take norms.
\end{proof}

\begin{lemma}\label{pbg:indexmap}
Let $D\subseteq B'$ and $K$, $K'$ be lattices, and $\phi$ a linear map with $\phi(D)\subseteq K$ and
$K'\subseteq\phi(B')+K$. If $[B':D]$ is finite, then $[K':K\cap K']$ divides $[B':D]$.
\end{lemma}
\begin{proof}
$K'/(K\cap K')$ embeds into $(\phi(B')+K)/K$, a quotient of $B'/D$ (as $\phi(D)\subseteq K$); the
order of a subgroup of a quotient of a finite group divides the order of the group.
\end{proof}

\subsubsection{The subframe minor and the gauge}\label{pbg:ss-minor}

Fix classes $c,c'$ and $\gamma\in\Gamma_{c,s,c'}$, and put $\Lambda=h_cL$ (this differs from the $\Lambda$ of \S\ref{ppi:sec}). As $h_c\in H(\auA_f)$ fixes
the $e_a$, $e_a^*$ and $z$, and $L=\bigoplus_a(\auO e_a\oplus\auO e_a^*)\perp L_{\mathrm{mid}}$ with
$L_{\mathrm{mid}}$ unimodular and $z\in L_{\mathrm{mid}}$ (proof of Lemma~\ref{pb:KR}), the lattice
$\Lambda$ is unimodular, $\Lambda\cap A=\bigoplus_a\auO e_a$, a vector of $\Lambda$ orthogonal to $A$
is $a+m$ with $a\in\Lambda\cap A$, $m\in\Lambda\cap(X_0\oplus Ez)$, and $z\in\Lambda$. As $k_\gamma=
h_c^{-1}\gamma h_{c'}\in K$, $\gamma h_{c'}L=\Lambda$; in particular $\gamma f_b\in\Lambda$ for all $b$.

Let $j=j(\gamma)$ be the cell rank and $j'=\min(\varrho,j)$, and assume $\delta_j(\gamma)\ne0$ and
$j\le\varrho\iff\pbg(\gamma)=0$ (by Proposition~\ref{ppi:cell} this holds for $\gamma\in\Gamma^*$). Put
\[
\mu(\gamma)=\auNm\delta_j(\gamma)\,N^{-2(j-\varrho)_+f},
\]
so that \eqref{ppi:eq-gauge} reads $\mathsf G(\gamma)=M^{2j'f}\max(\auNm\pbg(\gamma),N^f\mu(\gamma))/(
i_\gamma i'_\gamma)$. Let $S_j=\{0,\dots,j'-1\}\cup\{\varrho\}$ and let $\delta_*(\gamma)=\det Z(\gamma)_{
S_jS_j}$ be the \emph{subframe minor}.

\begin{lemma}[Shared vectors]\label{pbg:shared}
Let $I=A\cap\gamma A$.
\begin{enumerate}
\item[(a)] If $j\le\varrho$, there is a common vector $f\hat x=\gamma f\hat y$ with $\hat y_\varrho=1$
and $\hat y_a=0$ for $j\le a<\varrho$; it satisfies $\hat x_\varrho\aubar{\hat x_\varrho}=1$ and
$(Z(\gamma)\hat y)_b=0$ for $b<\varrho$.
\item[(b)] $I\cap\langle f_0,\dots,f_{j'-1}\rangle=0$, and $I$ has a basis $u_b=f_b+\sum_{a<j'}c_{ab}
f_a$, $j'\le b<\varrho$. So $\dim I=\varrho-j'$.
\item[(c)] A vector $a\in A$ with $h(\gamma f_b,a)=0$ for $0\le b\le\varrho$ lies in $I$.
\end{enumerate}
\end{lemma}
\begin{proof}
(a) By the definition of $j$ there is $\zeta$,
$\zeta\aubar\zeta=1$, such that for each $a$ with $j\le a\le\varrho$ there is a common vector $fx^{(a)}=
\gamma fy^{(a)}$ with top index $a$ and $(x^{(a)}_a,y^{(a)}_a)=(\zeta,1)$. Starting from $a=\varrho$ and
subtracting suitable multiples of the vectors with $a=\varrho-1,\dots,j$, we get a common vector
$f\hat x=\gamma f\hat y$ with $\hat y_\varrho=1$ and $\hat y_a=0$ for $j\le a<\varrho$. As $A$ is
isotropic and orthogonal to $z$, $-N\hat x_\varrho\aubar{\hat x_\varrho}=h(f\hat x,f\hat x)=h(f\hat y,f
\hat y)=-N$; and for $b<\varrho$, $(Z\hat y)_b=h(f_b,\gamma f\hat y)=h(f_b,f\hat x)=0$.

(b) Let $fx=\gamma fy\in I$ with $x$ supported on $\{0,\dots,j'-1\}$; then $y_\varrho=0$. For every
$b$, $h(\gamma f_b,fx)=h(f_b,fy)=0$ ($fy\in A$). For $b<j'$ this says that the transpose conjugate of
the leading block of order $j'$ of $Z$ kills $x$. If $j\le\varrho$ this block has determinant
$\delta_j\ne0$; if $j=\varrho+1$ use all $b\le\varrho$ and $\delta_n\ne0$ (with $x$ extended by $0$).
So $x=0$. For $j'\le a<\varrho$ (so $j\le\varrho$) the common vector $fx^{(a)}=\gamma fy^{(a)}$ of (a)
has top index $a<\varrho$, so it lies in $A\cap\gamma A=I$ and has coefficient $\zeta\ne0$ at $f_a$.
Triangular elimination gives the $u_b$; by the first claim they span $I$.

(c) Subtract from $a$ the combination of the $u_b$ with the same coefficients at $f_b$, $b\ge j'$;
the difference $a''$ lies in $\langle f_0,\dots,f_{j'-1}\rangle$ and is still orthogonal to all
$\gamma f_b$ (as $I\subseteq\gamma A$ is orthogonal to $\gamma A+E\gamma z$). The argument of (b) gives
$a''=0$.
\end{proof}

\begin{lemma}[The subframe minor]\label{pbg:subframe}
$\auNm\delta_*(\gamma)=N^{2f}\mu(\gamma)$.
\end{lemma}
\begin{proof}
If $j=\varrho+1$, $S_j=\{0,\dots,\varrho\}$ and $\delta_*=\delta_n=\delta_j$, $\mu=\auNm\delta_n/N^{2f}$.
Let $j\le\varrho$, and let $P$ be the matrix whose rows $a<j$ are those of $Z=Z(\gamma)$ and whose other
rows are those of the identity; $\det P=\delta_j\ne0$. By Lemma~\ref{pbg:shared}(a), $P\hat y=
\epsilon_\varrho$ (the last unit vector). Expanding along the identity rows, $\delta_*$ is the
determinant of $P$ with its last row replaced by the last row $Z_\varrho$ of $Z$; by Cramer's rule
this is $(Z_\varrho\hat y)\det P$. And $Z_\varrho\hat y=h(z,\gamma f\hat y)=h(z,f\hat x)=-N\hat x_\varrho$.
So $\delta_*=-N\hat x_\varrho\delta_j$ and, as $|\tau(\hat x_\varrho)|=1$ at every $\tau$, $\auNm\delta_*
=N^{2f}\auNm\delta_j$.
\end{proof}

\begin{remark}[The gauge]\label{pbg:rem-gauge}
(i) The gauge may also be written as $M^{2j'f}\max(\mathsf g^\sharp,\auNm\pbg)$ with $\mathsf g^\sharp=
\auNm\delta_*/(i_\gamma i'_\gamma N^f)$. By Lemma~\ref{pbg:subframe}, $\mathsf g^\sharp=N^f\mu/(i_\gamma
i'_\gamma)$; and if $\pbg(\gamma)\ne0$ then $j=\varrho+1$ and $i_\gamma=i'_\gamma=1$ (see (ii)), so the
expression with $\mathsf g^\sharp$ is \eqref{ppi:eq-gauge}. (ii) The indices are, exactly: $i_\gamma=[R_A:R_{<j'}]$, where
$R_A$ (resp.\ $R_{<j'}$) is the $\auO$-span in $E^n$ of the rows $a<\varrho$ (resp.\ $a<j'$) of
$Z(\gamma)$, and $i'_\gamma$ is the same for $Z(\gamma^{-1})=Z(\gamma)^*$. (In lattice coordinates the
rows are the products of $e_a$ with the frame $k_\gamma f$, as $h(e_a,k_\gamma f_b)=h(e_a,\gamma f_b)$.)
They are finite: by Lemma~\ref{pbg:shared}(b),(c), the kernel of $a\mapsto(h(\gamma f_b,a))_b$ on $A$ is
$I$, of codimension $j'$, so rows $a<j'$ span the row space over $E$. For $j=\varrho+1$, $j'=\varrho$
and $i_\gamma=i'_\gamma=1$.
\end{remark}

\subsubsection{The lattice of a link}\label{pbg:ss-link}

For a pair $P=(y_0+\lambda'z,\ y_0'+\lambda z)\in Q^2$ ($y_0,y_0'\in X_0$) put $\pgnsrc P=y_0+\lambda z^+
\in Q^+$. On $Q^2$ we use two forms: the \emph{link metric}
$\pgnlm PR=\tfrac12\bigl(h^+(\pgnsrc P,\pgnsrc R)+h^+(\pgnsrc\bar P,\pgnsrc\bar R)\bigr)$, where $\bar
P=(y_0'+\lambda z,y_0+\lambda'z)$ is the swapped pair, which is positive definite at every complex
place; and the \emph{link form} $\pgnlf{(p,p')}{(r,r')}=h(p,r)-h(p',r')$. By Lemma~\ref{pbe:link}(b),
for $P,R\in\Lambda_\gamma$ one has $\pgnsrc\bar P=\Upsilon_\gamma\pgnsrc P$, so $\pgnlm PR=h^+(\pgnsrc
P,\pgnsrc R)$, and $\pgnsrc:\Lambda_\gamma\to Q^+$ is an isometric isomorphism.

\begin{lemma}\label{pbg:linkforms}
Let $\gamma,\gamma'\in\Gamma$.
\begin{enumerate}
\item[(a)] For $t,t'\in(A+\gamma A)^\perp$ with images $P,R\in\Lambda_\gamma$, $\pgnlf PR=0$, and
$\Lambda_\gamma$ is its own orthogonal for $\pgnlf\cdot\cdot$ in $Q^2$.
\item[(b)] For $t\in(A+\gamma A)^\perp$, $t'\in(A+\gamma'A)^\perp$ with images $P\in\Lambda_\gamma$,
$R\in\Lambda_{\gamma'}$, $\pgnlf PR=h(t,t')-h(t,\gamma\gamma'^{-1}t')$.
\item[(c)] For $P\in\Lambda_\gamma$, $R\in\Lambda_{\gamma'}$, $\pgnlf PR=h^+(a,b-ub)$, where $a=\pgnsrc P$,
$b=\pgnsrc R$ and $u=\Upsilon_\gamma^{-1}\Upsilon_{\gamma'}$.
\item[(d)] For $t\in(A+\gamma A)^\perp$ with image $P$, $\pgnlm PP=h(t,t)+N^{-1}(|h(z,t)|^2+|h(\gamma z,
t)|^2)$; and the two linear forms $\zeta_1(p,p')=h(z,p)$, $\zeta_2(p,p')=h(z,p')$ take at $P$ the values
$h(z,t)$, $h(\gamma z,t)$. On $\Lambda_\gamma$, $\zeta_1=\pgnlm{\varepsilon_1}\cdot$ and $\zeta_2=
\pgnlm{\varepsilon_2}\cdot$, where $\varepsilon_1,\varepsilon_2\in\Lambda_\gamma$ have sources
$-\Upsilon_\gamma^{-1}z^+$ and $-z^+$.
\end{enumerate}
\end{lemma}
\begin{proof}
Classes in $Q=A^\perp/A$ carry the induced form. (b) $h([t],[t'])=h(t,t')$ as $t,t'\perp A$, and
$h([\gamma^{-1}t],[\gamma'^{-1}t'])=h(\gamma^{-1}t,\gamma'^{-1}t')=h(t,\gamma\gamma'^{-1}t')$ as
$\gamma^{-1}t,\gamma'^{-1}t'\perp A$. (a) is (b) with $\gamma'=\gamma$; as $\dim\Lambda_\gamma=k+1$ is
half of $\dim Q^2$ and $\pgnlf\cdot\cdot$ is nondegenerate, $\Lambda_\gamma$ is its own orthogonal. (c) Write
$P=(y_0+\lambda'z,y_0'+\lambda z)$, $R=(w_0+\mu'z,w_0'+\mu z)$. Then $h(p,r)-h(p',r')=(h(y_0,w_0)-
N\aubar{\lambda'}\mu')-(h(y_0',w_0')-N\aubar\lambda\mu)=h^+(a,b)-h^+(\Upsilon_\gamma a,\Upsilon_{\gamma'}b)$,
and $h^+(\Upsilon_\gamma a,\Upsilon_{\gamma'}b)=h^+(a,ub)$. (d) With $t=\alpha+y_0+\lambda'z$, $\alpha\in A$,
and $\gamma^{-1}t=\alpha'+y_0'+\lambda z$: $h(t,t)=|y_0|^2-N|\lambda'|^2$, $h(z,t)=-N\lambda'$, $h(\gamma
z,t)=h(z,\gamma^{-1}t)=-N\lambda$, and $\pgnlm PP=|y_0|^2+N|\lambda|^2$. For the last claim,
$\pgnlm{\varepsilon_1}P=-h^+(\Upsilon_\gamma^{-1}z^+,a)=-h^+(z^+,\Upsilon_\gamma a)=-N\lambda'=h(z,p)$, and
$\pgnlm{\varepsilon_2}P=-h^+(z^+,a)=-N\lambda=h(z,p')$.
\end{proof}

Let $\gamma\in\Gamma_{c,s,c'}$ be as in \S\ref{pbg:ss-minor}. Define the lattices of $\Lambda_\gamma$
(for the link metric)
\[
L^1_\gamma=\{([t],[\gamma^{-1}t]):t\in\Lambda\cap(A+\gamma A)^\perp\},\qquad
L^+_\gamma=\{([t],[\gamma^{-1}t]):t\in\Lambda\cap(A+\gamma A)^\perp,\ t-m\in M(\Lambda\cap A)\},
\]
where in the second $t=\alpha+m$ with $\alpha\in A$, $m\in X_0\oplus Ez$. They are lattices of rank $k+1$
($t\mapsto([t],[\gamma^{-1}t])$ maps $(A+\gamma A)^\perp$ onto $\Lambda_\gamma$ with kernel $I$, and
$L^+_\gamma\supseteq ML^1_\gamma$). We write $\pgnd$ for $\pgnd_{\pgnlm\cdot\cdot}$ on $\Lambda_\gamma$.
Let $t(\gamma)=[\Lambda\cap(A+\gamma A):\Lambda\cap A+\gamma(\Lambda\cap A)]$, the torsion index of
Proposition~\ref{ppi:G} written in the coordinates of $\Lambda$.

\begin{proposition}[The lattice $L^1_\gamma$]\label{pbg:L1}
$\pgnd(L^1_\gamma)\,i_\gamma\,i'_\gamma\,t(\gamma)\le\mu(\gamma)$.
\end{proposition}
\begin{proof}
\emph{The quotient.} $I=A\cap\gamma A$ is isotropic and $\Lambda\cap I$ is primitive in $\Lambda$. Let
$\bar Q=(\Lambda\cap I^\perp)/(\Lambda\cap I)$, a lattice in $I^\perp/I$ with the induced form. It is
unimodular: a linear form $\phi$ on $\bar Q$ lifts to $\Lambda\cap I^\perp$, and extends to $\Lambda$
(as $\Lambda\cap I^\perp$ is primitive); since $\Lambda$ is unimodular, $\phi=h(w,\cdot)$ with $w\in
\Lambda$, and $w\perp\Lambda\cap I$, so $w\in\Lambda\cap I^\perp$ (as $\Lambda\cap I$ spans $I$) and
$\phi=h([w],\cdot)$; and $\bar Q$ is integral. With $n_{\mathrm{tot}}=\dim Y_0$, its rank is $\dim I^\perp/I=n_{\mathrm{tot}}-2(\varrho-j')=2j'+k+1$.

\emph{The two lists.} The vectors $f_a,\gamma f_a$ ($a<\varrho$) lie in $\Lambda\cap I^\perp$ (as $A$,
$\gamma A$ are isotropic and contain $I$), and so do $z$ and $\gamma z$ ($z\perp A$, $\gamma z\perp\gamma
A$). Let $F$ be the classes in $\bar Q$ of $f_a,\gamma f_a$ for $a<j'$, and $p=([z],[\gamma z])$. By
Lemma~\ref{pbg:shared}(b), the class of $f_b$ for $b\ge j'$ is that of $-\sum_{a<j'}c_{ab}f_a$, and
likewise for $\gamma f_b$ (apply (b) to $\gamma^{-1}$, which has the same cell rank, the same $\pbg$
and $\delta_l(\gamma^{-1})=\aubar{\delta_l(\gamma)}$, since $Z(\gamma^{-1})=Z(\gamma)^*$, and use $A\cap
\gamma^{-1}A=\gamma^{-1}I$); hence the span $W$ of $F$ is the image of
$A+\gamma A$, which has dimension $(\varrho+j')-(\varrho-j')=2j'$, and $F$ is linearly independent.

\emph{The complement.} A vector $t\in\Lambda\cap I^\perp$ is orthogonal to $A+\gamma A$ if and only if
$[t]\perp W$. So $t\mapsto[t]$ maps $\Lambda\cap(A+\gamma A)^\perp$ onto $P^\perp=\{x\in\bar Q:x\perp
F\}$ (onto, because $\Lambda\cap I^\perp\to\bar Q$ is onto), with kernel $\Lambda\cap I$; composing
with the isomorphism induced by $t\mapsto([t],[\gamma^{-1}t])$, whose kernel is also $I$, we get an
isomorphism of $P^\perp$ onto $L^1_\gamma$, which by Lemma~\ref{pbg:linkforms}(d) is isometric from
$\tilde h=h+N^{-1}(|h(p_1,\cdot)|^2+|h(p_2,\cdot)|^2)$ to the link metric. Lemma~\ref{pbg:jacobi} with
$c=N$ gives
\[
\pgnd(L^1_\gamma)^2\,\iota^2\,N^{4f}=\auNm\det\mathsf G,\qquad\iota=[F^{\mathrm{sat}}:\textstyle\sum
\auO F_a].
\]

\emph{The bordered determinant.} $\mathsf G$ is the Gram matrix of the two lists $(f_a\ (a<j'),\ z+w_1)$
and $(\gamma f_a\ (a<j'),\ \gamma z+w_2)$ in $I^\perp/I\perp Ew_1\perp Ew_2$, with $h(w_i,w_i)=N$. Each
list is isotropic ($h(z+w_1,z+w_1)=-N+N=0$), and the products of the first with the second form
the matrix $Z(\gamma)_{S_jS_j}$. So $\det\mathsf G=\pm\delta_*\aubar{\delta_*}$ and, by
Lemma~\ref{pbg:subframe}, $\auNm\det\mathsf G=(N^{2f}\mu)^2$. Hence $\pgnd(L^1_\gamma)\iota=\mu$.

\emph{The index.} It remains to show that $i_\gamma i'_\gamma t(\gamma)$ divides $\iota$. Let $L_1$ be
the $\auO$-span of $F$, $L_2$ that of the classes of $f_a$ ($a<\varrho$) and $\gamma f_a$ ($a<j'$), $L_3$
that of the classes of all $f_a,\gamma f_a$ ($a<\varrho$), and $L_4=\bar Q\cap W=F^{\mathrm{sat}}$; then
$L_1\subseteq L_2\subseteq L_3\subseteq L_4$ and $\iota=[L_4:L_1]$. (i) The map $\phi(x)=(h(\gamma f_b,
x))_{b\le\varrho}$ is defined on $\bar Q$ (as $\gamma f_b\perp I$); it sends the class of $f_a$ to the
conjugate of row $a$ of $Z$, and that of $\gamma f_a$ to $(h(f_b,f_a))_b=0$. By
Lemma~\ref{pbg:indexmap} (with $B'=L_2$, $D=L_1$, $K=\phi(L_1)$, $K'=\phi(L_2)$), $i_\gamma$ divides
$[L_2:L_1]$. (ii) In the same way, with $\phi'(x)=(h(f_b,x))_b$, which kills the classes of the
$f_a$ and sends that of $\gamma f_a$ to the conjugate of row $a$ of $Z(\gamma^{-1})$, $i'_\gamma$
divides $[L_3:L_2]$. (iii) Choose a section $\beta:\bar Q\to\Lambda\cap I^\perp$, $\auO$-linear, of the
projection. For $w\in\Lambda\cap(A+\gamma A)$, $w-\beta([w])\in\Lambda\cap I\subseteq\Lambda\cap A$; so
$\beta$ maps $L_3$ into $T_0=\Lambda\cap A+\gamma(\Lambda\cap A)$ and $L_4$ into $\Lambda\cap(A+\gamma A)$,
and $\Lambda\cap(A+\gamma A)\subseteq\beta(L_4)+T_0$. Lemma~\ref{pbg:indexmap} gives $t(\gamma)\mid[L_4:
L_3]$. Multiply.
\end{proof}

\begin{proposition}[The lattice $L^+_\gamma$]\label{pbg:Lplus}
$t(\gamma)\,\pgnd(L^+_\gamma)\le M^{2j'f}\mu(\gamma)/(i_\gamma i'_\gamma)$.
\end{proposition}
\begin{proof}
Let $\Phi(\alpha)=(h(\gamma f_b,\alpha))_{b\le\varrho}$ for $\alpha\in A$, and $R=\Phi(\Lambda\cap A)$, a
lattice of rank $\varrho-\dim I=j'$ (Lemma~\ref{pbg:shared}(c)). For $P=([t],[\gamma^{-1}t])\in
L^1_\gamma$ with $t=\alpha+m$, the class of $\Phi(\alpha)$ in $R/MR$ depends only on $P$ ($t$ is
determined modulo $\Lambda\cap I$, and $\Phi(I)=0$). If it vanishes, $\Phi(\alpha)=M\Phi(\omega)$ with
$\omega\in\Lambda\cap A$, so $\alpha-M\omega\in\Lambda\cap I$ by Lemma~\ref{pbg:shared}(c), and $t'=t-(
\alpha-M\omega)$ lies in $\Lambda\cap(A+\gamma A)^\perp$, has the same image $P$ and $A$-part $M\omega$:
$P\in L^+_\gamma$. So $L^+_\gamma$ contains the kernel of $L^1_\gamma\to R/MR$, and $[L^1_\gamma:
L^+_\gamma]\le\#R/MR=\pgnN(M\auO)^{j'}=M^{2j'f}$. By Lemma~\ref{pbg:basic}(a) and
Proposition~\ref{pbg:L1}, $t(\gamma)\pgnd(L^+_\gamma)\le M^{2j'f}t(\gamma)\pgnd(L^1_\gamma)\le M^{2j'f}
\mu/(i_\gamma i'_\gamma)$.
\end{proof}

\subsubsection{The cut along a second link}\label{pbg:ss-cut}

\begin{lemma}[The vectors $z^+$ and $\Upsilon_\gamma^{-1}z^+$]\label{pbg:linkvec}
Let $\varkappa_\gamma$ be the $z^+$-coordinate of $\Upsilon_\gamma^{-1}z^+$, and $\mathsf M(\gamma)=\max(
\auNm\pbg(\gamma),N^f\mu(\gamma))$. The Gram matrix of $(z^+,\Upsilon_\gamma^{-1}z^+)$ for $h^+$ is
$\bigl(\begin{smallmatrix}N&N\varkappa_\gamma\\N\aubar{\varkappa_\gamma}&N\end{smallmatrix}\bigr)$. If $j=\varrho+1$,
then $\pbg(\gamma)=\delta_n\aubar{\delta_n}(1-\varkappa_\gamma\aubar{\varkappa_\gamma})$. Consequently, for every
principal minor $\mathsf m$ of order $r\le1+(j-\varrho)_+$ of this Gram matrix, $\auNm\mathsf m\le(\mathsf
M(\gamma)/\mu(\gamma))^2$.
\end{lemma}
\begin{proof}
The diagonal entries are $h^+(z^+,z^+)=N$ and $h^+(\Upsilon^{-1}z^+,\Upsilon^{-1}z^+)=N$; the
off-diagonal one is $h^+(z^+,\Upsilon^{-1}z^+)=N\varkappa_\gamma$. Let $j=\varrho+1$, so $\delta=\delta_n\ne0$;
let $y=\operatorname{adj}(Z)\epsilon_\varrho$, so that $Zy=\delta\epsilon_\varrho$ and $y_\varrho=
\delta_{n-1}$. Put $t=\gamma fy$. For $b<\varrho$, $h(f_b,t)=(Zy)_b=0$ and $h(\gamma f_b,t)=h(f_b,fy)=0$; so
$t\perp A+\gamma A$. Now $[\gamma^{-1}t]=[fy]=\delta_{n-1}z$, and $h(z,t)=(Zy)_\varrho=\delta$, so the
$z$-coefficient of $[t]$ is $-\delta/N$. By Lemma~\ref{pbe:link}(b), $\Upsilon_\gamma(y_0+\delta_{n-1}z^+)=
-(\delta/N)z^+$, where $y_0$ is the $X_0$-part of $t$; so $\Upsilon_\gamma^{-1}z^+=-(N/\delta)(y_0+
\delta_{n-1}z^+)$ and $\varkappa_\gamma=-N\delta_{n-1}/\delta$. Hence $1-\varkappa_\gamma\aubar{\varkappa_\gamma}=(\delta
\aubar\delta-N^2\delta_{n-1}\aubar{\delta_{n-1}})/\delta\aubar\delta=\pbg/\delta\aubar\delta$.

For the minors: $r=0$ gives $1\le(N^f\mu/\mu)^2$; $r=1$ gives $\auNm N=N^{2f}=(N^f\mu/\mu)^2$; $r=2$
occurs only if $j=\varrho+1$, and then $\mathsf m=N^2(1-\varkappa_\gamma\aubar{\varkappa_\gamma})=N^2\pbg/\delta_n
\aubar{\delta_n}$, so $\auNm\mathsf m=N^{4f}(\auNm\pbg)^2/(\auNm\delta_n)^2=(\auNm\pbg/\mu)^2$, as $\mu=
\auNm\delta_n/N^{2f}$ and $\pbg\in E^+$.
\end{proof}

The factor $\mathsf{mn}(s)$ rests on the following residual form of Lemma~\ref{ppi:climb}. Recall
(\S\ref{pb:sec}) that $s\in\auG$ meets at a place $\mathfrak P\in\{w,\aubar w\}$ above a chosen $w$ if
the spans of the frame and of its image under $s$ modulo $\mathfrak P$ meet nontrivially, and
$\mathsf{mn}(s)=\prod\auNm\mathfrak P$ over these places.

\begin{lemma}[Climbing modulo a level place]\label{pbg:climbp}
Let $s$ agree, let $\kappa\in K$ reduce to $s$, and let $s$ meet at $\mathfrak P$. Then there are $x,y\in
\auO^n$ with $fx-\kappa fy\in\mathfrak PL$ and $(x_\varrho,y_\varrho)\not\equiv(0,0)\bmod\mathfrak P$.
\end{lemma}

\begin{proof}
Write $\mathfrak P$ as a pair $(w,\epsilon)$: $\mathfrak P=w$ or $\mathfrak P=\aubar w$. An element of
$K$ is determined modulo $J$ by its reductions modulo the places $(w',\epsilon)$, $w'$ chosen, with the
\emph{same} $\epsilon$ (the reduction modulo $\aubar{w'}$ is the inverse adjoint of that modulo $w'$);
so for $u,u'\in K$, $[u]\,s=s\,[u']$ in $\auG$ (where $[u]$ denotes the image of $u$ in $\auG$) if and only if $u\kappa\equiv \kappa u'$ modulo every place
$(w',\epsilon)$. Write $\tilde\cdot$ for reduction modulo $\mathfrak P$, $\mathbb F=\auO/\mathfrak P$, and
$\tilde\kappa$ for the matrix of $\kappa$ modulo $\mathfrak P$.

\emph{Lifts.} By the Chinese remainder theorem, every $\theta\in\mathbb F$ has a lift $\hat\theta\in
\widehat\auO_E$ with $\hat\theta\equiv0$ modulo every place above $M$ other than $\mathfrak P$ (in
particular modulo $\aubar{\mathfrak P}$), and $\hat0=0$. Then $\aubar{\hat\theta}\equiv0$ modulo every
place above $M$ other than $\aubar{\mathfrak P}$, in particular modulo $\mathfrak P$ and modulo all
$(w',\epsilon)$.

\emph{One step.} Let $\tilde f x=\tilde\kappa\tilde f y$ with $x,y\in\mathbb F^n$ vanishing above $a<\varrho$,
$(x_a,y_a)=(\xi,\xi')$. As in the proof of Lemma~\ref{ppi:climb}, $(\xi,\xi')$ occurs at $a+1$ modulo
$\mathfrak P$ as soon as every $\mathbb F$-linear form $\lambda$ on $L/\mathfrak PL$ that vanishes on
$\tilde f_c,\tilde\kappa\tilde f_c$ ($c\le a$) satisfies $\xi\lambda(\tilde f_{a+1})=\xi'\lambda(\tilde\kappa\tilde f_{a+1})$.
The frame vectors $f_c$ belong to a basis $\beta$ of $L$. Define $\widehat\auO$-linear forms $m,m'$ on
$L\otimes\widehat\auO$ by $m(\beta_l)=\widehat{\lambda(\tilde\beta_l)}$ and $m'(\beta_l)=\widehat{\lambda(\tilde
\kappa\tilde\beta_l)}$ (lifts as above). Then $m\equiv\lambda$, $m'\equiv\lambda\circ\tilde\kappa$ modulo $\mathfrak P$,
both vanish modulo every place above $M$ other than $\mathfrak P$, and $m(f_c)=m'(f_c)=0$ exactly for
$c\le a$ (as $\hat0=0$). As $L$ is unimodular, $m=h(q,\cdot)$ and $m'=h(q',\cdot)$ with $q,q'\in L\otimes
\widehat\auO$, orthogonal to $f_0,\dots,f_a$. Let $\hat x,\hat y$ be the coordinatewise lifts of $x,y$,
$p=f\hat x$, $p'=f\hat y\in\langle f_0,\dots,f_a\rangle$, fix $\omega$ as in Lemma~\ref{ppi:eichler} and
$r\in\widehat\auO_E$, and let
\[
u=\mathsf E_{rp,q},\qquad u'=\mathsf E_{rp',q'}
\]
be the Eichler elements of Lemma~\ref{ppi:eichler} (with $j=a<\varrho$; the lemma is pure algebra and
holds over $\auA_f$). They are integral, so $u,u'\in U(\auA_f)\cap K$, and $\sigma(u)=r\hat x_am(f_{a+1})$,
$\sigma(u')=r\hat y_am'(f_{a+1})$. Now
\[
u(t)=t+h(q,t)\,rp-\aubar r\,h(p,t)\,(q+\varsigma rp),\qquad h(p,t)=\textstyle\sum_b\aubar{\hat x_b}\,h(f_b,t).
\]
\emph{This is where both $w$ and $\aubar w$ enter:} the last term carries the conjugates
$\aubar{\hat x_b}$, which vanish modulo $\mathfrak P$ because $\hat x_b\equiv0$ modulo $\aubar{\mathfrak
P}$, whereas the second term carries $\hat x_b$ and $m$. Hence modulo $\mathfrak P$, $\tilde u=1+\tilde r
\,(\tilde fx)\,\lambda$, and modulo every place $(w',\epsilon)\ne\mathfrak P$ both terms vanish ($\hat x$,
$\aubar{\hat x}$ and $m$ vanish there), so $u\equiv1$. Likewise $\tilde u'=1+\tilde r(\tilde fy)(\lambda\circ
\tilde\kappa)$ and $u'\equiv1$ at the other places $(w',\epsilon)$. Modulo $\mathfrak P$,
\[
\tilde u\tilde\kappa=\tilde\kappa+\tilde r(\tilde fx)(\lambda\circ\tilde\kappa)=\tilde\kappa+\tilde r(\tilde\kappa\tilde fy)(\lambda\circ\tilde\kappa)=
\tilde\kappa\tilde u' ,
\]
and $u\kappa\equiv \kappa\equiv \kappa u'$ at the other places. So $[u]\,s=s\,[u']$ in $\auG$ and, $s$ agreeing, $\psi_f(u)=
\psi_f(u')$, that is $\psi_{E,f}(r(\Theta-\Theta')/(M\varpi))=1$ with $\Theta=\hat x_am(f_{a+1})$,
$\Theta'=\hat y_am'(f_{a+1})$, for every $r\in\widehat\auO_E$. Since $\varpi$ generates the different of
$E$, the $\alpha\in\auA_{E,f}$ with $\psi_{E,f}(r\alpha)=1$ for all $r\in\widehat\auO_E$ form $\varpi^{-1}
\widehat\auO_E$ (the dual of $\widehat\auO_E$ for $\psi_{E,f}$ is the inverse different, (Cv2) of \S\ref{pb:ss-setting}). Hence $\Theta-\Theta'\in
M\widehat\auO_E$, and modulo $\mathfrak P$: $\xi\lambda(\tilde f_{a+1})=\xi'\lambda(\tilde\kappa\tilde f_{a+1})$.

\emph{Climbing.} As $s$ meets at $\mathfrak P$, there are $x,y\in\mathbb F^n$, $x\ne0$, with $\tilde fx=
\tilde\kappa\tilde fy$. Let $a$ be the top index (the largest index with $(x_a,y_a)\ne(0,0)$; it exists as
$x\ne0$). If $a<\varrho$, by the step applied $\varrho-a$ times, $(x_a,y_a)$ occurs at $a+1,\dots,\varrho$. Lift the resulting $x,y$ to $\auO^n$.
\end{proof}

\begin{lemma}[Meeting places]\label{pbg:meet}
Let $s$ agree, $\gamma\in\Gamma_{c,s,c'}$, and $t_1,t_2\in\Lambda\cap(A+
\gamma A)^\perp$. Then $\Delta=\det\bigl(\begin{smallmatrix}h(z,t_1)&h(z,t_2)\\h(\gamma z,t_1)&h(\gamma
z,t_2)\end{smallmatrix}\bigr)\in\auO$ lies in $\aubar{\mathfrak P}$ for every place $\mathfrak P$ at
which $s$ meets; so every ideal generated by such determinants has norm at least $\mathsf{mn}(s)$ if it
is not $0$.
\end{lemma}
\begin{proof}
Put $\kappa=k_\gamma$ and $\tilde t_l=h_c^{-1}t_l\in L$. As $h_c$, $h_{c'}$ fix the frame and $h_c$ is an
isometry, $h(z,t_l)=h(z,\tilde t_l)$, $h(\gamma z,t_l)=h(\kappa z,\tilde t_l)$, and $\tilde t_l\perp A+\kappa A$. Take
$x,y$ as in Lemma~\ref{pbg:climbp}. Since $\tilde t_l\perp A$ and $\tilde t_l\perp \kappa A$, $h(fx,\tilde t_l)=
\aubar{x_\varrho}h(z,\tilde t_l)$ and $h(\kappa fy,\tilde t_l)=\aubar{y_\varrho}h(\kappa z,\tilde t_l)$, and $h(fx-\kappa fy,
\tilde t_l)\in h(\mathfrak PL,L)\subseteq\aubar{\mathfrak P}$. So $(\aubar{x_\varrho},-\aubar{y_\varrho})$, which
is nonzero modulo $\aubar{\mathfrak P}$, is in the left kernel of the matrix of $\Delta$ modulo
$\aubar{\mathfrak P}$, and $\Delta\in\aubar{\mathfrak P}$. The entries are integral as $z,\gamma z,t_l
\in\Lambda$. The places $\aubar{\mathfrak P}$ are distinct primes with the norms $\auNm\mathfrak P$.
\end{proof}

\begin{proposition}[The cut]\label{pbg:cut}
Let $s$ agree, $\gamma\in\Gamma_{c,s,c'}\cap\Gamma^*$ and $\gamma'\in\Gamma$. Let $\Lambda_\cap=
\Lambda_\gamma\cap\Lambda_{\gamma'}$, $i=k+1-\dim\Lambda_\cap$, $V_y$ the orthogonal complement of
$\Lambda_\cap$ in $\Lambda_\gamma$ for the link metric, $\pi_y$ the orthogonal projection onto $V_y$, and
$L_y=\pi_y(L^+_\gamma)$, a lattice of rank $i$ in $V_y$. Let $C_\gamma=L_\gamma\cap\pbeCL(\gamma,\gamma')$.
Let $x$ be an integer with $1\le x\le\mathsf{mn}(s)$, and $x=1$ unless $i=0$ and $\pbg(\gamma)\ne0$. Then $C_\gamma$ is a lattice of rank $\dim\pbeCL(\gamma,\gamma')$
and
\[
t(\gamma)\,\pgnd(C_\gamma)\,x\,\pgnd(L_y)\le\mathsf G(\gamma).
\]
\end{proposition}
\begin{proof}
\emph{Three pieces.} Let $K_z=\ker\zeta_1\cap\ker\zeta_2$, the pairs of vectors of $X_0$; $V_m=
\Lambda_\cap\cap K_z$; $V_u$ the orthogonal complement of $V_m$ in $\Lambda_\cap$, of dimension $r$. The
first projection maps $V_m$ isometrically (link metric to $h$) onto $\pbeCL(\gamma,\gamma')$, by the
definition of $\pbeCL$ (a pair in $\Lambda_\gamma$ is determined by its source, which for a pair in
$K_z$ is its first member). The space $\Lambda_\gamma\cap K_z$ is the graph of $\Upsilon_\gamma$ on
$X_0^\gamma$, of dimension $k-1$ if $\pbg(\gamma)\ne0$ and $k$ otherwise (Lemma~\ref{pbe:link}(c) and
Proposition~\ref{ppi:cell}), that is, $k-(j-\varrho)_+$. As $\Lambda_\gamma\cap K_z+\Lambda_\cap\subseteq
\Lambda_\gamma$, $\dim V_m\ge k-(j-\varrho)_++\dim\Lambda_\cap-(k+1)$, so $r\le1+(j-\varrho)_+$; and if
$i=0$ and $j=\varrho+1$, then $\Lambda_\cap=\Lambda_\gamma$, $V_m=\Lambda_\gamma\cap K_z$ and $r=2$. The
forms $\zeta=(\zeta_1,\zeta_2)$ are injective on $V_u$, as $V_u\cap K_z\subseteq V_u\cap V_m=0$.

Apply Lemma~\ref{pbg:basic}(b) to $\Lambda_\gamma=\Lambda_\cap\perp V_y$ and $L^+_\gamma$, and then to
$\Lambda_\cap=V_m\perp V_u$ and $L^+_\gamma\cap\Lambda_\cap$: with $L_m=L^+_\gamma\cap V_m$ and $L_u$ the
projection of $L^+_\gamma\cap\Lambda_\cap$ to $V_u$,
\[
\pgnd(L^+_\gamma)=\pgnd(L_m)\,\pgnd(L_u)\,\pgnd(L_y).
\]

\emph{The common part.} Let $P=([t],[\gamma^{-1}t])\in L_m$, $t=\alpha+m\in\Lambda\cap(A+\gamma A)^\perp$
with $\alpha\in M(\Lambda\cap A)$. As $\zeta_1(P)=h(z,t)=0$, $m\in X_0$, and $m=t-\alpha\in\Lambda$, so
$m\in\ppiLc$; it is the first member of $P$, so $m\in\pbeCL(\gamma,\gamma')$. As $\zeta_2(P)=h(\gamma
z,t)=0$ and $t\perp\gamma A$, the vector $m+\alpha\in m+M(\Lambda\cap A)$ is orthogonal to $\gamma f_0,
\dots,\gamma f_\varrho$, so $m\in L_\gamma$. Hence the first projection maps $L_m$ isometrically into
$C_\gamma\subseteq\ppiLc\cap\pbeCL(\gamma,\gamma')$; its image has rank $\dim\pbeCL$. So $C_\gamma$ is a
lattice of that rank and, by Lemma~\ref{pbg:basic}(a), $\pgnd(C_\gamma)\le\pgnd(L_m)$.

\emph{The part seen by the $z$-forms.} Choose $\rho\subseteq\{1,2\}$ with $\#\rho=r$ such that
$\zeta_\rho$ is an isomorphism of $V_u$ onto $E^r$ ($\rho=\{1,2\}$ if $r=2$). By
Lemma~\ref{pbg:linkforms}(d), $\zeta_\rho=\pgnlm{\varepsilon_\rho}\cdot$ on $\Lambda_\gamma$. For $P\in
L_u$, $P=P'-v$ with $P'\in L^+_\gamma\cap\Lambda_\cap$ and $v\in V_m\subseteq K_z$, so $\zeta(P)=
\zeta(P')=(h(z,t),h(\gamma z,t))$ for a lift $t\in\Lambda$ of $P'$, which is integral. So the ideal
$\mathfrak p$ generated by the $\det\zeta_\rho(\ell)$, $\ell$ an $r$-tuple of $L_u$, is integral and
nonzero, and $\pgnN\mathfrak p\ge x$: trivially if $x=1$, and by Lemma~\ref{pbg:meet} if $x\ne1$ (then
$r=2$). Apply Lemma~\ref{pbg:basic}(c) to the free lattice $\sum_{l\in\rho}\auO\varepsilon_l$ (whose
$r$-tuples give the same ideal $\mathfrak p$) and $L_u$; the Gram matrix of $\varepsilon_\rho$ is a
principal minor of that of $(\Upsilon_\gamma^{-1}z^+,z^+)$, so by Lemma~\ref{pbg:linkvec} (for $r=0$ all
is trivial)
\[
x\,\mu(\gamma)\le\mathsf M(\gamma)\,\pgnd(L_u).
\]

\emph{Conclusion.} With Proposition~\ref{pbg:Lplus},
\[
t(\gamma)\pgnd(C_\gamma)\,x\,\pgnd(L_y)\le t(\gamma)\pgnd(L_m)\pgnd(L_u)\pgnd(L_y)\frac{\mathsf M}{\mu}=
t(\gamma)\pgnd(L^+_\gamma)\frac{\mathsf M}{\mu}\le\frac{M^{2j'f}\mathsf M(\gamma)}{i_\gamma i'_\gamma}=
\mathsf G(\gamma). \qedhere
\]
\end{proof}

\subsubsection{The congruence and the index}\label{pbg:ss-congr}

\begin{proposition}[The congruence]\label{pbg:congr}
Let $\gamma,\gamma'\in\Gamma_{c,s,c'}\cap\Gamma^*$, and let $L_y\subseteq V_y\subseteq\Lambda_\gamma$ and
$L'_y\subseteq V'_y\subseteq\Lambda_{\gamma'}$ be as in Proposition~\ref{pbg:cut}, for $(\gamma,\gamma')$
and for $(\gamma',\gamma)$, both of rank $i$. Then $1\le(2/M)^{2if}\,\pgnd(L_y)\,\pgnd(L'_y)$.
\end{proposition}
\begin{proof}
If $i=0$ both sides are $1$. \emph{Nondegeneracy.} Let $R\in V'_y$ with $\pgnlf PR=0$ for all $P\in V_y$.
As $\Lambda_\cap\subseteq\Lambda_{\gamma'}$, $\pgnlf PR=0$ also for $P\in\Lambda_\cap$
(Lemma~\ref{pbg:linkforms}(a)), so $R$ is orthogonal to $\Lambda_\gamma$, hence $R\in\Lambda_\gamma$ by
Lemma~\ref{pbg:linkforms}(a); so $R\in\Lambda_\cap\cap V'_y=0$. Thus $\pgnlf\cdot\cdot$ pairs $V_y$ and
$V'_y$ nondegenerately.

\emph{Integrality.} $\delta=\gamma\gamma'^{-1}$ satisfies $h_c^{-1}\delta h_c=k_\gamma k_{\gamma'}^{-1}\in K$
with reduction $ss^{-1}=1$, so $h_c^{-1}\delta h_c\in J=J_M$ and $(\delta-1)\Lambda\subseteq M\Lambda$.
For $P\in L^1_\gamma$, $R\in L^1_{\gamma'}$ with lifts $t,t'\in\Lambda$, Lemma~\ref{pbg:linkforms}(b)
gives $\pgnlf PR=h(t,(1-\delta)t')\in M\auO$. An element of $L_y$ is $P-C$ with $P\in L^+_\gamma\subseteq
L^1_\gamma$ and $C\in\Lambda_\cap$, and an element of $L'_y$ is $R-C'$ likewise; as $C,C'\in\Lambda_
\gamma\cap\Lambda_{\gamma'}$, Lemma~\ref{pbg:linkforms}(a) gives $\pgnlf{P-C}{R-C'}=\pgnlf PR\in M\auO$.

\emph{The norm bound.} Let $u=\Upsilon_\gamma^{-1}\Upsilon_{\gamma'}$, $S=M^{-1}\pgnsrc(L_y)$ and $S'=(1-u)
\pgnsrc(L'_y)$, lattices in $Q^+$. By Lemma~\ref{pbg:linkforms}(c), $\pgnlf PR=h^+(\pgnsrc P,(1-u)\pgnsrc
R)$ for $P\in V_y$, $R\in V'_y$. So $1-u$ is injective on $\pgnsrc V'_y$ (by the nondegeneracy), and the
ideal generated by the determinants $\det h^+(\ell,\ell')$, $\ell$, $\ell'$ running over the $i$-tuples
of $S$, $S'$, is $M^{-i}$ times the ideal generated by the $\det\pgnlf{\ell_1}{\ell'_1}$ over the
$i$-tuples of $L_y$, $L'_y$; it is integral and nonzero. Lemma~\ref{pbg:basic}(c) with $q=h^+$ and $x=1$,
then Lemma~\ref{pbg:basic}(d) and $\pgnd_{h^+}(\pgnsrc L)=\pgnd(L)$, give
\[
1\le\pgnd_{h^+}(S)^2\,\pgnd_{h^+}(S')^2\le M^{-4if}\pgnd(L_y)^2\cdot4^{2if}\pgnd(L'_y)^2 ,
\]
as scaling by $M^{-1}$ multiplies the Gram determinants of $i$-tuples by $M^{-2i}$. Take square roots.
\end{proof}

\begin{proposition}[The index]\label{pbg:index}
Let $\gamma,\gamma'\in\Gamma_{c,s,c'}$, and let $C_\gamma=L_\gamma\cap\pbeCL(\gamma,\gamma')$. Then
$[C_\gamma:C_\gamma\cap L_{\gamma'}]$ divides $t(\gamma')$; note $C_\gamma\cap L_{\gamma'}=L_{\gamma\gamma'}$.
\end{proposition}
\begin{proof}
$L_{\gamma\gamma'}=L_\gamma\cap L_{\gamma'}\cap\pbeCL=C_\gamma\cap L_{\gamma'}$. Let $\phi(v)=(h(\gamma'f_b,v)
)_{b\le\varrho}$ and $R'=\phi(\Lambda\cap A)$. A vector $v\in\ppiLc$ lies in $L_{\gamma'}$ if and only if
$\phi(v)\in MR'$ ($v+M\omega\perp\gamma'f$ means $\phi(v)=-M\phi(\omega)$). So $[C_\gamma:C_\gamma\cap
L_{\gamma'}]=[\phi(C_\gamma):\phi(C_\gamma)\cap MR']$.

\emph{Claim: $\phi(C_\gamma)\subseteq M\phi(\Lambda\cap(A+\gamma'A))+MR'$.} Let $v\in C_\gamma$, and
$\omega\in\Lambda\cap A$ with $y=v-M\omega\perp\gamma f_0,\dots,\gamma f_\varrho$. Then $y\in\Lambda$,
$y\perp A$ (as $v\in X_0$), so $y\in(A+\gamma A)^\perp$ and $y\perp\gamma z$: the pair $([y],[\gamma^{-1}
y])=(v,v')$ lies in $\Lambda_\gamma$ with $v'\in X_0$. As $v\in\pbeCL(\gamma,\gamma')$, some pair $(v,v'')\in
\Lambda_\gamma\cap\Lambda_{\gamma'}$ has $v''\in X_0$; both pairs lie in $\Lambda_\gamma$ and have the
source $v$, so $v''=v'$ and $(v,v')\in\Lambda_{\gamma'}$: there is $t'\perp A+\gamma'A$ with $[t']=v$,
$[\gamma'^{-1}t']=v'$. Then $t'-y\in A$ and $\gamma^{-1}y-\gamma'^{-1}t'\in A$, and with $m=\gamma'\gamma^{-1}$,
\[
(m-1)y=(t'-y)+\gamma'(\gamma^{-1}y-\gamma'^{-1}t')\in A+\gamma'A .
\]
As in Proposition~\ref{pbg:congr}, $(m-1)\Lambda\subseteq M\Lambda$, so $w=-M^{-1}(m-1)y\in\Lambda\cap(A+
\gamma'A)$. Since $my\perp m\gamma f=\gamma'f$, $\phi(y)=\phi(my)-\phi((m-1)y)=M\phi(w)$; and $\phi(y)=
\phi(v)-M\phi(\omega)$. So $\phi(v)=M\phi(w)+M\phi(\omega)$.

\emph{Conclusion.} $\phi(\gamma'(\Lambda\cap A))=0$ ($h(\gamma'f_b,\gamma'e_a)=h(f_b,e_a)=0$), so $M\phi$
maps $T'_0=\Lambda\cap A+\gamma'(\Lambda\cap A)$ into $MR'$; by the claim, $\phi(C_\gamma)\subseteq
M\phi(\Lambda\cap(A+\gamma'A))+MR'$. Lemma~\ref{pbg:indexmap} (with $B'=\Lambda\cap(A+\gamma'A)$,
$D=T'_0$, $K=MR'$, $K'=\phi(C_\gamma)$) shows that $[\phi(C_\gamma):\phi(C_\gamma)\cap MR']$ divides $[B':
T'_0]=t(\gamma')$, which is finite as $T'_0$ spans $A+\gamma'A$.
\end{proof}

\subsubsection{Statement (G$'$)}\label{pbg:ss-main}

\begin{theorem}\label{pbg:main}
Let $s$ agree, $\gamma,\gamma'\in\Gamma_{c,s,c'}\cap\Gamma^*$, $C_\gamma=L_\gamma\cap\pbeCL(\gamma,\gamma')$,
$C_{\gamma'}=L_{\gamma'}\cap\pbeCL(\gamma',\gamma)$, $i=i(\gamma,\gamma')$, and $\mathsf x=\mathsf{mn}(s)$ if
$i=0$ and $\pbg(\gamma)\ne0$, $\mathsf x=1$ otherwise. Then $C_\gamma$ is a lattice of rank $\dim\pbeCL(
\gamma,\gamma')$, $[C_\gamma:L_{\gamma\gamma'}]$ divides $t(\gamma')$, and
\[
t(\gamma)\,\ppidisc C_\gamma\cdot t(\gamma')\,\ppidisc C_{\gamma'}\cdot\mathsf x^2\le(2/M)^{2if}\,\mathsf
G(\gamma)\,\mathsf G(\gamma').
\]
That is, statement \textup{(G$'$)} of Proposition~\ref{ppi:G} holds.
\end{theorem}
\begin{proof}
The first two claims are Propositions~\ref{pbg:cut} and~\ref{pbg:index}. $\pbeCL$ and $i$ are symmetric
in $\gamma,\gamma'$. If $i=0$ then $\Lambda_\gamma=\Lambda_{\gamma'}$, so $X_0^\gamma=X_0^{\gamma'}$, and by
Lemma~\ref{pbe:link}(c) and Proposition~\ref{ppi:cell} $\pbg(\gamma)\ne0$ if and only if $\pbg(\gamma')\ne
0$. So $x=\mathsf x$ is admissible in Proposition~\ref{pbg:cut} for $(\gamma,\gamma')$ and for $(\gamma',
\gamma)$, and multiplying the two bounds,
\[
t(\gamma)\ppidisc C_\gamma\,t(\gamma')\ppidisc C_{\gamma'}\,\mathsf x^2\,\pgnd(L_y)\pgnd(L'_y)\le\mathsf G(
\gamma)\mathsf G(\gamma').
\]
Multiply by $(2/M)^{2if}$ and use Proposition~\ref{pbg:congr}. Here $\ppidisc=\pgnd$ on lattices of
$\ppiLc$, and $t(\gamma)$ is the torsion index of (G$'$) as defined in
Proposition~\ref{ppi:G}; the two expressions given there agree because $h_c,h_{c'}$ fix $A$ pointwise,
so $\gamma A=h_ck_\gamma A$ and $\gamma(h_cL\cap A)=h_ck_\gamma(L\cap A)$.
\end{proof}



% ===== B-count.tex =====
% proofs/B-count.tex
% Statements (C1) and (C2) of Proposition ppi:C (proofs/B-pillars.tex): the count of near terms over
% one lattice, and the sum over the lattices of a given corank. Labels of this file start with pbc:.
\providecommand{\auO}{\mathcal O}
\providecommand{\auNm}{\operatorname{Nm}}
\providecommand{\auA}{\mathbb A}
\providecommand{\aubar}[1]{\overline{#1}}
\providecommand{\auf}[1]{\texttt{#1}}
\providecommand{\pbR}{\overline R_0}
\providecommand{\pbet}{\beta_*}
\providecommand{\pbeGs}{\Gamma^*}
\providecommand{\pbeAp}{A^{+}}
\providecommand{\pbel}{\ell}
\providecommand{\ppiLc}{\Lambda^{X}_{c}}
\providecommand{\ppidisc}{\operatorname{disc}}
% macros of this file (prefix \pbc)
\providecommand{\pbcD}{\mathfrak D}
\providecommand{\pbcB}{\mathcal B}
\providecommand{\pbcadj}{\operatorname{adj}}
\providecommand{\pbcspan}{\operatorname{span}}
\providecommand{\pbcrd}{\mathrm{rd}_E}
\providecommand{\pbcHom}{\operatorname{Hom}}

\subsection{The count over one lattice and the sum over lattices}\label{pbc:sec}

This section proves the two statements (C1) and (C2) of Proposition~\ref{ppi:C}, from which that
proposition deduces the count (C) of Theorem~\ref{ppi:pillars}. We use the notation of
\S\S\ref{pbe:sec} and~\ref{ppi:sec}: the frame $f_0,\dots,f_\varrho$ ($f_a=e_{a+1}$ for $a<\varrho$,
$f_\varrho=z$), $n=\varrho+1$, $A=\langle e_1,\dots,e_\varrho\rangle$, the decomposition $Y_0=A\oplus
X_0\oplus Ez\oplus A^*$ with $h(e_a,e_b^*)=\delta_{ab}$, the lattice $L$ of $Y_0$, the classes
$c\in\mathcal C$ with representatives $h_c\in H(\auA_f)$, the sets $\Gamma_{c,s,c'}$, the groups $U_c$,
$U_c^0=U_c\cap U^0$, the level $M$, the prime $N=\pi_N\aubar{\pi_N}$, the number $f$ of places
$\tau\in\auPl$ (so $[E:\Q]=2f$), $\pbet=b^{kf/4}$, $\ppiLc=h_cL\cap X_0$, the lattices $L_\gamma$, the
discriminant $\ppidisc$, the labels, and the cell rank $j(\gamma)$ with the leading minors
$\delta_l(\gamma)$ of the matrix $Z(\gamma)=(h(f_a,\gamma f_b))_{0\le a,b\le\varrho}$. For $x\in E$,
$\auNm x=\prod_{\tau}|\tau(x)|^2=|N_{E/\Q}(x)|$, and for a nonzero ideal $\mathfrak a$ of $\auO_E$,
$\auNm\mathfrak a=\#(\auO_E/\mathfrak a)$. Lattices are finitely generated $\auO_E$-submodules; a
lattice in $X_0$ given by local data (such as $\ppiLc$, or $u\auO_E^k$ for $u\in H(\auA_f)$) is the
$\auO_E$-lattice with these completions.

In this section, as in \S\ref{pbg:sec} and \S\ref{ppi:ss-cell}, hermitian forms are written linear in the second variable
and conjugate-linear in the first, the opposite of convention (Cv1) of \S\ref{pb:ss-setting}: we write $h$ for the form
$(x,y)\mapsto h(y,x)$. This is the reading under which the map $u$ of Lemma~\ref{pbc:moves} is linear, $\beta_{ab}=h(e_a,\gamma f_b)$
is the $e_a^*$-coordinate of $\gamma f_b$, the map $v\mapsto\alpha$ in Step~2 of the proof of Theorem~\ref{pbc:C1} is
$\auO_E$-linear, and \eqref{pbc:eq-lift} holds as written. Orthogonality, integrality, lattices and their duals, indices,
discriminants, isometries, $U_c^0$, and the norms $\auNm$ of the entries and minors of $Z(\gamma)$ (which is replaced by its
entrywise conjugate) do not change under this replacement (\S\ref{pbv:ss-weak}(iv)); so (C1) and (C2), which are counts,
and Lemmas~\ref{pbc:box}, \ref{pbc:indep}, which concern norms, hold in either convention, while the explicit formulas
inside the section (such as the map $u$ of Lemma~\ref{pbc:moves}) are to be read in the convention just stated.

We define the reduced and boxed points (Definition~\ref{pbc:def-reduced}) and prove the bound
$\|y_\tau^{-1}\tilde f_a\|\le b^{37n}$ at boxed points used in Lemma~\ref{ppi:gauge} (Lemma~\ref{pbc:box}).
We then prove (C1) (Theorem~\ref{pbc:C1}), under $\varrho\ge280$ and the sizes of {\cite[Lemma~13.4]{P46}},
and (C2) (Theorem~\ref{pbc:C2}), using Minkowski's second theorem for lattices in
$\R^D$~\cite[Chapter~VIII, \S4.3, Theorem~V]{Cassels} (see also \cite[Theorem~1.3]{HenkMinima}), under the
size conditions of Remark~\ref{pbc:rem-sizes}, which are deduced there from {\cite[Lemma~13.4]{P46}}.

\subsubsection{Reduced and boxed points}\label{pbc:ss-box}

At a place $\tau$ and for $y\in G_\tau$ let $w_a=y^{-1}\tilde f_a$ and let $o_a$, $t_{a'a}$ be as in
Lemma~\ref{pbe:len}(c): $w_a=o_a+\sum_{a'<a}t_{a'a}(y)\,o_{a'}$. For $y=(y_\tau)\in G_\infty$ and
$a'<a$ put
\[
T_{a'a}(y)=\Bigl(t_{a'a}(y_\tau)\,\frac{s_{a'}}{s_a}\Bigr)_{\tau\in\auPl}\in\C^{\auPl},
\]
with the scales $s_a=s_{a,\tau}$ of \eqref{pbe:eq-scale}. (These are the coefficients of the
Gram--Schmidt process read in the frame $f$ instead of the normalised frame $\tilde f$; Step~1 of
Proposition~\ref{pbe:count} shows how they change under $U_{c'}$.) Let $M\auO_E\subset\C^{\auPl}$ be the
lattice of the vectors $(\tau(Mx))_\tau$, $x\in\auO_E$.

\begin{lemma}\label{pbc:domain}
There is a Borel set $\pbcD\subset\C^{\auPl}$ that contains exactly one point of each coset of
$M\auO_E$ and whose points have all coordinates of absolute value at most $b^{26}$.
\end{lemma}
\begin{proof}
Let $Q=\{w\in\C^{\auPl}:|w_\tau|\le b^{26}\ \forall\tau\}$. By Step~1 of Proposition~\ref{pbe:count}
every coset of $M\auO_E$ meets $Q$; it meets $Q$ in a finite set, because $M\auO_E$ is discrete. Let
$\lambda_1,\lambda_2,\dots$ be an enumeration of $M\auO_E$, let $P$ be a half-open fundamental
parallelepiped of $M\auO_E$ (a Borel set of representatives), and for $x\in P$ let $i(x)$ be the least
$i$ with $x+\lambda_i\in Q$. Then $\pbcD=\{x+\lambda_{i(x)}:x\in P\}$ is a set of representatives in
$Q$, and it is Borel: $\pbcD=\bigcup_i\bigl(((P\cap(Q-\lambda_i))\setminus\bigcup_{i'<i}(Q-
\lambda_{i'}))+\lambda_i\bigr)$.
\end{proof}

\begin{definition}[Reduced and boxed points]\label{pbc:def-reduced}
Fix $\pbcD$ as in Lemma~\ref{pbc:domain}. A point $y\in G_\infty$ is \emph{reduced} if $T_{a'a}(y)\in
\pbcD$ for all $0\le a'<a\le\varrho$, and \emph{boxed} if moreover $\pbel_\tau(y_\tau)\le b^{6n+2}$ for
every $\tau$. An element $\gamma\in\Gamma$ is \emph{near} if there is a boxed $y$ with $\gamma^{-1}y$
boxed (here $\gamma$ acts through its image in $G_\infty$).
\end{definition}

This is the set $\mathfrak D$ of reduced points used in \S\ref{ppi:ss-statement}. The set of boxed points
is Borel (the $t_{a'a}$ and $\pbel_\tau$ are continuous on $G_\tau$), and both sets are invariant under
left multiplication by every $s\in\prod_\tau G_\tau$ with $s_\tau^{-1}f_a=f_a$ for all $a$ and
$\tau$, and lying in $\prod_\tau R_\tau$: such $s$ do not change the $w_a$, nor $\pbel_\tau$ by
Lemma~\ref{pbe:len}(a). In particular they are invariant under $U^0(E^+)$.

\begin{lemma}[Box]\label{pbc:box}
Let $y$ be boxed. Then $|t_{a'a}(y_\tau)|\le b^{26}$ for $a'<a$, and $\|y_\tau^{-1}\tilde f_a\|\le
b^{27}\pbel_\tau(y_\tau)\le b^{37n}$ for all $\tau$ and $a$.
\end{lemma}
\begin{proof}
By \eqref{pbe:eq-scale}, $|s_a/s_{a'}|\le1$ for $a'<a$, so $|t_{a'a}|=|\tau\text{-coordinate of }T_{a'a}|
\cdot|s_a/s_{a'}|\le b^{26}$. By Lemma~\ref{pbe:len}(c), $\|o_{a'}\|<\pbel_\tau(y_\tau)$, hence $\|w_a\|\le
\sum_{a'\le a}|t_{a'a}|\,\|o_{a'}\|\le n\,b^{26}\pbel_\tau(y_\tau)\le b^{27}b^{6n+2}\le b^{37n}$, using
$n\le b$ and $t_{aa}=1$.
\end{proof}

\begin{remark}\label{pbc:rem-box}
Lemma~\ref{pbc:box} is the input of Lemma~\ref{ppi:gauge}. Thus, for near $\gamma$, $|\tau(Z(\gamma)_{
ab})|\le b^{134n}$ and $\auNm\delta_l(\gamma)\le b^{270n^2f}$ for all $\tau,a,b,l$; these two bounds are
the only properties of nearness used in (C1), besides Lemma~\ref{pbc:box} itself.
\end{remark}

\subsubsection{Tools}\label{pbc:ss-tools}

\paragraph*{Lattices and discriminants}
For a lattice $L'$ of rank $\rho$ in a subspace $V\subseteq X_0$ with $h(L',L')\subseteq\auO_E$, let
$\mathfrak d(L')$ be the ideal of $\auO_E$ generated by the Gram determinants $\det(h(v_i,v_{i'}))_{i,
i'\le\rho}$ of the $\rho$-tuples of $L'$, let $L'^\vee=\pbcHom_{\auO_E}(L',\auO_E)$, and let $L'^h
\subseteq L'^\vee$ be the image of $L'$ under $x\mapsto h(x,\cdot)|_{L'}$ (injective, as $h$ is definite).

\begin{lemma}\label{pbc:disc}
Let $L'\subseteq\Lambda$ be lattices in $X_0$, $\Lambda$ integral, $V=EL'$ of dimension $\rho$, $\Lambda_V=
\Lambda\cap V$.
\begin{enumerate}
\item[(a)] $\ppidisc L'=\auNm\mathfrak d(L')^{1/2}$; in particular $\ppidisc L'=[\Lambda_V:L']\,\ppidisc
\Lambda_V$, $\ppidisc L'$ does not depend on $\Lambda$, $\ppidisc(\eta L')=\ppidisc L'$ for $\eta\in
H(E^+)$, and $\ppidisc L'\ge1$.
\item[(b)] $[L'^\vee:L'^h]=(\ppidisc L')^2$, and $[L'^\vee:\Lambda_V^\vee]=[\Lambda_V:L']$ (restriction
$\Lambda_V^\vee\hookrightarrow L'^\vee$).
\item[(c)] If $\Lambda$ is unimodular, then $\ppidisc(\Lambda\cap V)=\ppidisc(\Lambda\cap V^\perp)$, where
$V^\perp$ is the orthogonal of $V$ in $X_0$.
\end{enumerate}
\end{lemma}
\begin{proof}
All indices are products of local indices, so we may localise at a prime $\mathfrak p$ of $E^+$; then
$\auO=\auO_{E,\mathfrak p}$ is a principal ideal ring, all lattices are free, and for an ideal $\mathfrak
a$, $\auNm\mathfrak a$ is the order of the quotient. (a) If $v$ is a basis of $L'$ with Gram matrix $G$,
the Gram determinant of the tuple $vC$ ($C$ a $\rho\times\rho$ matrix over $\auO$) is $\det(C)^-\det(C)
\det G$, so $\mathfrak d(L')=(\det G)$. If $v'$ is a basis of $\Lambda_V$ and $v=v'C$, then $[\Lambda_V:L']=
\auNm(\det C)$ and $\auNm\mathfrak d(L')=\auNm(\det C)^2\auNm\mathfrak d(\Lambda_V)$; this is the first
formula, given the definition of $\ppidisc$ in \S\ref{ppi:ss-statement}. An isometry preserves Gram
determinants, and $\mathfrak d(L')$ is integral and nonzero. (b) In the basis $v$ and the dual basis
of $L'^\vee$, the map $L'\to L'^\vee$ has the matrix ${}^tG$ (up to the conjugation, which preserves
norms), so its cokernel has order $\auNm(\det G)$. The second claim: in the bases $v'$ and $v=v'C$
the restriction has the matrix ${}^tC$. (c) Put $\Lambda_1=\Lambda\cap V$ and $\Lambda_2=\Lambda\cap
V^\perp$. As $\Lambda/\Lambda_1$ is torsion free, $\Lambda_1$ is a direct summand of $\Lambda$, and as
$\Lambda$ is unimodular, $x\mapsto h(x,\cdot)|_{\Lambda_1}$ maps $\Lambda$ onto $\Lambda_1^\vee$, with
kernel $\Lambda_2$. Its composite with $\Lambda_1^\vee\to\Lambda_1^\vee/\Lambda_1^h$ has kernel $\Lambda_1
\oplus\Lambda_2$ (if $h(x,\cdot)=h(x_1,\cdot)$ on $\Lambda_1$ with $x_1\in\Lambda_1$, then $x-x_1\in
\Lambda_2$). So $[\Lambda:\Lambda_1\oplus\Lambda_2]=(\ppidisc\Lambda_1)^2$ by (b), and by symmetry ($V=
(V^\perp)^\perp$, $\Lambda_2$ is also a direct summand) it equals $(\ppidisc\Lambda_2)^2$.
\end{proof}

\begin{lemma}[Packing]\label{pbc:pack}
\textup{(a)} Let $m\ge1$ and $R\ge0$. At most $(1+2R)^{2mf}$ vectors $v\in\auO_E^m$ satisfy $|\tau
(v_i)|\le R$ for all $i$ and $\tau$.
\textup{(b)} Let $\Sigma\subset\R^D$ be a set with $\|x\|\ge1$ and $\|x-x'\|\ge1$ for all $x\ne x'$ in
$\Sigma$ (Euclidean norm). Then $\sum_{x\in\Sigma}\|x\|^{-s}\le5^s$ for every real $s\ge D+2$.
\end{lemma}
\begin{proof}
(a) If $v\ne v'$, some $v_i-v'_i$ is a nonzero algebraic integer, so $\auNm(v_i-v'_i)\ge1$ and
$|\tau(v_i-v'_i)|\ge1$ for some $\tau$. Hence the polydiscs of radius $\frac12$ in $\C^{m\auPl}$ around
the points $(\tau(v_i))_{i,\tau}$ are disjoint; they lie in the polydisc of radius $R+\frac12$, and
comparing volumes gives the bound. (b) The balls of radius $\frac12$ around the points of $\Sigma$
are disjoint, so for $T\ge1$ the number $\nu(T)$ of $x\in\Sigma$ with $\|x\|\le T$ is at most
$(2T+1)^D\le(3T)^D$. Hence $\sum\|x\|^{-s}=\int_1^\infty sT^{-s-1}\nu(T)\,dT\le s3^D\int_1^\infty
T^{D-s-1}dT=s3^D/(s-D)\le\frac s2\,3^s\le5^s$.
\end{proof}

\paragraph*{Moves}
The following elements of $U_c^0$ are used to move frames.

\begin{lemma}[Moves]\label{pbc:moves}
Let $c\in\mathcal C$, let $x_1,\dots,x_\varrho\in\ppiLc$, and let $S=(S_{ba})$ be a $\varrho\times\varrho$
matrix over $\auO_E$ with ${}^t\aubar S=-S$. Let $G=(h(x_b,x_a))_{b,a}$, let $\omega\in\auO_E$ be as in
Lemma~\ref{ppi:eichler} ($\omega+\aubar\omega=1$), put $W=M\omega G-S$, and
\[
u(t)=t+M\sum_{a=1}^\varrho h(x_a,t)\,e_a-M\sum_{a=1}^\varrho h(e_a,t)\Bigl(x_a+\sum_{b=1}^\varrho W_{ba}
e_b\Bigr)\qquad(t\in Y_0).
\]
Then $u\in U_c^0$. For $\gamma\in\Gamma$ write $\gamma f_b=\sum_a\alpha_{ab}e_a+\xi_b+\zeta_bz+\sum_a
\beta_{ab}e_a^*$ with $\xi_b\in X_0$, and let $\alpha(\gamma)=(\alpha_{ab})$, $\xi(\gamma)=(\xi_0,\dots,
\xi_\varrho)$, $\beta(\gamma)=(\beta_{ab})$ ($1\le a\le\varrho$, $0\le b\le\varrho$). Then $Z(u\gamma)=Z(
\gamma)$, $\beta(u\gamma)=\beta(\gamma)$, $\zeta(u\gamma)=\zeta(\gamma)$, and
\[
\xi(u\gamma)=\xi(\gamma)-M\,X\beta(\gamma),\qquad\text{and, if }X=0,\qquad\alpha(u\gamma)=\alpha(\gamma)+
M\,S\beta(\gamma),
\]
where $X\beta$ is the row $(\sum_ax_a\beta_{ab})_b$.
\end{lemma}
\begin{proof}
Write $u=1+M\nu$. \emph{Unitarity.} Put $\mathsf a_a=h(x_a,t)$, $\mathsf b_a=h(e_a,t)$ and similarly
$\mathsf a'_a,\mathsf b'_a$ for $t'$. Using that $A$ is isotropic and orthogonal to $X_0$, and $h(e_a,
e_b^*)=\delta_{ab}$, one finds $h(\nu t,t')+h(t,\nu t')=-\sum_{a,b}\aubar{\mathsf b_a}\mathsf b'_b(W+
{}^t\aubar W)_{ab}$ (the terms with $\mathsf a$ cancel in pairs) and $h(\nu t,\nu t')=\sum_{a,b}
\aubar{\mathsf b_a}\mathsf b'_bG_{ab}$. As $W+{}^t\aubar W=M(\omega+\aubar\omega)G=MG$, $h(ut,ut')=h(t,
t')+M\bigl(h(\nu t,t')+h(t,\nu t')+Mh(\nu t,\nu t')\bigr)=h(t,t')$. \emph{Unipotence.} $\nu$ vanishes on
$A$, maps $A^\perp=A\oplus X_0\oplus Ez$ into $A$ and $Y_0$ into $A\oplus X_0\subseteq A^\perp$. So $u$
is the identity on $A$, preserves the flag of $A$, acts trivially on $A^\perp/A$ and on $Y_0/A^\perp$;
as in the proof of Lemma~\ref{ppi:eichler}, $u\in U(E^+)$; also $\nu^3=0$. \emph{The frame.} $h(x_a,
f_b)=0$ and $h(e_a,f_b)=0$, so $\nu f_b=0$ and $u\in U^0(E^+)$. \emph{Level.} Let $L^\sharp=h_cL$; as
in the proof of Lemma~\ref{pbe:eichler}, $L^\sharp=\bigoplus_a(\widehat\auO e_a\oplus\widehat\auO
e_a^*)\oplus L^\sharp_{\rm mid}$ is unimodular, so $h(x_a,t),h(e_a,t)\in\widehat\auO$ for $t\in
L^\sharp$, and $x_a,e_a\in L^\sharp$, $W$ is integral; hence $\nu L^\sharp\subseteq L^\sharp$ and $(u-1)
L^\sharp\subseteq ML^\sharp$, $(u^{-1}-1)L^\sharp=(-M\nu+M^2\nu^2)L^\sharp\subseteq ML^\sharp$. By the
description of $U_c$ in the proof of Lemma~\ref{pbe:eichler}, $u\in U_c$.
\emph{Action.} $h(f_a,u\gamma f_b)=h(u^{-1}f_a,\gamma f_b)=h(f_a,\gamma f_b)$, and $\beta_{ab}=h(e_a,
\gamma f_b)$, $-N\zeta_b=h(z,\gamma f_b)$ are entries of $Z$. Finally $\nu(\gamma f_b)=\sum_a h(x_a,
\xi_b)e_a-\sum_a\beta_{ab}(x_a+\sum_bW_{ba}e_b)$ has $X_0$-part $-\sum_a\beta_{ab}x_a$ and, if $X=0$
(so $W=-S$), $A$-part $\sum_{a,b'}S_{b'a}\beta_{ab}e_{b'}$.
\end{proof}

\paragraph*{Uniform independence at boxed points, and shared vectors}

\begin{lemma}\label{pbc:indep}
Let $y\in G_\tau$ with $\|y^{-1}\tilde f_a\|\le b^{37n}$ for all $a$ and $\pbel_\tau(y)\le b^{7n}$ (by
Lemma~\ref{pbc:box} this holds for every component of a boxed point). Then for every $c\in\C^n$ and every $a$,
\[
|c_a|\le b^{46n^2}\Bigl\|\sum_{a'}c_{a'}\,y^{-1}\tilde f_{a'}\Bigr\| .
\]
\end{lemma}
\begin{proof}
Let $v=\sum c_{a'}y^{-1}\tilde f_{a'}$. As $\pbel_\tau(y)\le b^{7n}$ there is $s\in R_\tau$ with $\|sy\|\le
b^{7n}$. Let $T$ be the matrix of $s$ on the normalised frame: $s\tilde f_{a'}=\sum_{a''}T_{a''a'}
\tilde f_{a''}$; it is unitriangular (\S\ref{pbe:ss-real}), so $\det T=1$ and $c=\pbcadj(T)\,Tc$. The
normalised basis is orthonormal, so $|(Tc)_{a''}|\le\|\sum_{a''}(Tc)_{a''}\tilde f_{a''}\|=\|sy\,v\|\le
b^{7n}\|v\|$, and likewise $|T_{a''a'}|\le\|sy\|\,\|y^{-1}\tilde f_{a'}\|\le b^{44n}$. The entries of
$\pbcadj T$ are minors of order $\varrho$, at most $\varrho!\,b^{44n\varrho}\le b^{(1+44n)\varrho}$ in
absolute value. Hence $|c_a|\le n\,b^{(1+44n)\varrho}b^{7n}\|v\|\le b^{46n^2}\|v\|$, as $1+(1+44n)
\varrho+7n\le46n^2$.
\end{proof}

For $\gamma\in\Gamma$ of cell rank $j=j(\gamma)$ with $\delta_j(\gamma)\ne0$, let $Y=Y_j(\gamma)$ be the
$n\times n$ matrix whose rows $a<j$ are those of $Z(\gamma)$ and whose rows $a\ge j$ are those of the
identity matrix; $\det Y=\delta_j(\gamma)$. For a coefficient vector $x\in E^n$ we write $fx=\sum x_b f_b$
and $\gamma fx=\gamma(fx)$, as in \S\ref{ppi:sec}.

\begin{lemma}[Shared vectors]\label{pbc:shared}
Let $\gamma\in\Gamma$ have cell rank $j$ with $\delta_j(\gamma)\ne0$, and $Y=Y_j(\gamma)$.
\begin{enumerate}
\item[(a)] For $j\le k\le\varrho$, $\gamma f(Y^{-1}\epsilon_k)$ lies in $\langle f_0,\dots,f_\varrho\rangle$
($\epsilon_k$ the unit vectors).
\item[(b)] If $j=\varrho+1$, the $X_0$-part $\xi$ of $\gamma f(Y^{-1}\epsilon_\varrho)$ satisfies
$\tau(h(\xi,\xi))\le1/N$ for every $\tau$.
\item[(c)] For $u\in U^0(E^+)$ the element $u\gamma$ has the same cell rank, $Z(u\gamma)=Z(\gamma)$,
$Y_j(u\gamma)=Y$, and \textup{(a)}, \textup{(b)} hold for $u\gamma$.
\end{enumerate}
\end{lemma}
\begin{proof}
(a) By definition of the cell rank there are $\zeta$ and, for each $a\ge j$, vectors $x^{(a)},y^{(a)}$
supported on $\{0,\dots,a\}$ with $fx^{(a)}=\gamma fy^{(a)}$ and $y^{(a)}_a=1$. Subtracting from
$y^{(k)}$ suitable multiples of the $y^{(a)}$, $j\le a<k$, by descending induction on the index, we
get $y$ with $y_a=\delta_{ak}$ for $a\ge j$ and $\gamma fy=fx\in\langle f\rangle$ for some $x$. For
$a<j$ we have $a<\varrho$, so $(Zy)_a=h(f_a,fx)=0$ ($A$ is isotropic and orthogonal to $z$); and $(Yy)_a=
y_a=\delta_{ak}$ for $a\ge j$. So $Yy=\epsilon_k$.
(b) Now $Y=Z$. Let $x=Z^{-1}\epsilon_\varrho$ and $t=\gamma fx$. Then $h(e_a,t)=0$ for all $a$, so $t$
has no $A^*$-part, and $h(z,t)=1$, so its $z$-coordinate is $-1/N$. Hence $h(t,t)=h(\xi,\xi)-N\cdot N^{-2}$,
while $h(t,t)=h(fx,fx)=-N|x_\varrho|^2$. So $h(\xi,\xi)=N^{-1}-N|x_\varrho|^2$, and $\tau(h(\xi,\xi))\le
1/N$.
(c) The cell rank is $U^0(E^+)$-bi-invariant (\S\ref{ppi:ss-statement}) and $Z(u\gamma)=Z(\gamma)$ as
in Lemma~\ref{pbc:moves}.
\end{proof}

\subsubsection{The count over one lattice: (C1)}\label{pbc:ss-C1}

\begin{theorem}[(C1)]\label{pbc:C1}
Assume $\varrho\ge280$, the sizes $M\le b^{23}$, $N\le b^{60n}$, $n\le b$ of {\cite[Lemma~13.4\textup{(S5)}]{P46}}, \textup{(S7)},
and $1\le|\kappa_\tau|\le b^{23}$ \eqref{pbe:eq-kappa} (Lemma~\ref{ppi:gauge} uses only $|\kappa_\tau|\ge1$). Let $c\in
\mathcal C$, $t\in\pbR$, $0\le l\le k$, $r=k-l$, and let $L'\subseteq\ppiLc$ be a lattice of rank $r$ with
$\ppidisc L'\le b^{n^3f}$. Let $\pbcB$ be the set of the pairs $(c',\gamma)$ with $c'\in\mathcal C$ and
$\gamma\in\Gamma_{c,s,c'}\cap\pbeGs$ for some $s$, such that $\gamma$ is near, chosen, of label $t$, and
$L'\subseteq L_\gamma$. Then $\#\pbcB\le b^{2000n^6(l+1)f}$.
\end{theorem}

The proof occupies the rest of this subsection. For $(c',\gamma)\in\pbcB$ we use the coordinates
$\alpha(\gamma),\xi(\gamma),\zeta(\gamma),\beta(\gamma)$ of Lemma~\ref{pbc:moves}; the \emph{frame} of
$\gamma$ is the tuple $\gamma f=(\gamma f_0,\dots,\gamma f_\varrho)$. Two facts about the lattice
$L^\sharp=h_cL=\bigoplus_a(\widehat\auO e_a\oplus\widehat\auO e_a^*)\oplus L^\sharp_{\rm mid}$ (proof of
Lemma~\ref{pbe:eichler}) are used: since $h_c^{-1}\gamma h_{c'}\in K$ and $h_{c'}$ fixes the $f_b$,
$\gamma f_b\in L^\sharp$, so $\alpha(\gamma)$ and $\beta(\gamma)$ are integral; and $L^\sharp_{\rm mid}\cap
X_0=L^\sharp\cap X_0=\ppiLc$ (by the proof of Lemma~\ref{pb:KR} the lattice of $X_0\perp Ez$ is
spanned by $x_1,\dots,x_{k-1},x_k/\aubar{\pi_N}+z/N,z$, and $h_c$ acts on $X_0$ only).

\emph{Step 0: frames modulo moves separate.} Let $(c',\gamma),(c'',\gamma')\in\pbcB$ and $u,u'\in U_c^0$
with $u\gamma f_b=u'\gamma'f_b$ for all $b$. Then $(c',\gamma)=(c'',\gamma')$.

Indeed, $h_c^{-1}uh_c\in J$ has reduction $1$, so $u\gamma\in\Gamma_{c,s,c'}$ and $u'\gamma'\in\Gamma_{c,
s',c''}$ where $s,s'$ are the reductions of $\gamma,\gamma'$. The element $g=(u\gamma)^{-1}u'\gamma'$ fixes
the frame, so (Step~2 of Proposition~\ref{pbe:count}) $g\in R_0(E^+)$, and $h_{c'}^{-1}gh_{c''}=(h_c^{-1}
u\gamma h_{c'})^{-1}(h_c^{-1}u'\gamma'h_{c''})\in K\cap R_0(\auA_f)=K_R$ has reduction $s^{-1}s'\in\pbR$. So
$s\pbR=s'\pbR$; as the labels (the components in $\pbR$ for the fixed decomposition $\auG=\bigsqcup_\rho
\rho\pbR$) agree, $s=s'$. Now Step~2 of Proposition~\ref{pbe:count}, applied to $u\gamma,u'\gamma'\in
\Gamma_{c,s,\cdot}$, gives $c''=c'$ and $g\in U_{c'}$; $g$ fixes the frame, so $g\in U_{c'}^0$. Then $\gamma'
\in U_c^0\gamma U_{c'}^0$, and both are chosen, so $\gamma'=\gamma$.

\emph{Step 1: cell rank and products.} By Proposition~\ref{ppi:cell}, every $(c',\gamma)\in\pbcB$ has a
cell rank $j(\gamma)\in\{0,\dots,\varrho+1\}$ with $\delta_{j(\gamma)}(\gamma)\ne0$; there are $\varrho+2\le
b^2$ values. The entries of $Z(\gamma)$ are integral (they are products of vectors of $L^\sharp$) and
$|\tau(Z_{ab})|\le b^{134n}$ (Remark~\ref{pbc:rem-box}); by Lemma~\ref{pbc:pack}(a) with $m=n^2$ there are
at most $(1+2b^{134n})^{2n^2f}\le b^{270n^3f}$ values of $Z(\gamma)$. We now fix $j$ and $Z$ with $\delta:=
\delta_j(Z)\ne0$, and bound the number of elements of the set $\pbcB_{j,Z}$ of $(c',\gamma)\in\pbcB$ with
$j(\gamma)=j$, $Z(\gamma)=Z$. Let $Y=Y_j(Z)$, $P=\beta(\gamma)$ (rows $0,\dots,\varrho-1$ of $Z$; it does not
depend on $\gamma$), $d=M\delta$ and $d'=M\delta\aubar\delta=M\det(Y\,{}^t\aubar Y)$. By
Remark~\ref{pbc:rem-box}, $\auNm\delta\le b^{270n^2f}$, so with $M\le b^{23}$ and $n\ge280$:
\begin{equation}\label{pbc:eq-moduli}
\auNm d=M^{2f}\auNm\delta\le b^{2n^3f},\qquad\auNm d'=M^{2f}(\auNm\delta)^2\le b^{2n^3f}.
\end{equation}

\emph{Step 2: lifts of $L'$.} Let $(c',\gamma)\in\pbcB_{j,Z}$ and let $K_\gamma\subseteq\auO_E^\varrho$ be the
set of $\alpha$ with $\sum_a\alpha_ae_a\perp\gamma f_b$ for all $b$. For $v\in L'\subseteq L_\gamma$ there is
$\alpha\in\auO_E^\varrho$ ($h_cL\cap A=\bigoplus\auO_Ee_a$) with $v+M\sum_a\alpha_ae_a\perp\gamma f_b$ for all
$b$, and $\alpha$ is unique modulo $K_\gamma$. So $v\mapsto\alpha\bmod K_\gamma$ is an $\auO_E$-linear map
$L'\to\auO_E^\varrho/K_\gamma$; as $L'$ is projective, it lifts to a linear map $\varphi_\gamma=(\varphi_{\gamma,
1},\dots,\varphi_{\gamma,\varrho})\in(L'^\vee)^\varrho$ with
\begin{equation}\label{pbc:eq-lift}
h(v,\xi_b(\gamma))+M\sum_a\aubar{\varphi_{\gamma,a}(v)}\,P_{ab}=0\qquad(v\in L',\ 0\le b\le\varrho),
\end{equation}
since $h(v,\gamma f_b)=h(v,\xi_b)$ and $h(e_a,\gamma f_b)=P_{ab}$. Colour $\pbcB_{j,Z}$ by the class of
$\varphi_\gamma$ in $(L'^\vee/L'^h)^\varrho$: by Lemma~\ref{pbc:disc}(b) there are at most $(\ppidisc
L')^{2\varrho}\le b^{2n^3\varrho f}$ colours. Fix one colour $\mathcal C_1$ and $\gamma_1\in\mathcal C_1$. For
$\gamma\in\mathcal C_1$ there are $\upsilon_a\in L'$ with $\varphi_{\gamma,a}-\varphi_{\gamma_1,a}=h(\upsilon_a,\cdot)$; let
$u_\gamma\in U_c^0$ be the move of Lemma~\ref{pbc:moves} with $x_a=-\upsilon_a$ ($\in\ppiLc$) and $S=0$. Then
$\xi(u_\gamma\gamma)=\xi(\gamma)+Y'P$ with $Y'=(M\upsilon_a)_a$, and by \eqref{pbc:eq-lift}, for $v\in L'$,
\[
h(v,\xi_b(u_\gamma\gamma))=-M\sum_a\bigl(\aubar{\varphi_{\gamma,a}(v)}-h(v,\upsilon_a)\bigr)P_{ab}=-M\sum_a
\aubar{\varphi_{\gamma_1,a}(v)}P_{ab},
\]
as $\aubar{h(\upsilon_a,v)}=h(v,\upsilon_a)$. So $h(v,\xi_b(u_\gamma\gamma))$ depends only on $v,b$ and the colour.

\emph{Step 3: the part in $X_0$.} Let $\gamma_0\in\mathcal C_1$ and put $\tilde\gamma=u_\gamma\gamma$ for
$\gamma\in\mathcal C_1$. By Step~2 the columns of $\xi(\tilde\gamma)-\xi(\tilde\gamma_0)$ are orthogonal to
$L'$; they lie in $\ppiLc$, because $\tilde\gamma f_b,\tilde\gamma_0f_b\in L^\sharp$ have the same
$z$-coordinate $\zeta_b$ (Lemma~\ref{pbc:moves}), so their difference of middle parts lies in
$L^\sharp_{\rm mid}\cap X_0=\ppiLc$. So they lie in $P'=\ppiLc\cap(EL')^\perp$, a lattice of rank $l$.
Colour $\mathcal C_1$ by the class of $\xi(\tilde\gamma)-\xi(\tilde\gamma_0)$ in $(P'/dP')^n$: there are at
most $(\auNm d)^{ln}\le b^{2n^4lf}$ colours. Fix one, $\mathcal C_2$, and $\gamma_2\in\mathcal C_2$; for
$\gamma\in\mathcal C_2$ the columns of $v=\xi(\tilde\gamma)-\xi(\tilde\gamma_2)$ lie in $d\,\ppiLc$.

We claim that $v=M\mathsf H P$ for a row $\mathsf H=(\eta_1,\dots,\eta_\varrho)$ of vectors of $\ppiLc$. Put $\mathsf H'=d^{-1}
v\,\pbcadj(Y)$, whose columns lie in $\ppiLc$ ($\pbcadj Y$ is integral), so that $vY^{-1}=M\mathsf H'$. Let $i$
satisfy $i\ge j$ or $i=\varrho$. The column $i$ of $vY^{-1}$ is the difference of the $X_0$-parts of
$\tilde\gamma f(Y^{-1}\epsilon_i)$ and $\tilde\gamma_2f(Y^{-1}\epsilon_i)$. By Lemma~\ref{pbc:shared}
(applied to $\tilde\gamma,\tilde\gamma_2$, which have cell rank $j$ and products $Z$), each of these
$X_0$-parts is $0$ (if $i\ge j$) or has $\tau(h(\cdot,\cdot))\le1/N$ (if $i=\varrho<j$); so the column
$x\in\ppiLc$ satisfies $\tau(h(x,x))\le4/N<1$ for all $\tau$. Since $h(x,x)$ is a totally nonnegative
algebraic integer, $h(x,x)=0$ and $x=0$. Hence the columns $i\ge j$ and the column $\varrho$ of $\mathsf H'$
vanish, and $v=M\mathsf H'Y=M\sum_{i<j,\,i<\varrho}\mathsf H'_i\,(\text{row }i\text{ of }Y)=M\mathsf HP$, with $\mathsf H$ the first
$\varrho$ columns of $\mathsf H'$ (rows $i<j$ of $Y$ are rows of $P$; the columns $j\le i<\varrho$ of $\mathsf H$ vanish).

Let $u'_\gamma\in U_c^0$ be the move with $x_a=\eta_a$ and $S=0$. Then $\xi(u'_\gamma\tilde\gamma)=\xi(
\tilde\gamma)-M\mathsf HP=\xi(\tilde\gamma_2)$. Put $\hat\gamma=u'_\gamma\tilde\gamma=u'_\gamma u_\gamma\gamma$ for
$\gamma\in\mathcal C_2$ (and $u'_{\gamma_2}=1$). All $\hat\gamma$ have the products $Z$, the cell rank $j$,
the same $\zeta$ and $\beta=P$, and the same $\xi$.

\emph{Step 4: shared vectors.} For $\gamma\in\mathcal C_2$ and $j\le i\le\varrho$ let $\sigma_i(\gamma)=
\alpha(\gamma)\,\pbcadj(Y)\epsilon_i\in\auO_E^\varrho$, the $A$-part of $\gamma f(\pbcadj(Y)\epsilon_i)=\delta\,
\gamma f(Y^{-1}\epsilon_i)$; put $\sigma_i=0$ for $i<j$. By Lemma~\ref{pbc:shared}(a) this vector lies in
$\langle f\rangle$, so it is fixed by every move, and $\sigma(\hat\gamma)=\sigma(\gamma)$. We bound
$\sigma(\gamma)$. Let $x=\pbcadj(Y)\epsilon_i$; the entries of $Y$ are entries of $Z$ or $0,1$, so $|\tau(
x_b)|\le\varrho!\,b^{134n\varrho}\le b^{(1+134n)\varrho}$. Write $\gamma fx=fw$ ($w\in E^n$; $\sigma_i=(w_0,
\dots,w_{\varrho-1})$). Let $y$ and $\gamma^{-1}y$ be boxed. At $\tau$, as $\tilde f_b=s_bf_b$,
\[
\sum_b\frac{\tau(w_b)}{s_b}\,y_\tau^{-1}\tilde f_b=y_\tau^{-1}(fw)=y_\tau^{-1}\gamma fx=\sum_b
\frac{\tau(x_b)}{s_b}\,(\gamma^{-1}y)_\tau^{-1}\tilde f_b .
\]
By Lemma~\ref{pbc:box} at $\gamma^{-1}y$, the right side has norm at most $n\,b^{(1+134n)\varrho}b^{37n}
\max_b|s_b|^{-1}$; by Lemma~\ref{pbc:indep} at $y$, $|\tau(w_a)|/|s_a|\le b^{46n^2}$ times this norm; and
$|s_a/s_b|\le|\kappa_\tau|^\varrho\le b^{23\varrho}$ by \eqref{pbe:eq-kappa}. So $|\tau(w_a)|\le b^{46n^2+1+(1+134n)\varrho+37n+
23\varrho}\le b^{182n^2}$ (with $\varrho=n-1$ the exponent is $180n^2-73n-23$). The entries of $\sigma$ are integral ($\gamma f_b\in L^\sharp$, $\pbcadj Y$ integral),
so by Lemma~\ref{pbc:pack}(a) with $m=\varrho n$ there are at most $(1+2b^{182n^2})^{2\varrho nf}\le
b^{366n^4f}$ values of $\sigma$. Colour $\mathcal C_2$ by $\sigma$ and fix a colour $\mathcal C_3$.

\emph{Step 5: the part in $A$.} Colour $\mathcal C_3$ by the class of $\alpha(\hat\gamma)$ modulo $d'$ (an
integral $\varrho\times n$ matrix): at most $(\auNm d')^{\varrho n}\le b^{2n^4\varrho f}$ colours. We claim
that each colour has at most one element; with Steps~1--4 and \eqref{pbc:eq-moduli} this gives
\[
\#\pbcB\le b^{2}\,b^{270n^3f}\,b^{2n^3\varrho f}\,b^{2n^4lf}\,b^{366n^4f}\,b^{2n^4\varrho f}\le
b^{2000n^6(l+1)f}.
\]
Let $\gamma,\gamma'$ have the same colour and put $\Xi=\alpha(\hat\gamma')-\alpha(\hat\gamma)=d'\Theta$, $\Theta$
integral. The frames $\hat\gamma f$ and $\hat\gamma'f$ differ only in their $A$-parts. (i) \emph{${}^t\aubar
\Theta P$ is skew-hermitian.} Both frames have the Gram matrix of $f$; expanding $h(\hat\gamma'f_b,\hat\gamma'
f_{b'})$ and using that $A$ is isotropic and $h(e_a,\hat\gamma f_{b'})=P_{ab'}$ gives ${}^t\aubar\Xi P+
{}^t\aubar P\Xi=0$; divide by $d'=\aubar{d'}$. (ii) Let $\alpha^\circ$ be the $n\times\varrho$ matrix of the
first $\varrho$ columns of $\pbcadj Y$. \emph{The columns $i\ge j$ and the column $\varrho$ of $\Theta\pbcadj Y$
vanish}: for $i\ge j$, $d'\Theta\pbcadj(Y)\epsilon_i=\sigma_i(\gamma')-\sigma_i(\gamma)=0$. If $j=\varrho+1$ (so
$Y=Z$), the rows $a<\varrho$ of $Y\pbcadj Y=\delta\cdot1$ give $P\pbcadj(Y)\epsilon_\varrho=0$ and $P\alpha^\circ
=\delta\cdot1$; then ${}^t\aubar P\Theta\pbcadj(Y)\epsilon_\varrho=-{}^t\aubar \Theta P\pbcadj(Y)\epsilon_\varrho=0$ and
$\aubar\delta \Theta\pbcadj(Y)\epsilon_\varrho={}^t\aubar{\alpha^\circ}\,{}^t\aubar P\Theta\pbcadj(Y)\epsilon_\varrho=0$.
(iii) Hence $\delta \Theta=\Theta\pbcadj(Y)Y=\Theta\alpha^\circ P$, since only the columns $i<\min(j,\varrho)$ of
$\Theta\pbcadj Y$ survive and the rows $i<j$ of $Y$ are rows of $P$. (iv) Put
\[
S=\aubar\delta\,\Theta\alpha^\circ-\delta\,{}^t\aubar{(\Theta\alpha^\circ)}-{}^t\aubar{\alpha^\circ}\,({}^t\aubar
P\Theta)\,\alpha^\circ .
\]
By (i), ${}^t\aubar S=-S$; $S$ is integral; and by (i), (iii),
\[
SP=\aubar\delta\delta \Theta-\delta\,{}^t\aubar{\alpha^\circ}\,{}^t\aubar\Theta P-{}^t\aubar{\alpha^\circ}\,{}^t
\aubar P\,\delta \Theta=\delta\aubar\delta \Theta-\delta\,{}^t\aubar{\alpha^\circ}({}^t\aubar \Theta P+{}^t\aubar P\Theta)=\delta
\aubar\delta \Theta .
\]
The move $u''\in U_c^0$ of Lemma~\ref{pbc:moves} with $x=0$ and this $S$ changes only the $A$-part, by
$MSP=d'\Theta=\Xi$. So $u''\hat\gamma f=\hat\gamma'f$, that is, $u''u'_\gamma u_\gamma\gamma f=u'_{\gamma'}
u_{\gamma'}\gamma'f$, and Step~0 gives $(c',\gamma)=(c'',\gamma')$. \hfill$\square$

\begin{remark}\label{pbc:rem-C1}
The exponent: $2+270n^3+2n^3\varrho+366n^4+2n^4\varrho+2n^4l\le2000n^6(l+1)$. The theorem is the
statement (C1) of Proposition~\ref{ppi:C}: the elements of $\pbeGs$ have a cell rank with non-vanishing
leading minor by Proposition~\ref{ppi:cell}, and no other property of $\pbeGs$ is used.
\end{remark}

\subsubsection{The sum over lattices of a given corank: (C2)}\label{pbc:ss-C2}

Let $H=\auU(X_0)$, $K^\circ\subseteq H(\auA_f)$ the stabiliser of $\auO_E^k$, $K_H\subseteq K^\circ$ the group of
Lemma~\ref{pb:KR}, and $L_1=L\cap X_0=\auO_Ex_1+\dots+\auO_Ex_{k-1}+\pi_N\auO_Ex_k$, so that $K_H$ stabilises
$L_1$ and $J_X=K_H\cap J\subseteq K_H$. Recall that $\mathcal C=H(E^+)\backslash H(\auA_f)/J_X$, that $c=
H(E^+)h_cJ_X$, and that the level is free: $H(E^+)\cap hJ_Xh^{-1}=1$ for every $h\in H(\auA_f)$
(\S\ref{pb:sec}). Then $\ppiLc=h_cL_1$. For $g\in H(\auA_f)$ write $[g]\in\mathcal C$ for its class.

\begin{lemma}[Averaging]\label{pbc:average}
Let $\mathcal A$ be a finite nonempty set, $a\colon\mathcal A\to H(\auA_f)$, $F\colon\mathcal C\to[0,
\infty)$ and $B\ge0$. If $\sum_{\phi\in\mathcal A}F([ga(\phi)])\le\#\mathcal A\cdot B$ for every $g\in
H(\auA_f)$, then $\sum_{c\in\mathcal C}F(c)\le\#\mathcal C\cdot B$.
\end{lemma}
\begin{proof}
Let $J'=J_X\cap\bigcap_\phi a(\phi)J_Xa(\phi)^{-1}$, an open compact subgroup, and $\mathcal X=H(E^+)
\backslash H(\auA_f)/J'$, a finite set ($H(E^+)\backslash H(\auA_f)$ is compact, $H$ being anisotropic).
The function $g\mapsto F([ga(\phi)])$ is right $J'$-invariant, so the hypothesis summed over $\mathcal
X$ gives $\sum_\phi\sum_{x\in\mathcal X}F([xa(\phi)])\le\#\mathcal X\#\mathcal A B$. For a fixed $\phi$,
$x\mapsto xa(\phi)$ is a bijection of $\mathcal X$ onto $H(E^+)\backslash H(\auA_f)/J''$, $J''=a(\phi)^{-1}
J'a(\phi)\subseteq J_X$, and the fibre of $H(E^+)\backslash H(\auA_f)/J''\to\mathcal C$ over $c$ is
$(H(E^+)\cap h_cJ_Xh_c^{-1})\backslash h_cJ_X/J''\cong J_X/J''$, of order $\pbvol(J_X)/\pbvol(J'')=
\pbvol(J_X)/\pbvol(J')=:\iota$ (the level is free, and $H(\auA_f)$ is unimodular). So $\sum_{x\in\mathcal
X}F([xa(\phi)])=\iota\sum_cF(c)$; likewise $\#\mathcal X=\iota\#\mathcal C$. Dividing by $\iota\#\mathcal A$
gives the claim.
\end{proof}

\begin{lemma}[Cosets]\label{pbc:index}
Let $\mathcal A=K^\circ/K_H$. For $y\in K^\circ$ let $\lambda_y$ be the last row of $y^{-1}$. Then
$yK_H\mapsto\lambda_y\bmod\pi_N$ is well defined and injective, $yL_1=\{v\in\auO_E^k:\lambda_yv\in\pi_N
\widehat\auO\}$, and $\#\mathcal A\ge N^{kf}/2$ if $4f\le5^k$.
\end{lemma}
\begin{proof}
By Lemma~\ref{pb:KR}, $y^{-1}y'\in K_H$ if and only if $\lambda_y y'\equiv(0,\dots,0,1)=\lambda_{y'}y'$
modulo $\pi_N$, that is, $\lambda_y\equiv\lambda_{y'}$. The description of $yL_1$ is that of $L_1$
transported by $y$. As in the proof of Lemma~\ref{pb:KR}, at each prime $w\mid\pi_N$ the group $H$ is
$\auGL_k(E_w)$, with $K^\circ_w=\auGL_k(\auO_w)$, and the congruence modulo $\pi_N$ is the congruence
modulo every $w\mid\pi_N$; as $K^\circ$ is a product of local groups, the image contains all rows that are
nonzero modulo every $w$. With $q_w=\auNm w\ge N\ge5$, $\prod_wq_w=\auNm(\pi_N)=N^f$ and at most $2f$
primes $w$, $\#\mathcal A\ge\prod_w(q_w^k-1)\ge N^{kf}(1-2f\cdot5^{-k})\ge N^{kf}/2$.
\end{proof}

\begin{lemma}[Forms vanishing on a lattice]\label{pbc:forms}
Let $\Lambda$ be a unimodular lattice of $X_0$, $L''\subseteq\Lambda$ a lattice of rank $r=k-l$, $V=EL''$
and $i=[\Lambda\cap V:L'']$. Then at most $i\,N^{fl}$ classes $\lambda\in\Lambda^\vee/\pi_N\Lambda^\vee$
satisfy $\lambda(L'')\subseteq\pi_N\auO_E$.
\end{lemma}
\begin{proof}
$\Lambda_V=\Lambda\cap V$ is a direct summand: $\Lambda=\Lambda_V\oplus\Lambda_2$, $\Lambda_2$ of rank $l$, and
$\Lambda^\vee/\pi_N\Lambda^\vee=\Lambda_V^\vee/\pi_N\Lambda_V^\vee\times\Lambda_2^\vee/\pi_N\Lambda_2^\vee$. The
second factor has $\auNm(\pi_N)^l=N^{fl}$ elements. In the first, the admissible $\lambda_1$ lie in $\pi_N
L''^\vee\cap\Lambda_V^\vee$ (restriction to $L''$), which contains $\pi_N\Lambda_V^\vee$ with index at most
$[\pi_NL''^\vee:\pi_N\Lambda_V^\vee]=[L''^\vee:\Lambda_V^\vee]=i$ by Lemma~\ref{pbc:disc}(b).
\end{proof}

\begin{lemma}[Sublattices]\label{pbc:zeta}
Let $\Lambda'$ be a lattice of rank $\rho$ and $s>\rho$. Then $\sum_{M'}[\Lambda':M']^{-s}=\prod_{i=0}^{\rho-1}
\zeta_E(s-i)$, the sum over the sublattices $M'$ of finite index. If $s\ge\rho+2\varrho+4$ and $2f\le
2^{2\varrho}$, the sum is at most $2^\rho$.
\end{lemma}
\begin{proof}
A sublattice of finite index is the same as a family $(M'_{\mathfrak p})$ of sublattices of finite index
of the completions, almost all equal to $\Lambda'_{\mathfrak p}$, and the index is the product of the
local indices; so the sum is the product over the primes $\mathfrak p$ of $\auO_E$ of the local sums.
Let $\varpi$ be a uniformiser and $q=\auNm\mathfrak p$; $\Lambda'_{\mathfrak p}\cong\auO_{\mathfrak p}^\rho$.
Each $M'_{\mathfrak p}$ has a unique basis matrix in Hermite form: upper triangular with diagonal
$\varpi^{a_1},\dots,\varpi^{a_\rho}$ and with the entries to the right of the diagonal in row $i$ taken in
a fixed set of representatives of $\auO_{\mathfrak p}/\varpi^{a_i}$. There are $\prod_iq^{a_i(\rho-i)}$ such
matrices for given $a_i$, and the index is $q^{\sum a_i}$; so the local sum is $\prod_{i=1}^\rho(1-q^{
\rho-i-s})^{-1}$, and the product over $\mathfrak p$ is $\prod_{i=0}^{\rho-1}\zeta_E(s-i)$. For the bound,
each factor is at most $\zeta_E(\sigma)$ with $\sigma=2\varrho+5$. Above each rational prime $p$ there are
at most $2f$ primes, of norm at least $p$, so $\zeta_E(\sigma)\le\zeta(\sigma)^{2f}$; and $\zeta(\sigma)\le
1+2^{-\sigma}+\int_2^\infty x^{-\sigma}dx\le1+2^{1-\sigma}$. Hence $\zeta_E(\sigma)\le\exp(2f\,2^{1-\sigma})
=\exp(2f\,2^{-2\varrho-4})\le\exp(1/16)\le2$.
\end{proof}

For $x\in X_0$ put $\|x\|^2=\sum_{\tau}\tau(h(x,x))$, the Euclidean norm of the image of $x$ in $X_{0,
\infty}=\prod_\tau X_0\otimes_{E,\tau}\C\cong\C^{kf}\cong\R^{2kf}$ (coordinates in the orthonormal
basis $x_1,\dots,x_k$). If $x\ne0$ lies in an integral lattice, $h(x,x)$ is a totally positive integer of
$E^+$, so $\|x\|^2\ge f\,(N_{E^+/\Q}h(x,x))^{1/f}\ge1$. Let $D_E$ be the discriminant of $E$ and
$\pbcrd=|D_E|^{1/2f}$ its root discriminant.

\begin{lemma}[Covolume of $\auO_E$]\label{pbc:tb-covol}
Let $\auPl$ be a set of representatives of the $f$ pairs of complex conjugate embeddings of the CM
field $E$ ($[E:\Q]=2f$). The image of $\auO_E$ under $a\mapsto(\tau(a))_{\tau\in\auPl}\in\C^f\cong\R^{2f}$
is a lattice of covolume $2^{-f}|D_E|^{1/2}$ for the Lebesgue measure $\prod_\tau d\Rel\,d\Iml$.
\end{lemma}
\begin{proof}
Let $\alpha_1,\dots,\alpha_{2f}$ be a $\Z$-basis of $\auO_E$ and $A$ the $2f\times2f$ matrix with rows
$(\sigma(\alpha_i))_\sigma$, $\sigma$ running over the $2f$ embeddings, ordered as $\tau_1,\aubar\tau_1,\dots,
\tau_f,\aubar\tau_f$; by definition $D_E=\det(A)^2$. The image lattice has the basis whose $i$-th vector has
coordinates $(\Rel\tau(\alpha_i),\Iml\tau(\alpha_i))_\tau$, so its covolume is $|\det B|$ for the matrix
$B$ of these rows. $B$ is obtained from $A$ by replacing, for each $\tau$, the pair of columns
$(\tau,\aubar\tau)$ by $(\tfrac12(\tau+\aubar\tau),\tfrac1{2i}(\tau-\aubar\tau))$, i.e.\ $B=AM$ with $M$
block diagonal with blocks $\bigl(\begin{smallmatrix}1/2&1/(2i)\\1/2&-1/(2i)\end{smallmatrix}\bigr)$
of determinant $-1/(2i)$. Hence $|\det B|=2^{-f}|\det A|=2^{-f}|D_E|^{1/2}$; in particular $\det B\ne0$
and the image is a lattice. (See also \cite[Chapter~I, (5.1) and Proposition~(5.2)]{Neukirch}.)
\end{proof}

\begin{lemma}[Covolume]\label{pbc:covol}
Let $\Lambda'$ be an integral lattice ($h(\Lambda',\Lambda')\subseteq\auO_E$) of rank $\rho$, spanning a subspace $W\subseteq X_0$. Its image in $W_\infty=\prod_\tau
W\otimes_{E,\tau}\C\cong\R^{2\rho f}$ (with the norm above) is a lattice of covolume $2^{-\rho f}|D_E|^{\rho/2}
\ppidisc\Lambda'$.
\end{lemma}
\begin{proof}
For $\auO_E\subset\C^f$ (embedded by the $\tau\in\auPl$) the covolume for the Lebesgue measure is
$2^{-f}|D_E|^{1/2}$ (Lemma~\ref{pbc:tb-covol}). Let $w_1,\dots,w_\rho\in\Lambda'$ be an $E$-basis of $W$ with Gram
matrix $G_w$ and $F=\bigoplus\auO_Ew_i\subseteq\Lambda'$, an integral lattice of finite index in $\Lambda'$. The map $\C^{\rho f}\to W_\infty$ given at $\tau$ by $\tau(w)$ (in
orthonormal coordinates) has real Jacobian $\prod_\tau|\det\tau(w)|^2=\prod_\tau\tau(\det G_w)=\auNm(\det
G_w)^{1/2}=\ppidisc F$ (Lemma~\ref{pbc:disc}(a); $\det G_w\in E^+$ is totally positive). So the claim holds
for $F$. Since $F\subseteq\Lambda'$ has finite index, $\operatorname{covol}(F)=[\Lambda':F]\operatorname{covol}(\Lambda')$
and, by Lemma~\ref{pbc:disc}(a), $\ppidisc F=[\Lambda':F]\ppidisc\Lambda'$; so the claim holds for $\Lambda'$.
\end{proof}

\begin{lemma}[Siegel's lemma]\label{pbc:siegel}
Let $\Lambda$ be a unimodular lattice of $X_0$ and $V\subseteq X_0$ a subspace of codimension $l\ge1$. There
are $E$-linearly independent $X_1,\dots,X_l\in\Lambda\cap V^\perp$ with $V=\{v:h(X_i,v)=0\ \forall i\}$ and
\[
\prod_{i=1}^l\|X_i\|^{2f}\le(lf\,\pbcrd)^{lf}\,\ppidisc(\Lambda\cap V).
\]
\end{lemma}
\begin{proof}
Let $\Lambda'=\Lambda\cap V^\perp$, of rank $l$ and integral as $\Lambda$ is, viewed as a lattice of rank $D=2lf$ over $\Z$ in $V^\perp_\infty
\cong\R^D$. By Lemmas~\ref{pbc:covol} and~\ref{pbc:disc}(c) its covolume is $2^{-lf}|D_E|^{l/2}\ppidisc
(\Lambda\cap V)$. Let $\lambda_1\le\dots\le\lambda_D$ be its successive minima for the Euclidean unit ball, of
volume $\pi^{lf}/(lf)!$. Minkowski's second theorem (\cite[Chapter~VIII, \S4.3, Theorem~V]{Cassels}; see also \cite[Theorem~1.3]{HenkMinima}) gives
\[
\lambda_1\cdots\lambda_D\le\frac{2^D\,2^{-lf}|D_E|^{l/2}\ppidisc(\Lambda\cap V)}{\pi^{lf}/(lf)!}\le(lf)^{lf}
\pbcrd^{\,lf}\ppidisc(\Lambda\cap V),
\]
as $|D_E|^{1/2}=\pbcrd^{\,f}$ and $(2/\pi)^{lf}\le1$. Let $v_1,\dots,v_D$ be $\Z$-independent with $\|v_j\|=
\lambda_j$. An $E$-subspace of dimension $e$ contains at most $2fe$ $\Z$-independent vectors, so
$v_1,\dots,v_{2f(i-1)+1}$ span an $E$-subspace of dimension at least $i$; choose inductively
$E$-independent $X_i$ among them with $\|X_i\|\le\lambda_{2f(i-1)+1}$. Then $\|X_i\|^{2f}\le\prod_{t=1}^{2f}
\lambda_{2f(i-1)+t}$, and the product over $i$ is at most $\lambda_1\cdots\lambda_D$. The $X_i$ span
$V^\perp$, so their orthogonal is $V$.
\end{proof}

\begin{lemma}[Height zeta function]\label{pbc:height}
Let $\Lambda$ be a unimodular lattice of $X_0$ and $0\le l\le k$. Then $\sum_V\ppidisc(\Lambda\cap V)^{-(k+1)}
\le(25\,lf\,\pbcrd)^{f(k+1)l}$ (read as $1$ for $l=0$), the sum over the subspaces $V\subseteq X_0$ of
codimension $l$.
\end{lemma}
\begin{proof}
For $l=0$ the only term is $V=X_0$, with $\ppidisc\Lambda=1$. Let $l\ge1$. Attach to each $V$ the tuple
$(X_1,\dots,X_l)$ of Lemma~\ref{pbc:siegel}; $V$ is determined by it. Raising the inequality of that
lemma to the power $k+1$,
\[
\ppidisc(\Lambda\cap V)^{-(k+1)}\le(lf\pbcrd)^{lf(k+1)}\prod_i\|X_i\|^{-2f(k+1)} .
\]
Summing over all tuples of nonzero vectors of $\Lambda$, the total is at most $(lf\pbcrd)^{lf(k+1)}
\bigl(\sum_{0\ne X\in\Lambda}\|X\|^{-2f(k+1)}\bigr)^l$. The image of $\Lambda\smallsetminus0$ in $\R^{2kf}$
satisfies the hypotheses of Lemma~\ref{pbc:pack}(b) (differences of lattice vectors are lattice vectors),
and $2f(k+1)\ge2kf+2$; so the inner sum is at most $5^{2f(k+1)}$.
\end{proof}

\begin{theorem}[(C2)]\label{pbc:C2}
Assume the size conditions of Remark~\ref{pbc:rem-sizes}. Let $0\le l\le n^3$ and $r=k-l$. Then
\[
\sum_{c\in\mathcal C}\ \sum_{L'}\ (\ppidisc L')^{-k}\le\frac{\#\mathcal C}{N^{kf}}\,\pbet^{\,l+1},
\]
the inner sum over the lattices $L'\subseteq\ppiLc$ of rank $r$ with $\ppidisc L'\le b^{n^3f}$.
\end{theorem}
\begin{proof}
Put $\mathsf D=b^{n^3f}$ and let $F(c)$ be the inner sum. Let $\mathcal A=K^\circ/K_H$ and fix
representatives $y_\phi\in K^\circ$ of the cosets $\phi$. By Lemma~\ref{pbc:average} it suffices to show,
for every $g\in H(\auA_f)$,
\begin{equation}\label{pbc:eq-goal}
\sum_{\phi\in\mathcal A}F([gy_\phi])\le\frac{\pbet^{\,l+1}}2\le\#\mathcal A\cdot\frac{\pbet^{\,l+1}}{N^{kf}},
\end{equation}
the second inequality being Lemma~\ref{pbc:index} ($4f\le2^{2\varrho+1}\le5^k$). Fix $g$ and put $\Lambda=g
\auO_E^k$, a unimodular lattice of $X_0$. For $\phi\in\mathcal A$ let $c=[gy_\phi]$; then $h_c=\eta g y_\phi j$
with $\eta\in H(E^+)$ and $j\in J_X\subseteq K_H$, so $\ppiLc=h_cL_1=\eta gy_\phi L_1$. By
Lemma~\ref{pbc:index}, $\eta^{-1}\ppiLc=\{v\in\Lambda:\lambda_\phi(v)\in\pi_N\widehat\auO\}$, where
$\lambda_\phi=\lambda_{y_\phi}\circ g^{-1}\in\Lambda^\vee$, and the classes of the $\lambda_\phi$ modulo $\pi_N
\Lambda^\vee$ are distinct. As $L'\mapsto\eta^{-1}L'$ preserves rank and discriminant
(Lemma~\ref{pbc:disc}(a)),
\[
\sum_\phi F([gy_\phi])\le\sum_{L''}\#\{\phi:\lambda_\phi(L'')\subseteq\pi_N\auO_E\}\,(\ppidisc L'')^{-k},
\]
the sum over the lattices $L''\subseteq\Lambda$ of rank $r$ with $\ppidisc L''\le\mathsf D$. For such $L''$ let
$V=EL''$, $d_V=\ppidisc(\Lambda\cap V)\ge1$ and $i=[\Lambda\cap V:L'']\ge1$, so that $\ppidisc L''=d_Vi\le\mathsf
D$. By Lemma~\ref{pbc:forms} the count is at most $iN^{fl}$. With $\mu=2\varrho+4$,
\[
i\,(d_Vi)^{-k}=i(d_Vi)^{\mu+1}d_V^{-(k+\mu+1)}i^{-(k+\mu+1)}\le\mathsf D^{\mu+2}\,d_V^{-(k+1)}\,i^{-(r+\mu)},
\]
using $i\le\mathsf D$, $d_Vi\le\mathsf D$, $d_V,i\ge1$ and $r\le k$. Group the $L''$ by $V$ (a subspace of
codimension $l$) and sum over the sublattices of finite index of $\Lambda\cap V$ by Lemma~\ref{pbc:zeta}
($2f\le2^{2\varrho}$), then over $V$ by Lemma~\ref{pbc:height}:
\[
\sum_\phi F([gy_\phi])\le N^{fl}\mathsf D^{2\varrho+6}\,2^r\,(25\,lf\,\pbcrd)^{f(k+1)l}
\le\bigl(2^{k}\mathsf D^{2\varrho+6}\bigr)\bigl(N^f(25n^3f\pbcrd)^{f(k+1)}\bigr)^l .
\]
By the size conditions, $2^{k+1}\mathsf D^{2\varrho+6}\le b^{(k+1)/8+n^3(2n+4)f}\le\pbet$ and $N^f(25n^3f
\pbcrd)^{f(k+1)}\le b^{60nf+f(k+1)/1000}\le\pbet$. This gives \eqref{pbc:eq-goal}.
\end{proof}

\begin{remark}[Size conditions]\label{pbc:rem-sizes}
Theorem~\ref{pbc:C2} uses: (Z1) $\log b\ge8\log2$, $N\ge5$; (Z2) $k\ge8n^3(2n+4)+1$, $k\ge2\varrho$,
and $60n+(k+1)/1000\le k/4$; (Z3) $N\le b^{60n}$; (Z4) $2f\le2^{2\varrho}$; (Z5) $(25n^3f\pbcrd)^{1000}\le
b$; and the description of $K_H$ at the primes $w\mid\pi_N$ in the proof of Lemma~\ref{pb:KR}.
Theorem~\ref{pbc:C1} uses $\varrho\ge280$, $M\le b^{23}$, $N\le b^{60n}$, $1\le|\kappa_\tau|\le b^{23}$, $n\le b$.
All follow from {\cite[Lemma~13.4]{P46}}: (Z1) from (S1) and $\log N\ge m\log b$ (S5); (Z2), (Z3),
$\varrho\ge280$ from (S7) ($k\ge512000n^9$, $N\le b^{60n}$, $\varrho\ge10^5$); $M\le b^{22}$ from (S5);
the bound on $\kappa_\tau$ is \eqref{pbe:eq-kappa}; $n\le b$ from $18\varrho f\le b^8$ (S7) or (S1). For (Z4) and (Z5): by the
proof of (S2), $\sqrt x\le m\le2\sqrt x$, so $n=m/2\le\sqrt x$ and $\varrho\ge\sqrt x/2-1$; by (S3),
(S4), $f\le x^{12}$ and $\pbcrd\le x^{12}$. Since $\log x\ge10^{14}$ (S1), $2^{2\varrho}\ge2^{\sqrt x-2}\ge2x^{12}
\ge2f$, and $25n^3f\pbcrd\le x^{26}$, while $x^{26000}\le x^{10^6}\le x_1\le b$ (S1). (Here $f$ is the
number of places $\tau$, as in the rest of the paper.)
\end{remark}



% ===== B-size.tex =====
% proofs/B-size.tex
% Statement (S) (size) of Theorem ppi:pillars of proofs/B-pillars.tex: the bound of one block of
% amplitudes by the sizes of its near orbits. Labels of this file start with pbs:.
\providecommand{\auO}{\mathcal O}
\providecommand{\auNm}{\operatorname{Nm}}
\providecommand{\auA}{\mathbb A}
\providecommand{\aubar}[1]{\overline{#1}}
\providecommand{\auf}[1]{\texttt{#1}}
\providecommand{\panot}[1]{\par\smallskip{\raggedright\noindent\textbf{[Not proved here:} #1\textbf{]}\par}\smallskip}
\providecommand{\pbR}{\overline R_0}
\providecommand{\pbDid}{\mathcal D_{\mathrm{id}}}
\providecommand{\pbDmain}{\mathcal D_{\mathrm{main}}}
\providecommand{\pbvol}{\operatorname{vol}}
\providecommand{\pbdeg}{\operatorname{deg}}
\providecommand{\pbg}{\mathsf g}
\providecommand{\pbeN}{\mathcal N}
\providecommand{\pbet}{\beta_*}
\providecommand{\pbenp}{\mathsf n_{\mathrm p}}
\providecommand{\pbesz}{\operatorname{size}}
\providecommand{\pbeGs}{\Gamma^*}
\providecommand{\pbeAp}{A^{+}}
\providecommand{\pbel}{\ell}
\providecommand{\pbeL}{\mathsf L}
% macros of this file (prefix \pbs)
\providecommand{\pbsv}{\mathsf v}
\providecommand{\pbsK}{\mathsf K}
\providecommand{\pbsF}{\mathcal F}
\providecommand{\pbsO}{\mathcal O}
\providecommand{\pbsrb}{\rho_\square}
\providecommand{\pbsth}{\theta_\square}
\providecommand{\pbsp}{\mathsf p}
\providecommand{\pbsi}{\mathsf i}
\providecommand{\pbsmn}{\mathsf m}

\subsection{The size of a block}\label{pbs:sec}

This section proves statement \textup{(S)} of Theorem~\ref{ppi:pillars}. We use the notation of
\S\S\ref{pb:sec}, \ref{pbe:sec} and~\ref{ppi:sec}: $n=\varrho+1$, the frame $f_0,\dots,f_\varrho$
($f_a=e_{a+1}$ for $a<\varrho$, $f_\varrho=z$), $A=\langle f_0,\dots,f_{\varrho-1}\rangle$, the matrix
$Z(\gamma)=(h(f_a,\gamma f_b))_{a,b}$ with leading minors $\delta_l(\gamma)$, the discriminant
$\pbg(\gamma)=\delta_n\aubar{\delta_n}-N^2\delta_\varrho\aubar{\delta_\varrho}$, the groups $U_c$, $U^0$,
$U_c^0=U_c\cap U^0$, the cell rank $j(\gamma)$ (Definition~\ref{ppi:def-cell}), the gauge $\mathsf
G(\gamma)$ \eqref{ppi:eq-gauge}, near and chosen elements, agreeing reductions, the set $\pbeN(c,c')$
(Definition~\ref{ppi:def-N}) and $\pbesz(\gamma)=\pbet^2(\mathcal D(\gamma)/(\pbenp\mathsf
G(\gamma)^k))^{1/2}$. At a real place $\tau$ we use the normalised basis, the norm $\|\cdot\|$, the
scales $s_a$ \eqref{pbe:eq-scale}, the length $\pbel_\tau$ and the cut of \S\ref{pbe:ss-real}.
In this section, as in \S\S\ref{pbg:sec}, \ref{ppi:ss-cell} and \ref{pbc:sec}, hermitian forms are
written conjugate-linear in the first variable and linear in the second, the opposite of convention
(Cv1); $Z(\gamma)$ is the matrix of \S\ref{ppi:ss-statement} in this reading.
Put
\[
\pbsrb=b^{37n},\qquad\pbsth=b^{6n+2},\qquad\mathsf j'=\min(\varrho,j),
\]
the radius of the boxes and the radius of the cut.

\emph{What is proved.}
\begin{itemize}
\item[(1)] The exact statement is Theorem~\ref{pbs:S}. Its right-hand side is that of
Theorem~\ref{ppi:pillars}(S). The constant in front of the sum is $\nu=\pbvol(U_c\backslash
U_\infty)$ for the Haar measure $du$ of convention (Cv3) of \S\ref{pb:ss-setting}; its explicit value
$\nu_{\mathrm{src}}=\nu_*$ is Lemma~\ref{pb:covol}, used in Lemma~\ref{pbs:open-arith}
(Remark~\ref{pbs:rem-nu}).
\item[(2)] The reduction of (S) to one orbit of one cell rank, the vanishing for non-agreeing $s$,
and the choice of representatives are \textbf{proved in full} (\S\ref{pbs:ss-reduction}).
\item[(3)] The bound for the orbits that are summed by Schur's test (all orbits except those on the
open cell with large discriminant) is \textbf{proved in full} (\S\S\ref{pbs:ss-schur}--%
\ref{pbs:ss-orbit}): the weighted Cauchy--Schwarz inequality, the computation of the weight at one
real place by a triangular basis, the Ostrowski--Taussky inequality and Hua's integral, the
colouring of an orbit by a lattice index, packing, and the final comparison of sizes.
\item[(4)] For the orbits on the open cell the bound is deduced (\S\ref{pbs:ss-open}) from two
statements: the open-cell pairing bound (O1) for the Bessel functions, stated exactly and
proved in \S\ref{pbo:sec}, and the comparison (O2) of the Siegel constant with the main second
moment, \textbf{proved here} (Lemma~\ref{pbs:O2}) from the mass formula of
Lemma~\ref{pb:massratio}.
\item[(5)] The size conditions on the data that are used are listed in \S\ref{pbs:ss-sizes};
all are in {\cite[Lemma~13.4]{P46}}.
\end{itemize}

\subsubsection{The statement}\label{pbs:ss-statement}

For a class $c$ let $\Lambda_c=h_cL\cap X_0$, a lattice of the positive definite space $X_0$
($\dim X_0=k$). For $\gamma\in\Gamma_{c,s,c'}$ let $k_\gamma=h_c^{-1}\gamma h_{c'}$, and for $0\le
j\le n$ let $i_\gamma=\pbsi_j(k_\gamma)$, $i'_\gamma=\pbsi_j(k_\gamma^{-1})$ be the indices of
\eqref{ppi:eq-gauge}: $\pbsi_j(k)$ is the index, in the $\widehat\auO$-module spanned by the
vectors $\xi_a(k)=(h(f_a,kf_b))_{0\le b\le\varrho}\in\auA_{E,f}^{\,n}$, $0\le a<\varrho$, of the
submodule spanned by the $\xi_a(k)$ with $a<\mathsf j'$. (These modules are finitely generated
$\widehat\auO$-modules of the same rank by Lemma~\ref{pbs:indep} below, so the index is finite; it
is $1$ for $j=n$.) For a cell rank $j$ put
\[
\pbsmn_j(\gamma)=\frac{|\auNm\delta_j(\gamma)|}{N^{2(j-\varrho)_+f}},\qquad
\mathsf G_j(\gamma)=\frac{\max\bigl(\auNm\pbg(\gamma),\ N^f\pbsmn_j(\gamma)\bigr)}{i_\gamma
i'_\gamma},
\]
where $\auNm x=\prod_{\tau\in\auPl}|\tau(x)|$ for $x\in E^+$ and $\auNm x=\prod_\tau|\tau(x)|^2$ for
$x\in E$. Then the gauge \eqref{ppi:eq-gauge} is $\mathsf G(\gamma)=M^{2\mathsf j'f}\mathsf
G_{j(\gamma)}(\gamma)$, and for $\gamma$ of cell rank $j$
\[
\pbesz(\gamma)=\mathsf S_j\bigl(\mathsf G_j(\gamma)\bigr),\qquad\mathsf S_j(x):=\pbet^2\Bigl(
\frac{\mathcal D_j}{\pbenp\,(M^{2\mathsf j'f}x)^k}\Bigr)^{1/2},\qquad\mathcal D_j=\begin{cases}
\pbDmain&j=n,\\ \pbet^{-16}\pbDid&j\le\varrho,\end{cases}
\]
since $\mathcal D(\gamma)=\pbDmain$ exactly when $\pbg(\gamma)\ne0$, that is when $j(\gamma)=n$
(Proposition~\ref{ppi:cell}).

\begin{theorem}[Size]\label{pbs:S}
Assume the size conditions of \S\ref{pbs:ss-sizes}, Theorem~\ref{pbe:ds}, and the statement
\textup{(O1)} of \S\ref{pbs:ss-open} (Theorem~\ref{pbo:O1}). Let $W^\flat$ be as in
Theorem~\ref{pbe:bessel}. For every stable $T\subseteq\pbeGs$ and every block $(c,s,c')$,
\[
\bigl\|A^\flat_T(c,s,c')\bigr\|\le\nu\sum_{\gamma\in\pbeN(c,c')\cap T\cap\Gamma_{c,s,c'}}
\pbesz(\gamma),
\]
and $A^\flat_T(c,s,c')=0$ if $s$ does not agree. This is \textup{(S)} of
Theorem~\ref{ppi:pillars}.
\end{theorem}

The proof occupies the rest of the section; it is assembled in \S\ref{pbs:ss-proof}.

\subsubsection{Measures, cutoffs and reduced points}\label{pbs:ss-measures}

\paragraph*{Measures} On $G_\infty=\prod_\tau G_\tau$ we use the product $dg$ of the Haar
measures of \S\ref{pbe:ss-real}. At each $\tau$ the group $U_\tau$ is the group of real points of a
split unipotent group; on it $du_\tau$ is the Haar measure of convention (Cv3) of
\S\ref{pb:ss-setting}, on $U_\infty$ the product $du$, and on $U_\tau\backslash G_\tau$ the quotient
measure (Remark~\ref{pbs:rem-nu}). On the subgroup $U^0_\tau$ of the elements of $U_\tau$ that fix
every $\tilde f_a$, also a unipotent group, we fix a Haar measure $du^0$. Unipotent groups are
unimodular. Recall $\nu=\pbvol(U_c\backslash U_\infty)$, finite and independent of $c$
(Lemma~\ref{pb:ortho}).

A \emph{cutoff} for a discrete group $\Delta$ acting on the left on $G_\infty$ is a measurable
$p\colon G_\infty\to[0,\infty)$ with $\sum_{\delta\in\Delta}p(\delta y)=1$ for all $y$. For a
discrete subgroup of $G_\infty$ the indicator of a measurable set of representatives of
$\Delta\backslash G_\infty$ is one. If $p$ is a cutoff for $\Delta$ and $F\ge0$ is measurable and
left invariant under $\Delta$, then $\int_{G_\infty}pF\,dg=\int_{\Delta\backslash G_\infty}F$; in
particular the left side does not depend on the cutoff.

\paragraph*{Reduced points}
Let $\mathsf D=\pbcD\subseteq\prod_\tau\C$ be the Borel set of Lemma~\ref{pbc:domain}: it contains
exactly one representative of each coset of the lattice $M\auO_E$ (embedded by $x\mapsto(\tau(x))_\tau$),
and all coordinates of its points have absolute value at most $b^{26}$. For $y\in G_\infty$ and each $\tau$ let $w_a=y_\tau^{-1}\tilde f_a$,
let $o_a$ be their Gram--Schmidt orthogonalisation and $w_a=o_a+\sum_{a'<a}t_{a'a}o_{a'}$ as in
Lemma~\ref{pbe:len}(c). The point $y$ is \emph{reduced} if for all $a'<a$ the point
$(t_{a'a}(y_\tau)s_{a'}/s_a)_\tau$ lies in $\mathsf D$; let $\omega$ be the indicator of the set of
reduced points. A point is \emph{boxed} if it is reduced and $\pbel_\tau(y_\tau)\le\pbsth$ for all
$\tau$. (These are the reduced and boxed points of Definition~\ref{pbc:def-reduced}.)

\begin{lemma}[Reduced points]\label{pbs:reduced}
\begin{enumerate}
\item[(a)] $\omega(uy)=\omega(y)$ for every $u=(u_\tau)$ with $u_\tau$ in the stabiliser $F_\tau$
of $\tilde f_0,\dots,\tilde f_\varrho$ in $G_\tau$.
\item[(b)] For every $y$, the set of $x\in U_c$ with $xy$ reduced is a single coset $U_c^0x_0$.
Hence $\sum_{x\in U_c}a(x)\omega(xy)=1$ for every function $a\colon U_c\to[0,\infty]$ whose sum
over each coset $U^0_cx$ is $1$.
\item[(c)] If $y$ is reduced and $\pbel_\tau(y_\tau)\le\pbsth R$ with $R\ge1$, then
$\|y_\tau^{-1}\tilde f_a\|\le\pbsrb R$ for all $a$.
\end{enumerate}
\end{lemma}
\begin{proof}
(a) $(u_\tau y_\tau)^{-1}\tilde f_a=y_\tau^{-1}\tilde f_a$. (b) For $x\in U_c$ write $xf_a=f_a+\sum_{a'
<a}C_{a'a}f_{a'}$; then $C_{a'a}\in M\auO_E$, because $U_c$ preserves $h_cL$, acts trivially on it
modulo $M$, and the $f_a$ are part of a basis of $h_cL$ adapted to the frame (proof of
Lemma~\ref{pbe:eichler}). As in Step~1 of the proof of Proposition~\ref{pbe:count}, replacing $y$
by $xy$ does not change the $o_a$ and replaces $t_{a'a}s_{a'}/s_a$ by $t_{a'a}s_{a'}/s_a+\tau(
C_{a'a})+\sum_{a'<a''<a}t_{a'a''}(s_{a'}/s_{a''})\tau(C_{a''a})$ (for $x^{-1}$ in place of $x$, which
is the same statement). By Lemma~\ref{pbe:eichler} every strictly upper triangular matrix over
$M\auO_E$ occurs as $(C_{a'a})$; choosing its entries by induction on $a-a'$ we reach a reduced
point, and at each step the choice is unique because $\mathsf D$ is a set of representatives. So
the $x$ with $xy$ reduced are those with one prescribed matrix $(C_{a'a})$, and two elements of
$U_c$ with the same matrix differ by an element of $U_c$ fixing the frame, that is of $U_c^0$. For
the second claim, the sum is the sum of $a$ over the coset $U^0_cx_0$. (c) The coefficients of a
reduced point satisfy $|t_{a'a}|\le|t_{a'a}s_{a'}/s_a|\le b^{26}$ by \eqref{pbe:eq-scale}, and
$\|o_a\|<\pbel_\tau(y_\tau)$ by Lemma~\ref{pbe:len}(c), so $\|w_a\|\le n\,b^{26}\pbsth R\le\pbsrb R$.
\end{proof}

\paragraph*{Class cutoffs} For each class $c$ fix a cutoff $p_c$ for $U_c^0$ and put
$\beta_c=p_c\,\omega$. By Lemma~\ref{pbs:reduced}(b) applied to $a(x)=p_c(xy)$, $\beta_c$ is a
cutoff for $U_c$:
\begin{equation}\label{pbs:eq-cutoff}
\sum_{x\in U_c}\beta_c(xy)=1\qquad(y\in G_\infty).
\end{equation}
Hence, for $F\ge0$ measurable and left invariant under $U_\infty$,
\begin{equation}\label{pbs:eq-nu}
\int_{G_\infty}\beta_c\,F\,dg=\int_{U_c\backslash G_\infty}F=\nu\int_{U_\infty\backslash G_\infty}F .
\end{equation}

\begin{remark}[The constant $\nu$]\label{pbs:rem-nu}
Put $\nu_{\mathrm{src}}=(M^{2d+\varrho}/N^k)^f\pbvol(\widehat\auO)^{-d}|D_E|^{\varrho/4}/\Lambda$
with $d=\sum_{a<\varrho}(2a+k+1)=\varrho(\varrho+k)$ and $\Lambda=\prod_\tau\prod_{a\le\varrho}
(\kappa_\tau^{2(\varrho-a)}/N)^{k+2(\varrho-a)}$; this is the constant $\nu_*$ of Lemma~\ref{pb:covol}
($M^{f(2d+\varrho)}=\#\pbUb$, $|D_E|^{\varrho/4}=\aurd^{f\varrho/2}$), so $\nu=\nu_{\mathrm{src}}$ for the
measure (Cv3) by that lemma. The statement of Theorem~\ref{pbs:S} itself does not depend on the
choice of $du$: if $du$ is multiplied by $\lambda>0$, the quotient measure on $U_\infty\backslash G_\infty$ is
divided by $\lambda$, the normalisation $\int_{U_\tau\backslash G_\tau}\|W_\tau\|^2=d_\tau$ multiplies
$W$ and $W^\flat$ by $\lambda^{1/2}$ (one factor $\lambda^{1/2f}$ at each place), so $A^\flat_T$ is
multiplied by $\lambda$, and so is $\nu$, while $\pbesz(\gamma)$ does not change. Its proof, however,
uses (O1), which is proved for the measure (Cv3), and the explicit value of $\nu$
(Lemmas~\ref{pbs:O2}, \ref{pbs:open-arith}).
\end{remark}

\paragraph*{Amplitudes with weights}
For $\gamma\in\Gamma$ and $y\in G_\infty$ put $\pbsK^\flat_\gamma(y)=W^\flat(y)W^\flat(\gamma^{-1}y)
^*$ and $\pbsK_\gamma(y)=W(y)W(\gamma^{-1}y)^*$, matrices indexed by $I\times I$ with the Euclidean
norm; $\|\pbsK^\flat_\gamma(y)\|=\|W^\flat(y)\|\,\|W^\flat(\gamma^{-1}y)\|$. For $u\in U_c$, $u'\in
U_{c'}$, $\pbsK^\flat_{u\gamma u'^{-1}}(uy)=\pbsK^\flat_\gamma(y)$, because $W^\flat(uy)=\psi_\infty(u)
W^\flat(y)$ and $\psi_\infty$ is trivial on $U_c$ and on $U_{c'}$ (\S\ref{pb:ss-layer2}).

\begin{lemma}[Change of weight]\label{pbs:weight}
Let $T$ be stable, $X=T\cap\Gamma_{c,s,c'}$, and let $\phi\colon X\times G_\infty\to[0,\infty)$ be
measurable in $y$ with
\[
\sum_{(u,u')\in U_c\times U_{c'}}\phi(u\gamma u'^{-1},uy)=1\qquad(\gamma\in X,\ y\in G_\infty).
\]
Then $\sum_{\gamma\in X}\int_{G_\infty}\phi(\gamma,y)\|\pbsK_\gamma(y)\|\,dy<\infty$,
\[
A^\flat_T(c,s,c')=\sum_{\gamma\in X}\int_{G_\infty}\phi(\gamma,y)\,\pbsK^\flat_\gamma(y)\,dy,\qquad
\sum_{\gamma\in X}\int_{G_\infty}\phi(\gamma,y)\|\pbsK_\gamma(y)-\pbsK^\flat_\gamma(y)\|\,dy\le
\frac\nu2\,b^{-k^3f}.
\]
\end{lemma}
\begin{proof}
The group $\Delta=U_c\times U_{c'}$ acts on $X$ by $(u,u')\gamma=u\gamma u'^{-1}$ ($X$ is stable
under it because $T$ is stable and $\Gamma_{c,s,c'}$ is stable under $U_c\times U_{c'}$). The
stabiliser of $\gamma$ is $\{(u,\gamma^{-1}u\gamma):u\in\Lambda_\gamma\}$ with $\Lambda_\gamma=U_c\cap
\gamma U_{c'}\gamma^{-1}$. Let $F\colon X\times G_\infty\to[0,\infty]$ be measurable in $y$ with
$F(u\gamma u'^{-1},uy)=F(\gamma,y)$. For an orbit $\Delta\gamma$,
\[
\sum_{\delta\in\Delta\gamma}\int\phi(\delta,y)F(\delta,y)\,dy=\sum_{(u,u')\in\Delta/\mathrm{Stab}}
\int\phi(u\gamma u'^{-1},uy)F(\gamma,y)\,dy=\int_{\Lambda_\gamma\backslash G_\infty}F(\gamma,y)\,dy,
\]
by the substitution $y\mapsto uy$ and the cutoff property (the function $y\mapsto\sum_{\Delta/
\mathrm{Stab}}\phi(u\gamma u'^{-1},uy)$ is a cutoff for $\Lambda_\gamma$, since summing it over
$\Lambda_\gamma$ gives the full sum over $\Delta$). So the sum over $X$ of $\int\phi F$ is the same
for all admissible $\phi$, namely $\sum_{\Delta\gamma}\int_{\Lambda_\gamma\backslash G_\infty}F$. With
$F=\|\pbsK_\gamma\|$ this is $\sum_{\Delta\gamma}\int_{\Lambda_\gamma\backslash G_\infty}\|W(y)W(\gamma^{-1}y)^*\|\,dy$,
the sum of the integrals of the norms in \eqref{pb:eq-amp} (with $W^\flat=W$), which is finite because
Proposition~\ref{pb:unfold} states that the sums \eqref{pb:eq-amp} for $W$ converge absolutely for every
stable $T$; and then the same computation without absolute values (justified by Fubini) gives the
expression \eqref{pb:eq-amp} for $A^\flat_T$, using $\|W^\flat\|\le\|W\|$. For the last claim take
$\phi(\gamma,y)=\beta_c(y)\beta_{c'}(\gamma^{-1}y)$, which is admissible by \eqref{pbs:eq-cutoff};
then the sum is at most the left side of the first display in the proof of
Proposition~\ref{pbe:cut} (for the block $(c,s,c')$ only), and the proof there gives the bound
$\nu\,b^{-2r^3f+12n^3rf+400r^2f+2kf+1}\le\frac\nu2b^{-k^3f}$; the last inequality is
$12n^3r+400r^2+2k+2+k^3\le2r^3$. Indeed $2r^3-k^3\ge r^3\ge(12n^3+500)r^2\ge12n^3r+400r^2+2k+2$,
as $k<r$ and $12n^3+500\le k<r$ ({\cite[Lemma~13.4(S7)(a)]{P46}}).
\end{proof}

\begin{lemma}[Non-agreeing reductions]\label{pbs:agree}
If $T$ is stable and $s$ does not agree, then $A^\flat_T(c,s,c')=0$.
\end{lemma}
\begin{proof}
Let $u,u'$ be integral elements of $U(\auA_f)$ (in the stabiliser of the lattice) with reductions
$\bar u s=s\bar u'$ and $\psi_f(u)\ne\psi_f(u')$. By strong approximation for the split unipotent
group $U$ there are $x,x'\in U(E^+)$ with $h_c^{-1}xh_c\in u\,(K_U\cap J)$ and $h_{c'}^{-1}x'h_{c'}\in
u'(K_U\cap J)$, where $K_U$ is the integral subgroup of $U(\auA_f)$; as $\psi_f$ is trivial on $U(\auA_f)
\cap J$ and invariant under conjugation by $h_c$, $\psi_f(x)=\psi_f(u)$ and $\psi_f(x')=\psi_f(u')$.
The map $\gamma\mapsto x\gamma x'^{-1}$ is a bijection of $T\cap\Gamma_{c,s,c'}$ onto $T\cap\Gamma_{c,
\bar us\bar u'^{-1},c'}=T\cap\Gamma_{c,s,c'}$ ($T$ is stable). By the substitution $y\mapsto xy$ in
Lemma~\ref{pbs:weight} with the weight $\phi(\delta,y)=\beta_c(xy)\beta_{c'}(x'\delta^{-1}
y)$ for $\delta=x^{-1}\gamma x'$ (admissible: $x$ normalises $U_c=U(E^+)\cap h_cJh_c^{-1}$ because
$h_c^{-1}xh_c\in K$ and $J$ is normal in $K$, so $\sum_{u\in U_c}\beta_c(xuy)=\sum_{u\in U_c}\beta_c(u\,
xy)=1$, and likewise for $x'$), and $W^\flat(xy)=\psi_\infty(x)W^\flat(y)$,
we get $A^\flat_T(c,s,c')=\psi_\infty(x)\aubar{\psi_\infty(x')}A^\flat_T(c,s,c')=\aubar{\psi_f(u)}
\psi_f(u')A^\flat_T(c,s,c')$, using $\psi_\infty\psi_f=1$ on $U(E^+)$. As $\psi_f(u)\ne\psi_f(u')$, the
block vanishes.
\end{proof}

\subsubsection{Reduction to one orbit}\label{pbs:ss-reduction}

For $\gamma\in\Gamma$ and classes $c,c'$ the \emph{amplitude majorant} is
\[
\mathsf a_{c,c'}(\gamma)=\int_{G_\infty}\beta_c(y)\,\beta_{c'}(\gamma^{-1}y)\,\|W^\flat(y)\|\,
\|W^\flat(\gamma^{-1}y)\|\,dy\in[0,\infty].
\]
By Lemma~\ref{pbs:weight} with $\phi(\gamma,y)=\beta_c(y)\beta_{c'}(\gamma^{-1}y)$,
\begin{equation}\label{pbs:eq-majorant}
\|A^\flat_T(c,s,c')\|\le\sum_{\gamma\in T\cap\Gamma_{c,s,c'}}\mathsf a_{c,c'}(\gamma).
\end{equation}

We say that $\gamma$ \emph{has cell rank $j$} ($0\le j\le n$) if (i) there is $\zeta\in E$ with
$\zeta\aubar\zeta=1$ such that $(\zeta,1)$ occurs at every $a$ with $j\le a\le\varrho$
(\S\ref{ppi:ss-statement}); (ii) $\delta_j(\gamma)\ne0$ ($\delta_0=1$); (iii) $j=n$ if and only if
$\pbg(\gamma)\ne0$. We never need that such a $j$ is unique; we always use $j=j(\gamma)$ of
Definition~\ref{ppi:def-cell}, by the following consequence of Proposition~\ref{ppi:cell}.

\begin{lemma}[Cell ranks on $\pbeGs$]\label{pbs:cells}
Every $\gamma\in\pbeGs$ has cell rank $j(\gamma)$, and $j(\gamma)\ge1$. The cell rank, the minors
$\delta_l$ and $\pbg$ are invariant under $\gamma\mapsto u\gamma u'$ for $u,u'\in U(E^+)$, and the
matrix $Z(\gamma)$ is invariant under $\gamma\mapsto u\gamma u'$ for $u,u'\in U^0(E^+)$. The sets
$\{\gamma\in T:j(\gamma)=j\}$ and $\{\gamma\in T:\auNm\pbg(\gamma)>N^f\pbsmn_j(\gamma)\}$ are stable
when $T$ is.
\end{lemma}
\begin{proof}
For $\gamma\in\pbeGs$, (i) holds at $j(\gamma)$ by definition, and (ii), (iii) by
Proposition~\ref{ppi:cell}. If $j(\gamma)=0$, then $\pbg(\gamma)=0$ and, for each $a$, there is a
common vector $fx^{(a)}=\gamma fy^{(a)}$ with $x^{(a)},y^{(a)}$ of top index $a$ and $y^{(a)}_a=1$;
by induction on $a$, $\langle f_0,\dots,f_a\rangle=\gamma\langle f_0,\dots,f_a\rangle$ (both sides have
dimension $a+1$ and the left contains $fx^{(a)}$ and the right side for $a-1$). For $a=\varrho-1$
this is $\gamma A=A$, which is excluded in $\pbeGs$. For $u\in U(E^+)$ the vectors $uf_a$ are $f_a$
plus a combination of the $f_{a'}$, $a'<a$; so $Z(u\gamma u')=C_{u^{-1}}^*Z(\gamma)C_{u'}$ with
unitriangular $C$, which preserves the leading minors, hence $\pbg$, $\pbsmn_j$ and $\auNm\pbg$; and
it preserves the occurrence of $(\zeta,1)$ at $a$ ($fx=\gamma fy$ becomes $u fx=u\gamma u'(u'^{-1}fy)$,
and $u$, $u'^{-1}$ act unitriangularly on coefficient vectors). For $u,u'\in U^0$ the matrices $C$
are $1$. The last claim follows.
\end{proof}

\begin{lemma}[Near orbits]\label{pbs:near}
If $\mathsf a_{c,c'}(\gamma)\ne0$ then $\gamma$ is near. If $\gamma$ is near, so is $u\gamma u'$ for
$u,u'\in U^0(E^+)$.
\end{lemma}
\begin{proof}
If $\mathsf a_{c,c'}(\gamma)\ne0$, some $y$ has $\beta_c(y)\beta_{c'}(\gamma^{-1}y)W^\flat(y)\otimes
W^\flat(\gamma^{-1}y)\ne0$. Then $y$ and $\gamma^{-1}y$ are reduced, and at every place $W^\flat_\tau
\ne0$ at their components, so $\pbel_\tau\le b^{6n+2}$ there (Proposition~\ref{pbe:local}(e)): both
are boxed. If $y$ and $\gamma^{-1}y$ are boxed, so are $uy$ and $(u\gamma u')^{-1}uy=u'^{-1}\gamma^{-1}
y$, by Lemma~\ref{pbs:reduced}(a) and Lemma~\ref{pbe:len}(a) ($U^0\subseteq F_\tau\cap R_\tau$).
\end{proof}

The heart of the proof is the following bound for one orbit, proved in
\S\S\ref{pbs:ss-schur}--\ref{pbs:ss-orbit}.

\begin{theorem}[One orbit]\label{pbs:orbit}
Let $\gamma\in\Gamma_{c,s,c'}$ have cell rank $j\ge1$ and satisfy $\auNm\pbg(\gamma)\le N^f
\pbsmn_j(\gamma)$, and let $\pbsO=U^0_c\gamma U^0_{c'}$. Then
\[
\sum_{\delta\in\pbsO}\mathsf a_{c,c'}(\delta)\le\nu\,\mathsf S_j\bigl(\mathsf G_j(\gamma)\bigr).
\]
\end{theorem}

For the orbits on the open cell we use the following bound, deduced in \S\ref{pbs:ss-open} from
(O1) and Lemma~\ref{pbs:O2}.

\begin{proposition}[Open cell]\label{pbs:open-block}
Let $T$ be stable with $\auNm\pbg(\gamma)>N^f\pbsmn_n(\gamma)$ on $T$. Then $\|A^\flat_T(c,s,c')\|\le
\nu\sum_\gamma\mathsf S_n(\auNm\pbg(\gamma))$, the sum over the near and chosen
$\gamma\in T\cap\Gamma_{c,s,c'}$.
\end{proposition}

\begin{proposition}[One cell rank]\label{pbs:rank}
Let $T\subseteq\pbeGs$ be stable, all of whose elements have cell rank $j$. Then
\[
\|A^\flat_T(c,s,c')\|\le\nu\sum_{\gamma}\mathsf S_j\bigl(\mathsf G_j(\gamma)\bigr),
\]
the sum over the near and chosen $\gamma\in T\cap\Gamma_{c,s,c'}$.
\end{proposition}
\begin{proof}
Let $\mathcal K=\{\gamma:\auNm\pbg(\gamma)>N^f\pbsmn_j(\gamma)\}$; $T\cap\mathcal K$ and $T\smallsetminus
\mathcal K$ are stable (Lemma~\ref{pbs:cells}) and $A^\flat_T=A^\flat_{T\cap\mathcal K}+A^\flat_{T
\smallsetminus\mathcal K}$ (the sums converge absolutely). On $\mathcal K$, $\pbg\ne0$, so $j=n$, the
indices are $1$ and $\mathsf G_n(\gamma)=\auNm\pbg(\gamma)$; Proposition~\ref{pbs:open-block} bounds
the first block. For the second, by \eqref{pbs:eq-majorant} it suffices to bound $\sum\mathsf a_{c,c'
}(\gamma)$ over $X=(T\smallsetminus\mathcal K)\cap\Gamma_{c,s,c'}$. Group $X$ into the orbits $\pbsO$
of $U^0_c\times U^0_{c'}$ (they lie in $X$: $U^0_c\subseteq U_c$ and $T$ is stable); each orbit has
one chosen element $\gamma_\pbsO$. If $\sum_{\delta\in\pbsO}\mathsf a_{c,c'}(\delta)\ne0$, some $\delta
\in\pbsO$ is near (Lemma~\ref{pbs:near}), hence so is $\gamma_\pbsO$, and by Theorem~\ref{pbs:orbit}
applied to $\gamma_\pbsO$ (which has cell rank $j\ge1$ and lies outside $\mathcal K$) the orbit
contributes at most $\nu\mathsf S_j(\mathsf G_j(\gamma_\pbsO))$. Adding the two bounds gives the
claim, as the near chosen elements of $T\cap\mathcal K$ and of $T\smallsetminus\mathcal K$ are
disjoint sets.
\end{proof}

\subsubsection{The weight of Schur's test}\label{pbs:ss-schur}

Fix a real place $\tau$. In the normalised basis the hermitian form is $h(v,w)=\langle\mathsf J_1v,
w\rangle$ (the form of \S\ref{pbe:ss-real} with its variables exchanged). Let $V_-$ be the $-1$-eigenspace of $\mathsf J_1$, spanned by the
$\tilde e_a-\tilde e_a^*$ and $\tilde z$; it has dimension $n$ and $h(v,v)=-\|v\|^2$ on it. For $y\in
G_\tau$ the space $P_y=yV_-$ is a maximal negative definite subspace, and $h(yv,yv)=-\|v\|^2$ for
$v\in V_-$. Let $N_0=\langle\tilde f_0,\dots,\tilde f_\varrho\rangle=A_\tau\oplus\C\tilde z$; then
$h(\sum_ax_a\tilde f_a,\sum_ax_a\tilde f_a)=-|x_\varrho|^2\le0$. Let $\pi_y$ be the $h$-orthogonal
projection onto $P_y$ (the form is non-degenerate on $P_y$). Let $F_\tau$ be the stabiliser of
$\tilde f_0,\dots,\tilde f_\varrho$ in $G_\tau$, and $U^0_\tau=U_\tau\cap F_\tau$.

\begin{lemma}\label{pbs:proj}
$\pi_y$ is injective on $N_0$, hence an isomorphism $N_0\to P_y$.
\end{lemma}
\begin{proof}
If $w\in N_0$ and $\pi_yw=0$, then $w\in P_y^\perp$, on which $h$ is positive definite (the signature
is $(k+\varrho,n)$ and $P_y$ is negative definite of dimension $n$); as $h(w,w)\le0$, $w=0$.
\end{proof}

\begin{definition}[Volume ratio]\label{pbs:def-v}
For $g,y\in G_\tau$ let $T_{g,y}=(\pi_y|_{N_0})^{-1}\circ\pi_y\circ g|_{N_0}\in\operatorname{End}(N_0)$ and
$\pbsv_y(g)=|\det T_{g,y}|$. (Up to the common factor $|\det(\pi_y|_{N_0})|$ this is the ratio of the
volumes of the projections of $gN_0$ and of $N_0$ to $P_y$.)
\end{definition}

\begin{lemma}\label{pbs:v-props}
Let $g,y\in G_\tau$.
\begin{enumerate}
\item[(a)] $0<\pbsv_y(g)<\infty$, and $y\mapsto\pbsv_y(g)$ is continuous.
\item[(b)] $\pbsv_{g^{-1}y}(g^{-1})=\pbsv_y(g)^{-1}$.
\item[(c)] For $u\in F_\tau$, $\pbsv_{uy}(g)=\pbsv_y(u^{-1}g)$ and $\pbsv_y(gu)=\pbsv_y(g)$.
\end{enumerate}
\end{lemma}
\begin{proof}
(a) $\pi_y g|_{N_0}$ is injective by Lemma~\ref{pbs:proj} applied to $gN_0$ in place of $N_0$ (the
form is $\le0$ on $gN_0$ too); the continuity follows from that of $y\mapsto\pi_y$. (b) $\pi_{g^{-1}y}
=g^{-1}\pi_yg$, so $T_{g^{-1},g^{-1}y}=(g^{-1}\pi_yg|_{N_0})^{-1}g^{-1}\pi_y|_{N_0}=(\pi_yg|_{N_0})^{-1}
\pi_y|_{N_0}=T_{g,y}^{-1}$. (c) $\pi_{uy}=u\pi_yu^{-1}$ and $u|_{N_0}=1$, so $T_{g,uy}=(\pi_yu^{-1}|_{N_0}
)^{-1}\pi_yu^{-1}g|_{N_0}=T_{u^{-1}g,y}$; and $gu|_{N_0}=g|_{N_0}$.
\end{proof}

For $y\in G_\infty$ and $\gamma\in\Gamma$ put $\pbsv_y(\gamma)=\prod_\tau\pbsv_{y_\tau}(\gamma_\tau)$. Let
$\beta^\square_c=\beta_c\cdot\mathbf 1_{\text{boxed}}$. For an integer $i\ge0$ the \emph{Schur
integral} is
\[
I^{(i)}_{c,c'}(\gamma)=\int_{G_\infty}\beta^\square_c(y)\,\|W^\flat(y)\|^2\,\beta^\square_{c'}(
\gamma^{-1}y)\,\pbsv_y(\gamma)^{-(k+i)}\,dy .
\]

\begin{lemma}[Cauchy--Schwarz with a weight]\label{pbs:cs}
For every $i\ge0$, $\mathsf a_{c,c'}(\gamma)^2\le I^{(i)}_{c,c'}(\gamma)\,I^{(i)}_{c',c}(\gamma^{-1})$.
\end{lemma}
\begin{proof}
Where $\beta_c(y)\beta_{c'}(\gamma^{-1}y)W^\flat(y)\otimes W^\flat(\gamma^{-1}y)\ne0$, both points are
boxed (proof of Lemma~\ref{pbs:near}), so $\beta_c=\beta_c^\square$ there. Write $e=k+i$ and
$\pbsv=\pbsv_y(\gamma)$. The integrand of $\mathsf a_{c,c'}(\gamma)$ is the product of $(\beta^\square_c
(y)\|W^\flat(y)\|^2\beta^\square_{c'}(\gamma^{-1}y)\pbsv^{-e})^{1/2}$ and $(\beta^\square_{c'}(\gamma^{-1}y)
\|W^\flat(\gamma^{-1}y)\|^2\beta^\square_c(y)\pbsv^{e})^{1/2}$. By the Cauchy--Schwarz inequality,
$\mathsf a_{c,c'}(\gamma)^2$ is at most $I^{(i)}_{c,c'}(\gamma)$ times the integral of the square of
the second factor. In that integral substitute $y=\gamma y'$; by Lemma~\ref{pbs:v-props}(b) (with
$g=\gamma^{-1}$, applied at each place), $\pbsv_{\gamma y'}(\gamma)=\pbsv_{y'}(\gamma^{-1})^{-1}$, so the
integral is $I^{(i)}_{c',c}(\gamma^{-1})$.
\end{proof}

\begin{lemma}[Integration in stages]\label{pbs:stages}
Let $\Phi\colon G_\infty\to[0,\infty]$ be measurable and left invariant under $U^0_c$. Let $(X,m)$
be a $\sigma$-finite measure space, $t\colon X\to U^0_\infty$ measurable, $B\subseteq X$ with $0<
m(B)<\infty$, and $C\ge0$ such that
\[
\int_B\Phi(u\,t(x)\,y)\,dm(x)\le m(B)\,C\qquad\text{for all }u\in U^0_\infty\text{ and all boxed }y.
\]
Then $\int_{G_\infty}\beta^\square_c(y)\|W^\flat(y)\|^2\Phi(y)\,dy\le C\,\nu\,d_\infty$, where $d_\infty=
\dim\pbW=\prod_\tau d_\tau$.
\end{lemma}
\begin{proof}
Disintegrate $dg$ along the orbits of $U^0_\infty$: $\int_{G_\infty}F=\int_{U^0_\infty\backslash G_\infty}
\int_{U^0_\infty}F(ua)\,du^0\,da$ for $F\ge0$. On an orbit $U^0_\infty a$ the functions $\omega$,
$\mathbf 1_{\text{boxed}}$ and $\|W^\flat\|$ are constant (Lemma~\ref{pbs:reduced}(a),
Lemma~\ref{pbe:len}(a), and $W^\flat(uy)=\psi_\infty(u)W^\flat(y)$), so $\beta^\square_c(ua)=p_c(ua)\,
\omega(a)\mathbf 1_{\text{boxed}}(a)$. It therefore suffices to show, for boxed $a$,
\begin{equation}\label{pbs:eq-stage}
\int_{U^0_\infty}p(u)\,\Phi_a(u)\,du^0\le C\int_{U^0_\infty}p(u)\,du^0,\qquad p(u)=p_c(ua),\ \Phi_a(u)=
\Phi(ua);
\end{equation}
for then the left side of the lemma is at most $C\int\beta_c\|W^\flat\|^2=C\nu\int_{U_\infty\backslash
G_\infty}\|W^\flat\|^2$ by \eqref{pbs:eq-nu}, and $\int_{U_\tau\backslash G_\tau}\|W^\flat_\tau\|^2\le
\int\|W_\tau\|^2=d_\tau$ (Proposition~\ref{pbe:local}(e)). Now $p$ is a cutoff for the left action of
the discrete group $U^0_c$ on $U^0_\infty$, and $\Phi_a$ is invariant under it. For $s\in U^0_\infty$,
$\int p(u)\Phi_a(us)\,du^0=\int_{U^0_c\backslash U^0_\infty}\Phi_a(us)=\int_{U^0_c\backslash U^0_\infty}
\Phi_a(u)=\int p\,\Phi_a$, because right translation by $s$ preserves the invariant measure of
$U^0_c\backslash U^0_\infty$ ($U^0_\infty$ is unimodular). Integrating over $x\in B$ with $s=t(x)$ and
using Tonelli,
\[
m(B)\int p\,\Phi_a=\int_{U^0_\infty}p(u)\int_B\Phi(u\,t(x)\,a)\,dm(x)\,du^0\le m(B)\,C\int p ,
\]
which is \eqref{pbs:eq-stage}.
\end{proof}

\subsubsection{The weight at one real place}\label{pbs:ss-oneplace}

Fix $\tau$ and drop it from the notation. For $g\in G_\tau$ let $\tilde Z(g)=(h(\tilde f_a,g\tilde
f_b))_{a,b}$, and let $\Phi(g)=(\Phi(g)_b)_{0\le b\le\varrho}$ with $\Phi(g)_b\in X_{0,\tau}$ the
component in $X_{0,\tau}$ of $g\tilde f_b$ for the decomposition $Y_\tau=A_\tau\oplus X_{0,\tau}\oplus
\C\tilde z\oplus A^*_\tau$, $A^*_\tau=\langle\tilde e^*_a\rangle$. We say that $(\zeta,1)$ \emph{occurs at
$a$ for $g$} if there are $x,x'\in\C^n$ supported on $\{0,\dots,a\}$ with $\sum x_b\tilde f_b=g\sum
x'_b\tilde f_b$, $x_a=\zeta$, $x'_a=1$, and that \emph{norm-one pairs occur from $j$} if for one $\zeta$
with $|\zeta|=1$ this holds at every $a$ with $j\le a\le\varrho$. If $\gamma\in\Gamma$ has cell rank
$j$, norm-one pairs occur from $j$ for $\gamma_\tau$ (rescale the common vectors of
Definition~\ref{ppi:def-cell}: $f_b=\tilde f_b/s_b$, and multiply by $s_a$; the leading coefficients
become $\tau(\zeta)$ and $1$), and $\tilde Z(\gamma_\tau)_{ab}=s_as_b\,\tau(Z(\gamma)_{ab})$.

Let $0\le j\le n$, $i=\min(\varrho,j)$. The \emph{subframe} is $\mathsf s=(\tilde f_0,\dots,\tilde
f_{i-1},\tilde f_\varrho)$; for a matrix $Z$ indexed by the frame, $B(Z)$ is its $(i+1)\times(i+1)$
block on the indices of the subframe, and $B_<(Z)$ the first $i$ rows of $B(Z)$. For $g\in G_\tau$
write $g\mathsf s$ for $(g\tilde f_0,\dots,g\tilde f_{i-1},g\tilde f_\varrho)$ and $\Phi_{\mathsf s}(g)$ for
the corresponding columns of $\Phi(g)$.

\begin{lemma}[Common vectors]\label{pbs:common}
Let norm-one pairs occur from $j$ for $g$, $Z=\tilde Z(g)$.
\begin{enumerate}
\item[(a)] For each $b$ with $j\le b<\varrho$ there are $x^{(b)},x'^{(b)}$, supported on $\{0,\dots,
b\}$, with $x'^{(b)}_b=1$, $|x^{(b)}_b|=1$ and $c_b:=\sum_ax^{(b)}_a\tilde f_a=g\sum_ax'^{(b)}_a\tilde
f_a=:gc'_b$; and $c_b\in A_\tau\cap gA_\tau$ is $h$-orthogonal to $gN_0$, and $x^{(b)*}Z=0$.
\item[(b)] If $j\le\varrho$ there are $x\in\C^n$ with $|x_\varrho|=1$ and $\mathsf c\in\C^{i+1}$ with
$\mathsf c_{i}=1$ (last entry) such that $g\mathsf s\,\mathsf c=\sum_ax_a\tilde f_a$, and $B(Z)\,\mathsf c
=(0,\dots,0,-x_\varrho)^{\mathsf T}$.
\item[(c)] $|\det B(Z)|=|\det(Z_{ab})_{a,b<j}|$ (for $j=n$ this is $\det Z$).
\end{enumerate}
\end{lemma}
\begin{proof}
(a) The vectors exist by definition. As $c_b=gc'_b$ and $g$ is an isometry, $|x^{(b)}_\varrho|=|x'^{(b)
}_\varrho|$ ($h(c,c)=-|x_\varrho|^2$ on $N_0$), and both vanish as $b<\varrho$. So $c_b\in A_\tau$, $c'_b\in
A_\tau$. For $w\in N_0$, $h(c_b,gw)=h(c'_b,w)=0$ because $A_\tau\perp N_0$; this is $x^{(b)*}Z=0$.
(b) Take the common vector at $a=\varrho$: $\sum x_a\tilde f_a=g\sum x'_a\tilde f_a$ with $x'_\varrho=1$,
$|x_\varrho|=1$. Subtracting the common vectors of (a) (multiplied by $x'_b$, for $b=\varrho-1,\dots,j$
in turn) we may assume $x'_b=0$ for $i\le b<\varrho$; then $x'$ is supported on the subframe and
$\mathsf c$ is its restriction (the norm equality still gives $|x_\varrho|=1$). The products of the
subframe with $\sum x_a\tilde f_a$ are $h(\tilde f_a,\sum x_b\tilde f_b)=-x_\varrho\delta_{a\varrho}$.
(c) If $j=n$ the subframe is the whole frame. If $j\le\varrho$, $i=j$; by Cramer's rule, replacing the
last column of $B(Z)$ by $B(Z)\mathsf c$ multiplies the determinant by $\mathsf c_i=1$; expanding the new
matrix along its last column, whose only non-zero entry is $-x_\varrho$ in the last row, gives $\pm
x_\varrho\det(Z_{ab})_{a,b<j}$.
\end{proof}

\begin{lemma}[The weight modulo common vectors]\label{pbs:mod}
Let norm-one pairs occur from $j$ for $g$, let $y\in G_\tau$, and let $\omega_0,\dots,\omega_i\in P_y$
be orthogonal to $A_\tau\cap gA_\tau$, with $\det(h(\mathsf s_a,\omega_l))_{a,l}\ne0$. Then
\[
\pbsv_y(g)=\frac{|\det(h(g\mathsf s_a,\omega_l))_{a,l}|}{|\det(h(\mathsf s_a,\omega_l))_{a,l}|}.
\]
\end{lemma}
\begin{proof}
Let $\lambda\colon Y_\tau\to\C^{i+1}$, $\lambda(w)=(h(w,\omega_l))_l$. As $\omega_l\in P_y$,
$\lambda\circ\pi_y=\lambda$, hence $\lambda\circ T_{g,y}=\lambda\circ g$ on $N_0$. The map $\lambda$ is
an isomorphism on $\langle\mathsf s\rangle$ and vanishes on the vectors $c_b$ ($j\le b<\varrho$, and
$i\le b$) of Lemma~\ref{pbs:common}(a). Two bases of $N_0$: $\mathcal B'=(\mathsf s,(c'_b)_{i\le b<
\varrho})$ and $\mathcal B=(\mathsf s,(c_b)_{i\le b<\varrho})$. In the frame order, the matrix of
$\mathcal B'$ (resp.\ $\mathcal B$) in the basis $\tilde f$ is triangular with diagonal entries $1$
(resp.\ $1$ or $x^{(b)}_b$, of modulus $1$); so $|\det T_{g,y}|$ is the modulus of the determinant of
the matrix of $T_{g,y}$ from $\mathcal B'$ to $\mathcal B$. As $T_{g,y}c'_b=c_b$ (because $\pi_ygc'_b=
\pi_yc_b$), this matrix is block triangular $\bigl(\begin{smallmatrix}\mathsf X&0\\\mathsf Y&1
\end{smallmatrix}\bigr)$, with $T_{g,y}\mathsf s_l=\sum_a\mathsf X_{al}\mathsf s_a+\sum_b\mathsf Y_{bl}c_b$.
Applying $\lambda$: $\lambda(g\mathsf s_l)=\lambda(T_{g,y}\mathsf s_l)=\sum_a\mathsf X_{al}\lambda(\mathsf
s_a)$. Taking determinants, $\det(\lambda(g\mathsf s))=\det\mathsf X\cdot\det(\lambda(\mathsf s))$ and
$\pbsv_y(g)=|\det\mathsf X|$.
\end{proof}

\begin{lemma}[Triangular bases]\label{pbs:tri}
Let $Z$ be a matrix indexed by the frame with $\delta:=\det B(Z)\ne0$, and $y\in G_\tau$. Put $D=Z_{
\cdot,\mathsf s}B(Z)^{-1}$ ($n\times(i+1)$). There are $\omega_0,\dots,\omega_i\in P_y$ and a lower
triangular invertible $L$ with $h(\omega_l,\omega_m)=-\delta_{lm}$ and $(h(\tilde f_a,\omega_l))_{a,l}=
DL$. If $g$ has $\tilde Z(g)=Z$ and norm-one pairs occur from $j$ for $g$, then:
\begin{enumerate}
\item[(a)] the $\omega_l$ are orthogonal to $A_\tau\cap gA_\tau$ and $(h(\mathsf s_a,\omega_l))=L$; hence
$\pbsv_y(g)=|\det(h(g\mathsf s_a,\omega_l))|/|\det L|$;
\item[(b)] with $\omega'=g\mathsf s\,B(Z)^{-1}L$ and $v=(v_0,\dots,v_i)$, $v_l\in X_{0,\tau}$ the
$X_0$-component of $\omega'_l-\omega_l$:
\[
|\delta|\det(1+vv^*)\le2^{i+1}|\det L|^2\,\pbsv_y(g);
\]
\item[(c)] if $\|y^{-1}\tilde f_a\|\le\rho$ for all $a$, then $|\det L|\le\rho^{i+1}$;
\item[(d)] if $\pbel_\tau(y)\le\theta$, then $|L_{ii}|^2\le\theta^{2(j-i)}(1+\|v_i\|^2)$.
\end{enumerate}
Here $vv^*$ is the endomorphism $x\mapsto\sum_lv_l\langle x,v_l\rangle$ of $X_{0,\tau}$.
\end{lemma}
\begin{proof}
\emph{Construction.} By Lemma~\ref{pbs:proj}, $w\mapsto(h(\tilde f_a,w))_a$ is an isomorphism $P_y\to
\C^n$ ($P_y\cong N_0^*$ through $h$, as $\pi_y|_{N_0}$ is an isomorphism and $h$ is non-degenerate on
$P_y$). Let $\omega^0_l\in P_y$ have products $D_{\cdot l}$; they are linearly independent because
$D$ has the left inverse $(\text{rows of }\mathsf s)$: the rows of $D$ at the subframe indices form
$B(Z)B(Z)^{-1}=1$. The matrix $\mathsf G_0=(-h(\omega^0_l,\omega^0_m))_{l,m}$ is positive definite, since $-h$ is
positive definite on $P_y$; by the Cholesky decomposition $\mathsf G_0^{-1}=LL^*$ with $L$ lower
triangular and invertible, i.e.\ $L^*\mathsf G_0L=1$. Put $\omega_m=\sum_l\omega^0_lL_{lm}$; then
$h(\omega_l,\omega_m)=-\delta_{lm}$. Then the products with the frame are $DL$.

(a) Let $c=\sum x_a\tilde f_a\in A_\tau\cap gA_\tau$; as in Lemma~\ref{pbs:common}(a), $x^*Z=0$
(write $c=gc'$ with $c'\in A_\tau$; then $h(c,gw)=h(c',w)=0$ for $w\in N_0$). Then $h(c,\omega_l)=(x^*DL)_l=(x^*Z_{\cdot,\mathsf s}B^{-1}L)_l=0$. The
products of the subframe with $\omega$ are the subframe rows of $DL$, that is $L$. Lemma~\ref{pbs:mod}
gives the formula.

(b) The products of the frame with $\omega'$ are $Z_{\cdot,\mathsf s}B^{-1}L=DL$, the same as with
$\omega$. So $\omega'_l-\omega_l\perp N_0$, i.e.\ $\omega'_l-\omega_l\in(A_\tau\oplus\C\tilde z)^\perp=
A_\tau\oplus X_{0,\tau}$, and as $A_\tau$ is isotropic and orthogonal to $A_\tau\oplus X_{0,\tau}$,
$h(\omega'_l-\omega_l,\omega'_m-\omega_m)=\langle v_m,v_l\rangle$ (with $h=\langle\ ,\ \rangle$ on
$X_{0,\tau}$). Let $\mathsf K=(h(\omega_l,\omega'_m))_{l,m}$ and $\mathsf G'=(h(\omega'_l,\omega'_m))$.
Expanding the Gram matrix of $\omega'-\omega$ gives $v^\dagger v=\mathsf G'-\mathsf K-\mathsf K^*+(-1)$,
where $v^\dagger v=(\langle v_m,v_l\rangle)_{l,m}$; so $(-\mathsf K)+(-\mathsf K)^*-(1+v^\dagger v)=-\mathsf G'$,
which is positive semidefinite since $\omega'_l\in gN_0$ and $h\le0$ on $gN_0$. By the
Ostrowski--Taussky inequality (Lemma~\ref{pbs:OT} below), $\det(1+v^\dagger v)\le2^{i+1}|\det\mathsf K|$.
Now $\mathsf K=(h(\omega_l,g\mathsf s_m))\,B^{-1}L$, and $(h(\omega_l,g\mathsf s_m))_{l,m}$ is the
conjugate transpose of $(h(g\mathsf s_m,\omega_l))_{m,l}$, so by (a) $|\det\mathsf K|=\pbsv_y(g)|\det L|
^2/|\delta|$. Finally $\det(1+v^\dagger v)=\det(1+vv^*)$.

(c) $L_{al}=h(\mathsf s_a,\omega_l)=h(y^{-1}\mathsf s_a,y^{-1}\omega_l)$, and the $y^{-1}\omega_l\in V_-$
are orthonormal for $\|\cdot\|$ (since $h=-\|\cdot\|^2$ on $V_-$). By \eqref{pbe:eq-norm}, $|L_{ll}|\le
\|y^{-1}\mathsf s_l\|\le\rho$; and $\det L=\prod_lL_{ll}$.

(d) \emph{Case $j\le\varrho$} ($i=j$). Let $x,\mathsf c$ be as in Lemma~\ref{pbs:common}(b). The last
column of $L$ is $L_{ii}e_i$ ($L$ is lower triangular), and $B^{-1}e_i=-\mathsf c/x_\varrho$, so $\omega'_i=
L_{ii}\,g\mathsf sB^{-1}e_i=t\sum_ax_a\tilde f_a$ with $t=-L_{ii}/x_\varrho$. Thus $\omega'_i\in N_0$,
hence $\omega'_i\perp\omega'_i-\omega_i$, and $h(\omega'_i,\omega'_i)=-|tx_\varrho|^2=-|L_{ii}|^2$. So
$-1=h(\omega_i,\omega_i)=h(\omega'_i,\omega'_i)+h(\omega'_i-\omega_i,\omega'_i-\omega_i)=-|L_{ii}|^2+\|v_i
\|^2$, i.e.\ $|L_{ii}|^2=1+\|v_i\|^2$. \emph{Case $j=n$} ($i=\varrho$, $j-i=1$). Choose $s\in R_\tau$
with $\|sy\|\le\theta$ (possible since $\pbel_\tau(y)=\inf_{s\in R_\tau}(\|sy\|+1)\le\theta$). Then $(sy)^{-1}\tilde z=y^{-1}s^{-1}\tilde z=y^{-1}(\tilde z+\sum_{a<\varrho}
\mu_a\tilde f_a)$ (elements of $R_\tau$ are unitriangular on the frame), and $h(\tilde f_a,\omega_\varrho)
=L_{a\varrho}=0$ for $a<\varrho$. Hence $L_{\varrho\varrho}=h((sy)^{-1}\tilde z,y^{-1}\omega_\varrho)$, of
modulus at most $\|(sy)^{-1}\|=\|sy\|\le\theta$; and $1\le1+\|v_\varrho\|^2$.
\end{proof}

\begin{lemma}[Ostrowski--Taussky]\label{pbs:OT}
Let $H$ and $\Theta+\Theta^*-H$ be positive semidefinite $m\times m$ complex matrices. Then $|\det H|\le
2^m|\det\Theta|$.
\end{lemma}
\begin{proof}
If $\det H=0$ there is nothing to prove; otherwise $H$ is positive definite. Let $\ell$ be an
eigenvalue of $2H^{-1}\Theta$ with eigenvector $v\ne0$: $2\Theta v=\ell Hv$. Then $\ell\,v^*Hv=2v^*\Theta
v$, and taking real parts, $(\operatorname{Re}\ell)\,v^*Hv=v^*(\Theta+\Theta^*)v\ge v^*Hv>0$; so
$\operatorname{Re}\ell\ge1$ and $|\ell|\ge1$. Hence $|\det(2H^{-1}\Theta)|=2^m|\det\Theta|/\det H\ge1$.
\end{proof}

\begin{lemma}[Hua's integral]\label{pbs:hua}
Let $k\ge7$, $\mathsf G$ a positive definite $k\times k$ matrix, $\mathsf S$ an invertible $i\times i$
matrix and $\mathsf c$ an $i\times k$ matrix. With Lebesgue measure on $\C^{i\times k}$,
\[
\int_{\C^{i\times k}}\det\bigl(\mathsf G+(\mathsf S\mathsf A+\mathsf c)^*(\mathsf S\mathsf A+\mathsf
c)\bigr)^{-(k+i)}d\mathsf A\le|\det\mathsf S|^{-2k}(\det\mathsf G)^{-k}.
\]
\end{lemma}
\begin{proof}
The real-linear map $\mathsf A\mapsto\mathsf S\mathsf A$ has determinant $|\det\mathsf S|^{2k}$, so after
the affine substitution it suffices to treat $\mathsf S=1$, $\mathsf c=0$. Then $\mathsf A^*\mathsf A=
\sum_l w_l^*w_l$ over the rows $w_l\in\C^k$ (row vectors). We show by induction on the number $i$ of
rows: for $\mathsf M>0$ and $r\ge k$,
\[
\int\det\Bigl(\mathsf M+\sum_{l<i}w_l^*w_l\Bigr)^{-(i+r)}dw\le\Bigl(\frac{\pi^k}{k!}\Bigr)^i(\det
\mathsf M)^{-r}.
\]
For one row: write $\mathsf M=\mathsf C^*\mathsf C$ and $w=u\mathsf C$; then $\det(\mathsf M+w^*w)=
\det\mathsf M\,(1+\|u\|^2)$, $dw=\det\mathsf M\,du$, and $\int_{\C^k}(1+\|u\|^2)^{-(r+1)}du\le\int(1+\|u\|^2)
^{-(k+1)}du=\pi^k/k!$; so the integral is at most $(\pi^k/k!)(\det\mathsf M)^{-r}$. For $i$ rows,
integrate first over the rows $l\ge1$ with $\mathsf M'=\mathsf M+w_0^*w_0$ and $r+1$ in place of $r$
(induction hypothesis), then over $w_0$ with the one-row case. Finally $\pi^k/k!\le1$ for $k\ge7$.
\end{proof}

\begin{theorem}[One real place]\label{pbs:oneplace}
Let $k\ge7$, $y\in G_\tau$ with $\|y^{-1}\tilde f_a\|\le\rho$ for all $a$ and $\pbel_\tau(y)\le\theta$,
where $\rho,\theta\ge1$. Let $1\le j\le n$, $i=\min(\varrho,j)$, let $g_0\in G_\tau$ be such that
norm-one pairs occur from $j$ for $g_0$, and let $Z=\tilde Z(g_0)$ with $\delta=\det B(Z)\ne0$. There is
a measurable $\varphi\colon X_{0,\tau}^{\,i}\to[0,\infty)$ with
\[
\int_{X_{0,\tau}^i}\varphi\le\frac{(2\rho^2)^{(i+1)i}\bigl(2^{i+1}\theta^{2(j-i)}\bigr)^k}{|\delta|^{k+i}}
\]
such that $\pbsv_y(g)^{-(k+i)}\le\varphi(a_0,\dots,a_{i-1})$ for every $g\in G_\tau$ with $\tilde
Z(g)=Z$, norm-one pairs occurring from $j$, and $\Phi_{\mathsf s}(g)=\Phi_{\mathsf s}(g_0)+\sum_{l<i}a_l
\,B_<(Z)_{l,\cdot}$ (that is, $\Phi_{\mathsf s}(g)_m=\Phi_{\mathsf s}(g_0)_m+\sum_{l<i}B(Z)_{lm}a_l$).
\end{theorem}
\begin{proof}
Take $\omega,L$ from Lemma~\ref{pbs:tri} for $Z$ and $y$; they do not depend on $g$. Let $v^{(0)}$ be
the vectors of Lemma~\ref{pbs:tri}(b) for $g_0$. For $g$ as in the statement, $\omega'(g)=g\mathsf s\,
B^{-1}L$, so the $X_0$-components satisfy $v(g)=\Phi_{\mathsf s}(g)B^{-1}L-(\omega)_{X_0}$ (in the
obvious matrix notation, columns indexed by the subframe), whence
\[
v(g)=v^{(0)}+\Bigl(\sum_{l<i}a_lB_{l\cdot}\Bigr)B^{-1}L=v^{(0)}+\Bigl(\sum_{l<i}a_lL_{l\cdot}\Bigr),
\]
because the first $i$ rows of $B\,B^{-1}L$ are those of $L$. As $L$ is lower triangular, $L_{li}=0$ for
$l<i$: the last vector $v_i(g)=v^{(0)}_i$ does not depend on $g$, and the others are $v_m(g)=v^{(0)}_m+
\sum_{l<i}L_{lm}a_l$ ($m<i$). Put $\mathsf G=1+v^{(0)}_iv^{(0)*}_i$, $L_1=(L_{lm})_{l,m<i}$ and, in
coordinates of an orthonormal basis of $X_{0,\tau}$, $\mathsf w(\mathsf a)=L_1^{\mathsf T}\mathsf a+
\mathsf c$ ($i\times k$; rows $w_m=v_m(g)$). Then $1+v(g)v(g)^*=\mathsf G+\mathsf w^*\mathsf w$, and by
Lemma~\ref{pbs:tri}(a),(b)
\[
\pbsv_y(g)^{-(k+i)}\le\varphi(\mathsf a):=\Bigl(\frac{2^{i+1}|\det L|^2}{|\delta|}\Bigr)^{k+i}\det(\mathsf
G+\mathsf w(\mathsf a)^*\mathsf w(\mathsf a))^{-(k+i)} .
\]
By Lemma~\ref{pbs:hua} ($\det L_1^{\mathsf T}=\det L_1\ne0$, $|\det L|=|\det L_1||L_{ii}|$, $\det\mathsf G=1+
\|v^{(0)}_i\|^2$),
\[
\int\varphi\le\frac{(2^{i+1}|\det L|^2)^{k+i}}{|\delta|^{k+i}|\det L_1|^{2k}(\det\mathsf G)^k}=\frac{(2^{i+1}
|\det L|^2)^i\,\bigl(2^{i+1}|L_{ii}|^2/(1+\|v^{(0)}_i\|^2)\bigr)^k}{|\delta|^{k+i}}.
\]
Lemma~\ref{pbs:tri}(c),(d) (for $g_0$) bound this by $(2^{i+1}\rho^{2(i+1)})^i(2^{i+1}\theta^{2(j-i)})^k/
|\delta|^{k+i}$, which is the claim.
\end{proof}

\subsubsection{An orbit: colouring, shifts and packing}\label{pbs:ss-orbit}

Fix $\gamma\in\Gamma_{c,s,c'}$ of cell rank $j\ge1$, $i=\mathsf j'=\min(\varrho,j)\ge1$, and its orbit
$\pbsO=U^0_c\gamma U^0_{c'}$. For $\delta\in\Gamma$ let $\Phi(\delta)=(\Phi(\delta)_b)_{b\le\varrho}\in
X_0^n$, $\Phi(\delta)_b$ the $X_0$-component of $\delta f_b$ for $Y_0=A\oplus X_0\oplus Ez\oplus A^*$.

\begin{lemma}\label{pbs:indep}
Let $\mathfrak R$ (resp.\ $\mathfrak R'$) be the $\auO_E$-module spanned by the rows $Z(\gamma)_{a,\cdot}
\in E^n$ with $a<\varrho$ (resp.\ $a<i$). The rows with $a<i$ are linearly independent and span the
same $E$-space as those with $a<\varrho$; so $[\mathfrak R:\mathfrak R']$ is finite, and it equals
$i_\gamma$ of \S\ref{pbs:ss-statement}.
\end{lemma}
\begin{proof}
Independence: the minor $\delta_j(\gamma)$ (for $j\le\varrho$), resp.\ $\det Z(\gamma)$ (for $j=n$), is
non-zero. If $j\le a<\varrho$, the common vector $fx=\gamma fy$ with top index $a$ and $x_a=\zeta\ne0$
has $y_\varrho=0$, so $\sum_{a'\le a}\aubar{x_{a'}}Z_{a'b}=h(fx,\gamma f_b)=h(fy,f_b)=0$ for all $b$:
the row $a$ is a combination of the rows $a'<a$. For the index: $h_c$ and $h_{c'}$ fix the $f_a$, so
$h(f_a,k_\gamma f_b)=h(f_a,\gamma f_b)$ and the vectors $\xi_a(k_\gamma)$ of \S\ref{pbs:ss-statement}
are the rows of $Z(\gamma)$; the index of the completions $\widehat\auO\otimes\mathfrak R'\subseteq
\widehat\auO\otimes\mathfrak R$ is the index $[\mathfrak R:\mathfrak R']$.
\end{proof}

\begin{lemma}[Moving inside an orbit]\label{pbs:move}
Let $\delta=u\gamma u'$ with $u\in U^0_c$, $u'\in U^0_{c'}$. Then $Z(\delta)=Z(\gamma)$, $\delta$ has cell
rank $j$, $\delta\in\Gamma_{c,s,c'}$, and $\delta f_b-\gamma f_b\in A\oplus X_0$ with
\[
\Phi(\delta)_b=\Phi(\gamma)_b+\sum_{a<\varrho}Z(\gamma)_{ab}\,q_a(u),\qquad q_a(u)\in M\Lambda_c,
\]
where $q_a(u)$ is the $X_0$-component of $(u-1)e^*_{a+1}$.
\end{lemma}
\begin{proof}
The first three claims are Lemma~\ref{pbs:cells} and the stability of $\Gamma_{c,s,c'}$. Now $\delta f_b
=u\gamma f_b$. The element $u\in U^0(E^+)$ fixes $A$ and $z$; it maps $X_0$ into $X_0+A$ (it is trivial
on $A^\perp/A$); and $ue^*_a=e^*_a+q+\alpha$ with $q\in X_0$, $\alpha\in A$ (its products with $z$ and the
$e_b$ are those of $e^*_a$). For $w\in Y_0$ the $A^*$-component of $w$ is $\sum_ah(e_a,w)e^*_a$. Hence
$uw-w\in A\oplus X_0$ with $X_0$-component $\sum_ah(e_a,w)q_a$; for $w=\gamma f_b$ this is the formula.
For integrality: $u\in U_c$ means $(u-1)h_cL\subseteq Mh_cL$ (at every finite place), and $e^*_a\in h_cL$
($h_c$ fixes $e^*_a$). Write $h_cL=L_{\mathrm{hyp}}\oplus L'_{\mathrm{mid}}$ as in the proof of
Lemma~\ref{pbe:eichler}; the vector $(u-1)e^*_a=q_a+\alpha$ has $L_{\mathrm{hyp}}$-component $\alpha$ and
$L'_{\mathrm{mid}}$-component $q_a$, so $q_a\in ML'_{\mathrm{mid}}\cap X_0=M\Lambda_c$ locally everywhere,
hence $q_a\in M\Lambda_c$.
\end{proof}

\begin{lemma}[Colouring]\label{pbs:colour}
The orbit $\pbsO$ is the union of at most $i_\gamma^{\,k}$ subsets $\mathcal B$ with the following
property: for $p_0\in\mathcal B$ and every $p\in\mathcal B$ there is $l_p=(l_{p,0},\dots,l_{p,i-1})\in
\Lambda_c^{\,i}$ with
\[
\Phi(p)_b=\Phi(p_0)_b+M\sum_{a<i}Z(\gamma)_{ab}\,l_{p,a}\qquad(0\le b\le\varrho).
\]
\end{lemma}
\begin{proof}
Let $\mathfrak M=\Lambda_c\otimes_{\auO_E}\mathfrak R$ and $\mathfrak M'=\Lambda_c\otimes\mathfrak R'$,
regarded inside $X_0\otimes_EE^n=X_0^n$ via $l\otimes\xi\mapsto(\xi_bl)_b$ (injective, as $\Lambda_c$ is
projective, hence flat). Then $[\mathfrak M:\mathfrak M']=|\Lambda_c\otimes\mathfrak R/\mathfrak R'|=
[\mathfrak R:\mathfrak R']^k=i_\gamma^k$ (locally $\Lambda_c\cong\auO^k$). By Lemma~\ref{pbs:move},
$\Phi(\delta)-\Phi(\gamma)\in M\mathfrak M$ for $\delta\in\pbsO$. Colour $\delta$ by the class of $M^{-1}
(\Phi(\delta)-\Phi(\gamma))$ in $\mathfrak M/\mathfrak M'$. Within a colour, $\Phi(p)-\Phi(p_0)\in
M\mathfrak M'$, which is the claim.
\end{proof}

\paragraph*{Shifts}
At a place $\tau$, every $u\in U^0_\tau$ satisfies, as in Lemma~\ref{pbs:move}, $\Phi(u^{-1}g)_b=\Phi(
g)_b+\sum_{a<\varrho}\tilde Z(g)_{ab}\,t_a(u)$ with $t_a(u)$ the $X_{0,\tau}$-component of $u^{-1}\tilde
e^*_{a+1}$; conversely, for every $t=(t_a)_{a<\varrho}\in X_{0,\tau}^\varrho$ there is $u_t\in U^0_\tau$ with
$t_a(u_t)=t_a$, depending continuously on $t$: put $u_t^{-1}=1$ on $N_0$, $u_t^{-1}x=x-\sum_ah(t_a,x)\tilde
e_{a+1}$ on $X_{0,\tau}$, and $u_t^{-1}\tilde e^*_{a+1}=\tilde e^*_{a+1}+t_a-\frac12\sum_{a'}h(t_{a'},t_a)
\tilde e_{a'+1}$; a direct check shows that $u_t^{-1}$ preserves $h$, and it lies in $U^0_\tau$. Let
$V_{\mathrm{col}}=\prod_\tau X_{0,\tau}^{\,i}$, with Lebesgue measure for orthonormal coordinates
(the basis $x_1,\dots,x_k$ of $X_0$ is orthonormal at every $\tau$), and for $\mathsf x=(\mathsf x_{\tau,
a})\in V_{\mathrm{col}}$ let $u_{\mathsf x}\in U^0_\infty$ have components $u_{t(\tau)}$ with $t(\tau)_a=
(M/s_a)\mathsf x_{\tau,a}$ for $a<i$ and $t(\tau)_a=0$ for $a\ge i$. Let $L_{\mathrm{col}}\subseteq
V_{\mathrm{col}}$ be the image of $\Lambda_c^{\,i}$ under $l\mapsto(\tau(l_a))_{\tau,a}$.

\begin{lemma}\label{pbs:col}
$L_{\mathrm{col}}$ is a lattice in $V_{\mathrm{col}}$ with a measurable fundamental domain $\mathsf D$
with $1\le\pbvol(\mathsf D)<\infty$.
\end{lemma}
\begin{proof}
$\Lambda_c^{\,i}$ is a projective $\auO_E$-module of rank $ik$, and $l\mapsto(\tau(l))_\tau$ maps it onto a
lattice of $V_{\mathrm{col}}$ (real dimension $2ikf=[E:\Q]\,ik$). For $l,l'\in\Lambda_c^{\,i}$ the
Euclidean product of their images is $\sum_\tau\operatorname{Re}\tau(\alpha)$ with $\alpha=\sum_ah(l_a,
l'_a)\in\auO_E$ ($h_cL$ is an integral hermitian lattice). Since $E$ contains the $5040$-th roots of
unity, $i=\sqrt{-1}\in E$, and $\sum_\tau\operatorname{Re}\tau(\alpha)=\frac12\operatorname{Tr}_{E/\Q}
\alpha=\frac12\operatorname{Tr}_{\Q(i)/\Q}(\operatorname{Tr}_{E/\Q(i)}\alpha)=\operatorname{Re}\beta\in
\Z$ for $\beta=\operatorname{Tr}_{E/\Q(i)}\alpha\in\Z[i]$. So the Gram matrix of a $\Z$-basis of
$L_{\mathrm{col}}$ is an integral positive definite matrix, its determinant is $\ge1$, and the
covolume, its square root, is $\ge1$.
\end{proof}

\begin{lemma}[Shifted points]\label{pbs:shifted}
Let $\mathcal B$, $p_0$, $l_p$ be as in Lemma~\ref{pbs:colour}, $u\in U^0_\infty$, $\mathsf x\in V_{\mathrm{
col}}$. Then for every $p\in\mathcal B$ and every $\tau$: $\tilde Z((uu_{\mathsf x})^{-1}_\tau p_\tau)=
\tilde Z(\gamma_\tau)$, norm-one pairs occur from $j$ for $(uu_{\mathsf x})_\tau^{-1}p_\tau$, and
\[
\Phi\bigl((uu_{\mathsf x})^{-1}_\tau p_\tau\bigr)_b=\Phi(u_\tau^{-1}p_{0,\tau})_b+\sum_{a<i}\tilde Z(\gamma_
\tau)_{ab}\,\frac M{s_a}\bigl(\tau(l_{p,a})+\mathsf x_{\tau,a}\bigr).
\]
\end{lemma}
\begin{proof}
Elements of $F_\tau$ do not change $\tilde Z$ ($\tilde Z(w^{-1}g)_{ab}=h(w\tilde f_a,g\tilde f_b)$) nor the
occurrence of norm-one pairs; and $\tilde Z(p_\tau)=\tilde Z(\gamma_\tau)$ by Lemma~\ref{pbs:move}. The
normalised components are $\Phi(p_\tau)_b=s_b\,\tau(\Phi(p)_b)$, and $\tilde Z_{ab}=s_as_b\tau(Z_{ab})$, so
by Lemma~\ref{pbs:colour} $\Phi(p_\tau)_b=\Phi(p_{0,\tau})_b+\sum_{a<i}\tilde Z_{ab}(M/s_a)\tau(l_{p,a})$. Apply
the shift formula for $u^{-1}$ (same $t_a(u)$ for $p$ and $p_0$, same $\tilde Z$) and then for $u_{\mathsf
x}^{-1}$.
\end{proof}

\begin{lemma}[Packing]\label{pbs:packing}
Let $Y\in G_\infty$ and let $\mathcal P\subseteq\Gamma_{c,s,c'}$ be a set such that $p^{-1}Y$ is boxed for
every $p\in\mathcal P$ and $pf_b-p'f_b\in A$ for all $p,p'\in\mathcal P$ and all $b$. Then $\sum_{p\in
\mathcal P}\beta_{c'}(p^{-1}Y)\le b^{12n^3\varrho f}\le\pbet$.
\end{lemma}
\begin{proof}
Group $\mathcal P$ by the frame $(pf_b)_b$. If $pf_b=p'f_b$ for all $b$, then by Step~2 of the proof of
Proposition~\ref{pbe:count} $p^{-1}p'\in U_{c'}$; so the elements with a given frame are of the form
$pu$ with distinct $u\in U_{c'}$, and their contribution is at most $\sum_{u\in U_{c'}}\beta_{c'}(u^{-1}
p^{-1}Y)=1$. To count the frames attach to $p$ the point $\mathsf F(p)=(Y_\tau^{-1}p\tilde f_b)_{\tau,b}$ of
$(\C^r)^{nf}$ with the norm $\max_{\tau,b}\|\cdot\|$. By Lemma~\ref{pbs:reduced}(c) (with $R=1$) the
points lie in the ball of radius $\pbsrb$, and by hypothesis in one translate of the real subspace
$\mathsf V=\bigoplus_\tau(Y_\tau^{-1}A_\tau)^n$, of real dimension at most $2\varrho nf$. If two
elements have different frames, their points are at distance at least $(N\pbsrb)^{-1}$: this is
Step~3 of the proof of Proposition~\ref{pbe:count}, with $y=Y$, $\rho=\pbsrb$, and $\pbel_\tau\le\pbsth$
in place of \eqref{pbe:eq-lenR}; the last inequality there becomes $b^{23\varrho+6n+2}(Nb^{37n})^{-1}<1$,
which holds. Comparing volumes in $\mathsf V$, the number of frames is at most $(1+4N\pbsrb^2)^{2\varrho
nf}\le(b^{6n^2})^{2\varrho nf}$, as $1+4N\pbsrb^2\le b^{134n+1}\le b^{6n^2}$. Finally $12n^3\varrho f\le
kf/4$.
\end{proof}

\paragraph*{The kernel bound}
For a place $\tau$ let $\theta=\pbsth$, $\rho=\pbsrb$, $\delta_\tau=|\tau(\delta_j(\gamma))|$, and
\[
\kappa_\gamma(\tau)=\prod_{a<i}\Bigl|\frac{s_a}M\Bigr|^{2k}\cdot\frac{(2\rho^2)^{(i+1)i}(2^{i+1}\theta^{2(j-
i)})^k}{|\det B(\tilde Z(\gamma_\tau))|^{k+i}} .
\]

\begin{lemma}\label{pbs:kernel}
$\prod_\tau\kappa_\gamma(\tau)\le\pbet\,(b^{13}/M^2)^{ikf}\,\pbsmn_j(\gamma)^{-k/2}$.
\end{lemma}
\begin{proof}
By Lemma~\ref{pbs:common}(c) and $\tilde Z_{ab}=s_as_b\tau(Z_{ab})$, $|\det B(\tilde Z)|=w^2N^{-j}\delta_\tau$
with $w=|\kappa_\tau|^{\sum_{a<j}(\varrho-a)}\ge1$; and $\prod_{a<i}|s_a/M|=w/(M\sqrt N)^i$ (the index
$a=\varrho$, present when $j=n$, contributes $\varrho-a=0$). Hence
\[
\kappa_\gamma(\tau)=\Bigl(\frac{(2\rho^2)^{i+1}N^j}{w^2\delta_\tau}\Bigr)^i\Bigl(\frac{2^{i+1}\theta^{2(j-i)}
N^{j-i}}{M^{2i}\delta_\tau}\Bigr)^k\le\Bigl(\frac{(2\rho^2)^{i+1}N^j}{\delta_\tau}\Bigr)^i\Bigl(\frac{(b^{13}/
M^2)^iN^{(j-\varrho)_+}}{\delta_\tau}\Bigr)^k,
\]
using $w\ge1$, $j-i=(j-\varrho)_+$ and $2^{i+1}\theta^{2(j-i)}\le b^{13i}$ (for $j\le\varrho$ this is $2^{i+1}
\le b^{13i}$; for $j=n$, $i=\varrho$ and $2^{\varrho+1}b^{12n+4}\le b^{13\varrho}$). Let $P=\prod_\tau
\delta_\tau$. The entries of $Z(\gamma)$ are algebraic integers ($h(f_a,\gamma f_b)=h(f_a,k_\gamma f_b)$ with
$k_\gamma\in K$), so $\delta_j(\gamma)$ is a non-zero algebraic integer and $P\ge1$; also $P^2=|\auNm\delta_j|$,
so $P/N^{(j-\varrho)_+f}=\pbsmn_j(\gamma)^{1/2}$. The product of the first factors is at most $((2\rho^2)^{
i+1}N^j)^{fi}\le b^{135n^2f\varrho}\le\pbet$ (as $2\rho^2\le b^{74n+1}$, $N\le b^{60n}$, $135n^2\varrho\le k/4$),
and the product of the second factors is $((b^{13}/M^2)^{if}N^{(j-\varrho)_+f}/P)^k=(b^{13}/M^2)^{ikf}
\pbsmn_j^{-k/2}$.
\end{proof}

\begin{proposition}[Sum of Schur integrals over an orbit]\label{pbs:orbit-sum}
$\displaystyle\sum_{\delta\in\pbsO}I^{(i)}_{c,c'}(\delta)\le\nu\,d_\infty\,\pbet^{\,2}\,(b^{13}/M^2)^{ikf}
\bigl(i_\gamma\,\pbsmn_j(\gamma)^{-1/2}\bigr)^k$.
\end{proposition}
\begin{proof}
Let $\Phi_\pbsO(y)=\sum_{p\in\pbsO}\beta^\square_{c'}(p^{-1}y)\pbsv_y(p)^{-(k+i)}$, so that the left side is
$\int\beta^\square_c\|W^\flat\|^2\Phi_\pbsO$. It is left invariant under $U^0_c$: for $w\in U^0_c$,
$\pbsv_{wy}(p)=\pbsv_y(w^{-1}p)$ by Lemma~\ref{pbs:v-props}(c), $(w^{-1}p)^{-1}y=p^{-1}wy$, and $p\mapsto
w^{-1}p$ permutes $\pbsO$. We apply Lemma~\ref{pbs:stages} with $X=V_{\mathrm{col}}$, $t(\mathsf x)=u_{
\mathsf x}$, $B=\mathsf D$ (Lemma~\ref{pbs:col}) and $C=i_\gamma^k\,\pbet\prod_\tau\kappa_\gamma(\tau)$; by
Lemma~\ref{pbs:kernel} this gives the proposition. Fix $u\in U^0_\infty$ and a boxed $y$; we must show
$\int_{\mathsf D}\Phi_\pbsO(uu_{\mathsf x}y)\,d\mathsf x\le\pbvol(\mathsf D)\,C$; as $\pbvol(\mathsf D)\ge1$ it
suffices that the integral is at most $C$, and by Lemma~\ref{pbs:colour} that $\int_{\mathsf D}\sum_{p\in
\mathcal B}(\cdots)\le\pbet\prod_\tau\kappa_\gamma(\tau)$ for each colour class $\mathcal B$ (with base point
$p_0$ and the $l_p$).

At each $\tau$ apply Theorem~\ref{pbs:oneplace} to the point $y_\tau$ (boxed: $\|y_\tau^{-1}\tilde f_a\|\le
\rho$ by Lemma~\ref{pbs:reduced}(c), $\pbel_\tau\le\theta$), $g_0=u_\tau^{-1}p_{0,\tau}$ and $Z=\tilde Z(
\gamma_\tau)$ ($\det B(Z)\ne0$ by Lemma~\ref{pbs:common}(c) and $\delta_j\ne0$; $k\ge7$). Let $\varphi_\tau$ be
the function obtained, and $\varphi(\mathsf x)=\prod_\tau\varphi_\tau\bigl(((M/s_a)\mathsf x_{\tau,a})_a\bigr)$
on $V_{\mathrm{col}}$; by the substitution $a_a=(M/s_a)\mathsf x_{\tau,a}$, $\int_{V_{\mathrm{col}}}\varphi\le
\prod_\tau\kappa_\gamma(\tau)$. By Lemma~\ref{pbs:v-props}(c), $\pbsv_{uu_{\mathsf x}y}(p)=\pbsv_y((uu_{\mathsf
x})^{-1}p)$, and by Lemma~\ref{pbs:shifted} and Theorem~\ref{pbs:oneplace},
\[
\pbsv_{uu_{\mathsf x}y}(p)^{-(k+i)}\le\varphi(\lambda(l_p)+\mathsf x),\qquad\lambda(l)=(\tau(l_a))_{\tau,a}
\in L_{\mathrm{col}} .
\]
Group $\mathcal B$ by $l_p$. For fixed $l$, the $p$ with $l_p=l$ have equal $\Phi(p)$, equal $Z(p)$ and
(Lemma~\ref{pbs:move}) the same components along $Ez$ and $A^*$ (these are read off from $Z$), so
$pf_b-p'f_b\in A$; by Lemma~\ref{pbs:packing} with $Y=uu_{\mathsf x}y$, $\sum_{l_p=l}\beta^\square_{c'}(p^{-1}
Y)\le\pbet$. Hence $\sum_{p\in\mathcal B}\beta^\square_{c'}(p^{-1}Y)\pbsv_Y(p)^{-(k+i)}\le\pbet\sum_{w\in
L_{\mathrm{col}}}\varphi(w+\mathsf x)$, and $\int_{\mathsf D}\sum_w\varphi(w+\mathsf x)\,d\mathsf x=\int_{V_{
\mathrm{col}}}\varphi$.
\end{proof}

\begin{proof}[Proof of Theorem~\ref{pbs:orbit}]
By Lemma~\ref{pbs:cs} and the Cauchy--Schwarz inequality for sums,
\[
\Bigl(\sum_{\delta\in\pbsO}\mathsf a_{c,c'}(\delta)\Bigr)^2\le\sum_{\delta\in\pbsO}I^{(i)}_{c,c'}(\delta)\cdot
\sum_{\delta\in\pbsO}I^{(i)}_{c',c}(\delta^{-1}).
\]
The elements $\delta^{-1}$ form the orbit $U^0_{c'}\gamma^{-1}U^0_c$ of $\gamma^{-1}\in\Gamma_{c',s^{-1},c}$.
Since $Z(\gamma^{-1})=Z(\gamma)^*$, $\gamma^{-1}$ has the same $\pbsmn_j$; it has cell rank $j$ (a common
vector $fx=\gamma fy$ gives $fy=\gamma^{-1}fx$, and $(1,\zeta)$ is proportional to $(\aubar\zeta,1)$); and
$i_{\gamma^{-1}}=i'_\gamma$ as $k_{\gamma^{-1}}=k_\gamma^{-1}$. Proposition~\ref{pbs:orbit-sum} for both orbits
gives
\[
\Bigl(\sum_{\delta\in\pbsO}\mathsf a_{c,c'}(\delta)\Bigr)^2\le\bigl(\nu d_\infty\pbet^2(b^{13}/M^2)^{ikf}
\bigr)^2\Bigl(\frac{i_\gamma i'_\gamma}{\pbsmn_j(\gamma)}\Bigr)^k=\bigl(\nu d_\infty\pbet^2(b^{13}/M^2)^{ikf}
\bigr)^2\Bigl(\frac{N^f}{\mathsf G_j(\gamma)}\Bigr)^k,
\]
because $\auNm\pbg\le N^f\pbsmn_j$ gives $\mathsf G_j(\gamma)=N^f\pbsmn_j/(i_\gamma i'_\gamma)$. It remains to
show that the right side is at most $\nu^2\mathsf S_j(\mathsf G_j)^2=\nu^2\pbet^4\mathcal D_j/(\pbenp M^{2ifk}
\mathsf G_j^k)$, i.e.
\[
d_\infty^2\,(b^{13}/M^2)^{2ikf}\,N^{fk}\le\mathcal D_j/(\pbenp M^{2ifk}).
\]
We use: $d_\infty\le b^{2kf}=\pbet^8$; $b^{13}/M^2\le(b^4M)^{-1}$ (as $b^{17}\le M$); $\mathcal D_j\ge\pbet^{8-
32i}\pbDid$ (for $j\le\varrho$, $\mathcal D_j=\pbet^{-16}\pbDid$ and $i\ge1$; for $j=n$, $\pbDmain\ge b^{11kf}
\pbDid=\pbet^{44}\pbDid$); $\pbDid=\#\pbR\cdot\#\mathcal C\cdot d_\infty$ and $\pbenp=\#\pbR\cdot\#\mathcal C\,
N^{-kf}$. Then the left side is at most $d_\infty\pbet^8b^{-8ikf}M^{-2ikf}N^{fk}$ and the right side at
least $\pbet^{8-32i}d_\infty N^{kf}M^{-2ifk}$; these are equal since $b^{8ikf}=\pbet^{32i}$.
\end{proof}

\subsubsection{The open cell}\label{pbs:ss-open}

For $\gamma\in\Gamma$ and a place $\tau$ let $\sigma_\tau=\prod_{a\le\varrho}s_a^2$ and $\pbg_\tau(\gamma)=
|\det\tilde Z(\gamma_\tau)|^2-|\tilde D_\varrho|^2$, $\tilde D_\varrho$ the leading minor of order $\varrho$ of
$\tilde Z(\gamma_\tau)$. As $\tilde Z_{ab}=s_as_b\tau(Z_{ab})$ and $s_\varrho^2=1/N$, $\pbg_\tau(\gamma)=\sigma_
\tau^2\,\tau(\pbg(\gamma))$. Let $\mathrm{vol}\,\auU(l)=\prod_{i<l}2\pi^{i+1}/i!$ and
\[
c_G=\frac{\mathrm{vol}\,\auU(k)}{\mathrm{vol}\,\auU(r)}\cdot\frac{2\pi^k}{(k-1)!}=\frac{2\pi^k}{(k-1)!}
\prod_{i=k}^{r-1}\frac{i!}{2\pi^{i+1}}
\]
and $\pbdeg^{\mathrm B}_\tau:=\pbdeg_\tau$, the formal
degree of \S\ref{pbe:ss-ds}. Put
\[
\pbsp_\tau(\gamma)=\Bigl(\frac{\pbsrb^{9n}}{\pbg_\tau(\gamma)}\Bigr)^{10^5n^4}\frac{d_\tau^2}{\pbdeg^{
\mathrm B}_\tau\,c_G\,(\pbg_\tau(\gamma)/2)^{k-1}} .
\]
We use the following statement, proved in \S\ref{pbo:sec} (Theorem~\ref{pbo:O1}).

\begin{itemize}
\item[(O1)] \emph{(Open-cell pairing.)} Let $\gamma\in\Gamma$ with $\pbg_\tau(\gamma)>0$ for all $\tau$, and
suppose some $y_0$ has $y_0$ and $\gamma^{-1}y_0$ boxed. Let $\phi\colon G_\infty\to[0,1]$ be measurable,
left invariant under $\prod_\tau F_\tau$, and such that $\phi(y)\ne0$ and $\pbel_\tau(y_\tau)\le\pbsth$
imply $\|y_\tau^{-1}\tilde f_a\|\le\pbsrb$ for all $a$. Then, for $W,W^\flat$ of Theorem~\ref{pbe:bessel},
\[
\Bigl\|\int_{G_\infty}\phi\,\pbsK^\flat_\gamma\Bigr\|\le\Bigl(\prod_\tau\pbsp_\tau(\gamma)\Bigr)^{1/2}+\int_{
G_\infty}\phi\,\|\pbsK_\gamma-\pbsK^\flat_\gamma\| .
\]
\end{itemize}

The second input is proved here. Let $\pbvol(\widehat\auO)=(2/\mathrm{rd}_E)^f$ (the normalisation of
Remark~\ref{pbs:rem-nu}), $c_S=(\pi^k/((k-1)!N^k))^f|D_E|^{1/4}\pbvol(\widehat\auO)^k$, $\#\pbUb=M^{f
\varrho(2k+2\varrho+1)}$, $a_*=\#\pbR\cdot\#\pbUb\cdot\#\mathcal C\cdot c_S$, $\pbdeg^{\mathrm B}=\prod_\tau
\pbdeg^{\mathrm B}_\tau=\pbdeg$, and $\Lambda$, $\nu_{\mathrm{src}}$ as in Remark~\ref{pbs:rem-nu}.

\begin{lemma}[(O2): the Siegel constant]\label{pbs:O2}
$a_*\le(1+b^{-9})\,\pbDmain\,\pbdeg\,c_G^{\,f}\,\pbvol(\widehat\auO)^{-n}\,\bigl(\nu_{\mathrm{src}}\Lambda/
d_\infty\bigr)^2$.
\end{lemma}
\begin{proof}
Write $\mathrm{rd}=\mathrm{rd}_E$, so $|D_E|=\mathrm{rd}^{2f}$ and $\pbvol(\widehat\auO)=(2/\mathrm{rd})^f$.
\emph{Step 1 (frame mass).} Put $\mathsf{fr}=N^{kf}\prod_{i=k}^{r-1}\bigl((2\pi)^{i+1}/(\mathrm{rd}^{\,i+1/2}
\,i!)\bigr)^f$. We claim
\begin{equation}\label{pbs:eq-fr}
c_S=\frac{\mathsf{fr}\cdot c_G^{\,f}\,|D_E|^{\varrho/2}}{N^{2kf}\,\pbvol(\widehat\auO)^{\varrho+1+2\varrho(
\varrho+k)}} .
\end{equation}
Both sides are positive; compare logarithms divided by $f$. With $\sum_{i=k}^{r-1}i=(2\varrho+1)(k+\varrho)$
and $\sum_{i=k}^{r-1}(i+\frac12)=(2\varrho+1)(k+\varrho+\frac12)$, the logarithm of $\mathsf{fr}\,c_G^f
|D_E|^{\varrho/2}$ divided by $f$ is
\[
k\log N-\log(k-1)!+k\log\pi+\bigl(1+(2\varrho+1)(k+\varrho)\bigr)\log2-\bigl((2\varrho+1)(k+\varrho+\tfrac12)-\varrho
\bigr)\log\mathrm{rd}
\]
(the factorials $i!$ and the powers $\pi^{i+1}$ cancel between $\mathsf{fr}$ and $c_G$), and that of
$c_SN^{2kf}\pbvol(\widehat\auO)^{\varrho+1+2\varrho(\varrho+k)}$ is
\[
k\log N+k\log\pi-\log(k-1)!+(k+\varrho+1+2\varrho(\varrho+k))\log2-(k-\tfrac12+\varrho+1+2\varrho(\varrho+k))
\log\mathrm{rd};
\]
the coefficients agree.

\emph{Step 2 (covolume).} From the definition, with $d'=\varrho(\varrho+k)$ and $M^{(2d'+\varrho)f}=\#\pbUb$,
$\nu_{\mathrm{src}}\Lambda=\#\pbUb\,\allowbreak N^{-kf}\pbvol(\widehat\auO)^{-d'}|D_E|^{\varrho/4}$. With
\eqref{pbs:eq-fr} and $\#\pbR=\#\pbUb\cdot\#\pbHb$ \eqref{pb:eq-orders},
\[
a_*=\#\pbHb\,\#\pbUb^{\,2}\,\#\mathcal C\,c_S=\#\pbHb\cdot\#\mathcal C\cdot\mathsf{fr}\cdot c_G^{\,f}\,
\pbvol(\widehat\auO)^{-n}\,(\nu_{\mathrm{src}}\Lambda)^2 ,
\]
as $n=\varrho+1$. Since $\pbDmain=(\#\mathcal C\,d_\infty)^2/(\pbdeg\,\pbmass_r)$ (\eqref{pb:eq-D},
\eqref{pb:eq-DX}), the lemma is equivalent to $\#\pbHb\cdot\mathsf{fr}\le(1+b^{-9})\,\#\mathcal C/\pbmass_r$, and
by $\#\mathcal C=\#\pbHb\,(N^k-1)^f\pbmass_k$ (proof of Lemma~\ref{pb:sizes}) to
\begin{equation}\label{pbs:eq-fr2}
\mathsf{fr}\cdot\pbmass_r\le(1+b^{-9})\,(N^k-1)^f\,\pbmass_k .
\end{equation}

\emph{Step 3 (masses).} By Lemma~\ref{pb:massratio}, for $k\le i<r$, $\pbmass_{i+1}/\pbmass_i=\bigl(\mathrm
{rd}^{\,i+1/2}i!/(2\pi)^{i+1}\bigr)^fL_i$ with $|\log L_i|\le f2^{1-i}$. Telescoping, $\mathsf{fr}\,\pbmass_r=
N^{kf}\pbmass_k\prod_iL_i$, and $\sum_i|\log L_i|\le f2^{2-k}$. Also $\log(N^k/(N^k-1))\le1/(N^k-1)\le2N^{-k}
\le2^{1-k}$. So \eqref{pbs:eq-fr2} follows from $f(2^{2-k}+2^{1-k})\le\log(1+b^{-9})$, and $\log(1+b^{-9})\ge
b^{-9}/2$. By {\cite[Lemma~13.4]{P46}}, $f\le x^{12}$ (S3), $k\ge r/2$ (S2) and $r\ge x^5$ ({\cite[Definition~13.3]{P46}}), so $k\ge x^5/2$, and $\log b\le x^{1/100}$
(S1, as $(\log b)^{100}\le x$); hence $f\,2^{3-k}\le8x^{12}2^{-x^5/2}\le\frac12e^{-9x^{1/100}}\le b^{-9}/2$.
\end{proof}

\begin{lemma}[Arithmetic of the open cell]\label{pbs:open-arith}
Assume $\nu=\nu_{\mathrm{src}}$ (Lemma~\ref{pb:covol}). Let $\gamma\in\Gamma_{c,s,c'}$ be near with $\pbg(\gamma)
\ne0$. Then $\auNm\pbg(\gamma)\ge1$, $\pbg_\tau(\gamma)>0$ for all $\tau$,
\[
\Bigl(\prod_\tau\pbsp_\tau(\gamma)\Bigr)^{1/2}\le\frac\nu2\,\mathsf S_n\bigl(\auNm\pbg(\gamma)\bigr),
\qquad\frac\nu2\,b^{-k^3f}\le\frac\nu2\,\mathsf S_n\bigl(\auNm\pbg(\gamma)\bigr).
\]
\end{lemma}
\begin{proof}
$\pbg(\gamma)\in E^+$ is a non-zero algebraic integer (the entries of $Z(\gamma)$ are integral), so
$\auNm\pbg=|\mathrm N_{E^+/\Q}\pbg|\ge1$. Fix $\tau$. By the proof of Lemma~\ref{pbe:link}(a), the Gram
matrix of the $2n$ vectors $f_0,\dots,f_\varrho,\gamma f_0,\dots,\gamma f_\varrho$ is $\bigl(\begin{smallmatrix}
\mathsf D&Z^*\\Z&\mathsf D\end{smallmatrix}\bigr)$, and exchanging the two block columns (sign $(-1)^{n^2}
=(-1)^n$) and applying the determinant formula there shows that its determinant is $(-1)^n\pbg(\gamma)$.
As $\pbg\ne0$ the $2n$ vectors span a non-degenerate subspace $\mathsf V_\tau$ of $Y_\tau$ of dimension $2n$.
It contains $N_0$, of dimension $n$, on which $h\le0$; in a non-degenerate space a subspace on which
$h\le0$ has dimension at most the negative index, so the negative index of $\mathsf V_\tau$ is $\ge n$,
hence $=n$ (that of $Y_\tau$). The determinant of a hermitian Gram matrix of a basis of $\mathsf V_\tau$
has the sign $(-1)^n$. Hence $\tau(\pbg(\gamma))>0$ and $\pbg_\tau(\gamma)=\sigma_\tau^2\tau(\pbg)>0$.

\emph{Second inequality.} $\mathsf S_n(x)=\pbet^2(\pbDmain/(\pbenp M^{2\varrho fk}x^k))^{1/2}\ge(M^{2\varrho fk}x^k)
^{-1/2}$, using $\pbet\ge1$ and $\pbenp\le\pbDid\le\pbDmain$. By Lemma~\ref{ppi:gauge}, $x=\auNm\pbg\le
b^{(121n+270n^2)f}\le b^{271n^2f}$, and $M\le b^{22}$; so $M^{2\varrho fk}x^k\le b^{(44\varrho+271n^2)kf}\le
b^{2k^3f}$.

\emph{First inequality.} We have $\prod_\tau\pbsp_\tau=C_{\mathrm{open}}/(\auNm\pbg)^{10^5n^4+k-1}\le C_{
\mathrm{open}}/(\auNm\pbg)^k$ with $C_{\mathrm{open}}=\prod_\tau(\pbsrb^{9n}/\sigma_\tau^2)^{10^5n^4}\allowbreak d_\tau^2/
(\pbdeg^{\mathrm B}_\tau c_G(\sigma_\tau^2/2)^{k-1})$. Each $s_a^2=\kappa_\tau^{2(\varrho-a)}/N$ lies in $[b^{-60n},
b^{46n}]$, so $b^{-60n^2}\le\sigma_\tau\le b^{46n^2}$ and $\Lambda_\tau:=\prod_as_a^{2(k+2(\varrho-a))}\le\sigma_\tau^k
b^{92n^3}$. Hence the factor of $C_{\mathrm{open}}$ at $\tau$ is at most
\[
b^{453n^2\cdot10^5n^4}\cdot2^k\sigma_\tau^2\frac{d_\tau^2}{\pbdeg^{\mathrm B}_\tau c_G\sigma_\tau^{2k}}\le2^kb^{
453\cdot10^5n^6+92n^2+184n^3}\frac{d_\tau^2}{\pbdeg^{\mathrm B}_\tau c_G\Lambda_\tau^2},
\]
and, as $453\cdot10^5n^6+92n^2+184n^3\le10^5n^9\le k/4$, $C_{\mathrm{open}}\le2^{kf}\pbet\,d_\infty^2/(\pbdeg^{
\mathrm B}c_G^f\Lambda^2)$. By Lemma~\ref{pbs:O2}, $d_\infty^2/(\pbdeg^{\mathrm B}c_G^f\Lambda^2)\le(1+b^{-9})\pbDmain\nu^2/(a_*
\pbvol(\widehat\auO)^n)$. Next, $2^{kf}\pbenp M^{2\varrho fk}\le\pbet\,a_*\pbvol(\widehat\auO)^n$: the factors
$\#\pbR\#\mathcal C$ and $N^{-kf}$ cancel, $M^{2\varrho fk}\le M^{f\varrho(2k+2\varrho+1)}$, and with
$x_0=200fk\,\mathrm{rd}_E$ we have $2^{kf}\le x_0^{kf}$, $(k-1)!\le x_0^k$, $\pi^k\ge1$, $|D_E|\ge1$ and
$\pbvol(\widehat\auO)=(2/\mathrm{rd}_E)^f\ge x_0^{-f}$, so it suffices that $x_0^{kf}\cdot x_0^{kf}\cdot
x_0^{(k+n)f}\le\pbet$, which follows from $x_0^{1000}\le b$ (size condition (Z3)). Altogether
\[
C_{\mathrm{open}}\le\pbet^2(1+b^{-9})\nu^2\frac{\pbDmain}{\pbenp M^{2\varrho fk}}\le\Bigl(\frac{\nu\pbet^2}2
\Bigr)^2\frac{\pbDmain}{\pbenp M^{2\varrho fk}},
\]
as $4(1+b^{-9})\le\pbet^2$; dividing by $(\auNm\pbg)^k$ gives $(\nu\mathsf S_n(\auNm\pbg)/2)^2$.
\end{proof}

\begin{lemma}[Trivial stabilisers]\label{pbs:stab}
If $\pbg(\gamma)\ne0$, $w\in U(E^+)$, $p\in U^0(E^+)$ and $p\gamma=\gamma w$, then $w=1$ (and $p=1$).
\end{lemma}
\begin{proof}
$w$ fixes the vectors $\gamma^{-1}f_b$, since $w\gamma^{-1}f_b=\gamma^{-1}pf_b=\gamma^{-1}f_b$. By
Proposition~\ref{ppi:cell}(a), $\det Z(\gamma)\ne0$, so the $\varrho\times n$ matrix $(h(e_a,\gamma^{-1}f_b))$
has rank $\varrho$ and there are $v_1,\dots,v_\varrho$ in the span of the $\gamma^{-1}f_b$ with $h(e_a,v_b)=
\delta_{ab}$; $w$ fixes them. Now $w\in U$ acts on $A$ by a unitriangular matrix $C$ and on $Y_0/A^\perp
\cong A^*$ by $(C^*)^{-1}$; the $A^*$-components of $wv_b$ and $v_b$ agree, so $C=1$ and $w$ fixes $A$.
Then $we^*_b=e^*_b+q_b+\alpha_b$ ($q_b\in X_0\oplus Ez$, $\alpha_b\in A$) and $w-1$ maps $A^\perp$ into $A$;
comparing $X_0\oplus Ez$-components of $wv_b=v_b$ gives $q_b=0$. For $x\in X_0\oplus Ez$, $wx=x+\sum_a
\lambda_a(x)e_a$ and $0=h(x,e^*_b)=h(wx,we^*_b)=h(x+\sum\lambda_a(x)e_a,e^*_b+\alpha_b)$, which equals
$\aubar{\lambda_b(x)}$ since $h(x,\alpha_b)=h(e_a,\alpha_b)=0$ and $h(e_a,e^*_b)=\delta_{ab}$; so $w$ is trivial on $X_0\oplus Ez$, and then
$wv_b=v_b$ gives $\alpha_b=0$. So $w=1$, and $p=\gamma w\gamma^{-1}=1$.
\end{proof}

\begin{proof}[Proof of Proposition~\ref{pbs:open-block}]
Let $X=T\cap\Gamma_{c,s,c'}$; $\pbg\ne0$ on $X$. In each double coset $U^0_c\delta U_{c'}$ ($\delta\in X$)
choose a representative $\mathsf r(\delta)$, near and chosen if the double coset contains a near and chosen
element. Let $\phi(\delta,y)=\mathbf 1[\delta=\mathsf r(\delta)]\,\omega(y)$. \emph{$\phi$ is admissible in
Lemma~\ref{pbs:weight}.} For $\delta\in X$ and $x\in U_c$ let $\mathsf n(x)=\#\{u'\in U_{c'}:x\delta u'^{-1}=
\mathsf r(x\delta)\}$. For each coset $U^0_cx$, $\sum_{v\in U^0_c}\mathsf n(vx)$ is the number of $(v,u')\in
U^0_c\times U_{c'}$ with $vx\delta u'^{-1}=\mathsf r(x\delta)$; there is at least one, and if $(v,u')$,
$(v_1,u'_1)$ are two, then $(v_1^{-1}v)(x\delta)=(x\delta)(u_1'^{-1}u')$, so $u'=u'_1$ and $v=v_1$ by
Lemma~\ref{pbs:stab} ($\pbg(x\delta)=\pbg(\delta)\ne0$). By Lemma~\ref{pbs:reduced}(b),
$\sum_{(x,u')}\phi(x\delta u'^{-1},xy)=\sum_x\mathsf n(x)\omega(xy)=1$.

Hence $A^\flat_T(c,s,c')=\sum_{\delta\in X,\ \delta=\mathsf r(\delta)}\int\omega\,\pbsK^\flat_\delta$. If $\delta=
\mathsf r(\delta)$ is not near and chosen, then $\omega\pbsK^\flat_\delta=0$ identically: otherwise some reduced
$y$ has $W^\flat(y)\ne0\ne W^\flat(\delta^{-1}y)$, so $y$ is boxed and $\pbel_\tau\le\pbsth$ at $\delta^{-1}y$;
choose $u\in U_{c'}$ with $\beta_{c'}(u\delta^{-1}y)\ne0$ \eqref{pbs:eq-cutoff}; then $u\delta^{-1}y$ is boxed,
$\delta u^{-1}$ is near, and the chosen element of $U^0_c\delta u^{-1}U^0_{c'}$ is near
(Lemma~\ref{pbs:near}) and lies in $U^0_c\delta U_{c'}$, contradicting the choice of $\mathsf r(\delta)$. If
$\delta$ is near and chosen, (O1) applies with $\phi=\omega$ (Lemma~\ref{pbs:reduced}(a),(c) and
Lemma~\ref{pbs:open-arith}). Summing, with Lemma~\ref{pbs:weight} for the error terms and
Lemma~\ref{pbs:open-arith} for the main terms,
\[
\|A^\flat_T(c,s,c')\|\le\sum_{\delta}\frac\nu2\mathsf S_n(\auNm\pbg(\delta))+\frac\nu2b^{-k^3f},
\]
the sum over the near and chosen $\delta\in X$. If there is no such $\delta$, the block is $0$; otherwise
the last term is at most $\frac\nu2\mathsf S_n(\auNm\pbg(\delta_0))$ for any one of them, and the
proposition follows.
\end{proof}

\subsubsection{Size conditions}\label{pbs:ss-sizes}

The proof uses the following conditions on the data, besides those of \S\ref{pbe:sec}:
\begin{itemize}
\item[(Z1)] $\varrho\ge10^5$ and $k\ge1000(2n)^9$ \eqref{ppi:eq-sizes}. Used for: $k\ge7$;
$12n^3\varrho\le k/4$ (Lemma~\ref{pbs:packing}); $135n^2\varrho\le k/4$ (Lemma~\ref{pbs:kernel});
$10^5n^9\le k/4$ (Lemma~\ref{pbs:open-arith}); $44\varrho+271n^2\le2k^2$; $2^{\varrho+1}b^{12n+4}\le
b^{13\varrho}$.
\item[(Z2)] From \S\ref{pbe:sec} and Lemma~\ref{pb:sizes}: $b\ge1000$; $b^{20}\le M\le b^{22}$ (we use
$b^{17}\le M$); $N\le b^{60n}$; $1\le|\kappa_\tau|\le b^{23}$; $d_\infty\le b^{2kf}$; $b^{11kf}\pbDid\le
\pbDmain$; and $\pbDid=\#\pbR\cdot\#\mathcal C\cdot d_\infty$ (the dimension of the space of source forms
is $\#\mathcal C\cdot d_\infty$).
\item[(Z3)] $(200\,fk\,\mathrm{rd}_E)^{1000}\le b$ (Lemma~\ref{pbs:open-arith}). It holds for the data: this is
{\cite[Lemma~13.4(S7)(d)]{P46}}.
\end{itemize}

\subsubsection{Proof of Theorem~\ref{pbs:S}}\label{pbs:ss-proof}

If $s$ does not agree, $A^\flat_T(c,s,c')=0$ by Lemma~\ref{pbs:agree}, and the sum is empty by
Definition~\ref{ppi:def-N}. Let $s$ agree. By Lemma~\ref{pbs:cells}, $T$ is the disjoint union of the
stable sets $T_j=\{\gamma\in T:j(\gamma)=j\}$, $1\le j\le n$, all of whose elements have cell rank $j$, and
$A^\flat_T=\sum_jA^\flat_{T_j}$. By Proposition~\ref{pbs:rank},
\[
\|A^\flat_T(c,s,c')\|\le\sum_j\nu\sum_{\gamma}\mathsf S_j(\mathsf G_j(\gamma))=\nu\sum_\gamma\pbesz(\gamma),
\]
the inner sums over the near and chosen $\gamma\in T_j\cap\Gamma_{c,s,c'}$; together these are the elements
of $\pbeN(c,c')\cap T\cap\Gamma_{c,s,c'}$ (they lie in $\pbeGs$, are near and chosen, and $s$ agrees), and
$\mathsf S_{j(\gamma)}(\mathsf G_{j(\gamma)}(\gamma))=\pbesz(\gamma)$ by \S\ref{pbs:ss-statement}. \qed



% ===== B-opencell.tex =====
% proofs/B-opencell.tex
% Statement (O1) of proofs/B-size.tex (the open-cell pairing bound for the cut Bessel functions).
% Labels of this file start with pbo:.
\providecommand{\auO}{\mathcal O}
\providecommand{\aubar}[1]{\overline{#1}}
\providecommand{\pbeq}{\mathsf q}
\providecommand{\pbvol}{\operatorname{vol}}
\providecommand{\pbdeg}{\operatorname{deg}}
\providecommand{\pbg}{\mathsf g}
\providecommand{\pbel}{\ell}
\providecommand{\pbsrb}{\rho_\square}
\providecommand{\pbsth}{\theta_\square}
\providecommand{\pbsK}{\mathsf K}
\providecommand{\pbsp}{\mathsf p}
% macros of this file (prefix \pbo)
\providecommand{\pboZ}{\mathsf Z}
\providecommand{\pboF}{\mathsf F}
\providecommand{\pboFd}{\mathsf F'}
\providecommand{\pboO}{\Omega}
\providecommand{\pbonu}{\nu}
\providecommand{\pboU}{U^0}
\providecommand{\pboR}{R^0}
\providecommand{\pbocG}{c_G}
\providecommand{\pboV}{\operatorname{vol}\mathrm U}
\providecommand{\pboq}{\mathsf q}
\providecommand{\pboS}{\mathcal S}
\providecommand{\pboB}{\mathsf B}
\providecommand{\pbod}{\mathsf d}

\subsection{The open-cell pairing}\label{pbo:sec}

This section proves statement \textup{(O1)} of \S\ref{pbs:ss-open}, the bound for the pairing of
the cut Bessel functions over an element $\gamma$ whose frame matrix lies in the open cell. We use
the notation of \S\S\ref{pbe:sec} and~\ref{pbs:sec}: at a real place $\tau$ the group $G_\tau$, the
normalised basis $\tilde e_a,x_i,\tilde z,\tilde e^*_a$ with Gram matrix $\mathsf J_1$, the
Euclidean norm for which it is orthonormal, the operator norm $\|g\|$, the normalised frame
$\tilde f_0,\dots,\tilde f_\varrho$, the matrix $\tilde Z(g)=(h(\tilde f_a,g\tilde f_b))_{a,b}$, the
groups $U_\tau$, $H_\tau=\mathrm U(X_{0,\tau})$, $R_\tau=U_\tau H_\tau$, the length $\pbel_\tau$, the
Bessel functions $W_\tau$, $W^\flat_\tau=\chi_\tau W_\tau$ of Proposition~\ref{pbe:local}, the formal
degree $\pbdeg_\tau$ and $d_\tau=\dim\mathbb W_\tau$. As in \S\ref{pbs:sec}, $n=\varrho+1$,
$r=k+2\varrho+1$, $\pbsrb=b^{37n}$, $\pbsth=b^{6n+2}$.

The exact statement is Theorem~\ref{pbo:O1}; it is \textup{(O1)} of \S\ref{pbs:ss-open} with the two
constants made explicit (the degree $\pbdeg^{\mathrm B}_\tau$ is the formal degree $\pbdeg_\tau$ of
\S\ref{pbe:ss-ds}, and $c_G$ is the number \eqref{pbo:eq-cG}), and with $\phi$ allowed complex valued,
$|\phi|\le1$. It is deduced from properties \textup{(BS1)--(BS5)} (Definition~\ref{pbo:def-system}), which
the Bessel functions satisfy by Theorem~\ref{pbo:system}; the latter uses Schmid's theorem on the lowest
$K$-type, Harish-Chandra's value of the Casimir operator, and the existence of $\beta\in C^\infty_c(G)$
with $\pi(\beta)v_0=v_0$. The normalisation of Haar measure as a Gram fibre measure (Lemma~\ref{pbo:GF})
uses Remark~\ref{pds:cayley}.

\subsubsection{The statement}\label{pbo:ss-statement}

Fix a real place and drop $\tau$ from the notation until \S\ref{pbo:ss-proof}. Let $G=G_\tau$, and
let $\pboF$, $\pboFd$ be the $r\times n$ matrices whose columns are $(\tilde e_1,\dots,\tilde
e_\varrho,\tilde z)$ and $(\tilde e^*_1,\dots,\tilde e^*_\varrho,x_{k_0})$, for a fixed index $k_0$.
For an $n\times n$ complex matrix $Z$ let $D_\varrho(Z)$ be its leading minor of order $\varrho$ and
\[
\pbg(Z)=|\det Z|^2-|D_\varrho(Z)|^2,\qquad \xi(Z)=(Z^{-1})_{\varrho\varrho},\qquad
\pbonu(Z)^2=1-|\xi(Z)|^2 .
\]
The \emph{open cell} is $\pboO=\{Z:\det Z\ne0,\ |\xi(Z)|<1\}$. By Cramer's rule $\xi=D_\varrho/\det Z$,
so for invertible $Z$
\begin{equation}\label{pbo:eq-gnu}
\pbg(Z)=|\det Z|^2\,\pbonu(Z)^2,\qquad\text{hence}\qquad\pboO=\{Z:\pbg(Z)>0\}.
\end{equation}
For $g\in G$, $\tilde Z(g)=\pboF^*\mathsf J_1g\pboF$; we write $\pbg(g)=\pbg(\tilde Z(g))$, which is
$\pbg_\tau(\gamma)$ of \S\ref{pbs:ss-open} for $g=\gamma_\tau$. Let $\pboU=U^0_\tau$ be the elements of
$U_\tau$ fixing every $\tilde f_a$ and $\pboR=\pboU H$ the fixer of the frame in $R_\tau$ (it is the
stabiliser $F_\tau\cap R_\tau$; in fact $F_\tau=\pboR$, Lemma~\ref{pbo:witt}). Let
\[
\pboV(j)=\prod_{i=1}^{j}\frac{2\pi^i}{(i-1)!},\qquad
\pbocG=\frac{\pboV(k)}{\pboV(r)}\cdot\frac{2\pi^k}{(k-1)!}.
\tag{\theequation}\label{pbo:eq-cG}
\]
\stepcounter{equation}
For a matrix $Z$ with $\pbg(Z)>0$, numbers $\rho\ge1$, $D>0$, $d\ge1$ put
\[
\pbsp(Z;D,\rho,d)=\Bigl(\frac{\rho^{9n}}{\pbg(Z)}\Bigr)^{10^5n^4}\frac{d^2}{D\,\pbocG\,(\pbg(Z)/2)^{k-1}}.
\]

\begin{theorem}[Open-cell pairing]\label{pbo:O1}
Assume the size conditions of \S\ref{pbo:ss-sizes}, Theorem~\ref{pbe:ds}. Let $W=\bigotimes_\tau W_\tau$ and $W^\flat=\bigotimes_\tau
W^\flat_\tau$ be as in Theorem~\ref{pbe:bessel}, and for $\gamma\in G_\infty$ let $\pbsK_\gamma(y)=
W(y)W(\gamma^{-1}y)^*$, $\pbsK^\flat_\gamma(y)=W^\flat(y)W^\flat(\gamma^{-1}y)^*$ (Hilbert--Schmidt
norm). Let $\gamma=(\gamma_\tau)\in G_\infty$ with $\pbg(\gamma_\tau)>0$ for all $\tau$, and let $\rho_\tau\ge
\pbsth$ be numbers such that some $y_0$ has $\|y_{0,\tau}^{-1}\tilde f_a\|\le\rho_\tau$ and
$\|y_{0,\tau}^{-1}\gamma_\tau\tilde f_a\|\le\rho_\tau$ for all $\tau$ and $a$. Let $\phi\colon G_\infty\to
\C$ be measurable with $|\phi|\le1$, with $\phi(uy)=\phi(y)$ for $u\in\prod_\tau\pboR_\tau$, and such
that $\phi(y)\pbsK^\flat_\gamma(y)\ne0$ implies $\|y_\tau^{-1}\tilde f_a\|\le\rho_\tau$ for all $\tau,a$.
Then $\pbsK^\flat_\gamma$ is continuous, $\phi\pbsK^\flat_\gamma$ is bounded, measurable and vanishes outside a compact set (so it is integrable), and
\[
\Bigl\|\int_{G_\infty}\phi\,\pbsK^\flat_\gamma\Bigr\|\le\Bigl(\prod_\tau\pbsp\bigl(\tilde Z(\gamma_\tau);
\pbdeg_\tau,\rho_\tau,d_\tau\bigr)\Bigr)^{1/2}+\int_{G_\infty}|\phi|\,\|\pbsK_\gamma-\pbsK^\flat_\gamma\| .
\]
\end{theorem}

\begin{remark}[Relation with \textup{(O1)} of \S\ref{pbs:ss-open}]\label{pbo:rem-O1}
\textup{(O1)} is the case $\rho_\tau=\pbsrb$, $\phi$ with values in $[0,1]$, with $\pbdeg^{\mathrm B}_\tau=
\pbdeg_\tau$ and $c_G$ as in \eqref{pbo:eq-cG}. Its hypotheses imply ours: if $y_0$ and $\gamma^{-1}y_0$
are boxed, Lemma~\ref{pbs:reduced}(c) with $R=1$ gives $\|y_{0,\tau}^{-1}\tilde f_a\|\le\pbsrb$ and
$\|(\gamma_\tau^{-1}y_{0,\tau})^{-1}\tilde f_a\|\le\pbsrb$; and if $\phi(y)\pbsK^\flat_\gamma(y)\ne0$ then
$W^\flat_\tau(y_\tau)\ne0$, so $\pbel_\tau(y_\tau)\le\pbsth$ (Proposition~\ref{pbe:local}(e)) and the
hypothesis of \textup{(O1)} on $\phi$ gives $\|y_\tau^{-1}\tilde f_a\|\le\pbsrb$. (Invariance under
$\prod F_\tau$ implies invariance under $\prod\pboR_\tau$.) The exponent $9n$ and the constants are
those of \S\ref{pbs:ss-open}, so Lemma~\ref{pbs:open-arith} applies unchanged.
\end{remark}

The proof occupies \S\S\ref{pbo:ss-system}--\ref{pbo:ss-proof}. Its structure: at one place, every
$g$ with $\pbg(g)>0$ is $s\gamma_Zt^{-1}$ with $s,t$ in the fixer of the frame and $\gamma_Z$ an
explicit smooth section of $g\mapsto\tilde Z(g)$ (\S\ref{pbo:ss-section}); Haar measure in the
coordinates $(s',Z,s)\mapsto s'\gamma_Zs$ has density at least $\pbocG\pbg(Z)^{k-1}$
(\S\ref{pbo:ss-GF}); a Sobolev inequality on a small cube of matrices near $\tilde Z(g)$, together
with Schur's orthogonality bound for $W$ and its derivatives, bounds the mean square over shifts
$t\in\pboR$ of the pairing (\S\S\ref{pbo:ss-sobolev}--\ref{pbo:ss-onplace}); the bounds multiply over
the places, one shift is no worse than the mean, and the cut is compared with $W$
(\S\ref{pbo:ss-proof}).

\subsubsection{Bessel pairs}\label{pbo:ss-system}

For a word $w=(\ell_1,\dots,\ell_m)$ in the alphabet $\mathcal E$ of pairs $(a,b)$ of basis indices
let $|w|=m$ and $(a,b)w$ the word with the letter $(a,b)$ prepended. A function $V\colon G\to\mathbb W$
into a finite-dimensional unitary representation $\mathbb W$ of $H$ is a \emph{Bessel function} if it
is continuous and bounded and $V(uhy)=\psi(u)\,h\cdot V(y)$ for $u\in U_\tau$, $h\in H$; then
$\|V(sy)\|=\|V(y)\|$ for $s\in R_\tau$, and more generally
\begin{equation}\label{pbo:eq-inv}
\Bigl\|\int_G\beta(y)V(sgy)\,dy\Bigr\|=\Bigl\|\int_G\beta(y)V(gy)\,dy\Bigr\|\qquad(s\in R_\tau,\ \beta\in
L^1(G)),
\end{equation}
because left translation by $s$ acts on the values by the unitary operator $\psi(u)h$, which
commutes with the integral. We write $L^2(U\backslash G)$ for the quotient measure of
\S\ref{pbe:ss-real}.

\begin{definition}[Bessel pair]\label{pbo:def-system}
Let $D>0$ and $b\ge1$. A \emph{Bessel pair of degree $D$} is a pair $(W,W^\flat)$ of functions $G\to
\mathbb W$, $d=\dim\mathbb W$, such that:
\begin{itemize}
\item[(BS1)] $W$ and $W^\flat$ are Bessel functions;
\item[(BS2)] $\int_{U\backslash G}\|W\|^2=d$;
\item[(BS3)] $W^\flat(y)\ne0$ implies $\pbel(y)\le b^{6n+2}$;
\item[(BS4)] (\emph{Schur bound with derivatives.}) for every $P\ge b^5$ there is a family
$(V_w)$ of Bessel functions indexed by the words, with $V_\varnothing=W$, such that (i) for every
$C^1$ curve $c$ in $G$ with $c'(t)=c(t)X$, $(V_w\circ c)'(t)=\sum_{a,b}X_{ab}V_{(a,b)w}(c(t))$; (ii)
for $|w|\le2n^2$ and every $\varphi\in L^1(G)$,
\[
\int_{U\backslash G}\Bigl\|\int_G\varphi(y)V_w(gy)\,dy\Bigr\|^2dg\le P^{2|w|}\,\frac dD\int_G|\varphi|^2 ;
\]
\item[(BS5)] $b\ge8r$.
\end{itemize}
\end{definition}


\begin{theorem}[The Bessel functions form Bessel pairs]\label{pbo:system}
Assume Theorem~\ref{pbe:ds} and the size conditions of \S\ref{pbo:ss-sizes}. For
every real place $\tau$, $(W_\tau,W^\flat_\tau)$ of Proposition~\ref{pbe:local} is a Bessel pair of
degree $\pbdeg_\tau$, with values in $\mathbb W_\tau$, $d=d_\tau$.
\end{theorem}

The proof is in \S\ref{pbo:ss-pairs}. Theorem~\ref{pbo:O1} will be deduced from the following local
statement, proved in \S\ref{pbo:ss-onplace}.

\begin{proposition}[Open-cell bound at one place]\label{pbo:local}
Let $(W,W^\flat)$ be a Bessel pair of degree $D$, $g\in G$ with $\pbg(g)>0$, $\rho\ge b$, and let some
$y$ have $\|y^{-1}\tilde f_a\|\le\rho$ and $\|y^{-1}g\tilde f_a\|\le\rho$ for all $a$. There are a compact
set $A\subseteq G$, a probability space $(T,m)$, a measurable map $t\colon T\to\pboR$ and a finite
measure $\mu$ on $G$ with $\mu(G)\le\pbsp(\tilde Z(g);D,\rho,d)$, such that
\begin{enumerate}
\item[(a)] for every bounded measurable $\varphi$ on $G$,
$\displaystyle\int_T\sum_{i,j}\Bigl|\int_A\varphi(z)\,W_i(gt(x)z)\,\aubar{W_j(t(x)z)}\,dz\Bigr|^2dm(x)\le
\int_G|\varphi|^2d\mu$, the coordinates being taken in an orthonormal basis of $\mathbb W$;
\item[(b)] for $m$-almost every $x$, $A$ contains every $z$ with $\|z^{-1}\tilde f_a\|\le\rho$ for all $a$
and $\pbel(gt(x)z)\le\rho$.
\end{enumerate}
\end{proposition}

\subsubsection{The section of the frame matrix over the open cell}\label{pbo:ss-section}

Norms of matrices are operator norms. For $x\in\C$ let $\Delta(x)=\operatorname{diag}(1,\dots,1,x)$
($n\times n$). The Gram matrix of the frame is $\pboF^*\mathsf J_1\pboF=\Delta(0)-1$, and
$\pboF^*\mathsf J_1\pboFd=\Delta(0)$; the columns of $(\pboF\,|\,\pboFd)$ are orthonormal.

\begin{definition}[The section]\label{pbo:def-section}
For $Z\in\pboO$ let $\zeta$ be the last row of $Z^{-1}$, $\pbonu=\pbonu(Z)>0$, and let $Y(Z)$ be the
$n\times n$ matrix with last row $(0,\dots,0,-1)$ and, for $a<\varrho$, entries $Y_{ab}=-c_b\,\aubar{
\zeta_a}\zeta_b$ with $c_b=\frac12$ for $b<\varrho$ and $c_\varrho=1$. Put $F^\vee(Z)=\pboF Y(Z)+\pboFd
\Delta(\pbonu)$, an $r\times n$ matrix. Let $\gamma_Z$ be the linear map that is the identity on the
orthogonal complement of the columns of $(\pboF\,|\,\pboFd)$ (spanned by the $x_i$, $i\ne k_0$) and
maps the columns of $\pboF Z^{-1}$ to those of $F^\vee(Z)$ and the columns of $F^\vee(Z^*)$ to those
of $\pboF(Z^*)^{-1}$.
\end{definition}

\begin{lemma}[The section]\label{pbo:section}
For $Z\in\pboO$: (a) $\gamma_Z$ is well defined and lies in $G$; (b) $\gamma_Z\pboF=F^\vee(Z)Z$, and
$\tilde Z(\gamma_Z)=Z$; (c) the matrix of $\gamma_Z$ is a real-analytic function of $Z$ on $\pboO$; it
is the restriction to $W=Z^*$ of a function $M(Z,W)$ holomorphic on $\pboO\times\pboO$, given by
rational expressions in the entries of $Z,Z^{-1},W^{-1}$ and in $\sigma=(1-\xi(Z)\xi(W))^{1/2}$
(principal branch; $|\xi(Z)\xi(W)|<1$), with $\sigma(Z,Z^*)=\pbonu(Z)$.
\end{lemma}
\begin{proof}
The vectors $\pboF Z^{-1}$ and $F^\vee(Z^*)$ (together $2n$ vectors) lie in the span $S$ of the columns
of $(\pboF|\pboFd)$, and so do $F^\vee(Z)$ and $\pboF(Z^*)^{-1}$. In the basis $(\pboF|\pboFd)$ of $S$
the two families are the columns of the block matrices $x=\bigl(\begin{smallmatrix}Z^{-1}&Y(Z^*)
\Delta(\pbonu)^{-1}\\0&1\end{smallmatrix}\bigr)\cdot\operatorname{diag}(1,\Delta(\pbonu))$ and $y=
\bigl(\begin{smallmatrix}Y(Z)&(Z^*)^{-1}\\\Delta(\pbonu)&0\end{smallmatrix}\bigr)$ up to the same
right factor, and $x$ is invertible (block triangular with invertible diagonal blocks), so $\gamma_Z|_S
=yx^{-1}$ is well defined. On $S$ the form is given by $h(\tilde e_a,\tilde e^*_b)=\delta_{ab}$, $h(\tilde z,\tilde z)=-1$, $h(x_{k_0},x_{k_0})=1$, all other products of basis vectors
of $S$ zero. One checks entry by entry that $y^*J_Sy=x^*J_Sx$: both are polynomial identities in the
entries of $Z^{-1}$ once $\pbonu^2=1-|\zeta_\varrho|^2$ is used. (For instance the Gram matrix of
$F^\vee(Z)$ is $Y^*(\Delta(0)-1)Y+Y^*\Delta(0)\Delta(\pbonu)+\Delta(\pbonu)\Delta(0)Y+\Delta(\pbonu)^2
\cdot e_{\varrho\varrho}$, and with the explicit $Y$ this is $-\zeta^*\zeta$, which is the Gram matrix
$(Z^{-1})^*(\Delta(0)-1)Z^{-1}$ of $\pboF Z^{-1}$; the products of $F^\vee(Z)$ with $\pboF$ are
$(\Delta(0)-1)Y+\Delta(0)\Delta(\pbonu)=1$, those of $\pboF Z^{-1}$ with $F^\vee(Z^*)$ are
$(Z^{-1})^*\cdot1$, and so on.) Hence $yx^{-1}$ is an isometry of $S$, and extended by the identity
on $S^\perp$ it is an isometry of $Y_\tau$: (a). (b) $\gamma_Z\pboF=\gamma_Z(\pboF Z^{-1})Z=F^\vee(Z)Z$,
and $\tilde Z(\gamma_Z)=\pboF^*\mathsf J_1F^\vee(Z)Z=((\Delta(0)-1)Y+\Delta(0)\Delta(\pbonu))Z=Z$, as
$(\Delta(0)-1)Y$ keeps the negated last row $(0,\dots,0,1)$ of $Y$ and $\Delta(0)\Delta(\pbonu)=
\Delta(0)$. (c) Replacing $Z^*$ by an independent $W\in\pboO$ and $\pbonu$ by $\sigma$ in the formula
$yx^{-1}$ gives $M(Z,W)$; $1-\xi(Z)\xi(W)$ lies in the disc of centre $1$ and radius $1$, where the
principal square root is holomorphic, and $\sigma(Z,Z^*)=(1-|\xi|^2)^{1/2}=\pbonu$.
\end{proof}

\begin{lemma}[Double cosets over the open cell]\label{pbo:witt}
The stabiliser $F_\tau$ of $\tilde f_0,\dots,\tilde f_\varrho$ in $G$ is $\pboR=\pboU H$. If $g\in G$ and
$Z=\tilde Z(g)\in\pboO$, there are $s,s'\in\pboR$ with $g=s\gamma_Zs'^{-1}$.
\end{lemma}
\begin{proof}
An element fixing the frame preserves $N_0^\perp$, where $N_0=\langle\tilde f_a\rangle$, hence acts on
$N_0^\perp/(N_0\cap N_0^\perp)=N_0^\perp/A_\tau\cong X_{0,\tau}$ by some $h\in H$; then $h^{-1}$ times it
fixes the frame and acts trivially on $N_0^\perp/A_\tau$, and such an element is unipotent and lies in
$U_\tau$, hence in $\pboU$. For the second claim, consider the $2n$ vectors $(\pboF|g\pboF)$ and
$(\pboF|\gamma_Z\pboF)$. Their Gram matrices are both $\bigl(\begin{smallmatrix}\Delta(0)-1&Z\\Z^*&\Delta
(0)-1\end{smallmatrix}\bigr)$ (as $g,\gamma_Z$ are isometries and $\tilde Z(g)=\tilde Z(\gamma_Z)=Z$), and
its determinant is $\pm\pbg(Z)\ne0$ (proof of Lemma~\ref{pbs:open-arith}, with $N=1$ in the normalised
basis). So both families are bases of non-degenerate subspaces, and by Witt's theorem there is
$x\in G$ with $x\pboF=\pboF$ and $xg\pboF=\gamma_Z\pboF$. Then $x\in\pboR$ and $\gamma_Z^{-1}xg\in\pboR$, and
$g=x^{-1}\gamma_Z(\gamma_Z^{-1}xg)$.
\end{proof}

\begin{lemma}[Sizes at and near a point of the open cell]\label{pbo:sizes}
Say that $Z$ is \emph{$P$-bounded} if $Z\in\pboO$ and $\|Z\|,\|Z^{-1}\|,\pbonu(Z)^{-1}\le P$.
\begin{enumerate}
\item[(a)] If $n\le M$, $|Z_{ab}|\le M^2$ for all $a,b$ and $\pbg(Z)\ge M^{-1}$, then $Z$ is
$T$-bounded for every $T\ge M^{8n}$.
\item[(a$'$)] If $n\le M$ and $|Z_{ab}|\le M^2$ for all $a,b$, then $\pbg(Z)\le|\det Z|^2\le M^{8n}$.
\item[(b)] If $Z$ is $T$-bounded with $T\ge4$ and $\|Z'-Z\|<T^{-6}$, then $Z'$ is $2T$-bounded; if also
$T\ge\max(8,n)$, then $\pbg(Z')\ge\pbg(Z)/2$.
\item[(c)] If $Z$ is $T$-bounded, $T\ge4$, then $\|\gamma_Z\|\le T^{18}$ and the real derivative of
$Z\mapsto\gamma_Z$ at $Z$ has norm at most $T^{25}$.
\item[(d)] If $s\in U_\tau$ and $\tilde Z(y)$, $\tilde Z(sy)$ are $P$-bounded, the entries of the matrix
$m(s)$ of $s$ on the frame ($s\tilde f_b=\sum_am(s)_{ab}\tilde f_a$) have modulus at most $P^2$.
\end{enumerate}
\end{lemma}
\begin{proof}
(a) Every square submatrix $B$ of $Z$ of size $l\le n$ has $\|B\|\le lM^2\le M^4$ (bound the operator
norm by $l$ times the largest entry), so $|\det B|\le M^{4l}$. From $M^{-1}\le\pbg=|\det Z|^2\pbonu^2$
and $\pbonu\le1$, $|\det Z|\ge M^{-1/2}$ and $\pbonu\ge M^{-1/2}|\det Z|^{-1}\ge M^{-1/2-4n}$; so
$\pbonu^{-1}\le M^{4n+1}$. By Cramer's rule $\|Z^{-1}\|\le|\det Z|^{-1}\cdot n\max|\text{minors}|\le
M\cdot M^2M^{4\varrho}=M^{4n-1}$. And $\|Z\|\le M^4$. (a$'$) $\pbg\le|\det Z|^2\le M^{8n}$ by the bound on $|\det B|$ with $B=Z$,
which uses only $n\le M$ and $|Z_{ab}|\le M^2$. (b) $Z'=Z(1+Z^{-1}
(Z'-Z))$ with $\|Z^{-1}(Z'-Z)\|\le T^{-5}\le\frac12$, so $Z'$ is invertible, $\|Z'^{-1}\|\le2T$ and
$\|Z'^{-1}-Z^{-1}\|=\|Z'^{-1}(Z-Z')Z^{-1}\|\le2T^{-4}$. Hence $|\xi(Z')|\le|\xi(Z)|+2T^{-4}$, and as
$\pbonu(Z)^2\ge T^{-2}$ one gets $\pbonu(Z')^2\ge(1-8T^{-2})\pbonu(Z)^2\ge\frac14T^{-2}$ (for $T\ge4$),
i.e.\ $\pbonu(Z')^{-1}\le2T$; and $\|Z'\|\le T+1\le2T$. If $T\ge\max(8,n)$, then $\pbonu(Z')^2\ge\frac78
\pbonu(Z)^2$, and from $Z=(1-(Z'-Z)Z'^{-1})Z'$, $|\det Z|\le(1+2T^{-5})^n|\det Z'|\le\frac98|\det Z'|$;
by \eqref{pbo:eq-gnu}, $\pbg(Z')\ge\frac78(\frac89)^2\pbg(Z)\ge\frac12\pbg(Z)$. (c) Let $Z$ be $P$-bounded,
$P\ge8$, and $W\in\pboO$ with $\|W^{-1}\|\le P$. Write $Y(Z,W)=\Delta(0)-1-\Delta(0)W^{-1}(1-\Delta(\frac12))
Z^{-1}\Delta(2)$, so that $Y(Z)=Y(Z,Z^*)$ (compare entries with Definition~\ref{pbo:def-section}), and
$F^\vee(Z,W)=\pboF Y(Z,W)+\pboFd\Delta(\sigma)$, $\sigma=(1-\xi(Z)\xi(W))^{1/2}$. Unwinding $yx^{-1}$ in the proof of
Lemma~\ref{pbo:section}, the matrix of $\gamma_Z$ is $M(Z,Z^*)$ with
\[
M(Z,W)=1+(F^\vee(Z,W)Z-\pboF)\pboF^*+\bigl((\pboF W^{-1}-F^\vee(Z,W)Z\,Y(W,Z))\Delta(\sigma)^{-1}-\pboFd\bigr)\pboFd^* .
\]
Now $\|\pboF\|=\|\pboFd\|=1$, $\|\Delta(x)\|=\max(1,|x|)$, $\|\Delta(0)-1\|=1$, $\|1-\Delta(\frac12)\|\le1$, $\|\Delta(2)\|=2$, so
$\|Y(Z,W)\|,\|Y(W,Z)\|\le1+2P^2$; $|\sigma|^2\le1+|\xi(Z)\xi(W)|\le2$; and $|1-\xi(Z)\xi(W)|\ge1-|\xi(Z)|\ge\frac12\pbonu(Z)^2$,
so $|\sigma^{-1}|\le\sqrt2\,\pbonu(Z)^{-1}\le2P$. Hence $\|F^\vee(Z,W)Z\|\le(1+2P^2+2)P\le3P^3$ (as $P\ge8$),
$\|(\pboF W^{-1}-F^\vee ZY(W,Z))\Delta(\sigma)^{-1}\|\le(P+3P^3(1+2P^2))\cdot2P\le2P\cdot7P^5=14P^6$, and
\[
\|M(Z,W)\|\le1+(3P^3+1)+(14P^6+1)\le15P^6\le P^{12}\qquad(P\ge8).
\] By (b), for $Z$ $T$-bounded and $T\ge4$,
every $Z'$ with $\|Z'-Z\|<T^{-6}$ is $2T$-bounded, and the same holds for $W'$ near $Z^*$; so $M$ is
holomorphic and bounded by $(2T)^{12}$ on the product of the operator-norm balls of radius $T^{-6}$ about $Z$ and $Z^*$ (Cauchy's estimate in one variable along each direction; \S\ref{pbv:ss-weak}(iii)). The Cauchy
estimate gives a derivative of norm at most $2(2T)^{12}T^6\le T^{25}$, and $\|\gamma_Z\|\le(2T)^{12}\le
T^{18}$. (d) $\tilde Z(y)=m(s)^*\tilde Z(sy)$ ($s$ maps the frame to itself by $m(s)$ and is an
isometry; this is the computation of Lemma~\ref{pbs:cells}), so $m(s)^*=\tilde Z(y)\tilde Z(sy)^{-1}$
has norm at most $P^2$.
\end{proof}

\begin{lemma}[Bounds from the length]\label{pbo:norm-z}
Let $8n\le M$, $M^{9n}\le R$, $Z\in\pboO$ with $\pbg(Z)\ge M^{-1}$ and $\|Z^{-1}\|\le R$, $s\in R_\tau$,
$v\in\pboU$, $h\in H$, and let $z\in G$ satisfy $\|z^{-1}\tilde f_a\|\le M$ for all $a$ and $\pbel(s\gamma_Z
(vh)^{-1}z)\le M$. Then $\|v^{-1}z\|\le R^5$.
\end{lemma}
\begin{proof}
Put $z_1=(vh)^{-1}z$; it has the same frame vectors $z_1^{-1}\tilde f_a=z^{-1}\tilde f_a$ (as $vh$ fixes
the frame), $\pbel(\gamma_Zz_1)\le M$ (Lemma~\ref{pbe:len}(a)) and $\|v^{-1}z\|=\|hz_1\|=\|z_1\|$. So we
may assume $s=v=h=1$. Choose $w\in U_\tau$ with $\|w\gamma_Zz\|\le M$ (the infimum defining $\pbel$ over
$R_\tau=U_\tau H$ may be taken over $U_\tau$, as $H$ preserves norms) and put $y=w\gamma_Zz$, $Z'=\tilde Z(
w\gamma_Z)=m(w)^{-*}Z$ (Lemma~\ref{pbo:sizes}(d)); $\pbg(Z')=\pbg(Z)$, as $m(w)$ is unitriangular. The
entries $Z'_{ab}=h(y^{-1}\tilde f_a,y^{-1}w\gamma_Z\tilde f_b)=h(y^{-1}\tilde f_a,z^{-1}\tilde f_b)$ have
modulus $\le M^2$, so by Lemma~\ref{pbo:sizes}(a) $\|Z'^{-1}\|\le M^{8n}\le R$. Reading off the
column at $\tilde e^*_a$ of the identity $z^{-1}=(y^{-1}w)\gamma_Z$, expanded with the formula for
$\gamma_Z$ (only the second block column of $yx^{-1}$ contributes, and $m(w)$ cancels against $Z$),
\[
z^{-1}\tilde e^*_a=\sum_{i\le\varrho}\bigl((Z'^*)^{-1}\bigr)_{ia}\,y^{-1}\tilde f_i+\sum_{i<\varrho}\tfrac12
(Z^{-1})_{i\varrho}\bigl((Z^*)^{-1}\bigr)_{\varrho a}\,z^{-1}\tilde f_i .
\]
Hence $\|z^{-1}\tilde e^*_a\|\le nR\cdot M+\varrho R^2M/2\le R^4$, and $\|z^{-1}\tilde f_a\|\le M\le R^4$.
Finally, for $A=z^{-1}$, the relation $A\mathsf J_1A^*=\mathsf J_1$ gives $\|A\|^2=\|\mathsf J_1+A(1-\mathsf
J_1)A^*\|\le1+\|A(1-\mathsf J_1)\|\,\|A\|$, and the columns of $1-\mathsf J_1$ are $0$, $2\tilde z$ and
$\pm(\tilde e_a-\tilde e^*_a)$, so $\|A(1-\mathsf J_1)\|\le(4\varrho+2)R^4$; as $\|A\|\ge1$, $\|z\|=\|A\|
\le1+(4\varrho+2)R^4\le R^5$.
\end{proof}

\subsubsection{Haar measure in the coordinates of the open cell}\label{pbo:ss-GF}

On complex matrices we use Lebesgue measure for the real and imaginary parts of the entries, and on
hermitian $m\times m$ matrices $T$ Lebesgue measure $dT$ in the real coordinates $T_{aa}$, $\Rel T_{ab}$,
$\Iml T_{ab}$ ($a<b$). For a tuple $E$ of $m$ vectors of $Y_\tau$ (an $r\times m$ matrix) let $\Phi(E)=
E^*\mathsf J_1E$. On the open set of tuples with $\Phi$ a submersion (e.g.\ $E$ of rank $m$ spanning a
non-degenerate subspace) Lebesgue measure disintegrates as $dE=\int\lambda_T\,dT$ with fibre measures
$\lambda_T$ on $\{\Phi=T\}$ (coarea formula). For a measure $\mu$ on $G$ and a tuple $E_0$ with
$\Phi(E_0)=T_0$, write $\mu\preceq\lambda_{E_0}$ if the image of $\mu$ under $g\mapsto gE_0$ is at most
$\lambda_{T_0}$, in the weak sense: $\int K(gE_0)\,d\mu(g)\le C$ whenever $K\ge0$ is lower semicontinuous
and $\int_{\Phi(E)\in S}K\,dE\le C\,\pbvol(S)$ for all open $S$ in a neighbourhood of $T_0$.

\begin{lemma}[Gram fibre normalisation \textup{(GF)}]\label{pbo:GF}
Let $dg$ be the Haar measure of \textup{(DS1)(iii)} and $\sigma$ a real-analytic section of $\Phi$ near
$\mathsf J_1$ with $\sigma(\mathsf J_1)=1$ and invertible values (e.g.\ $\sigma(T)=(\mathsf J_1T)^{1/2}$, the
power series square root near $1$). Then $(g,T)\mapsto g\sigma(T)$ carries $\frac{\pboV(r)}{2^r}\,dg\,dT$ to
Lebesgue measure $dY$ on a neighbourhood of $G=\{\Phi=\mathsf J_1\}$. Consequently every measure
$\mu\preceq\lambda_1$ satisfies $\mu\le\frac{\pboV(r)}{2^r}dg$.
\end{lemma}
\begin{proof}
\emph{Reduction to a constant.} The map $(g,T)\mapsto g\sigma(T)$ is a local diffeomorphism near
$G\times\{\mathsf J_1\}$ (its differential is computed below and is bijective) and commutes with
$Y\mapsto g_0Y$, $g_0\in G$, which preserves $dY$ ($|\det g_0|=1$) and $dg\,dT$. Hence the image of
$dg\,dT$ is $c(T)^{-1}dY$ for a function $c$ of $T$ alone. It does not depend on $T$ either: for
$A\in GL_r(\C)$ near $1$ and $T$ near $\mathsf J_1$ put $g'(T)=\sigma(T)A\,\sigma(A^*TA)^{-1}$; since
$\Phi(\sigma(T)A)=A^*TA=\Phi(\sigma(A^*TA))$, we get $\Phi(g'(T))=\mathsf J_1$, i.e.\ $g'(T)\in G$, and
$g\sigma(T)A=gg'(T)\,\sigma(A^*TA)$. Thus $Y\mapsto YA$ corresponds to $(g,T)\mapsto(gg'(T),A^*TA)$. The
first map multiplies $dY$ by $|\det A|^{2r}$; the second multiplies $dg\,dT$ by $|\det A|^{2r}$, because
it is triangular, right translation by $g'(T)$ preserves $dg$ ($G$ is unimodular) and $T\mapsto A^*TA$
multiplies $dT$ by $|\det A|^{2r}$. Hence $c(T)=c(A^*TA)$, and every $T$ near $\mathsf J_1$ is
$A^*\mathsf J_1A$ with $A=\sigma(T)$ near $1$. So $c$ is a constant $C$, the Jacobian at $(1,\mathsf J_1)$. Writing $T=\mathsf J_1+H$,
the differential there is $(X,H)\mapsto X+\frac12\mathsf J_1H$ on $\mathfrak g_0\times\mathrm{Herm}_r$, because
$\Phi(1+X+\frac12\mathsf J_1H)=\mathsf J_1+H+O(2)$ ($X^*\mathsf J_1+\mathsf J_1X=0$, $\mathsf J_1^2=1$). With
$X_0=\mathsf J_1X\in\mathfrak u(r)$ it is $\mathsf J_1\circ\bigl((X_0,H)\mapsto X_0+\frac12H\bigr)$, and $\mathsf J_1$ is
unitary. By Remark~\ref{pds:cayley}, the density of $dg$ at $1$ in the coordinate $X_0$ equals that of the
Haar probability measure of $\mathrm U(r)$ at $1$ in the coordinate $X_0$ (the Cayley density at $0$ divided
by the same normalising constant, and $X\mapsto\mathsf J_1X$ has Jacobian $1$). So $C$ is the constant of the
case $\mathsf J_1=1$, $G=\mathrm U(r)$, $dg$ of mass one.

\emph{The compact case.} For $\mathsf J_1=1$ the same argument applies at every point of
$\mathrm U(r)\times\{T>0\}$ with the global diffeomorphism $(g,T)\mapsto gT^{1/2}$ onto the invertible
matrices (polar decomposition). For $A\in GL_r(\C)$, $Y\mapsto YA$ multiplies $dY$ by $|\det A|^{2r}$ and
$T\mapsto A^*TA$ multiplies $dT$ by $|\det A|^{2r}$; and $gT^{1/2}A=gu_0(A^*TA)^{1/2}$ with $u_0=T^{1/2}A(A^*TA)^{-1/2}$
unitary, so the fibre over $T$ is carried to the fibre over $A^*TA$ by $g\mapsto gu_0$, which preserves $dg$.
Hence the density of $dY$ against $dg\,dT$ is the same at every $T>0$, namely $C$. So $dY=C\,dg\,dT$ globally. Integrate $e^{-\operatorname{tr}Y^*Y}$: $\pi^{r^2}=C\int_{T>0}
e^{-\operatorname{tr}T}dT=C\,\pi^{r(r-1)/2}\prod_{j=1}^r(j-1)!$ (complex Siegel gamma integral at $s=r$:
$\int_{T>0}e^{-\operatorname{tr}T}(\det T)^{s-r}dT=\pi^{r(r-1)/2}\prod_{j=1}^r\Gamma(s-j+1)$, proved by induction on
$r$ with $T=\bigl(\begin{smallmatrix}T'&v\\v^*&t\end{smallmatrix}\bigr)$, $t=t_1+v^*T'^{-1}v$, $\det T=t_1\det T'$, and
the Gaussian integral in $v$ giving $\pi^{r-1}\det T'$). Hence $C=\prod_{j=1}^r\pi^j/(j-1)!=\pboV(r)/2^r$.

\emph{The consequence.} If $\mu\preceq\lambda_1$ and $O\subseteq G$ is open, test with $K$ the indicator
of the open set $\{g\sigma(T):g\in O,T\in S_0\}$ for a small open $S_0\ni\mathsf J_1$; for open $S\subseteq S_0$,
$\int_{\Phi\in S}K\,dY=C\pbvol(O)\pbvol(S)$, so $\mu(O)\le C\pbvol(O)$; outer regularity of $dg$ gives
$\mu\le C\,dg$.
\end{proof}


\begin{lemma}[Appending vectors]\label{pbo:append}
\begin{enumerate}
\item[(a)] Let $E=(E_1|E_2)$ be a tuple split into two parts, let $\mu\preceq\lambda_{E_1}$, and let
$\nu$ be a measure on the stabiliser $G_{E_1}$ of $E_1$ such that the image of $\nu$ under $u\mapsto
uE_2$ is at most the fibre measure, over $(E_1^*\mathsf J_1E_2,\Phi(E_2))$, of the map $X\mapsto(E_1^*
\mathsf J_1X,\Phi(X))$ on tuples $X$ of the size of $E_2$. Then $\mu*\nu\preceq\lambda_E$.
\item[(b)] Let $\mu\preceq\lambda_{\pboF}$ and let $ds=du^0\,dh$ be the measure on $\pboR$, $du^0$ the
Haar measure on $\pboU$ that is the image of Lebesgue measure on the parameters $q=(\beta,\sigma)$
($\beta\in M_{\varrho\times k}(\C)$, $\sigma$ real skew-hermitian coordinates)  $q$ of
$u_q$ in \eqref{pbo:eq-uq} below, and $dh$ of mass one. Then $\mu*ds\le\frac{\pboV(r)}{2^{\varrho+1}\pboV(k)}\,dg$.
\end{enumerate}
\end{lemma}
\begin{proof}
(a) For $K$ as in the definition, disintegrate $dE=dE_1\,dE_2$ first along $\Phi(E_1)$ and then, for
fixed $E_1$, along $(E_1^*\mathsf J_1E_2,\Phi(E_2))$; the Gram matrix of $E$ is determined by these
blocks. Moving $E_1$ to $gE_1$ by $g$ (in the support of $\mu$) and using $\nu$ on its stabiliser,
$\int K(gE)\,d(\mu*\nu)(g)=\int\int K(guE)\,d\nu(u)\,d\mu(g)$, and the inner integral is bounded through
the hypothesis on $\nu$ by the integral of $K(g\,\cdot)$ over the fibre in the second variable, then
the outer one through $\mu\preceq\lambda_{E_1}$. (b) Append to $\pboF$ first the vectors $\tilde
e^*_1,\dots,\tilde e^*_\varrho$, then $x_1,\dots,x_k$. The tuples $X=(\tilde e^*_a)+\pboF'' p$, with $\pboF''$
the columns $\tilde e_a,x_i$ (orthonormal, $h$-orthogonal to $\pboF$), are exactly those with $\pboF^*
\mathsf J_1X=\pboF^*\mathsf J_1(\tilde e^*_a)$; writing such a tuple as $u_q(\tilde e^*+\tilde e\,T/2)$
with $T$ its Gram matrix, Lebesgue measure is $dp=2^{-\varrho}dT\,dq$ (translate the $\tilde e$-part
by the $q$-dependent quadratic term and split a complex matrix into hermitian and skew-hermitian parts,
which costs $2^{-\varrho}$ in the diagonal; this proves the formula, which is quoted from here in \eqref{pbo:eq-dual} below for the
first $\varrho$ vectors). So (a) applies with $\nu=2^{-\varrho}du^0$. For the last $k$ vectors the stabiliser is $H$,
and the image of $\frac{\pboV(k)}{2^k}dh$ under $h\mapsto(hx_i)$ is the fibre measure of $X\mapsto
(X^*X)$ on $k$-tuples of $X_0$ over $1$ (\textup{(GF)} for $\mathrm U(k)$); tuples of $Y_\tau$
orthogonal to the first $2\varrho+1$ vectors are tuples in $X_0$, and the linear part of the map costs
nothing (orthonormal coordinates). By (a) twice, $\mu*(2^{-\varrho}du^0)*(\frac{\pboV(k)}{2^k}dh)\preceq
\lambda_1$, so by \textup{(GF)} $\mu*ds\le2^\varrho\frac{2^k}{\pboV(k)}\frac{\pboV(r)}{2^r}dg$, and $r=k+
2\varrho+1$.
\end{proof}

We make the parametrisation of $\pboU$ explicit. For $\beta\in M_{\varrho\times k}(\C)$ and real
skew-hermitian $\sigma\in M_\varrho(\C)$ let $u_q$ ($q=(\beta,\sigma)$) fix $\tilde e_a$ and $\tilde z$, map
$x\mapsto x+\tilde e\,\beta$ (that is, $x_i\mapsto x_i+\sum_a\beta_{ai}\tilde e_a$), and
\begin{equation}\label{pbo:eq-uq}
u_q(\tilde e^*+\tilde e\,a)=\tilde e^*+\tilde e\,(\sigma+a-\tfrac12\beta\beta^*)-x\,\beta^*,
\end{equation}
in matrix notation for tuples ($\tilde e$, $x$, $\tilde e^*$ the rows of basis vectors). One checks directly
that $u_q$ preserves $h$ (the only non-trivial products are $h(u_qx_i,u_q\tilde e^*_b)=\aubar{\beta_{bi}}-\aubar{\beta_{bi}}=0$
and the Gram matrix $(\sigma+\sigma^*)-\beta\beta^*+\beta\beta^*=0$ of $u_q\tilde e^*$). Every element of $\pboU$ is a unique $u_q$, and $q\mapsto u_q$ carries
Lebesgue measure to a Haar measure $du^0$ (it is a two-step nilpotent group in exponential-type
coordinates). We fix this $du^0$; nothing below depends on the choice except through quantities in
which it cancels.

\begin{proposition}[Mean value inequality]\label{pbo:meanvalue}
Assume \textup{(GF)}. For measurable $f\ge0$ on $G$,
\[
\int_{\pboO}\pbocG\,\pbg(Z)^{k-1}\int_{\pboR}\int_{\pboR}f(s'\gamma_Zs)\,ds\,ds'\,dZ\le\int_Gf(g)\,dg .
\]
\end{proposition}
\begin{proof}
\emph{Step 1: tuples dual to the frame.} An $n$-tuple $E$ with $\pboF^*\mathsf J_1E=Z$ invertible is
$E=E'Z$ with $\pboF^*\mathsf J_1E'=1$; such $E'$ are $E'=\mathsf J_1\pboF+\pboF''p$ with $p$ a
$(\varrho+k)\times n$ matrix, and $T'=\Phi(E')=(Z^{-1})^*\Phi(E)Z^{-1}$. For a hermitian $T'$ with
$1+T'_{\varrho\varrho}>0$ let $E'(T')$ be the tuple with columns $\tilde e^*_b+\frac12\sum_aT'_{ab}\tilde
e_a$ ($b<\varrho$) and $-\tilde z+\sum_aT'_{a\varrho}\tilde e_a+(1+T'_{\varrho\varrho})^{1/2}x_{k_0}$; it is dual
to the frame with Gram matrix $T'$. Every dual tuple is $s\,E'(T')$ with $s\in\pboR$: the first $\varrho$
vectors are moved by $\pboU$ (as in Lemma~\ref{pbo:append}(b)), and the $X_0$-part of the last vector,
of squared length $1+T'_{\varrho\varrho}$, by $H$ (polar coordinates in $X_0\cong\C^k$). In these
coordinates
\begin{equation}\label{pbo:eq-dual}
dp=2^{-\varrho}\,\mathsf m_k(1+T'_{\varrho\varrho})\,dT'\,ds,\qquad\mathsf m_k(t)=\frac{\pi^kt^{k-1}}{(k-1)!}
\ (t>0),
\end{equation}
where $\mathsf m_k(t)$ is the mass of the sphere $\{\|x\|^2=t\}$ of $\C^k$ for $dx/d\|x\|^2$ (polar
coordinates: $dx=\mathsf m_k(t)\,dt\,dh$ with $dh$ of mass one), and the factor $2^{-\varrho}$ comes from
the first $\varrho$ vectors, by the formula $dp=2^{-\varrho}dT\,dq$ established in the proof of Lemma~\ref{pbo:append}(b). The substitution $E=E'Z$ multiplies $dp$ by $|\det Z|^{2(\varrho+k)}$ and
$T=Z^*T'Z$ multiplies $dT'$ by $|\det Z|^{-2n}$; so on the tuples over $Z$,
\[
dp_E=|\det Z|^{2(k-1)}2^{-\varrho}\mathsf m_k\bigl(1+((Z^{-1})^*TZ^{-1})_{\varrho\varrho}\bigr)\,dT\,ds .
\]
\emph{Step 2: the fibre over the frame.} At $T=\Phi(\pboF)=\Delta(0)-1$ we have $1+((Z^{-1})^*T Z^{-1})_{
\varrho\varrho}=1-|\xi(Z)|^2=\pbonu(Z)^2$, and the tuple $E'((Z^{-1})^*TZ^{-1})Z$ is $\gamma_Z\pboF$ (compare
coefficients with Definition~\ref{pbo:def-section}: the coefficients of $\tilde e_a$, $x_{k_0}$, $\tilde
z$, $\tilde e^*_a$ are those of $Y(Z)$ and $\Delta(\pbonu)$). By \eqref{pbo:eq-gnu},
$|\det Z|^{2(k-1)}2^{-\varrho}\mathsf m_k(\pbonu^2)=2^{-\varrho-1}\frac{2\pi^k}{(k-1)!}\pbg(Z)^{k-1}$. As the
tuples $\mathsf J_1\pboF Z+\pboF''p$ ($Z\in\pboO$) are a subset of all tuples and $(Z,p)\mapsto\mathsf J_1
\pboF Z+\pboF''p$ preserves Lebesgue measure, the measure $\mu_0$, image of $2^{-\varrho-1}\frac{2\pi^k}{(
k-1)!}\pbg(Z)^{k-1}\mathbf 1_{\pboO}(Z)\,dZ\,ds'$ under $(Z,s')\mapsto s'\gamma_Z$, satisfies $\mu_0\preceq
\lambda_{\pboF}$ (the coordinates $(T,Z,s')$ are continuous at $T_0$ and the density is continuous
there, which is what the weak definition needs).

\emph{Step 3.} By Lemma~\ref{pbo:append}(b), $\mu_0*ds\le\frac{\pboV(r)}{2^{\varrho+1}\pboV(k)}dg$.
Integrating $f$ and dividing by the constant gives the claim, since $\frac{2^{\varrho+1}\pboV(k)}{\pboV(r)}
\cdot2^{-\varrho-1}\frac{2\pi^k}{(k-1)!}=\pbocG$.
\end{proof}

\subsubsection{Unfolding, Sobolev's inequality, and the bound at the section}\label{pbo:ss-sobolev}

\begin{lemma}[Unfolding along $U_\tau$]\label{pbo:unfold}
Let $w,\varphi\ge0$ be measurable on $G$, $w$ left invariant under $R_\tau$, and $K,R\ge1$ such that
$\int_{\pboU}\varphi(uy)\,du^0\le K$ for all $y$, and such that $\varphi(y)\ne0\ne\varphi(sy)$ with $s\in U_\tau$
implies $|m(s)_{ab}|\le R$ for all $a,b$. Then $\int_G\varphi w\,dg\le K(2R)^{2n^2}\int_{U\backslash G}w$,
for the Haar measure $du=du^0\,dx$ on $U_\tau$, $dx$ Lebesgue measure in the strictly upper triangular
entries $x$ of $m(s)$.
\end{lemma}
\begin{proof}
The map $s\mapsto m(s)$ is a homomorphism from $U_\tau$ onto the unitriangular group of the frame with
kernel $\pboU$, so every $s\in U_\tau$ is $s=u\,\ell(x)$ with $u\in\pboU$ and $\ell$ a measurable section
with $m(\ell(x))=1+x$. The measure $du^0dx$ is the measure $du$ of convention \textup{(Cv3)}: the
\textup{(Cv3)} coordinates of $s$ in the frame columns are the entries of $m(s)-1$, and on $\pboU$ the
other \textup{(Cv3)} coordinates are $\beta$, the entries $(\sigma-\frac12\beta\beta^*)_{ab}$ ($a<b$) and
$\Iml\sigma_{aa}$, which differ from $q$ by a translation depending on $\beta$ (so Lebesgue measure is
unchanged); by Weil's formula for the normal subgroup $\pboU$ of $U_\tau$ the two Haar measures
coincide (the measure of (Cv3) is a Haar measure by Lemma~\ref{pbv:haar}, and the two have the same
density at $1$ in these coordinates). By unfolding, $\int_G\varphi w=\int_{U\backslash G}w(y)\int_{U_\tau}\varphi(sy)\,ds\,dy$.
Fix $y$ and some $s_0$ with $\varphi(s_0y)\ne0$ (if none, the inner integral is $0$); substituting $s\mapsto
ss_0$, the inner integral is $\int dx\int_{\pboU}\varphi(u\ell(x)s_0y)\,du^0$, whose inner integral is $\le K$
and vanishes unless $|x_{ab}|\le R$ (apply the hypothesis to $s_0y$ and $u\ell(x)$, whose frame matrix is
$1+x$). The box of strictly upper triangular complex $x$ with $|x_{ab}|\le R$ has volume $\le(\pi R^2)^{
n^2/2}\le(2R)^{2n^2}$.
\end{proof}

\begin{lemma}[Integral over the section]\label{pbo:section-int}
Assume \textup{(GF)}. Let $P\ge1$, $c>0$, $Q\subseteq\pboO$ measurable with $\pbg\ge c$ and every
$Z\in Q$ $P$-bounded, and $w\ge0$ measurable and left invariant under $R_\tau$. Then
\[
\int_Q\int_{H\times\pboU}w\bigl(\gamma_Z(vh)^{-1}\bigr)\,\pbocG c^{k-1}\,dh\,dv\,dZ\le(2P^2)^{2n^2}\int_{U
\backslash G}w .
\]
\end{lemma}
\begin{proof}
Let $\chi\ge0$ be a cutoff for $\pboU$ ($\int_{\pboU}\chi(uy)du^0=1$ for all $y$) and $\chi_Q=\chi\cdot\mathbf
1[\tilde Z(\cdot)\in Q]$. By Lemma~\ref{pbo:unfold} with $K=1$, $R=P^2$ and Lemma~\ref{pbo:sizes}(d),
$\int_G\chi_Qw\le(2P^2)^{2n^2}\int_{U\backslash G}w$. By Proposition~\ref{pbo:meanvalue} applied to $\chi_Qw$
and restricted to $Q$, $\int_G\chi_Qw\ge\int_Q\pbocG\pbg^{k-1}\int\int\chi_Q(s'\gamma_Zs)w(s'\gamma_Zs)\,ds\,
ds'\,dZ$. For $s,s'\in\pboR$, $\tilde Z(s'\gamma_Zs)=Z$ and $w(s'\gamma_Zs)=w(\gamma_Zs)$; integrating $\chi$
over $s'=uh$ gives $1$ (the measure $ds$ is invariant under inversion, as $(uh)^{-1}=(h^{-1}u^{-1}h)h^{-1}$
and conjugation by $H$ preserves $du^0$). So the right side is $\int_Q\pbocG\pbg^{k-1}\int w(\gamma_Z(vh)^{-1})
\,dh\,dv\,dZ$, and $\pbg\ge c$ on $Q$.
\end{proof}

\begin{lemma}[Sobolev's inequality on a cube]\label{pbo:sobolev}
Let $f(x,\omega)$ be functions of $x\in[0,\varepsilon]^N$, $\varepsilon\le1$, and of $\omega$ in a measure space,
with values in a finite-dimensional Hilbert space, indexed by words $w$ in a finite alphabet $L$, such
that for each coordinate $j$, each $w$, $x$ and $\omega$, $t\mapsto f_w(x+te_j,\omega)$ is differentiable on
the cube with derivative of norm $\le\Lambda\sum_{l\in L}\|f_{lw}(x+te_j,\omega)\|$. Then
\[
\int\|f_\varnothing(0,\omega)\|^2d\omega\le a^N\max_{|w|\le N}\int_{[0,\varepsilon]^N}\int\|f_w(x,\omega)\|^2\,d\omega\,
dx,\qquad a=\frac2\varepsilon\bigl(1+(|L|\Lambda)^2\bigr).
\]
\end{lemma}
\begin{proof}
For a function $g$ on $[0,\varepsilon]$ with values in a Hilbert space, differentiable, $g(0)=g(t)-\int_0^tg'$ gives
$\|g(0)\|^2\le2\|g(t)\|^2+2\varepsilon\int_0^\varepsilon\|g'\|^2$, and averaging over $t$, $\|g(0)\|^2\le\frac2\varepsilon
\int_0^\varepsilon\|g\|^2+2\varepsilon\int_0^\varepsilon\|g'\|^2$. If $\|g'\|\le\Lambda\sum_{l\in L}\|g_l\|$, then $\|g'\|^2\le
\Lambda^2|L|\sum_l\|g_l\|^2$. For $0\le j\le N$ and a word $w$ let $I_j(w)$ be the integral of $\|f_w(x,\omega)\|^2$
over $\omega$ and over $x$ with $x_1,\dots,x_j\in[0,\varepsilon]$ and $x_{j+1}=\dots=x_N=0$. Applying the one-variable
inequality in $x_{j+1}$ with the other variables fixed and integrating,
\[
I_j(w)\le\frac2\varepsilon I_{j+1}(w)+2\varepsilon\Lambda^2|L|\sum_{l\in L}I_{j+1}(lw).
\]
Iterating from $I_0(\varnothing)$, each of the $N$ steps either keeps the word (factor $2/\varepsilon$) or prepends
one of $|L|$ letters (factor $2\varepsilon\Lambda^2|L|$ each); all words reached have length $\le N$, so with
$K=\max_{|w|\le N}I_N(w)$,
\[
I_0(\varnothing)\le\sum_{m=0}^N\binom Nm\Bigl(\frac2\varepsilon\Bigr)^{N-m}(2\varepsilon\Lambda^2|L|^2)^mK=\Bigl(\frac2\varepsilon+
2\varepsilon(|L|\Lambda)^2\Bigr)^NK\le a^NK,
\]
as $\varepsilon\le1$.
\end{proof}

\begin{theorem}[$L^2$ bound at the section]\label{pbo:schur-section}
Assume \textup{(GF)}. Let $(W,W^\flat)$ be a Bessel pair of degree $D$, $P\ge\max(8r,b^5)$, $c>0$, and
$Z$ such that every $Z'$ with $\|Z'-Z\|<P^{-6}$ is $P$-bounded with $\pbg(Z')\ge c$. Let $\varphi\in
L^1(G)$ vanish where $\|y\|>P^{20}$. Then, with $\mathcal B=\{(h,v)\in H\times\pboU:\|v\|\le P^{20}\}$,
\[
\int_{\mathcal B}\Bigl\|\int_G\varphi(y)W\bigl(\gamma_Z(vh)^{-1}y\bigr)dy\Bigr\|^2\pbocG c^{k-1}dh\,dv\le
P^{2000n^2}\frac dD\int_G|\varphi|^2 .
\]
\end{theorem}
\begin{proof}
Take $(V_w)$ from (BS4) for this $P$. For words $i$ in the alphabet $\mathcal E^{(2)}$ of letters
$((a,b),c,d)$ define weights $\psi_\varnothing=\varphi$, $\psi_{((a,b),c,d)i}(y)=\psi_i(y)\,(y^{-1})_{ad}\,y_{cb}/
P^{40}$; then $|\psi_i|\le|\varphi|$ since $|y_{cb}|,|(y^{-1})_{ad}|\le\|y\|\le P^{20}$ on the support. Put
$\mathcal W_i(g)=\int_G\psi_i(y)V_{\bar i}(gy)\,dy$, $\bar i$ the word of first components. If $c$ is a $C^1$
curve with $c'=cX$, then $(c(t)y)'=c(t)y\,(y^{-1}Xy)$, so by (BS4)(i) and differentiation under the
integral (the $V_w$ are bounded, $\psi_i$ integrable),
\[
\frac{d}{dt}\mathcal W_i(c(t))=\sum_{a,b,c,d}P^{40}X_{dc}\,\mathcal W_{((a,b),c,d)i}(c(t)),
\]
since $(y^{-1}Xy)_{ab}=\sum_{c,d}(y^{-1})_{ad}X_{dc}y_{cb}$. Hence $\|\frac d{dt}\mathcal W_i(c(t))\|\le P^{40}
\|X\|\sum_\ell\|\mathcal W_{\ell i}\|$. Choose real linear coordinates $x\mapsto Z+e(x)$ on $M_n(\C)$
($N=2n^2$ coordinates, the real and imaginary parts of the entries), measure preserving, with $\|e(x)\|
\le2n^2\|x\|_\infty$; with $\varepsilon=P^{-8}$ the cube maps into the ball of radius $P^{-6}/8$ about $Z$, and
$4\|e(e_j)\|\le P^2$. For $(h,v)\in\mathcal B$ and $x$ in the cube, $Z'=Z+e(x)$ is $P$-bounded, so by
Lemma~\ref{pbo:sizes}(c) (with $T=P$; the ball hypothesis gives the bounds at every point of the segment)
the curve $t\mapsto\gamma_{Z+e(x+te_j)}(vh)^{-1}$ has logarithmic derivative of norm $\le\|\gamma_{Z'}^{-1}\|\,
\|(vh)^{-1}\|\,P^{25}\cdot\frac{P^2}4\cdot\|vh\|\le P^{18+20+25+2+20}=P^{85}$. Lemma~\ref{pbo:sobolev} with
$f_i(x,(h,v))=\mathcal W_i(\gamma_{Z+e(x)}(vh)^{-1})$, $\Lambda=P^{40}P^{85}$, $|L|=r^4\le P^4$, gives the
factor $a^N$ with $a(2P^2)\le\frac2{P^{-8}}(1+P^{258})2P^2\le P^{275}$. For each word $i$ with $|i|\le2n^2$,
by Lemma~\ref{pbo:section-int} with $Q$ the ball of radius $P^{-6}$ about $Z$ and $w=\|\mathcal W_i\|^2$
(left invariant under $R_\tau$ by \eqref{pbo:eq-inv}), and by (BS4)(ii),
\[
\int_{\rm cube}\int_{\mathcal B}\|f_i\|^2\pbocG c^{k-1}\le(2P^2)^{2n^2}\int_{U\backslash G}\|\mathcal W_i\|^2\le
(2P^2)^{2n^2}P^{4n^2}\frac dD\int|\varphi|^2 .
\]
Altogether the left side of the theorem is at most $(P^{275})^{2n^2}P^{4n^2}\frac dD\int|\varphi|^2\le P^{2000n^2}
\frac dD\int|\varphi|^2$.
\end{proof}

\subsubsection{The bound at one place}\label{pbo:ss-onplace}

\begin{lemma}[Tubes]\label{pbo:tube}
For $T\ge8r$ there are a compact $A\subseteq G$ and a compact $\Omega_0\subseteq\pboU$ of positive measure
with: $\|y\|\le T^{10}$ on $A$; $vy\in A$ whenever $v\in\Omega_0$ and $\|y\|\le T$; and $\pbvol\{v\in\pboU:vy\in
A\}\le3\pbvol(\Omega_0)$ for every $y$.
\end{lemma}
\begin{proof}
\emph{A gauge on $\pboU$.} For $u=u_q$ let $A(u)=\sigma-\frac12\beta\beta^*$ be the $\tilde e$-coefficient matrix in
\eqref{pbo:eq-uq} and $N(u)=\max(\|\beta\|^2,\|A(u)\|)$. From \eqref{pbo:eq-uq} and $u_qx=x+\tilde e\beta$,
$u_qu_{q'}$ has parameters $\beta+\beta'$ and $A=A(u)+A(u')-\beta\beta'^*$; hence $\sqrt{N(uu')}\le\sqrt{N(u)}+
\sqrt{N(u')}$. Also $A(u^{-1})=A(u)^*$ (from $A+A^*=-\beta\beta^*$), so $N(u^{-1})=N(u)$. As $\beta$ and $A$ are
blocks of $u-1$, $N(u)\le\max(\|u-1\|^2,\|u-1\|)\le(\|u\|+1)^2$; and $\|u\|\le1+2\|\beta\|+\|A\|\le(1+\sqrt{N(u)})^2$.
Conjugation by the element $a_t$ of $G$ with $a_t\tilde e_a=t\tilde e_a$, $a_t\tilde e^*_a=t^{-1}\tilde e^*_a$, fixing
$x_i,\tilde z$, maps $u_{(\beta,\sigma)}$ to $u_{(t\beta,t^2\sigma)}$; it multiplies $N$ by $t^2$ and $du^0$ by $t^Q$,
$Q=2\varrho k+\varrho^2$ (real dimension of the parameters). So the balls $B(\rho)=\{N\le\rho\}$ satisfy
$\pbvol B(\rho')=(\rho'/\rho)^{Q/2}\pbvol B(\rho)$, and they are compact with non-empty interior.

\emph{The tube.} Let $\rho=Q^2(T^2+1)^2$, $\Omega_0=B(\rho)$, $A=\Omega_0\cdot\{y:\|y\|\le T\}$, compact. On $A$,
$\|vy\|\le(1+\sqrt\rho)^2T\le(1+Q(T^2+1))^2T\le T^9\le T^{10}$, using $Q\le r^2/2\le T^2/128$. Fix $y$ and $v_0$ with
$v_0y=v_1y_1\in A$ ($v_1\in\Omega_0$, $\|y_1\|\le T$); if also $vy=v_2y_2\in A$, then $w:=v_2^{-1}vv_0^{-1}v_1=y_2y_1^{-1}$
lies in $\pboU$ and has $\|w\|\le T^2$, so $\sqrt{N(w)}\le T^2+1$, and $\sqrt{N(vv_0^{-1}v_1)}=\sqrt{N(v_2w)}\le\sqrt\rho+
T^2+1$. Thus $\{v:vy\in A\}\subseteq B(\rho')g^{-1}$ with $g=v_0^{-1}v_1$, $\sqrt{\rho'}=\sqrt\rho+T^2+1$. By right
invariance of $du^0$, its measure is at most $(\rho'/\rho)^{Q/2}\pbvol(\Omega_0)=(1+1/Q)^Q\pbvol(\Omega_0)\le
e\,\pbvol(\Omega_0)\le3\pbvol(\Omega_0)$.
\end{proof}

\begin{proof}[Proof of Proposition~\ref{pbo:local}]
Let $Z_0=\tilde Z(g)$. As $|Z_{0,ab}|=|h(y^{-1}\tilde f_a,y^{-1}g\tilde f_b)|\le\rho^2$ and $n\le\rho$, put
$M=\rho^{9n}/\pbg(Z_0)$; by Lemma~\ref{pbo:sizes}(a$'$) with $M=\rho$, $\pbg(Z_0)\le\rho^{8n}$, so $M\ge\rho$, and $M^{-1}\le
\pbg(Z_0)$, $|Z_{0,ab}|\le M^2$. Let $M_0=M^{9n}$, $P=M_0^{5n}$; then $P\ge M\ge\rho\ge b\ge8r$ and $P\ge
b^5$. By Lemma~\ref{pbo:witt} write $g=s\gamma_{Z_0}s'^{-1}$ with $s,s'\in\pboR$, and for $x=(h,v)\in H\times
\pboU$ put $t(x)=s'h^{-1}v^{-1}$, so that $gt(x)=s\gamma_{Z_0}(vh)^{-1}$. By Lemma~\ref{pbo:sizes}(a),(b)
(with $T=P/2\ge M^{8n}$), every $Z'$ with $\|Z'-Z_0\|<P^{-6}$ is $P$-bounded with $\pbg(Z')\ge c:=\pbg(Z_0)/2$.
Let $A,\Omega_0$ be a tube for $T=M_0^5$ (Lemma~\ref{pbo:tube}); then $\|y\|\le P^{20}$ on $A$ and $\Omega_0
\subseteq A$. Let $\mu_c=\pbocG c^{k-1}dh\,dv$, $m=\mu_c(\cdot\,|\,H\times\Omega_0)$, a probability measure, and
\[
\mu=\frac{P^{2000n^2}d}{D\,\mu_c(H\times\Omega_0)}\,\mathbf 1_A\|W\|^2\,dg .
\]
(a) Fix $\varphi$ bounded measurable. By \eqref{pbo:eq-inv} and the Bessel property, for fixed $x$ the
Hilbert--Schmidt norm squared of $\int_A\varphi(z)W(gt(x)z)W(t(x)z)^*dz$ equals $\sum_j\|\int_G\beta_j(z)
W(\gamma_{Z_0}(vh)^{-1}z)dz\|^2$ with $\beta_j=\mathbf 1_A\varphi\,\aubar{W_j}$: expand in the coordinates $j$
of the second factor, by \textup{(BS1)}, $W(t(x)z)=\pi_R(t(x))W(z)$ and $W(gt(x)z)=\pi_R(s)W(\gamma_{Z_0}(vh)^{-1}z)$ with
$\pi_R(t(x))$, $\pi_R(s)$ unitary (where $\pi_R(uh)=\psi(u)h$), and the Hilbert--Schmidt norm is invariant
under unitary factors on either side. As $\beta_j$ vanishes off $A$, Theorem~\ref{pbo:schur-section} applies (with
$\mathcal B\supseteq H\times\Omega_0$, since $\|v\|\le T^{10}\le P^{20}$ on $\Omega_0$):
\[
\int_T\cdots dm\le\frac1{\mu_c(H\times\Omega_0)}P^{2000n^2}\frac dD\sum_j\int|\beta_j|^2=\int|\varphi|^2d\mu .
\]
(b) Let $x=(h,v)$ with $v\in\Omega_0$, and $z$ with $\|z^{-1}\tilde f_a\|\le\rho\le M$ and $\pbel(gt(x)z)\le\rho
\le M$. By Lemma~\ref{pbo:norm-z} (with $R=M_0^{\,}\ge M^{9n}$; $\|Z_0^{-1}\|\le M^{8n}\le M_0$) $\|v^{-1}z\|\le
M_0^5=T$, so $z=v(v^{-1}z)\in A$.

\emph{The mass of $\mu$.} By Lemma~\ref{pbo:unfold} with $\varphi=\mathbf 1_A$, $w=\|W\|^2$ (left invariant), $K=3
\pbvol(\Omega_0)$ (tube) and $R=P^{40}$ (if $y,sy\in A$ then $\|s\|\le P^{40}$, so $|m(s)_{ab}|\le P^{40}$), and
(BS2), $\int_A\|W\|^2\le3\pbvol(\Omega_0)(2P^{40})^{2n^2}d$. Since $\mu_c(H\times\Omega_0)=\pbocG c^{k-1}\pbvol(
\Omega_0)$, the volume of $\Omega_0$ cancels:
\[
\mu(G)\le\frac{3(2P^{40})^{2n^2}P^{2000n^2}d^2}{D\,\pbocG\,(\pbg(Z_0)/2)^{k-1}}\le M^{10^5n^4}\frac{d^2}{D\pbocG(
\pbg(Z_0)/2)^{k-1}}=\pbsp(Z_0;D,\rho,d),
\]
because $3(2P^{40})^{2n^2}P^{2000n^2}\le M_0^5(M P^{40})^{2n^2}P^{2000n^2}=M^{93600n^4+2n^2+45n}\le M^{10^5n^4}$
(with $P=M^{45n^2}$, $M_0^5=M^{45n}$, $3\le M_0^5$, $2\le M$). Finally $A$ is compact, $t$ is measurable
and $m$ is a probability measure.
\end{proof}

\subsubsection{Proof of Theorem~\ref{pbo:O1}}\label{pbo:ss-proof}

By Theorem~\ref{pbo:system}, $(W_\tau,W^\flat_\tau)$ is a Bessel pair of degree $\pbdeg_\tau$ at every $\tau$,
and $\rho_\tau\ge\pbsth\ge b$. The kernel $\pbsK^\flat_\gamma(y)$ is the adjoint of $\pbsK^{\flat\prime}(y):=
W^\flat(\gamma^{-1}y)W^\flat(y)^*$, and the adjoint is an isometry of the Hilbert--Schmidt norm commuting
with integration; and $\pbg(\gamma_\tau^{-1})=\pbg(\gamma_\tau)$ since $\tilde Z(g^{-1})=\tilde Z(g)^*$. So it
suffices to prove the bound for $\int\bar\phi\,\pbsK^{\flat\prime}$ with $g_\tau=\gamma^{-1}_\tau$ and the point
$y_1=\gamma^{-1}y_0$ in place of $y_0$ (the hypotheses are symmetric).

Apply Proposition~\ref{pbo:local} at each $\tau$ to $g_\tau$, $\rho_\tau$, obtaining $A_\tau,(T_\tau,m_\tau),t_\tau,
\mu_\tau$. For $\mathbf x=(x_\tau)\in\prod T_\tau$ let $t(\mathbf x)=(t_\tau(x_\tau))$ and
\[
\mathcal N(\mathbf x)=\Bigl\|\int_{\prod A_\tau}\bar\phi(z)\,W(g\,t(\mathbf x)z)\,W(t(\mathbf x)z)^*\,dz\Bigr\|^2 .
\]
\emph{Products of contractions.} If kernels $k_\tau(x_\tau,z_\tau)$ (matrix valued) satisfy
$\int\|\int\varphi k_\tau\|^2dm_\tau\le\int|\varphi|^2d\mu_\tau$ for all bounded $\varphi$, then the tensor product
kernel satisfies $\int\|\int\varphi\bigotimes k_\tau\|^2\prod dm_\tau\le\int|\varphi|^2\prod d\mu_\tau$ for bounded
$\varphi$ on the product: by induction on the number of factors, integrate first over $z_{\tau_1}$ against the
kernel of $\tau_1$ with the other variables frozen as a vector-valued function ($\varphi$ replaced by the
partial integral, which is a bounded function of $z_{\tau_1}$ with values in a Hilbert space; the
contraction property for scalar $\varphi$ extends to Hilbert-space-valued $\varphi$ coordinatewise), then apply
the inductive hypothesis. Hence $\int\mathcal N\,d\prod m_\tau\le\int|\phi|^2d\prod\mu_\tau\le\prod_\tau\mu_\tau(
G_\tau)$.

\emph{One good shift.} The set of $\mathbf x$ for which (b) fails at some $\tau$ is null; as $\prod m_\tau$ is a
probability measure, there is $\mathbf x$ outside it with $\mathcal N(\mathbf x)\le\prod_\tau\mu_\tau(G_\tau)$.
Let $t=t(\mathbf x)\in\prod\pboR_\tau$.

\emph{Support.} If $\phi(tz)W^\flat(gtz)W^\flat(tz)^*\ne0$, then (hypothesis on $\phi$, via $y=tz$ after the
symmetrisation) $\|(tz_\tau)^{-1}\tilde f_a\|=\|z_\tau^{-1}\tilde f_a\|\le\rho_\tau$, and $W^\flat_\tau(g_\tau t_\tau
z_\tau)\ne0$ so $\pbel_\tau(g_\tau t_\tau z_\tau)\le b^{6n+2}\le\rho_\tau$ by (BS3); so $z\in\prod A_\tau$ by (b).
Hence $z\mapsto\bar\phi(tz)\pbsK^{\flat\prime}(tz)$ vanishes off the compact set $\prod A_\tau$, it is continuous
there up to the measurable factor $\phi$, and $\int_{G_\infty}\bar\phi\,\pbsK^{\flat\prime}=\int_{\prod A_\tau}\bar\phi(tz)
\pbsK^{\flat\prime}(tz)\,dz$ by left invariance of Haar measure; also $\phi(tz)=\phi(z)$.

\emph{Comparison with the uncut kernel.} Writing $\pbsK^{\flat\prime}=\pbsK'-(\pbsK'-\pbsK^{\flat\prime})$ on
$\prod A_\tau$,
\[
\Bigl\|\int\bar\phi\,\pbsK^{\flat\prime}\Bigr\|\le\mathcal N(\mathbf x)^{1/2}+\int_{G_\infty}|\phi|\,\|\pbsK'-
\pbsK^{\flat\prime}\|\le\Bigl(\prod_\tau\pbsp(\tilde Z(g_\tau);\pbdeg_\tau,\rho_\tau,d_\tau)\Bigr)^{1/2}+\int|\phi|\,
\|\pbsK_\gamma-\pbsK^\flat_\gamma\| ,
\]
using $\pbsp(\tilde Z(g^{-1}))=\pbsp(\tilde Z(g))$ ($\pbg$ is unchanged) and that the adjoint preserves the norm of
the difference. That $\phi\pbsK^\flat_\gamma$ vanishes outside a compact set follows from the support statement (translated back by $t$ and by the symmetrisation); it is bounded since $W$, $W^\flat$ are bounded, and measurable as a product of measurable functions; $\pbsK^\flat_\gamma$ is continuous by \textup{(BS1)}. \qed

\subsubsection{The Bessel functions form Bessel pairs}\label{pbo:ss-pairs}

Fix $\tau$; $c=c_\tau$, $\pbdeg=\pbdeg_\tau$, $\pi=\pi_\tau$, $v_0$ as in \S\ref{pbe:ss-ds}, $da=dg$ the Haar measure
of \textup{(DS1)}. Let $\mathfrak g_0$ be the Lie algebra of $G$ (skew-hermitian for $\mathsf J_1$), with the
operator norm. For $X\in\mathfrak g_0$ write $dX$ for the right derivative $(dXf)(g)=\frac d{dt}f(ge^{tX})|_{t=0}$.
Let $L_\tau=\max_i|2\varsigma_{\tau,i}|$.

\begin{lemma}[Derivatives of $K$-finite vectors]\label{pbo:K}
Let $\mathcal H^s$ be the span of the vectors $d\pi(X_1)\cdots d\pi(X_j)v_0$, $j\le s$, $X_i\in\mathfrak g_0$
(these are $K$-finite smooth vectors). For $x\in\mathcal H^s$ and $X\in\mathfrak g_0$,
\[
\|d\pi(X)x\|\le(r+1)^4\,(L_\tau+r+1+s)\,\|X\|\,\|x\|,\qquad d\pi(X)x\in\mathcal H^{s+1}.
\]
\end{lemma}
\begin{proof}
$G=\mathrm U(k+\varrho,\varrho+1)$ is connected, $K=\mathrm U(k+\varrho)\times\mathrm U(\varrho+1)$ is connected and
contains the diagonal torus $T_c$, a compact Cartan subgroup; $\pi$ is a discrete series representation with
Harish-Chandra parameter $\varsigma$ (regular, entries in $\frac12\Z$; the chamber is fixed in \S\ref{pbe:ss-ds}).
Use the trace form $B(X,Y)=\operatorname{tr}(XY)$ on $\mathfrak g$, the Cartan decomposition $\mathfrak g_0=
\mathfrak k_0\oplus\mathfrak p_0$, and coordinates on $\mathfrak t^*$ in which the roots are $\epsilon_i-\epsilon_j$.
We use two published results.

(i) \emph{Lowest $K$-type} (Schmid \cite[Lemma~(2.22), Lemma~(3.8), Corollary~(3.9), Corollary~(4.19) and the proof of Theorem~(1.3) on p.~560]{Schmid75}; see also \cite[Theorem~(1.3), p.~130]{HechtSchmid75}; for groups outside their standing hypotheses these apply to $\auU(p,q)=Z\cdot\mathrm{SU}(p,q)$, $Z$ the compact centre acting by a character, through $\mathrm{SU}(p,q)$, see the paragraph ``On (D3)'' in {\cite[\S3.1.3]{P49}}): for $G$ connected semisimple, or reductive with compact centre component as here, with a
compact Cartan subgroup, the discrete series with Harish-Chandra parameter $\varsigma$ contains the $K$-type of
highest weight $\mu_0=\varsigma+\rho_n-\rho_c$ with multiplicity one, $\rho_n,\rho_c$ the half sums of the
noncompact and compact roots positive for $\varsigma$. Its hypotheses hold by the above. As $v_0$ is a highest
weight vector of this $K$-type and the entries of $\rho_n$, $\rho_c$ have absolute value $\le r/2$, $|\mu_{0,i}|\le
L_\tau/2+r$.

(ii) \emph{Casimir} (Harish-Chandra; see \cite[Theorem~8.18, (8.32) and the Example on p.~225]{Knapp}): on the $(\mathfrak g,K)$-module of $\pi$ the Casimir element $\Omega_G=\sum
X_iX^i$ (dual bases for $B$) acts by $|\varsigma|^2-|\rho|^2$; on an irreducible representation of the
connected group $K$ of highest weight $\mu$, $\Omega_K$ acts by $|\mu+\rho_K|^2-|\rho_K|^2$, with $\rho_K=\rho_c$ for
the compact positive roots. (Norms of the standard inner product in the coordinates $\epsilon_i$.)

\emph{$K$-types of $\mathcal H^s$.} The adjoint representation of $K$ on $\mathfrak p_\C$ has weights $\pm(\epsilon_i-
\epsilon_j)$ (one index in each block), with entries in $\{0,\pm1\}$. If $x$ lies in a sum of $K$-types of highest
weights $\mu$, then $d\pi(X)x$ ($X\in\mathfrak p$) lies in the image of $\mathfrak p_\C\otimes V_\mu$, whose
irreducible constituents have highest weights $\mu+\beta$ with $\beta$ a weight of $\mathfrak p_\C$ (a constituent's
highest weight is a weight of the tensor product, hence of the form (weight of $V_\mu$)$+\beta$, and dominance
bounds it entrywise by $\max|\mu_i|+1$); $d\pi(X)$, $X\in\mathfrak k$, preserves isotypic components. So every
$K$-type $\mu$ of $\mathcal H^s$ has $\max_i|\mu_i|\le L_\tau/2+r+s=:\Lambda_s$.

\emph{$\mathfrak k$-part.} For $X\in\mathfrak k_0$, conjugate by $K$ to $X=\operatorname{Ad}(k)H$, $H\in\mathfrak t_0$ with
$\|H\|=\|X\|$; on $V_\mu$, $d\pi(H)$ is diagonal in a weight basis with eigenvalues $\langle\nu,H\rangle$, $|\langle\nu,
H\rangle|\le\sum_i|\nu_i||H_i|\le r\Lambda_s\|X\|$ (the weights $\nu$ of $V_\mu$ satisfy $|\nu_i|\le\max|\mu_j|$). Hence $\|d\pi(X)
x\|\le r\Lambda_s\|X\|\|x\|$.

\emph{$\mathfrak p$-part.} Let $(Y_j)$ be a basis of $\mathfrak p_0$ orthonormal for $B$ (positive on $\mathfrak p_0$) and
$(Z_i)$ a basis of $\mathfrak k_0$ with $B(Z_i,Z_{i'})=-\delta_{ii'}$; then $\Omega_G=\sum Y_j^2-\sum Z_i^2$ and $\Omega_K=-\sum Z_i^2$, so
$-\sum_jd\pi(Y_j)^2=\Omega_K-\Omega_G$ on $K$-finite vectors. For $X=\sum a_jY_j\in\mathfrak p_0$, $(\sum a_j^2)^{1/2}=\|X\|_{HS}\le
\sqrt r\|X\|$, and by skew-adjointness of $d\pi(Y_j)$ on smooth vectors,
\[
\|d\pi(X)x\|^2\le\|X\|_{HS}^2\sum_j\|d\pi(Y_j)x\|^2=\|X\|_{HS}^2\langle(\Omega_K-\Omega_G)x,x\rangle\le r\|X\|^2(|\mu+\rho_K|^2_{\max}+|\rho|^2)\|x\|^2,
\]
where $|\mu+\rho_K|^2\le r(\Lambda_s+r)^2$ and $|\rho|^2\le r\cdot r^2$; so $\|d\pi(X)x\|\le\sqrt2\,r(\Lambda_s+r)\|X\|\|x\|$.

\emph{Total.} For $X\in\mathfrak g_0$, the components $X_{\mathfrak k}=\frac12(X+\theta X)$, $X_{\mathfrak p}=\frac12(X-\theta X)$ have norm $\le\|X\|$, so
$\|d\pi(X)x\|\le(r+\sqrt2r)(\Lambda_s+r)\|X\|\|x\|\le3r\cdot(L_\tau+r+1+s)\cdot2\cdot\|X\|\|x\|$, and $6r\le(r+1)^4$. The same
holds for the complex matrix units $E_{ab}=X+iX'$ with $\|X\|,\|X'\|\le1$, at the cost of a factor $2$, still within
$(r+1)^4$.
\end{proof}

\begin{lemma}[Schur orthogonality for translates]\label{pbo:translates}
Let $\psi=\pbdeg\cdot c$. Then $\psi(g^{-1})=\aubar{\psi(g)}$, $\psi(1)=\pbdeg$, and for all $t,g\in G$
\begin{equation}\label{pbo:eq-schur}
\int_G\psi(s)\,\psi(t^{-1}s^{-1}g)\,ds=\frac{\psi(t^{-1})}{\pbdeg}\,\psi(g);
\end{equation}
in particular ($t=1$) $\psi*\psi=\psi$. For a Bessel function $F$ with $\int_{U\backslash G}\|F\|<\infty$ and $F*\psi=F$, where
$(F*k)(y)=\int k(a^{-1})F(ya)\,da$, the translates satisfy
\[
\int_{U\backslash G}\langle F(ya),F(yb)\rangle\,dy=\frac{\psi(a^{-1}b)}{\pbdeg}\int_{U\backslash G}\|F\|^2 .
\]
\end{lemma}
\begin{proof}
The first two identities are \textup{(DS1)(i)}. \eqref{pbo:eq-schur} is Schur's orthogonality relation
$\int_G\langle u,\pi(s)v\rangle\,\aubar{\langle u',\pi(s)v'\rangle}\,ds=\pbdeg^{-1}\langle u,u'\rangle\aubar{\langle v,v'\rangle}$
(conjugate-linear first variable; $\pi$ square integrable and irreducible, $\int|c|^2=\pbdeg^{-1}$), applied to
$u=v_0$, $v=v_0$, $u'=\pi(t)v_0$, $v'=\pi(g)v_0$ after writing $\psi(t^{-1}s^{-1}g)=\pbdeg\,\aubar{\langle\pi(g)v_0,
\pi(s)\pi(t)v_0\rangle}$; it contains \textup{(DS1)(ii)}.
Let $k(g)=\int_{U\backslash G}\langle F(y),F(yg)\rangle dy$; by the transformation law the integrand is left
$U_\tau$-invariant, and the quotient measure is right invariant, so the left side is $k(a^{-1}b)$; $k$ is
bounded, $k(g^{-1})=\aubar{k(g)}$, and $F*\psi=F$ gives $k(g)=\int\psi(t^{-1})k(gt)\,dt$ and $k(g)=\int\psi(t)
k(t^{-1}g)\,dt$. Combining, $k(g)=\int\int\psi(s)\psi(t^{-1})k(s^{-1}gt)\,ds\,dt$, and (substitute $s^{-1}gt=u$ and use \eqref{pbo:eq-schur} in the form $\int\psi(s)\psi(u^{-1}s^{-1}g)\,ds=\psi(u^{-1})\psi(g)/
\pbdeg$) yields $k(g)=\frac{\psi(g)}{\pbdeg}\int\psi(u^{-1})k(u)\,du=\frac{\psi(g)}{\pbdeg}k(1)$.
\end{proof}

\begin{proof}[Proof of Theorem~\ref{pbo:system}]
(BS1) is Proposition~\ref{pbe:local}(a) and Lemma~\ref{pbe:len}(a) (for $W^\flat=\chi W$, $\chi$ continuous and
$R_\tau$-invariant). (BS2) is the normalisation of $W_\tau$. (BS3) is Proposition~\ref{pbe:local}(e).
(BS5) is a size condition (\S\ref{pbo:ss-sizes}).

(BS4). Let $\psi=\pbdeg\,c$. By Theorem~\ref{pbe:bessel}(B7), $\int e(a)W(ya)\,da=W(y)$ with $e=\pbdeg\,\bar c$,
i.e.\ $W*\psi=W$, since $\psi(a^{-1})=\pbdeg\,\aubar{c(a)}=e(a)$. Choose $\beta\in C^\infty_c(G)$ with $\pi(\beta)
v_0=v_0$, where $\pi(\beta)=\int\beta(g)\pi(g)\,dg$ (it exists because $v_0$ is $K$-finite and
$Z(\mathfrak g)$-finite: by Harish-Chandra's theorem \cite[Theorem~1, p.~18]{HC66}, used in this form in
\cite[Ch.~I, \S4, proof of Lemma~14, p.~15]{HarishChandraAF},
a smooth function $f$ on $G$ that is left $K$-finite and $Z(\mathfrak g)$-finite satisfies $f=\alpha*f$ for some
$\alpha\in C^\infty_c(G)$, $(\alpha*f)(x)=\int\alpha(y)f(y^{-1}x)\,dy$ (the published statement is $f^\vee*\alpha'=f^\vee$
for the right $K$-finite function $f^\vee(x)=f(x^{-1})$; put $\alpha(y)=\alpha'(y^{-1})$, as in (D5) of
{\cite[\S4.2.3]{P53}}); apply this to $f(x)=\langle u,\pi(x^{-1})v_0\rangle$ for a fixed $u\ne0$, which is
linear in $v_0$ (the inner product is conjugate-linear in the first variable), smooth, left $K$-finite
($f(kx)=\langle u,\pi(x^{-1})\pi(k^{-1})v_0\rangle$) and $Z(\mathfrak g)$-finite (the left regular action of
$z\in Z(\mathfrak g)$ on $f$ is the action of an element of $Z(\mathfrak g)$ on $v_0$, $Z(\mathfrak g)$ being stable
under the principal antiautomorphism), because $v_0$ is smooth, $K$-finite and $Z(\mathfrak g)$-finite. Then
$(\alpha*f)(x)=\int\alpha(y)\langle u,\pi(x^{-1})\pi(y)v_0\rangle dy=\langle u,\pi(x^{-1})\pi(\alpha)v_0\rangle$, so
$\pi(\alpha)v_0-v_0$ is orthogonal to $\pi(G)u$, which spans a dense subspace as $\pi$ is unitary and
irreducible; take $\beta=\alpha$). For a word $w=(\ell_1,\dots,\ell_m)$ of matrix units $E_{ab}$ let $x_w=d\pi(E_{\ell_1})\cdots
d\pi(E_{\ell_m})v_0$ (complexified action) and $X_w(g)=\pbdeg\,\langle v_0,\pi(g)x_w\rangle$; $X_\varnothing=\psi$.
Then $X_w=\psi*X_w$ (as $x_w$ lies in the irreducible $\pi$: Schur orthogonality), $X_w$ is integrable (it is
$\psi*(\text{derivative of }\check\beta)$, a convolution of $L^1$ functions), and by Schur orthogonality and
Lemma~\ref{pbo:K},
\[
\frac1{\pbdeg}\int_G|X_w|^2=\|x_w\|^2\le\bigl((r+1)^4(L_\tau+r+1+|w|)\bigr)^{2|w|}\le P^{2|w|}
\]
for $|w|\le2n^2$ and $P\ge b^5$, by the size condition (Z4). Put $V_w=W*X_w$. These are Bessel functions
(left translation commutes with right convolution; bounded as $W$ is bounded and $X_w\in L^1$). For a curve
with $c'=cX$, $X=\sum X_{ab}E_{ab}$, $V_w(c(t))=\int W(a)X_w(a^{-1}c(t))\,da$ and the right derivative of
$X_w$ along $E_{ab}$ is $X_{(a,b)w}$ (as $\frac d{dt}\pi(ge^{tE})x=\pi(g)d\pi(E)x$); differentiation under the
integral is justified as follows. From $v_0=\pi(\beta)v_0$ we get $x_w=\pi(\beta_w)v_0$ with $\beta_w\in
C^\infty_c(G)$ (left derivatives of $\beta$), so $X_w(h)=\int\beta_w(a)\psi(ha)\,da$. For $c(t)$ in a compact
set $C$, $\sup_t|X_{(a,b)w}(a'^{-1}c(t))|\le\|\beta_{(a,b)w}\|_1\sup\{|\psi(a'^{-1}cb')|:c\in C,\ b'\in
\operatorname{supp}\beta_{(a,b)w}\}\le C'\bigl(2/\|a'\|\bigr)^{\pbeq}$ by \textup{(DS2)} with $j=0$ and
$\|a'^{-1}cb'\|\ge\|a'\|/(\|c\|\,\|b'\|)$, for a constant $C'$ depending on $C$ and $w$; this bound is
integrable in $a'$ by \textup{(DS3)}, and $W$ is bounded, so dominated convergence applies. This is
(BS4)(i).

For (BS4)(ii): by Lemma~\ref{pbo:translates} the vectors $\omega_a=F(\cdot\,a)$ of $L^2(U\backslash G)$, $F=W$, have
$\langle\omega_a,\omega_b\rangle=\frac{\psi(a^{-1}b)}{\pbdeg}d$. Let $\mathcal Z=X_w*\check\varphi$; then $g\mapsto\int\varphi(y)
V_w(gy)dy=W*\mathcal Z$, and in $L^2(U\backslash G)$ this function $\Phi$ satisfies $\|\Phi\|^2=\langle\Phi,\int
\mathcal Z(a^{-1})\omega_a\,da\rangle\le\|\Phi\|\,\|\int\mathcal Z(a^{-1})\omega_a\,da\|$. The vector $V=\int
\mathcal Z(a^{-1})\omega_a\,da=\int\varphi(g)\,\Big(\int X_w(a^{-1})\omega_{ga}\,da\Big)dg$ lies in the closed span of the
$\omega_a$; for $u=\sum\lambda_i\omega_{a_i}$, $\langle u,\int X_w(a^{-1})\omega_{ga}da\rangle=\frac d{\pbdeg}\sum\bar\lambda_iX_w(a_i^{-1}g)$
(by $\psi*X_w=X_w$), whose $L^2(dg)$-norm squared is $\frac{d^2}{\pbdeg^2}\cdot\frac{\pbdeg\|u\|^2}{d}\cdot\frac{\int|X_w|^2}{\pbdeg}$
by Schur orthogonality ($\int X_w(a^{-1}g)\aubar{X_w(b^{-1}g)}dg=\frac{\psi(a^{-1}b)}{\pbdeg}\int|X_w|^2$). By
Cauchy--Schwarz in $g$, $|\langle u,V\rangle|^2\le\int|\varphi|^2\cdot\frac{d}{\pbdeg}\frac{\int|X_w|^2}{\pbdeg}\|u\|^2$, and by
density $\|V\|^2\le\frac{\int|X_w|^2}{\pbdeg}\cdot\frac d{\pbdeg}\int|\varphi|^2\le P^{2|w|}\frac d{\pbdeg}\int|\varphi|^2$.
This is (BS4)(ii) with $D=\pbdeg$.
\end{proof}

\subsubsection{Size conditions}\label{pbo:ss-sizes}

Besides those of \S\S\ref{pbe:sec}, \ref{pbs:ss-sizes}, the proof uses:
\begin{itemize}
\item[(Z4)] $8r\le b$ (BS5), and $(r+1)^4(L_\tau+r+1+2n^2)\le b^5$ for every $\tau$, where $L_\tau$ is the
largest absolute value of an entry of $2\varsigma_\tau$. The first is in \S\ref{pbe:sec} ($r\le b$ is listed there;
$8r\le b$ holds since $r\le3x^5$ and $b\ge x^{10^{12}}$ by {\cite[Lemma~13.4]{P46}}). The second holds as soon
as $L_\tau\le b^4$, by the same comparison; and $L_\tau\le b^2$ holds: the entries of $\varsigma_\tau$ are $k-1$ numbers of absolute value $<k\le3x^5$ and
the numbers $-\kappa(\tau)-a\Delta$ with $|\kappa(\tau)|\le b$, $1\le a\le m\le2\sqrt x$, $\Delta=2r^8\le2\cdot3^8x^{40}$
({\cite[\S14.1]{P46}}, {\cite[Lemma~13.4(S2)]{P46}},(S5)), so $L_\tau\le2(b+m\Delta)\le b^2$.
\end{itemize}


\appendix

\section{Statements used from \texorpdfstring{\cite{P46}}{[Row D]}, \texorpdfstring{\cite{P49}}{[Theta]}, \texorpdfstring{\cite{P52}}{[Frames]}}\label{B1:quoted}

For the convenience of the reader we reproduce here, without proofs, the statements of other papers that are used in this one, in the notation of the paper quoted. Each is headed by its place in that paper.

\begin{quoteddef}{\cite[Definition~13.3]{P46}}{ (admissible data)}
Let $(q_0,d)$ be parameters. \emph{Admissible data} $D$ for $(q_0,d)$ consist
of the following objects, subject to the conditions \textup{(D1)--(D8)}.
\begin{enumerate}
\item[(D1)] Integers $1\le x\le x_1\le b$, the \emph{scales}.
\item[(D2)] They satisfy
\begin{gather*}
 10^{12}\,\anCP\le\log\log b,\qquad 100\log\log b\le\log x,\qquad 10^{6}\log x\le\log x_1,\\
 10^{6}\log x_1\le\log b,\qquad 6\log\log x_1\le\log\log b .
\end{gather*}
\item[(D3)] Integers $m$, $r$, $s$: $m$ is a power of $2$ with
$\sqrt x\le m\le2\sqrt x$; $r$ is a power of $3$ with $x^{5}\le r\le3x^{5}$;
$2^{s}=m^{22}$; and $s\le32\log x$. The field of the data is $E=E_s$, with
$q_1=q_1(s)$, $q_2=q_2(s)$, $f=[E:K]$, $\anPl=\anPl_s$.
\item[(D4)] $q_1x_1$ divides $b$.
\item[(D5)] A set $\anPb\subseteq\anPl$ (the \emph{$b$-places}) with
$\ancl_{q_0}(\sigma)=-1$ for all $\sigma\in\anPb$, with
$T(s)\le\#\anPb\cdot\log b$, and sparse for the primes up to $5r(r+1-m)$.
\item[(D6)] The weight $\kappa\colon\anPl\to\mathbb N$: $\kappa(\tau)=b$ for
$\tau\in\anPb$ and $\kappa(\tau)=xb/x_1$ otherwise.
\item[(D7)] Two distinct primes $p_0,p_1\equiv1\pmod{q_2^{2}}$ in the interval
$[b^{10},b^{11}]$ and an integer $\pi_M$ of $K$ with
$\pi_M\overline{\pi_M}=p_0p_1=:M$.
\item[(D8)] A prime $N\equiv1\pmod{q_2^{2}}$ with
\[
 \log N_-\le\log N\le\log N_-+\sqrt{\log b},\qquad
 \log N_-:=\frac mf\sum_{\tau\in\anPl}\log\kappa(\tau)+(m-1)\log\frac{M\,\anrd_E}{2\pi}-\frac{15}2\log b ,
\]
where $\anrd_E=|\operatorname{disc}E|^{1/[E:\mathbb Q]}$ is the root
discriminant, and an integer $\pi_N$ of $E$ with $\pi_N\pi_N^{*}=N$.
\end{enumerate}
We write $a=xb$ and $k=r+1-m$; the sets $\anPa$ and $\anPb$ are disjoint.
\end{quoteddef}

\begin{quotedthm}{\cite[Lemma~13.4]{P46}}{}
For admissible data the following hold.
\begin{enumerate}
\item[(S1)] $x^{10^{6}}\le x_1$, $x_1^{10^{6}}\le b$,
$4\log x_1\le(\log x_1)^{6}\le\log b$,
$(\log b)^{100}\le x$, $\log b\ge128$, and
$\log x\ge10^{14}(q_0+2d+|\aneu|)$.
\item[(S2)] $x^{11}\le2^{s}\le2^{22}x^{11}$; $l_0(s)\le2592\log x$;
$2m\le r$, so $r/2\le k\le r$.
\item[(S3)] $x^{11}\le T(s)=\#\anPa\le f\le q_1\le x^{12}$, and
$f=\varphi(q_0)\,\#\anPa$; moreover $1\le\#\anPa/\#\anPb\le\log b$.
\item[(S4)] $\anrd_E=\sqrt d\;q_1\prod_{p\mid q_0}p^{-1/(p-1)}$, and
$x^{10}\le\anrd_E\le x^{12}$.
\item[(S5)] $1\le xb/x_1\le\kappa(\tau)\le b\le a\le b^{2}$ for all $\tau$;
$b^{20}\le M\le b^{22}$;
$20\log b\le\log(M\anrd_E/2\pi)\le23\log b$; and
$m\log b\le\log N\le25\,m\log b$.
\item[(S6)] With $\Omega:=e^{\aneu}\cdot2\pi/\anrd_E$: $0<\Omega\le1$ and
$20\log b\le\log(M/\Omega)\le23\log b$.
\item[(S7)] With $n=m/2$ and $\varrho=n-1$ (so that $r=k+2\varrho+1$): $r$ is odd; every prime
dividing $q_1$ divides $q_0$; $q_0^{\,q_0}\le x$; $\varrho\ge10^{5}$;
$k\ge512000\,n^{9}$ and $k\ge1000\,(2\varrho+2)^{9}$; $k\ge2\varrho^{2}+3$; $18\varrho f\le b^{8}$;
$f\le b/1000$; $\anrd_E\le b$; and $N\le b^{60n}$. Moreover, with $\Delta=2r^{8}$:
\begin{itemize}
\item[(a)] $\varrho\ge10^{8}$ (so $n\ge100$); $n^{7}\le k$; $12n^{3}+500\le k$; $k-3n-6>10^{9}n^{3}$;
\item[(b)] $r\ge100$ (so $2r^{8}\ge500r^{3}+10r^{5}$); $3r^{3}f\le r^{7}$; $8nr\le b$; $kf\ge40$; $f\ge4000\,n$;
\item[(c)] $2f\le2^{2\varrho}$ and $2^{-\varrho}\le b^{-6}$;
\item[(d)] $(200\,f\,k\,\anrd_E)^{1000}\le b$ and $(25\,n^{3}f\,\anrd_E)^{1000}\le b$;
\item[(e)] $\sqrt d\le\anrd_E$, and $n,r,f,\sqrt d\le b$; $M\le b^{23}$;
\item[(f)] $\binom{a+k-1}{k-1}^{f}\le b^{2kf}$ (this bounds the dimension $d_\infty=\dim W_{\lambda^X}$ of
{\cite[\S14.3]{P46}}, a product of at most $f$ such binomial coefficients);
\item[(g)] $m\Delta+r\le\kappa(\tau)$ for all $\tau$; $k+2b+2m\Delta\le b^{2}$ (this bounds every
entry of $2\varsigma_\tau$, so that $L_\tau:=\max_i|2\varsigma_{\tau,i}|\le b^{2}$); and
$(r+1)^{4}(b^{2}+r+1+2n^{2})\le b^{5}$.
\end{itemize}
\end{enumerate}
\end{quotedthm}

\begin{quotedthm}{\cite[Lemma~14.5]{P46}}{ (Free levels)}
Let $n\le2n_m$ and let $J\subseteq K^\circ$ consist of elements congruent to $1$ modulo every
chosen place. Then $\Gamma\cap gJg^{-1}=1$ for all $g\in G_f$.
\end{quotedthm}

\begin{quotedthm}{\cite[Theorem~14.14]{P46}}{ (the Spanning Theorem)}
$\audT>0$, and there is a finite set $\mathsf t\subseteq\auG$ with
\[
\#\mathsf t\cdot\dim\auSram\le5\,\audT\log(1+2b^2)\qquad\text{and}\qquad
\dim\sum_{s\in\mathsf t}s\aupp(\auSram)\ \ge\ \Bigl(1-\frac4{b^2}\Bigr)\audT .
\]
\end{quotedthm}

\begin{quotedthm}{\cite[Lemma~14.29]{P46}}{ (Count of the ramified forms)}
$\dim\auSram\ge\bigl(1-f/(N-1)\bigr)\auDX\ge(1-b^{-8})\auDX$.
\end{quotedthm}

\begin{quotedthm}{\cite[Lemma~14.30]{P46}}{ (One translate)}
Let a finite group $\auG$ act unitarily on $T$, let $T_0\subseteq T$ be a stable subspace with
projection $\mathrm{pr}_0$, let $B\colon S\to T$ be linear and $\varepsilon\ge0$. Suppose
\begin{enumerate}
\item[(H1)] $\|\mathrm{pr}_0B\,\mathrm{pr}_U\|^2\ge\dim U-\varepsilon\dim S$ for every subspace
$U\subseteq S$;
\item[(H2)] $\dim T_0\cdot\sum_{s\in\auG}\|B^*sB\|^2\le\#\auG\,(1+\varepsilon)(\dim S)^2$.
\end{enumerate}
Then for every subspace $F\subseteq T_0$, with $\alpha=\dim F/\dim T_0$, there is $s\in\auG$ with
$\dim(\mathrm{pr}_0sB)^{-1}(F)\le\bigl(\tfrac12(1+\alpha)+\tfrac32\varepsilon\bigr)\dim S$.
\end{quotedthm}

\begin{quotedthm}{\cite[Lemma~14.31]{P46}}{ (Few translates)}
Let $B_s\colon V\to W$, $s\in\Sigma$, be linear maps with $V\ne0$, $\dim W>0$, and
$0<\varepsilon<\tfrac1{16}$. Suppose the $B_s$ \emph{spread with error $\tfrac32\varepsilon$}: for
every subspace $F\subseteq W$, with $\alpha=\dim F/\dim W$, some $s$ has
$\dim B_s^{-1}(F)\le(\tfrac12(1+\alpha)+\tfrac32\varepsilon)\dim V$. Then there is a finite
$\mathsf t\subseteq\Sigma$ with
\[
\#\mathsf t\cdot\dim V\le\tfrac52\bigl(1+\log\tfrac1{16\varepsilon}\bigr)\dim W,\qquad
\dim\sum_{s\in\mathsf t}B_s(V)\ge(1-16\varepsilon)\dim W .
\]
Moreover, for a subspace $U\subseteq V$,
$\dim\sum_{s\in\mathsf t}B_s(V)\le\dim\sum_{s\in\mathsf t}B_s(U)+\#\mathsf t(\dim V-\dim U)$.
\end{quotedthm}

\begin{quotedthm}{\cite[Proposition~14.32]{P46}}{ (Second moment of the Poincar\'e map)}
Let $\deg$ be the product over $\tau\in\auPl$ of the formal degrees of the discrete series
$\varsigma_\tau$, and $\mathrm{mass}_r$ the mass of the genus of the standard hermitian lattice of
rank $r$. There are an inner product on $\auS$ and a scalar $c_P>0$ such that $P=c_P\auP$
satisfies:
\begin{enumerate}
\item[(i)] $0<\audT\le\#\auG\cdot\deg\cdot\mathrm{mass}_r$;
\item[(ii)] $\|\varrho(e_0)P\,\mathrm{pr}_U\|^2\ge\dim U-5b^{-4}\auDX$ for every subspace
$U\subseteq\auS$;
\item[(iii)] $\sum_{s\in\auG}\|P^*\varrho(s)P\|^2\le(1+9b^{-4})\,\auDX^2/(\deg\cdot
\mathrm{mass}_r)$.
\end{enumerate}
\end{quotedthm}

\begin{quotedthm}{\cite[Proposition~14.34]{P46}}{ (Spreading)}
$\audT>0$, and for every subspace $F\subseteq\auT$, with $\alpha=\dim F/\audT$, there is
$s\in\auG$ with $\dim\auP_s^{-1}(F)\le\bigl(\tfrac12(1+\alpha)+\tfrac32(8b^2)^{-1}\bigr)\auDX$.
\end{quotedthm}

\begin{quotedthm}{\cite[Proposition~14.36]{P46}}{ (Equality of multiplicities)}
For $1\le j\le m$ and every element $h$ of the spherical Hecke algebra away from $M$ (the
compactly supported functions on the finite adelic group outside the chosen places that are
bi-invariant under the stabiliser of the lattice, at all places, split or not),
$\operatorname{tr}(h\,|\,e_0\mathcal A(W_j))=\operatorname{tr}(h\,|\,\auT)$. In particular
$\dim e_0\mathcal A(W_j)=\audT$, and each system of eigenvalues of the Hecke operators at the split
places prime to $M$ has the same multiplicity in $e_0\mathcal A(W_j)$ and in $\auT$. (The
hypothesis of Proposition~\ref{pm:mult} holds for the data: at every $\tau$, two entries of
$\varsigma_\tau$ on coordinates of different signs \eqref{pm:eq-signs} differ by at least
$\Delta=2r^8\ge500r^3+10r^5$.)
\end{quotedthm}

\begin{quotedthm}{\cite[Proposition~2.19]{P49}}{}
Proposition~\ref{pds:haar} holds, with the same statement, for every parameter to which it is
applied in the paper. More precisely: for $(\lambda,s)$ satisfying the gap condition
\eqref{pbr:eq-gap} and $dg$ the Haar measure compatible with the Haar measure of mass $1$ on
$\auU(N)$, $\int_G|c|^2dg=\bigl(\prod_{i<j}(\lambda_i-\lambda_j)/(j-i)\bigr)^{-1}$ for the coefficient
$c$ of a unit vector of the lowest $K$-type; and the parameters of \S\ref{pds:sec} satisfy
\eqref{pbr:eq-gap}.
\end{quotedthm}

\begin{quotedthm}{\cite[Theorem~3.17]{P49}}{}
Assume $h\ge n^2+1$. Then $F=c_\lambda$ on $G$; that is, \eqref{pds:eq-FJ} holds.
\end{quotedthm}

\begin{quotedthm}{\cite[Lemma~3.18]{P49}}{}
Let $w\in\Z^n$ and let $x$ be invertible.
\begin{enumerate}
\item[(a)] $\pcom_{-w}(x)=\pcom_w(x^{-1})$, and $\overline{\pcom_w(x)}=\pcom_w(x^*)$.
\item[(b)] Let $b$ be invertible, with $\Delta_S(b)=\prod_{a\in S}b_{aa}$ for every $S$. If
$b_{ab}=0$ whenever $w_a<w_b$, then $\pcom_w(xb)=b^w\pcom_w(x)$. If $b_{ab}=0$ whenever
$w_b<w_a$, then $\pcom_w(bx)=b^w\pcom_w(x)$. The condition on the minors holds for diagonal $b$,
for $b=1+sE_{ab}$ \textup($a\ne b$, $s\in\C$\textup), and for $b\in\pcoB\cup\pcoB^-$.
\item[(c)] If $\varepsilon_a=\varepsilon_b$ and $\lambda_a>\lambda_b$, then $\Lambda_a\ge\Lambda_b$.
For all $a,b$: $\mu_a>\mu_b$ if and only if $\lambda_a>\lambda_b$.
\item[(d)] If $Z$ is accretive and $y$ invertible, then $y^*Zy$ and $Z^{-1}$ are accretive, and
every principal minor of $Z$ is nonzero; so $\pcom_w(Z)$ is defined and nonzero.
\end{enumerate}
\end{quotedthm}

\begin{quotedthm}{\cite[Corollary~3.19]{P49}}{}
Let $Y$ be invertible.
\begin{enumerate}
\item[(a)] For $b\in\pcoB$: $\pcom_\Lambda(Yb)=b^\Lambda\pcom_\Lambda(Y)$, $\pcom_\mu(Yb)=
b^\mu\pcom_\mu(Y)$, $\pcom_\Lambda(b^*Y)=\overline{b^\Lambda}\pcom_\Lambda(Y)$ and $\pcom_\mu(b^*Y)
=\overline{b^\mu}\pcom_\mu(Y)$.
\item[(b)] If $\varepsilon_a=\varepsilon_b$ and $\lambda_a>\lambda_b$, then $\pcom_\Lambda(Y(1+sE_{ab}))
=\pcom_\Lambda(Y)=\pcom_\Lambda((1+sE_{ba})Y)$ for $s\in\C$. For $t$ diagonal,
$\pcom_\Lambda(Yt)=t^\Lambda\pcom_\Lambda(Y)=\pcom_\Lambda(tY)$.
\item[(c)] Put $\psi(Y)=\pcom_{-\mu}(Y)^{-1}$ where $\pcom_{-\mu}(Y)\ne0$. If $\lambda_a<\lambda_b$
\textup(any signs\textup), $\psi(Y(1+sE_{ab}))=\psi(Y)$; for $t$ diagonal, $\psi(Yt)=t^\mu\psi(Y)$.
\end{enumerate}
\end{quotedthm}

\begin{quotedthm}{\cite[Proposition~3.21]{P49}}{}
For every $x\in\mathrm{GL}_n(\C)$, $a\mapsto I(x;a)$ is constant on $A_x$. Moreover the function
$\mathcal J(p,r)$ defined in the proof is locally constant on the open set $U_x$ defined there; in
particular $\mathcal J(p,r)=\mathcal J(1,1)$ whenever $(p,r)$ is joined to $(1,1)$ by a path in $U_x$.
\end{quotedthm}

\begin{quotedthm}{\cite[Lemma~3.23]{P49}}{ ($K$-behaviour)}
Let $x\in G$, $a\in A_x$ and $k_1\in K$. Then
\begin{align*}
F(xk_1)&=d_\Lambda\int_K\frac{\pcom_\Lambda(au)\,\pcom_\Lambda(u^*k_1)}{\pcom_{-\mu}(u^{-1}a^{-1}xu)}\,du,
\tag*{{\cite[(16)]{P49}}}\\
F(k_1x)&=d_\Lambda\int_K\frac{\pcom_\Lambda(k_1au)\,\overline{\pcom_\Lambda(u)}}{\pcom_{-\mu}(u^{-1}a^{-1}xu)}
\,du.\tag*{{\cite[(17)]{P49}}}
\end{align*}
Consequently $F(xt)=t^\Lambda F(x)=F(tx)$ for $t\in T$, and if $\varepsilon_a=\varepsilon_b$,
$\lambda_a>\lambda_b$, then $R(E_{ab})F=0$.
\end{quotedthm}

\begin{quotedthm}{\cite[Lemma~3.25]{P49}}{ (Decay)}
For every integer $0\le c\le h$ and $g\in G$, $|F(g)|\le(2/\|g\|)^c$. If $h\ge n^2+1$, then $F\in
L^1(G)\cap L^2(G)\cap L^\infty(G)$.
\end{quotedthm}

\begin{quoteddef}{\cite[Definition~8.1]{P49}}{}
Let $e=e_j$ or $\overline{e_j}$. For a function $F$ on $G_j(\auA)$ put $(R(e)F)(g)=
\int_{G_{j,\infty}}e(x)F(gx)\,dx$ when the integral converges absolutely. A \emph{form of type $e$}
on $G_j$ is a bounded continuous function $y$ on $G_j(\auA)$, left invariant under $G_j(E^+)$, right
invariant under an open subgroup of $G_j(\auA_f)$ (a \emph{level}), with $R(e)y=y$. The
\emph{forms of the archimedean type of the chain} on $G_j$ are the forms of type $e_j$; their
complex conjugates are the forms of type $\overline{e_j}$. The space $\auAz$ of
{\cite[\S14.3]{P46}} is the space of the forms of type $e_0$ and level $J_M$; this is the precise
meaning of ``the fixed vector of the discrete series with parameter $\varsigma_\tau$''. For $j=m$
the forms of type $e_m$ are the functions $g\mapsto l\bigl(g_\infty^{-1}\varphi(g_f)\bigr)$ with
$\varphi$ an algebraic modular form of weight $\lambda^\circ$ and $l$ the linear form given by
$v_m$ (Peter--Weyl); in this way the last lift has values in algebraic modular forms.
\end{quoteddef}

\begin{quotedthm}{\cite[Proposition~8.4]{P49}}{ (The lower chain is well defined)}
Assume {\rm(A1)--(A3)}. Let $0\le j\le m$.
\begin{enumerate}
\item[(a)] Every form $y$ of type $e_j$ or $\overline{e_j}$ on $G_j$ is smooth, $K_{j,\infty}$-finite
and finite under the centre of the enveloping algebra; it is a cusp form; it is rapidly
decreasing on Siegel sets, in the form: for every Siegel set $\mathfrak S_\infty\subset G_{j,\infty}$, finite set $\mathsf c\subset G_j(\auA_f)$, open subgroup $K'$ of $G_j(\auA_f)$ and $N\ge0$ there is $C_N$ with $|y(g)|\le C_N\|g_\infty\|^{-N}$ for $g\in\mathfrak S_\infty\mathsf cK'$; and the forms of a given level form a finite-dimensional space.
\item[(b)] Let $j<m$, $f\in\psHS_f(W)$ and let $y$ be a form of type $e_j$ on $G_j$. Then the
integral $F(h)=\int_{[G_j]}\theta\bigl(\omega_{\mu_{j+1}}(g,h)(\Phi_j\otimes f)\bigr)y(g)\,dg$ converges
absolutely and $F$ is a form of type $e_{j+1}$ on $G_{j+1}$.
\item[(c)] Let $j<m$, $f\in\psHS_f(W)$ and let $z$ be a form of type $\overline{e_{j+1}}$ on
$G_{j+1}$. Then $F^*(g)=\int_{[G_{j+1}]}\theta\bigl(\omega_{\mu_{j+1}}(g,h)(\Phi_j\otimes f)\bigr)z(h)
\,dh$ converges absolutely and $F^*$ is a form of type $\overline{e_j}$ on $G_j$.
\item[(d)] The space $\auAz$ is finite-dimensional and stable under the right translations by
the stabiliser of the lattice at the chosen places, so $\auG$ acts on it; the Poincar\'e map
$\auP$ of {\cite[\S14.3]{P46}} is well defined; and $\iota_\vartheta$ maps $\auAz$ to algebraic
modular forms of weight $\lambda^\circ$ on $\auU(Y_m)$.
\end{enumerate}
In particular the assumption made in {\cite[\S5.1.3]{P49}} holds: the lifts $\Theta^Y_{j,\vartheta}$ and
their transposes $\Theta^{Y,*}_{j,\vartheta}$ preserve the types, and all the forms concerned are
cusp forms and rapidly decreasing.
\end{quotedthm}

\begin{quotedthm}{\cite[Theorem~2.17]{P52}}{ (Theorem~\ref{pbm:inputs}\,\textup{(R)})}
Assume \textup{(H0)}--\textup{(H5)}. There are functions $J_i=J^\tau_i$, $0\le i\le n$, on pairs of
sequences of non-zero complex numbers such that:
\begin{enumerate}
\item[(R-a)] $(J_i)_{0\le i\le n}$ is a sieve system (Definition~\ref{pbm:system}) at the place
$(\C,\text{Lebesgue},c\mapsto e^{c-\bar c},w\mapsto w^{-1})$ with mass $(1+b^{-10})/\pbdeg_\tau$,
weights $\Lambda^\infty_i$ and decay $\beta^\infty$ of \eqref{pbm:eq-beta}:
\[
\Lambda^\infty_i(p)=\prod_{l=1}^i\bigl[|\delta_l|\le b^{60n^2}\bigr]\max(1,|\delta_l|^{-1})^{n^6},\quad
\delta_l=p_0\cdots p_{l-1};\qquad\beta^\infty(\xi)=\pbmTh\,b^{2000n^9}\,\pfrell_{k-n^7}(|\xi|^{-1}),
\]
$\pfrell_K(t)=t^K\min(1,t)^6$;
\item[(R-b)] $J_n(x,x')=\pfrip{\pfrX(x)}{\pfrX(x')}$; and for every $\sigma\in\pfrG^{(\varrho)}$ and all $x,x'$
with $x_\varrho=x'_\varrho=h(\sigma z,z)$, $J_n(x,x')=\pfrip{\pfrN_\sigma(x)}{\pfrN_\sigma(x')}$.
\end{enumerate}
\end{quotedthm}

\begin{quotedthm}{\cite[Lemma~2.19]{P52}}{ (Haar measure)}
The group measure $\pfrM$ of the place datum at $\tau$ is $\lambda$ times the Haar measure $dx$ of
$\pfrG$ for a number $0<\lambda<\infty$ with $\lambda\,\pbvol\auU(k)/\pbvol\auU(r)\le1$; and for every
measurable $\phi\ge0$ on $\pfrG$ invariant under left translation by $\pfrU$,
$\int\phi\,d\theta_{\ge0}=\lambda\int_{\pfrU\backslash\pfrG}\phi$.
\end{quotedthm}

\begin{quotedthm}{\cite[Lemma~2.41]{P52}}{}
$\pfrm_j=(\eta_0*\dots*\eta_{j-1})*\theta_{[0,j)}$.
\end{quotedthm}

\begin{quotedthm}{\cite[Lemma~4.1]{P52}}{}
\textup{(a)} For $\nu\in\pfrU_{(<a)}$, $a<\varrho$, and $n=n(v,it)\in\pfrN_a$:
$G_{a+1}(\nu n)=e^{-\pi(|v|^2+t^2)}\,G_a(\nu)$.
\textup{(b)} $\int G_\varrho\,d\eta_\pfrU=1$.
\textup{(c)} The map $(n_0,\dots,n_{\varrho-1})\mapsto n_0\cdots n_{\varrho-1}$ is a homeomorphism of
$\pfrN_0\times\dots\times\pfrN_{\varrho-1}$ onto $\pfrU$, and $\eta_\pfrU$ is a Haar measure of $\pfrU$.
\end{quotedthm}

\begin{quotedthm}{\cite[Corollary~5.12]{P52}}{}
If $b\ge1000$ and $2^{-\varrho}\le b^{-6}$, then $\pffm(\pffK)\le(1+b^{-6})\,\pbvol(\pffK)\,
\pbvol(\pffKM)$.
\end{quotedthm}

\begin{quotedthm}{\cite[Lemma~5.13]{P52}}{ (The basis (L1))}
For the data, $L_T$ has a basis as in \textup{(L1)}, consisting of integral vectors with integral
dual basis.
\end{quotedthm}

\begin{quotedthm}{\cite[Theorem~5.24]{P52}}{ (The gain for fixed vectors)}
Assume (L1)--(L3) at $T$ and at its subsets stable under conjugation, $b\ge1000$, $4f\le b$,
$k\ge3\cdot10^{10}n^4$, and that $T$ contains the places above $M$. For every set $\mathsf A$ of
chosen places above $p_0$ and every family $\mathsf d$ of frame vectors there is an intermediate
group $\pffZ=\pffZ(\mathsf A,\mathsf d)$ such that
\begin{enumerate}
\item[(a)] $\pffZ$ contains the component at $T$ of every $\kappa\in K^{\mathrm{ad}}$ with
$\mathrm{red}(\kappa)\in\pbFix(\mathsf A,\mathsf d)$;
\item[(b)] $\pffm(\pffZ)\,\pbvol(\pffK)\,b^{6k\#\mathsf A}\le\pffm(\pffK)\,\pbvol(\pffZ)$.
\end{enumerate}
\end{quotedthm}

\begin{quotedthm}{\cite[Theorem~5.28]{P52}}{ (Local system at $T$; Theorem~\ref{pbm:inputs}\,\textup{(F)})}
Assume the hypotheses on the data of Theorem~\ref{pbm:L2}. Let $T$ be a finite set of finite places
of $E$, stable under conjugation, containing the places above $MN$. Let $\mathsf A$ be a set of
chosen places above $p_0$ and $\mathsf d$ a family of frame vectors, and let $\pffZ=\pffZ(\mathsf A,
\mathsf d)$ be the intermediate group of {\cite[Theorem~5.24]{P52}}. Then the forms
{\cite[(12)]{P52}} for $\pffZ$ satisfy:
\begin{enumerate}
\item[(a)] $\pffZ$ contains the component at $T$ of every $\kappa\in K^{\mathrm{ad}}$ with
$\mathrm{red}(\kappa)\in\pbFix(\mathsf A,\mathsf d)$;
\item[(b)] $(J^T_i)_{0\le i\le n}$ is a sieve system at $\mathfrak P_T$
(Definition~\ref{pbm:system}) with mass $(1+b^{-6})\,b^{-6k\#\mathsf A}$, weights $\Lambda^T$ and
decay $\beta^T$.
\end{enumerate}
\end{quotedthm}

\begin{quotedthm}{\cite[Theorem~5.29]{P52}}{ (Siegel's step; Theorem~\ref{pbm:inputs}\,\textup{(S)})}
Assume the hypotheses of {\cite[Theorem~5.28]{P52}}, and in addition that $T$ contains the places at which
an element of the set $\mathsf B$ of \S\ref{pbm:ss-cell} is not a unit. Let $J^\tau_n$ be the top
forms at the real places of {\cite[Theorem~2.17]{P52}}. Let $(c_j,c'_j,s_j,\sigma_j)_{j\in\mathsf I}$ be a
finite family of pairwise distinct quadruples of classes $c_j,c'_j$, elements $s_j\in\pbFix(\mathsf
A,\mathsf d)$ and $\sigma_j\in\pbmGf(E^+)$ such that $t_j=t(\sigma_j)\in\pbmT$, $|\tau(t_j)|>N$
for all $\tau$, and $h_{c_j}^{-1}\sigma_jh_{c'_j}$ stabilises the lattice at every finite place
outside $T$. Then
\[
\pbmrS^{\,2}\sum_{j\in\mathsf I}\Bigl\|\sum_{p\in\pbmT^\varrho}\pbmOrb_{c_j,s_j,c'_j}(w_p\sigma_j)
\Bigr\|^2\le(1+b^{-6})\,\pbmaS\,\Rel\sum_{t\in\{t_j\}}\ \sum_{p,q\in\pbmT^\varrho}\prod_\tau J^\tau_n
\bigl([p,t]_\tau,[q,t]_\tau\bigr)\,J^T_n\bigl((p,t),(q,t)\bigr),
\]
with $J^T_n$ the top form {\cite[(12)]{P52}} for the group $\pffZ(\mathsf A,\mathsf d)$ of
{\cite[Theorem~5.28]{P52}}.
\end{quotedthm}

\begin{quotedthm}{\cite[Proposition~6.35]{P52}}{ ({\cite[Theorem~5.31]{P52}})}
{\cite[Theorem~5.31]{P52}} holds as stated, with the measures and the normalisation of {\cite[\S5.7]{P52}}; moreover
the $z$-integral defining $\delta(\Phi,\alpha)$ converges absolutely.
\end{quotedthm}

\begin{thebibliography}{L-values}
\bibitem[Row D]{P46} D.~A.~Trejo Pizzo, \emph{The sign of the radial intersection form on the idele class group}, preprint, 2026, \url{https://github.com/dtrejopizzo/riemann/tree/main/04-papers/46-row-d-sign-radial-intersection-form}.
\bibitem[Theta]{P49} D.~A.~Trejo Pizzo, \emph{The Short Forms Theorem: theta lifts, complex multiplication and divisibility}, preprint, 2026. % LINK-PLACEHOLDER
\bibitem[Hecke]{P50} D.~A.~Trejo Pizzo, \emph{Elementary divisors of an edge Hecke operator on a Frobenius cycle}, preprint, 2026. % LINK-PLACEHOLDER
\bibitem[Frames]{P52} D.~A.~Trejo Pizzo, \emph{The Spanning Theorem: sieve systems at the real and finite places and Siegel's step}, preprint, 2026. % LINK-PLACEHOLDER
\bibitem[Exchange]{P53} D.~A.~Trejo Pizzo, \emph{The Covolume Theorem: exchange identities and the determinant of a chain of theta lifts}, preprint, 2026. % LINK-PLACEHOLDER
\bibitem[L-values]{P54} D.~A.~Trejo Pizzo, \emph{The Covolume Theorem: one theta lift, Hecke $L$-values and test forms of least valuation}, preprint, 2026. % LINK-PLACEHOLDER
\bibitem[Galois]{P55} D.~A.~Trejo Pizzo, \emph{The Covolume Theorem: regular families of twists, Selmer groups and Ribet's method}, preprint, 2026. % LINK-PLACEHOLDER
\bibitem[Hodge]{P56} D.~A.~Trejo Pizzo, \emph{Explicit cup exponents for local Galois representations beyond the Fontaine--Laffaille range}, preprint, 2026. % LINK-PLACEHOLDER
\bibitem{BorelHC1962} A.~Borel and Harish-Chandra, \emph{Arithmetic subgroups of algebraic groups},
Ann.\ of Math.\ (2) \textbf{75} (1962), 485--535.
\bibitem{Green1955} J.~A.~Green, \emph{The characters of the finite general linear groups},
Trans.\ Amer.\ Math.\ Soc.\ \textbf{80} (1955), 402--447.
\bibitem{Labesse} J.-P.~Labesse, \emph{Changement de base CM et s\'eries discr\`etes}, in: On the
Stabilization of the Trace Formula, Int.\ Press, 2011, 429--470.
\bibitem{HarishChandraAF} Harish-Chandra, \emph{Automorphic Forms on Semisimple Lie Groups}, Lecture
Notes in Mathematics 62, Springer, Berlin, 1968.
\bibitem{FJ} M.~Flensted-Jensen, Discrete series for semisimple symmetric spaces, Ann. of Math. (2) 111 (1980), 253--311.
\bibitem{Knapp} A.~W.~Knapp, Representation Theory of Semisimple Groups: An Overview Based on Examples, Princeton Univ. Press, 1986.
\bibitem{Schmid75} W.~Schmid, Some properties of square-integrable representations of semisimple Lie groups, Ann. of Math. (2) 102 (1975), 535--564.
\bibitem{KnappBeyond} A.~W.~Knapp, Lie Groups Beyond an Introduction, 2nd ed., Birkh\"auser, 2002.
\bibitem{GHY} W.~T.~Gan, J.~P.~Hanke and J.-K.~Yu, \emph{On an exact mass formula of Shimura},
Duke Math.\ J.\ \textbf{107} (2001), 103--133.
\bibitem{Macdonald} I.~G.~Macdonald, \emph{Symmetric functions and Hall polynomials}, 2nd ed.,
Oxford University Press, 1995.
\bibitem{BZ76} I.~N.~Bernstein and A.~V.~Zelevinsky, \emph{Representations of the group
$GL(n,F)$ where $F$ is a non-archimedean local field}, Russian Math.\ Surveys \textbf{31}:3
(1976), 1--68.
\bibitem{BorelFin} A.~Borel, \emph{Some finiteness properties of adele groups over number
fields}, Publ.\ Math.\ Inst.\ Hautes \'Etudes Sci.\ \textbf{16} (1963), 5--30.
\bibitem{FlickerKazhdan} Y.~Z.~Flicker and D.~A.~Kazhdan, \emph{A simple trace formula},
J.\ Analyse Math.\ \textbf{50} (1988), 189--200.
\bibitem{Godement63} R.~Godement, \emph{Domaines fondamentaux des groupes arithm\'etiques},
S\'eminaire Bourbaki, vol.~8 (1962--1964), expos\'e 257, 201--225 (Th\'eor\`eme~7, p.~211);
\url{http://www.numdam.org/item/SB_1962-1964__8__201_0}.
\bibitem{HC66} Harish-Chandra, \emph{Discrete series for semisimple Lie groups II. Explicit
determination of the characters}, Acta Math.\ \textbf{116} (1966), 1--111.
\bibitem{Helgason} S.~Helgason, \emph{Groups and Geometric Analysis}, Academic Press, 1984.
\bibitem{Jacobowitz} R.~Jacobowitz, \emph{Hermitian forms over local fields}, Amer.\ J.\ Math.\
\textbf{84} (1962), 441--465.
\bibitem{Landherr} W.~Landherr, \emph{\"Aquivalenz Hermitescher Formen \"uber einem beliebigen
algebraischen Zahlk\"orper}, Abh.\ Math.\ Sem.\ Univ.\ Hamburg \textbf{11} (1935), 245--248.
\bibitem{MW} C.~M\oe glin and J.-L.~Waldspurger, \emph{Spectral Decomposition and Eisenstein
Series}, Cambridge Tracts in Math.\ \textbf{113}, Cambridge University Press, 1995.
\bibitem{Neukirch} J.~Neukirch, \emph{Algebraic Number Theory}, Grundlehren Math.\ Wiss.\
\textbf{322}, Springer, 1999.
\bibitem{Scharlau} W.~Scharlau, \emph{Quadratic and Hermitian Forms}, Grundlehren Math.\ Wiss.\
\textbf{270}, Springer, 1985.

% refs-B-orbital.tex: new references of proofs/B-orbital.tex.
% None: HC66, Knapp (refs-B-multiplicities), KnappBeyond (refs-A-coefficient), Schmid75
% (refs-A-archimedean) and HornJohnson (refs-B-discrete-series) are already in the bibliography.
\bibitem{HornJohnson} R.~A.~Horn and C.~R.~Johnson, \emph{Matrix analysis}, 2nd ed., Cambridge University
Press, 2013.
\bibitem{WallachCT} N.~R.~Wallach, \emph{On the constant term of a square integrable automorphic
form}, in: Operator algebras and group representations, Vol.~II (Neptun, 1980), Monogr.\ Stud.\
Math.~18, Pitman, 1984, 227--237.

% refs-B-pillars.tex: new references cited by proofs/B-pillars.tex.
% None: the only citation, BombieriVaaler, is already in refs-B-bessel-expansion.tex.

% refs-B-period.tex: new references cited by proofs/B-period.tex.
% None: \cite{Scharlau} (Witt's theorem for hermitian forms) is in refs-B-multiplicities.tex.

% refs-B-gain.tex: new references cited by proofs/B-gain.tex.
% None.
\bibitem{Cassels} J.~W.~S.~Cassels, \emph{An Introduction to the Geometry of Numbers}, Grundlehren
der mathematischen Wissenschaften 99, Springer, Berlin, 1959.
\bibitem{HenkMinima} M.~Henk, \emph{Successive minima and lattice points}, Rend.\ Circ.\ Mat.\ Palermo (2) Suppl.\ \textbf{70} (2002), 377--384; arXiv:math/0204158.

% refs-B-size.tex: new references cited by proofs/B-size.tex.
% None: the Ostrowski--Taussky inequality and Hua's integral are proved in the text.
\bibitem{WallachRRG1} N.~R.~Wallach, \emph{Real Reductive Groups I}, Pure and Applied Mathematics 132, Academic Press, 1988.
\bibitem{HechtSchmid75} H.~Hecht and W.~Schmid, \emph{A proof of Blattner's conjecture}, Invent.\ Math.\ \textbf{31} (1975), 129--154.

% refs-B-frames-real.tex: no new bibitems so far.

% refs-B-frames-shared.tex: new bibitems (cited by author and title in the text of B-frames-shared.tex).

% refs-B-frames-realplace.tex: no new bibitems.
\end{thebibliography}

\end{document}
